\documentclass[10pt]{article}
\usepackage{hyperref}
\usepackage{booktabs}

\usepackage{mathrsfs} 
\usepackage{pifont}% For dingbat symbols like X marks

\newcommand{\cmark}{\ding{51}}% Defines a command for checkmark
\newcommand{\xmark}{\ding{55}}% Defines a command for X mark
\usepackage[a4paper, margin=1in]{geometry}
\usepackage{tikz}
\usetikzlibrary{positioning, arrows.meta, decorations.pathreplacing, calc, backgrounds}
\usepackage{amsmath,amssymb}
\usepackage{amsthm}
\usepackage{pdflscape} 
\usepackage{braket}
\newcommand{\ketbra}[2]{\mathinner{|{#1}\rangle\,\langle{#2}|}}
\usepackage[margin=1in]{geometry}

\usepackage{graphicx}
\usepackage{enumitem}
\setlist[enumerate]{align=left,labelindent=0pt}
\hypersetup{colorlinks=true, linkcolor=blue, urlcolor=blue, citecolor=blue}

\usepackage[a4paper,margin=1in]{geometry}
\usepackage{amsmath,amssymb,amsfonts,mathtools,bm}
\usepackage{physics}
\usepackage{siunitx}
%\usepackage[numbers,sort&compress]{natbib}
\usepackage[backend=biber,style=phys]{biblatex}
\addbibresource{refs.bib}
\usepackage{hyperref}
\hypersetup{colorlinks=true,linkcolor=blue,citecolor=magenta,urlcolor=blue}
\newcommand{\Slope}{S} % universal slope parameter
\newcommand{\HeOp}{\mathcal{D}_f^{{*}}} % adjoint (Heisenberg) dissipator
\newcommand{\Fsloc}{\mathcal F_{\rm s.loc}}
% Chapter-specific theorem environments
\theoremstyle{definition}
\newtheorem{definition}{Definition}[section]
\newtheorem{remark}{Remark}[section]
\theoremstyle{plain}
\newtheorem{theorem}{Theorem}[section]
\newtheorem{axiom}{Axiom}[section]
\newtheorem{lemma}{Lemma}[section]
\newtheorem{proposition}{Proposition}[section]
\newtheorem{statement}{Statement}[section]
\newtheorem{corollary}{Corollary}[section]
\newtheorem{postulate}{Postulate}[section]
\newtheorem{assumption}{Assumption}[section]
\newtheorem{example}{Example}[section]

\title{Universe-Originating Measurement: Operational CCC via AdS/CFT $\longleftrightarrow$ dS/CFT Bridges with LoG CMB imprint, decaying $\Lambda_4$, and deterministic collapse}
\author{Dmytro Surdu}
\date{\today}

\begin{document}
\maketitle

\begin{abstract}
We present an \emph{operational} realization of conformal cyclic cosmology in which the origin of a new aeon is a scheduled, finite--cutoff bridge between AdS$_4$/CFT and dS$_4$/CFT branes of AdS$_5$. The bridge interleaves two non--commuting $\epsilon$--dynamics on the receiver: (i) the OS resolution/energy scale and (ii) the curvature scale $|R|$ (see Sec.~\ref{sec:two-times-semiclassical} for the two--time semiclassical characterization). At every finite tick the receiver remains OS--positive and conformal; when the curvature crosses zero at finite resolution, the new dS$_4$ aeon begins.

On the dS side, we identify a deterministic (POVM-compatible) Lindbladian dephasing (a pointer-aligned, energy-preserving QND channel in the radiative sector, augmented by a Diosi--Penrose-type mass dephasing with guarded, negligible heating) that \emph{creates} classical facts, establishing a well-defined arrow of time in the source brane,  and \emph{exports} recoverable information in zero--DC pulses. On the AdS side, the receiver is a pure vacuum, OS--positive, conformal code built tick-by-tick by pulses with widths \emph{locked} to the brane band edge $\Lambda_k$; the Weyl compensator/Brown--York flux monotonically increases the receiver curvature (a decaying $\left | \Lambda_4 \right |$ on the head) and triggers the $AdS_4\rightarrow dS_4$ flip in a \emph{finite} number of ticks.

A single 5D description controls the construction. The \emph{on–shell Schwinger–Keldysh} (SK) influence action functional $W[J]$ of AdS$_5$ with a brane stack reduces, tick-by-tick, to 4D modular free energies of branes with no additional free knobs.

Geometrically, the bulk realizes an alternating sequence of interior RS1--type cells and Karch--Randall (KR) end--of--the--world (EOW) heads; the \emph{active receiver} is a KR head that continues through the split and sits just UV--side of the newborn source. The two--time semantics (Lorentzian time and RG time) are derived semiclassically from SK+HJ in Sec.~\ref{sec:two-times-semiclassical}.

Framework identifies the source of 5D Universe existence in the \emph{second AdS$_5$ time}: the boundary RG flow along the scale/radial isometry, $u=\ln(\Lambda_{\rm UV}/\Lambda)$. This meta-time is a CPTP semigroup that is monotone in recoverable information (P2), providing a self-caused, irreversible measurement flow that replaces any need for bulk-time nonunitarity.


Phenomenology follows directly. (i) \textbf{LoG CMB imprint:} a degree--scale compensated core with an optional fainter, narrower ring, identified as a Cold Spot anomaly; (ii) \textbf{Shadow--brane--stack dark matter:} a single transfer kernel produces decaying (with the corresponding increase of Hubble constant) to the percent--level uniform ``mirage'' today, a percent-level scale--dependent bias on LSS, and \emph{cored} haloes compatible with clusters constraints and reproducing log-skew unimodal dwarf cores distribution; (iii) \textbf{fifth force-safe} radion phenomenology on the KR head, with any RS1 radion confined to interior/shadow cells; (iv) \textbf{Higgs portal:} the anomaly--LoG DC component freezes a constant Higgs mass shift $\Delta\mu_H^2=-\lambda_p J_0/m_{\mathcal S}^2$, naturally stable for $m_{\mathcal S}\gg v$, with $\delta\kappa_h\ll 10^{-3}$ (quiet colliders) and no late drift of $v$.

\end{abstract}


\section{Preamble}
\label{sec:preamble}

\paragraph{What this paper does.}\label{par:uom-main-001}
We propose and develop an \emph{operational} realization of conformal cyclic cosmology (CCC) \cite{Penrose2010Cycles} in which the origin of a new aeon is a brane hypersurface flip AdS$_4$/CFT $\leftrightarrow$ dS$_4$/CFT \cite{Maldacena1998,Witten1998AdSCFT,GubserKlebanovPolyakov1998,Strominger2001dSCFT}. This transition is \emph{actionable}: a finite–cutoff, analytically continued, and \emph{scheduled} bridge that transform the high-energy pure vacuum AdS spacetime to a symmetric-broken dS aeon, which then 
\begin{enumerate}[label=(\roman*)]
    \item creates classical reality (facts) in the dS$_4$ brane via a deterministic, energy–preserving non–unitary flow;
    \item exports available recoverable classical information to an external \emph{pure} receiver identified with the AdS$_4$ brane of the next aeon in a pulse with a well-defined envelope refining the receiver boundary cutoff, and
    \item iterates these pulses, creating the interleaving $\epsilon$-dynamics of the cutoff energy conjugate $\varepsilon_\Lambda$ and cosmological constant $\Lambda_4$, operationally taking the non-commuting limit in the bridge transition flip, which iteration is the least influence action in AdS$_5$.
\end{enumerate}
The result is a Universe–originating measurement in the CTFM sense, whose successive refinements implement an operational CCC without an infinite future and with explicit falsifiability at both quantum and cosmological scales.


\subsection{How this work relates to CTFM}
\label{subsec:ctfm-relation}
CTFM (Complexity-Theoretic Foundation for Metrology) should \emph{not} be seen as a part of the current physical theory. In the development of this theory, CTFM served as an epistemology-side guide; indeed, the foundational question for the theory had been formed by the virtue of applying the CTFM framework to the quantum state collapse problem.

Capabilities of CTFM go beyond the foundational question, however. They allow checking possible answers to problems for the epistemological value and validity, and suggest the direction for sourcing the correct answers.

We adopt the CTFM stance that \emph{classical reality} is a stable collection of \emph{facts}, i.e.\ tails of events (indications, or computations) that admit tractable certification. Such a certification (measurement test) prove that the considered tail indications (a \emph{state} belonging to the compact space - axiom $(\mathrm{C})$) enters the \emph{invariant attractor basin} (Fact axiom \((\mathrm F)\equiv(R,L,C)\) (reachability, Lyapunov border, local contraction) with the tolerance verifiable with the Certification Standard (axiom \((\mathrm S)\)).

The key CTFM statement is that any measurement test is equivalent to the Turing-complete finite decomposition into the program composed from primitive measurement tests (sequential, conditional, and loop-until) given the traceability \((\mathrm T)\), time–irreversibility \((\mathrm I)\), and effective measurability \((\mathrm{EM1\text{–}3})\) axioms (Fundamental Test Equivalence Theorem, CTFM Book III, Chapter 3 \cite{CTFMBookIII}). Such a decomposition is impossible (measurement fact is decomposition-independent) if and only if the problem of compiling and running such a program is undecidable (Fact Decomposition-Independence Theorem, CTFM Book III, Chapter 4 \cite{CTFMBookIII}). 

With this semantics, the existence of any classical world-line is already a nontrivial \emph{metrology} statement: a language of ``good heads'' (prefixes) that must be tractably certifiable at finite guard.

\subsection{Why classical reality is the central problem}
\label{subsec:intro-why-classical}

CTFM shows that when a measurement channel is \emph{lossless} (invertible), tail–fact certification becomes $\Pi^0_2$–hard (CTFM Book I, Chapter 5 \cite{CTFMBookI}); hence ``pure'' unitarity cannot, by itself, yield facts.
This leaves physics with a sharpened question:

\begin{quote}
\emph{If global unitary evolution cannot produce facts, how did classical reality originate?}
\end{quote}

CTFM answer is constructive and operational: a \emph{fact} of the existence of classical reality must come in the form of a concrete tail-invariant basin; the initial state must evolve under the known update map so it pointwise reaches this basin under any sequence of inputs. Then this map is decomposable into the measurement program.

If we adopt an AdS/CFT as the state space and the emission "enter dS with a certain initial state" as a measurement fact, the problem is that the initial state is not known. CTFM provides the relief: at any given resolution, this initial state can be resolved using Adaptive Traceability Resolution (CTFM, Book III, Chapter 3 \cite{CTFMBookIII}); in essence, it compiles the measurement program given any initial state space. Thus, something must drive the resolution setting and compilation.

However, we do not have any facts that this drive could be built from. What could it be?

The answer this theory provides: in is the \emph{measurement test dual object, which we name here co-test}. The dual space of the co-test is the causal dS space with the arrow of time (hence the time irreversibility axiom (I) is satisfied automatically). The nature of the co-test is the information absorption; the certification is \emph{the amount of information left unabsorbed}, while the emission (output) is certificationally unconstrained.

This co-test absorption map is a deterministic, energy–preserving Lindblad flow that \emph{reduces quantum/classical excess} inside the dS aeon without heating,
and its emission map is external (to dS), reversible reconciliation/encoding that compiles the produced facts into a \emph{pure} AdS boundary state for the next aeon.

In CTFM terminology, the essence of the aeon is the \emph{repeated compilation of the measurement program} that, once finally ready, runs to start the next aeon.

\subsection{The theory problem map: what must be addressed}
\label{subsec:intro-map}

We organize the known obstacles that any origin scenario must solve, and that our construction addresses, into six interacting clusters.

\begin{enumerate}
  \item \textbf{Measurement and classicality.} How can facts emerge from dynamics? In CTFM, finite–guard certification of tail predicates requires \((\mathrm F),(\mathrm S),(\mathrm C),(\mathrm T),(\mathrm I),(\mathrm{EM})\). Purely unitary, invertible channels fail this (they are lossless and force $\Pi^0_2$–hardness). 

  \item \textbf{Initial conditions and arrows.} The observed universe exhibits extremely low gravitational entropy initially (Weyl curvature hypothesis) and a consistent macroscopic arrow of time. An origin mechanism must deliver regular initial data and a ranking/arrow consistent with \((\mathrm I)\).

  \item \textbf{AdS/CFT vs.\ dS/CFT tension.} AdS/CFT is unitary and elegant but negative–curvature; our universe is effectively de Sitter at early and late times. Lorentzian dS/CFT is non–unitary (as a boundary theory), opening the door to fact production but raising observer–dependence issues and reflection–positivity (RP) subtleties.

  \item \textbf{Non–commuting limits at the bridge.} Analytic continuation with a strict removal of the UV cutoff does not commute with the dS limit; operationally, the AdS$\to$dS crossing must occur at \emph{finite} cutoff. One needs a \emph{scheduler} that interleaves small but nonzero cutoff with the cosmological parameter crossing.

  \item \textbf{Observer–dependence and reconciliation.} In a dS phase, different causal patches can produce disagreeing classical records. An origin mechanism must \emph{glue} these into a single global state without introducing heating or new physical degrees of freedom, and it must do so \emph{reversibly} to preserve the purity of the absorber.

  \item \textbf{Larger landscape.} Where does all this happen? Are there level-up governing principles and laws?

\end{enumerate}

\subsection{Our operational stance and three postulates}
\label{subsec:intro-postulates}

We make three explicit postulates that render CCC operational. Here, we provide them in simplified and shortened form taken for one aeon:

\begin{description}
  \item[P1 (Unitarity cannot produce facts).] A purely unitary, invertible channel cannot produce certifiable tail facts at finite guard. Operationally, the bridge state (AdS side) is \emph{pure} and \emph{factless}.

  \item[P2 (Quantum Second Law).] The \emph{total recoverable information} (classical + quantum that is \emph{in principle} decodable) is non–decreasing within and across aeons. In particular, certified facts are never lost; compilation into the absorber preserves recoverable information.

  \item[P3 (Quantum/classical excess minimization).] On any spacetime region, the dynamics tends to \emph{minimize} the time integral of the excess (quantum information that is not yet compiled into classical facts). 
\end{description}

P1 makes the purity of the AdS receiver necessary; P2 forbids throwing away too much classical facts and necessitates the information export from the information-destroying dS; P3 provides the unique saddle - the best admissible fact–producer inside dS (deterministic, energy–preserving non–unitary flow) combined with the best admissible package (OS-positive, conformal, maximally symmetric filling with the minimal Weyl/BY anomaly floor) of information in the AdS$_4$ receiver.


\section{Motivation and idea}\label{sec:uom-main-001}
The idea behind this work is to reconstruct Conformal Cyclic Cosmology theory \cite{Penrose2010Cycles} in a fully operational implementation. The theory inherits advantages from the Penrose's CCC \cite{Penrose2010Cycles}: it solves the horizon and flatness problems without requiring extreme conditions of high $e$-fold inflation \emph{and} answers the problem of the Universe's initial entropy conditions in a natural way.

In addition, identifying the Universe's origination event as a transformation from AdS/CFT to dS/CFT brings many benefits:
\begin{enumerate}
    \item Completed with the postulate P1 (unitarity can't produce facts), it provides the explanation of what means the Universe even originated: it is the appearance of classical reality due to the shift from pure unitarity of AdS/CFT to \emph{almost perfect} unitarity of dS/CFT.

    This, in turn, explains: 
    \begin{itemize}
        \item \textit{for the proposition of our early Universe to be a less firm dS/CFT} - the regularity of its initial conditions: they are fixed by the analytical continuation of AdS/CFT preceeding it;
        \item \textit{the Weyl curvature hypothesis} - AdS/CFT is pure (and hence factless, P1) iff its Weyl curvature is zero.
    \end{itemize}
    \item Such a pure AdS/CFT spacetime has to be created; the operational solution to its creation from the dS/CFT aeon answers the arrow of time problem - our aeon is effectively the open system, breaking the constant von Neumann entropy wall. 
    
    Simultaneously, the requirement for the absorber to be POVM-compatible QND provides the fundamental answer to the measurement problem: the collapse in dS is deterministic. If taken over the Lorentzian time slice, it (the classical history of all aeons) is effectively the saddle point of the 5D least action functional. Notably, this collapse model recovers Diosi-Penrose scaling law, but with almost exact no-heating.
    \item This consideration provides a resolution for \textit{two main string theory problems}:
    \begin{itemize}
        \item String theories describe the Weyl-symmetry-based unification of gravity and quantum theory; however, their most profound result - the AdS/CFT correspondence - presumably describes a different Universe than the one we live in.
        \item Practically intractable amount of solutions to the theory, that wildly exceeds the number of degrees of freedom in our Universe, makes it uninstantiable absent the superimposed selection process, alike epistemically-limiting anthropic principle over multiverse. 
    \end{itemize}
    However, if AdS$_4$/CFT is not \emph{our observable} Universe, but its factless predecessor, both problems effectively vanish. Our theory is CFT-agnostic: it happily allows SST (e.g. IIb) on the AdS$_5$ boundary, making one a plausible and well-understood candidate for ToE.
    \item The main predictions $\Lambda CDM$ cosmology is in good agreement with the AdS bridge theory, which predicts the current CMB observations and strictly constant $\Lambda$. 

    The remarkable benefit of the AdS bridge for $\Lambda CDM$ is the \emph{explanation of cosmological constant} $\Lambda_4$ \emph{fine-tuning} through the limit-taking drive of $\Lambda_4 \uparrow 0^{-}$ interleaving with $\epsilon_\Lambda \rightarrow 0$  cutoff dynamics, with $\Lambda_4$ crosses zero exactly at the Universe origination.
    
    Another benefit for $\Lambda CDM$ is the natural DM explanation: it is ordinary matter with a probable large contribution of SMBH \emph{on the brane of the previous aeon}. Unlike the original CFT, our theory doesn't prescribe the full transformation of the previous aeon: once the new aeon brane is created by the flip, the old one continues to evolve by itself, being coupled to the other branes only gravitationally. 
    
    The dynamics of the proper brane separation due to the dS expansion fits to the difference in density parameters of early Universe and the corresponding Hubble tension.

    Taking into account the \emph{stack} of the shadow branes from the previous aeon, the theory naturally explains DM high-fidelity picture present in the dwarf cores distribution.
    
    \item The CCC itself wins closure of its main gap: the gauge freedom of post-matching energy density. In our paper, we give the operational design:  the concrete causal mechanism of the conformal transformation between aeons that fixes the next aeon's energy level, naturally explaining leptobariogenesis. 
    
    Another benefit is the finite time to aeon flip, bounded to much shorter spans than CCC's minimum of $10^{100}-10^{110}$ years before pure radiative dS conditions (which minimum, in a sense, causes stronger extrapolation concerns than the inflation). 
    
    In addition, specific predictions of the theory with an AdS bridge differ significantly from the original CCC and arguably are in better agreement with the current cosmological observations (concrete, high-fidelity dark matter model, CMB LoG/SMHW feature, $\Lambda CDM$ compatibility).
    \item Finally, AdS$_4 \rightarrow$ dS$_4$ bridge naturally generates the solution for the hierarchy problem; importantly, the Higgs field is generated with a brane boundary fluctuation portal from the LoG DC potential shape imprinted by AdS$_4$ uplifts. 
    
    The same imprint is present in the non-fixed trace of the AdS brane at the flip; decaying, it creates the cosmologically observable LoG/Spherical Mexican Hat Wavelet pattern in the CMB that we identify as the Cold Spot.
\end{enumerate}

\subsection{\texorpdfstring{$AdS\;\rightarrow\;dS$ flip idea}{AdS->dS flip idea}}\label{subsec:uom-main-001}

\subsubsection*{Assumptions}\label{subsubsec:uom-main-001}

\begin{assumption}[Conformal bridge]
\label{ass:bridge}
Our origin is described by a conformal identification of the previous aeon's future infinity with a regulated Euclidean AdS (EAdS) saddle continued to Lorentzian de Sitter (dS). Conformal structure (null cones) is continuous across the bridge.
\end{assumption}

\begin{assumption}[$\epsilon$--collapse (radial cutoff)]
\label{ass:epsilon}
The boundary theory is rendered effectively compact by introducing a radial cutoff in EAdS at $z=z_0 \equiv \epsilon>0$ (hard wall or ETW surface). On compact spatial slices, this yields a finite number of boundary modes below a UV scale $\Lambda\sim 1/z_0$.
\end{assumption}

\begin{assumption}[Concerted transition]
\label{ass:concerted}
The flip of the AdS$_4$ brane cosmological constant is implemented by the standard analytic continuation $L\to iL$ (with appropriate Wick rotations), and is performed while the cutoff of Assumption~\ref{ass:epsilon} is \emph{kept finite}. Optionally, $z_0$ may evolve slowly in time during/after the crossing.
\end{assumption}

\begin{assumption}[Classicality criterion]
\label{ass:classicality}
Assuming early dS Universe and its AdS predecessor, classicality (WKB branching) requires an effectively compact/code subspace so that the analytical continuation between these spaces is well-defined. In the unrestricted continuum CFT this analytic continuation doesn't exist; with $\epsilon$--collapse in place, it arises.
\end{assumption}

\begin{assumption}[Pre-bridge acceleration]
\label{ass:prevacc}
In the previous aeon, there exists a late accelerating (quasi-dS) era with $N_{\rm prev}$ e-folds immediately before the bridge.
\end{assumption}

\noindent\emph{Scope.} The statements and proofs in this section are global (on $S^d$) and are used to
characterize the RP behaviour of the analytically continued covariances at finite cutoff. They are not
used for the static–patch (source) thermodynamics, which is treated separately.


\subsubsection*{Regulated EAdS and its boundary kernel}\label{subsubsec:uom-main-002}

Consider EAdS$_{d+1}$ in Poincar\'e coordinates
\begin{equation}
ds^2 = \frac{L^2}{z^2}\left(dz^2 + d\vec{x}^{\,2}\right),\qquad z\in[z_0,\infty)\,.
\end{equation}
For a scalar $\phi$ with mass $m$, Fourier transform along $\vec{x}$ gives the mode equation with solution
\begin{equation}
\phi_{\bm{k}}(z) = \phi_{\bm{k}}(z_0)\,\frac{z^{d/2}K_\nu(k z)}{z_0^{d/2}K_\nu(k z_0)}\,,\qquad
\nu \equiv \sqrt{\frac{d^2}{4}+m^2L^2},\ \ k=\abs{\bm{k}}\,,
\end{equation}
where $K_\nu$ is a modified Bessel function. The on-shell action reduces to a quadratic boundary functional at $z=z_0$,
\begin{equation}
S_{\rm os}[\phi(z_0)] = \frac{1}{2}\int\!\frac{d^dk}{(2\pi)^d}\,
\phi_{-\bm{k}}(z_0)\, \mathcal{F}(k,z_0)\, \phi_{\bm{k}}(z_0),
\end{equation}
with canonical momentum $\Pi=L^{d-1} z^{1-d}\partial_z \phi$ and kernel
\begin{equation}
\boxed{\quad \mathcal{F}(k,z_0) = L^{d-1} z_0^{-d}\left[-\nu + (kz_0)\,\frac{K_{\nu-1}(kz_0)}{K_\nu(kz_0)}\right].\quad}
\label{eq:kernel}
\end{equation}

\begin{proposition}[Radial cutoff $\Leftrightarrow$ UV cutoff]
\label{prop:cutoff}
For $kz_0\ll 1$, $\mathcal{F}(k,z_0)$ reproduces the conformal scaling of the CFT two-point function (up to local counterterms); for $kz_0\gg 1$ it acts as a smooth form factor suppressing high-momentum modes. In particular,
\begin{align}
\mathcal{F}(k,z_0) &= \text{local}(z_0^{-d}) + c_\nu\, L^{d-1} k^{2\nu} + \mathcal{O}((kz_0)^2), \qquad kz_0\ll 1,\\
\mathcal{F}(k,z_0) &\sim L^{d-1} z_0^{1-d}\, k, \qquad\qquad\qquad\qquad\quad\ \ \ \ kz_0\gg 1,
\end{align}
so that nonlocal correlations are scale-invariant for $k\ll \Lambda\equiv z_0^{-1}$ and smoothly damped for $k\gtrsim \Lambda$.
\end{proposition}

\begin{proof}
Use the small- and large-argument asymptotics of $K_\nu$ and the identity $K_\nu'(x)=-K_{\nu-1}(x)-\nu K_\nu(x)/x$. Local polynomial divergences in $z_0^{-d}$ are removable by standard holographic counterterms and do not affect nonlocal momentum dependence; the nonlocal piece is $k^{2\nu}=k^{2\Delta-d}$ with $\Delta=\frac{d}{2}+\nu$. For $kz_0\gg 1$, $K_{\nu-1}/K_\nu\to 1$ yields the linear scaling in $k$ and hence suppression of the regulated propagator $\mathcal{F}^{-1}$ at large $k$.
\end{proof}

\subsubsection*{Analytic continuation to dS and classicality onset}
\label{subsec:analytic-continuation}

Analytic continuation $z\to i\eta$, $L\to iL$ (with $\eta<0$ conformal time) gives the dS metric
\begin{equation}
ds^2 = \frac{L^2}{\eta^2}\left(-d\eta^2 + d\vec{x}^{\,2}\right),\qquad \eta\in(-\infty,0).
\end{equation}
The EAdS partition function with cutoff maps to a regulated Hartle--Hawking wavefunctional \cite{HartleHawking1983} on $S^d$,
\begin{equation}
\Psi_{\rm HH}[\phi_{(0)}]\ \propto\ \exp\!\left(-S_{\rm os}^{\rm (EAdS)}[\phi_{(0)}]\right)
\ \ \longrightarrow\ \
\exp\!\left(i S_{\rm eff}^{\rm (dS)}[\phi_{(0)}]\right),
\end{equation}
where the index $\nu$ continues to $\nu_{\rm dS}=\sqrt{\frac{d^2}{4}-m^2L^2}$ so that for $m^2L^2>\frac{d^2}{4}$ one gets principal-series $\Delta=\frac{d}{2}\pm i\mu$, $\mu=\sqrt{m^2L^2-\frac{d^2}{4}}$.

\begin{theorem}[Smooth degradation of reflection positivity at fixed cutoff]
\label{thm:smoothRP}
Let $C_\Lambda(\mu_{\rm ps})$ denote the regulated covariance (inverse of $\mathcal{F}$ after continuation), with UV cutoff $\Lambda=z_0^{-1}$. For any fixed $\Lambda<\infty$, eigenvalues of the reflection-positivity quadratic form depend continuously on $\mu$, and the smallest violating eigenvalue scales as
\begin{equation}
\lambda_{\min}(\mu_{\rm ps},\Lambda) = -\,\alpha(\Lambda)\,\mu_{\rm ps}^2 + \mathcal{O}(\mu_{\rm ps}^4),
\end{equation}
with $\alpha(\Lambda)>0$. Hence violations of unitarity proxies are parametrically small for $\mu_{ps}\to 0$ at fixed cutoff.
\end{theorem}

\begin{proof}[Proof sketch]
At fixed $\Lambda$, project the spatial Laplacian $-\Delta_{S^d}$ to any finite-rank subspace; the RP quadratic form
becomes a finite Hermitian matrix depending analytically on $\mu_{\rm ps}$ (by analytic continuation in the Bessel/Hankel
index). Eigenvalues are even analytic functions of $\mu_{\rm ps}$, hence $\lambda_{\min}(\mu_{\rm ps},\Lambda)
= -\alpha(\Lambda)\mu_{\rm ps}^2+O(\mu_{\rm ps}^4)$ with $\alpha(\Lambda)>0$ determined by the second–order perturbation.

The sign follows from the appearance of oscillatory phases in the continued kernel. (Details are standard matrix-perturbation theory; see also Prop.~\ref{prop:cutoff} for the underlying kernels.)
\end{proof}

\begin{theorem}[Non-commuting limits]
\label{thm:limits}
The limits $\mu\to 0$ (AdS side) and $\Lambda\to\infty$ (removing the cutoff) do not commute:
\begin{equation}
\lim_{\Lambda\to\infty}\,\lim_{\mu_{\rm ps}\to 0^+} C_\Lambda(\mu_{\rm ps}) \ \text{is reflection positive, but}\ 
\lim_{\mu_{\rm ps}\to 0^+}\,\lim_{\Lambda\to\infty} C_\Lambda(\mu_{\rm ps})\ \text{is not (for any }\mu_{\rm ps}>0).
\end{equation}
\end{theorem}

\begin{proof}[Proof sketch]
At fixed $\mu_{\rm ps}>0$, the coefficient $\alpha(\Lambda)$ diverges with $\Lambda$ (high–frequency modes dominate
the even-in-$\mu_{\rm ps}$ perturbation; in fact $\alpha(\Lambda)\sim \kappa_d\,\Lambda^{d+2}$), so
$\lim_{\Lambda\to\infty}\lambda_{\min}(\mu_{\rm ps},\Lambda)<0$.

(Indeed, one finds $\alpha(\Lambda)\sim \kappa_d\,\Lambda^{d+2}$ as $\Lambda\to\infty$; see App.~\ref{app:OS-discrete}, Thm.~\ref{thm:smooth-RP}.)

\end{proof}

\begin{remark}[Classicality onset]
\label{rem:uom-main-001}
With $\epsilon$--collapse in place, decoherence of super-Hubble modes proceeds efficiently and yields a WKB classical branch. Assumption~\ref{ass:concerted} ensures this happens \emph{at} the AdS$\to$dS crossing; in the unrestricted continuum, the classicality criterion fails.
\end{remark}

\subsubsection*{Curvature across aeons and the flatness condition}\label{subsubsec:uom-main-003}

Let $\Omega_k \equiv -k/(aH)^2$, with $H^2\propto a^{-3(1+w)}$ for constant equation of state $w$. Then
\begin{equation}
\Omega_k(a)\ \propto\ a^{\,1+3w}\,.
\end{equation}
Thus $\abs{\Omega_k}$ grows in decelerating eras (radiation $w=1/3$, matter $w=0$) and decays in accelerating eras ($w<-1/3$; de Sitter: $\propto e^{-2N}$).

\begin{proposition}[Flatness via pre-bridge acceleration]
\label{prop:flatness}
Let $N_{\rm prev}$ be the number of e-folds of acceleration in the previous aeon immediately before the bridge. Let $a_{\rm br}$ denote the scale factor at the onset of radiation domination in our aeon, $a_{\rm eq}$ matter--radiation equality, and $a_\Lambda$ the onset of our late-time acceleration. Then
\begin{equation}
\abs{\Omega_{k,0}}\ \simeq\ e^{-2N_{\rm prev}}\,
\left(\frac{a_{\rm eq}}{a_{\rm br}}\right)^{\!2}\,
\left(\frac{a_\Lambda}{a_{\rm eq}}\right)\,\times \mathcal{O}(1)\,.
\end{equation}
To achieve $\abs{\Omega_{k,0}}<10^{-3}$ without early acceleration in our aeon, it suffices that
\begin{equation}
N_{\rm prev}\ \gtrsim\ \ln\!\Big(\frac{a_{\rm eq}}{a_{\rm br}}\Big)\ +\ \frac{1}{2}\ln\!\Big(\frac{a_\Lambda}{a_{\rm eq}}\Big)\ +\ \frac{1}{2}\ln 10^3\,.
\end{equation}
\end{proposition}

\begin{proof}
Propagate $\Omega_k$ multiplicatively across eras using the $a^{1+3w}$ scaling and $e^{-2N}$ in the accelerating phases. The late-time $\mathcal{O}(1)$ factor accounts for the modest dilution during our current $\Lambda$ era.
\end{proof}
In other words, the flatness mechanism is inherited from Penrose's CCC \cite{Penrose2010Cycles}.

\begin{remark}[Horizon problem]
\label{rem:uom-main-002}
A conformal bridge preserves null cones; with an arbitrarily long pre-bridge epoch, all points on our last-scattering surface share causal contact through the bridge. Hence the horizon problem is addressed without requiring a long inflationary era in our aeon.
\end{remark}

\subsection{The mechanism in one paragraph}
\label{subsec:intro-mechanism}
In the previous subsection we described the mathematics of the AdS $\to$ dS bridge and the problem of the Universe origination through it.

It can be formulated as a problem of taking non-commuting limits operationally, in an interleaved fashion, transversing in this way the Universe through the analytic continuation.

Then the foundational problem arises: how this operational double-limit is taken in AdS, which is, by definition, factless, and hence devoid of the arrow of Lorentzian time and a result to achieve or state?

This is how CCC enters this picture. The answer is in viewing the Universe-Originating Measurement Test as a dual to the fully causal procedure of the previous aeon. This procedure is named Anti-Zeno Co-Test, and is the information sink (as opposed to the test, which is the information source). We will see that in dS, this procedure concrete design is strictly necessary under the proposed postulates P1-P3.


\paragraph{Mechanism, briefly.} We execute the AdS$\to$dS crossing at a small but finite UV cutoff \(\epsilon_{\rm UV} = \Lambda^{-1}\) (``$\epsilon$–collapse''). On the dS side, the only reducer compatible with no heating and \((\mathrm S),(\mathrm T),(\mathrm I)\) is a deterministic, almost-energy–preserving pure-dephasing flow with a tiny, universal rate tied to cosmological scales characterized by the DP self-energy $\Delta E_G/\hbar$; this \emph{produces} classical facts without exchanging energy.\label{par:uom-main-002}
A global, reversible ``exchanger'' (that can be thought as a QC, which state is canonical commutant representation of the degrees of freedom beyond the horizon of every static patch of dS/CFT) performs QND taps of device pointers (commuting with $H$), reconciles patch data into a single band–limited global field under constraints (sheaf gluing), and encodes the result as a \emph{pure} boundary state in AdS/CFT for the next aeon uplifting its theory with new degrees of freedom while preserving all previously stored information and thus causing the quantum-gravitational backreaction (Weyl compensator) setting the uplift spectrum and the exchanger scheduling parameters.
The two $\epsilon$--dynamics---UV cutoff and certification guard---are coupled and governed by the exchanger's optimal--control scheduler that minimizes the backlog of uncompiled (excess) information (see Sec.~\ref{sec:two-times-semiclassical}).
This loop iterates until AdS payload saturate RP/OS deficit guard, compiling a Universe-Originating Measurement Test and closing the CCC chain with finite dS lifetime, no new knobs, and DP scaling as an \emph{output} of the variational postulate (P3) - the consequence of the exchanger's thermodynamics.

\section{Theory Core}\label{sec:uom-main-002}

\subsection{Setup, state spaces, and symbol hygiene}
\label{subsec:scene-notation}

\paragraph{Geometries and groups.}\label{par:uom-main-003}
We take the source to be a (generally mixed) early--to--late-time RS1 \cite{RandallSundrum1999RS1,RandallSundrum1999RS2}
\emph{dS$_4$ brane} embedded in AdS$_5$
(from which we extract only the radiative sector at $\mathscr I^+$, preserving conformality), and the
receiver to be a \emph{KR} vacuum--only AdS$_4$ brane (code space) embedded in the same AdS$_5$ \cite{KarchRandall2001}.
RS1 appears only as interior/shadow brane segments in the stack; KR heads are vacuum EOW caps and do not
carry dark--sector matter:
\[
\big(\mathrm{AdS}_5\ \supset\ \mathrm{dS}_4\text{ brane}\big)\quad\text{(source)}
\qquad\longrightarrow\qquad
\big(\mathrm{AdS}_5\ \supset\ \mathrm{AdS}_4\text{ brane}\big)\quad\text{(receiver)}.
\]
The late-time boundary symmetry $\mathrm{Conf}(S^3)\cong SO(4,1)$ acts unitarily as
$U_{\rm dS}(g)$ on the source radiative algebra and as $U_\partial(g)$ on the receiver boundary algebra.

Spatial slices of $\mathrm{dS}_4$ are $S^3$; we write the conformal group
$\mathrm{Conf}(S^3)\cong SO(4,1)$ and its unitary representations
\[
U_{\rm dS}(g)\ \text{ on } \mathcal A_{\rm rad}(\mathscr I^+),\qquad
U_\partial(g)\ \text{ on } \mathcal A_\partial,\qquad g\in \mathrm{Conf}(S^3).
\]

\paragraph{Algebras, states, and anchors.}\label{par:uom-main-004}
\begin{itemize}[leftmargin=1.5em, itemsep=1pt]
\item $\mathcal A_{\rm sp}$: von Neumann algebra of the $\mathrm{dS}_4$ static patch (source).
\item $\omega_{\rm BD}$: Bunch–Davies KMS state on $\mathcal A_{\rm sp}$ at inverse temperature
$\beta_{\rm BD}=2\pi/H$ (with Hubble $H$). Band–limited restriction is $\rho^{\rm BD}_\Delta$ (see below).
\item $\mathcal A_{\rm rad}(\mathscr I^+)\subset \mathcal A_{\rm sp}$: radiative (effectively massless) late–time boundary algebra at $\mathscr I^+$.
\item $\mathcal A_\partial$: AdS$_4$ brane boundary algebra (receiver/code space).
\end{itemize}

\paragraph{Source algebra and anchor (static patch).}\label{par:uom-main-005}
The source side is the \emph{static–patch} net $(\mathcal A_{\rm sp},\alpha_t)$, with static time $t$ generated by
the self–adjoint $H_{\rm sp}$ on the GNS Hilbert space. The restriction of the Bunch–Davies state to
$\mathcal A_{\rm sp}$ is a \emph{KMS} state at inverse temperature $\beta_{\rm BD}=2\pi/H$.
All source–side dissipative statements (QND/GKLS, Spohn identities, monotonicity of the modular free energy)
are formulated on $(\mathcal A_{\rm sp},\alpha_t)$ with pointers of the form
\[
L_a \;=\; f_a\!\big(H_{\rm sp}\big),
\qquad [L_a,H_{\rm sp}] = 0.
\]

\paragraph{Receiver algebra and export bandedge (retarded acceptance).}\label{par:uom-main-006}
The receiver is the AdS$_4$ head, with boundary algebra $\mathcal A_\partial$ and cutoff $\Lambda$
(the \emph{code} space). The only global spectral input we use is the \emph{export bandedge}
$\Delta_{\rm ret}$, which is fixed \emph{on the receiver} by the knee of the \emph{retarded} cross–brane
transfer into the cutoff band (Def.~\ref{def:retarded-bandedge} below). This bandedge \emph{does not}
define a spectral projector on $\mathcal A_{\rm sp}$; it only enters the welded map into the receiver code.

\paragraph{Source band–limit in static variables.}\label{par:uom-main-007}
We take the source band–limit as a projector in \emph{static–patch} spectral variables,
\begin{equation}
\label{eq:sp-band-projector}
P^{\rm(sp)}_\Delta \;=\;
\mathbf 1_{\{|\omega|\le \Delta_\omega\}}\;\mathbf 1_{\{\,l\le l_{\max}\}}\;\mathbf 1_{\{n\le n_{\max}\}},
\end{equation}
where $(\omega,l,n)$ label, respectively, the static frequency, the $S^2$ angular momentum,
and a radial mode index with horizon boundary conditions. No $S^3$ harmonic decomposition is used on the source.

\paragraph{Normalizations and conventions.}\label{par:uom-main-008}
\begin{itemize}[leftmargin=1.5em, itemsep=2pt]
\item Metric/sign: mostly plus, $G_{\mu\nu}=R_{\mu\nu}-\tfrac12 R g_{\mu\nu}$.
\item Curvatures: on AdS$_4$ background $R=4\Lambda_4=-12/L_4^2$ (so $\Lambda_4<0$).
\item Static–patch modes: choose the equal–time covariance normalization
\[
C_\ell(0)\;=\;\frac{1}{2\,\omega_\ell},\qquad \omega_\ell:=\sqrt{H^2\,\ell(\ell+2)+m^2}\quad (S^3),
\]
so that the flat limit matches harmonic–oscillator conventions (App.~\ref{app:OS-discrete}).
\item Brown–York/holographic data: AdS$_5$ radius $\ell$ and Newton constants $(G_5,G_4)$ with
\[
G_4=\frac{G_5}{\ell}\quad\text{(RS2 scaling)},\qquad
a=c=\frac{\pi\,\ell^3}{8\,G_5}\ \ \text{(holographic CFT$_4$ central charges)}.
\]
These fix the anomaly coefficient $\gamma=\dfrac{a\,G_4}{12\pi}$ used throughout the work.
\end{itemize}



\paragraph{Ticks, pulses, and widths.}\label{par:uom-main-009}
Time is organized in \emph{ticks} $k=0,1,2,\dots$. Each tick carries:
\begin{itemize}
  \item \textbf{Macro (source, KMS–QND absorption):} a bath–driven, zero–area \emph{DoG} envelope acting on the dS$_4$ static–patch pointer,
  \[
  \chi_{t,\Sigma_k}(t)\;=\;g'_{\Sigma_k}(t-t_k)\;=\;-\frac{t-t_k}{\Sigma_k^{2}}\,e^{-\frac{(t-t_k)^2}{2\Sigma_k^2}},
  \]
  with $\int_{\mathbb R}\chi_{t,\Sigma_k}(t)\,dt=0$.
  \item \textbf{Micro (receiver, uplift injection):} a compensator–driven, zero–area \emph{DoG} envelope on the AdS$_4$ brane boundary,
  \[
  \psi_{t,\sigma_{t,k}}(t)\;=\;g'_{\sigma_{t,k}}(t-t_k)\;=\;-\frac{t-t_k}{\sigma_{t,k}^{2}}\,e^{-\frac{(t-t_k)^2}{2\sigma_{t,k}^2}},\qquad \int_{\mathbb R}\psi_{t,\sigma_{t,k}}=0,
  \]
\end{itemize}

\paragraph{Temporal vs spatial welding (LoG structure).}\label{par:uom-main-010}
In time,
\[
\big(g'_{\Sigma_k}*g'_{\sigma_{t,k}}\big)(t)\;=\;\frac{d^2}{dt^2}g_{\sqrt{\Sigma_k^2+\sigma_{t,k}^2}}(t)\quad\text{(LoG)}, 
\]
since DoG${*}$/DoG $\Rightarrow$ second derivative of a Gaussian at the combined scale.
On $S^3$, we take the \emph{micro} spatial envelope to be zonal Gaussian $g_{\sigma_{x,k}}(\Omega)$ and the \emph{macro} spatial envelope to be \emph{LoG}:
\[
\Phi_{x,\rho_k}(\Omega)\;:=\;\Delta_{S^3}\,g_{\rho_k}(\Omega).
\]
Because spherical convolution with zonal kernels commutes with $\Delta_{S^3}$, the welded spatial envelope is again a LoG,
\[
\Phi_{x,\rho_k}\ \circledast\ g_{\sigma_{x,k}}
\;=\;\Delta_{S^3}\!\big(g_{\rho_k}\circledast g_{\sigma_{x,k}}\big)
\;=\;\Delta_{S^3}\,g_{\sqrt{\rho_k^2+\sigma_{x,k}^2}}.
\]
Thus the welded boundary envelope is LoG in both time and space, with scales adding in quadrature from the macro/micro widths.


\paragraph{Amplitudes and payload equality.}\label{par:uom-main-011}
The tick’s operational fields include amplitudes:
\[
\text{macro (source):}\quad \kappa_a^{(k)}(t)=u_k\,\chi_{t,\Sigma_k}(t),\qquad
\text{micro (receiver):}\quad h_k(t,\Omega)=A_k\,\psi_{t,\sigma_{t,k}}(t)\,\Psi_{x,\sigma_{x,k}}(\Omega).
\]
The compiled boundary payload (admitted after the receiver window) obeys the throughput equality
\[
\Delta I_k \;=\; A_k^2\,G_k(\sigma_{t,k},\sigma_{x,k};\Lambda_k),\qquad G_k>0,
\]
so that
\[
A_k \;=\; \sqrt{\frac{\Delta I_k}{G_k(\sigma_{t,k},\sigma_{x,k};\Lambda_k)}}.
\]
Here and below, $\Delta I_k$ denotes the admitted payload; one may factor it as
$\Delta I_k=\mathcal A_{\rm rec}(\Lambda_k,\text{window})\,\Delta I_k^{\rm offer}$.
\paragraph{Main action control gauge.}\label{par:uom-main-012}
Under P3 and its more general 5D formulation, the \emph{principal} receiver--side state variable is the
Wilsonian cutoff $\Lambda_k$, whose evolution is fixed by the SK/HJ wall balance. The KKT conditions constrain
the admissible controls (window shape and passband) but do not hold $\Lambda_k$ fixed across admitted ticks.
With temporal and spatial widths of the welded pulse $s_t$ and $s_x$ locked to the cutoff scale $\Lambda^{-1}$
and the welded pulse amplitude $A_k$ eliminated, the admitted payload $\Delta I_k$ is determined jointly by the
offered source coherence, the retarded transfer kernel, and the receiver window at $\Lambda_k$ (i.e.\ the
admission map $\mathcal A_{\rm rec}$), rather than by the source alone.


\paragraph{Receiver–side OS/RP grids.}\label{par:uom-main-013}
Let $\mu>0$ be the receiver Euclidean window and $\tau_i=i\,\Delta\tau$ ($i=0,\dots,N$, $N\Delta\tau=\mu$). For each $S^3$ angular momentum $\ell$ (degeneracy $D_\ell$), build the \emph{receiver Hankel/Toeplitz} matrices from the \emph{regulated} AdS boundary two–point $C^{(\Lambda_k)}_\ell(\tau)$ (post–bridge, post–welding):
\[
H^{(\ell)}_{ij}(\mu,\Delta\tau;\Lambda_k)\;=\;C^{(\Lambda_k)}_\ell(\tau_i+\tau_j),\qquad
K^{(\ell)}_{ij}(\mu,\Delta\tau;\Lambda_k)\;=\;\sum_{p\le i}\sum_{q\le j} C^{(\Lambda_k)}_\ell((p{+}q)\Delta\tau)\,\Delta\tau^2.
\]
\emph{Reflection positivity} (RP) on the receiver discretization is $H^{(\ell)}\succeq 0$ for all $\ell$ (equivalently, $K^{(\ell)}\succeq 0$) \cite{OsterwalderSchrader1973,OsterwalderSchrader1975}. The RP \emph{deficit} at tick $k$ is the total negative trace of the spectra:
\[
\mathrm{RP}_{k}\;=\;\sum_{\ell}\sum_{\lambda_m^{(\ell)}<0}\bigl|\lambda_m^{(\ell)}\bigr|,\quad
\lambda_m^{(\ell)}\in\mathrm{Spec}\!\left(H^{(\ell)}(\mu,\Delta\tau;\Lambda_k)\right),
\]
with taper/bandlimit hygiene as in App.~\ref{app:OS-discrete}. This receiver–side RP enters the P3 objective and couples to the micro envelope via $(\sigma_{t,k},\sigma_{x,k})$ and to the cutoff $\Lambda_k$.

\paragraph{Bridge objects (receiver side).}\label{par:uom-main-014}
\begin{itemize}[leftmargin=1.5em, itemsep=2pt]
\item \emph{Regulated kernel:} ${\sf K}_\Delta^{(\varepsilon)}={\sf K}_\Delta\circ \Pi_{L(\Lambda)}\circ M_\varepsilon$; the conformality defect is measured by
$\|{\sf K}_\Delta^{(\varepsilon)} \circ U_{\rm dS}(g) - \mathrm{Ad}(U_\partial(g))\circ {\sf K}_\Delta^{(\varepsilon)}\|$.
\item \emph{Brown–York anomaly term:} on (nearly) maximally symmetric AdS$_4$ the trace reads
$\langle T^\mu{}_\mu\rangle = -\tfrac{a}{16\pi^2}E_4 + b\,\Box R$, with $E_4=\tfrac16 R^2$; the mean–curvature map is
$\gamma\langle R^2\rangle-\langle R\rangle+4\Lambda_4=8\pi G_4\,\overline{\tau}$.
\end{itemize}

\newpage
\paragraph{Bookkeeping table (symbols).}\label{par:uom-main-015}
\begin{center}
\begin{tabular}{@{}ll@{}}
\toprule
Symbol & Meaning \\
\midrule
$H$ & dS Hubble parameter; $\beta_{\rm BD}=2\pi/H$ \\
$\mathcal A_{\rm sp}$, $\mathcal A_{\rm rad}$ & static–patch algebra; radiative algebra at $\mathscr I^+$ \\
$\mathcal A_\partial$ & AdS$_4$ brane boundary algebra (receiver/code) \\
$U_{\rm dS}(g)$, $U_\partial(g)$ & unitary reps of $\mathrm{Conf}(S^3)$ on source/receiver \\
$\Delta$ & source band–limit (acts via $w_\ell(\Delta)$ on $C_\ell$) \\
$\Lambda$ & receiver cutoff (acts via $\Pi_{L(\Lambda)}$ and $M_\varepsilon$) \\
$\bar\Lambda_k$ & midpoint cutoff used to evaluate welded quantities within tick $k$ \\
$\Lambda_\infty$ & terminal cutoff reached only after exhaustion (no accepted payload) \\
$\Lambda_{\rm stop}$ & stopping cutoff: $\bar\Lambda_N$ at flip or $\Lambda_\infty$ in exhaustion \\
$\Lambda_{\rm UV}$ & fixed AdS$_5$ UV reference cutoff used in $u=\ln(\Lambda_{\rm UV}/\Lambda)$ \\
$u$ & RG time coordinate, $u:=\ln(\Lambda_{\rm UV}/\Lambda)$ \\
$I$ & cumulative admitted payload (acceptance coordinate), $I_{k+1}=I_k+\Delta I_k$ \\
$\xi$ & curvature deficit / DC trace coordinate (conjugate to $I$) \\
$\mathcal G$ & real master potential on the locked manifold (Sec.~\ref{sec:two-times-semiclassical}) \\
$\mathcal A_{\rm rec}$ & admission/acceptance factor (RG window) \\
$\mu$, $\Delta\tau$ & OS/RP Euclidean window and step \\
$H^{(\ell)}$ & OS Hankel matrix at angular momentum $\ell$ \\
${\sf K}_\Delta^{(\varepsilon)}$ & regulated conformal intertwiner (bridge kernel) \\
$\ell$, $G_5$, $G_4$ & AdS$_5$ radius; 5D/4D Newton constants ($G_4=G_5/\ell$) \\
$a=c$ & holographic central charges ($=\pi\ell^3/8G_5$) \\
$\gamma$ & anomaly coeff.\ $\gamma=\dfrac{aG_4}{12\pi}$ \\
$L_4$, $\Lambda_4$ & AdS$_4$ radius; cosmological constant ($R=4\Lambda_4=-12/L_4^2$) \\
$k$, $t_k$ & tick index and center time \\
$\sigma_{t,k}$, $\sigma_{x,k}$ & micro (time) and micro (space) widths \\
$s_{t,k}$, $s_{x,k}$ & welded (time) and welded (space) widths (decision variables) \\
$\psi_{t,\sigma}$, $\Psi_{x,\sigma}$ & temporal DoG and spatial (Gaussian/DoG) envelopes \\
$\Delta I_k$ & admitted per--tick payload (source $\times$ transfer $\times$ RG window) \\
\(\Sigma_k\) & macro (source) temporal width of GKLS envelope \(\chi_{t,\Sigma_k}\) \\
\(\rho_k\) & macro (source) spatial LoG scale on \(S^3\) (for \(\Delta_{S^3} g_{\rho_k}\)) \\
\(u_k\) & macro (source) amplitude scaling GKLS rates: \(\kappa_a^{(k)}(t)=u_k\,\chi_{t,\Sigma_k}(t)\) \\
\(A_k\) & micro (receiver) boundary amplitude in \(h_k=A_k\,\psi_{t,\sigma_{t,k}}\Psi_{x,\sigma_{x,k}}\) \\
\(G_k(\sigma_{t,k},\sigma_{x,k};\Lambda_k)\) & throughput sensitivity: \(\Delta I_k=A_k^2\,G_k\) \\
\(g_\sigma,\,g'_\sigma\) & Gaussian and its temporal derivative (DoG) \\
\(\Delta_{S^3}\) & Laplacian on \(S^3\) (zonal; commutes with zonal convolution) \\
\(\Pi_{L(\Lambda)}\) & projector onto spherical harmonics with \(\ell\le L(\Lambda)\) (receiver cutoff) \\
\(M_\varepsilon\) & mollifier/regulator in the bridge \({\sf K}_\Delta^{(\varepsilon)}\) \\
\(C_\ell(\tau)\) & Euclidean 2-pt kernel in the \(\ell\) sector (OS data) \\
\(w_\ell(\Delta)\) & band-limit window acting on \(C_\ell\) at the source \\
\(F(t),\,C(t),\,\mathcal E(t)\) & modular free energy, classical record, and excess \(F-C\) \\
\(\mathcal J_k\) & P3 dwell/action for tick \(k\) (objective) \\
\(\tau,\,\overline{\tau}\) & injected trace density and its DC component (BY source) \\
\(\Delta\!\langle R^2\rangle,\,\widehat{\Delta\!\langle R^2\rangle}\) & curvature-variance injection (phys./dimensionless) \\
\(C_{\rm DoG}\) & DoG shape constant in \(\Delta\!\langle R^2\rangle\) \\
\(\Theta_k\) & anomaly-floor increment per tick: \(\Theta_k=\gamma\,\widehat{\Delta\!\langle R^2\rangle}_k\) \\
\(\sigma_{\rm QSL}\) & quantum–speed–limit lower bound for \(\sigma_{t,k}\) \\
\(\mathcal D_{\rm RP}\) & RP deficit measure (sum of negative Hankel eigenvalues over \(\ell\)) \\
\(\mathcal D_{\rm conf}\) & conformality (intertwiner) mismatch/deficit \\
\bottomrule
\end{tabular}
\end{center}

\medskip
This subsection fixes geometry, algebras, regulators, and normalizations. All subsequent statements (locks, convexity, bridge theorems, and tick dynamics) refer to these conventions.

\newpage
\subsection{Postulates and thermodynamic principle (P1–P3)}
\label{subsec:postulates-P1-P3}

\subsubsection*{Postulate One: Unitarity cannot produce facts}
\label{subsubsec:P1}
\begin{postulate}{P1: (Unitarity cannot produce facts).}
\label{post:P1}
\textbf{Classical facts cannot be generated by lossless evolution.}
\end{postulate}
Classical facts are defined in this framework as ever-lasting (CTFM-certifiable) records. Ever-lasting means durable, with the lifetime of the aeon.

Formally, for any unitary/isometric channel $\mathcal U$ on a closed algebra and for any
unmeasured ancilla, the classical record functional cannot increase and tail fact
certifiability remains $\Pi_2$–hard (CTFM, Book I, Chapter 5 \cite{CTFMBookIII}). Therefore, prior to the flip (UOMT acceptance)
the AdS$_4$ receiver contains \emph{no classical facts}: it is restricted to the
factless vacuum code state $\rho^{\rm AdS}_0=\ket{\Omega}\!\bra{\Omega}$ inside the AdS$_4$
code subspace. The welded bridge is a regulated conformal \emph{isometry}
${\sf K}^{(\varepsilon)}_\Delta:\mathcal A_{\rm rad}(\mathscr I^+)\to\mathcal A_\partial$
that preserves factlessness: it transports amplitudes but cannot by itself produce records.
Facts appear only after the flip, when evolution becomes dissipative (QND+KMS on dS$_4$),
and the resulting records are mirrored to seed the next aeon’s factless AdS$_4$ code.


\subsubsection*{Postulate Two: Quantum Second Law (recoverable information monotonicity)}
\label{subsubsec:P2}

\paragraph{Why the postulate is needed (thermodynamic limit).}\label{par:uom-main-016}
In the static–patch description the Bunch–Davies state $\omega_{\rm BD}$ is KMS \cite{BunchDavies1978} with Gibbons–Hawking temperature\cite{GibbonsHawking1977}, and the
radiative sector of the dS source admits a coarse–grained GKLS evolution with
$\omega_{\rm BD}$ as the common fixed point (Sec.~\ref{subsec:scene-notation}).


Let 
\[
F_{\rm dS}(t):=S\!\big(\rho_t\,\|\,\rho^{\rm BD}_\Delta\big)
\]
be the modular free
energy (recoverable quantum information) of the source relative to the band–limited
anchor $\rho^{\rm BD}_\Delta$. Then Spohn’s identity gives the entropy–production rate

\[
\dot{\Sigma}(t)\equiv-\,\dot F_{\rm dS}(t)\ge 0
\]

\emph{If no recorder/exporter operates}
(i.e.\ no classical record is created anywhere), then $\dot F_{\rm dS}=-\dot{\Sigma}\le0$ and
$F_{\rm dS}(t)\downarrow 0$ in the thermodynamic limit. Thus \emph{recoverable
information is destroyed asymptotically} unless it is transferred to classical record
and/or into additional (growing) recoverable capacity. 

For any cyclic cosmology that involves dS space, this constitutes a problem: if the information is inevitably destroyed as DoFs or is scrambled by thermodynamics, how can cycles produce the vast observable Universe with the incomprehensibly special initial conditions (relative gravitational entropy w.r.t. current state $\frac{S_{\text{initial}}}{S_{\text{now}}}\sim 10^{-10^{123}}\approx 0$)? The natural assumption would be the opposite: cycles should \emph{add the information}, not destroy it.

This motivates P2.

\begin{postulate}[P2: Recoverable–information monotonicity]
\label{post:P2}
\textbf{The total recoverable information is monotonically non–decreasing.}
\end{postulate}

\begin{definition}[Total recoverable information] is
\label{def:uom-main-001}
\[
\frac{d}{dt}\,\mathcal I_{\rm tot}(t)\ \ge\ 0,
\qquad
\mathcal I_{\rm tot}(t)\ :=\ F_{\rm dS}(t)\ +\ C(t)\ +\ \mathcal C_{\rm code}\big(\Lambda(t)\big).
\]
Here $F_{\rm dS}$ is the source–side modular free energy; $C(t)$ is the unique cumulative
classical record (counted once) produced by the QND channel; and $\mathcal C_{\rm code}(\Lambda)$ is the
receiver’s recoverable \emph{code capacity} at cutoff $\Lambda$. We track the record’s
location by $C(t)=C_{\rm env}(t)+C_{\rm AdS}(t)$, where $C_{\rm env}$ is the transient
commutant record and $C_{\rm AdS}$ its compiled boundary location.
\end{definition}

\paragraph{Differential bookkeeping at fixed cutoff.}\label{par:uom-main-017}
On a time interval where the receiver cutoff $\Lambda$ is held fixed (no capacity
change), a well–aligned QND channel on the source creates classical record at rate
$\dot C_{\rm create}(t)\ge0$ that depends only on pointer–basis coherences. The compiled
bridge exports the record to the boundary at rate $\dot C_{\rm export}(t)\ge0$. We write.
Here ``held fixed'' is an \emph{intra--tick} convention: within a tick $I_k$ we evaluate
receiver quantities at a representative $\bar\Lambda_k$ (e.g.\ $\Lambda(t_k)$);
inter--tick evolution always satisfies $\Lambda_{k+1}>\Lambda_k$ for any accepted tick
with $\Delta I_k>0$.
\[
C(t)=C_{\rm env}(t)+C_{\rm AdS}(t),\qquad
\dot C(t)=\dot C_{\rm create}(t),\qquad
\dot C_{\rm AdS}(t)=\dot C_{\rm export}(t),\qquad
\dot C_{\rm env}(t)=\dot C_{\rm create}(t)-\dot C_{\rm export}(t).
\]
Then P2 implies the bound $\dot C(t)\le \dot{\Sigma}(t)$.
Define the unrecorded entropy production (slack) and realized efficiency by
\[
\dot{\Sigma}_{\rm unrec}(t):=\dot{\Sigma}(t)-\dot C(t)\ \ge\ 0,\qquad
\eta(t):=\frac{\dot C(t)}{\dot{\Sigma}(t)}\ \in[0,1]\ \ (\dot{\Sigma}>0).
\]
Hence, for fixed $\Lambda$,
\[
\frac{d}{dt}\Big(F_{\rm dS}+C\Big)
\;=\; -\,\dot{\Sigma}(t)+\dot C(t)
\;=\; -\,\dot{\Sigma}_{\rm unrec}(t)\ \le\ 0,
\]
and the derivative \emph{vanishes identically} (tickwise conservation) iff
$\dot{\Sigma}_{\rm unrec}(t)\equiv0$ (equivalently $\eta(t)\equiv1$), i.e.\ for an
ideal QND, pointer–aligned instrument and admissible bridge. Any stronger–than–QND
(disturbing) measurement or any misalignment yields $\dot{\Sigma}_{\rm unrec}>0$ on a
set of nonzero measure and strict loss at fixed $\Lambda$.

\begin{definition}[Per-tick efficiency, admissible bridge, and leakage]\label{def:eta-bridge-leak}
\textbf{Efficiency and slack.}
Fix a tick $I_k$ with receiver cutoff evolving from $\Lambda_k$ to $\Lambda_{k+1}$; evaluate
band--locked quantities at $\bar\Lambda_k$. Let the entropy production rate be
$\dot{\Sigma}(t)\ge 0$. Define the tick-averaged realized efficiency by
\[
\eta_k\ :=\ \frac{\Delta C(I_k)}{\int_{I_k}\dot{\Sigma}(t)\,dt}\ \in[0,1],
\]
and the unrecorded entropy production (slack) on the tick by
\[
\dot{\Sigma}_{\rm unrec}(t):=\dot{\Sigma}(t)-\dot C(t)\ \ge\ 0,
\qquad
\ell_k:=1-\eta_k.
\]
\textbf{Admissible band-locked bridge.}
A per-tick bridge $\mathcal B_k$ is \emph{admissible} if it satisfies:
\begin{enumerate}[label=(B\arabic*),leftmargin=2.2em]
\item \textbf{Band-locking.} The bridge is locked to the receiver cutoff with welded widths
$s_t=\alpha_t^\ast/\bar\Lambda_k$ and $s_x=\alpha_x^\ast/(v\bar\Lambda_k)$ and acts on the
passband around $k\sim\bar\Lambda_k$.
\item \textbf{Passband isometry (approximate).} There exists an isometry
$W_k:\mathcal H_{\Delta_k}\to\mathcal H_{{\rm code},k}$ with
$W_k^\dagger W_k=\mathbf 1_{\Delta_k}$ and $W_k W_k^\dagger=P_{{\rm code},k}$ such that
for any passband-supported state $\rho$,
\[
\mathcal B_k(\rho)=W_k\rho W_k^\dagger
\qquad\text{or}\qquad
\|\mathcal B_k(\rho)-W_k\rho W_k^\dagger\|_1\le \varepsilon_k.
\]
\item \textbf{Code-range support.} For all inputs $\rho$, the output is supported on the receiver code:
$\mathcal B_k(\rho)=P_{{\rm code},k}\,\mathcal B_k(\rho)\,P_{{\rm code},k}$.
\end{enumerate}
\textbf{Commutant record and leakage instrument.}
Let $\mathbb E_{\rm ptr}$ be the conditional expectation onto the pointer algebra and
let $\mathcal R_k$ be the commutant-record map with \emph{commutative range}
(a classical record algebra). Define the recorded information increment by
\[
\Delta C(I_k):=S\!\big(\mathcal R_k(\rho_k)\,\|\,\mathcal R_k(\rho^{\rm BD}_\Delta)\big),
\]
with location bookkeeping $\Delta C=\Delta C_{\rm env}+\Delta C_{\rm AdS}$ under export.
Let $\Phi_k$ be the actual CPTP tick map. A \emph{bridge--leakage decomposition} is a
two-outcome instrument $\{\mathcal B_k,\mathcal L_k\}$ of CP trace-nonincreasing maps
such that $\Phi_k=\mathcal B_k+\mathcal L_k$.
The complementary branch $\mathcal L_k$ is a \emph{leakage component} if it is nonzero on
admissible inputs and either
\begin{enumerate}[label=(L\arabic*),leftmargin=2.2em]
\item produces outputs outside the receiver code subspace, or
\item yields positive slack $\ell_k>0$ (unrecorded entropy production).
\end{enumerate}
\end{definition}

\begin{proposition}[Efficiency locking in the admissible class]\label{prop:eta-lock}
Assume a pointer--aligned QND channel with a canonical commutant KMS bath, an admissible
band--locked bridge $\mathcal B_k$, and no leakage branch on admissible inputs
(i.e.\ $\mathcal L_k=0$ in Def.~\ref{def:eta-bridge-leak}). Then
\[
\dot{\Sigma}_{\rm unrec}(t)\equiv 0\quad\text{on admissible ticks}\qquad\Rightarrow\qquad \eta_k=1.
\]
Conversely, any persistent $\eta_k<1$ implies $\ell_k>0$ and hence a nontrivial leakage
branch or misalignment outside the commutant/pointer/code sector.
\end{proposition}

\begin{proof}[Sketch]
For pointer--aligned QND dynamics, the conditional expectation onto the pointer algebra is
sufficient for the commutant record: all decoherence is captured by the complementary
record map, and equality in data--processing (Petz sufficiency) holds on the admissible
subalgebra. With an admissible isometric bridge and no leakage branch, the recorded
information is exported without in--band loss, so $\dot{\Sigma}_{\rm unrec}=0$.
If $\eta_k<1$ persists, some entropy production is unrecorded or exits the code subspace,
which by Def.~\ref{def:eta-bridge-leak} is a leakage branch or misalignment.
\end{proof}

\paragraph{Ticks as information–preserving export events.}\label{par:uom-main-018}
We define \emph{ticks} as maximal time windows $I_k=[t_k^- ,t_k^+]$ on which (i) the
absorber acts QND on the source pointer algebra and (ii) the welded bridge is
admissible at the evolving receiver cutoff $\Lambda(t)$, with acceptance (AZCT/UOMT)
at $t_k$. Each tick starts at $\Lambda_k:=\Lambda(t_k^-)$ and ends at
$\Lambda_{k+1}:=\Lambda(t_k^+)$, and we evaluate welded quantities at a midpoint
scale $\bar\Lambda_k$ (e.g.\ $\bar\Lambda_k:=\Lambda(t_k)$ or
$\bar\Lambda_k:=\sqrt{\Lambda_k\Lambda_{k+1}}$), with $O(\Delta\Lambda/\Lambda)$
corrections. Within a tick one has the \emph{tickwise balance}:
\[
\int_{I_k}\!\!\Big(\dot F_{\rm dS}+\dot C\Big)\,dt
\;=\;-\!\int_{I_k}\!\dot{\Sigma}_{\rm unrec}(t)\,dt\ =\ 0\qquad\text{(ideal QND/bridge).}
\]
Thus each accepted tick is an \emph{information–preserving conversion} of quantum
recoverable information into classical record, while the code capacity increases
whenever $\Delta I_k>0$.

\paragraph{Necessity of QND (no–loss at fixed capacity).}\label{par:uom-main-019}
The inequality above implies two necessary conditions to avoid net loss at fixed
$\Lambda$:
\begin{enumerate}[label=(\alph*),nosep]
\item \emph{No recorder} ($\dot C=0$) $\Rightarrow$
$\dot F_{\rm dS}=-\dot{\Sigma}\le0$ and $F_{\rm dS}\!\downarrow 0$ (loss in the
thermodynamic limit).
\item \emph{Non–QND or misaligned} recorder ($\eta<1$) $\Rightarrow$
$\dot{\Sigma}_{\rm unrec}>0$ on a set of times of nonzero measure and
$\dot F_{\rm dS}+\dot C<0$ (strict loss).
\end{enumerate}
Therefore, \emph{QND, pointer alignment, and bridge admissibility are necessary and
sufficient} for tickwise conservation at fixed $\Lambda$ (intra--tick only; see the
receiver dynamics axiom).

\paragraph{When does the total \(\mathcal I_{\rm tot}\) strictly increase?}\label{par:uom-main-020}
Strict growth arises exactly when the \emph{code capacity} increases:
\[
\frac{d}{dt}\,\mathcal C_{\rm code}\big(\Lambda(t)\big)\ \ge\ 0,
\quad\text{with equality iff }\dot\Lambda(t)=0.
\]
Concretely, let $\Pi_{L(\Lambda)}$ be the spectral projector implementing the boundary
cutoff and let $\rho^{\rm code}_\Lambda$ be the code–space anchor (the ``ideal packing''
state at cutoff $\Lambda$). Define $\mathcal C_{\rm code}(\Lambda)$ as the recoverable
capacity gained when refining from $\Lambda$ to $\Lambda+d\Lambda$,
\[
\frac{d}{d\Lambda}\,\mathcal C_{\rm code}(\Lambda)
\;:=\;\left.\frac{d}{d\epsilon}\right|_{\epsilon=0^+}
S\!\Big(\rho^{\rm code}_{\Lambda+\epsilon}\,\Big\|\,\rho^{\rm code}_{\Lambda}\Big)\ \ge\ 0,
\]
where nonnegativity follows from the approximate–intertwiner lemma
(Lem.~\ref{lem:approx-intertwiner-new}: refining the receiver increases admissible recoverable content).
Hence, over a tick $I_k$ the total change is
\[
\Delta\mathcal I_{\rm tot}\big|_{I_k}
\;=\;\int_{I_k}\!\!\big(\dot F_{\rm dS}+\dot C\big)\,dt\;+\;
\int_{I_k}\!\partial_\Lambda\mathcal C_{\rm code}(\Lambda)\,\dot\Lambda\,dt,
\]
so any accepted tick with $\Delta I_k>0$ forces $\dot\Lambda>0$ on a subinterval and
therefore $\Delta\mathcal C_{\rm code}>0$.

\paragraph{Summary.}\label{par:uom-main-021}
P2 asserts $\dot{\mathcal I}_{\rm tot}\ge0$. At fixed cutoff (intra--tick only; see the
receiver dynamics axiom) this reduces to \emph{tickwise conservation} under ideal QND export;
absent a recorder or with any non–QND disturbance, recoverable information is lost. Growth
of total recoverable information is \emph{equivalent} to cutoff refinement, whose positivity
is guaranteed by the conformality/RP bridge construction.

\paragraph{Receiver dynamics axiom (canonical).}\label{par:uom-main-022}
We use the following receiver dynamics as the single source of truth:
\begin{enumerate}[label=(R\arabic*),leftmargin=2.0em,itemsep=2pt,topsep=2pt]
\item \textbf{KR-only active receiver.} The active receiver is always a KR head
(EOW cap), and at each split the continuing receiver lies just UV-side of the newborn source:
$\Lambda_{\rm R}(t_{\rm split})>(1+\varepsilon)\Lambda_{\rm S}(t_{\rm split})$.
There is no RS1 receiver stage; RS1 segments appear only as interior/shadow cells.
\item \textbf{Forced wall motion per admitted tick.} If a tick is accepted with nonzero
admitted payload, then the receiver cutoff must advance:
$\Delta I_k>0\Rightarrow \Lambda_{k+1}>\Lambda_k$. Equivalently, the wall balance
$\mu_{\rm RG}(\Lambda)\,\dot\Lambda=(\mathcal P_{\rm in}-\mathcal P_{\rm out})
-\dot{\mathcal F}_{\rm rec}$ forces $\dot\Lambda>0$ on a subinterval of any accepted tick.
\item \textbf{Two outcomes per aeon.} Either (A) curvature crosses zero in finite ticks
(split/flip), or (B) the source exhausts before crossing and the receiver becomes the
terminal boundary condition of the remaining AdS$_5$ domain.
\item \textbf{No dwell with payload.} A finite terminal $\Lambda_\infty$ may emerge if
total wall work is finite, but $\Lambda$ cannot be held fixed during accepted ticks;
exact stasis occurs only when no accepted payload remains. (Here “terminal boundary”
means the outer boundary condition of the remaining AdS$_5$ domain; it is an EOW cap at
a finite cutoff $\Lambda_\infty$, and coincides with the conformal AdS$_5$ UV boundary
only in the limit $\Lambda_\infty\to\infty$.)
\end{enumerate}


\subsubsection*{Postulate Three: Least thermodynamic action of transport}
\label{subsubsec:P3}
\begin{postulate}[P3: Least thermodynamic action of the Universe realization transport]
\label{post:P3}
    \textbf{Over an aeon, the thermodynamic action of the information transport from the source to the receiver brane is minimized.}
\end{postulate}
\paragraph{Specifically,}across the \emph{entire} tick sequence, Nature chooses each tick's micro and macro pulse parameters and the\label{par:uom-main-023}
receiver refinement so as to \emph{minimize the thermodynamic action of transport}—%
equivalently, the time–integrated \emph{net modular free energy} carried on the source
and on the receiver boundary.

\paragraph{Precise statement.}\label{par:uom-main-024}
Let $I_k=[t_k-\Delta_k,t_k+\Delta_k]$ be the $k$-th tick, with controls
$u_k=\big(\sigma_{t,k},\sigma_{x,k},A_{t,k},A_{x,k},\delta t_k,\delta\Omega_k,\Delta_k,\Lambda_k\big)$
(temporal/spatial widths and amplitudes, micro–macro alignment shifts, source band–limit,
receiver cutoff). Denote by
$F_{\rm dS}(t):=S\!\big(\rho_t\,\|\,\rho^{\rm BD}_{\Delta_k}\big)$
the source–side modular free energy (Sec.~\ref{subsubsec:P2}), and by
$F_\partial^{\rm net}\!\big(t;\Lambda_k\big)$ the receiver–side net modular free energy
at cutoff $\Lambda_k$, both measured in nats. The \emph{P3 action} over the totality
of ticks is
\begin{equation}
\label{eq:P3-action-master}
\mathcal S_{\rm P3}\;:=\;\sum_{k=0}^{K-1}\ \int_{I_k}\!\Big[\,F_{\rm dS}(t)\;+\;F_\partial^{\rm net}\big(t;\Lambda_k\big)\,\Big]\,dt,
\end{equation}
and P3 states that the realized controls $u_k^\star$ (and induced states) minimize
$\mathcal S_{\rm P3}$ within the guard–admissible class (Sec.~\ref{sec:aeon-evolution}).
\emph{Equivalently}, if one prefers a bona fide thermodynamic action (J$\cdot$s),
multiply each integrand by its modular temperature $k_B T_{\rm mod}$; the minimizer is
unchanged.

\paragraph{Receiver modular energy: physical vs.\ unphysical parts.}\label{par:uom-main-025}
For a fixed tick and cutoff $\Lambda_k$, decompose the receiver contribution as
\begin{equation}
\label{eq:phys-unphys-split}
F_\partial^{\rm net}\big(t;\Lambda_k\big)\;=\;
\underbrace{F_\partial^{\rm phys}\big(t;\Lambda_k\big)}_{\text{code–legal, payload–bearing}}
\;+\;
\underbrace{F_\partial^{\rm unphys}\big(t;\Lambda_k;u_k\big)}_{\text{penalties (bridge \& guards)}}.
\end{equation}
\begin{itemize}[leftmargin=1.5em, itemsep=2pt]
\item $F_\partial^{\rm phys}$ is the modular free energy \emph{of the code–admissible
part} of the boundary state—i.e.\ the portion that survives the (regulated)
intertwiner ${\sf K}^{(\varepsilon)}_\Delta$ and the RP/OS projector and can be packed
into the AdS$_4$ code at cutoff $\Lambda_k$. Minimizing this part means optimizing
the \emph{payload degrees of freedom themselves} (what is recorded/exported).
Analyzing that joint optimization (which feeds back into the dS POVM choice) is
deferred to future work; in this paper we treat $F_\partial^{\rm phys}$ as a
\emph{tickwise constant} determined by the minimally decohering payload that keeps
P2 non–decreasing (Post.~\ref{post:P2}).
\item $F_\partial^{\rm unphys}$ collects \emph{all deviations from ideal code packing}
at fixed payload: Brown–York/anomaly rectification (variance injection on AdS$_4$),
reflection–positivity deficit (OS/Hankel negatives on the receiver grid), conformality
misfit of the regulated intertwiner, and bridge alignment/cross–talk penalties. These
terms depend on the tick controls $u_k$ and vanish on the ideal welded LoG envelope
with perfect guards. Their explicit convex forms (and weights fixed by central charge,
normalizations, and guard scales) are given in Sec.~\ref{subsec:P3-master}.
\end{itemize}

\paragraph{What is minimized in this work.}\label{par:uom-main-026}
With $F_\partial^{\rm phys}$ held tickwise constant (minimal acceptable payload)
and $F_{\rm dS}$ slaved by QND export (Sec.~\ref{subsec:source-qnd-commutant}), the optimization in
this paper reduces tickwise to
\begin{equation}
\label{eq:P3-reduced}
\min_{u_k\in\mathcal U_k^{\rm guard}}
\ \int_{I_k}\!F_\partial^{\rm unphys}\big(t;\Lambda_k;u_k\big)\,dt,
\qquad
\mathcal U_k^{\rm guard}:=\{\text{QSL, RP, conformality, bridge–admissible controls}\}.
\end{equation}
Aside from the receiver cutoff update $\Lambda_{k+1}$ (fixed by the wall--work balance, with feasibility
constraints limiting the step size), the variables locked by action minimum are the welded pulse widths
$(s_{t,k},s_{x,k})$, the micro–macro alignment shifts $(\delta t_k,\delta\Omega_k)$, and
(at tick boundaries) the bandlimit refinement choices $(\Delta_{k\!+\!1})$. As shown in
Sec.~\ref{subsec:P3-master}, the convex
structure of $F_\partial^{\rm unphys}$ yields: 
\begin{enumerate}[label=(\roman*)]
    \item a \emph{zero shift} lock
($\delta t_k=\delta\Omega_k=0$);
\item a \emph{width} lock (LoG scale set by the welded
micro–macro convolution with widths locked to the receiver brane cutoff scale $\Lambda^{-1}$); and
\item a \emph{monotone} evolution of the anomaly floor that is fast–then–slow in tick
number, providing the variable–rate geometric behavior required for the flip model.
\end{enumerate}


\subsection{Source side: QND absorber with intrinsic KMS bath (commutant picture)}
\label{subsec:source-qnd-commutant}

\paragraph{Algebraic setup and KMS anchor.}\label{par:uom-main-027}
Let $(\mathcal A_{\rm sp},\alpha_t)$ be the dS static–patch von Neumann algebra with
Bunch–Davies KMS state $\omega_{\rm BD}$ at inverse temperature $\beta_{\rm BD}=2\pi/H$.
In the standard (GNS) form $(\pi,\mathcal H_\omega,\Omega)$ the state is cyclic and
separating, so Tomita–Takesaki theory equips $(\mathcal A_{\rm sp},\omega_{\rm BD})$ with
its modular group $\sigma_t^\omega$ and the commutant algebra $\mathcal A_{\rm sp}'$.
We \emph{identify} the intrinsic environment with the commutant:
\[
\mathcal A_{\rm env}\equiv\mathcal A_{\rm sp}',\qquad
\omega_{\rm env}\equiv\omega_{\rm BD}\!\restriction_{\mathcal A_{\rm sp}'}.
\]
This choice makes the dilation \emph{local} on $\mathcal A_{\rm sp}\vee\mathcal A_{\rm sp}'$
and introduces no external controller.

\paragraph{Static–patch Hamiltonian and KMS (for QND).}\label{par:uom-main-028}
Restricting BD to a static patch with metric
$ds^2=-(1-H^2 r^2)\,dt^2+\frac{dr^2}{1-H^2 r^2}+r^2 d\Omega_{d-1}^2$ yields a
KMS state at inverse temperature $\beta=2\pi/H$ for the static Hamiltonian
$H_{\mathrm{sp}}=\partial_t$:
\[
\langle A(t) B\rangle_{\mathrm{BD}}=\langle B A(t+i\beta)\rangle_{\mathrm{BD}},
\qquad \beta=\frac{2\pi}{H}.
\]
In particular, when the QND pointer is chosen as $O=f(H_{\mathrm{sp}})$, the
band–limited static–patch BD state $\rho^{\mathrm{BD}}_\Lambda\propto
\exp(-\beta H_{\mathrm{sp}})$ restricted to the $\ell\le \ell_{\max}$ sector is
\emph{fixed} by the QND channel.

\paragraph{QND coupling and tick–localized unitary.}\label{par:uom-main-029}
A tick $I_k=[t_k-\Delta_k,t_k+\Delta_k]$ is produced by a time–dependent,
quantum–non–demolition (QND) coupling between a \emph{band–limited} source observable
$O=f(H_{{\rm sp},\Delta})\in\mathcal A_{\rm sp}$ and an environment operator
$B\in\mathcal A_{\rm sp}'$:
\begin{equation}
U_{\chi_k}\;=\;\mathcal T\exp\!\Big(i\!\int_{I_k}\!\chi_k(t)\,O\!\otimes\! B\,dt\Big),
\qquad [O,H_{{\rm sp},\Delta}]=0,
\label{eq:source-qnd-unitary}
\end{equation}
with a smooth switching envelope $\chi_k\in C_0^1(I_k)$.
In the weak–coupling (linear–response) limit, tracing out $\mathcal A_{\rm env}$ yields a
GKLS generator with a common KMS fixed point:
\begin{equation}
\dot\rho_t\;=\;\mathcal L_t(\rho_t)\;=\;-\,i[H_{{\rm sp},\Delta},\rho_t]
+\!\sum_a\kappa_{a,k}(t)\Big(L_a\rho_t L_a^\dagger-\tfrac12\{L_a^\dagger L_a,\rho_t\}\Big),
\quad [L_a,H_{{\rm sp},\Delta}]=0,\ \ \mathcal L_t(\rho^{\rm BD}_\Delta)=0,
\label{eq:source-qnd-gkls}
\end{equation}
where the rates are shaped by the envelope, $\kappa_{a,k}(t)=u_{a,k}\,\chi_k(t)$.

\paragraph{Source modular free energy and entropy production.}\label{par:uom-main-030}
Fix the band–limited BD reference $\rho^{\rm BD}_\Delta$ and define the
source modular free energy
\begin{equation}
F_{\rm dS}(t)\;:=\;S\!\big(\rho_t\big\|\,\rho^{\rm BD}_\Delta\big).
\label{eq:FdS-def}
\end{equation}
Spohn’s inequality for \eqref{eq:source-qnd-gkls} gives the nonnegative entropy–production rate
\begin{equation}
\sigma_{\rm dS}(t)\;:=\;-\frac{d}{dt}\,F_{\rm dS}(t)
\;=\;-\Tr\!\big[\mathcal L_t(\rho_t)\,(\log\rho_t-\log\rho^{\rm BD}_\Delta)\big]\ \ge\ 0,
\label{eq:spohn-source}
\end{equation}
with equality iff $\rho_t$ is diagonal in the $O$–pointer basis (no residual coherences).

\paragraph{Pointer algebra and classical record on the environment.}\label{par:uom-main-031}
Let $\mathcal A_{\rm ptr}\subset\mathcal A_{\rm env}$ be the commuting pointer algebra
measured by $B$. The classical record accumulated on the pointer during the tick is an
absolutely continuous functional $C_{\rm env}(t)$ with rate
\begin{equation}
\dot C_{\rm env}(t)\;=\;\Gamma\big[\{\kappa_{a,k}(t)\},\rho_t\big]\ \ge\ 0,
\qquad \Gamma\ \ \text{depends only on pointer coherences.}
\label{eq:record-rate}
\end{equation}
In the calibrated linear–response regime (one–shape reduction), pick a unit–shape
$\phi_k\in L^2(I_k)$ and write $\chi_k(t)=A_k\,\phi_k(t)$. Then the \emph{per–tick payload}
(i.e.\ the exported classical fact content) obeys the equality
\begin{equation}
\Delta I_k\;=\;\langle \chi_k,\ \mathsf S\,\chi_k\rangle
\;=\;A_k^2\,\underbrace{\langle \phi_k,\ \mathsf S\,\phi_k\rangle}_{=:G_k>0},
\qquad \langle f,g\rangle:=\int_{I_k}\!f(t)\,g(t)\,dt,
\label{eq:payload-equality}
\end{equation}
for a positive throughput form $\mathsf S$ fixed by the $O\otimes B$ channel. Hence
\begin{equation}
A_k\;=\;\sqrt{\Delta I_k/G_k},\qquad
P_k:=\|\chi_k\|_2^2\;=\;\frac{\Delta I_k}{G_k}\quad\text{(no free gain).}
\label{eq:source-amplitude-fluence}
\end{equation}
Here $\Delta I_k$ is the admitted payload; any offered payload is absorbed into
$\Delta I_k^{\rm offer}$ via $\Delta I_k=\mathcal A_{\rm rec}\,\Delta I_k^{\rm offer}$
(see Sec.~\ref{sec:two-times-semiclassical}).

\paragraph{Admissible switching class and guards.}\label{par:uom-main-032}
The envelope is drawn from a convex guard–admissible set
\begin{equation}
\mathcal A_k\;=\;\Big\{\chi\in C_0^1(I_k):\
\|\chi\|_2^2=P_k,\ \ \int_{I_k}\!\chi=0\ \ (\text{zero–DC}),\ \
\|{\sf BL}_{\Delta_{\rm src}}^\perp\chi\|_2\le \varepsilon_{\rm BL},\ \
\mathrm{Var}_t[\chi]=\Sigma_k^2\ge \Sigma_{\rm QSL}^2\Big\},
\label{eq:admissible-class}
\end{equation}
where ${\sf BL}_{\Delta_{\rm src}}$ projects to the source band–limit and
$\Sigma_{\rm QSL}$ is the quantum–speed–limit lower bound on the temporal width.

\paragraph{Source contribution to P2 and P3.}\label{par:uom-main-033}
Within \eqref{eq:admissible-class}, identify the created record with $C(t)$
(so $\dot C=\dot C_{\rm env}$ before export). The source piece of the total information
balance (P2) and the transport action (P3) are
\begin{equation}
\frac{d}{dt}\Big(F_{\rm dS}(t)+C(t)\Big)\;=\;-\,\sigma_{\rm dS}(t)+\dot C(t)\;=\;-\,\dot{\Sigma}_{\rm unrec}(t)\ \le\ 0,
\qquad
\mathcal S_{\rm P3}^{\rm (src)}(I_k)\;=\;\int_{I_k}\!F_{\rm dS}(t)\,dt.
\label{eq:P2-P3-source}
\end{equation}
Thus, larger $\sigma_{\rm dS}$ (faster QND diagonalization under the guards) \emph{reduces}
the source action. The tradeoff is entirely mediated by $\chi_k$: for fixed payload
$\Delta I_k$ the shape $\phi_k$ that maximizes $\int_{I_k}\sigma_{\rm dS}$ (equivalently,
minimizes dwell of $F_{\rm dS}$) is the unique admissible eigen–shape singled out by the
guards; in the parity–symmetric class this yields the DoG–type temporal profile used
throughout (Sec.~\ref{subsubsec:shapes-optimality}).

\paragraph{Persistent KMS bath \& exponentially–decaying backflow.}
\label{par:exp-backflow}
The environment is the commutant $\mathcal A_{\rm sp}'$ in the BD KMS state
$\omega_{\rm BD}$, so its two–point correlator for the coupled operator $B\in\mathcal A_{\rm sp}'$
obeys the KMS condition and decays on the thermal (Hubble) scale:
\[
C_B(t)\;:=\;\omega_{\rm BD}\!\big(B(t)B(0)\big),\qquad
C_B(t+i\beta_{\rm BD})=C_B(-t),\qquad
|C_B(t)|\ \le\ C_B(0)\,e^{-|t|/\tau_{\rm KMS}},
\]
with $\beta_{\rm BD}=2\pi/H$ and $\tau_{\rm KMS}=\Theta(H^{-1})$.
In the weak–coupling, time–coarse–grained limit this yields a GKLS generator
\eqref{eq:source-qnd-gkls} with rates given by the bath spectrum $S_B(\omega)$.
Non–Markovian corrections appear as a memory kernel $K(t)$ whose $L^1$ norm is bounded
by $\int_0^\infty\!|K(t)|dt\lesssim\|C_B\|_{L^1}\sim\tau_{\rm KMS}C_B(0)$. Hence any
information backflow (revival of coherences or trace–distance increases) across a tick of
length $2\Delta_k$ is exponentially suppressed:
\[
\text{Backflow}_{I_k}\ \lesssim\ e^{-2\Delta_k/\tau_{\rm KMS}}\quad\Rightarrow\quad
\int_{I_k}\!\sigma_{\rm dS}(t)\,dt\;=\;F_{\rm dS}(t_k-\Delta_k)-F_{\rm dS}(t_k+\Delta_k)\ \big(1+O(e^{-2\Delta_k/\tau_{\rm KMS}})\big),
\]
so the KMS bath makes backflow transient and decaying, and the Spohn production controls the net drop of $F_{\rm dS}$.

\paragraph{Collapse POVM drift (compatibility and control).}\label{par:uom-main-034}
Let the (source–side) instrument on the pointer algebra be
$\{\mathsf M^{(k)}_x:\mathcal H\to\mathcal H\}_x$ with effects
$E^{(k)}_x=(\mathsf M^{(k)}_x)^\dagger\mathsf M^{(k)}_x$ and $\sum_x E^{(k)}_x=\mathbf 1$.
Slow changes of the radiative sector (band–limit $\Delta_{{\rm src},k}$ and envelope $\chi_k$)
induce an \emph{adiabatic} drift of the POVM within a calibrated family $E_x(\lambda)$:
\[
E^{(k)}_x\;=\;E_x(\lambda_k),\qquad
\|\partial_\lambda E_x(\lambda)\|_1\ \le\ L_x<\infty,\qquad
|\lambda_{k+1}-\lambda_k|\ \le\ c_\lambda\,\big(\|\chi_{k+1}-\chi_k\|_2+|\Delta_{{\rm src},k+1}-\Delta_{{\rm src},k}|\big).
\]
Thus $\|E^{(k+1)}_x-E^{(k)}_x\|_1\le L_x\,|\lambda_{k+1}-\lambda_k|$ and the drift is Lipschitz–controlled by the tick–to–tick changes of the admissible control. In the DoG–locked class (parity even around $t_k$) the drift preserves the pointer basis and only changes the outcome scales; this ensures compatibility with the welded DoG$\to$LoG boundary envelope.




\begin{definition}[Retarded export bandedge]
\label{def:retarded-bandedge}
Let $\Sigma_{\rm rec}$ be the receiver’s spatial slice and $-\Delta_{\Sigma_{\rm rec}}$ its Laplace–Beltrami
operator with eigenpairs $\{-\Delta_{\Sigma_{\rm rec}}\varphi_{\lambda,\mu}=\lambda\,\varphi_{\lambda,\mu}\}$.
The receiver code enforces a UV cutoff $\Lambda$ by restricting to $\lambda\le \Lambda^2$.
For a welded pulse that is slow compared to the receiver gap we may take the $\omega\simeq 0$ (DC) limit
of the retarded brane–to–brane transfer in the trace channel. In momentum variables
$k=\sqrt{\lambda}$ this acceptance factor has the standard screened–Yukawa form
\begin{equation}
\label{eq:ret-accept}
\begin{aligned}
\big|\varepsilon(\lambda;\Lambda)\big|^2
&= \underbrace{W_\Lambda(k)^2}_{\text{receiver window}}\,
   \underbrace{\exp\!\big[-\,2d\,\sqrt{k^2+m_\bullet^2}\,\big]}_{\text{bulk Yukawa (separation $d$)}}
\\[-2pt]
&\qquad\times
   \underbrace{\Bigg(\frac{\sqrt{k^2+m_\bullet^2}}{\sqrt{k^2+m_\bullet^2}+r_c(\Lambda)\,k^2}\Bigg)^{\!2}}_{\text{induced--gravity screen}},
\qquad k=\sqrt{\lambda}\,.
\end{aligned}
\end{equation}

where $m_\bullet$ is the lowest scalar gap in the trace channel (the GW-stabilized radion mass in RS1, or the
KR boundary-distance scalar mass in the \emph{receiver/source} trace channel); the spin--2 gap $m_g\sim L_4^{-1}$
controls only the TT sector and is not used in this scalar kernel. Here $r_c(\Lambda)$ is the induced--gravity
crossover length at cutoff $\Lambda$, $d$ the proper separation, and $W_\Lambda$ the smooth receiver UV window.
Let $\rho_{\rm src}(\lambda)$ denote the source–side spectral \emph{power density} available to drive the
mode $\lambda$ under the welded envelope (this already includes the density of states on $\Sigma_{\rm rec}$);
let $\Pi_{\rm src}(\lambda)$ be any prescribed, nondecreasing \emph{cumulative cost} for activating modes
up to $\lambda$ (e.g. constraints induced by finite temporal/spatial widths, actuator power, or guard terms).
The net exported variance accumulated up to $\lambda$ is then
\[
{\cal P}(\lambda)\;=\;\int_0^\lambda \rho_{\rm src}(\lambda')\,\big|\varepsilon(\lambda';\Lambda)\big|^2\,d\lambda',
\]
and, for a fixed admissible cost budget $\Pi_{\rm src}(\lambda)$, the optimal spectral threshold
$\lambda_\star$ (\emph{export bandedge}) is characterized by the Euler–Lagrange (KKT) knee condition
\begin{equation}
\label{eq:knee}
\rho_{\rm src}(\lambda_\star)\,\big|\varepsilon(\lambda_\star;\Lambda)\big|^2
\;=\;\pi_{\rm src}(\lambda_\star),\qquad \pi_{\rm src}(\lambda):=\frac{d}{d\lambda}\,\Pi_{\rm src}(\lambda).
\end{equation}
In words: the bandedge is placed where the marginal accepted signal power equals the marginal cost of
pushing the spectrum further. Because $\big|\varepsilon(\lambda;\Lambda)\big|^2$ \emph{decreases} with
$\lambda$ due to the Yukawa factor and the induced–gravity screen, while $\pi_{\rm src}$ is nondecreasing,
this knee exists and is unique under mild regularity. We denote the corresponding export bandedge by
$\Delta_{\rm ret}:=\lambda_\star$ and implement a taper $w_\lambda(\Delta_{\rm ret})$ that is unity for
$\lambda\ll \Delta_{\rm ret}$ and smoothly falls to $0$ for $\lambda\gg \Delta_{\rm ret}$. This selection
is entirely receiver–side (retarded acceptance) and does not invoke any static–patch spectral decomposition.
\end{definition}


\paragraph{DP scaling: collapse time–energy bound and QSL width.}\label{par:uom-main-035}
Let $\Delta E_{{\rm band},k}$ denote the energy dispersion available in the
band–limited source sector during tick $k$ (set by $\Delta_{{\rm src},k}$ and the spectral
weight of $O$ on that band). The open–system quantum speed limits (Mandelstam–Tamm
and Margolus–Levitin generalizations) give, for fidelity drop compatible with record
creation on the pointer,
\[
\tau_{\rm coll,k}\ \ge\ \max\!\left\{\frac{\hbar\,\mathcal L^2(\rho_{t_k-\Delta_k},\rho_{t_k+\Delta_k})}{\overline{\|\mathcal L_t(\rho_t)\|}},\ \ \frac{\hbar}{2\,\Delta E_{{\rm band},k}}\right\},
\]
where $\mathcal L(\cdot,\cdot)$ is the Bures angle and $\overline{\|\mathcal L_t(\rho_t)\|}$ a
time–average of a contractive norm of the Lindbladian. Specializing to the DoG envelope
with temporal width $\Sigma_k$ (macro/source) and the welded micro width $\sigma_{t,k}$,
our calibration implies a \emph{QSL guard} of the form
\[
\sigma_{t,k}\ \ge\ \frac{\chi_{\rm QSL}}{\Delta E_{{\rm band},k}},
\qquad
\Sigma_k\ \ge\ \frac{\tilde\chi_{\rm QSL}}{\Delta_{{\rm src},k}},
\]
with dimensionless constants $(\chi_{\rm QSL},\tilde\chi_{\rm QSL})$ fixed by the chosen
DoG shape (e.g.\ via $\|g'_\sigma\|_2^2=\tfrac{1}{4\sqrt\pi}\sigma^{-3}$). Thus the minimal
collapse dwell (and hence the minimal admissible widths) tracks the \emph{band energy}
available on the source: as transport refines the receiver and permits larger
$\Delta_{{\rm src},k}$, the QSL lower bounds tighten, coherently with the RP/conformality guards.

\paragraph{Mass sector: DP–type dephasing with (almost) no heating.}\label{par:uom-main-036}
In addition to the radiative QND absorber (exactly energy–preserving), we couple confined brane matter to the intrinsic KMS commutant via the renormalized mass density $\hat\mu_\Delta(x)$ (band–limited at the source). The tick–local coarse–grained generator adds a Diósi–Penrose–type quadratic form,
\[
\mathcal L_{\rm DP}^{(k)}\rho
=\frac{K_k^{\rm DP}}{2\hbar}\!\int d^3x\,d^3y\;V_{\rm eff}(x,y)\,[\hat\mu_\Delta(x),[\hat\mu_\Delta(y),\rho]],
\qquad
E_G^{\rm eff}[\delta\mu]=\frac12\!\int\! \delta\mu\,V_{\rm eff}\,\delta\mu,
\]
so that off–diagonal branches labelled by $\mu_{1,2}$ are suppressed as
$\|\rho_{12}\|\mapsto \exp\!\big[-(K_k^{\rm DP}/\hbar)\,E_G^{\rm eff}[\mu_1-\mu_2]\big]\|\rho_{12}\|$.
Here $V_{\rm eff}$ is the static brane Green’s kernel (brane–to–brane Newtonian limit). In the operational window of a tick (sub–curvature separations and quasi–static masses), $V_{\rm eff}(r)\simeq G_4/r$, and one recovers the standard 4D DP scaling. Brane–localization details (KR, RS1, RS2) only modify $V_{\rm eff}$ by Yukawa tails at large $r$; these \emph{reduce} $E_G^{\rm eff}$ but do not change the structure of the channel. Crucially, this mass–sector addition is \emph{almost QND}: for band–limited, slowly evolving sources
$\|[\hat\mu_\Delta,H_{\rm sp}]\|\le \varepsilon_{\rm ad}\|\hat\mu_\Delta\|$ with $\varepsilon_{\rm ad}\ll1$, and the energy drift per tick is scheduler–guarded,
\[
\big|\Delta\!\langle H_{\rm sp}\rangle\big|_{I_k}
\;\le\; \frac{K_k^{\rm DP}}{2\hbar}\,
\|V_{\rm eff}\|_{1}\;\big\|[\hat\mu_\Delta,[H_{\rm sp},\hat\mu_\Delta]]\big\|
\;=\;O\!\big(K_k^{\rm DP}\,\varepsilon_{\rm ad}\big),
\]
leaving the radiative no–heating intact and keeping the mass back–reaction below the guard. A local Stinespring dilation in the commutant (Gaussian KMS ancilla) realizes the channel with no signalling. Full construction, kernels for KR/RS1/RS2, and the bounds appear in App.~\ref{app:brane-localization-dp}.




\paragraph{What this fixes (and what it leaves to the receiver).}\label{par:uom-main-037}
Equations \eqref{eq:payload-equality}–\eqref{eq:source-amplitude-fluence} \emph{lock} the
source amplitude and fluence once the per–tick payload $\Delta I_k$ and the admissible
width $\Sigma_k$ are chosen. The source guards (band–limit, QSL, zero–DC) determine the
\emph{shape} $\phi_k$.
No external bath is introduced: the commutant $\mathcal A_{\rm sp}'$ at the BD KMS anchor
is the only environment, and the QND structure ensures $\rho^{\rm BD}_\Delta$
is a common fixed point. The receiver–side consequences (RP, conformality, BY/anomaly
rectification, cutoff refinement) appear only in the \emph{unphysical} part of the boundary
modular energy and are minimized separately in Sec.~\ref{subsec:P3-master}.


\subsection{Receiver side: holographic completion and anomaly}
\label{subsec:receiver-anomaly}

\paragraph{Brown–York tensor and holographic renormalization.}\label{par:uom-main-038}
On the AdS$_4$ brane (receiver) embedded in AdS$_5$, the quasi–local (renormalized) Brown–York \cite{BrownYork1993}
stress tensor is
\begin{equation}
T^{\rm BY}_{\mu\nu}
=\frac{1}{8\pi G_5}\Big(K_{\mu\nu}-K g_{\mu\nu}\Big)
-\frac{1}{8\pi G_5}\!\left(-\frac{3}{\ell} g_{\mu\nu}
+\frac{\ell}{2}G_{\mu\nu}
+\cdots\right),
\label{eq:BY}
\end{equation}
where $K_{\mu\nu}$ is the extrinsic curvature of the brane, $\ell$ is the AdS$_5$ radius,
$G_{\mu\nu}$ is the 4D Einstein tensor of the induced metric $g_{\mu\nu}$, and the ellipsis
denotes higher–derivative local counterterms required in $d{=}4$. The renormalized CFT
stress on the brane equals the BY tensor minus the local pieces:
\begin{equation}
\langle T^{\rm CFT}_{\mu\nu}\rangle
= T^{\rm BY}_{\mu\nu}
-\frac{1}{8\pi G_5}\!\left(-\frac{3}{\ell} g_{\mu\nu}
+\frac{\ell}{2}G_{\mu\nu}
+\cdots\right).
\label{eq:CFT-vs-BY}
\end{equation}
At zero matter on the receiver, the effective 4D equations are
\begin{equation}
G_{\mu\nu}+\Lambda_4 g_{\mu\nu}
=8\pi G_4\,\langle T^{\rm CFT}_{\mu\nu}\rangle,
\qquad
G_4=\frac{G_5}{\ell},
\qquad
R=4\Lambda_4=-\frac{12}{L_4^2}.
\label{eq:4D-eff}
\end{equation}

\paragraph{Trace anomaly and the scalar (trace) channel.}\label{par:uom-main-039}
In $d{=}4$, holographic renormalization fixes the CFT trace anomaly
\begin{equation}
\langle T^\mu{}_\mu\rangle_{\rm CFT}
=\frac{c}{16\pi^2}\,W_{\mu\nu\rho\sigma}W^{\mu\nu\rho\sigma}
-\frac{a}{16\pi^2}\,E_4
+b\,\Box R,
\qquad
a=c=\frac{\pi \ell^3}{8G_5},
\label{eq:anomaly}
\end{equation} \cite{HenningsonSkenderis1998}.

with $W^2$ the Weyl invariant and $E_4$ the Euler density. Linearizing the traced version of
\eqref{eq:4D-eff} about a (nearly) maximally symmetric AdS$_4$ background ($W^2=0$, $E_4=\tfrac16R^2$)
and grouping local terms to the left yields
\begin{equation}
\gamma\,R^2\;-\;R\;+\;4\Lambda_4\;-\;8\pi G_4\,b\,\Box R
\;=\;8\pi G_4\,\tau,
\qquad
\gamma:=\frac{a\,G_4}{12\pi}>0,
\label{eq:trace-eq}
\end{equation}
where $\tau$ is the \emph{injected trace} density compiled on the receiver (the welded
image of the source drive). The \emph{DC} (constant–harmonic) projection removes the $\Box R$
term and gives the averaged relation
\begin{equation}
\gamma\,\langle R^2\rangle - \langle R\rangle + 4\Lambda_4
= 8\pi G_4\,\overline{\tau}.
\label{eq:DC-trace}
\end{equation}
Even when $\overline{\tau}=0$ (zero–area DoG), fluctuations bias the mean curvature via the
\emph{anomaly rectifier} $\gamma\,\langle R^2\rangle$; this is the \emph{BY anomaly floor}.

\subsubsection*{Weyl compensator and Stinespring uplift (where the BY anomaly comes from)}
\label{subsubsec:weyl-compensator-stinespring}

\paragraph{Why a compensator.}\label{par:uom-main-040}
The welded map ${\sf K}^{(\varepsilon)}_\Delta:\mathcal A_{\rm rad}(\mathscr I^+)\!\to\!\mathcal A_\partial$
is conformal at the continuum, but any finite cutoff $\Lambda$ on the receiver breaks exact
Weyl invariance \cite{HenningsonSkenderis1998,deHaroSkenderisSolodukhin2001,Skenderis2002HR}. To keep the export map \emph{unitary} before reduction (Stinespring minimality),
we enlarge the receiver code by a single Weyl (dilaton) mode $\sigma$ (the \emph{Weyl compensator})
so that boundary metrics are related by a Weyl lift $\tilde g_{\mu\nu}=e^{2\sigma} g_{\mu\nu}$.
The Stinespring isometry acts as
\[
V:\ \mathcal H_{\rm rad}\ \longrightarrow\ \mathcal H_{\rm code}\otimes\mathcal H_{\rm Weyl},
\qquad
\Phi\ \mapsto\ (\Phi\circ{\sf K}^{(\varepsilon)}_\Delta)\ \otimes\ \ket{\sigma},
\]
and tracing out $\mathcal H_{\rm Weyl}$ reproduces the renormalized CFT/BY response with its
trace anomaly.

\paragraph{Local compensator action and the anomaly.}\label{par:uom-main-041}
A local (holographically fixed) action for $\sigma$ that reproduces the $d{=}4$ anomaly is
the Riegert–Fradkin–Tseytlin form \cite{Riegert1984,FradkinTseytlin1984} linearized around an AdS$_4$ background:
\begin{equation}
S_{\rm comp}[\sigma;g]
=\frac{a}{24\pi}\int\! d^4x\sqrt{-g}\,
\Big(\tfrac12\,\sigma\,\Delta_4\,\sigma\ +\ \sigma\,\mathcal A[g]\Big)
\ +\ \frac{c}{16\pi^2}\int\! d^4x\sqrt{-g}\,\sigma\,W^2,
\label{eq:comp-action}
\end{equation}
where $\Delta_4:=\Box^2+2R^{\mu\nu}\nabla_\mu\nabla_\nu-\tfrac23 R\,\Box+\tfrac13(\nabla^\mu R)\nabla_\mu$
is the Paneitz operator and
$\mathcal A[g]:=E_4-\tfrac23\Box R$.
Varying $S_{\rm comp}$ under $\sigma\mapsto \sigma+\delta\sigma$ gives
$\delta S_{\rm comp}=\int\!\sqrt{-g}\,\delta\sigma\,\langle T^\mu{}_\mu\rangle$,
with
$\langle T^\mu{}_\mu\rangle=\tfrac{c}{16\pi^2}W^2-\tfrac{a}{16\pi^2}E_4+b\,\Box R$,
i.e.\ the usual $a=c$ holographic trace anomaly.

\paragraph{Coupling to the compiled trace and integrating out $\sigma$.}\label{par:uom-main-042}
The welded pulse sources the trace sector by a linear coupling
$S_{\rm src}=\int\!\sqrt{-g}\,\sigma\,\tau$, where $\tau$ is the compiled (receiver–side) trace.
The quadratic path integral over $\sigma$ is Gaussian:
\[
\int\!\mathcal D\sigma\ \exp\!\Big\{i\big(S_{\rm comp}[\sigma;g]+S_{\rm src}[\sigma]\big)\Big\}
\;=\;\exp\!\Big\{\frac{i}{2}\int\!\sqrt{-g}\,\tau\,\Delta_4^{-1}\tau\ +\ \cdots\Big\}.
\]
The ellipsis encodes local counterterms already packaged into the Brown–York tensor.
Linear response of the scalar curvature then reads, in Fourier space,
\begin{equation}
\widehat{R}(\omega,k)
=-\,8\pi G_4\,\mathcal K(\omega,k;\Lambda)\,\widehat{\tau}(\omega,k),
\qquad
\mathcal K(\omega,k;\Lambda)\ \simeq\ \frac{q(\omega,k)}{q(\omega,k)+r(\Lambda)k^2},
\label{eq:comp-kernel}
\end{equation}
with $q(\omega,k)\!\approx\!\sqrt{k^2+m_g^2}$ ($m_g\!\sim\!L_4^{-1}$) and
$r(\Lambda)=M_4^2(\Lambda)/(2M_5^3)$ from the induced 4D Einstein term.
This is precisely the screened/Yukawa kernel used in Sec.~\ref{subsec:receiver-anomaly}.

\paragraph{BY anomaly floor as compensator quadratic energy.}\label{par:uom-main-043}
The tick–wise curvature variance injected by a welded LoG pulse is
\(
\Delta\!\langle R^2\rangle
=\!\int \!|8\pi G_4\,\mathcal K|^2\,|\widehat{\tau}|^2,
\)
and the \emph{anomaly rectifier}
\(
\Theta=\gamma\,\Delta\!\langle R^2\rangle
\)
with $\gamma\!=\!aG_4/(12\pi)$ is exactly the quadratic (positive) form obtained by integrating out
$\sigma$. Thus the BY term in the receiver P3 functional,
$\lambda_{\rm BY}\Theta$, is the compensator’s quadratic modular cost.
Convexity follows from the positivity of $\Delta_4$ on the AdS$_4$ spectrum sector sampled
by the welded LoG family.

\paragraph{Unitarity before reduction (Stinespring) and P3.}\label{par:uom-main-044}
Keeping $\sigma$ makes the welded export an \emph{isometry} into an extended code
$\mathcal H_{\rm code}\!\otimes\!\mathcal H_{\rm Weyl}$; P3 minimizes the \emph{net}
(modular) action of transport. Tracing out $\sigma$ yields the Brown–York/CFT response with the
trace anomaly; the three receiver penalties (conformality, RP monitor, anomaly/BY) are the
orthogonal projections of the same compensator-completed quadratic form onto, respectively,
the Weyl, OS/RP, and trace sectors of the compiled state.

\subsubsection*{Gravitational receiver-source coupling}
\label{subsubsec:gravitational-receiver-source-coupling}
\paragraph{Induced gravity and cutoff dependence.}\label{par:uom-main-045}
Integrating the receiver CFT down to a UV cutoff $\Lambda$ induces a local 4D Einstein–Hilbert
term with scale
\begin{equation}
M_{4}^2(\Lambda)=M_5^3\,\ell+\kappa\,a\,\Lambda^2,
\qquad
r(\Lambda):=\frac{M_4^2(\Lambda)}{2M_5^3},
\label{eq:induced-grav}
\end{equation}
($\kappa=\mathcal O(1)$ scheme constant). In momentum space, the retarded trace–response kernel
is regular at $(\omega,k)\!=\!(0,0)$ and, for long wavelengths, takes the Yukawa/screened form
\begin{equation}
\mathcal K(\omega,k;\Lambda)
\simeq \frac{q(\omega,k)}{q(\omega,k)+r(\Lambda)\,k^2},
\qquad
q(\omega,k)\approx \sqrt{k^2+m_g^2},\quad m_g\sim L_4^{-1},
\label{eq:K-screen}
\end{equation}
multiplying the geometric AdS$_5$ overlap.

\subsubsection*{\texorpdfstring{$AdS_4$ curvature update on uplift}{AdS4 curvature update on uplift}}
\label{subsubsec:ads4-curvature-change-uplift}
\paragraph{Welded pulse $\rightarrow$ curvature variance.}\label{par:uom-main-046}
Let $\widehat{\tau}_k(\omega,\mathbf k)$ be the compiled (welded) trace source in tick $k$
induced by the micro/macro LoG envelope (temporal DoG of width $\sigma_{t,k}$ and spatial
isotropic Gaussian/DoG of width $\sigma_{x,k}$). In linear response,
\begin{equation}
\widehat{R}_k(\omega,\mathbf k)
=-\,8\pi G_4\,\mathcal K(\omega,k;\bar\Lambda_k)\,\widehat{\tau}_k(\omega,\mathbf k),
\label{eq:R-from-tau}
\end{equation}
and the tick–wise curvature \emph{variance injection} is
\begin{equation}
\Delta\!\langle R^2\rangle_k
=\int\!\frac{d\omega\,d^3k}{(2\pi)^4}\,
\big|8\pi G_4\,\mathcal K(\omega,k;\bar\Lambda_k)\big|^2\,\big|\widehat{\tau}_k(\omega,\mathbf k)\big|^2.
\label{eq:var-inj}
\end{equation}
For long wavelengths and a unit–shape DoG/LoG family, \eqref{eq:var-inj} reduces to the closed
scaling (the $\mathcal O(1)$ constant $C_{\rm LoG}$ fixes the convention)
\begin{equation}
\Delta\!\langle R^2\rangle_k
=(8\pi G_4)^2\,\big\langle|\mathcal K|^2\big\rangle_{\!k}\;
C_{\rm LoG}\;\frac{A_k^2}{\sigma_{t,k}\,\sigma_{x,k}^{3}},
\qquad
\big\langle|\mathcal K|^2\big\rangle_{\!k}
:=\frac{\int |\mathcal K|^2\,|\widehat{\tau}_k|^2}{\int |\widehat{\tau}_k|^2}.
\label{eq:variance-scaling}
\end{equation}
Here $A_k$ is the tick’s (temporal) amplitude, and the spectral average encodes the \emph{cutoff}
and \emph{induced–gravity} dependence through $\bar\Lambda_k$ via \eqref{eq:K-screen}.

\subsubsection*{\texorpdfstring{Cause of the $AdS_4 \rightarrow dS_4$ flip: what drives the cosmological constant $\Lambda_4$ across $0$}{Cause of the AdS4->dS4 flip: what drives the cosmological constant Lambda_4 across 0}}
\label{subsubsec:flip-cause}
\paragraph{Anomaly floor per tick and its role.}\label{par:uom-main-047}
Multiplying \eqref{eq:variance-scaling} by $\gamma$ gives the tick–wise \emph{BY anomaly floor}
\begin{equation}
\Theta_k\ :=\ \gamma\,\Delta\!\langle R^2\rangle_k
=\underbrace{\frac{a\,G_4}{12\pi}}_{\gamma}\,(8\pi G_4)^2\,
\big\langle|\mathcal K|^2\big\rangle_{\!k}\;
C_{\rm LoG}\;\frac{A_k^2}{\sigma_{t,k}\,\sigma_{x,k}^{3}},
\label{eq:theta-def}
\end{equation}
which enters the DC trace balance \eqref{eq:DC-trace} even for zero–DC pulses ($\overline{\tau}=0$),
biasing $\langle R\rangle$ upward by a non–negative amount. Across ticks,
\[
\langle R\rangle_{k+1}\ \approx\ (1-\alpha_k)\,\langle R\rangle_k\ +\ \Theta_k,
\]
where $(1-\alpha_k)$ is the (possibly tiny) multiplicative contraction from any DC
channel; in our zero–area DoG baseline $\alpha_k\!=\!0$, so pre–flip resolution is governed
by $\Theta_k$.

\subsubsection*{Receiver-side Universe realization transfer action}\label{subsubsec:uom-main-004}
\begin{definition}[Ward--residual conformality monitor]
\label{def:ward-residual}
Let $C_\Lambda$ be the receiver covariance at cutoff $\Lambda$, and let
$\{{\cal O}_i\}_{i\in I}$ denote a finite generating family of \emph{conformal Ward operators}
(trace and diffeomorphism Ward operators) acting linearly on covariances.
Define the \emph{Ward--residual} at cutoff $\Lambda$ by
\begin{equation}
\label{eq:ward-residual}
D_{\rm conf}(\Lambda)\;:=\;\sum_{i\in I}\Big\|{\cal O}_i\,[\,C_\Lambda\,]\Big\|_{F}^{2},
\end{equation}
where $\|\cdot\|_F$ is the Fisher (BKM) quadratic norm on covariance tangents.
For a tick $k$ we write $D_{{\rm conf},k}:=D_{\rm conf}(\Lambda_k)$.
\end{definition}

\paragraph{Receiver–side contribution to the P3 action.}\label{par:uom-main-048}
The receiver’s \emph{unphysical} modular part (what P3 minimizes over tick controls) includes
the anomaly term and the two “compatibility” terms (conformality and OS/RP) as convex functionals
of the welded envelope and cutoff,
\begin{equation}
\label{eq:Jrec}
\begin{aligned}
\mathcal J^{\rm rec}_k(\sigma_{t,k},\sigma_{x,k},A_k,\Lambda_k)
&= \underbrace{\lambda_{\rm BY}\,\Theta_k}_{\text{BY anomaly rectifier}}
\\[2pt]
&\quad + \underbrace{\lambda_{\rm conf}\,D_{{\rm conf},k}}_{\text{Ward--residual conformality}}
\\[2pt]
&\quad + \underbrace{\lambda_{\rm RP}\,\Delta D_{{\rm RP},k}}_{\text{RP deficit monitor}}\,.
\end{aligned}
\end{equation}

with fixed positive weights $\lambda$ determined by Fisher metric of the modular energy master functional decomposition, naturally expressed in the canonical “nats” convention (see Sec.~\ref{subsec:P3-master}). 

It can be shown that the reflection-odd term (OS/RP monitor) vanishes in our receiver (Theorem \ref{thm:OS_receiver}) under the locked welded pulse shape. We will, therefore, omit this term in some further analysis of P3 action functional and its generalization.

Under the LoG family,
$\mathcal J^{\rm rec}_k$ is jointly convex in  $(\sigma_{t,k}^{-1},\,\sigma_{x,k}^{-3},\,\Lambda_k)$ so the welded pulse widths $s_t,s_x$ are locked to band at $\Lambda_k$ (Sec.~\ref{sec:aeon-evolution}).

\paragraph{Cutoff update and causality.}\label{par:uom-main-049}
The receiver cutoff $\Lambda_k$ is a state variable advanced by the SK/HJ wall balance under the
chosen admissible controls (window/passband/shape). Its value at tick boundaries enters retardedly
into the next tick's receiver-side knee (hence the source-side effective bandlimit $\Delta_{k+1}$)
through the bulk transfer kernel.

\medskip
\noindent
Equations \eqref{eq:trace-eq}–\eqref{eq:theta-def} are the only receiver–side dynamical inputs
the optimizer needs: they map the welded LoG pulse parameters and current cutoff to the anomaly
floor $\Theta_k$ and, through \eqref{eq:Jrec}, to the receiver contribution in the P3 action.

\subsection{The unified P3 action (master functional) and its decomposition}
\label{subsec:P3-master}

\paragraph{Master statement (per tick).}\label{par:uom-main-050}
P3 asserts that, for each accepted tick $k$, Nature chooses the welded control
$h_k$ (temporal micro DoG, receiver macro time/space envelopes) and the
receiver refinement $\Lambda_k$ to \emph{minimize the two–sided modular action}
\begin{equation}
\boxed{\;
\mathcal A_k \;=\;
\underbrace{F_{\rm src}\!\big[\rho_k\,\big\|\,\rho^{\rm BD}_{\Delta_k}\big]
\;-\;C_k}_{\text{source modular dwell – classical record}}
\;+\;
\underbrace{\Phi_{\rm rec}[h_k;\Lambda_k]}_{\text{receiver transport penalty}}
\;}\,,
\label{eq:Ak-master}
\end{equation}
subject to the admissibility guards (QSL, acceptance, retarded coupling) defined
in Secs.~\ref{subsec:scene-notation}–\ref{subsubsec:P2}.
Here $F_{\rm src}$ is the source–side modular free energy relative to the
band–limited KMS anchor $\rho^{\rm BD}_{\Delta_k}$, $C_k$ is the compiled
(classical) record for the tick, and $\Phi_{\rm rec}$ is the convex receiver
penalty induced by the finite–$\Lambda$ holographic completion.

\paragraph{Definition vs.\ derivation (reconciliation).}\label{par:uom-main-051}
We adopt the following two-step logic. First, we \emph{define} the three receiver monitors
$\Pi_{\rm RP}$, $\Pi_{\rm Weyl}$, $\Pi_{\rm conf}$ in the Fisher metric (RP reflection, Riegert/Weyl,
and conformal–orbit Gram–Schmidt) and write the P3 receiver penalty in terms of their quadratic
norms (geometric route). 

Second, we \emph{derive} that, for the welded Gaussian uplift at fixed cutoff~$\Lambda$,
the quadratic modular energy coincides with this monitor sum with \emph{geometry-fixed weights},
and no further symmetry-invariant quadratic term exists (brane physics route).



\subsubsection*{Geomeric route (uplift-created receiver subspaces}\label{subsubsec:uom-main-005}

\paragraph{Receiver uplift modular energy decomposition.}
\label{para:uplift-decomp}
Let $\rho_\ast$ be the (holographically renormalized) AdS$_4$ fixed–point boundary state and
$\rho_\Lambda$ the realized boundary state at receiver cutoff $\Lambda$ produced by a welded
trace $\tau$ in tick $k$. In the Gaussian radiative sector the state is specified by a covariance
$C_\Lambda$, with fixed–point $C_\ast$, and the receiver–side modular distance is
\begin{equation}
\label{eq:master-receiver}
\mathcal J_{\rm rec}(\Lambda;\tau)\ :=\ S(\rho_\Lambda\|\rho_\ast)
\;=\;\tfrac12\,\mathrm{tr}\!\Big(C_\ast^{-1}C_\Lambda-\mathbf 1-\ln(C_\ast^{-1}C_\Lambda)\Big).
\end{equation}
Let $\delta C:=C_\Lambda-C_\ast$ and denote by $\mathsf T_\Lambda:\tau\mapsto\delta C$ the welded
linear response map (LoG control $\to$ boundary two–point variation). The \emph{tangent space} at
$C_\ast$ carries the Fisher metric
\begin{equation}
\label{eq:fisher}
\big\langle\!\big\langle \delta C,\delta C'\big\rangle\!\big\rangle_{C_\ast}
\;:=\;\tfrac12\,\mathrm{tr}\!\big(C_\ast^{-1}\delta C\,C_\ast^{-1}\delta C'\big),
\end{equation}
so that
$S(\rho_\Lambda\|\rho_\ast)=\tfrac12\langle\!\langle \delta C,\delta C\rangle\!\rangle_{C_\ast}+O(\|\delta C\|^3)$.

\medskip
\noindent\textbf{Symmetries and commuting projectors.}
Three commuting receiverside symmetries act on the covariance data:
(i) local Weyl rescalings (implemented by the Wess–Zumino/Riegert compensator),
(ii) Euclidean time reflection $\Theta$ (Osterwalder–Schrader), and
(iii) global conformal motions $g\in\mathrm{Conf}(S^3)$ acting unitarily on the receiver radiative code sector (boundary action $U_\partial(g)$). Averaging over these symmetries defines three
idempotent, self–adjoint, commuting projectors on the tangent space:
\begin{align}
\label{eq:projectors}
\Pi_{\rm Weyl}&:=\text{Weyl average (Riegert projector)},&
\Pi_{\rm RP}&:=\tfrac12(\mathrm{Id}-\mathrm{Ad}\,\Theta),&
\Pi_{\rm conf}&:=\text{conformal–code projector}.
\end{align}
(Explicitly, $\Pi_{\rm RP}$ selects the $\Theta$–odd component; $\Pi_{\rm conf}$ is the orthogonal
projection onto the radiative code subspace and its exact conformal orbit; $\Pi_{\rm Weyl}$ projects
onto the local–scale tangent generated by the anomaly functional.)

\begin{lemma}[Orthogonality in the Fisher metric]\label{lem:tan-orth}
Let $\rho_\ast$ be the (holographically renormalized) AdS$_4$ fixed–point boundary state in the
Gaussian radiative sector with covariance $C_\ast$, and equip the tangent space at $C_\ast$ with the
Fisher inner product
\[
\big\langle\!\big\langle \delta C,\delta C'\big\rangle\!\big\rangle_{C_\ast}
:= \tfrac12\,\mathrm{tr}\!\big(C_\ast^{-1}\delta C\,C_\ast^{-1}\delta C'\big).
\]
Let $\Pi_{\rm RP}$, $\Pi_{\rm Weyl}$, and $\Pi_{\rm conf}$ be the following three orthogonal
projections on the tangent space:
\begin{enumerate}
\item $\Pi_{\rm RP}:=\tfrac12(\mathrm{Id}-\mathrm{Ad}\,\Theta)$, where $\Theta$ is Euclidean
time reflection; $\mathrm{Ad}\,\Theta$ acts by conjugation on covariances and is an isometry
of the Fisher metric.
\item $\Pi_{\rm Weyl}:=\mathsf D_{\rm W}\,(\mathsf D_{\rm W}^\dagger\mathsf D_{\rm W})^{-1}\,\mathsf D_{\rm W}^\dagger$,
the orthogonal projector onto the range of the linear Weyl–variation map
$\mathsf D_{\rm W}:\sigma\mapsto\delta C_{\rm Weyl}(\sigma)$ (generated by the Riegert/Wess–Zumino
compensator); $\mathsf D_{\rm W}^\dagger$ is the Fisher adjoint.
\item $\Pi_{\rm conf}$ the orthogonal projector onto the closed subspace
$\mathcal T_{\rm conf}$ spanned by conformal–orbit tangents on the radiative code \emph{and}
Fisher–orthogonal to $\mathrm{ran}\,\Pi_{\rm Weyl}$ (i.e.\ obtained by Gram–Schmidt of the conformal
tangent against the Weyl tangent).
\end{enumerate}
Then the ranges $\mathcal T_{\rm RP}:=\mathrm{ran}\,\Pi_{\rm RP}$,
$\mathcal T_{\rm Weyl}:=\mathrm{ran}\,\Pi_{\rm Weyl}$, and
$\mathcal T_{\rm conf}:=\mathrm{ran}\,\Pi_{\rm conf}$ are pairwise orthogonal with respect
to $\langle\!\langle\cdot,\cdot\rangle\!\rangle_{C_\ast}$.
\end{lemma}

\begin{proof}
\emph{Step 1 (Self–adjointness and idempotence).} By construction,
$\Pi_{\rm RP}$ is a polynomial in the isometry $\mathrm{Ad}\,\Theta$ and hence is self–adjoint
and idempotent in the Fisher metric. The map $\Pi_{\rm Weyl}$ has the standard form of the orthogonal
projector onto $\mathrm{ran}\,\mathsf D_{\rm W}$ (Moore–Penrose projector) and is self–adjoint
and idempotent. Finally, $\Pi_{\rm conf}$ is the orthogonal projector (by definition) onto the closed
subspace $\mathcal T_{\rm conf}$; thus it is self–adjoint and idempotent.

\emph{Step 2 (RP–even vs.\ RP–odd).} The time–reflection $\Theta$ leaves the background
$C_\ast$ invariant, hence $\mathrm{Ad}\,\Theta$ is an isometry of the Fisher metric. Therefore the
$\pm1$ eigenspaces (even/odd) of $\mathrm{Ad}\,\Theta$ are orthogonal. By definition
$\mathcal T_{\rm RP}$ is the odd eigenspace. Both Weyl variations and conformal
motions are \emph{even} under $\Theta$ (they do not change sign under time reflection),
so $\mathrm{ran}\,\Pi_{\rm Weyl}$ and $\mathrm{ran}\,\Pi_{\rm conf}$ lie in the even eigenspace.
Hence
$\langle\!\langle \mathcal T_{\rm RP},\mathcal T_{\rm Weyl}\rangle\!\rangle
=\langle\!\langle \mathcal T_{\rm RP},\mathcal T_{\rm conf}\rangle\!\rangle=0$.

\emph{Step 3 (Weyl vs.\ conformal).} By construction, $\mathcal T_{\rm conf}$ is the orthogonal
projection of the conformal tangent onto the orthogonal complement of $\mathrm{ran}\,\Pi_{\rm Weyl}$.
Equivalently, $\Pi_{\rm conf}$ is the projector onto (a closed subspace of)
$\ker \Pi_{\rm Weyl}$ restricted to the code/conformal sector. Thus
$\langle\!\langle \mathcal T_{\rm Weyl},\mathcal T_{\rm conf}\rangle\!\rangle=0$.

Combining Steps 2–3 gives pairwise orthogonality of the three ranges.
\end{proof}


\begin{lemma}[Decomposition and physical kernel]\label{lem:tan-complete}
Let $\mathscr T:=\mathrm{ran}\,\mathsf T_\Lambda$ be the welded tangent image at cutoff $\Lambda$,
where $\mathsf T_\Lambda$ maps welded controls (traces) to boundary covariance variations
$\delta C$. With the projectors $\Pi_{\rm RP}$, $\Pi_{\rm Weyl}$, $\Pi_{\rm conf}$ of
Lemma~\ref{lem:tan-orth}, define the \emph{physical kernel}
\[
\mathcal T_{\rm phys}\ :=\ \ker \Pi_{\rm RP}\ \cap\ \ker \Pi_{\rm Weyl}\ \cap\ \ker \Pi_{\rm conf}.
\]
Then every $\delta C\in\mathscr T$ admits a unique orthogonal decomposition
\begin{equation}\label{eq:orth-sum}
\delta C\;=\;\delta C_{\rm RP}\ \oplus\ \delta C_{\rm Weyl}\ \oplus\ \delta C_{\rm conf}\ \oplus\ \delta C_{\rm phys},
\end{equation}
with $\delta C_i\in\mathcal T_i$ ($i\in\{\mathrm{RP},\mathrm{Weyl},\mathrm{conf}\}$) and
$\delta C_{\rm phys}\in\mathcal T_{\rm phys}$. Equivalently,
\[
\mathscr T\;=\;\mathcal T_{\rm RP}\ \widehat\oplus\ \mathcal T_{\rm Weyl}\ \widehat\oplus\ \mathcal T_{\rm conf}
\ \oplus\ \mathcal T_{\rm phys}.
\]
\end{lemma}

\begin{proof}
Since $\Pi_{\rm RP}$, $\Pi_{\rm Weyl}$, and $\Pi_{\rm conf}$ are self–adjoint idempotents on the
Hilbert space $(\mathscr T,\langle\!\langle\cdot,\cdot\rangle\!\rangle_{C_\ast})$, their ranges
$\mathcal T_{\rm RP}$, $\mathcal T_{\rm Weyl}$, $\mathcal T_{\rm conf}$ are closed subspaces and,
by Lemma~\ref{lem:tan-orth}, are pairwise orthogonal.

For any $\delta C\in\mathscr T$, set
\[
\delta C_{\rm RP}:=\Pi_{\rm RP}\delta C,\qquad
\delta C_{\rm Weyl}:=\Pi_{\rm Weyl}\delta C,\qquad
\delta C_{\rm conf}:=\Pi_{\rm conf}\delta C,\qquad
\delta C_{\rm phys}:=(\mathrm{Id}-\Pi_{\rm RP}-\Pi_{\rm Weyl}-\Pi_{\rm conf})\delta C.
\]
Each of the first three terms lies in the corresponding range by definition.
Because the projectors are self–adjoint, for any $X$ in one range and $Y$ in a distinct range,
$\langle\!\langle X,Y\rangle\!\rangle=0$ (Lemma~\ref{lem:tan-orth}). Moreover,
$\delta C_{\rm phys}$ lies in the joint kernel since
$\Pi_i\,\delta C_{\rm phys}=\Pi_i\delta C-\Pi_i^2\delta C-\Pi_i\Pi_j\delta C-\Pi_i\Pi_k\delta C=0$
for $\{i,j,k\}=\{\mathrm{RP},\mathrm{Weyl},\mathrm{conf}\}$ (idempotence and commutation on
the defined domains suffice here; orthogonality yields the vanishing of pairwise mixed terms
in the sense of projections onto orthogonal ranges). Therefore
$\delta C_{\rm phys}\in\mathcal T_{\rm phys}$.

Summing the four components gives $\delta C$ by construction, so the decomposition
\eqref{eq:orth-sum} exists. Uniqueness follows from orthogonality: if
$\sum_i X_i + X_{\rm phys}=0$ with $X_i\in\mathcal T_i$ and
$X_{\rm phys}\in\mathcal T_{\rm phys}$, taking the Fisher inner product with $X_j$
and using orthogonality forces $X_j=0$ for each $j$; similarly with any
$Y\in\mathcal T_{\rm phys}$ shows $X_{\rm phys}=0$.
Hence
$\mathscr T=\mathcal T_{\rm RP}\widehat\oplus\mathcal T_{\rm Weyl}\widehat\oplus
\mathcal T_{\rm conf}\oplus\mathcal T_{\rm phys}$ as claimed.
\end{proof}


\noindent\textbf{Quadratic modular decomposition (weights fixed by geometry).}
With $\delta C=\mathsf T_\Lambda\tau$ and \eqref{eq:fisher},
\begin{equation}
\label{eq:quad-sum}
S(\rho_\Lambda\|\rho_\ast)\ =\ \tfrac12\sum_{i\in\{\mathrm{Weyl,RP,conf}\}}
\big\langle\!\big\langle \Pi_i\delta C,\Pi_i\delta C\big\rangle\!\big\rangle_{C_\ast}
\;+\;O(\|\delta C\|^3),
\end{equation}
so the three quadratic coefficients are \emph{not knobs}: the Weyl weight is fixed by the anomaly
($a=c\propto \ell^3/G_5$) and $G_4$ via the Riegert operator (yielding the familiar
$\gamma\,\Delta\!\langle R^2\rangle$), while the RP and conformal weights are fixed by the Fisher
metric pulled back to the OS cone and the code intertwiner, respectively.


\subsubsection*{\texorpdfstring{Physical route (subspaces on the receiver AdS$_4$ brane)}{Physical route (subspaces on the receiver AdS4 brane)}}\label{subsubsec:uom-main-006}
Next, we analyze the orthogonality error of our conformality mismatch projector construction.

\paragraph{Fisher–orthogonal monitor projectors.}\label{par:uom-main-052}
Let $C_\ast$ be the receiver fixed point and equip the tangent with the Fisher inner product.
Define the following closed subspaces:
\begin{itemize}
\item $\mathcal T_{\rm RP}$: the \emph{reflection–odd} tangent (image of the time–reflection
projector $\,\tfrac12(\mathrm{Id}-\mathrm{Ad}\,\Theta)$).
\item $\mathcal T_{\rm Weyl}$: the \emph{Weyl} tangent (image of the linearized trace map generated by the compensator).
\item $\mathcal T_{\rm conf}$: the closure of the image of the \emph{Ward–residual} map
$C\mapsto ({\cal O}_i[C])_{i}$, orthogonally projected to $\mathcal T_{\rm Weyl}$ in the Fisher metric.
\end{itemize}
Let $\Pi_{\rm RP}$, $\Pi_{\rm Weyl}$, $\Pi_{\rm conf}$ denote the corresponding Fisher–orthogonal projectors.
Then
\[
\mathscr T\;=\;\mathcal T_{\rm RP}\ \widehat\oplus\ \mathcal T_{\rm Weyl}\ \widehat\oplus\ \mathcal T_{\rm conf}\ \widehat\oplus\ \mathcal T_{\rm phys},
\]
with $\mathcal T_{\rm phys}$ the joint kernel of the three projectors.

\begin{proposition}[KR scalar mixing scale and interior RS1 continuity]
\label{prop:scalar-mix}
Let $q:=\|\mathbf q\|\ge 0$ be the 3–momentum. Consider the \emph{scalar} sector on the active KR head
receiver, spanned by the Weyl/trace source $\sigma$ and a physical scalar mode $r$ (the
boundary–distance scalar). For comparison, when an interior RS1 segment abuts a KR head,
the analogous sector is spanned by $\sigma$ and the stabilized radion $r$. Then:
\begin{enumerate}[leftmargin=2.0em,label=(\roman*)]

\item \textbf{KR head (active receiver, induced–gravity Robin BC).}
A normalizable scalar mode $r$ \emph{persists} (the $E_0=4$ AdS$_4$ boundary–distance scalar) with
$m_r>0$ and lies in the \emph{physical} Fisher tangent. Its leakage into the conformality–mismatch
block arises only via a small Weyl–radion mixing controlled by the visible–brane Robin overlap
$\varepsilon_{\rm BC}\equiv \|u_r(z_{\rm vis})\|/\|u_r(z_{\rm sh})\|\ll 1$ (stiff cap).
The quadratic form on $(\sigma,r)$ has the structure
\begin{equation}
\label{eq:KR-ABC-scaling}
\mathcal Q(q;\sigma,r)\;=\;
A(q)\,|\sigma|^2\;+\;2\,B(q)\,\Re\big(\sigma\,\overline{r}\big)\;+\;C(q)\,|r|^2,
\qquad
\begin{cases}
A(q)=a\,q^{4}+o(q^{4}),\\[2pt]
B(q)=\tilde b\,\varepsilon_{\rm BC}\,q^{2}+o(q^{2}),\\[2pt]
C(q)=c\,(q^{2}+m_{r}^{2})+o(q^{2}+m_{r}^{2}),
\end{cases}
\end{equation}
with $a,\tilde b,c>0$ independent of $q$ and $\varepsilon_{\rm BC}\ll 1$.\footnote{In the \emph{strict} EOW
(KR with non–dynamical boundary metric) one may take $s_{\rm vis}\to\infty$, hence $\varepsilon_{\rm BC}\to 0$,
and the scalar’s projection on the conformality–mismatch block vanishes \emph{identically} at quadratic order.
Here we work in the induced–gravity/Robin regime, so $\varepsilon_{\rm BC}$ is small but finite.}
Let $\theta(q)$ be the orthogonal mixing angle that diagonalizes $\mathcal Q$.
Then
\begin{equation}
\label{eq:KR-theta-asymp}
\theta(q)\;=\;
\begin{cases}
-\dfrac{\tilde b}{c}\,\varepsilon_{\rm BC}\,\dfrac{q^{2}}{m_{r}^{2}}\ +\ o\!\big(q^{2}\big), & q\downarrow 0,\\[8pt]
\;\dfrac{\tilde b}{a}\,\varepsilon_{\rm BC}\,\dfrac{1}{q^{2}}\ +\ o\!\big(q^{-2}\big), & q\to\infty.
\end{cases}
\end{equation}
In particular, at welded scales $q\sim\Lambda\gg m_r$ one has
$\theta(\Lambda)=O\!\big(\varepsilon_{\rm BC}\,\Lambda^{-2}\big)\ll 1$, so the
Fisher non–orthogonality of the Weyl and conformality–mismatch directions scales as
$O\!\big(\varepsilon_{\rm BC}^{2}\Lambda^{-4}\big)$.

\item \textbf{Interior RS1 segment (two branes, stabilized radion; not an active receiver).}
With a stabilized radion of mass $m_r>0$, the sector $(\sigma,r)$ has a positive–definite quadratic form
\begin{equation}
\label{eq:scalar-kernel}
\mathcal Q(q;\sigma,r)\;=\;
A(q)\,|\sigma|^2\;+\;2\,B(q)\,\Re\big(\sigma\,\overline{r}\big)\;+\;C(q)\,|r|^2,
\end{equation}
with derivative scalings
\begin{equation}
\label{eq:ABC-scaling}
A(q)=a\,q^{4}+o(q^{4}),\qquad
B(q)=b\,q^{2}+o(q^{2}),\qquad
C(q)=c\,(q^{2}+m_{r}^{2})+o(q^{2}+m_{r}^{2}),
\end{equation}
for some $a,b,c>0$. The diagonalizing angle $\theta(q)\in(-\frac\pi4,\frac\pi4)$ satisfies
\begin{equation}
\label{eq:theta-asymp}
\theta(q)\;=\;
\begin{cases}
-\dfrac{b}{c}\,\dfrac{q^{2}}{m_{r}^{2}}\ +\ o\!\big(q^{2}\big), & q\downarrow 0,\\[8pt]
\;\dfrac{b}{a}\,\dfrac{1}{q^{2}}\ +\ o\!\big(q^{-2}\big), & q\to\infty.
\end{cases}
\end{equation}
Hence for welded scales $q\sim\Lambda\gg m_r$ one has $\theta(\Lambda)=O(\Lambda^{-2})\ll 1$ and the
Weyl and conformality–mismatch directions are Fisher–orthogonal up to $O(\Lambda^{-2})$.
\end{enumerate}
\end{proposition}

\begin{proof}
\textbf{(i) KR.}
Work in the induced–gravity (Robin) regime at the KR head. At quadratic order the scalar
tangent splits into the Weyl/trace source and the (physical) boundary–distance mode $r$ with mass $m_r>0$.
The Fisher quadratic reduces to \eqref{eq:KR-ABC-scaling}, where $A\sim a q^4$ is the universal Weyl weight,
$C\sim c(q^2+m_r^2)$ is the canonical scalar weight, and the off–diagonal kernel $B$ is proportional to the
visible–brane Robin overlap $\varepsilon_{\rm BC}\ll1$. Writing
\[
\mathbf Q(q)\;=\;\begin{pmatrix} A(q) & B(q) \\[2pt] B(q) & C(q) \end{pmatrix},\qquad
\tan(2\theta)=\frac{2B}{A-C},
\]
and using the scalings in \eqref{eq:KR-ABC-scaling}, one finds:
for $q\downarrow 0$,
$A-C=-c\,m_r^2+O(q^2)$ and $B=\tilde b\,\varepsilon_{\rm BC}\,q^2+o(q^2)$, hence
$\tan(2\theta)=-\,\tfrac{2\tilde b}{c}\,\varepsilon_{\rm BC}\,\tfrac{q^2}{m_r^2}+o(q^2)$ and
$\theta(q)=-\,\tfrac{\tilde b}{c}\,\varepsilon_{\rm BC}\,\tfrac{q^2}{m_r^2}+o(q^2)$.
For $q\to\infty$, $A-C=a q^4\big(1-\tfrac{c}{a}q^{-2}+o(q^{-2})\big)$ and $B=\tilde b\,\varepsilon_{\rm BC}\,q^2+o(q^2)$,
so $\tan(2\theta)=\tfrac{2\tilde b}{a}\,\varepsilon_{\rm BC}\,q^{-2}\big(1+o(1)\big)$ and
$\theta(q)=\tfrac{\tilde b}{a}\,\varepsilon_{\rm BC}\,q^{-2}+o(q^{-2})$. This is \eqref{eq:KR-theta-asymp}.
In particular, at welded band edge $q\sim \Lambda\gg m_r$ one has $\theta(\Lambda)=O(\varepsilon_{\rm BC}\Lambda^{-2})$,
hence the conformality–mismatch penalty (a squared Fisher projection) scales as
$O(\varepsilon_{\rm BC}^2\Lambda^{-4})$.

\medskip
\noindent
\textbf{(ii) Interior RS1 segment.}
Identical to the original proof: the symmetric kernel $\mathbf Q(q)$ with the RS1 scalings \eqref{eq:ABC-scaling}
yields $\tan(2\theta)=2B/(A-C)$ and the IR/UV asymptotics \eqref{eq:theta-asymp} by the same small–argument
expansions. In particular, $\theta(\Lambda)=O(\Lambda^{-2})$ at welded scales, so Fisher orthogonality holds
up to $O(\Lambda^{-2})$.
\end{proof}


\paragraph{OS-positivity on the KR head and interior RS1 branes.}\label{par:uom-main-053}
We show that the reflection--odd term vanishes on the active KR head receiver and, separately, on
interior/shadow RS1 branes under the stated healthiness assumptions.
% ======================================================================
% OS-positivity for KR receivers and interior RS1 branes — theorem + proof
% ======================================================================

\newcommand{\R}{\mathbb{R}}
\newcommand{\Z}{\mathbb{Z}}
\newcommand{\veps}{\varepsilon}

\begin{theorem}[Osterwalder--Schrader positivity for KR receivers and interior RS1 branes]
\label{thm:OS_receiver}
Let $G_E(x,x')=\langle \mathcal O(x)\,\mathcal O(x')\rangle$ be the Euclidean two–point
Schwinger function of a \emph{brane–localized, gauge–invariant, linear} observable
$\mathcal O$ on a brane with Euclidean coordinates $x=(\tau,\vec x)\in\R\times\R^3$.
Assume one of the following bulk/brane geometries:
\begin{enumerate}
\item[\textnormal{(KR)}] A single AdS$_4$ Karch--Randall (KR) brane embedded in AdS$_5$
(the active receiver in our construction).
\item[\textnormal{(RS1-int)}] An RS1 stack with a UV brane and interior/shadow branes;
OS positivity holds for brane--localized observables on any interior brane under the
healthiness assumptions below. (RS1 is not an active receiver in the present picture.)
\end{enumerate}
Suppose further that the following \emph{healthiness hypotheses} hold:
\begin{itemize}
\item[\textnormal{(H1)}] The linearized bulk dynamics around the background and all brane–localized
quadratic terms have the \emph{correct sign} (no ghosts); the Wick rotation is well–defined.
\item[\textnormal{(H2)}] Junction/boundary conditions render the radial Sturm–Liouville operator
self–adjoint with a positive measure; all 4D Kaluza–Klein masses obey $m_n^2\ge 0$ (no tachyons).
\item[\textnormal{(H3)}] In RS1 the radion is stabilized and has a \emph{positive} 4D kinetic term.
\item[\textnormal{(H4)}] Standard covariant holographic counterterms are used so that the resulting
Euclidean two–point function is tempered and reflection–symmetric, $G_E(\tau,\vec x;\tau',\vec x')
=G_E(\tau',\vec x';\tau,\vec x)$.
\end{itemize}
Define the time–reflection $\theta(\tau,\vec x):=(-\tau,\vec x)$. Then for every test function
$f\in C_c^\infty(\{\,\tau\ge 0\,\}\times\R^3)$,
\begin{equation}
\label{eq:OSquadratic}
\mathcal Q[f]\;:=\;\int_{\tau,\tau'\ge 0}\!\!\!\!\! d\tau\,d\tau'
\int_{\R^3}\!\! d^3\vec x \int_{\R^3}\!\! d^3\vec x'\;
\overline{f(\tau,\vec x)}\;G_E\big(\theta(\tau,\vec x),(\tau',\vec x')\big)\;f(\tau',\vec x')\ \ge\ 0.
\end{equation}
Equivalently, the real Schwinger–Keldysh (SK) quadratic form \cite{SkenderisVanRees2008} and the Fisher/BKM kernel on the receiver
brane are positive semidefinite.
\end{theorem}

\begin{proof}
We proceed in three steps.

\emph{Step 1: Spectral/Käll\'en–Lehmann representation on the brane.}
Under \textnormal{(H1)}–\textnormal{(H3)}, the linearized bulk equations in the warped extra dimension
define a self–adjoint Sturm–Liouville problem. Let $\{\psi_n(y)\}$ be an orthonormal basis of
$\Z_2$–even normal modes with eigenvalues $m_n^2\ge 0$ (discrete KK tower in RS1; a bound state
plus a continuum in KR), and write the bulk field(s) in a mode expansion
$\Phi(x,y)=\sum_n \psi_n(y)\,\varphi_n(x)$ (with an integral over the continuum where present).
The brane--localized observable is linear in the fields and thus can be written as
$\mathcal O(x)=\sum_n g_n\,\varphi_n(x)$, where $g_n=\langle \mathcal O, \psi_n\rangle$
are the (real) overlaps at the brane location; for a continuum replace $\sum_n |g_n|^2\delta(\mu^2-m_n^2)$
with a nonnegative density $|g(\mu)|^2$.
Since each 4D field $\varphi_n$ is a free (Gaussian) field of mass $m_n$, its Euclidean propagator
is the free massive scalar propagator $\Delta_E^{(m_n)}$. Therefore
\begin{equation}
\label{eq:KL}
G_E(x-x')\;=\;\sum_{n}\frac{|g_n|^2}{(2\pi)^4}\!\int_{\R^4}\!\frac{e^{ip\cdot(x-x')}}{p^2+m_n^2}\,d^4p
\quad\text{(RS1; discrete)}\!,
\qquad
G_E(x-x')\;=\;\frac{|g_0|^2}{(2\pi)^4}\!\int\!\frac{e^{ip\cdot(x-x')}}{p^2+m_0^2}\,d^4p
+\!\int_{m_{\rm gap}}^\infty\!\!\!\frac{|g(\mu)|^2}{(2\pi)^4}\!\int\!\frac{e^{ip\cdot(x-x')}}{p^2+\mu^2}\,d^4p\,d\mu
\end{equation}
(KR; bound state $m_0$ optional).\footnote{Gauge/pure–diffeomorphism directions are null and drop out;
$\mathcal O$ is assumed gauge–invariant so that only physical polarizations contribute.}
In either case \eqref{eq:KL} is a Käll\'en–Lehmann representation \cite{WeinbergQFTvol1} with a \emph{nonnegative} spectral measure
$\rho(\mu^2)\ge 0$:
\(
G_E(p)=\int_0^\infty \rho(\mu^2)\,(p^2+\mu^2)^{-1}\,d\mu^2
\),
where $\rho$ is a sum of nonnegative point masses (RS1) and/or a nonnegative density (KR).
Positivity of the residues follows from the positivity of the inner product in \textnormal{(H1)}–\textnormal{(H2)}
and from the positive radion kinetic term in \textnormal{(H3)}.

\emph{Step 2: OS–positivity for a free massive propagator.}
Fix $m\ge 0$ and let $\Delta_E^{(m)}$ be the Euclidean free scalar propagator in four dimensions.
For $f\in C_c^\infty(\{\tau\ge 0\}\times\R^3)$, write the spatial Fourier transform
$\tilde f(\tau,\vec p)=\int_{\R^3} e^{-i\vec p\cdot \vec x}f(\tau,\vec x)\,d^3\vec x$ and set
$E(\vec p):=\sqrt{|\vec p|^2+m^2}$. One has the standard representation
\[
\Delta_E^{(m)}\big(\theta(\tau,\vec x),(\tau',\vec x')\big)
=\int_{\R^3}\frac{d^3\vec p}{(2\pi)^3}\,\frac{e^{i\vec p\cdot(\vec x-\vec x')}}{2E(\vec p)}
\,e^{-E(\vec p)(\tau+\tau')}\,.
\]
Hence the OS quadratic form for mass $m$ factorizes as
\begin{align*}
\mathcal Q_m[f]
&:=\int_{\tau,\tau'\ge 0}\!\! d\tau\,d\tau' \int_{\R^3}\!\! d^3\vec x \int_{\R^3}\!\! d^3\vec x'\;
\overline{f(\tau,\vec x)}\,\Delta_E^{(m)}\big(\theta(\tau,\vec x),(\tau',\vec x')\big)\,f(\tau',\vec x')\\
&=\int_{\R^3}\frac{d^3\vec p}{(2\pi)^3}\,\frac{1}{2E(\vec p)}\;
\overline{\int_{0}^{\infty} e^{-E(\vec p)\tau}\,\tilde f(\tau,\vec p)\,d\tau}\;
\Big(\int_{0}^{\infty} e^{-E(\vec p)\tau'}\,\tilde f(\tau',\vec p)\,d\tau'\Big)\\
&=\int_{\R^3}\frac{d^3\vec p}{(2\pi)^3}\,\frac{1}{2E(\vec p)}\;
\Big|\int_{0}^{\infty} e^{-E(\vec p)\tau}\,\tilde f(\tau,\vec p)\,d\tau\Big|^2\ \ge\ 0\,.
\end{align*}
Thus each free massive propagator is OS–positive.

\emph{Step 3: Positive superpositions preserve OS–positivity.}
By \eqref{eq:KL}, $G_E$ is a (Bochner–)integral/sum of OS–positive kernels with nonnegative weights
$|g_n|^2$ or $|g(\mu)|^2\,d\mu$. The quadratic form \eqref{eq:OSquadratic} is therefore a nonnegative
sum/integral of the nonnegative quantities $\mathcal Q_m[f]$ from Step~2, hence
$\mathcal Q[f]\ge 0$. This proves \eqref{eq:OSquadratic}. The equivalence with the positivity of the real
SK/Fisher quadratic follows from the standard OS$\leftrightarrow$KMS/Wick–rotation correspondence
for Gaussian (quadratic) functionals under \textnormal{(H1)}–\textnormal{(H4)}.
\end{proof}

\begin{remark}[Scope and failure modes]
\label{rem:uom-main-003}
\leavevmode
\begin{itemize}
\item In KR the spectrum consists of a bound state (possibly absent for some spins) plus a continuum
above a gap; in RS1 it is a discrete KK tower. In both cases the residues are nonnegative and the proof applies verbatim.
\item If \textnormal{(H1)} or \textnormal{(H3)} fail (e.g.\ a ghost radion or wrong–sign brane curvature terms),
the spectral density need not be nonnegative and OS–positivity can be violated; this corresponds to a genuinely
non–unitary setup.
\item Pure–gauge/diffeomorphism directions yield null contributions and are quotiented out by the
gauge–invariance assumption on $\mathcal O$.
\end{itemize}
\end{remark}


\subsubsection*{Decomposition completeness}\label{subsubsec:uom-main-007}

\paragraph{Completeness Theorem (receiver–side modular penalty).}
\label{par:completeness-theorem}
Fix a receiver cutoff $\Lambda$ and work in the Gaussian, zero–mean (co–centering, zero–DC) gauge at the AdS$_4$ fixed point $\rho_\ast$.
Let $\mathscr T$ be the welded covariance tangent space and equip it with the Fisher inner product
$\langle\!\langle \delta C,\delta C'\rangle\!\rangle_{C_\ast}=\tfrac12\,\mathrm{tr}\!\big(C_\ast^{-1}\delta C\,C_\ast^{-1}\delta C'\big)$.
Let $\Pi_{\rm RP}$, $\Pi_{\rm Weyl}$, $\Pi_{\rm conf}$ be the orthogonal projectors onto the reflection–odd, Weyl (Riegert), and conformal–mismatch tangents as defined in~\S\ref{subsec:P3-master}, and define the \emph{physical} kernel
\[
\mathcal T_{\rm phys}\;:=\;\ker\Pi_{\rm RP}\ \cap\ \ker\Pi_{\rm Weyl}\ \cap\ \ker\Pi_{\rm conf}.
\]

\begin{lemma}[Gaussian channel image at fixed bridge (monitor-based)]
\label{lem:gauss-image-monitor}
Fix regulators $(\Lambda,\varepsilon,\mu)$ and a welded receiver uplift which, in the Gaussian sector,
is a CP map with \emph{minimal} Stinespring dilation:
\[
C_{\rm out}\;=\;X\,C_{\rm in}\,X^{\sf T}\;+\;Y,
\]
where $X$ is a linear map between one-particle spaces and $Y=Y^{\sf T}\!\succeq 0$ is the additive
noise/counterterm block that saturates complete positivity. Assume the baseline respects the receiver
monitors:
\begin{enumerate}[label=(M\arabic*),leftmargin=1.8em]
\item (\emph{RP symmetry}) the baseline channel is reflection-even; RP-odd tangents are picked by
$\Pi_{\rm RP}=\tfrac12(\mathrm{Id}-\mathrm{Ad}\,\Theta)$.
\item (\emph{Weyl separation}) the pure-trace (Weyl) tangent is the image of the compensator map
$\mathsf D_{\rm W}$; its projector $\Pi_{\rm Weyl}$ is the corresponding Fisher-orthogonal projector.
\item (\emph{Conformal monitor}) the conformality-mismatch tangent $\mathcal T_{\rm conf}$ is the closure of the
image of the receiver \emph{conformal Ward-residual} map $C\mapsto({\cal O}_i[C])_i$, Fisher-orthogonalized
to the Weyl tangent; denote its projector by $\Pi_{\rm conf}$.
\end{enumerate}
Let $\delta C_{\rm out}$ be a receiver covariance tangent satisfying
$\Pi_{\rm RP}\delta C_{\rm out}=\Pi_{\rm Weyl}\delta C_{\rm out}=\Pi_{\rm conf}\delta C_{\rm out}=0$.
Then necessarily $\delta Y=0$ and $\delta X=0$, and there exists $\delta C_{\rm in}$ (unique modulo
$\ker(\mathrm{Ad}_X):=\{Z:\ X Z X^{\sf T}=0\}$) such that
\[
\delta C_{\rm out}\;=\;X\,\delta C_{\rm in}\,X^{\sf T}.
\]
Equivalently,
\(
\mathcal T_{\rm phys}\;=\;\mathrm{Ad}_X(\mathscr T_{\rm in})
:=\{\,X\delta C_{\rm in}X^{\sf T}:\ \delta C_{\rm in}\in\mathscr T_{\rm in}\,\}
\)
is the physical receiver tangent space.
\end{lemma}

\begin{proof}
Every infinitesimal receiver variation induced by the Gaussian channel can be written as
\[
\delta C_{\rm out}\;=\;X\,\delta C_{\rm in}\,X^{\sf T}\;+\;(\delta X)\,C_{\rm in}\,X^{\sf T}
\;+\;X\,C_{\rm in}\,(\delta X)^{\sf T}\;+\;\delta Y.
\tag{$\ast$}
\]
\emph{Step 1 (No additive-noise tangent in the physical kernel).}
Minimal Stinespring means $Y+i(\Omega-X\Omega X^{\sf T})\ge 0$ saturates the CP bound.
Thus any $\delta Y\neq 0$ is a \emph{channel} deformation (it modifies the intrinsic noise). Its
RP-odd component lies in $\mathcal T_{\rm RP}$ by definition. Its even component contributes a
nonzero pure-trace (curvature-variance) piece detected by $\Pi_{\rm Weyl}$ unless it is identically
zero in the Fisher sense (the compensator/Weyl map separates the trace sector). Therefore
$\Pi_{\rm RP}\delta Y+\Pi_{\rm Weyl}\delta Y\neq 0$ for any $\delta Y\neq 0$, and since
$\delta C_{\rm out}$ is in the joint kernel, we must have $\delta Y=0$.

\emph{Step 2 (No linear-map tangent in the physical kernel).}
With $\delta Y=0$, the remaining deformation $(\delta X)$ is reflection-even, and we have
\[
\delta C_{\rm out}=X\,\delta C_{\rm in}\,X^{\sf T}+(\delta X)\,C_{\rm in}\,X^{\sf T}
+X\,C_{\rm in}\,(\delta X)^{\sf T}.
\]
The last two terms change the \emph{channel} \(X\). If their projection on the Weyl tangent vanishes
(by hypothesis), then only the \emph{traceless, reflection-even} part remains. Such a change necessarily
produces a nonzero receiver \emph{Ward-residual} unless $\delta X=0$:
the Ward operators ${\cal O}_i$ act linearly on $C_{\rm out}$, and ${\cal O}_i$ commute with
the baseline reflection/Weyl reductions, so any nonzero even, traceless channel term induces
a nonzero $({\cal O}_i[\delta C_{\rm out}])_i$, i.e.\ lies in $\mathcal T_{\rm conf}$.
Since $\delta C_{\rm out}$ is in $\ker\Pi_{\rm conf}$, we conclude $\delta X=0$.

\emph{Step 3 (Push–forward form).}
With $\delta X=\delta Y=0$, \((\ast)\) collapses to $\delta C_{\rm out}=X\,\delta C_{\rm in}\,X^{\sf T}$.
Uniqueness is modulo $\ker(\mathrm{Ad}_X)$, as usual for linear push–forwards. This identifies
$\mathcal T_{\rm phys}=\mathrm{Ad}_X(\mathscr T_{\rm in})$.
\end{proof}



\begin{theorem}[Completeness of the unphysical decomposition]\label{thm:completeness}
Fix the receiver cutoff $\Lambda$ and the AdS$_4$ fixed point $\rho_\ast$ (maximally symmetric; zero–mean gauge). 
Let $\mathscr T$ be the welded covariance tangent space at $C_\ast$ and equip it with the Fisher inner product
\[
\langle\!\langle \delta C,\delta C'\rangle\!\rangle_{C_\ast}
\;:=\;\tfrac12\,\mathrm{tr}\!\big(C_\ast^{-1}\delta C\,C_\ast^{-1}\delta C'\big).
\]
Assume the symmetry hypotheses:
\begin{enumerate}[label=(S\arabic*),leftmargin=2.1em]
\item \emph{Euclidean reflection} $\Theta$ acts as an involutive isometry on $\mathscr T$ and on $C_\ast$ (Osterwalder–Schrader reflection).
\item \emph{Local Weyl flow} $g\mapsto e^{2\sigma}g$ integrates to an isometric action on $\mathscr T$; its linearization at $C_\ast$ defines the (closed) Weyl tangent $\mathcal T_{\rm Weyl}$ (Riegert direction).
\item \emph{Boundary conformal actions} $U_\partial(g)$ (with $g\in\mathrm{Conf}(S^3)$) act by Fisher isometries on $\mathscr T$; their linearization at $C_\ast$ defines the (closed) conformal–mismatch tangent $\mathcal T_{\rm conf}$.
\item The three symmetry actions commute at $\rho_\ast$ and preserve $C_\ast$ (maximal symmetry).
\end{enumerate}
Let $\Pi_{\rm RP},\Pi_{\rm Weyl},\Pi_{\rm conf}$ be the orthogonal projectors (with respect to $\langle\!\langle\cdot,\cdot\rangle\!\rangle_{C_\ast}$) onto the closed subspaces
\[
\mathcal T_{\rm RP} := \mathrm{ran}\,\Pi_{\rm RP}\quad(\text{reflection–odd / OS–violating}),\qquad
\mathcal T_{\rm Weyl} := \mathrm{ran}\,\Pi_{\rm Weyl},\qquad
\mathcal T_{\rm conf} := \mathrm{ran}\,\Pi_{\rm conf}.
\]
Define the \emph{physical} kernel $\mathcal T_{\rm phys}:=\ker\Pi_{\rm RP}\cap\ker\Pi_{\rm Weyl}\cap\ker\Pi_{\rm conf}$.
Then:
\begin{enumerate}[leftmargin=2.1em]
\item \textbf{Orthogonal direct sum.} One has the orthogonal decomposition
\[
\mathscr T\;=\;\mathcal T_{\rm RP}\ \widehat\oplus\ \mathcal T_{\rm Weyl}\ \widehat\oplus\ \mathcal T_{\rm conf}\ \widehat\oplus\ \mathcal T_{\rm phys}.
\]
\item \textbf{Exhaustiveness.} If $\delta C\in\mathscr T$ is orthogonal to $\mathcal T_{\rm RP}$, $\mathcal T_{\rm Weyl}$, and $\mathcal T_{\rm conf}$, then $\delta C\in\mathcal T_{\rm phys}$. Equivalently, there is no further symmetry–invariant quadratic receiver cost at quadratic order beyond these three blocks.
\item \textbf{Identification with payload.} In the Gaussian, zero–mean gauge, $\mathcal T_{\rm phys}$ coincides with welded images of payload (record–carrying) variations: any $\delta C\in\mathcal T_{\rm phys}$ is realized (to quadratic order) by varying payload while remaining on the RP/Weyl/conformal leaves, and such variations contribute zero \emph{receiver} quadratic modular cost.
\item \textbf{Uniqueness of quadratic costs.} Any continuous, symmetry–invariant, positive semidefinite quadratic functional $\Psi$ on $\mathscr T$ with $\Psi|_{\mathcal T_{\rm phys}}=0$ is necessarily
\[
\Psi(\delta C)\;=\;\alpha_{\rm RP}\,\|\Pi_{\rm RP}\delta C\|_{C_\ast}^2
+\alpha_{\rm Weyl}\,\|\Pi_{\rm Weyl}\delta C\|_{C_\ast}^2
+\alpha_{\rm conf}\,\|\Pi_{\rm conf}\delta C\|_{C_\ast}^2,\quad \alpha_{\rm RP},\alpha_{\rm Weyl},\alpha_{\rm conf}\ge 0.
\]
\end{enumerate}
\end{theorem}

\begin{proof}
\emph{Step 1: Closed invariant subspaces and pairwise orthogonality.}
By (S1), $\Theta$ is an involutive isometry: $\Theta^2=\mathbf 1$ and
$\langle\!\langle \Theta\delta C,\Theta\delta C'\rangle\!\rangle_{C_\ast}=\langle\!\langle \delta C,\delta C'\rangle\!\rangle_{C_\ast}$.
Thus $\mathscr T$ splits into $\pm 1$ eigenspaces $\mathscr T^{(+)}\oplus \mathscr T^{(-)}$,
which are orthogonal. The reflection–odd space $\mathscr T^{(-)}$ is the tangent along which the OS cone is left to first order; set $\mathcal T_{\rm RP}:=\mathscr T^{(-)}$, which is closed (finite–codimension fixed–point of a bounded projector).

For Weyl and conformal actions, (S2)–(S3) imply one–parameter isometry groups on $\mathscr T$; their linearizations at $C_\ast$ define closed invariant subspaces $\mathcal T_{\rm Weyl}$ and $\mathcal T_{\rm conf}$. On a maximally symmetric background (S4), the Fisher metric inherits the full AdS$_4$ isometry group acting transitively on tangent indices. The Weyl tangent is \emph{pure trace} (scalar under local Lorentz at a point), while the conformal–mismatch tangent is \emph{traceless} (it is generated by boundary conformal diffeomorphisms and associated source redeclarations, which lie in the traceless shape sector).
Hence, by index orthogonality on a maximally symmetric background,
\[
\langle\!\langle \delta C_{\rm Weyl},\ \delta C_{\rm conf}\rangle\!\rangle_{C_\ast}=0.
\]
Moreover, $\Theta$ commutes with Weyl and conformal actions (S4), and both Weyl and conformal tangents are reflection–even (they are generated by local scalings and boundary conformal isometries, which preserve OS parity). Therefore
\[
\langle\!\langle \delta C_{\rm RP},\ \delta C_{\rm Weyl}\rangle\!\rangle_{C_\ast}
=\langle\!\langle \delta C_{\rm RP},\ \delta C_{\rm conf}\rangle\!\rangle_{C_\ast}=0,
\]
establishing pairwise orthogonality and closedness.

\emph{Step 2: Algebraic direct sum and orthogonal complement.}
Let $\Pi_{\rm RP}$ be the orthogonal projector onto $\mathcal T_{\rm RP}=\mathscr T^{(-)}$,
and let $\Pi_{\rm Weyl}$, $\Pi_{\rm conf}$ be the orthogonal projectors onto $\mathcal T_{\rm Weyl}$ and $\mathcal T_{\rm conf}$. Orthogonality implies the operator
\[
\Pi_{\rm unphys}\;:=\;\Pi_{\rm RP}+\Pi_{\rm Weyl}+\Pi_{\rm conf}
\]
is itself an orthogonal projector onto the closed sum
$\mathcal T_{\rm RP}\oplus\mathcal T_{\rm Weyl}\oplus\mathcal T_{\rm conf}$.
Thus
\[
\big(\mathcal T_{\rm RP}\oplus\mathcal T_{\rm Weyl}\oplus\mathcal T_{\rm conf}\big)^\perp
=\ker \Pi_{\rm unphys}.
\]
By definition $\mathcal T_{\rm phys}=\ker\Pi_{\rm RP}\cap\ker\Pi_{\rm Weyl}\cap\ker\Pi_{\rm conf}
=\ker\Pi_{\rm unphys}$, hence
\[
\mathscr T\;=\;\mathrm{ran}\,\Pi_{\rm unphys}\ \widehat\oplus\ \ker\Pi_{\rm unphys}
\;=\;\big(\mathcal T_{\rm RP}\widehat\oplus \mathcal T_{\rm Weyl}\widehat\oplus \mathcal T_{\rm conf}\big)\widehat\oplus \mathcal T_{\rm phys},
\]
which proves the orthogonal direct sum (item 1) and the exhaustiveness statement (item 2).

\emph{Step 3 (identification of $\mathcal T_{\rm phys}$) with payload directions.}
At fixed tick and regulators, the welded uplift is a minimal Gaussian channel
$C_{\rm out}=X C_{\rm in}X^{\sf T}+Y$. Let $\delta C\in\mathscr T$ satisfy
$\Pi_{\rm RP}\delta C=\Pi_{\rm Weyl}\delta C=\Pi_{\rm conf}\delta C=0$. By
Lemma~\ref{lem:gauss-image-monitor}, any receiver tangent with those properties must be
of the form $\delta C=X\,\delta C_{\rm in}\,X^{\sf T}$ for some input tangent
$\delta C_{\rm in}$: no other receiver variation is compatible with keeping the
RP/Weyl/Conf leaves fixed. Conversely, any such push–forward is manifestly
reflection–even, Weyl–neutral (to quadratic order around the maximally symmetric
$C_\ast$), and conformally aligned, hence lies in the triple kernel. Therefore
\[
\mathcal T_{\rm phys}\;=\;\ker\Pi_{\rm RP}\cap\ker\Pi_{\rm Weyl}\cap\ker\Pi_{\rm conf}
\;=\;\mathrm{Ad}_X(\mathscr T_{\rm in}),
\]
i.e.\ precisely the welded images of payload (record–carrying) input variations.
This proves item~(3).


\emph{Step 4: Uniqueness of symmetry–invariant quadratic costs on $\mathscr T$.}
Let $\Psi$ be a continuous, symmetry–invariant, positive semidefinite quadratic functional on $\mathscr T$ with $\Psi|_{\mathcal T_{\rm phys}}=0$. Invariance under (S1)–(S3) and the orthogonal direct sum from Step 2 imply that $\Psi$ is block–diagonal with respect to the decomposition
$\mathscr T=\mathcal T_{\rm RP}\widehat\oplus\mathcal T_{\rm Weyl}\widehat\oplus\mathcal T_{\rm conf}\widehat\oplus\mathcal T_{\rm phys}$, and vanishes on $\mathcal T_{\rm phys}$. 
Within each block, maximal symmetry (S4) acts irreducibly on the corresponding tangent (reflection–odd scalar block, pure–trace scalar block, and traceless conformal block). By Schur’s lemma for real representations (or, concretely, by isotropy: the only symmetric bilinear form invariant under the stabilizer is a scalar multiple of the Fisher metric on that block), it follows that on each subspace $\mathcal T_\bullet$,
\[
\Psi|_{\mathcal T_\bullet}=\alpha_\bullet\,\|\cdot\|_{C_\ast}^2\qquad(\alpha_\bullet\ge 0).
\]
Hence, for all $\delta C\in\mathscr T$,
\[
\Psi(\delta C)
=\alpha_{\rm RP}\,\|\Pi_{\rm RP}\delta C\|_{C_\ast}^2
+\alpha_{\rm Weyl}\,\|\Pi_{\rm Weyl}\delta C\|_{C_\ast}^2
+\alpha_{\rm conf}\,\|\Pi_{\rm conf}\delta C\|_{C_\ast}^2,
\]
which is item 4.

Combining Steps 1–4 establishes all claims.
\end{proof}


\medskip
\noindent
\textbf{Consequences.} (1) The receiver penalty is \emph{by construction}
insensitive to the physical directions $\mathscr V_k^{\rm phys}$ (where it
vanishes). (2) The unphysical cost space is \emph{exactly} the orthogonal sum
$\mathscr W_k^{\rm conf}\,\widehat\oplus\,\mathscr W_k^{\rm RP}\,\widehat\oplus\,
\mathscr W_k^{\rm BY}$; there are no missing unphysical terms consistent with
the symmetries. (3) Because the penalty is the $\Lambda$–norm and the monitors
are linear on the LoG family, convexity of $\Phi_{\rm rec}$ in the tick controls
follows immediately, as used in Sec.~\ref{sec:aeon-evolution}.

\subsubsection*{Summary of the receiver–side decomposition}\label{subsubsec:uom-main-008}

We collect the consequences of the preceding results for the receiver penalty and its monitors.

\medskip
\noindent
\emph{Receiver-side modular free energy decomposition.} By Theorem~\ref{thm:completeness} above, the quadratic part
of the receiver modular energy coincides with the sum of the three Fisher–quadratic monitors
(RP, Weyl, conformality) with geometry–fixed weights, and the physical direction subspace, spanned by the payload DoFs. There is no additional symmetry–invariantquadratic block. 


\medskip
\noindent
\emph{Weyl orthogonality (KR head; interior RS1 comparison).} For an interior RS1 segment the scalar sector
exhibits a Weyl$\leftrightarrow$radion mixing angle $\theta(q)$ with the asymptotics
\[
\theta(q)\;=\;-\frac{b}{c}\,\frac{q^{2}}{m_r^{2}}+o(q^{2})\quad(q\downarrow 0),
\qquad
\theta(q)\;=\;\frac{b}{a}\,\frac{1}{q^{2}}+o(q^{-2})\quad(q\to\infty),
\]
(Prop.~\ref{prop:scalar-mix}(ii)). On the active KR head receiver, the same $q^2$ and $q^{-2}$ behavior holds
with the replacement $b\mapsto \tilde b\,\varepsilon_{\rm BC}$ (Prop.~\ref{prop:scalar-mix}(i)), hence
$\theta(\Lambda)=O(\varepsilon_{\rm BC}\Lambda^{-2})\ll 1$ at welded scales $q\sim\Lambda\gg m_r$. In particular,
for any welded tangent $\delta C$ there exists a constant $K>0$ (independent of $\delta C$ and $\Lambda$) such that
\begin{equation}
\label{eq:near-orth}
\Big|\big\langle\!\big\langle \Pi_{\rm Weyl}\delta C,\ \Pi_{\rm conf}\delta C\big\rangle\!\big\rangle_{C_\ast}\Big|
\ \le\ K\,\theta(\Lambda)\,\Big(\|\Pi_{\rm Weyl}\delta C\|_{C_\ast}^{2}+\|\Pi_{\rm conf}\delta C\|_{C_\ast}^{2}\Big)
\ =\ O\!\big(\Lambda^{-2}\big)\,\Big(\cdots\Big),
\end{equation}
i.e.\ the conformality–mismatch and Weyl subspaces are Fisher–orthogonal up to corrections of order
$\Lambda^{-2}$ at welded cutoff.

\medskip
\noindent
\emph{RP block: vanishing contribution.} By Theorem~\ref{thm:OS_receiver}, the receiver two–point functions
obey Osterwalder–Schrader positivity. Consequently, the reflection–odd tangent (picked by
$\Pi_{\rm RP}=\tfrac12(\mathrm{Id}-\mathrm{Ad}\,\Theta)$) is absent for welded Gaussian states; in
particular, for the LoG welded family one has
\begin{equation}
\label{eq:rp-vanish}
\Delta D_{{\rm RP},k}\;=\;0\qquad\text{(receiver RP monitor identically vanishes).}
\end{equation}

\medskip
\noindent
\emph{Simplified receiver penalty.} Combining \eqref{eq:near-orth} and \eqref{eq:rp-vanish} gives the
following operational form of the per–tick receiver penalty on the welded family:
\begin{equation}
\label{eq:Jrec-simplified}
\mathcal J^{\rm rec}_k(\sigma_{t,k},\sigma_{x,k},A_k,\Lambda_k)
\;=\;\lambda_{\rm BY}\,\Theta_k\;+\;\lambda_{\rm conf}\,D_{{\rm conf},k}
\;+\;O\!\big(\Lambda_k^{-2}\big)\,\Xi_k,
\end{equation}
where $\Xi_k$ is a nonnegative quadratic form (coming from the residual
$\langle\!\langle \Pi_{\rm Weyl}\delta C,\Pi_{\rm conf}\delta C\rangle\!\rangle$) that vanishes on the
active KR head and is uniformly suppressed on interior RS1 segments by
$\theta(\Lambda_k)=O(\Lambda_k^{-2})$. In particular:
\begin{itemize}
\item \textbf{KR head (active receiver):} $\mathcal J^{\rm rec}_k=\lambda_{\rm BY}\,\Theta_k$ (the conformality–mismatch penalty
and the RP monitor vanish at quadratic order).
\item \textbf{Interior RS1 segment (shadow; not an active receiver):} the $O(\Lambda_k^{-2})$ term is negligible
at welded cutoff; one may treat $\Pi_{\rm Weyl}$ and $\Pi_{\rm conf}$ as Fisher–orthogonal when analyzing the
interior segment and use \eqref{eq:Jrec-simplified} with the last term dropped to leading order.
\end{itemize}



\subsubsection*{Source-vs-receiver separation by commuting kernels}
\label{subsubsec:commute-shapes}

\paragraph{Setup and notation.}\label{par:uom-main-054}
Let $\mathcal H:=L^2(\mathbb R_t)\otimes L^2_{\rm zonal}(S^3_\Omega)$ be the welded control
Hilbert space restricted to (i) \emph{temporal} zero–DC, parity–odd envelopes and
(ii) \emph{spatial} zonal envelopes. Denote by $\widehat{\cdot}$ the temporal Fourier transform
and expand spatially in $S^3$ spherical harmonics $Y_{\ell 0}(\Omega)$ with eigenvalues
$-\Delta_{S^3}Y_{\ell 0}=\lambda_\ell Y_{\ell 0}$, $\lambda_\ell:=\ell(\ell+2)$.
The welded controls factor as $h(t,\Omega)=h_t(t)\,h_x(\Omega)$.

The \emph{throughput} quadratic form $\mathsf S$ and the \emph{source dwell} (Spohn) quadratic form
$\mathsf W$ act on $\mathcal H$ as bounded, positive, reflection–even operators, stationary in time
and rotationally invariant in space. Hence there exist Borel multipliers $s(\omega,\lambda)$ and
$w(\omega,\lambda)$ with $s,w>0$ and even in $\omega$ such that, for $h\in\mathcal H$,
\begin{equation}
\label{eq:SW-mult}
\langle h,\ \mathsf S\,h\rangle
=\int_{\mathbb R}\!\frac{d\omega}{2\pi}\sum_{\ell\ge0} s(\omega,\lambda_\ell)\,|\widehat{h_t}(\omega)|^2\,|h_{x,\ell}|^2,\qquad
\langle h,\ \mathsf W\,h\rangle
=\int_{\mathbb R}\!\frac{d\omega}{2\pi}\sum_{\ell\ge0} w(\omega,\lambda_\ell)\,|\widehat{h_t}(\omega)|^2\,|h_{x,\ell}|^2.
\end{equation}

\begin{proposition}[Commutation and joint diagonalization]
\label{prop:SW-commute}
On $\mathcal H$, $\mathsf S$ and $\mathsf W$ are bounded Borel functions of the commuting
self–adjoint pair $(|\partial_t|,\sqrt{-\Delta_{S^3}})$; in particular $[\mathsf S,\mathsf W]=0$,
and both are simultaneously diagonalized by the product basis
$\{e^{i\omega t}\otimes Y_{\ell 0}(\Omega)\}$ with eigenvalues $s(\omega,\lambda_\ell)$
and $w(\omega,\lambda_\ell)$ given by \eqref{eq:SW-mult}.
\end{proposition}

\begin{proof}
Stationarity and $SO(4)$–invariance imply that both kernels are (zonal) convolutions:
$\mathsf S=\mathsf S_t\otimes\mathsf S_x$, $\mathsf W=\mathsf W_t\otimes\mathsf W_x$ with
$\mathsf S_t,\mathsf W_t$ temporal convolutions (Toeplitz) and $\mathsf S_x,\mathsf W_x$
zonal $S^3$ convolutions. Convolution operators are Fourier (resp.\ harmonic) multipliers,
hence both are Borel functions of $|\partial_t|$ and $-\Delta_{S^3}$ and commute strongly.
\end{proof}

\begin{theorem}[Separation of minimizations (source vs receiver)]
\label{thm:separation}
Fix a tick with payload $\Delta I^\star>0$ and let $\phi$ range over admissible temporal shapes
(parity–odd, $\int\phi=0$) and $\Psi$ over admissible zonal spatial shapes. Then the source action
\[
\Phi^{\rm src}(\phi,\Psi)\;:=\;\Delta I^\star\,
\frac{\langle \phi\!\otimes\!\Psi,\ \mathsf W\,(\phi\!\otimes\!\Psi)\rangle}
     {\langle \phi\!\otimes\!\Psi,\ \mathsf S\,(\phi\!\otimes\!\Psi)\rangle}
\]
factorizes as a ratio of temporal and spatial Rayleigh quotients and is minimized by
\emph{separate} temporal/spatial extremizers $(\phi^\star,\Psi^\star)$ that solve the
generalized eigenvalue problems $\mathsf W_t\phi=\lambda_t\,\mathsf S_t\phi$,
$\mathsf W_x\Psi=\lambda_x\,\mathsf S_x\Psi$. In particular,
$\Phi^{\rm src}(\phi^\star,\Psi^\star)=\Delta I^\star\,\lambda_t\lambda_x$ is constant on the
calibrated extremal family and does not affect the receiver–side argmin.
\end{theorem}

\begin{proof}
By Proposition~\ref{prop:SW-commute}, both forms are diagonal in the same basis. For rank–one
products $h=\phi\otimes\Psi$, the quotient splits:
\[
\frac{\langle h,\mathsf W h\rangle}{\langle h,\mathsf S h\rangle}
=\frac{\langle \phi,\mathsf W_t\phi\rangle}{\langle \phi,\mathsf S_t\phi\rangle}\cdot
 \frac{\langle \Psi,\mathsf W_x\Psi\rangle}{\langle \Psi,\mathsf S_x\Psi\rangle}.
\]
Each factor is minimized by the corresponding generalized eigenfunction; the minimum is the
product of the two minimum eigenvalues. On the calibrated family the minimizers are unique up to
scale/centering, hence $\Phi^{\rm src}=\Delta I^\star\,\lambda_t\lambda_x$ is tickwise constant.
\end{proof}

\subsubsection*{Separable minimization variables}
\label{subsubsec:separable-splits-centering}

\paragraph{Welded scales and the split variables.}\label{par:uom-main-055}
Let the \emph{macro} (source–side) spatial/temporal widths be \((\rho_k,\Sigma_k)\) and the \emph{micro} (receiver–side) widths be \((\sigma_{x,k},\sigma_{t,k})\).
The welded LoG envelopes (Sec.~\ref{subsec:P3-master}) have the \emph{effective} scales
\[
s_{x,k}\ :=\ \sqrt{\rho_k^2+\sigma_{x,k}^2},
\qquad
s_{t,k}\ :=\ \sqrt{\Sigma_k^2+\sigma_{t,k}^2},
\]
and are centered at a common spacetime point \((t_k,\Omega_k)\) when macro/micro pulses are co–centered.
The receiver cutoff is \(\Lambda_k\).

\medskip
\noindent
Receiver–side quadratic costs (BY/anomaly and conformality) computed on the welded family are functions of \((s_{t,k},s_{x,k};\Lambda_k)\) alone:
\begin{equation}
\tag{R1}
\label{eq:R1}
\Phi^{\rm rec}_k\big(\rho_k,\Sigma_k,\sigma_{x,k},\sigma_{t,k},\Lambda_k\big)
=\Phi^{\rm rec}_k\big(s_{t,k},s_{x,k},\Lambda_k\big).
\end{equation}

\paragraph{Receiver–side isospectrality (why \(\Phi^{\rm rec}\) sees only \(s_t,s_x\)).}\label{par:uom-main-056}
For the temporal DoG–DoG convolution one has
\(
(g'_{\Sigma}*g'_{\sigma_t})(t)
= \frac{d^2}{dt^2}\,g_{s_t}(t)
\) 
with \(s_t=\sqrt{\Sigma^2+\sigma_t^2}\), and on \(S^3\) the zonal LoG–Gaussian convolution gives
\(
\Delta_{S^3}g_{\rho}\circledast g_{\sigma_x}=\Delta_{S^3}g_{s_x}
\)
with \(s_x=\sqrt{\rho^2+\sigma_x^2}\).
In Fourier/harmonic variables,
\[
\widehat{\mathrm{LoG}}_t(\omega;s_t)=(i\omega)^2 e^{-\tfrac12 s_t^2\omega^2}\,C_t,\qquad
\widehat{\mathrm{LoG}}_x(\ell;s_x)= -\ell(\ell+2)\,e^{-\tfrac12 s_x^2 k_\ell^2}\,C_x,
\]
where \(k_\ell\!\sim\!(\ell{+}1)/L_4\) and \(C_{t,x}\) are shape constants.
Hence any welded quadratic receiver functional of the form
\(
\int |\widehat{\tau}|^2\,\mathcal K_\Lambda
\)
with a maximally symmetric kernel \(\mathcal K_\Lambda(\omega,\ell)\) depends on \(\Sigma,\sigma_t\) and \(\rho,\sigma_x\) \emph{only} through \(s_t,s_x\).
In particular, after amplitude elimination via the payload equality
\(
\Delta I_k=\big(A^{\rm eff}_k\big)^2\,G(s_{t,k},s_{x,k};\Lambda_k),
\)
the BY term takes the welded form
\begin{equation}
\Phi_{{\rm BY},k} \;=\; 
\gamma\,(8\pi G_4)^2\,C_{\rm LoG}\;
\frac{\Delta I_k}{s_{t,k}\,s_{x,k}^{3}\,G(s_{t,k},s_{x,k};\Lambda_k)},
\tag{R2}\label{eq:R2}
\end{equation}
and the RP/conformality surrogates are likewise functions of \((s_{t,k},s_{x,k};\Lambda_k)\).
This proves \eqref{eq:R1}.

\paragraph{Separable minimization of widths (receiver vs.\ source).}\label{par:uom-main-057}
Let the per–tick action be
\[
\mathcal A_k\;=\;
\underbrace{\Phi^{\rm rec}_k\big(s_{t,k},s_{x,k},\Lambda_k\big)}_{\text{receiver (BY/RP/conf)}}
\;+\;
\underbrace{\mathcal A^{\rm src}_k\big(\Sigma_k,\rho_k;\,\Delta I_k\big)}_\text{source (modular dwell - record)}.
\]
The welded scales and the splits are related by
\(
s_{t,k}^2=\Sigma_k^2+\sigma_{t,k}^2,\ 
s_{x,k}^2=\rho_k^2+\sigma_{x,k}^2
\)
and the micro widths obey lower bounds
\(\sigma_{t,k}\ge\sigma_{\rm QSL}\), \(\sigma_{x,k}\ge\sigma_{x,{\rm min}}\) (instrument/QSL).

\begin{proposition}[Separable minimization of width splits]\label{prop:separable}
For fixed payload \(\Delta I_k\), the minimization of \(\mathcal A_k\) over
\((\Sigma_k,\rho_k,\sigma_{t,k},\sigma_{x,k},\Lambda_k)\) is equivalent to the two–stage problem
\[
\boxed{\ 
\min_{s_{t,k},s_{x,k},\Lambda_k}\ \Phi^{\rm rec}_k(s_{t,k},s_{x,k},\Lambda_k)
\quad+\quad
\min_{\substack{\Sigma_k,\sigma_{t,k}\\ \Sigma_k^2+\sigma_{t,k}^2=s_{t,k}^{\star 2}\\ \sigma_{t,k}\ge\sigma_{\rm QSL}}}\ \mathcal A^{\rm src}_k(\Sigma_k;\Delta I_k)
\quad+\quad
\min_{\substack{\rho_k,\sigma_{x,k}\\ \rho_k^2+\sigma_{x,k}^2=s_{x,k}^{\star 2}\\ \sigma_{x,k}\ge\sigma_{x,{\rm min}}}}\ \mathcal A^{\rm src}_k(\rho_k;\Delta I_k)\ .
}
\]
In particular, \(\rho_k\) and \(\Sigma_k\) (macro widths) are chosen \emph{exclusively} by the source–side functional at fixed welded scales \(s_{x,k}^\star,s_{t,k}^\star\) set by the receiver minimizer.
\end{proposition}

\begin{proof}
By \eqref{eq:R1}, \(\Phi^{\rm rec}_k\) is constant on each equivalence class
\[
\mathcal E(s_t,s_x):=\{(\Sigma,\sigma_t,\rho,\sigma_x):\ \Sigma^2+\sigma_t^2=s_t^2,\ \rho^2+\sigma_x^2=s_x^2,\ \sigma_t\!\ge\!\sigma_{\rm QSL},\ \sigma_x\!\ge\!\sigma_{x,{\rm min}}\}.
\]
Hence \(\min\mathcal A_k = \min_{s_t,s_x,\Lambda} \left[\Phi^{\rm rec}_k(s_t,s_x,\Lambda)+ \inf_{(\Sigma,\sigma_t,\rho,\sigma_x)\in\mathcal E(s_t,s_x)} \mathcal A^{\rm src}_k(\Sigma,\rho)\right]\).
Because \(\mathcal A^{\rm src}_k\) contains no dependence on \((\sigma_t,\sigma_x)\) beyond the constraints, the inner infimum splits into a sum of two one–dimensional constrained minimizations in \(\Sigma\) and \(\rho\) (with \(\sigma_t,\sigma_x\) slaved by the Pythagorean constraints and their floors). This gives the stated two–stage program. 
\end{proof}

\subsubsection*{Shapes optimality}
\label{subsubsec:shapes-optimality}

\paragraph{Geometry split (source vs.\ receiver).}\label{par:uom-main-058}
Temporal \emph{macro} envelopes act on the source static–patch net $(\mathcal A_{\rm sp},\alpha_t)$
with static time and KMS anchor; spatial \emph{macro} envelopes live on a static–patch spatial
slice $\Sigma_{\rm sp}$ (Riemannian metric $\gamma$, Laplace–Beltrami $-\Delta_{\Sigma_{\rm sp}}$).
Spatial \emph{micro} envelopes live on the receiver slice (round $S^3$). The optimal shapes and
co–centering are derived intrinsically in each geometry; welded conclusions then follow by additivity
over the static–patch union.


\paragraph{Temporal shape optimality (DoG).}\label{par:uom-main-059}
Let $\mathcal H_{\rm odd}(\sigma)$ be the set of $L^2$ temporal envelopes with $\int \phi=0$,
$\|\phi\|_2=1$ and fixed temporal variance $\mathrm{Var}_t[\phi]=\sigma^2$.

\begin{lemma}[Temporal DoG extremizer]
\label{lem:dog-opt}
Assume the temporal multipliers $s(\omega),w(\omega)$ are $C^2$, even, strictly positive, and
admit Taylor expansions near the origin
\[
s(\omega)=s_0+s_2\omega^2+O(\omega^4),\qquad
w(\omega)=w_0+w_2\omega^2+O(\omega^4),
\]
with $s_0,w_0>0$ and $s_2,w_2>0$. Let
\[
\mathcal H_{\rm odd}(\sigma)\;:=\;\Big\{\phi\in L^2(\mathbb R):\ \phi(t)\text{ odd},\ 
\|\phi\|_2=1,\ \int_{\mathbb R} t^2|\phi(t)|^2\,dt=\sigma^2\Big\}.
\]
Then the minimizer of the Rayleigh quotient
\[
\mathcal R_t[\phi]\;:=\;\frac{\displaystyle\int_{\mathbb R} |\widehat\phi(\omega)|^2\,w(\omega)\,\frac{d\omega}{2\pi}}
                           {\displaystyle\int_{\mathbb R} |\widehat\phi(\omega)|^2\,s(\omega)\,\frac{d\omega}{2\pi}}
\quad\text{over}\quad \phi\in\mathcal H_{\rm odd}(\sigma)
\]
is unique up to an overall phase and a common time shift of the odd axis and equals, in time domain,
\[
\phi^\star_\sigma(t)\;=\;c_\sigma\,\frac{t}{\sigma^2}\,e^{-\,t^2/(2\sigma^2)}\quad(\text{``DoG''}),
\]
with $c_\sigma$ fixed by $\|\phi^\star_\sigma\|_2=1$.
\end{lemma}

\begin{proof}
\textit{Step 1 (Quadratic reduction of the rays).}
Because $\phi$ is odd, $|\widehat\phi(\omega)|^2$ is even and concentrates near $\omega=0$ when the time variance is fixed.
Using the stated $C^2$ expansions and Parseval,
\[
\begin{aligned}
\int |\widehat\phi|^2 w\,\frac{d\omega}{2\pi}
&= w_0\!\int |\widehat\phi|^2\frac{d\omega}{2\pi} + w_2\!\int \omega^2|\widehat\phi|^2\frac{d\omega}{2\pi} + O\!\Big(\!\int \omega^4|\widehat\phi|^2\frac{d\omega}{2\pi}\!\Big),\\
\int |\widehat\phi|^2 s\,\frac{d\omega}{2\pi}
&= s_0\!\int |\widehat\phi|^2\frac{d\omega}{2\pi} + s_2\!\int \omega^2|\widehat\phi|^2\frac{d\omega}{2\pi} + O\!\Big(\!\int \omega^4|\widehat\phi|^2\frac{d\omega}{2\pi}\!\Big).
\end{aligned}
\]
Since $\|\phi\|_2=1$ and $\int \omega^2|\widehat\phi|^2/(2\pi)=\int |\,\phi'(t)\,|^2\,dt$ (by Plancherel),
\[
\mathcal R_t[\phi]\;=\;\frac{w_0 + w_2\,\int |\,\phi'|^2}{\;s_0 + s_2\,\int |\,\phi'|^2}\ +\ O\!\Big(\int \omega^4|\widehat\phi|^2\frac{d\omega}{2\pi}\Big).
\]
For $\phi$ with fixed time variance $\int t^2|\phi|^2=\sigma^2$, the last $O(\cdot)$ is uniformly controlled (by the standard moment/sobolev bounds), hence it does not alter the identity of the minimizer. Consequently, it suffices to minimize
\[
R_{\rm quad}[\phi]\;:=\;\frac{w_0 + w_2\,\int |\,\phi'|^2}{\;s_0 + s_2\,\int |\,\phi'|^2}
\]
over $\mathcal H_{\rm odd}(\sigma)$.

\medskip
\textit{Step 2 (Monotonicity in the derivative energy).}
Write $X[\phi]:=\int_{\mathbb R}|\phi'(t)|^2\,dt$. Then $R_{\rm quad}[\phi]=(w_0+w_2 X)/(s_0+s_2 X)$ is a strictly monotone function of $X$ (its derivative equals $(w_2 s_0 - w_0 s_2)/(s_0+s_2 X)^2$ and $s_0,s_2>0$). Thus the minimizers of $R_{\rm quad}$ are exactly the minimizers (or maximizers, depending on the sign of $w_2 s_0 - w_0 s_2$) of $X[\phi]$ under the constraints $\phi\in\mathcal H_{\rm odd}(\sigma)$.

\medskip
\textit{Step 3 (Variational problem at fixed time variance).}
Minimize (or maximize) $X[\phi]$ subject to $\|\phi\|_2=1$, $\int t^2|\phi|^2=\sigma^2$, and oddness. Consider the constrained functional
\[
\mathcal L[\phi]=\int |\,\phi'|^2 + \alpha\!\left(\int |\phi|^2 - 1\right)
+ \beta\!\left(\int t^2|\phi|^2 - \sigma^2\right),
\]
with Lagrange multipliers $\alpha,\beta\in\mathbb R$. The Euler–Lagrange equation is
\[
-\,\phi''(t)\;+\;\beta\,t^2\,\phi(t)\;=\;\lambda\,\phi(t),\qquad \lambda:=-\alpha,
\]
which is the 1D harmonic-oscillator eigenvalue equation. Its $L^2$-solutions are Hermite functions; oddness selects the first excited state,
\[
\phi(t)\ \propto\ t\,e^{-\frac{\sqrt{\beta}}{2}\,t^2}.
\]
Imposing the variance constraint fixes $\sqrt{\beta}=\sigma^{-2}$, hence
\[
\phi^\star_\sigma(t)\;=\;c_\sigma\,\frac{t}{\sigma^2}\,e^{-\frac{t^2}{2\sigma^2}},
\]
with $c_\sigma$ chosen for $\|\phi^\star_\sigma\|_2=1$. This profile uniquely minimizes (and also maximizes, respectively, in the opposite monotonicity case) $X[\phi]$ in the odd class at fixed variance; uniqueness holds up to an overall phase and a shift of the origin (shifting the odd axis merely relocates the node).

\medskip
\textit{Step 4 (Conclusion).}
By Step 2, the extremizers of $R_{\rm quad}$ coincide with the extremizers of $X[\phi]$ on $\mathcal H_{\rm odd}(\sigma)$, hence with the DoG profile above. The $O(\omega^4)$ correction in Step 1 does not change the minimizer because it is dominated uniformly in the admissible class by the quadratic terms (the admissible set is compact in the natural Sobolev topology induced by the constraints). Therefore the minimizer of $\mathcal R_t$ is $\phi^\star_\sigma$ as claimed.
\end{proof}


\paragraph{Spatial shape optimality (Gaussian on $S^3$/LoG on$\Sigma_{\rm sp}$).}\label{par:uom-main-060}
Let $g_{\sigma}(\Omega)$ be the zonal heat kernel at time $\sigma^2$ on $S^3$,
with harmonic coefficients $\widehat g_\sigma(\ell)\propto e^{-\frac12\sigma^2\lambda_\ell}$.

\begin{lemma}[Micro spatial Gaussian extremizer]
\label{lem:spatial-gauss}
Let $-\Delta_{S^3}\,Y_{\ell 0}=\lambda_\ell\,Y_{\ell 0}$ with $\lambda_\ell=\ell(\ell+2)$ and let a zonal envelope $\Psi$ have spectral coefficients $\Psi_\ell:=\langle \Psi,Y_{\ell 0}\rangle$.
Assume the (operator) multipliers $s_x(\lambda),w_x(\lambda)$ are $C^1$ on $[0,\infty)$, strictly positive, and strictly decreasing in $\lambda$ on the spectral support. Fix the two quadratic constraints
\[
\sum_{\ell\ge 0} |\Psi_\ell|^2 = 1,\qquad
\sum_{\ell\ge 0} \lambda_\ell\,|\Psi_\ell|^2 = \Lambda_2\quad(\Lambda_2>0).
\]
Then the Rayleigh quotient
\[
\mathcal R_x[\Psi]\;:=\;\frac{\sum_{\ell\ge 0} w_x(\lambda_\ell)\,|\Psi_\ell|^2}{\sum_{\ell\ge 0} s_x(\lambda_\ell)\,|\Psi_\ell|^2}
\]
is minimized uniquely (up to an overall phase and a spatial rotation of the center) by the \emph{zonal heat kernel} (spherical Gaussian)
\[
g_{\sigma_x}\ \ \Longleftrightarrow\ \ \Psi_\ell\ \propto\ e^{-\frac12\sigma_x^2\,\lambda_\ell}\quad(\ell\ge 0),
\]
with $\sigma_x>0$ uniquely fixed by $\Lambda_2=\sum_\ell \lambda_\ell e^{-\sigma_x^2\lambda_\ell}/\sum_\ell e^{-\sigma_x^2\lambda_\ell}$.
\end{lemma}

\begin{proof}
\textit{Step 1 (Fractional programming $\Rightarrow$ linear surrogate).}
For any $\kappa\in\mathbb R$ define
\[
\Phi_\kappa[\Psi]\;:=\;\sum_{\ell\ge 0}\bigl(w_x(\lambda_\ell)-\kappa\,s_x(\lambda_\ell)\bigr)\,|\Psi_\ell|^2 .
\]
By Dinkelbach’s transform, $\kappa^\star:=\inf_{\Psi}\mathcal R_x[\Psi]$ is characterized by
\[
\inf_{\Psi}\ \Phi_{\kappa^\star}[\Psi]\quad\text{subject to}\quad \sum |\Psi_\ell|^2=1,\ \sum \lambda_\ell |\Psi_\ell|^2=\Lambda_2
\]
and the infimum equals $0$ at the optimizer.

\medskip
\textit{Step 2 (Lagrange multiplier system in the spectral variables).}
Introduce multipliers $\mu,\nu$ for the two constraints and consider
\[
\mathcal L(\{\Psi_\ell\};\mu,\nu)
=\sum_\ell \Big(w_x(\lambda_\ell)-\kappa s_x(\lambda_\ell)+\mu+\nu\,\lambda_\ell\Big)\,|\Psi_\ell|^2.
\]
Stationarity w.r.t.\ $|\Psi_\ell|^2$ (and positivity) implies:
\[
|\Psi_\ell|^2>0\quad\Rightarrow\quad
w_x(\lambda_\ell)-\kappa s_x(\lambda_\ell)+\mu+\nu\,\lambda_\ell=0,
\]
while if the left-hand side is positive, the coefficient must vanish. Because $w_x$ and $s_x$ are $C^1$ and strictly decreasing, the map
\[
\lambda\ \longmapsto\ w_x(\lambda)-\kappa s_x(\lambda)
\]
is strictly decreasing. Therefore the equation
\(
w_x(\lambda)-\kappa s_x(\lambda)+\mu+\nu\lambda=0
\)
can hold on a \emph{set of $\lambda$ with positive measure} only if the left-hand side is an affine function of $\lambda$; that happens precisely when the density $|\Psi_\ell|^2$ spans a continuum of $\lambda$ and the Kuhn–Tucker conditions are realized as a \emph{Gibbs law}.

\medskip
\textit{Step 3 (Gibbs/exponential family under moment constraints).}
To make this precise, relax the discrete $\ell$ to a continuum $\lambda\in[0,\infty)$ (large cutoff $L$ regime) and consider probability densities $\rho(\lambda)\ge0$ with $\int\rho=1$ and $\int \lambda\rho=\Lambda_2$. Minimizing $\int \big(w_x-\kappa s_x\big)\rho$ under these two moment constraints and a small entropic regulator $\epsilon\int \rho\log\rho$ yields, by standard convex duality, the unique solution
\[
\rho_\epsilon(\lambda)\;=\;\frac{1}{Z_\epsilon}\exp\!\big(-\beta_\epsilon\,\lambda\big),
\]
with $\beta_\epsilon$ fixed by $\int \lambda\rho_\epsilon=\Lambda_2$. Sending $\epsilon\downarrow 0$ preserves the exponential form (the feasible set is compact and the objective is continuous), hence the unregularized minimizer is the \emph{exponential density} $\rho(\lambda)\propto e^{-\beta\lambda}$. Pulling back to the discrete spectrum on $S^3$ gives
\[
|\Psi_\ell|^2\ \propto\ e^{-\beta\,\lambda_\ell}\qquad(\ell\ge 0).
\]
But the zonal spherical heat kernel at time $\tau=\sigma_x^2/2$ has precisely these weights:
\(
\widehat g_{\sigma_x}(\ell)\propto e^{-\frac12\sigma_x^2\lambda_\ell}.
\)

\medskip
\textit{Step 4 (Uniqueness and determination of $\sigma_x$).}
Strict convexity of the constrained problem (the feasible set is a simplex cut by a strictly convex moment map, and the objective is strictly convex along feasible directions thanks to strict monotonicity of $w_x,s_x$) implies uniqueness of the weights $\{|\Psi_\ell|^2\}$, hence uniqueness (up to phase and rotation) of the minimizer. The parameter $\sigma_x$ is fixed by the second-moment constraint:
\[
\Lambda_2\;=\;\frac{\sum_{\ell\ge 0}\lambda_\ell\,e^{-\sigma_x^2\lambda_\ell}}{\sum_{\ell\ge 0}e^{-\sigma_x^2\lambda_\ell}}.
\]
Therefore $\Psi$ minimizing $\mathcal R_x$ is the zonal heat kernel $g_{\sigma_x}$.
\end{proof}

\begin{lemma}[Macro spatial LoG extremizer on the static--patch slice (monopole--free)]
\label{lem:spatial-log-static}
Let $(\Sigma_{\rm sp},\gamma)$ be a static--patch spatial slice with Riemannian metric $\gamma$ and
self--adjoint Laplace--Beltrami operator $-\Delta_{\Sigma_{\rm sp}}$ on $L^2(\Sigma_{\rm sp},d\mathrm{vol}_\gamma)$
with domain chosen so that the heat semigroup $e^{-s(-\Delta_{\Sigma_{\rm sp}})}$ exists for $s>0$ and admits a
symmetric, strictly positive heat kernel $K_{\rm sp}(s;x,y)$ (e.g.\ Dirichlet or Neumann conditions at the
static horizon suffice). Fix a center $x_0\in\Sigma_{\rm sp}$ and define the class of \emph{zonal} envelopes
centered at $x_0$ by
\[
\mathcal Z_{x_0}\ :=\ \{\,\Phi\in L^2(\Sigma_{\rm sp}):\ \Phi\ \text{depends only on }d_\gamma(\,\cdot\,,x_0)\,\}.
\]
Consider the admissible set
\[
\mathcal A\ :=\ \Big\{\,\Phi\in\mathcal Z_{x_0}\cap L^2(\Sigma_{\rm sp})\ :\ \int_{\Sigma_{\rm sp}}\Phi\,d\mathrm{vol}_\gamma=0,\ \ \|\Phi\|_{L^2}=1,\ \
\int_0^\infty \lambda\,d\mu_\Phi(\lambda)=\Lambda_2\ (>0)\,\Big\},
\]
where $\mu_\Phi$ is the scalar spectral measure of $-\Delta_{\Sigma_{\rm sp}}$ induced by $\Phi$, i.e.
$\langle \Phi, f(-\Delta_{\Sigma_{\rm sp}})\,\Phi\rangle=\int_0^\infty f(\lambda)\,d\mu_\Phi(\lambda)$
for all bounded Borel $f$. Let $s_x,w_x:[0,\infty)\to(0,\infty)$ be $C^1$, strictly decreasing functions.
Define the Rayleigh quotient
\[
\mathcal R_x[\Phi]\;:=\;\frac{\big\langle \Phi,\,w_x(-\Delta_{\Sigma_{\rm sp}})\,\Phi\big\rangle}{\big\langle \Phi,\,s_x(-\Delta_{\Sigma_{\rm sp}})\,\Phi\big\rangle}
\;=\;\frac{\int_0^\infty w_x(\lambda)\,d\mu_\Phi(\lambda)}{\int_0^\infty s_x(\lambda)\,d\mu_\Phi(\lambda)}\;.
\]
Then $\mathcal R_x$ has a unique minimizer in $\mathcal A$, up to the trivial phase, given by the (normalized)
\emph{Laplacian--of--Gaussian} (LoG)
\[
\Phi_\rho(\,\cdot\,)\;=\;\frac{\Delta_{\Sigma_{\rm sp}} K_{\rm sp}(\rho^2;\,\cdot\,,x_0)}{\big\|\Delta_{\Sigma_{\rm sp}} K_{\rm sp}(\rho^2;\,\cdot\,,x_0)\big\|_{L^2}}\ ,
\]
for a unique scale $\rho>0$ fixed by the second--moment constraint. In particular, in Riemann normal
coordinates around $x_0$ one has
\[
\Phi_\rho(x)\;=\;c(\rho)\,\Delta_{\mathbb R^3}\exp\!\Big(-\frac{d_\gamma(x,x_0)^2}{2\rho^2}\Big)\ +\ O\!\Big(\frac{\rho^2}{\ell^2}\Big),
\]
where $\ell$ is the curvature radius of $(\Sigma_{\rm sp},\gamma)$ and $c(\rho)$ is chosen so that
$\|\Phi_\rho\|_{L^2}=1$.
\end{lemma}

\begin{proof}
\emph{Step 1 (Spectral representation and constraints).}
By the spectral theorem for $-\Delta_{\Sigma_{\rm sp}}$, there exists a projection--valued measure
$E(d\lambda)$ on $[0,\infty)$ such that $f(-\Delta_{\Sigma_{\rm sp}})=\int_0^\infty f(\lambda)\,E(d\lambda)$
for all bounded Borel $f$. For $\Phi\in L^2$, define $\mu_\Phi(B):=\langle \Phi,E(B)\Phi\rangle$; then
$\mu_\Phi$ is a finite positive measure on $[0,\infty)$ with $\mu_\Phi([0,\infty))=\|\Phi\|_{L^2}^2$ and
$\langle \Phi,f(-\Delta_{\Sigma_{\rm sp}})\Phi\rangle=\int f(\lambda)\,d\mu_\Phi(\lambda)$.

The constraints $\|\Phi\|_{L^2}=1$ and $\int \Phi\,d\mathrm{vol}_\gamma=0$ translate into
\[
\int_0^\infty d\mu_\Phi(\lambda)=1,\qquad \mu_\Phi(\{0\})=0,
\]
since the constant function is the (normalized) eigenfunction at $\lambda=0$ on the static patch, and
the mean--zero constraint enforces orthogonality to that eigenspace. The fixed spectral second moment is
\[
\int_0^\infty \lambda\,d\mu_\Phi(\lambda)=\Lambda_2\ > 0.
\]
Thus, minimizing $\mathcal R_x$ over $\mathcal A$ is equivalent to minimizing the ratio
$\int w_x\,d\mu / \int s_x\,d\mu$ over finite positive measures $\mu$ on $(0,\infty)$ (monopole suppressed)
subject to the two linear moment constraints
\(
\int d\mu=1,\ \ \int \lambda\,d\mu=\Lambda_2.
\)

\smallskip
\emph{Step 2 (Fractional programming $\Rightarrow$ linear surrogate).}
Let $\mathcal M$ denote the convex set of admissible measures and define, for $\kappa\in\mathbb R$,
\[
\Phi_\kappa(\mu)\;:=\;\int_0^\infty \bigl(w_x(\lambda)-\kappa\,s_x(\lambda)\bigr)\,d\mu(\lambda).
\]
By a standard Dinkelbach transform (see, e.g., fractional programming in measure spaces), the minimal value
$\kappa^\star:=\inf_{\mu\in\mathcal M}\ \mathcal R_x[\mu]$ is characterized by
\begin{equation}
\label{eq:dink}
\Phi_{\kappa^\star}(\mu)\ \ge\ 0\quad \text{for all }\mu\in\mathcal M,\qquad \inf_{\mu\in\mathcal M}\ \Phi_{\kappa^\star}(\mu)\;=\;0,
\end{equation}
and any minimizer of $\mathcal R_x$ is a minimizer of $\Phi_{\kappa^\star}$ (and conversely).

\smallskip
\emph{Step 3 (Duality and KKT conditions).}
Consider the Lagrangian
\[
\mathcal L(\mu;\alpha,\beta,\kappa)\;=\;\int_{(0,\infty)}\!\!\!\bigl(w_x(\lambda)-\kappa\,s_x(\lambda)\bigr)\,d\mu(\lambda)
+\alpha\Big(\int d\mu - 1\Big)+\beta\Big(\int \lambda\,d\mu-\Lambda_2\Big),
\]
with multipliers $\alpha,\beta\in\mathbb R$. For fixed $(\alpha,\beta,\kappa)$, minimizing $\mathcal L$
over $\mu\ge 0$ yields the pointwise Karush--Kuhn--Tucker (KKT) condition
\begin{equation}
\label{eq:kkt}
\bigl(w_x(\lambda)-\kappa\,s_x(\lambda)+\alpha+\beta\,\lambda\bigr)\ \ge\ 0
\quad\text{for all }\lambda>0,\qquad
\bigl(w_x(\lambda)-\kappa\,s_x(\lambda)+\alpha+\beta\,\lambda\bigr)\,d\mu(\lambda)=0.
\end{equation}
Since $w_x$ and $s_x$ are $C^1$ and strictly decreasing, the map
\(
\lambda\mapsto w_x(\lambda)-\kappa\,s_x(\lambda)+\alpha+\beta\,\lambda
\)
is strictly \emph{increasing} whenever $\beta>0$, and strictly \emph{decreasing} whenever $\beta<0$.
Because $\mu$ must have finite nonzero first moment $\Lambda_2$, the minimizer has $\beta<0$; hence
the integrand in \eqref{eq:kkt} is strictly decreasing, and the support of any minimizing $\mu$ is a
connected set on which the integrand is constant (equal to $0$). To rule out singular measures and ensure
uniqueness, we employ an entropic regularization (standard in convex duality): for $\epsilon>0$, let
\[
\mathcal L_\epsilon(\rho;\alpha,\beta,\kappa)\;=\;\int_{0}^{\infty}\!\Big(w_x-\kappa s_x + \alpha + \beta \lambda\Big)\rho(\lambda)\,d\lambda
\ +\ \epsilon\int_{0}^{\infty}\rho(\lambda)\log \rho(\lambda)\,d\lambda,
\]
where we restrict to absolutely continuous measures $d\mu=\rho(\lambda)\,d\lambda$ on $(0,\infty)$.
Minimizing over $\rho\ge 0$ yields the unique solution
\[
\rho_{\epsilon}(\lambda)\;=\;\frac{1}{Z_\epsilon}\,\exp\!\Big(-\frac{1}{\epsilon}\big(\alpha+\beta\lambda + w_x(\lambda)-\kappa s_x(\lambda)\big)\Big),
\qquad Z_\epsilon>0,
\]
whose parameters $(\alpha,\beta)$ are fixed by the two moment constraints. Sending $\epsilon\downarrow 0$
(Helly’s selection plus lower semicontinuity of the cost) produces an absolutely continuous limit
$d\mu^\star(\lambda)=\rho^\star(\lambda)\,d\lambda$ supported where $w_x-\kappa s_x+\alpha+\beta\lambda=0$,
and hence of the \emph{exponential family} form $\rho^\star(\lambda)\propto \exp(-\beta \lambda)$.
The monopole suppression enforces $\mu^\star(\{0\})=0$, which in the absolutely continuous case amounts to
no atom at the origin; in particular, any weak limit of $\rho_\epsilon$ satisfies $\rho^\star(0)=0$.

Thus, the unique minimizer of $\mathcal R_x$ has density of the form
\begin{equation}
\label{eq:opt-density}
d\mu^\star(\lambda)\ =\ \zeta(\beta)\,e^{-\beta \lambda}\,d\lambda,\qquad \beta>0,
\end{equation}
with $\beta$ fixed by $\int \lambda\,d\mu^\star(\lambda)=\Lambda_2$ and $\zeta(\beta)$ fixed by normalization.

\smallskip
\emph{Step 4 (Realization by a LoG and uniqueness).}
We claim that $\Phi_\rho:=\Delta_{\Sigma_{\rm sp}} K_{\rm sp}(\rho^2;\cdot,x_0)$ realizes precisely the optimal
density \eqref{eq:opt-density} in the sense that its spectral measure has the form $d\mu_{\Phi_\rho}(\lambda)
= c(\rho)\,e^{-\rho^2\lambda}\,d\lambda$ on $(0,\infty)$, for some $c(\rho)>0$ (after restricting to the cyclic
subspace generated by $x_0$). Indeed, write
\[
K_{\rm sp}(\rho^2;\cdot,x_0)\;=\;e^{-\rho^2(-\Delta_{\Sigma_{\rm sp}})}\,\delta_{x_0}
\;=\;\int_0^\infty e^{-\rho^2\lambda}\,E(d\lambda)\,\delta_{x_0},
\]
in the distributional sense (the heat semigroup applied to $\delta_{x_0}$). Applying $\Delta_{\Sigma_{\rm sp}}$
multiplies spectral components by $\lambda$:
\[
\Phi_\rho\;=\;\Delta_{\Sigma_{\rm sp}} K_{\rm sp}(\rho^2;\cdot,x_0)
\;=\;\int_0^\infty \lambda\,e^{-\rho^2\lambda}\,E(d\lambda)\,\delta_{x_0}.
\]
Hence, for any bounded $f$,
\[
\big\langle \Phi_\rho,\ f(-\Delta_{\Sigma_{\rm sp}})\,\Phi_\rho\big\rangle
\;=\;\int_0^\infty f(\lambda)\,\lambda^2\,e^{-2\rho^2\lambda}\,d\mu_{x_0}(\lambda),
\]
where $d\mu_{x_0}(\lambda)=\langle \delta_{x_0},E(d\lambda)\delta_{x_0}\rangle$ is the cyclic spectral
measure at $x_0$. Because $d\mu_{x_0}$ is absolutely continuous near $\lambda=0$ (by standard elliptic
regularity; in particular, there is no atom at $\lambda=0$ except for the constant mode, which is removed
by our mean--zero constraint), it follows that $d\mu_{\Phi_\rho}(\lambda)=|\widehat\Phi_\rho(\lambda)|^2\,d\lambda$
is absolutely continuous with density proportional to $e^{-2\rho^2\lambda}$ on $(0,\infty)$ after normalization
(the extra $\lambda^2$ factor is absorbed into the equivalent $d\lambda$-based description because $\lambda^2$ is
strictly positive on $(0,\infty)$ and plays no role in determining the exponential rate). Thus $\Phi_\rho$ has
the exponential spectral density \eqref{eq:opt-density} with $\beta=2\rho^2$, and the second--moment constraint
$\int \lambda\,d\mu_{\Phi_\rho}(\lambda)=\Lambda_2$ fixes $\rho$ uniquely (strict monotonicity of the map
$\rho\mapsto \int \lambda e^{-2\rho^2 \lambda}d\lambda$ on $(0,\infty)$).

Finally, strict convexity of the map $\mu\mapsto \int (w_x-\kappa s_x)\,d\mu$ on $\mathcal M$ and the strict
monotonicity of $w_x,s_x$ imply uniqueness of the minimizing measure (hence of the minimizing $\Phi$, up to
phase and center). Therefore $\Phi_\rho$ is the unique minimizer of $\mathcal R_x$ in $\mathcal A$, and the
Riemann--normal expansion yields the stated local flat asymptotic.
\end{proof}


\begin{corollary}[Welded LoG optimality (time and space)]
\label{cor:welded-log}
With $\phi^\star_\sigma$ from Lemma~\ref{lem:dog-opt}, $g_{\sigma_x}$ from
Lemma~\ref{lem:spatial-gauss}, and $\Phi_\rho$ from Lemma~\ref{lem:spatial-log-static},
the welded envelopes are
\[
\big(g'_{\Sigma}{*}\,g'_{\sigma_t}\big)(t)\;=\;\tfrac{d^2}{dt^2}g_{\sqrt{\Sigma^2+\sigma_t^2}}(t),
\qquad
\Phi_\rho\ \circledast\ g_{\sigma_x}\;=\;\Delta_{S^3}\,g_{\sqrt{\rho^2+\sigma_x^2}}(\Omega),
\]
i.e.\ \emph{LoG} in both time and space with scales adding in quadrature. This pair
minimizes the source Rayleigh quotient and hence renders the source action
$\Phi^{\rm src}=\Delta I^\star\lambda_t\lambda_x$ tickwise constant.
\end{corollary}

\subsubsection*{Pulses co-centering}\label{subsubsec:uom-main-009}

\paragraph{Macro–micro co–centering in time.}\label{par:uom-main-061}
Let \(t_{\rm mac}\) and \(t_{\rm mic}\) be macro/micro temporal centers, and define the relative shift \(\delta t:=t_{\rm mic}-t_{\rm mac}\).
For fixed welded scale \(s_t\), the receiver term is time–translation invariant and independent of \(\delta t\). The source term in the linear–response/QND regime is a quadratic functional
\(
\mathcal A^{\rm src}_k=\iint \chi_k(t)\,\mathsf K(t-t')\,\chi_k(t')\,dt\,dt'
\)
with an even kernel \(\mathsf K(\tau)=\mathsf K(-\tau)\) (KMS bath). Writing \(\chi_k(t)=A_{\rm mac}\,g'_{\Sigma}(t-t_{\rm mac})\) and expanding around the tick center \(t_k\), one finds that \(\mathcal A^{\rm src}_k\) is an even, strictly convex function of \(t_{\rm mac}-t_k\). Because the welded envelope center is \(t_{\rm mac}+t_{\rm mic}\) and we take \(t_k\) as the acceptance center, the total per–tick action is an even convex function of \(\delta t\) with unique minimum at \(\delta t=0\).
Hence:
\begin{equation}
\boxed{\ \text{Macro–micro temporal co–centering holds at the minimizer:}\quad t_{\rm mac}=t_{\rm mic}=t_k\ .\ }
\tag{C\textsubscript{t}}\label{eq:Ct}
\end{equation}

\paragraph{Macro–micro spatial co–centering on the static–patch slice.}\label{par:uom-main-062}
Let $(\Sigma_{\rm sp},\gamma)$ be a static–patch spatial slice with Laplace–Beltrami operator
$-\Delta_{\Sigma_{\rm sp}}$ and heat kernel $K_{\rm sp}(s;x,y)$ (so $\partial_s K_{\rm sp}=\Delta_y K_{\rm sp}$).
Fix scales $(\rho,\sigma)>0$. Define the zonal macro LoG and micro Gaussian about centers
$x,y\in\Sigma_{\rm sp}$ by
\[
\mathrm{LoG}_\rho(\cdot;x)\ :=\ \Delta_{\Sigma_{\rm sp}}\,K_{\rm sp}(\rho^2;\,\cdot\,,x),\qquad
G_\sigma(\cdot;y)\ :=\ K_{\rm sp}(\sigma^2;\,\cdot\,,y).
\]
Let the welded profile on the patch be the zonal convolution
\(
\tau_{x,y}:=\mathrm{LoG}_\rho(\cdot;x)\ast G_\sigma(\cdot;y)
\)
and define the (patch–wise) correlation functional
\begin{equation}
\label{eq:spatial-corr}
\mathcal C(x,y)\;:=\;\big\langle \mathrm{LoG}_\rho(\cdot;x),\ G_\sigma(\cdot;y)\big\rangle_{L^2(\Sigma_{\rm sp})}
\;=\;\int_{\Sigma_{\rm sp}}\Delta_{z}K_{\rm sp}(\rho^2;x,z)\,K_{\rm sp}(\sigma^2;y,z)\,d\mathrm{vol}_\gamma(z).
\end{equation}

\begin{lemma}[Macro–micro spatial co–centering on $\Sigma_{\rm sp}$]
\label{lem:spatial-cocenter-static}
For each fixed $(\rho,\sigma)$, the map $(x,y)\mapsto \mathcal C(x,y)$ depends only on the geodesic
distance $d_\gamma(x,y)$ and is strictly decreasing in $d_\gamma(x,y)$ on $(0,\mathrm{inj\,rad})$.
Equivalently, the welded action is uniquely minimized at co–centering $x=y$.
\end{lemma}

\begin{proof}
By symmetry of $K_{\rm sp}$ and integration by parts,
\[
\mathcal C(x,y)\;=\;\int_{\Sigma_{\rm sp}}K_{\rm sp}(\rho^2;x,z)\,\Delta_{z}K_{\rm sp}(\sigma^2;y,z)\,d\mathrm{vol}_\gamma(z)
\;=\;\partial_{\sigma^2}\!\int_{\Sigma_{\rm sp}} K_{\rm sp}(\rho^2;x,z)\,K_{\rm sp}(\sigma^2;y,z)\,d\mathrm{vol}_\gamma(z).
\]
By the heat–semigroup property,
\(
\int K_{\rm sp}(\rho^2;x,z)\,K_{\rm sp}(\sigma^2;y,z)\,dz=K_{\rm sp}(\rho^2+\sigma^2;x,y).
\)
Hence
\begin{equation}
\label{eq:kernel-derivative}
\mathcal C(x,y)\;=\;\partial_{\sigma^2}\,K_{\rm sp}(\rho^2+\sigma^2;\,x,y).
\end{equation}
The heat kernel on a Riemannian manifold is a strictly positive, radially decreasing function of
$d_\gamma(x,y)$ for each fixed time; its time–derivative preserves this strict radial monotonicity
for $(\rho^2+\sigma^2)$ in the parabolic regime (see standard heat kernel estimates and the parabolic
maximum principle). Therefore $\mathcal C(x,y)$ is a strictly decreasing function of the geodesic
distance $d_\gamma(x,y)$, and it is uniquely maximized at $x=y$. Since the welded quadratic action is,
up to constants, $\|\mathrm{LoG}_\rho\|^2+\|G_\sigma\|^2-2\mathcal C(x,y)$, it is uniquely minimized at $x=y$.
\end{proof}

\begin{equation}
\boxed{\ \text{Macro--micro spatial co--centering holds at the minimizer:}\quad x=y\ .\ }
\tag{C\textsubscript{x}}\label{eq:Cx}
\end{equation}

\noindent
In particular, in Riemann normal coordinates around the macro center, the conclusion reduces to the
classical flat–space fact that the correlation of a Laplacian–of–Gaussian with a Gaussian is strictly
peaked at zero shift. The same argument applies patch–wise on the static–patch union.

\subsubsection*{Welded ticks spatial co–centering lock across aeon.}
\label{subsubsec:welded-co-centering}

\begin{proposition}[History--induced spatial center locking]
\label{prop:history-center-lock}
Let $\Sigma$ be the receiver spatial slice (round $S^3$) with Laplace--Beltrami $-\Delta_\Sigma$
and heat kernel $K_\Sigma(s;\cdot,\cdot)$. Suppose that after $(k-1)$ ticks the receiver trace sector
carries a nonzero \emph{zonal LoG background}
\[
B(\Omega)\;=\;\Delta_\Sigma\,K_\Sigma(\eta^2;\,R_{\Omega_\star}^{-1}\Omega)\,,
\]
centered at $\Omega_\star\in S^3$, for some $\eta>0$ (the welded BY imprint of the history).
Let the welded spatial profile of the next tick be
\[
\tau_{\rho,\sigma}(\,\cdot\,;\Omega)\;:=\;\big(\Delta_\Sigma\,K_\Sigma(\rho^2;\,\cdot\, ,\Omega)\big)\ \ast\ K_\Sigma(\sigma^2;\,\cdot\,,\Omega),
\]
with fixed $(\rho,\sigma)>0$ (macro LoG scale $\rho$, micro Gaussian scale $\sigma$), centered at $\Omega\in S^3$.
Let $\mathcal K$ denote any real, $SO(4)$--invariant, positive semidefinite receiver kernel (e.g. the Fisher
or screened canonical kernel entering the quadratic part of the P3 receiver modular energy functional). Then the receiver
quadratic action around the background contains the cross term
\[
\mathcal C(\Omega)\;:=\;\big\langle \tau_{\rho,\sigma}(\,\cdot\,;\Omega),\ \mathcal K * B\big\rangle_{L^2(S^3)}\,,
\]
and $\Omega\mapsto \mathcal C(\Omega)$ depends only on the geodesic distance $\delta:=\mathrm{dist}_{S^3}(\Omega,\Omega_\star)$
and is \emph{strictly decreasing} in $\delta$ on $(0,\pi)$. Consequently, the (receiver--side) action for the $k$--th tick
is \emph{uniquely} minimized by centering $\tau_{\rho,\sigma}$ at $\Omega=\Omega_\star$, and the dependence on $\Omega$ is
strictly convex in any normal chart around $\Omega_\star$.
\end{proposition}

\begin{proof}
\emph{Step 1 (Radialization of the cross term).}
By $SO(4)$--invariance of $\mathcal K$ and the zonality of $B$, the function $\mathcal K*B$ is also zonal about $\Omega_\star$.
Therefore the map $\Omega\mapsto \mathcal C(\Omega)$ depends only on $\delta=\mathrm{dist}_{S^3}(\Omega,\Omega_\star)$.

\smallskip
\emph{Step 2 (Monotonicity via heat--kernel semigroup).}
Write $\mathrm{LoG}_\rho(\cdot;\Omega):=\Delta_\Sigma K_\Sigma(\rho^2;\cdot,\Omega)$ and $G_\sigma(\cdot;\Omega):=K_\Sigma(\sigma^2;\cdot,\Omega)$,
so $\tau_{\rho,\sigma}(\cdot;\Omega)=\mathrm{LoG}_\rho(\cdot;\Omega)\ast G_\sigma(\cdot;\Omega)$. By symmetry and integration by parts,
\[
\mathcal C(\Omega)\;=\;\big\langle \mathrm{LoG}_\rho(\cdot;\Omega),\ G_\sigma(\cdot;\Omega)\ast(\mathcal K*B)\big\rangle.
\]
Using that $\mathcal K$ is $SO(4)$--invariant and positive semidefinite, its action on a zonal LoG equals a radial convolution, and
the heat--semigroup property implies
\[
G_\sigma(\cdot;\Omega)\ast (\mathcal K*B)\;=\;\widetilde{\mathcal K}\ast K_\Sigma(\sigma^2+\eta^2;\cdot,\Omega_\star),
\]
for some radial convolution kernel $\widetilde{\mathcal K}$ (the precise spectral form is not needed).
It follows that
\[
\mathcal C(\Omega)\;=\;\partial_{\rho^2}\,\partial_{\eta^2}\,\Big(\widetilde{\mathcal K}\ast K_\Sigma(\rho^2+\sigma^2+\eta^2;\cdot,\cdot)\Big)\Big|_{(\cdot,\cdot)=(\Omega,\Omega_\star)}.
\]
Since $K_\Sigma(\cdot;\Omega,\Omega_\star)$ is strictly positive and strictly decreasing in the geodesic distance for each fixed time,
and radial convolution with a positive kernel preserves strict radial monotonicity, the map
$\delta\mapsto \mathcal C(\Omega)$ is strictly decreasing on $(0,\pi)$.

\smallskip
\emph{Step 3 (Uniqueness and convexity).}
Strict decrease in $\delta$ implies the unique global maximum of $\mathcal C$ is at $\Omega=\Omega_\star$, hence the quadratic action
(which contains a center--independent self--term minus $2\,\mathcal C(\Omega)$ as the cross term) is uniquely minimized at $\Omega_\star$.
Smoothness of $K_\Sigma$ and $\widetilde{\mathcal K}$ in normal coordinates implies the Hessian of $-\mathcal C$ at $\Omega_\star$ is
positive definite, giving strict convexity in $\Omega$ near $\Omega_\star$.
\end{proof}

\noindent
\emph{Co-centered ticks: conclusion.}
Lemma~\ref{lem:spatial-cocenter-static} establishes that, for each fixed tick and fixed scales $(\rho,\sigma_x)$,
\emph{any} relative spatial displacement between the macro and micro centers injects $m\neq 0$ components
outside the rank–one LoG manifold at every $\ell$, and the Fisher–quadratic penalty of that orthogonal
complement strictly increases. Hence \emph{per tick} the welded spatial envelope is uniquely minimized by
\emph{co–centering} the macro and micro profiles. By $SO(4)$–homogeneity of the round $S^3$, the \emph{absolute}
location of that common center is undetermined in the perfectly symmetric setting; one may fix it by a global
rotation (e.g.\ choose the brane reference origin) without cost.

Crucially, once the receiver trace sector carries a nonzero \emph{zonal LoG background} produced by the BY
flux of prior pulses, the next tick’s welded envelope acquires a center–dependent \emph{cross term} in the
P3 receiver-side action functional. By Proposition~\ref{prop:history-center-lock}, this cross term is a strictly decreasing
function of the geodesic separation from the background center, so the action is \emph{uniquely} minimized
when the new pulse is centered \emph{at} that background center. Consequently, after the first nonzero BY imprint,
all subsequent welded pulses are not only co–centered within each tick but also \emph{share the same absolute
center}. Equivalently, from that tick onward the sequence of pulses selects a single spatial axis, and by
homogeneity we may place its apex at the brane reference origin.

(In the perfectly symmetric, history–free limit the common center remains a flat direction: it can be fixed
\emph{by convention} but is not selected dynamically. The presence of any nonzero BY imprint removes this
degeneracy and locks the center uniquely by strict convexity in a normal chart around the background point. The allowable anisotropy origin extends to other branes as well, as shown in Proposition~\ref{prop:aeon-center-inheritance}.)


\subsubsection*{Practical optimization flow}\label{subsubsec:uom-main-010}
The simple algorithm of the one-tick separable minimization:
\begin{enumerate}[leftmargin=1.4em,itemsep=2pt]
\item \emph{Receiver stage:} minimize \(\Phi^{\rm rec}_k(s_{t},s_{x},\Lambda)\) to get \((s_{t,k}^\star,s_{x,k}^\star,\Lambda_k^\star)\). (Convex/GP; Eqs.~\eqref{eq:R2} and Sec.~\ref{subsec:P3-master}.)
\item \emph{Center locks:} set \(t_{\rm mac}=t_{\rm mic}=t_k\) and \(\Omega_{\rm mac}=\Omega_{\rm mic}=\Omega_k\) by \eqref{eq:Ct}–\eqref{eq:Cx}.
\item \emph{Source stage:} choose \((\Sigma_k,\rho_k)\) by minimizing \(\mathcal A^{\rm src}_k\) under
\(
\Sigma_k^2+\sigma_{t,k}^2=s_{t,k}^{\star 2},\ \sigma_{t,k}\ge\sigma_{\rm QSL}
\)
and
\(
\rho_k^2+\sigma_{x,k}^2=s_{x,k}^{\star 2},\ \sigma_{x,k}\ge\sigma_{x,{\rm min}}.
\)
\end{enumerate}

\paragraph{Consequences.}\label{par:uom-main-063}
(i) The \emph{receiver} fully fixes the welded scales that enter the anomaly floor and thus the \(\Lambda_4\) evolution. (ii) The \emph{source} fixes how those welded scales are realized macro vs.\ micro, trading QSL and bath/actuator costs. (iii) Co–centering is a universal lock: any macro–micro misalignment only increases the action.



\paragraph{Payload equality and amplitude elimination.}\label{par:uom-main-064}
Let the welded throughput form $\mathsf S$ be fixed. For any admissible widths
$(\Sigma,\sigma_t,\sigma_x)$ and the extremal shapes above, the compiled payload obeys
\begin{equation}
\label{eq:payload-eq-final}
\Delta I^\star\;=\;\langle h,\ \mathsf S\,h\rangle
=\big(A^{\rm eff}\big)^2\,G(\Sigma,\sigma_t,\sigma_x;\Lambda),
\qquad
\Rightarrow\quad
\big(A^{\rm eff}\big)^2=\frac{\Delta I^\star}{G(\Sigma,\sigma_t,\sigma_x;\Lambda)}.
\end{equation}
Thus all amplitudes are eliminated in favor of the widths and cutoff.

\paragraph{Unconstrained per–tick optimization (receiver side only).}\label{par:uom-main-065}
With $\Phi^{\rm src}$ tickwise constant and amplitudes eliminated by
\eqref{eq:payload-eq-final}, the per–tick minimization is \emph{unconstrained} (apart from
zero–DC built into the shapes) on the receiver side:
\begin{equation}
\label{eq:unconstrained-J}
\min_{\Sigma,\sigma_t,\sigma_x,\Lambda}\ \ \Phi_{\rm BY}(\Sigma,\sigma_t,\sigma_x,\Lambda)
\;+\;\Phi_{\rm conf}(\Sigma,\sigma_t,\sigma_x,\Lambda),
\end{equation}
where, for the welded LoG family,
\[
\Phi_{\rm BY}
=\gamma(8\pi G_4)^2\,C_{\rm LoG}\,\frac{\Delta I^\star}{\sigma_t\,\sigma_x^{3}\,G(\Sigma,\sigma_t,\sigma_x;\Lambda)},
\quad
\Phi_{\rm conf}=A_{\rm eff}^2\,\Xi_{\rm conf}(\Sigma,\sigma_t,\sigma_x;\Lambda),,
\]
with $A_{\rm eff}^2$ replaced by $\Delta I^\star/G$ via \eqref{eq:payload-eq-final}.
Assuming $G$ is log–concave in $(\Sigma,\sigma_t,\sigma_x)$ and the projectors are
represented by convex quadratic surrogates on the LoG family, the objective in
\eqref{eq:unconstrained-J} is convex in $(\Sigma,\sigma_t,\sigma_x,\Lambda)$ and admits
an equivalent geometric–program form in log–coordinates. The unique minimizer exhibits
(i) co–centering (zero shift), (ii) the DoG/LoG locks proven above, and (iii) the
fast–then–slow, variable–rate geometric contraction of the BY anomaly floor as $\Lambda$
refines and the welded scales move along the convex trade–off surface.


\subsection{Bulk–mediated brane–brane coupling and shadow–brane phenomenology}
\label{subsec:shadow-brane}

\paragraph{Scope and setting.}\label{par:uom-main-066}
We consider purely \emph{gravitational} coupling through AdS$_5$ between two post–flip
de Sitter branes: our aeon (dS$_4$) and a \emph{shadow} dS$_4$ brane (a relic from the
previous aeon). All matter-carrying shadow branes are interior/shadow RS1 segments deeper IR
than the (already RS1) source; KR heads are vacuum EOW caps and do not carry dark--sector matter.
Both are \emph{decoupled from the absorber} and obey ordinary 4D dS dynamics on their worldvolumes.
Their UV windows $(\Lambda_1,\Lambda_2)$ were fixed at their respective flips and remain constant thereafter.
The receiver AdS$_4$ brane is a pure vacuum and plays no role in the dark--sector phenomenology below.

\medskip
\noindent\textbf{Cross–brane transfer kernel (unified form).}
Let $\rho_1(t,\mathbf x)$ and $\rho_2(t,\mathbf x)$ denote comoving matter densities on
our brane (1) and the shadow brane (2). The AdS$_5$ bulk induces a linear, retarded
map in harmonic space (time Fourier $\omega$ and $S^3$ angular momentum $\ell$):
\[
\delta\rho^{\rm eff}_1(\omega,\ell)\;=\;\varepsilon(\omega,\ell;\Lambda_1,\Lambda_2,d,\theta)\;
\delta\rho_2(\omega,\ell),
\]
with \vspace{-0.25em}
\begin{equation}
\label{eq:eps-factorization}
\boxed{%
\begin{aligned}
\varepsilon(\omega,\ell;\cdots)\;=\;&
\underbrace{\mathcal C(\theta)}_{\text{tilt/orientation}}\;
\underbrace{W_{\Lambda_1}(\ell)\,W_{\Lambda_2}(\ell)}_{\text{UV windows}}\;
\underbrace{e^{-\,d\,q(\omega,k_\ell)}}_{\text{bulk propagation}}\;
\underbrace{\frac{q(\omega,k_\ell)}{q(\omega,k_\ell)+r(\Lambda)\,k_\ell^2}}_{\text{induced gravity screen}}
\\[2pt]
&k_\ell:=\frac{\sqrt{\ell(\ell+2)}}{L_4},\qquad
q(\omega,k):=\sqrt{k^2+m_\bullet^2-\omega^2}\,.
\end{aligned}}
\end{equation}

Here $d$ is the proper separation of the brane embeddings, $\theta$ a small tilt
between their normals, $m_\bullet$ the relevant trace-channel scalar gap (the GW-stabilized RS1 radion mass);
the spin--2 gap $m_g$ enters only the TT sector and is not used in this scalar kernel. Here
$r(\Lambda)=M_4^2(\Lambda)/(2M_5^3)$ is the induced--gravity crossover
scale on \emph{our} brane at cutoff $\Lambda$, and $W_{\Lambda_i}(\ell)$ are smooth high--$\ell$
windows reflecting the fixed cutoffs $\Lambda_i$ at flip. Causality is automatic:
only the \emph{retarded} bulk propagator contributes; the constant windows merely
implement UV softness.

\emph{DC/long–wavelength limit.} For the background ($\omega=0$, $\ell=0$),
\begin{equation}
\label{eq:eps-DC}
\varepsilon_{\rm DC}\;=\;\varepsilon(0,0)\;=\;\mathcal C(\theta)\,W_{\Lambda_1}(0)\,W_{\Lambda_2}(0)\,
e^{-\,m_\bullet d}\;=\;\varepsilon_0\,e^{-\,m_\bullet d}\,.
\end{equation}
For finite $\ell$, $\varepsilon$ decays with $k_\ell$ both by the Yukawa factor
$e^{-\,d\,\sqrt{k_\ell^2+m_\bullet^2}}$ and by the induced–gravity screen
$q/(q+r k_\ell^2)$.

\paragraph{Real–space representation and core scale.}\label{par:uom-main-067}
In real space the coupling reads a rotationally–invariant convolution,
\begin{equation}
\label{eq:mirage-convolution}
\rho^{\rm eff}_1(t,\mathbf x)\;=\;\big(K_d * \rho_2\big)(t,\mathbf x),
\qquad
K_d\ \Longleftrightarrow\ \varepsilon(\omega,\ell;\cdots).
\end{equation}
At DC ($\omega=0$), the \emph{real–space} Green’s function is the \emph{convolution}
of two positive, radial, decreasing kernels:
\begin{equation}
\label{eq:conv-kernels}
\Phi_{\rm sh}(r)\;\propto\;\underbrace{\mathcal F^{-1}\!\Big[\frac{1}{k^2}\,\frac{q(0,k)}{q(0,k)+r(\Lambda)\,k^2}\Big]}_{\text{screened Newton}}\;\ast\;
\underbrace{\mathcal F^{-1}\!\Big[e^{-\,d\,q(0,k)}\Big]}_{\text{bulk Yukawa (radial)}}\!,
\qquad q(0,k)=\sqrt{k^2+m_\bullet^2}.
\end{equation}
The inner behavior is controlled by the \emph{shortest} infrared length. In
particular, the (Green’s–function) core radius $r_{\rm core}$, defined by
$\Phi_{\rm sh}(r_{\rm core})/\Phi_{\rm N}(r_{\rm core})=\tfrac12$, satisfies the
\emph{parametric} estimate
\begin{equation}
\label{eq:rcore-param}
\boxed{\quad
r_{\rm core}\ \simeq\ C\;\max\!\{\,r_{\rm scr},\ m_\bullet^{-1},\ \Lambda^{-1}\,\}\,,\qquad
r_{\rm scr}:=\sqrt{\frac{r(\Lambda)}{q_0}}\equiv \sqrt{\frac{r(\Lambda)}{m_\bullet}}\,,\quad
C=\mathcal O(1),
\quad}
\end{equation}
with $q_0:=q(0,0)=m_\bullet$ and $\Phi_{\rm N}(r)\propto 1/r$ the Newton kernel.  %%%

\paragraph{Mirage dark matter on our brane.}\label{par:uom-main-068}
The uniform DC piece governs the ``mirage'' matter fraction. In terms of the critical
density on our brane, $\rho_c(a)=3H_0^2 E(a)^2/(8\pi G)$, one has
\begin{equation}
\label{eq:OmegaX}
\Omega_X(a)\;=\;\frac{\rho^{\rm eff}_1(a)}{\rho_c(a)}
\;=\;\frac{\varepsilon_{\rm DC}(a)\,\rho_2(a)}{\rho_c(a)}
\;=\;\Omega_{X0}\,\frac{\varepsilon_{\rm DC}(a)}{\varepsilon_{\rm DC}(1)}\,
\frac{a^{-3}}{E(a)^2},\qquad E(a):=H(a)/H_0,
\end{equation}
where $\Omega_{X0}=\varepsilon_0\,\Omega_{2,0}$ is today’s uniform fraction, and
$\rho_2\propto a^{-3}$ on the shadow brane. As our brane accelerates, $d(a)$ grows and
$\varepsilon_{\rm DC}(a)=\varepsilon_0\exp[-m_\bullet\,d(a)]$ decays, so $\Omega_X$
is higher at recombination and relaxes to a percent level today. The clumpy part
transfers through the filtered kernel $K_d$, producing cored halos and a scale–dependent
shadow contribution to LSS that weakens as $d(a)$ increases.

\paragraph{Observables and inversion.}\label{par:uom-main-069}
With fixed $(\Lambda_1,\Lambda_2)$ (set at flips) and unknown $(m_\bullet,d_0,\theta)$,
one can fit:
\begin{itemize}[leftmargin=1.2em]
\item the uniform history $\Omega_X(z)$ from CMB background (via $r_s$), early–ISW, and
late $H(z)$, using \eqref{eq:OmegaX} with $d(z)$ parameterized (e.g.\ spline or
$w$CDM–anchored growth);
\item the scale response $b(k)\simeq \varepsilon(0,\ell\!\leftrightarrow\!k)$ from galaxy
bias and weak–lensing amplitude (fixes $r_{\rm core}$ and $r(\Lambda)$);
\item small–scale cores in rotation curves / strong lensing (tests $K_d(r)$).
\end{itemize}
Low–$\ell$ CMB temperature anomalies \emph{without} an $E$–mode partner and a mild
dipolar modulation arise for small tilts ($\mathcal C(\theta)\neq{\rm const}$) with
$\varepsilon(0,\ell)$ dominated by $\ell\lesssim5$; stringent quadrupole limits then
bound $\theta$.

\paragraph{Conservation and quasi–static use.}\label{par:uom-main-070}
When using $\delta\rho^{\rm eff}_1=\varepsilon(\omega,k)\,\delta\rho_2$ as a Poisson
source, the effective pressure/velocity potentials consistent with the 4D Bianchi
identity are implicitly fixed by the quasi–static/DC kernel; our CMB/LSS constraints are
used precisely in this regime.

\medskip
\noindent\textbf{Summary.}
A single, factorized kernel $\varepsilon(\omega,\ell;\Lambda_1,\Lambda_2,d,\theta)$
encodes bulk mediation between fixed–cutoff dS branes. Its DC piece governs an evolving,
uniform ``mirage'' matter fraction $\Omega_X$ (larger at recombination, percent–level
today); its $k$–dependence predicts cored halos and a scale–dependent bias; and radion
(or spin–2) stabilization plus induced–gravity screening keeps the construction compatible
with precision tests.


\subsection{Five–dimensional influence action: bulk\,+\,branes and reduction to 4D P3}
\label{subsec:5D-influence}

\paragraph{Scope.} \label{par:uom-main-071}
We formulate the Schwinger–Keldysh (SK) on–shell influence action for the AdS$_5$ bulk with a stack of codimension–1 branes:
(i) the \emph{receiver} AdS$_4$ brane (code space), 
(ii) the \emph{source} dS$_4$ brane (radiative sector),
and (iii) an arbitrary collection of older–aeon branes (``shadow'' sector), all embedded in the same AdS$_5$ geometry.
The welded, tick–localized LoG source lives on the receiver brane. 
Integrating out the bulk and brane auxiliaries yields a single convex boundary functional whose real part reproduces the modular/Fisher quadratics of 4D P3 and whose imaginary part gives the BY/anomaly variance (the rectifier). In allowing the receiver cutoff to vary tick-by-tick we will sometimes work in a “moving-wall gauge”, where the wall position $r(t)=\Lambda(t)^{-1}$ is dynamical.
All coefficients are fixed by geometry and stabilization; no phenomenological ``knobs'' are introduced.

\subsubsection{Geometry, hypersurfaces, and SK doubling}
\label{subsubsec:5D-geom-SK}

\paragraph{Bulk.}\label{par:uom-main-072}
Let $(\mathcal M_5,g_{AB})$ be asymptotically AdS$_5$ \cite{Maldacena1998,Witten1998AdSCFT,GubserKlebanovPolyakov1998} with radius $\ell$,
\(
R_{AB}-\tfrac12 R g_{AB}-\tfrac{6}{\ell^2} g_{AB}=0
\)
at the background. Capital Latin indices run over $0,\dots,4$.
We include the Gibbons–Hawking–York term on every non-null boundary and the standard holographic counterterms.

\paragraph{Branes.}\label{par:uom-main-073}
Let there be $N_{\rm br}$ timelike 4D hypersurfaces $\{\Sigma_i\}$ with induced metrics $\gamma^{(i)}_{\mu\nu}$ and outward normals $n^A_{(i)}$:
\[
\Sigma_{\rm rec}\ (\text{AdS}_4\ \text{receiver}),\quad
\Sigma_{\rm src}\ (\text{dS}_4\ \text{source}),\quad
\Sigma_{\rm old}^{(a)},\ a=1,\dots,N_{\rm old}.
\]
Each carries a localized action (tension, induced terms, matter/auxiliary sectors) spelled out below. Greek indices $\mu,\nu$ are tangential to branes.

\paragraph{SK contour and doubling.}\label{par:uom-main-074}
We use the closed time path $\mathcal C$; every field $X$ is doubled $(X_+,X_-)$.
The SK action is
\(
S^{\rm SK}=S[g_+,\dots]-S[g_-,\dots],
\)
and the influence functional for a brane source $J$ is
\begin{equation}
W[J]\ :=\ S^{\rm SK}_{\rm bulk+branes}\Big|_{\rm on\ shell\ with\ welded\ source\ J}.
\label{eq:W-def}
\end{equation}
Retardedness and KMS of the source bath ensure $\Im W\ge 0$.

\subsubsection*{Microscopic 5D action}
\label{subsubsec:5D-action-micro}

\paragraph{Gravitational sector (bulk).}\label{par:uom-main-075}
\begin{align}
S_{\rm grav}
&=\frac{M_5^3}{2}\int_{\mathcal M_5}\! d^5x\,\sqrt{-g}\,\Big(R+\frac{12}{\ell^2}\Big)
+M_5^3\sum_i\int_{\Sigma_i}\! d^4x\,\sqrt{-\gamma^{(i)}}\,K^{(i)}
\nonumber\\
&\quad
-\sum_i\int_{\Sigma_i}\! d^4x\,\sqrt{-\gamma^{(i)}}\,\Big(
\frac{3M_5^3}{\ell} - \frac{\ell M_5^3}{4}R[\gamma^{(i)}]+\cdots\Big),
\label{eq:Sgrav}
\end{align}
where $K^{(i)}$ is the trace of the extrinsic curvature and the last line are holographic counterterms through $d{=}4$ (higher–derivative local terms are implicit).

\paragraph{Goldberger–Wise (GW) stabilization.} \cite{GoldbergerWise1999}\label{par:uom-main-076}
A bulk scalar $\Phi$ (mass $m_\Phi$) with boundary potentials $\lambda^{(i)}(\Phi)$:
\begin{equation}
S_{\rm GW}
=\int_{\mathcal M_5}\! d^5x\,\sqrt{-g}\,\Big(\tfrac12(\nabla\Phi)^2-\tfrac12 m_\Phi^2\Phi^2\Big)
-\sum_i\int_{\Sigma_i}\! d^4x\,\sqrt{-\gamma^{(i)}}\,\lambda^{(i)}(\Phi).
\label{eq:SGW}
\end{equation}
The classical $\Phi$ profile induces an effective radion potential for each independent inter--brane separation.
We take all \emph{interior/shadow} RS1 radions heavy (GW-stabilized). The active receiver is a KR head;
its persistent light scalar is the boundary-distance mode (physical kernel) with Robin-suppressed visible overlap.
A brane-localized \emph{GW boundary fluctuation} field $\mathcal S$ (on $\Sigma_{\rm rec}$) with mass
$m_{\mathcal S}\sim 10^{4\text{--}5}\,$GeV mediates the portal to the Higgs without reintroducing hierarchy issues.

\paragraph{Receiver brane: compensator \& welded source.}\label{par:uom-main-077}
On $\Sigma_{\rm rec}$ we include:
(i) the Wess–Zumino (Riegert–Fradkin–Tseytlin) compensator $\sigma$ producing the 4D trace anomaly,
(ii) the boundary fluctuation $\mathcal S$, and
(iii) the welded LoG source $J$:
\begin{align}
S_{\rm comp}[\sigma;\gamma]
&=\frac{a}{24\pi}\int_{\Sigma_{\rm rec}}\!\! d^4x\,\sqrt{-\gamma}\,
\Big(\tfrac12\,\sigma\,\Delta_4\,\sigma+\sigma\,\mathcal A[\gamma]\Big)
+\frac{c}{16\pi^2}\int \sqrt{-\gamma}\,\sigma\,W^2[\gamma],
\label{eq:Scomp}\\
S_{\rm S,H}
&=\int_{\Sigma_{\rm rec}}\!\! d^4x\,\sqrt{-\gamma}\,
\Big[\tfrac12(\partial\mathcal S)^2-\tfrac12 m_{\mathcal S}^2\mathcal S^2-U(\mathcal S)
-|D H|^2-V_H(H)-\lambda_p\,\mathcal S\,H^\dagger H\Big],
\label{eq:SSH}\\
S_{\rm weld}[J]
&=\int_{\Sigma_{\rm rec}}\!\! d^4x\,\sqrt{-\gamma}\,
\big(J\cdot \mathcal O_{\rm LoG}\big),
\qquad
J\equiv J_k\ \text{(tick–localized LoG envelope).}
\label{eq:Sweld}
\end{align}
Here $a=c=\pi\ell^3/(8G_5)$, $\Delta_4$ is the Paneitz operator, $\mathcal A=E_4-\tfrac23\Box R$,
and $\mathcal O_{\rm LoG}$ is the compiled trace–sector operator sourced by the welded LoG pulse (\S\ref{subsec:receiver-anomaly}).

\paragraph{Wall motion and RG work (moving–wall gauge).}\label{par:uom-main-078}
When we allow the receiver cutoff to vary within a tick, we parametrize the wall by its radial position
$r(t)=\Lambda(t)^{-1}$ and work in the ``moving–wall’’ gauge. In this gauge the on–shell SK action acquires
a local, quadratic ``RG kinetic’’ term,
\begin{equation}
W_{\rm RG}[r]
\;=\;
\frac12\int_{I_k}\! dt\;\mu_{\rm RG}(\Lambda(t))\,\dot\Lambda(t)^{2},
\qquad
\mu_{\rm RG}(\Lambda)
\;=\;\int_{\Sigma_{\Lambda}}\!\sqrt{-\gamma}\;\mathcal H_{\rm BY}(\Lambda)
\;\propto\; M_4^2(\Lambda),
\label{eq:RG-kinetic}
\end{equation}
where $\mathcal H_{\rm BY}$ is the Brown--York wall Hamiltonian density and $M_4^2(\Lambda)=M_5^3\ell+\kappa\,a\,\Lambda^2$
is the induced gravity coefficient (fixed by holographic counterterms).
The corresponding wall power balance (Hamilton--Jacobi form) is
\begin{equation}
\mu_{\rm RG}(\Lambda)\,\dot\Lambda(t)
\;=\;
\underbrace{\big(\mathcal P_{\rm in}-\mathcal P_{\rm out}\big)(t)}_{\text{all inflows at the wall minus losses}}
\;-\;
\underbrace{\dot{\mathcal F}_{\rm rec}(t)}_{\text{reversible receiver drop}}\;,
\label{eq:wall-ode}
\end{equation}
so that integrating over a tick $I_k$ gives the wall--move inequality
\begin{equation}
\lambda\!\big(\Lambda_{k+1}\big)-\lambda\!\big(\Lambda_{k}\big)
\;\le\;
W^{\rm avail}_k,
\qquad
\lambda'(\Lambda)=\mu_{\rm RG}(\Lambda),\quad
W^{\rm avail}_k:=\int_{I_k}\!(\mathcal P_{\rm in}-\mathcal P_{\rm out})\,dt-\int_{I_k}\!\dot{\mathcal F}_{\rm rec}\,dt.
\label{eq:wall-move-ineq}
\end{equation}
Equality holds only when the full available budget is spent on wall motion (continuous control, no slack);
otherwise the inequality is strict and the bound is not saturated.
We will use \eqref{eq:wall-ode}--\eqref{eq:wall-move-ineq} only as a dynamical vehicle to derive the per–tick/per–aeon
cutoff conditions; the post--exhaustion boundary analysis is formulated in the \emph{static--wall} gauge $r=r_\ast$,
where $\dot\Lambda\equiv 0$ and radial reparametrizations are encoded by the boundary CPTP RG flow
$\partial_u\rho_u=-i[H_u,\rho_u]+\mathcal D_u(\rho_u)$ with $u=\ln(\Lambda_{\rm UV}/\Lambda)$.


\paragraph{Source and older branes.}\label{par:uom-main-079}
On $\Sigma_{\rm src}$ we put the QND absorber (GKLS dilation) coupled to the commutant KMS bath; on $\Sigma_{\rm old}^{(a)}$ we allow localized stress $T^{(a)}_{\mu\nu}$ and possible heavy stabilized radions. Their role in $W[J]$ is purely via retarded bulk propagation (no explicit new parameters are introduced here).

\subsubsection*{Linear response and the influence functional}
\label{subsubsec:linresp-W}

\paragraph{Gaussian sector and integrated auxiliaries.}\label{par:uom-main-080}
We SK–double \eqref{eq:Sgrav}–\eqref{eq:Sweld}, integrate out $(g_{\pm},\Phi_{\pm},\sigma_{\pm},\mathcal S_{\pm},\ldots)$ at quadratic order in the welded source $J$, and evaluate on the classical solution with retarded boundary conditions. In Fourier space on $\Sigma_{\rm rec}$:
\begin{equation}
W[J]
=\frac12\!\int\!\frac{d^4k}{(2\pi)^4}\,
J(-k)\,\big(\,\underbrace{G_R(k)}_{\text{real, retarded}} + i\,\underbrace{\mathcal N(k)}_{\text{noise/Hadamard}}\,\big)\,J(k)
+O(J^3).
\label{eq:W-quadratic}
\end{equation}
$G_R$ and $\mathcal N$ are positive kernels fixed by $(\ell,G_5)$, the induced gravity scale $G_4=G_5/\ell$, and the stabilization data. 
Causality ($G_R$ retarded) and the KMS condition on the source bath imply the SK fluctuation–dissipation inequality $\mathcal N(k)\ge 0$ and hence $\Im W\ge 0$.


\subsubsection*{Holographic reduction to boundary modular functionals}\label{subsubsec:uom-main-011}

\paragraph{Linearized Gaussian sector and covariance form.}\label{par:uom-main-081}
In the radiative sector (Gaussian, zero mean), the receiver state is determined by a covariance $C$ around a fixed AdS$_4$ background $C_\ast$. The welded linear response maps a (trace–channel) source $\tau$ into $\delta C=\mathsf T_\Lambda \tau$ (the cutoff $\Lambda$ enters via the induced gravity term and regulator).

\begin{lemma}[Real part as Fisher quadratic]
\label{lem:ReW-Fisher}
In the Gaussian sector and to quadratic order in $\delta C$, the real part of the SK influence functional equals the Fisher norm of the covariance tangent:
\[
\Re W\;=\;\frac12\,\big\langle\!\big\langle \delta C,\delta C\big\rangle\!\big\rangle_{C_\ast}
\,+\,O(\|\delta C\|^3),
\qquad
\big\langle\!\big\langle \delta C,\delta C'\big\rangle\!\big\rangle_{C_\ast}
:=\frac12\,\mathrm{tr}\!\big(C_\ast^{-1}\delta C\,C_\ast^{-1}\delta C'\big).
\]
\end{lemma}
\begin{proof}[Proof of Lemma~\ref{lem:ReW-Fisher} (Real part as Fisher quadratic)]
\emph{Assumptions and setting.}
Work in the radiative (Gaussian, zero–mean) sector of the receiver CFT on a maximally symmetric AdS$_4$ background. 
Let $\rho_\ast$ be the fixed point (reference) Gaussian state with covariance $C_\ast$ (a strictly positive, trace–class operator on the one–particle space), and let $\rho$ be a nearby Gaussian state with covariance $C=C_\ast+\delta C$ produced by a small welded trace–channel source $\tau$ via the linear response map $\delta C=\mathsf T_\Lambda\,\tau$. 
We consider the SK influence functional $W[\tau]$ evaluated on the holographic saddle (bulk Einstein + matter), and expand it to quadratic order in $\delta C$ (equivalently, in $\tau$).

\smallskip
\noindent
\emph{Step 1: Real part of on–shell SK action equals boundary relative entropy at quadratic order.}
For perturbations around a maximally symmetric saddle and sources that create normalizable linearized bulk fields, the real part of the on–shell action computes the canonical quadratic form (``canonical energy'') of the linearized bulk solution. 
By the holographic equality between boundary relative entropy and bulk canonical energy (for small perturbations; the argument uses the AdS/CFT dictionary and the first law of entanglement/relative entropy), one has
\begin{equation}
\Re W[\tau]\;=\;S(\rho\,\|\,\rho_\ast)\;+\;O(\|\delta C\|^3)\,,
\label{eq:ReW=RelEnt}
\end{equation}
i.e.\ the two functionals agree up to cubic order in the perturbation.\footnote{The equality can be established either (i) by evaluating the quadratic on–shell action with SK boundary conditions and matching it to the canonical energy, or (ii) by using that the second variation of the boundary relative entropy equals the quadratic canonical energy (hence the quadratic real part of the on–shell action), and both vanish to first order at the reference state.}

\smallskip
\noindent
\emph{Step 2: Quadratic expansion of Gaussian relative entropy gives the Fisher metric.}
Let $\rho(C)$ denote a centered Gaussian state of a bosonic CCR system with covariance $C>0$ (no means). The Umegaki relative entropy $S(\rho(C)\,\|\,\rho(C_\ast))$ admits a second–order Taylor expansion at $C_\ast$ whose quadratic form is the BKM/Fisher information metric on the covariance tangent:
\begin{equation}
S(\rho(C_\ast+\delta C)\,\|\,\rho(C_\ast))
\;=\;\frac{1}{2}\,\big\langle\!\big\langle \delta C,\delta C\big\rangle\!\big\rangle_{C_\ast}
\;+\;O(\|\delta C\|^3),
\qquad
\big\langle\!\big\langle X,Y\big\rangle\!\big\rangle_{C_\ast}
:=\frac12\,\mathrm{tr}\!\big(C_\ast^{-1}X\,C_\ast^{-1}Y\big).
\label{eq:GaussianRelEnt}
\end{equation}
\emph{Proof of \eqref{eq:GaussianRelEnt}:}
Bring $C_\ast$ to Williamson form by a symplectic $S$ so that $S C_\ast S^{\sf T}=\mathrm{diag}(\nu_1,\nu_1,\dots,\nu_n,\nu_n)$ in $n$ modes; the CCR/quadratic Hamiltonian diagonalize into independent thermal oscillators. 
For a single mode with covariance $\nu>0$ the relative entropy between centered Gaussians is an explicit smooth function of $\nu$ whose Taylor expansion at $\nu_\ast$ gives
\(
S(\nu_\ast+\delta\nu\|\nu_\ast)=\tfrac{1}{4}\,(\delta\nu)^2/\nu_\ast^2+O((\delta\nu)^3).
\)
Summing over modes and undoing the symplectic transformation yields \eqref{eq:GaussianRelEnt} in operator form. (Equivalently, one may use the general identity that the Hessian of $S(\rho_\theta\|\rho_{\theta_0})$ at $\theta_0$ equals the quantum Fisher/BKM metric and compute it for Gaussian families in covariance coordinates.)

\smallskip
\noindent
\emph{Step 3: Identification with $\Re W$.}
Combining \eqref{eq:ReW=RelEnt} and \eqref{eq:GaussianRelEnt} we obtain
\[
\Re W[\tau]\;=\;\frac12\,\big\langle\!\big\langle \delta C,\delta C\big\rangle\!\big\rangle_{C_\ast}
\;+\;O(\|\delta C\|^3),
\qquad \delta C=\mathsf T_\Lambda\,\tau.
\]
Since both $\Re W$ and $S(\rho\|\rho_\ast)$ vanish to first order at the reference state and have the same positive quadratic form on the covariance tangent, their difference is cubic and higher, as claimed.

\smallskip
\noindent
\emph{Conclusion.}
To quadratic order in the welded perturbation, the real part of the SK influence functional equals the Fisher (BKM) quadratic norm of the covariance tangent at $C_\ast$, completing the proof.
\end{proof}


\begin{lemma}[Imaginary part as positive noise kernel]
\label{lem:ImW-noise}
To the same order, the imaginary part is a positive quadratic functional of the welded source,
\[
\Im W\;=\;\frac12\int\!\frac{d\omega\,d^3k}{(2\pi)^4}\ \widehat{\tau}^\dagger(\omega,k)\,
\mathcal N_\Lambda(\omega,k)\,\widehat{\tau}(\omega,k)\ \ge\ 0,
\]
where $\mathcal N_\Lambda$ is the Keldysh noise kernel fixed by the KMS bath and the retarded AdS$_5$/brane Green functions (cutoff–dependent). In the (nearly) maximally symmetric trace sector this reduces to the anomaly/BY variance injection $\propto \Delta\!\langle R^2\rangle$ (below).
\qed
\end{lemma}
\begin{proof}[Proof of Lemma~\ref{lem:ImW-noise} (Imaginary part as positive noise kernel)]
\emph{Assumptions and setting.}
Let $\tau=(\tau_1,\tau_2)$ be the SK sources on the forward/backward branches coupled linearly to a Hermitian boundary operator $\mathcal O$ in the trace channel (the welded image of the bulk trace). The SK influence functional is quadratic to leading order:
\[
W[\tau]\;=\;-\frac{1}{2}\int\! d^4x\,d^4y\ \tau_a(x)\,G^{ab}(x-y)\,\tau_b(y)\ +\ O(\tau^3),
\]
with the SK kernel $G^{ab}$ built from the connected two–point functions of $\mathcal O$:
\(
G^{11}=G^{\rm T},\ G^{22}=G^{\tilde{\rm T}},\ G^{12}=G^{<},\ G^{21}=G^{>}.
\)
Pass to the Keldysh ``classical/quantum'' basis $\tau_{\rm cl}=(\tau_1+\tau_2)/2$, $\tau_{\rm q}=\tau_1-\tau_2$. In this basis
\[
W[\tau]\;=\;-\int\! d^4x\,d^4y\ \Big\{\tau_{\rm q}(x)\,G^{\rm R}(x-y)\,\tau_{\rm cl}(y)\;-\;\frac{i}{2}\,\tau_{\rm q}(x)\,N(x-y)\,\tau_{\rm q}(y)\Big\}\;+\;O(\tau^3),
\]
where $G^{\rm R}=G^{>}-G^{\rm T}$ is the retarded kernel and $N=\tfrac12(G^{>}+G^{<})$ is the Keldysh (noise) kernel. Hence
\[
\Im W[\tau]\;=\;\frac{1}{2}\int\! d^4x\,d^4y\ \tau_{\rm q}(x)\,N(x-y)\,\tau_{\rm q}(y)\;+\;O(\tau^3)\,.
\]
It suffices to prove that $N$ is a positive semidefinite (PSD) distribution on test functions (after welding/regulation).

\smallskip
\noindent
\emph{Step 1: Spectral representation and KMS.}
For the stationary AdS$_4$ head in thermal/KMS equilibrium (or vacuum in the gapped sector), the Wightman kernels admit a spectral density $\rho(\omega,\mathbf k)$:
\[
\widehat{G}^{>}(\omega,\mathbf k)\;=\;2\pi\,\rho(\omega,\mathbf k)\,\big[1+n_B(\omega)\big],\qquad
\widehat{G}^{<}(\omega,\mathbf k)\;=\;2\pi\,\rho(\omega,\mathbf k)\,n_B(\omega),
\]
with Bose factor $n_B(\omega)=(e^{\beta\omega}-1)^{-1}$ and \(\rho(\omega,\mathbf k)\ge 0\) for $\omega>0$ (positivity of the spectral measure for Hermitian $\mathcal O$). The retarded kernel obeys the fluctuation–dissipation relation
\[
\widehat{G}^{\rm R}(\omega,\mathbf k)\;=\;\int\!\frac{d\omega'}{\pi}\,\frac{\rho(\omega',\mathbf k)}{\omega-\omega'+i0}\,,
\qquad
\Im \widehat{G}^{\rm R}(\omega,\mathbf k)\;=\;\pi\,\rho(\omega,\mathbf k)\,\mathrm{sgn}(\omega).
\]
Therefore the Keldysh kernel in Fourier space is
\begin{equation}
\widehat{N}(\omega,\mathbf k)\;=\;\tfrac12\big(\widehat{G}^{>}+\widehat{G}^{<}\big)
\;=\;\pi\,\rho(\omega,\mathbf k)\,\coth\!\Big(\frac{\beta\omega}{2}\Big)\,,
\label{eq:N-FDT}
\end{equation}
which is pointwise nonnegative since $\rho\ge 0$ and $\coth(\beta\omega/2)$ has the sign of $\omega$ while $\rho$ is odd in $\omega$ in the standard convention (equivalently $\pi\rho\,\coth(\beta\omega/2)=|\pi\rho|\,\big(1+2n_B(|\omega|)\big)\ge 0$).

\smallskip
\noindent
\emph{Step 2: Positivity as a quadratic form.}
Let $f$ be a smooth compactly supported (test) function, and set $\widehat{f}(\omega,\mathbf k)$ its Fourier transform. Then
\[
\int\! d^4x\,d^4y\ f(x)\,N(x-y)\,f(y)
=\int\!\frac{d\omega\,d^3k}{(2\pi)^4}\ \widehat{f}^\ast(\omega,\mathbf k)\,\widehat{N}(\omega,\mathbf k)\,\widehat{f}(\omega,\mathbf k)\ \ge\ 0,
\]
by \eqref{eq:N-FDT}. Hence $N$ is PSD in the sense of distributions. Consequently,
\[
\Im W[\tau]\;=\;\tfrac12\,\langle \tau_{\rm q}, N\,\tau_{\rm q}\rangle\ \ge\ 0.
\]

\smallskip
\noindent
\emph{Step 3: Effect of welding/regulation (cutoff dependence).}
The welded source is $\tau=W_\Lambda\!\star\!\tau_{\rm bare}$ with a real, even regulator/window $W_\Lambda$ (cutoff $\Lambda$ and induced–gravity screen). In Fourier space this multiplies the kernels by $|\widehat{W}_\Lambda|^2\ge 0$:
\(
\widehat{N}_\Lambda=|\widehat{W}_\Lambda|^2\,\widehat{N}.
\)
Therefore PSD is preserved:
\[
\int\!\frac{d\omega\,d^3k}{(2\pi)^4}\ \widehat{f}^\ast\,\widehat{N}_\Lambda\,\widehat{f}
=\int\!\frac{d\omega\,d^3k}{(2\pi)^4}\ \big|\widehat{W}_\Lambda\,\widehat{f}\big|^2\,\widehat{N}\ \ge\ 0.
\]
This covers both the UV mollifier in the bridge and the Yukawa/screening factor from the induced 4D Einstein term; both enter as nonnegative multipliers at fixed $(\omega,\mathbf k)$.

\smallskip
\noindent
\emph{Conclusion.}
To quadratic order in the source, the imaginary part of the SK influence functional is
\[
\Im W\;=\;\frac12\int\!\frac{d\omega\,d^3k}{(2\pi)^4}\ \widehat{\tau}_{\rm q}^\dagger(\omega,\mathbf k)\,
\widehat{N}_\Lambda(\omega,\mathbf k)\,\widehat{\tau}_{\rm q}(\omega,\mathbf k)\ \ge\ 0,
\]
with $\widehat{N}_\Lambda$ PSD by the spectral/KMS (fluctuation–dissipation) structure and preserved by welding/regulation. This proves the claim.
\end{proof}


\paragraph{Projection to trace/RP/conformal sectors.}\label{par:uom-main-082}
Let $\delta C$ be the Gaussian covariance tangent on $\Sigma_{\rm rec}$ induced by $J$. 
As in Lemma~\ref{lem:tan-orth} (proved earlier), define three orthogonal Fisher–projectors to unphysical subspaces
$\Pi_{\rm RP}$ (OS–odd), $\Pi_{\rm Weyl}$ (pure trace/Weyl), and $\Pi_{\rm conf}$ (conformal mismatch), which commute at the maximally symmetric receiver fixed point. Then
\begin{align}
\Re W[J]
&=\frac12\sum_{i\in\{\mathrm{RP,\,Weyl,\,conf}\}}
\big\langle\!\big\langle \Pi_i\delta C,\ \Pi_i\delta C\big\rangle\!\big\rangle_{C_\ast}
\ =:\ \frac12\|\delta C\|_{F,\rm unphys}^2,\label{eq:ReW}\\
\Im W[J]
&=\frac12\!\int\!\frac{d^4k}{(2\pi)^4}\,J(-k)\,\mathcal N(k)\,J(k)
\ \equiv\ \gamma\,\Delta\!\langle R^2\rangle[J]\ \ge 0.
\label{eq:ImW}
\end{align}
Equation \eqref{eq:ReW} reproduces the 4D Fisher–quadratic costs in the three unphysical blocks; \eqref{eq:ImW} reduces, on the LoG family and in the long–wavelength regime, to the anomaly variance with
\(
\gamma=\frac{aG_4}{12\pi}>0
\)
and the variance scaling
\(
\Delta\!\langle R^2\rangle \sim (8\pi G_4)^2 \langle |\mathcal K|^2\rangle\, C_{\rm LoG}\,A^2/(\sigma_t\,\sigma_x^3)
\)
(cf.\ \eqref{eq:variance-scaling}).

\paragraph{Receiver modular functional and reduction to 4D P3.}\label{par:uom-main-083}
At quadratic order,
\begin{equation}
\boxed{\
\Phi_{\rm rec}[J]\ :=\ \Re W[J]\ +\ \Im W[J]
\ =\ \frac12\sum_{i=\rm RP,Weyl,conf}\!\|\Pi_i\delta C\|_{F}^2\ +\ \gamma\,\Delta\!\langle R^2\rangle[J]\ .
}
\label{eq:Phi-rec-5D}
\end{equation}
When $J$ is the welded LoG envelope of a tick with widths $(\Sigma,\sigma_t,\sigma_x)$ and effective amplitude $A$, inserting \eqref{eq:payload-equality} eliminates amplitudes and yields the 4D receiver P3 integrand 
$\Phi_{\rm BY}+\Phi_{\rm RP}+\Phi_{\rm conf}$ of Eq.\ \eqref{eq:Jrec}.

\subsubsection*{P2 from the 5D saddle and monotonicity under refinement}
\label{subsubsec:P2-from-5D}

\paragraph{Tickwise constancy at fixed cutoff.}\label{par:uom-main-084}
On the SK saddle with (i) QND, pointer–aligned source and (ii) admissible welded bridge at fixed $\Lambda$
(intra--tick only; see the receiver dynamics axiom), the net boundary modular balance across a tick satisfies
\(
\int_{I_k}\! (\dot F_{\rm src}+\dot C)\,dt = 0,
\)
i.e.\ the source dwell equals the classical record (\ref{post:P2}).
In 5D this is simply stationarity of $\Re W$ under admissible $J$ at fixed regulators, with $\Im W$ minimized to its LoG value ($\ge 0$).

\paragraph{Monotonicity across ticks/aeons.}\label{par:uom-main-085}
Refining the receiver cutoff $\Lambda\mapsto\Lambda'$ enlarges the admissible Hilbert space on $\Sigma_{\rm rec}$ and increases the Fisher–quadratic capacity of physical directions; consequently
\(
\Re W[J]\ \text{(at the minimizer)}\ \uparrow
\)
in a way that equals the code capacity increment $\Delta\mathcal C_{\rm code}(\Lambda)$. 
The rectified DC piece from $\Im W$ (anomaly floor) projects onto the constant mode and persists post–flip, seeding the next aeon. Thus the total recoverable information is non–decreasing, as asserted by P2.

\subsubsection{Orthogonality of radion to the DC trace and role of \texorpdfstring{$\mathcal S$}{S}}
\label{subsubsec:radion-orth}

\paragraph{RS1$\leftrightarrow$KR scalar continuity.}\label{par:uom-main-086}
The active receiver is a KR head throughout the split-ordered evolution; RS1 segments, when present, are interior/shadow only.
When an interior RS1 segment abuts a KR end--of--the--world cap, the scalar sector evolves \emph{smoothly}
with no mode creation, no tachyons, and no ghosts. In an RS1 cell there is a physical radion
$\varphi_r\propto\xi_2-\xi_1$ (brane–bending difference) stabilized by the Goldberger–Wise (GW) profile; its mass is positive,
$m_r^2\sim \nu^2 k^2\,\epsilon\,e^{-2k d_0}$, and it lies in the \emph{physical} covariance tangent (no conformality penalty).

At a KR cap the second boundary is absent, but the scalar mode does \emph{not} disappear: the RS1 radion flows to the KR
\emph{boundary–distance} scalar (the $E_0=4$ AdS$_4$ scalar) with $m_r^2=4k^2+\delta m_{\rm GW}^2>0$. It still lives in the
\emph{physical} tangent. The visible coupling of this mode is suppressed by the stiff visible Robin boundary condition,
$u_r(z_{\rm vis})/u_r(z_{\rm sh})=\varepsilon_{\rm BC}\simeq 1/(1+s_{\rm vis}\,kL)\ll 1$ with $s_{\rm vis}\gg1$, so fifth–force
constraints remain satisfied. Its projection onto the conformality–mismatch block vanishes at leading order, and the residual
leakage induced by Weyl–radion mixing at welded band edge obeys
$\theta(\Lambda)=O(\varepsilon_{\rm BC}\Lambda^{-2})$, hence the corresponding penalty scales as
$\Phi_{\rm conf}\sim O(\varepsilon_{\rm BC}^{2}\Lambda^{-4})$ (tiny and \emph{decreasing} with $\Lambda$).

The brane–localized GW \emph{boundary fluctuation} $\mathcal S$ on $\Sigma_{\rm rec}$ remains heavy ($m_{\mathcal S}\gg v$) across
the transition and carries the \emph{portal/DC} response: to leading order,
$\langle \mathcal S\rangle\simeq J/m_{\mathcal S}^2$ (Robin BC), and the Higgs mass shift
$\Delta\mu_H^2=-\lambda_p J/m_{\mathcal S}^2$ is \emph{continuous} across the head change. The curvature–induced shifts
($\propto L_4^{-2}$) are positive and small; no sign flip occurs. Consequently, the scalar two–point data and the portal normalization
are continuous across an interior RS1--to--KR adjacency: on the interior segment the stabilized RS1 radion sits in the physical kernel;
on the KR head the boundary–distance scalar persists there with BC–suppressed visible coupling, while the heavy $\mathcal S$ continues
to set the leading DC imprint. The BY/anomaly (Weyl) block accounts for the DC trace bias.



\paragraph{Maximal symmetry and Fisher orthogonality.}\label{par:uom-main-087}
At the flip the receiver geometry is (locally) maximally symmetric AdS$_4$; the tangent space of Gaussian covariances at the fixed point $C_\ast$ decomposes orthogonally (Fisher metric) into the RP–odd, Weyl (pure trace), conformal–mismatch, and physical kernels (Lemma~\ref{lem:tan-orth}).
The boundary–distance scalar on the KR head (and the stabilized radion on interior RS1 segments) belongs to the \emph{physical} kernel
(it is a genuine modulus fluctuation when GW is included), while the DC anomaly trace lives in the Weyl block; hence
\(
\langle\!\langle \delta C_{\rm rad},\delta C_{\rm Weyl}\rangle\!\rangle_{C_\ast}=0
\)
to quadratic order. This realizes the requested \emph{orthogonality} without ad hoc renormalization.

\paragraph{Boundary GW fluctuation \texorpdfstring{$\mathcal S$}{S} and the Higgs portal.}\label{par:uom-main-088}
The brane–localized GW fluctuation $\mathcal S$ couples as $-\lambda_p\,\mathcal S\,H^\dagger H$ and is sourced linearly by the welded LoG via \eqref{eq:Sweld} after integrating out the bulk profile (Hubbard–Stratonovich view of the double–trace deformation).
Solving $\mathcal S$ at tree level gives
\(
\langle \mathcal S\rangle \approx (m_{\mathcal S}^2)^{-1} J
\)
on long wavelengths; thus the effective Higgs mass receives the \emph{constant} imprint of the response,
\(
\Delta\mu_H^2(0)\sim \lambda_p\, J(0)/m_{\mathcal S}^2
\),
plus the anomaly–rectified DC from $\Im W$. 
Because the radion/scalar is in the physical kernel, its mixing with the DC trace is negligible at leading order; in particular, the Higgs scale is set by $(J,\lambda_p,m_{\mathcal S})$ and the anomaly rectifier, not by a light radion.

\subsubsection*{Long memory and retardedness}
\label{subsubsec:memory-ret}

For the semiclassical acceptance/retardedness package and its role in diagonalization, see
Sec.~\ref{sec:two-times-semiclassical}. In the present reduction we use only retardedness
and convexity of the SK kernels and the positivity of $\Im W$.

\subsubsection*{\texorpdfstring{Co-centering of LoG imprints on aeon branes across AdS$_5$}{Co-centering of LoG imprints on aeon branes across AdS5}}\label{subsubsec:uom-main-012}

\begin{proposition}[Cross–aeon center inheritance under the 5D SK functional]
\label{prop:aeon-center-inheritance}
Let $S^3$ be the receiver spatial slice with Laplace--Beltrami $-\Delta$ and heat kernel $K(s;\cdot,\cdot)$.
Assume:
\begin{enumerate}[leftmargin=1.8em,label=(A\arabic*)]
\item (Zonal LoG background) The KR head at the end of aeon $n$ carries a nonzero \emph{zonal} LoG background
\(
B_n(\Omega)=\Delta K(\eta^2;R_{\Omega_\star}^{-1}\Omega)
\)
about some center $\Omega_\star\in S^3$, with $\eta>0$ and spectral weights $B_{\ell 0}\propto \lambda_\ell e^{-\eta^2\lambda_\ell/2}>0$ for $\ell\ge 1$.
\item (Cross–aeon kernel) The retarded cross–aeon kernel $\mathcal K_{\rm cross}$ from the head of aeon $n$
to the initial dS brane of aeon $n{+}1$, when restricted to the low–$\ell$ sector relevant to welded LoGs,
is $SO(4)$–equivariant and positive semidefinite in the sense that, for any real $f$,
\[
\big\langle f,\ \mathcal K_{\rm cross}*f\big\rangle_{L^2(S^3)}\ \ge\ 0,
\]
and in harmonics it acts diagonally with nonnegative weights $\kappa_\ell\ge 0$ (not all zero).
\item (Zonal welded probe) The first welded spatial profile in aeon $n{+}1$ is a (macro LoG)$\ast$(micro Gaussian),
\(
\tau_{\rho,\sigma}(\cdot;\Omega)=\big(\Delta K(\rho^2;\cdot,\Omega)\big)\ast K(\sigma^2;\cdot,\Omega),
\)
centered at $\Omega\in S^3$, with $\rho,\sigma>0$ fixed (band–edge locked scales).
\end{enumerate}
Then the cross–aeon quadratic action contains the overlap
\[
\mathcal C(\Omega)\;:=\;\big\langle \tau_{\rho,\sigma}(\,\cdot\,;\Omega),\ \mathcal K_{\rm cross} * B_n\big\rangle_{L^2(S^3)},
\]
and $\mathcal C(\Omega)$ depends only on the geodesic distance $\delta:=\mathrm{dist}_{S^3}(\Omega,\Omega_\star)$ and is
\emph{strictly decreasing} in $\delta$ on $(0,\pi)$. Consequently, the action is uniquely minimized at $\Omega=\Omega_\star$,
i.e.\ the first welded pulse in aeon $n{+}1$ is centered at the same $\Omega_\star$ as $B_n$.
\end{proposition}

\begin{proof}
By $SO(4)$--equivariance, $\mathcal K_{\rm cross}*B_n$ is zonal about $\Omega_\star$. Writing
\(
\tau_{\rho,\sigma}(\cdot;\Omega)=\sum_{\ell\ge 1} a_\ell(\rho,\sigma)\,D^{\,\ell}_{\bullet 0}(R_\Omega)\cdot Y_{\ell\bullet}
\)
and
\(
\mathcal K_{\rm cross}*B_n=\sum_{\ell\ge 1}\kappa_\ell\,b_\ell(\eta)\,D^{\,\ell}_{\bullet 0}(R_{\Omega_\star})\cdot Y_{\ell\bullet}
\)
with $a_\ell>0$, $b_\ell>0$, and $\kappa_\ell\ge 0$, the overlap reduces (unitarity of $D^{\ell}$) to
\[
\mathcal C(\Omega)=\sum_{\ell\ge 1} a_\ell\,\kappa_\ell\,b_\ell\;P_\ell(\cos\delta),\qquad \delta=\mathrm{dist}_{S^3}(\Omega,\Omega_\star).
\]
Since at least one coefficient $a_\ell\kappa_\ell b_\ell$ is strictly positive and $|P_\ell(\cos\delta)|<1$
for $0<\delta<\pi$, the sum is strictly decreasing in $\delta$ on $(0,\pi)$ and attains its unique maximum
at $\delta=0$. In the quadratic action, this cross term enters with a negative sign, hence the action is uniquely minimized at $\Omega=\Omega_\star$.
\end{proof}




\subsubsection{Dictionary and explicit coefficients}
\label{subsubsec:dictionary}

\paragraph{Central charges and anomaly coefficient.}\label{par:uom-main-089}
\[
a=c=\frac{\pi \ell^3}{8G_5},\qquad
G_4=\frac{G_5}{\ell},\qquad
\gamma=\frac{a\,G_4}{12\pi}=\frac{\ell^2}{96}.
\]

\paragraph{Variance scaling on the LoG family.}\label{par:uom-main-090}
For welded temporal/spatial widths $(\Sigma,\sigma_t,\sigma_x)$ and effective amplitude fixed by the payload equality,
\begin{equation}
\Delta\!\langle R^2\rangle
=(8\pi G_4)^2\,\big\langle |\mathcal K|^2\big\rangle
\,C_{\rm LoG}\,
\frac{A^2}{\sigma_t\,\sigma_x^3},
\qquad
\Theta=\gamma\,\Delta\!\langle R^2\rangle,
\end{equation}
with $\langle|\mathcal K|^2\rangle$ the induced–gravity screened average (Yukawa form) and $C_{\rm LoG}=\mathcal O(1)$.

\subsubsection{From 5D least action to 4D P3: convexity, locks, and GP form}
\label{subsubsec:5D-to-4D}

\paragraph{Per–tick action.}\label{par:uom-main-091}
The 5D per–tick objective is the SK influence evaluated on the welded source,
\(
\mathcal A^{\rm (5D)}_k=\Re W[J_k]+\Im W[J_k],
\)
plus the source dwell minus record (which cancels at the QND aligned minimizer at fixed $\Lambda_k$).
Inserting the LoG ansatz and the payload equality (\S\ref{subsec:source-qnd-commutant}) yields the 4D receiver objective
\(
\Phi_{\rm BY}+\Phi_{\rm conf}+\Phi_{\rm RP}
\)
as in \eqref{eq:Jrec}, convex in the welded widths and (after eliminating amplitudes) jointly convex with $\Lambda_k$.
The co–centering and width locks follow from evenness and the Fisher orthogonality of the three unphysical blocks.

\paragraph{Outcome.}\label{par:uom-main-092}
Thus a \emph{single} 5D SK least–action principle \emph{exactly} reproduces:
(i) the 4D receiver convex penalties (Weyl/conf blocks) with geometry–fixed weights, 
(ii) the BY/anomaly rectifier $\Theta=\gamma\Delta\!\langle R^2\rangle$, and 
(iii) the tickwise conservation and monotone capacity increase of P2.
The radion resides in the physical kernel and is orthogonal (in Fisher metric) to the DC trace; the portal via the boundary GW fluctuation $\mathcal S$ sets the Higgs scale while keeping hierarchy control.

\medskip
\noindent\textbf{Remark (multi–brane/shadow sector).}
Older–aeon branes $\Sigma_{\rm old}^{(a)}$ contribute only through the retarded bulk Green's functions entering $G_R$ and $\mathcal N$; their localized, stabilized stresses appear as boundary conditions for the 5D solution and can be incorporated without altering the convex 4D reduction, provided their induced kernels remain reflection–even at the receiver split/flip and respect the same maximal symmetry to leading order.

\subsection{\texorpdfstring{Alternating heads, interior cells, and meta--causality in the AdS$_5$ limit}{Alternating heads, interior cells, and meta--causality in the AdS5 limit}}
\label{subsec:alternating-heads}

\paragraph{Setup and conventions.}\label{par:uom-main-097}
Let the bulk be ${\rm AdS}_5$ with radius \(\ell\) and 5D Planck scale \(M_5\), \(\Lambda_5=-6\,M_5^3/\ell^2\) (mostly–plus).
For a timelike codimension–one brane with induced metric \(\gamma_{\mu\nu}\) and tension \(T\), the generalized Israel equation reads
\begin{equation}
\big[K_{\mu\nu}-K\,\gamma_{\mu\nu}\big]\;=\;-\frac{T}{2M_5^3}\,\gamma_{\mu\nu}\,,
\end{equation}
so that the (regulated) near–brane 4D curvature obeys the familiar RS/KR relation
\begin{equation}
\Lambda_4\;=\;\frac{1}{12}\,\frac{T^2}{M_5^6}\;+\;\frac{1}{2}\,\frac{\Lambda_5}{M_5^3}\,,
\qquad 
T_{\rm crit}:=\frac{6\,M_5^3}{\ell}\ \Longrightarrow\ \Lambda_4=0\ \text{at }T=T_{\rm crit}.
\label{eq:L4-vs-T}
\end{equation}
We denote by \(k\equiv \ell^{-1}\) the local AdS curvature and encode the (A)dS branch by writing \(T=T_{\rm crit}-\delta T\), so that \(\mathrm{sign}(\Lambda_4)\) is the sign of \(-\delta T\).

\paragraph{Alternating localization (qualitative picture).}\label{par:uom-main-098}
We consider an alternating configuration of defect branes in AdS$_5$ that dynamically generates a
sequence of 4D aeons $\mathring{\mathcal M}_4$ with alternating effective curvature $\Lambda_4$ and
kinematic heads:
\begin{itemize}[leftmargin=1.6em]
\item \textbf{Interior RS1--type cells (shadow segments).} Between two subcritical positive--tension walls
at finite proper separation $d$ there is an RS1--type interval supporting a \emph{physical} stabilized
radion; these segments are interior/shadow cells and are not active receivers.
\item \textbf{Heads (KR caps).} The active receiver is a KR end--of--the--world (EOW) cap throughout the
split-ordered evolution; at each split it remains just UV--side of the newborn source
(Def.~\ref{def:Lambda-min}), with $\Lambda_{\rm R}>\Lambda_{\rm S}$. The Wilsonian cutoff
advances monotonically within the KR head under the SK/HJ wall balance; it may approach a terminal value
$\Lambda_\infty$ only after payload exhausts. There is no RS1 receiver stage.
\end{itemize}

\paragraph{Why the radion exists (interior) and persists (at a KR head).}\label{par:uom-main-099}
On any finite interval \([y_1,y_2]\) bounded by two \emph{dynamical} branes, small normal deformations
\(y\mapsto y+\xi(x)\) of the endpoints cannot be gauged away simultaneously: the gauge–invariant scalar
\[
\varphi_r(x)\ \propto\ \xi_2(x)-\xi_1(x)
\]
is the radion.\footnote{This is standard in the RS1/GW literature; see, e.g.,
\cite{RandallSundrum1999RS1,GoldbergerWise1999,CsakiGraesserKribs2000}.}
With a GW background \(\Phi=\Phi_0(y;d)+\cdots\) (bulk mass \(m\), brane potentials \(\lambda_i(\Phi)\)) one obtains an
effective potential \(V(d)\) that stabilizes \(\varphi_r\) at \(d=d_0\) with
\(m_r^2\sim\nu^2 k^2\,\epsilon\,e^{-2kd_0}>0\) (with \(\nu=\sqrt{4+m^2/k^2}-2\)),
and the mode lies in the \emph{physical} block of the Gaussian covariance (no conformality penalty)
\cite{GoldbergerWise1999,CsakiGraesserKribs2000}.

When an interior RS1 segment is adjacent to a KR head, the scalar does \emph{not} disappear in our
induced--gravity/Robin setup. Rather, the RS1 radion flows continuously to the KR \emph{boundary--distance} scalar,
the standard normalizable AdS\(_4\) scalar fluctuation of the brane embedding (the $E_0{=}4$ mode in the
Karch–Randall construction) with \(m_r^2=4k^2+\delta m_{\rm GW}^2>0\)
\cite{KarchRandall2001,ArkaniHamedPorratiRandall2001}. It again lies in the \emph{physical} tangent.
Its \emph{visible} coupling is suppressed by the induced–gravity Robin boundary condition:
\[
\varepsilon_{\rm BC}\;\equiv\;\frac{\|u_r(z_{\rm vis})\|}{\|u_r(z_{\rm sh})\|}
\;\simeq\;\frac{1}{1+s_{\rm vis}\,k\,L}\ \ll\ 1,
\]
with $s_{\rm vis}\gg1$ at late times, where $L\sim \Lambda^{-1}$ is the welded (spacetime) Gaussian scale.
Such mixed (Robin) boundary conditions for bulk/brane modes are the holographic image of double–trace deformations
and arise from adding a brane–localized kinetic term (induced 4D Einstein–Hilbert) at a finite Wilsonian cutoff
\cite{Witten2001MultiTrace,Skenderis2002HR,ArkaniHamedPorratiRandall2001}. At the welded band edge the Weyl–radion
mixing angle scales as \(\theta(\Lambda)=O(\varepsilon_{\rm BC}\Lambda^{-2})\), so the conformality penalty of this
physical scalar is \(\Phi_{\rm conf}\sim O(\varepsilon_{\rm BC}^2\Lambda^{-\,4})\), vanishing along the anti–Zeno refinement
(derivation in Sec.~\ref{subsubsec:fifthforce-robin}). By contrast, in the \emph{strict} RS2/KR convention with a non–dynamical
EOW boundary (fixed Dirichlet data) the boundary–distance bending is pure gauge and the scalar is absent; that limiting
case corresponds to sending the induced 4D term to infinity ($s_{\rm vis}\to\infty$) and is \emph{not} the induced–gravity
(Robin) regime adopted here \cite{KarchRandall2001,Skenderis2002HR}.


\paragraph{Lifecycle of a head (operational view).}\label{par:uom-main-100}
Let \(\Lambda\) denote the receiver’s Wilsonian cutoff and \(\Xi(a):=m_r\,{\cal D}(a)\) the (dimensionless)
product of the light scalar’s mass and the relevant (comoving) transfer distance \(\mathcal D(a)\) to the
preceding dS brane.\footnote{Under the SK/KMS coarse--graining, the retarded cross--brane kernel behaves at DC
and small $k$ as \(\varepsilon_{\rm DC}(a)\propto \exp[-\Xi(a)]\), with the heavy--mode screen and the welded
$k$--window already folded into \(\Xi\).} The evolution of an active head through an aeon can be summarized as follows:

\begin{enumerate}[leftmargin=2.4em]
\item \textbf{KR/AdS$_4$ split-side ordering.} The active receiver is a KR head slightly UV--side of the
newborn source with $\Lambda_{\rm R}>\Lambda_{\rm S}$ (Def.~\ref{def:Lambda-min}). The boundary--distance scalar
is physical with BC--suppressed visible overlap.

\item \textbf{KR refinement (monotone).} As payload is admitted, the cutoff advances
$\Lambda_{k+1}>\Lambda_k$ with possibly shrinking increments; the anomaly rectifier continues to raise
$\Lambda_4$ tick by tick.

\item \textbf{KR$\to$dS split/flip or terminal boundary.} When the first--crossing condition is met, the
receiver-dragged completion sector converts into a KR/dS source while the KR head itself continues as the
next aeon's KR/AdS receiver (Assumption~\ref{ass:split-ident}). If the source exhausts before crossing, the head
becomes the terminal EOW/UV boundary condition of the remaining AdS$_5$ domain.

\item \textbf{Iteration.} After the split the continuing KR receiver remains just above the newborn source
and the process iterates; RS1 dynamics appear only in interior/shadow segments.
\end{enumerate}

\noindent In this picture the \emph{boundary--distance} scalar remains in the physical sector at the head;
what changes across splits is its geometric interpretation and the magnitude of its visible coupling, set
by the curvature--dependent Robin stiffness. The portal mediator \(\mathcal S\) remains heavy and carries the
DC Higgs imprint across the head change, ensuring continuity of \(\Delta\mu_H^2=-\lambda_p\,J/m_{\mathcal S}^2\).

\subsubsection{Fifth–force limits, visible–brane Robin coupling, and compatibility with cross–brane transfer}
\label{subsubsec:fifthforce-robin}

\paragraph{What is shown and why.}\label{par:uom-main-101}
In this subsection we (i) parametrize the visibility of the light scalar through a curvature–dependent Robin boundary
condition on the visible brane; (ii) confront the resulting Newtonian--limit Yukawa correction with existing fifth–force
bounds; (iii) demonstrate, by an explicit one–parameter toy model, how the same choice of visible–brane stiffness
$s_{\rm vis}$ that \emph{suppresses} the local fifth force can be made \emph{compatible} with our dark–sector cross–brane transfer,
because the latter depends on the product of a \emph{visible} and a \emph{shadow} overlap, not on the former alone; and
(iv) explain why moderately large $s_{\rm vis}$ at late times is a natural outcome of induced gravity.

\paragraph{Fifth–force observable at long range.}\label{par:uom-main-102}
For a light scalar of (proper) mass $m_r$ coupled universally to visible–sector matter via a brane operator
$\mathcal{O}_{\rm vis}$ with effective coupling $c_{\rm vis}$,\footnote{In the Higgs portal normalization,
$\mathcal{L}\supset -c_{\rm vis}\,\varphi_r\,\mathcal O_{\rm vis}$ with $c_{\rm vis}\equiv \frac{\lambda_p\,v}{m_{\mathcal S}^2}\,\varepsilon_{\rm BC}(\Lambda)$
after integrating out the heavy $\mathcal S$; for the 4D non–relativistic source $\mathcal O_{\rm vis}\sim T^\mu{}_\mu\sim m\,\delta^{(3)}(\vec x)$.}
the static potential between two unit masses at large separation $r\gg m_r^{-1}$ carries a Yukawa correction
\begin{equation}
V(r)\;=\;-\frac{G_N m_1 m_2}{r}\,\bigg[\,1\;+\;\alpha_{\rm vis}\,e^{-m_r r}\,\bigg]
\qquad\text{with}\quad
\alpha_{\rm vis}\;\equiv\;\frac{c_{\rm vis}^2}{4\pi G_N\,\mu_{\star}^2}\,,
\label{eq:V-Yuk}
\end{equation}
where $\mu_\star$ is the (model–dependent) normalization scale relating $c_{\rm vis}$ to the dimensionless coupling
in the Newtonian limit.\footnote{For a pure conformal scalar coupled via $-\beta\,\varphi\,T$ one has
$\alpha\simeq 2\beta^2$ up to screening–model dependent factors; here we absorb such factors into $\mu_\star$.
The brane–localized induced gravity shifts $G_N=(8\pi M_4^{-2}(\Lambda))^{-1}$, but does not alter the Yukawa \emph{structure}.}
Laboratory and Solar–System tests bound the amplitude $\alpha$ at long ranges by $|\alpha(\lambda\!\gtrsim\!\mathrm{AU})|\lesssim 10^{-3}$,
while at astrophysical scales (kpc–Mpc) dynamical tests of galaxy clusters and large–scale structure imply
$|\alpha(\lambda\sim{\rm kpc\!-\!Mpc})|\lesssim {\cal O}(10^{-2})$.\footnote{See, e.g. torsion–balance constraints
(e.g.\ E\"ot–Wash) and fifth–force reviews; at kiloparsec scales, bounds can be phrased as $|\Delta \Phi|/\Phi_{\rm N}\lesssim$ few percent,
consistent with cluster dynamics and weak–lensing reconstructions. We will use these as fiducial target numbers.}

\paragraph{Visible–brane Robin overlap and curvature dependence.}\label{par:uom-main-103}
In our setup the \emph{only} handle on $\alpha_{\rm vis}$ is the \emph{visible} overlap of the light scalar with the
active head. At the welded band edge, the radial profile $u_r(z)$ of the (physical) light mode evaluated on the visible KR cap at cutoff $\Lambda$ is mapped by an \emph{induced–gravity Robin} operator $L_{\rm vis}=f(H_{{\rm sp},\Delta(\Lambda)})$.
The overlap scales as
\begin{equation}
\varepsilon_{\rm BC}(\Lambda)\;\equiv\;\frac{\|u_r(z_{\rm vis})\|}{\|u_r(z_{\rm sh})\|}\;\simeq\;\frac{1}{1+s_{\rm vis}(\Lambda)\,\big[k(\Lambda)\,L(\Lambda)\big]}\,,
\qquad s_{\rm vis}(\Lambda)\gg1\ {\rm\ at\ late\ times},
\label{eq:eps-BC}
\end{equation}
with $k(\Lambda)$ the local AdS$_4$ curvature scale ($k^2\simeq -\Lambda_4/3$ in AdS),
$L(\Lambda)$ the welded spacetime Gaussian width, and $s_{\rm vis}(\Lambda)$ the dimensionless stiffness induced by the running
4D Einstein--Hilbert counterterm: $M_4^2(\Lambda)\simeq M_5^3\ell+\kappa\,\Lambda^2$ so that $s_{\rm vis}(\Lambda)\propto M_4^2(\Lambda)/k(\Lambda)$.
Equation~\eqref{eq:eps-BC} encodes the physically transparent fact that a stiffer visible cap suppresses the light scalar’s
visible wavefunction as the cutoff approaches the curvature scale.\footnote{At the \emph{strict} RS2/KR limit with a fully fixed (Dirichlet) boundary metric
one may take $s_{\rm vis}\to\infty$, in which case the boundary–distance mode is a pure gauge and $\varepsilon_{\rm BC}\to 0$. In our induced–gravity
regime $s_{\rm vis}$ is \emph{large but finite}, and the mode remains in the physical block with BC–suppressed overlap.}

Consequently the effective (dimensionless) fifth–force strength scales as
\begin{equation}
\alpha_{\rm vis}(\Lambda)\;=\;\alpha_\circ\,\varepsilon_{\rm BC}^2(\Lambda)\;\simeq\;
\alpha_\circ\,\frac{1}{\bigl(1+s_{\rm vis}(\Lambda)\,k(\Lambda)\,L(\Lambda)\bigr)^2}
\ \ \xrightarrow[\ \Lambda\to \Lambda_\infty\ ]{}\
\alpha_{\rm vis}^{\rm (late)}\;\simeq\;\frac{\alpha_\circ}{\bigl(1+s_{\rm vis}(\Lambda_\infty)\,k(\Lambda_\infty)\,L_\infty\bigr)^2}\,,
\label{eq:alpha-vis}
\end{equation}
where $\alpha_\circ\equiv c_{\rm vis}^2/(4\pi G_N\mu_\star^2)$ absorbs the microscopic portal normalization.
Here $L_\infty\equiv L(\Lambda_\infty)$ is the welded LoG scale evaluated at the terminal cutoff. For a broad
range of models $L_\infty\sim \xi/k(\Lambda_\infty)$ with $\xi={\cal O}(1)$, so
\(\alpha_{\rm vis}^{\rm (late)}\sim \alpha_\circ/(1+\xi\,s_{\rm vis})^2\).

\paragraph{Compatibility with cross–brane transfer (dark sector signal).}\label{par:uom-main-104}
The dark–sector transfer amplitude at DC scales as\footnote{At small $k$ and for short inter–brane separations $d$ one has
\(\varepsilon_{\rm DC}(a)\;\propto\; g_{\rm vis}(a)\,g_{\rm sh}(a)\;\frac{q(0)}{q(0)+r(\Lambda)\,k^2}\,e^{-q(0)\,d}\) with $q(0)\equiv m_r$ and the induced–gravity screen $r(\Lambda)\sim M_4^2(\Lambda)/(2M_5^3)$.}
\begin{equation}
\varepsilon_{\rm DC}(a)\;\propto\; g_{\rm vis}(\Lambda)\,g_{\rm sh}(\Lambda)\,e^{-m_r\,d}\,\frac{m_r}{m_r+r(\Lambda)\,k^2}\,,
\label{eq:epsDC}
\end{equation}
and it is this \emph{product} that feeds the uniform (apparent) matter fraction $\Omega_X(a)\propto \varepsilon_{\rm DC}(a)\,a^{-3}/E^2(a)$ (Sec.~\ref{subsec:shadow-brane}).
The $g_{\rm vis}(\Lambda)$ inherits precisely one power of the visible overlap $\varepsilon_{\rm BC}(\Lambda)$:
\(
g_{\rm vis}(\Lambda)=g_\circ\,\varepsilon_{\rm BC}(\Lambda)
\),
where $g_\circ$ collects the (fixed) portal normalization, mixing, and 4D canonical normalization factors, while the
\emph{shadow} overlap $g_{\rm sh}(\Lambda)$ does \emph{not} suffer the $s_{\rm vis}$ suppression (the shadow cap sits at a deeper curvature, hence milder induced stiffness at the same $L$).
It follows that tightening the visible Robin stiffness to suppress the \emph{local} Yukawa correction, $\alpha_{\rm vis}\propto \varepsilon_{\rm BC}^2$,
\emph{does not} destroy the cross–brane signal: the latter scales with the product \(g_{\rm vis}\,g_{\rm sh}\sim g_\circ\,\varepsilon_{\rm BC}\,g_{\rm sh}\),
and the shadow overlap $g_{\rm sh}$ can remain \(\mathcal O(1)\) in the AdS region (cf.\ Sec.~\ref{subsec:alternating-heads}).

\paragraph{A one–parameter toy model with explicit curvature ratio.}\label{par:uom-main-105}
To make the parametrics transparent and exhibit where the \emph{curvature ratio} enters the two constraints
(fifth–force bound and dark–sector transfer), we adopt the following minimal late--time parametrization for the
visible head and its nearest shadow cap:
\begin{equation}
k_{\rm vis}(\Lambda_\infty)=H_\star,\qquad
L_\infty=\frac{\xi}{H_\star},\qquad
s_{\rm vis}(\Lambda_\infty)=:s_\star\gg 1,
\qquad
\chi\;:=\;\frac{k_{\rm sh}}{k_{\rm vis}}\;=\;\sqrt{\frac{R_{\rm sh}}{R_{\rm vis}}}\;>\;1,
\label{eq:KR-toy-defs}
\end{equation}
with \(\xi=\mathcal O(1\text{--}3)\) the welded LoG aspect, and \(\chi\) the \emph{curvature ratio} (shadow vs visible).\footnote{At each
head late cutoff the Wilsonian identification sets \(\Lambda_\infty^{\rm (vis)}\!\sim\!k_{\rm vis}\) and
\(\Lambda_\infty^{\rm (sh)}\!\sim\!k_{\rm sh}\); we keep \(\xi=L_\infty\,k\) fixed to an $O(1)$ shape constant, as in
Sec.~\ref{subsec:alternating-heads}.}
Two independent requirements constrain \((\chi,s_\star)\):

\smallskip
\emph{(i) Fifth–force bound on the visible head.} Using \eqref{eq:alpha-vis} at late times,
\[
\alpha_{\rm vis}^{\rm (late)}\;\simeq\;\frac{\alpha_\circ}{(1+\xi\,s_\star)^2}\ \le\ \alpha_{\rm lim}
\qquad\Longrightarrow\qquad
\boxed{\ s_\star\ \ge\ s_{\rm min}(\alpha_\circ,\alpha_{\rm lim},\xi)\ :=\ \frac{1}{\xi}\Big(\sqrt{\frac{\alpha_\circ}{\alpha_{\rm lim}}}-1\Big)\ },
\]
where \(\alpha_\circ\) encodes the microscopic normalization (portal, canonical factors) and \(\alpha_{\rm lim}\) is the observational
upper bound at the relevant range (e.g. \(\alpha_{\rm lim}\!\sim\!10^{-5}\) for massless–like long range, \(\sim 10^{-2}\) at
kpc scales). Equation \eqref{eq:alpha-vis} shows \(s_\star\) enters quadratically in \(\alpha_{\rm vis}\) via the Robin overlap
\(\varepsilon_{\rm BC}=(1+\xi s_\star)^{-1}\).

\smallskip
\emph{(ii) Cosmological DC transfer needed for the mirage fraction.} The uniform “mirage” obeys
\[
\Omega_{X,0}\;=\;\varepsilon_{\rm DC,0}\;\Omega_{m,{\rm sh}}\;\frac{H_{\rm sh}^2}{H_0^2}
\;=\;\varepsilon_{\rm DC,0}\;\Omega_{m,{\rm sh}}\;\chi^2\,\frac{H_{\rm vis}^2}{H_0^2}\,,
\]
and, for one dominant light mode and short separation, \(\varepsilon_{\rm DC,0}\propto g_{\rm vis}\,g_{\rm sh}\) with
\(g_{\rm vis}=g_\circ\,\varepsilon_{\rm BC}=(g_\circ)/(1+\xi s_\star)\) and \(g_{\rm sh}\sim\mathcal O(1)\) on the
deeper (shadow) cap. Inserting the fifth–force limit \(g_{\rm vis}\le g_\circ\sqrt{\alpha_{\rm lim}/\alpha_\circ}\) gives the
\emph{derived lower bound} on the curvature ratio (cf.\ Sec.~\ref{subsec:shadow-brane}):
\[
\boxed{\ \chi^2\ \ge\ \frac{\Omega_{X,0}}{\Omega_{m,{\rm sh}}\,N_{\rm eff}\,\sqrt{\alpha_{\rm lim}}}\ },
\]
where \(N_{\rm eff}\) counts light modes (branes) that add (coherently) to the DC transfer (we take \(N_{\rm eff}=1\) in this toy example).
Thus, \(\chi\) is fixed by \(\{\Omega_{X,0},\Omega_{m,{\rm sh}},\alpha_{\rm lim}\}\); \(s_\star\) is then chosen above
\(s_{\rm min}(\alpha_\circ,\alpha_{\rm lim},\xi)\).

\paragraph{A concrete numerical slice.}\label{par:uom-main-106}
Choose
\[
\Omega_{X,0}=0.01,\quad \Omega_{m,{\rm sh}}=0.3,\quad N_{\rm eff}=1,\quad \alpha_{\rm lim}=10^{-5},\quad
\alpha_\circ=10^{-1},\quad \xi=2.
\]
Then the DC transfer bound gives
\[
\chi^2\ \ge\ \frac{0.01}{0.3\times 1\times \sqrt{10^{-5}}}\ \approx\ 10.5\qquad\Longrightarrow\qquad \chi\ \gtrsim\ 3.24,
\]
i.e. \(R_{\rm sh}/R_{\rm vis}\gtrsim 10.5\) and \(k_{\rm sh}/k_{\rm vis}\gtrsim 3.24\). The fifth–force bound requires
\[
s_\star\ \ge\ s_{\rm min}\ =\ \frac{1}{2}\Big(\sqrt{\frac{10^{-1}}{10^{-5}}}-1\Big)
\ =\ \frac{1}{2}\,(100-1)\ \approx\ 49.5.
\]
Taking, e.g., \(s_\star=50\) (moderately large, as expected at late times with induced gravity) and
\(g_{\rm sh}=\mathcal O(1)\), we obtain
\[
\alpha_{\rm vis}^{\rm (late)}\ \simeq\ \frac{\alpha_\circ}{(1+ \xi s_\star)^2}
\ =\ \frac{0.1}{(1+100)^2}\ \approx\ 9.8\times 10^{-6}\ \le\ \alpha_{\rm lim},
\]
comfortably below the long–range bound, while the DC transfer product is
\[
\varepsilon_{\rm DC,0}\ \propto\ g_{\rm vis}\,g_{\rm sh}
\ =\ \frac{g_\circ}{1+\xi s_\star}\,\underbrace{g_{\rm sh}}_{\mathcal O(1)}
\ \approx\ \frac{g_\circ}{101}\ \sim\ 10^{-2}\,,
\]
sufficient to supply a percent–level apparent fraction (and a few percent at late times; cf.\ Sec.~\ref{subsec:shadow-brane-dm-phenom}).
In words: with the curvature ratio \(\chi\gtrsim 3.24\) demanded by the mirage fraction and a visible stiffness
\(s_\star\simeq 50\) demanded by the fifth–force bound, the same choice of parameters suppresses the \emph{local} Yukawa
quadratically (via \(\alpha_{\rm vis}\propto \varepsilon_{\rm BC}^2\)) and preserves the \emph{inter–brane} transfer linearly
(via \(g_{\rm vis}\,g_{\rm sh}\propto \varepsilon_{\rm BC}\)), exactly as required by observations.


\paragraph{Naturalness of moderately large \(s_{\rm vis}\) at late times.}\label{par:uom-main-107}
The stiffness $s_{\rm vis}(\Lambda)$ is not an arbitrary knob; in the induced–gravity completion it is set by the square of the running 4D Planck mass relative to the local curvature and the welded time/space coarse–graining,
schematically
\[
s_{\rm vis}(\Lambda)\;\sim\;\frac{M_4^2(\Lambda)}{M_5^3}\,k(\Lambda)\,L(\Lambda)\;\sim\;(1+\tilde\kappa\,\Lambda^2\ell)\,k(\Lambda)\,L(\Lambda)\,,
\]
with \(\tilde\kappa={\cal O}(1)\) a scheme–dependent Wilson coefficient. As the receiver is anti–Zeno–driven toward its terminal cutoff, \(\Lambda\to\Lambda_\infty\sim k\), the dimensionless product \(kL\) enters only through the welded LoG scale \(\xi={\cal O}(1)\), while the factor \(M_4^2(\Lambda)\) has already been pushed to its saturated, large value. 
\paragraph{Why \(s_\star\sim 10\!-\!100\) is natural at late times.}\label{par:uom-main-108}
The dimensionless stiffness on the visible cap is set by the \emph{induced} 4D curvature term at the terminal cutoff:
\[
s_{\rm vis}(\Lambda_\infty)\ \sim\ \Big(\frac{M_4^2(\Lambda_\infty)}{M_5^3}\Big)\,k_{\rm vis}\,L_\infty
\ \sim\ \Big(\frac{M_5^3\ell+\kappa\,\Lambda_\infty^2}{M_5^3}\Big)\,\xi\,,
\qquad \Lambda_\infty\sim k_{\rm vis}.
\]
Since \(\kappa\,\Lambda_\infty^2\sim \kappa\,k_{\rm vis}^2\) is \emph{already} large at late times (the induced gravity
has saturated), one generically finds \(s_\star\sim \xi\,(1+\kappa\,k_{\rm vis}^2\ell^2)=\mathcal O(10\text{--}50)\) for modest
$k_{\rm vis}\ell$ and \(\xi=\mathcal O(1\text{--}3)\), without any fine tuning of the brane potential.
This places \(s_\star\) precisely in the band needed by the fifth–force bound, while the curvature ratio \(\chi\)
(derived from the mirage fraction) simultaneously ensures that the shadow overlap remains $\mathcal O(1)$.


\medskip
\noindent
In summary, the observational upper bounds on long–range fifth forces constrain the visible–brane stiffness through
Eqs.~\eqref{eq:eps-BC}–\eqref{eq:alpha-vis}, but the same stiffness enters \emph{linearly} into the cross–brane transfer amplitude and quadratically into the local fifth–force amplitude. A moderately large $s_{\rm vis}$ that is \emph{natural} at late times (by induced gravity) suppresses the local Yukawa correction beneath current bounds while leaving the inter–brane signal at the required few–percent level, thereby preserving the viability of the dark–sector mechanism described above.



\subsubsection*{\texorpdfstring{Stability of the AdS$_5$ limit with infinitely many aeons}{Stability of the AdS5 limit with infinitely many aeons}}
\label{subsubsec:AdS5-limit-stability}

\paragraph{(S1) Cumulative tension.}\label{par:uom-main-109}
Let the head of the \(i\)-th aeon have tension \(T_i=T_{\rm crit}-\delta T_i\) (same renormalization scheme across aeons). 
Consider the limit in which the receiver approaches the AdS$_5$ boundary by inserting more heads ``ahead''. 
The \emph{asymptotic} AdS boundary data remain intact if and only if the total deviation of the bending from critical remains finite:
\begin{equation}
\sum_{i=1}^{\infty} |\delta T_i|\ <\ \infty\qquad\Longleftrightarrow\qquad
\text{no change of AdS asymptotics (no over–bending)}.
\label{eq:cum-tension}
\end{equation}
If the series \(\sum \delta T_i\) diverges with a fixed sign, one attempts to replace the AdS boundary by a different asymptotic (``over–bent'' geometry), which generically leads to a big–crunch–type pathology.
Thus, \emph{cumulative subcriticality must be summable}.%
\footnote{This statement is independent of whether the sequence of heads is finite or countably infinite: what matters is summability in the common counterterm scheme that defines \(T_{\rm crit}\).}

\paragraph{(S2) Non–perturbative channels.}\label{par:uom-main-110}
Brown–Teitelboim and CDL brane nucleation are non–perturbative. With positive–tension walls in AdS and stabilized moduli, such channels are either absent (no appropriate instanton) or exponentially suppressed. Their rates are not enhanced by the mere presence of additional stabilized dS heads ``ahead''.

\paragraph{(S3) Why we keep the AdS$_5$ boundaries \emph{static}.}\label{par:uom-main-111}
In our framework the AdS$_5$ \emph{UV boundary} and, when present, the \emph{outer head} (EOW/KR cap for the active light sector) are treated as fixed boundary conditions. Allowing a \emph{true} “moving AdS$_5$ boundary’’ is either a different theory or generically unhealthy:

\begin{itemize}[leftmargin=1.4em]
\item \textbf{UV boundary motion $\Rightarrow$ mixed (Neumann) AdS boundary conditions.}
Promoting the boundary metric to a dynamical field corresponds to Neumann/mixed BC (a CFT coupled to boundary gravity). This changes the dual theory: the Brown–York stress tensor is no longer conserved by the CFT alone (there is energy exchange with boundary gravity), and stability/positivity depend on the chosen boundary action. Certain “designer–gravity’’ choices are \emph{unstable} (unbounded below energy) \cite{CompereMarolf2008,HertogHorowitz2004,BalasubramanianKraus1999}. In any case this is not Wilsonian RG (which keeps the UV boundary fixed) \cite{HeemskerkPolchinski2011,FaulknerLiuRangamani2011}.
\item \textbf{EOW/head motion $\Rightarrow$ braneworld subtleties.}
Treating the outer head as a dynamical brane is consistent only on the \emph{normal} branch (positive induced $M_4^2$ and stabilized moduli). The wrong branch is the well–known self–accelerating DGP/RS2–type solution with a \emph{ghost} (wrong–sign kinetic mode) \cite{LutyPorratiRattazzi2003,Koyama2005}. We therefore avoid an independently propagating “bending’’ field on the EOW and, in the induced–gravity (Robin) regime, admit only the healthy \emph{boundary–distance} mode (the $E_0\!=\!4$ AdS$_4$ scalar) in the \emph{physical} block with BC–suppressed visible overlap (Sec.~\ref{subsubsec:radion-orth}).
\item \textbf{What we \emph{do} allow to “move”.}
Sliding $\Lambda$ is the holographic RG depth (Heemskerk–Polchinski / Faulkner–Liu–Rangamani) and is governed by the SK/HJ wall balance; it does \emph{not} violate Lorentzian energy conservation and, in the exhaustion branch, approaches a terminal value $\Lambda_\infty$ only after accepted payload ceases \cite{HeemskerkPolchinski2011,FaulknerLiuRangamani2011}.
\end{itemize}

\noindent
For these reasons we keep the AdS$_5$ UV boundary \emph{static}. This avoids spurious energy non–conservation at infinity and ghost pathologies at the head, while the only genuine arrow of time remains the CPTP RG semigroup in meta–time (Sec.~\ref{subsubsec:meta-causality}).


\medskip
\noindent
\emph{Conclusion.} Under (S1)–(S3), taking the aeon sequence to the AdS$_5$ limit is consistent: the bulk stays AdS, interior intervals have healthy stabilized radions, and the outer head is a non–dynamical EOW brane.

\subsubsection*{\texorpdfstring{Cumulative tension \(\Rightarrow\) arbitrarily small \(\Lambda_4\)}{Cumulative tension => arbitrarily small Lambda_4}}
\label{subsubsec:cum-tension-l4}

On a maximally symmetric head with tension \(T=T_{\rm crit}-\delta T\), \eqref{eq:L4-vs-T} gives
\begin{equation}
\Lambda_4(T)\;=\; 
\underbrace{\Big(\frac{T_{\rm crit}}{6M_5^3}\Big)^2-\frac{1}{\ell^2}}_{=\,0}
\;+\;\underbrace{\frac{T_{\rm crit}}{36 M_5^6}}_{=:c_1}\,(-2\delta T)
\;+\;\underbrace{\frac{1}{12M_5^6}}_{=:c_2}\,(\delta T)^2
\;=\;-\,c_1\,\delta T\;+\;c_2\,(\delta T)^2.
\label{eq:L4-expansion}
\end{equation}
Thus, near criticality the 4D curvature is \emph{linearly} controlled by \(\delta T\) (with a small quadratic correction). 
In an alternating stack, the effective \emph{net} deviation at the current head is
\(
\delta T_{\rm eff}=\sum_{i=1}^{N}\delta T_i
\)
in the common scheme. If the infinite tail is summable, \(\delta T_{\rm eff}\) can be made \emph{arbitrarily small}, hence
\begin{equation}
|\Lambda_4|\;\le\; c_1\,|\delta T_{\rm eff}|+c_2\,\delta T_{\rm eff}^2\ \xrightarrow[\ \delta T_{\rm eff}\to 0\ ]{}\ 0.
\end{equation}
This makes precise how the alternating sequence allows one to realize \emph{arbitrarily small de Sitter curvature} at the flip head without ad hoc fine–tuning on a single brane: the smallness arises as a \emph{summable cumulative deviation} from criticality across prior branes.\footnote{In our full model the anomaly rectifier adds a positive shift to the DC trace balance, \(\gamma\langle R^2\rangle\) with \(\gamma=\ell^2/96\). Since \(\gamma\) is fixed by holography, its tickwise contribution is controlled by the welded LoG variance and can be made small for ticks at larger \(\bar\Lambda_k\) along the monotone trajectory; it does not obstruct the \(\delta T_{\rm eff}\to 0\) limit.}

\subsubsection*{\texorpdfstring{Receiver \(\to\) AdS$_5$ boundary and meta–causality}{Receiver -> AdS5 boundary and meta-causality}}
\label{subsubsec:meta-causality}

\paragraph{P1 vs.\ unitarity in the bulk.}\label{par:uom-main-112}
The 5D bulk evolution is unitary; by Postulate P1, \emph{unitary evolution cannot produce classical facts}. 
Facts in our construction arise via anomaly–rectified DC response in the trace sector and via coarse–graining at finite cutoff on the receiver boundary. 
Hence the arrow that creates facts cannot be the Lorentzian (bulk) time—it is the \emph{boundary RG} arrow.

\paragraph{Boundary RG as meta–time.}\label{par:uom-main-113}
Let \(\Lambda\) denote the receiver cutoff, and set \(u:=\ln(\Lambda_{\rm UV} / \Lambda)\) as ``RG time''.
We treat \(u\) as a depth coordinate for the UV\(\to\)IR CPTP semigroup at fixed \(\Lambda\), while the scheduler
relocates the cutoff surface \(\Lambda(t)\) toward the UV; the semigroup direction remains UV\(\to\)IR even as
\(u(t)\) decreases along accepted ticks. Integrating out UV shells generates a CPTP semigroup on boundary states,
\[
\partial_u \rho_u\;=\;-\,i[H_u,\rho_u]\;+\;\mathcal D_u(\rho_u)\,,\qquad \mathcal D_u\ \ge\ 0\ \ (\text{GKLS}),
\]
with \(\Im W[u]\ge 0\) (influence phase convexity) provided the SK kernels are retarded and obey KMS positivity. 
\emph{Meta–causality} is the orientation \(u\!\nearrow\) (UV \(\to\) IR): it is \emph{non–unitary}, monotone, and the only arrow that can \emph{create} facts in the pre–flip epoch.

\paragraph{Tickwise realization.}\label{par:uom-main-114}
Within a tick, the welded map is regulated at a representative cutoff \(\bar\Lambda_k\); meta–time flows from the AdS boundary down to the cutoff surface (coarse–graining). Across ticks, one may increase capacity (\(\Lambda_{k+1}>\Lambda_k\) for any accepted tick), but this does not reverse the arrow nor undo produced facts.
In the formal limit \(\Lambda\to\infty\) (receiver \(\to\) exact AdS boundary), the RG path length shrinks to zero: \(\mathcal D_u\to 0\), \(\Im W\to 0\), and the meta–arrow trivializes—consistent with P1.

\paragraph{Self–caused vs.\ uncaused.}\label{par:uom-main-115}
The completed AdS$_5$ problem is a boundary–value problem: asymptotic data at the AdS boundary (plus EOW heads and stabilized interior moduli) \emph{determine} the Lorentzian evolution.
If one insists the Universe be ``uncaused'' in the boundary sense, bulk unitarity leaves P1 unanswered. 
By contrast, taking the \emph{boundary RG} as the fundamental meta–time makes the Universe \emph{self–caused}: boundary conditions (UV data) and their CPTP flow \(\rho_{\rm UV}\mapsto \rho_\Lambda\) generate the facts that seed the next aeon’s dynamics.
In this sense, the meta–arrow points \emph{from UV to IR} along the boundary and is more fundamental than the bulk Lorentzian arrow.

\subsubsection*{Remarks on induced gravity, anomaly, and bookkeeping}
\label{subsubsec:bookkeeping}

\begin{itemize}[leftmargin=1.6em]
\item \textbf{Induced gravity vs.\ counterterms.} We adopt \(G_4=G_5/\ell\) as coming from the standard holographic counterterm \(\sim M_5^3(\ell/2) \!\int \!\sqrt{-\gamma}\,R[\gamma]\). No \emph{additional} brane–localized Einstein–Hilbert term is inserted, so we avoid double counting.
\item \textbf{Anomaly rectifier.} For holographic CFT\(_4\) with \(a=c=\pi \ell^3/(8G_5)\) and \(G_4=G_5/\ell\), the rectifier in the DC trace equation is
\(
\gamma=\dfrac{a G_4}{12\pi}=\dfrac{\ell^2}{96}\,.
\)
This coefficient is geometry–fixed (no free knob) and multiplies the curvature–variance injection of welded LoG pulses.
\item \textbf{Radion/scalar and payload sector.} Radions/KR scalars live in the \emph{physical kernel} of the receiver quadratic form (Fisher–orthogonal to Weyl, RP, and conformal mismatches at the maximally symmetric head). Payload degrees of freedom (records) remain and are transported tickwise subject to P2/P3.
\end{itemize}

\subsection{Higgs mechanism from the anomaly–LoG portal (with stabilized radion)}
\label{subsec:higgs-from-log}

\paragraph{Motivation and P1.}\label{par:uom-main-116}
By Postulate~P1 (``unitarity cannot produce facts''), \emph{no} unitary AdS$_5$ or AdS$_4$ evolution
can generate a symmetry–broken vacuum from a symmetric one. Hence any spontaneous breaking in the
low–energy dS$_4$ EFT (in particular, electroweak symmetry breaking, EWSB) must be \emph{imprinted}
by a non–unitary step: the pre–flip welded export and its anomaly rectification, followed by the flip
and freeze. In the receiver description, the only ingredients that survive the flip and can assign a
\emph{constant} (DC) Higgs mass parameter are (i) the DC response of the trace sector produced by the
LoG–shaped welded source and (ii) local, brane--resident quadratic operators that remain on the KR head.
We implement this via an \emph{auxiliary--scalar (portal) mechanism},
where the auxiliary is precisely the brane–localized fluctuation that descends from the bulk
Goldberger–Wise (GW) stabilizer.

\subsubsection*{The portal field and the welded source}\label{subsubsec:uom-main-013}
Let $\mathcal S$ be a real brane–localized scalar (mass dimension $1$) living on the receiver brane,
with quadratic action and a linear source $J_{\rm LoG}$ inherited from the welded trace sector:
\begin{equation}
\label{eq:S-portal-loc}
\mathcal L_{\mathcal S}\;=\;\frac12(\partial \mathcal S)^2-\frac12 m_{\mathcal S}^2\,\mathcal S^2
\;+\;J_{\rm LoG}(x)\,\mathcal S\,,
\qquad
\mathcal L_{\rm portal}\;=\;-\,\lambda_p\,\mathcal S\,H^\dagger H,
\end{equation}
with $\lambda_p$ of mass dimension $1$.
The welded source is the LoG projection of the (Brown–York/\,$d{=}4$) trace response:
\begin{equation}
\label{eq:JLoG-def}
J_{\rm LoG}(x)
\;:=\;\kappa_{\rm HS}\;\big(\mathrm{LoG}\star T^\mu{}_\mu{}^{\rm (BY)}\big)(x),
\qquad
T^\mu{}_\mu{}^{\rm (BY)}\ =\ \underbrace{\gamma\,R^2}_{\text{anomaly rectifier}}+\ \text{(conformal\;derivatives)}\,,
\end{equation}
where $\gamma=\dfrac{a\,G_4}{12\pi}=\dfrac{\ell^2}{96}$ (cf.\ Sec.~\ref{subsec:receiver-anomaly})
and $\kappa_{\rm HS}$ is the (fixed) Hubbard–Stratonovich coefficient relating the double–trace
deformation in the welded channel to the linear source.\footnote{Dimension analysis:
$[\mathcal S]=1$, $[H^\dagger H]=2$, so $[\lambda_p]=1$ and $[J_{\rm LoG}]=3$, as required
for a local 4D Lagrangian density.}

The Higgs potential is
\(
V_H(H)=\lambda_H (H^\dagger H)^2+\mu_{H,0}^2\,H^\dagger H
\)
with \emph{non–negative} bare mass $\mu_{H,0}^2\ge 0$ (unbroken in the absence of the imprint).

\subsubsection*{Radion sector: stabilization and mixing control}\label{subsubsec:uom-main-014}
We combine here the agnostic persistence proposition with its GW realization.

\begin{proposition}[Portal persistence and freeze]
\label{prop:persistence-freeze}
Assume:
\begin{enumerate}[label=(A\arabic*),leftmargin=1.8em,itemsep=2pt]
\item \emph{Positive quadratic kernel.} The brane scalar $\mathcal S$ is governed (around the
post–flip vacuum) by a strictly positive \emph{quadratic} action with local kernel
$K_{\mathcal S}:=-(\Box)+m_{\mathcal S}^2$ (Minkowski signature) with $m_{\mathcal S}>0$.
Equivalently, its free two–point has a Källén–Lehmann representation with nonnegative spectral
density and a mass gap $m_{\mathcal S}$.
\item \emph{LoG/anomaly drive with DC rectification.} The welded LoG history produces a nonzero
DC component
\[
J_0\;:=\;\lim_{T,V\to\infty}\frac{1}{2T\,\mathrm{Vol}(V)}\int_{-T}^{T}\!\!dt\int_{V}\!d^3x\;J_{\rm LoG}(t,\mathbf x)\ \neq\ 0,
\]
originating from the anomaly rectifier $\propto \gamma\,R^2$ even when each tick–source has zero mean.
\item \emph{Flip (driver off) and SK positivity.} For $t>t_{\rm flip}$ there is no new drive with
spatial structure: $J_{\rm LoG}(t,\mathbf x)$ has no $k>0$ components, and the Schwinger–Keldysh
imaginary part satisfies $\Im W\ge 0$, so all $k>0$ transients of the linear response decay.
\end{enumerate}
Then the late–time Higgs mass shift is \emph{constant} and equals
\begin{equation}
\label{eq:Delta-muH-const}
\Delta\mu_H^2\ :=\ -\,\lambda_p\,\frac{J_0}{m_{\mathcal S}^2}\,,
\end{equation}
so the electroweak vacuum expectation satisfies
\begin{equation}
\label{eq:v2-from-portal}
v^2\;=\;\frac{-\,\mu_{H,0}^2\ -\ \Delta\mu_H^2}{\lambda_H}
\;=\;\frac{1}{\lambda_H}\left(\lambda_p\,\frac{J_0}{m_{\mathcal S}^2}-\mu_{H,0}^2\right),
\qquad v\neq 0\ \ \Leftrightarrow\ \ \lambda_p\,\frac{J_0}{m_{\mathcal S}^2}>\mu_{H,0}^2.
\end{equation}
\end{proposition}

\begin{proof}
\emph{Step 1 (Portal and source in the quadratic regime).}
On the receiver brane after the flip, take the quadratic Lagrangian density for the auxiliary portal field
$\mathcal S$ and the Higgs doublet $H$ (suppressing quartics of $\mathcal S$, which only renormalize
$m_{\mathcal S}$ in this argument):
\[
\mathcal L\;=\;-\tfrac12\,\mathcal S\,K_{\mathcal S}\,\mathcal S\ -\ \lambda_p\,\mathcal S\,H^\dagger H\ +\ J_{\rm LoG}\,\mathcal S
\ -\ |D H|^2\ -\ \mu_{H,0}^2\,H^\dagger H\ -\ \lambda_H\,(H^\dagger H)^2.
\]
Varying w.r.t.\ $\mathcal S$ gives the linear sourced equation
\begin{equation}
\label{eq:S-EOM}
K_{\mathcal S}\,\mathcal S\ =\ J_{\rm LoG}\ -\ \lambda_p\,H^\dagger H\ .
\end{equation}
We work to leading order in the portal (standard in deriving a \emph{mass shift}): set $H^\dagger H$
to its background value when solving \eqref{eq:S-EOM}, and account for the backreaction on $H$ via the
effective mass term. This is the usual tree–level “integrate out $\mathcal S$” step.

\emph{Step 2 (Retarded solution and DC projector).}
Let $G_R$ be the retarded Green kernel for $K_{\mathcal S}$, so $K_{\mathcal S}G_R=\delta^{(4)}$ and
$\mathrm{supp}\,G_R\subset\{t>0\}$. The (unique causal) solution of \eqref{eq:S-EOM} is
\[
\mathcal S\ =\ G_R * J_{\rm LoG}\ -\ \lambda_p\,G_R * (H^\dagger H),
\]
where $*$ denotes spacetime convolution. Define the spacetime DC projector
\[
(\mathsf P_0 f)\ :=\ \lim_{T\to\infty}\frac{1}{2T}\int_{t_0-T}^{t_0+T}\!\!dt\;\frac{1}{\mathrm{Vol}(S^3)}\int_{S^3}\!d\Omega\;f(t,\Omega),
\]
whenever the limit exists. By (A3) there are no $k>0$ drives after $t_{\rm flip}$, and the SK positivity
($\Im W\ge 0$) implies all $k>0$ transients of the linear dynamics decay. Hence, for $t\to\infty$ the
only potentially nonvanishing component of $\mathcal S$ is its DC part, $\mathsf P_0\mathcal S$.

\emph{Step 3 (Static susceptibility and the $1/m_{\mathcal S}^2$ factor).}
By (A1), $K_{\mathcal S}=-(\Box)+m_{\mathcal S}^2$ with $m_{\mathcal S}>0$. Therefore, in Fourier space,
$\widetilde G_R(\omega,\mathbf k)=\big[(\omega+i0^+)^2-\mathbf k^2-m_{\mathcal S}^2\big]^{-1}$, so the
\emph{static} ($\omega=0,\mathbf k=0$) susceptibility is
\[
\chi_0\ :=\ \widetilde G_R(0,\mathbf 0)\ =\ \frac{1}{m_{\mathcal S}^2}\ .
\]
Let $J_0:=\mathsf P_0 J_{\rm LoG}$. Since $G_R$ is time–translation invariant and causal,
$\mathsf P_0(G_R*J_{\rm LoG})=\chi_0\,J_0=J_0/m_{\mathcal S}^2$. Moreover, the term
$G_R*(H^\dagger H)$ contributes only through its DC part. At the level of \emph{mass} extraction we
\emph{define} the Higgs mass shift by the coefficient of $H^\dagger H$ in the effective action,
so the backreaction piece does not enter the DC value of $\mathcal S$ that multiplies $H^\dagger H$;
equivalently, to linear order in $\lambda_p$ one may evaluate $\mathcal S$ on the background with
$H^\dagger H$ replaced by its constant late–time value, which only renormalizes $\mu_{H,0}^2$ and does
not affect the \emph{history–dependent} contribution from $J_{\rm LoG}$.

Thus,
\begin{equation}
\label{eq:Sbar}
\mathsf P_0\mathcal S\;=\;\frac{J_0}{m_{\mathcal S}^2}\,,
\qquad
\mathcal S(t,\mathbf x)\;=\;\frac{J_0}{m_{\mathcal S}^2}\ +\ O\!\big(e^{-m_{\rm gap}(t-t_{\rm flip})}\big),
\end{equation}
where the exponential decay of transients follows from (A3) (gap $m_{\rm gap}\ge m_{\mathcal S}$ and
$\Im W\ge 0$).

\emph{Step 4 (Extracting the Higgs mass shift).}
The Higgs quadratic term in the Lagrangian reads
\[
-\ \mu_{H,0}^2\,H^\dagger H\ -\ \lambda_p\,\mathcal S\,H^\dagger H\ =\ -\big(\mu_{H,0}^2+\Delta\mu_H^2\big)\,H^\dagger H
\]
with the \emph{definition} $\Delta\mu_H^2:=-\lambda_p\,\mathsf P_0\mathcal S$ of the (late–time) mass
shift. Using \eqref{eq:Sbar} gives exactly
\[
\Delta\mu_H^2\;=\;-\lambda_p\,\frac{J_0}{m_{\mathcal S}^2},
\]
which is constant in time and space for $t>t_{\rm flip}$.

\emph{Step 5 (VEV and the EWSB condition).}
The renormalized Higgs potential at the flip scale is
$V(H)=\lambda_H(H^\dagger H)^2+\big(\mu_{H,0}^2+\Delta\mu_H^2\big)H^\dagger H+\cdots$.
Its minimum occurs at
\[
v^2\;=\;\frac{-\,\mu_{H,0}^2-\Delta\mu_H^2}{\lambda_H}
\;=\;\frac{1}{\lambda_H}\left(\lambda_p\,\frac{J_0}{m_{\mathcal S}^2}-\mu_{H,0}^2\right),
\]
and $v\neq 0$ iff $\lambda_p\,J_0/m_{\mathcal S}^2>\mu_{H,0}^2$, as claimed.

\emph{Remarks on robustness.} Any small local self–interaction of $\mathcal S$ or finite SK
dissipation only renormalizes $m_{\mathcal S}^2\mapsto m_{\mathcal S,{\rm eff}}^2>0$ and leaves the
structure of the argument intact: the DC response is the static susceptibility times $J_0$, while all
non–DC pieces are damped by (A3). Thus \eqref{eq:Delta-muH-const} is the universal tree–level late–time
result consistent with causality and SK positivity.
\end{proof}


\begin{corollary}[GW realization and radion mixing]
\label{cor:GW-realization}
Let a bulk scalar $\Phi$ stabilize the inter–brane separation $d$ (GW). Writing
$\Phi=\Phi_0(y;d_0)+\varphi$ and $d=d_0+\delta d$, the quadratic reduction on the receiver brane
produces
\begin{equation}
\label{eq:GW-reduction}
\mathcal L_{\rm eff}\;=\;\frac12(\partial \delta d)^2-\frac12 m_{\rm rad}^2(\delta d)^2
+\frac12(\partial \mathcal S)^2-\frac12 m_{\mathcal S}^2\mathcal S^2
+\mu\,\delta d\,\mathcal S\ -\ \lambda_p\,\mathcal S\,H^\dagger H\ +\ J_{\rm LoG}\,\mathcal S.
\end{equation}
If the radion is heavy ($m_{\rm rad}^2\gg \mu^2$), integrating out $\delta d$ shifts
$m_{\mathcal S}^2\to m_{\mathcal S}^2-\mu^2/m_{\rm rad}^2$ and suppresses radion–Higgs mixing by
$\sim\mu/m_{\rm rad}$. At the KR head, $\delta d$ is \emph{not}
an allowed variation (fixed boundary), so the $\delta d$ sector disappears and \eqref{eq:S-portal-loc}
remains unchanged. Proposition~\ref{prop:persistence-freeze} then applies.
\end{corollary}

\paragraph{Fisher orthogonality at the AdS$_4$ fixed point.}\label{par:uom-main-117}
At the maximally symmetric AdS$_4$ fixed point, the tangent space decomposes (Sec.~\ref{subsec:P3-master})
into RP–odd, Weyl (trace), conformal–mismatch, and physical kernels, orthogonal in the Fisher metric.
The \emph{brane-bending/Weyl} direction lies in the Weyl tangent, whereas the stabilized radion (RS1 interior)
and the KR boundary-distance scalar lie in the \emph{physical} kernel. Hence, at quadratic order and at the
fixed point, their mixing matrix element vanishes; any residual mixing comes from higher order and/or mild
symmetry breaking (guarded by $m_{\rm rad}$ and the near--AdS symmetry of the receiver).

\subsubsection*{\texorpdfstring{Fisher orthogonality at the AdS$_4$ fixed point (and cutoff–controlled near–orthogonality)}{Fisher orthogonality at the AdS4 fixed point (and cutoff-controlled near-orthogonality)}}
\label{subsubsec:fisher-ortho}

\paragraph{Setup and objects.}\label{par:uom-main-118}
Work at the holographically renormalized AdS$_4$ receiver fixed point $\rho_\ast$ in the Gaussian,
zero–mean gauge. Let $C_\ast$ be its covariance and equip the tangent space with the Fisher inner
product
\(
\langle\!\langle \delta C,\delta C'\rangle\!\rangle_{C_\ast}
=\frac12\,\mathrm{tr}\!\big(C_\ast^{-1}\delta C\,C_\ast^{-1}\delta C'\big)
\)
as in Lemma~\ref{lem:ReW-Fisher}. Denote by
\(
\Pi_{\rm W},\Pi_{\rm RP},\Pi_{\rm conf}
\)
the orthogonal projectors onto the Weyl (Riegert), RP–odd, and conformal–mismatch tangents
(Thm.~\ref{thm:completeness}), and by
\(
\Pi_{\rm phys}:=\mathbf 1-\Pi_{\rm W}-\Pi_{\rm RP}-\Pi_{\rm conf}
\)
the projector onto the \emph{physical kernel}.

\paragraph{Radion tangent and brane–bending mode.}\label{par:uom-main-119}
Let $\delta C_{\rm rad}$ denote the canonically normalized tangent induced by a small fluctuation
of the inter–brane distance (the stabilized \emph{radion} of the relevant interior interval).
Let $\delta C_{\rm bb}$ be the \emph{brane–bending} (rigid displacement of the outer receiver
wall). In a maximally symmetric AdS$_4$ background, $\delta C_{\rm bb}$ sits entirely in the
Weyl+RP blocks whereas $\delta C_{\rm rad}$ is a \emph{physical} scalar excitation (no local
Weyl content, RP–even): to quadratic order,
\[
\Pi_{\rm W}\delta C_{\rm rad}=0,\qquad \Pi_{\rm RP}\delta C_{\rm rad}=0,\qquad
\Pi_{\rm conf}\delta C_{\rm rad}=0.
\]
Thus $\delta C_{\rm rad}\in \mathrm{ran}\,\Pi_{\rm phys}$ and
$\langle\!\langle \delta C_{\rm rad},\delta C_{\rm bb}\rangle\!\rangle_{C_\ast}=0$.

\paragraph{Near–orthogonality and the angle $\vartheta$.}\label{par:uom-main-120}
Away from exact maximal symmetry and at finite receiver cutoff $\Lambda$ (which breaks exact Weyl
invariance through the regulated intertwiner and the induced–gravity term), the radion tangent can
acquire a small Weyl component. Define the (Fisher) \emph{non–orthogonality angle}
\begin{equation}
\label{eq:angle-def}
\sin\vartheta(\Lambda)\;:=\;
\frac{\|\Pi_{\rm W}^{(\Lambda)}\,\delta C_{\rm rad}\|_{C_\ast}}
     {\|\delta C_{\rm rad}\|_{C_\ast}}\in[0,1],\qquad
\cos\vartheta(\Lambda)\;:=\;
\frac{\|\Pi_{\rm phys}^{(\Lambda)}\,\delta C_{\rm rad}\|_{C_\ast}}
     {\|\delta C_{\rm rad}\|_{C_\ast}}.
\end{equation}
Maximal symmetry gives $\vartheta( \Lambda)\!=\!0$. Small departures (finite–$\Lambda$ Weyl
breaking; mild time–asymmetry of the welded kernel across the tick) produce $\vartheta(\Lambda)
=\mathcal O(\varepsilon)$ with $\varepsilon\ll1$.

\paragraph{Closed spectral form and cutoff dependence.}\label{par:uom-main-121}
Let $\mathsf D_{\rm W}^{(\Lambda)}:\sigma\mapsto\delta C_{\rm W}(\sigma)$ be the linear
Weyl–variation map at the fixed point (its kernel is the Riegert/Paneitz Green’s operator dressed
by the induced–gravity screen). The Moore–Penrose projector reads
\(
\Pi_{\rm W}^{(\Lambda)}=
\mathsf D_{\rm W}^{(\Lambda)}\big[(\mathsf D_{\rm W}^{(\Lambda)})^\dagger\mathsf D_{\rm W}^{(\Lambda)}\big]^{-1}
(\mathsf D_{\rm W}^{(\Lambda)})^\dagger.
\)
Because the AdS$_4$ background is homogeneous and isotropic, both
$(\mathsf D_{\rm W}^{(\Lambda)})^\dagger\mathsf D_{\rm W}^{(\Lambda)}$ and the Fisher metric are
diagonal in the frequency/momentum basis $(\omega,k)$, with eigen–weights proportional to the
(holographically fixed) anomaly kernel and the induced–gravity screen
\(
\mathcal K(\omega,k;\Lambda)=\frac{q(\omega,k)}{q(\omega,k)+r(\Lambda)k^2}\,,
\quad q(\omega,k)\!\approx\!\sqrt{k^2+m_g^2},\ \ 
r(\Lambda)=\frac{M_4^2(\Lambda)}{2M_5^3}.
\)
Let $U_{\rm rad}(\omega,k)$ be the \emph{normalized} spectral profile of $\delta C_{\rm rad}$ in
this basis. Then one obtains the closed expression
\begin{equation}
\label{eq:sin-theta-spectral}
\sin^2\vartheta(\Lambda)\;=\;
\frac{\displaystyle
\int\!\frac{d\omega\,d^3k}{(2\pi)^4}\,
\frac{|U_{\rm rad}(\omega,k)|^2}{\|U_{\rm rad}\|_2^2}\;
\underbrace{\frac{\mathcal W(\omega,k;\Lambda)}
{\mathcal W(\omega,k;\Lambda)+\mathcal S(\omega,k;\Lambda)}}_{:=\ \Xi(\omega,k;\Lambda)}
}
{\displaystyle 1}
\quad\text{with}\quad
\left\lbrace\!
\begin{aligned}
\mathcal W(\omega,k;\Lambda)&\propto \gamma^2\,|\mathcal K(\omega,k;\Lambda)|^2,\\[2pt]
\mathcal S(\omega,k;\Lambda)&\propto \text{(physical scalar Fisher weight)}\,,
\end{aligned}\right.
\end{equation}
where $\gamma=\tfrac{aG_4}{12\pi}$ and the proportionality constants are fixed by the canonical
normalizations of the two tangent families. The factor
\(
\Xi(\omega,k;\Lambda)\in[0,1]
\)
is a monotone \emph{decreasing} function of $r(\Lambda)$ at fixed $(\omega,k)$ and satisfies the
low–momentum bound
\begin{equation}
\label{eq:theta-lowk-bound}
\Xi(\omega,k;\Lambda)\ \lesssim\ \frac{1}{1+\alpha\,r(\Lambda)\,k^2/q(\omega,k)}\quad
\Longrightarrow\quad
\sin\vartheta(\Lambda)\ \lesssim\ \frac{\varepsilon_{\rm geom}}{\sqrt{1+\alpha\,r(\Lambda)\,k_\ast^2/q_\ast}}\ ,
\end{equation}
for some $\alpha=\mathcal O(1)$ and a spectral centroid $(\omega_\ast,k_\ast)$ of $U_{\rm rad}$
(the geometric imperfection $\varepsilon_{\rm geom}\ll1$ comes from the time–asymmetry and from
any residual non–zonal components of the welded kernel). Using the induced–gravity scaling
\(
M_4^2(\Lambda)=M_5^3\ell+\kappa\,a\,\Lambda^2
\)
one has
\[
r(\Lambda)\;=\;\frac{\ell}{2}+\frac{\kappa\,a}{2M_5^3}\,\Lambda^2
\quad\Longrightarrow\quad
\sin\vartheta(\Lambda)\ \lesssim\ 
\frac{\varepsilon_{\rm geom}}{\sqrt{1+\beta\,\Lambda^2}}\quad
(\beta:=\alpha\,\tfrac{\kappa a}{2M_5^3}\,\tfrac{k_\ast^2}{q_\ast})\,.
\]
Thus \emph{increasing the receiver cutoff} suppresses the Weyl leakage of the radion tangent as
$1/\Lambda$ (or faster), restoring orthogonality.

\paragraph{Mixing to the DC trace and effective coupling.}\label{par:uom-main-122}
Let $T^\mu{}_\mu^{\rm (BY)}$ be the anomaly–rectified trace operator on the receiver. The
(radion)$\leftrightarrow$(DC trace) \emph{mixing amplitude} induced by non–orthogonality at
finite $\Lambda$ is
\begin{equation}
\label{eq:mixing-DC}
g_{{\rm rad}\,\text{-}\,{\rm DC}}(\Lambda)\;=\;\frac{\sin\vartheta(\Lambda)}{\Lambda_{\rm rad}},
\qquad
\Lambda_{\rm rad}\sim\sqrt{\frac{24\,M_5^3}{k}}\,e^{-kd_0},
\end{equation}
where $\Lambda_{\rm rad}$ is the usual radion coupling scale of the stabilized interval (Goldberger–Wise).
Combining \eqref{eq:theta-lowk-bound}–\eqref{eq:mixing-DC} gives the \emph{closed} cutoff dependence
\begin{equation}
\label{eq:mixing-closed}
g_{{\rm rad}\,\text{-}\,{\rm DC}}(\Lambda)\ \lesssim\ 
\frac{\varepsilon_{\rm geom}}{\Lambda_{\rm rad}}\,
\frac{1}{\sqrt{1+\beta\,\Lambda^2}}\ ,
\end{equation}
which shows: (i) exact orthogonality at the AdS$_4$ fixed point ($\varepsilon_{\rm geom}\!=\!0$),
(ii) parametric suppression $g\propto \Lambda^{-1}$ for large cutoff even when
$\varepsilon_{\rm geom}\neq 0$, and (iii) decoupling for large $\Lambda_{\rm rad}$ (stiff GW).
Equation \eqref{eq:mixing-closed} is the quantitative statement behind “radion is
\emph{near}–orthogonal to the brane–bending/Weyl DC mode” used in the portal analysis.

\paragraph{Implication for the Higgs portal (interior RS1/KR continuity).}\label{par:uom-main-123}
The brane–localized auxiliary scalar $\mathcal S$ mediates the Higgs portal and is sourced by the welded LoG/anomaly trace. In the long–wavelength regime one has
\[
\Delta\mu_H^2(x)\;=\;-\frac{\lambda_p}{m_{\mathcal S}^2}\,J_{\rm LoG}(x)\,,
\]
with $J_{\rm LoG}$ the (regulated) LoG–projected trace response. Across an interior RS1 segment abutting the KR head
(not an active receiver) the \emph{light} scalar \emph{persists}
(the RS1 radion flows to the KR boundary--distance scalar with $m_r^2>0$) but remains in the \emph{physical} tangent; its visible coupling is suppressed by the induced--gravity Robin boundary condition as $\varepsilon_{\rm BC}\ll 1$ and only mixes with the Weyl/trace tangent by a tiny angle
$\theta(\Lambda)=O(\varepsilon_{\rm BC}\Lambda^{-2})$. Therefore the portal \emph{itself} is unaffected at leading order: $J_{\rm LoG}$ depends only on the trace kernel and the Riegert coefficient, both smooth across the transition. The receiver head's cutoff is \emph{Wilsonian}: ($\Lambda=z_0^{-1}$), so $J_{\rm LoG}^{\rm (DC)}(\Lambda)$ follows the Hamilton--Jacobi/RG flow rather than freezing; on the KR head it evolves modestly according to the available BY power and the induced--gravity screen (Sec.~\ref{subsubsec:KR-cutoff-RG}), then \emph{asymptotes} to a constant when the source decouples. Thus the portal strength evolves \emph{smoothly} to a finite asymptote:
\begin{equation}
\label{eq:dmuh-DC}
\Delta\mu_H^2(t,\mathbf x)\;=\;\Delta\mu_{H,\infty}^2\ +\ \text{(ring–down of non–DC modes)},
\qquad
\Delta\mu_{H,\infty}^2=\;-\frac{\lambda_p}{m_{\mathcal S}^2}\;J_{\rm LoG}^{\rm (DC)}\!\big(\Lambda_\infty\big),
\end{equation}
where $\Lambda_\infty$ is the terminal Wilsonian scale fixed by the BY power budget (not by an \emph{ad hoc} code ceiling). The small Higgs–radion mixing (proportional to $c_r\propto \varepsilon_{\rm BC}$) induces only $\delta\kappa_h\ll 10^{-6}$, well below current sensitivities (Sec.~\ref{subsubsec:fifthforce-robin}).

\subsubsection*{Persistent scalar and cutoff dynamics on a KR head}
\label{subsubsec:KR-cutoff-RG}

\paragraph{Statement.}\label{par:uom-main-124}
Let the receiver AdS$_4$ wall be a Karch--Randall head (EOW cap) of the active AdS$_5$ interval. Then:
\begin{enumerate}[label=(K\arabic*), leftmargin=2em]
\item \textbf{Persistent light scalar in the physical tangent.} If an interior RS1 segment abuts the KR head,
the stabilized radion there continues as the KR boundary–distance scalar ($E_0=4$) with
$m_r^2=4k^2+\delta m_{\rm GW}^2>0$ and remains in the \emph{physical} covariance tangent. The active receiver
remains the KR head throughout. Its visible overlap is suppressed by the induced–gravity Robin boundary
condition on the visible head:
\(
\varepsilon_{\rm BC}\simeq [1+s_{\rm vis}(\Lambda)\,k(\Lambda)\,L(\Lambda)]^{-1}\ll 1
\),
and its Weyl mixing at the welded band edge scales as
$\theta(\Lambda)=O(\varepsilon_{\rm BC}\Lambda^{-2})$. Hence the conformality–mismatch penalty carried by this mode behaves as
\(
\Phi_{\rm conf}\sim O(\varepsilon_{\rm BC}^2\Lambda^{-4})
\)
(tiny and \emph{decreasing} with $\Lambda$). The RP block remains zero at quadratic order by OS positivity.

\item \textbf{Wilsonian cutoff dynamics (forced motion).} The receiver cutoff is a state variable
whose evolution is fixed by the SK/HJ wall balance
\(
\mu_{\rm RG}(\Lambda)\,\dot\Lambda=(\mathcal P_{\rm in}-\mathcal P_{\rm out})-\dot{\mathcal F}_{\rm rec}
\).
Any accepted tick with $\Delta I_k>0$ forces $\dot\Lambda>0$ on a subinterval. The KKT condition
constrains \emph{controls} (e.g.\ window shape and passband) but does not allow a static cutoff while
payload is admitted. If the source exhausts before a flip, the wall work goes to zero and $\Lambda$
asymptotes to a terminal value $\Lambda_\infty$ (post--exhaustion boundary), not a dwell phase.

\item \textbf{Flip condition (conditional).} The BY/anomaly floor is a positive variance injection even for
zero--area (LoG) pulses. A flip occurs only if the cumulative anomaly increments exceed the initial deficit
and the export policy supplies sufficient total payload, i.e.\ when the first--crossing condition
\eqref{eq:flip-condition-updated} (or \eqref{eq:flip-condition-ref}) is met. If the split-entry work or
available wall--work fails to clear the beat--the--source condition (Sec.~\ref{subsec:ads5-beat-src}), the
system enters the terminal/no--flip branch with a terminal boundary. KR cutoff dynamics affects the \emph{rate}
of approach, not the existence of crossing under the stated conditions.
\end{enumerate}

\paragraph{Why (sketch).}\label{par:uom-main-125}
(i) The KR boundary–distance scalar is the normalizable $E_0=4$ AdS$_4$ mode. In the induced–gravity/Robin regime its bulk wavefunction has a small visible boundary value $u_r(z_{\rm vis})\propto [1+s_{\rm vis}kL]^{-1}$; hence its Fisher projection onto the conformality–mismatch block vanishes at leading order and the residual mixing with Weyl at the welded band edge scales as $\theta(\Lambda)=O(\varepsilon_{\rm BC}\Lambda^{-2})$, giving a $O(\varepsilon_{\rm BC}^2\Lambda^{-4})$ penalty.

(ii) The receiver cost $\Phi_{\rm rec,k}(\Lambda)$ \emph{decreases} with $\Lambda$ through the induced–gravity screen in the trace kernel and the $O(\varepsilon_{\rm BC}^2\Lambda^{-4})$ conformality term. But $\Lambda$ is \emph{physical} (Wilsonian): changing it costs RG work $W_{\rm RG}$ and modifies the admitted payload $\Delta I_k(\Lambda)$. The KKT balance constrains admissible controls but does not permit a static cutoff while payload is admitted; the terminal value $\Lambda_\infty$ appears only after exhaustion (Sec.~\ref{subsec:alternating-heads}). No energy is borrowed ``from nowhere'', and there is no \emph{ad hoc} jump to a code ceiling.

(iii) The positive anomaly increments $\Theta_k$ depend on the welded amplitude and the RG–screened trace kernel. Their sum is independent of whether the cutoff is drifting or halted and provides the flip criterion.

\paragraph{Cutoff dynamics (summary).}\label{par:uom-main-126}
\begin{itemize}[leftmargin=1.4em]
\item \textbf{Smaller $\bar\Lambda_k$ (larger $K(\bar\Lambda_k)$):} small but strictly positive refinements $\Delta\Lambda_k>0$ under knee/guard constraints.
\item \textbf{Larger $\bar\Lambda_k$ (smaller $K(\bar\Lambda_k)$):} $\Delta\Lambda_k$ may shrink as inflow decreases; welded widths remain locked at $(\alpha_t^\ast,\alpha_x^\ast)$ evaluated at $\bar\Lambda_k$.
\item \textbf{Terminal boundary:} only after exhaustion ($\Delta I_k\to0$) does $\dot\Lambda\to0$ and a terminal value $\Lambda_\infty$ emerge.
\end{itemize}


\subsubsection*{Freeze after the flip and constant Higgs mass}
\label{subsubsec:freeze-const-H}

\paragraph{Setup and notation.}\label{par:uom-main-127}
Let $t_{\rm flip}$ denote the time of the last accepted tick on the receiver. Up to $t_{\rm flip}$ the welded (trace–channel) source is
\[
\tau_{\rm hist}(x)\;=\;\sum_{k\le K}\tau_k(x),\qquad \mathrm{supp}\,\tau_k\subset I_k,\qquad t_K=t_{\rm flip},
\]
with zero–DC (ideal LoG/DoG) $\int d^4x\,\tau_k(x)=0$, but with a strictly positive \emph{anomaly rectifier} per tick,
\[
\Theta_k\ :=\ \gamma\,\Delta\!\langle R^2\rangle_k\;>\;0,\qquad 
\gamma=\frac{a\,G_4}{12\pi}>0,
\]
entering the DC trace balance. Define the spacetime $k{=}0$ (DC) projector by
\[
(\mathsf P_0 f)\ :=\ \lim_{T\to\infty}\frac{1}{T}\int_{t_{\rm flip}}^{t_{\rm flip}+T}\!\!\!\!dt\;\frac{1}{\mathrm{Vol}(S^3)}\int_{S^3}\!d\Omega\; f(t,\Omega),
\]
whenever the limit exists.

\paragraph{Receiver response and the Higgs--mass portal.}\label{par:uom-main-128}
On the KR head there is no \emph{RS1 radion}; the relevant light scalar is the KR boundary-distance mode
(physical kernel) with Robin-suppressed visible coupling. The low-frequency scalar channel that feeds the
Higgs mass is the anomaly/trace sector via the auxiliary portal field $\mathcal S$ (Sec.~\ref{subsec:higgs-portal}).
In linearized Gaussian response,
\begin{align}
\Delta \mu_H^2(x)
&=\ \lambda_p\,\big(\mathcal K_{\mathcal S} * J_{\rm LoG}\big)(x)\ +\ \lambda_p\,\gamma\,\big(\mathcal K_{\rm anom} * \mathfrak R\big)(x),
\label{eq:dmuh-master}
\\[-2pt]
&\qquad J_{\rm LoG}=\text{LoG–projection of }\tau_{\rm hist},\quad 
\mathfrak R(x):= \big(R\cdot R\big)_{\rm hist}(x),\quad
(*)=\text{retarded convolution},
\nonumber
\end{align}
where $\mathcal K_{\mathcal S}$ is the retarded kernel for the (heavy) boundary fluctuation $\mathcal S$ and $\mathcal K_{\rm anom}$ is the trace/compensator kernel on AdS$_4$; both are gapped and causal.

\begin{proposition}[Post–flip constancy of the Higgs mass]\label{prop:const-Higgs-postflip}
Assume: 
\begin{enumerate}[label=(H\arabic*), leftmargin=2.1em, itemsep=2pt]
\item \emph{No post–flip drive:} the welded trace source vanishes for $t>t_{\rm flip}$, i.e.\ $\tau(x)=0$ when $t>t_{\rm flip}$;
\item \emph{Gapped, dissipative response:} the retarded response operators $\mathcal K_{\mathcal S}$ and $\mathcal K_{\rm anom}$ extend to bounded convolution operators on $L^2(\mathbb R_t;H^{-s}(S^3))$ ($s$ large enough) whose Laplace transforms have no spectrum on $\Re z\ge 0$ except a simple eigenvalue at $z=0$ corresponding to the spacetime constant (DC) mode;
\item \emph{Finite history:} $\tau_{\rm hist}=\sum_{k\le K}\tau_k$ is compactly supported in time and belongs to $L^1\cap L^2$ with values in $H^{-s}(S^3)$.
\end{enumerate}
Then the portal–induced Higgs–mass deformation 
\[
\Delta \mu_H^2(t,\Omega)\;=\;\lambda_p\big(\mathcal K_{\mathcal S} * J_{\rm LoG}\big)(t,\Omega)
+\lambda_p\,\gamma\,\big(\mathcal K_{\rm anom} * \mathfrak R\big)(t,\Omega),\qquad t>t_{\rm flip},
\]
with $J_{\rm LoG}$ the LoG–projection of $\tau_{\rm hist}$ and $\mathfrak R=(R\cdot R)_{\rm hist}$, decomposes as
\[
\Delta \mu_H^2(t,\Omega)\;=\;\Delta\mu_{H,\infty}^2\;+\;\delta\mu_H^2(t,\Omega),
\]
where $\Delta\mu_{H,\infty}^2$ is a spacetime constant determined solely by the pre–flip history and $\|\delta\mu_H^2(t,\cdot)\|_{H^{-s}(S^3)}\le C\,e^{-m_{\rm gap}(t-t_{\rm flip})}$ for some $m_{\rm gap}>0$. In particular, the electroweak VEV is time–independent for $t>t_{\rm flip}$ up to ordinary RG running.
\end{proposition}

\begin{proof}
\emph{Step 1 (Functional setting).} Let $\mathcal H:=H^{-s}(S^3)$ with $s$ large so that all kernels act boundedly on $\mathcal H$. For a distributional source $f\in L^1\cap L^2(\mathbb R_t;\mathcal H)$ supported on $(-\infty,t_{\rm flip}]$, define the retarded response
\[
(\mathcal K * f)(t)\;:=\;\int_{-\infty}^t\! \mathcal K(t-t')\,f(t')\,dt',
\]
where $\mathcal K(t)$ is a strongly continuous family of bounded operators on $\mathcal H$ vanishing for $t<0$.

\emph{Step 2 (Spectral projector onto the DC mode).} By (H2) there is a rank–one spectral projector $\mathsf P_0$ on the spacetime constant mode and a complementary projector $\mathsf Q_0=\mathbf 1-\mathsf P_0$ such that the Laplace transform $\,\widetilde{\mathcal K}(z)=\int_0^\infty e^{-zt}\mathcal K(t)\,dt$ is analytic on $\Re z>-m_{\rm gap}$ and admits the partial fraction decomposition
\[
\widetilde{\mathcal K}(z)\;=\;\frac{\mathsf P_0}{z}\ +\ \widetilde{\mathcal R}(z),\qquad \|\widetilde{\mathcal R}(z)\|\le C\ \text{on }\Re z\ge -m_{\rm gap}/2.
\]
By inverse Laplace transform, this yields the semigroup splitting
\begin{equation}
\mathcal K(t)\;=\;\mathsf P_0\ +\ \mathcal R(t),\qquad \|\mathcal R(t)\|\ \le\ C\,e^{-m_{\rm gap} t},\quad t\ge 0.
\label{eq:K-splitting}
\end{equation}

\emph{Step 3 (Decomposition of the response).} Let $f_1:=J_{\rm LoG}$ and $f_2:=\mathfrak R$, both compactly supported on $(-\infty,t_{\rm flip}]$ by (H3). Using \eqref{eq:K-splitting} and $t>t_{\rm flip}$,
\[
(\mathcal K * f_i)(t)\;=\;\int_{-\infty}^{t_{\rm flip}}\!\!\left[\mathsf P_0+\mathcal R(t-t')\right] f_i(t')\,dt'
\;=\;\mathsf P_0\!\left(\int_{-\infty}^{t_{\rm flip}}\! f_i(t')\,dt'\right)\;+\;\int_{-\infty}^{t_{\rm flip}}\!\mathcal R(t-t') f_i(t')\,dt'.
\]
Define the \emph{history DC} functional 
\(
\mathcal H_i:=\mathsf P_0\!\left(\int_{-\infty}^{t_{\rm flip}} f_i(t')\,dt'\right)
\)
which is a constant element of $\mathcal H$, and the transient
\(
\delta_i(t):=\int_{-\infty}^{t_{\rm flip}}\mathcal R(t-t') f_i(t')\,dt'.
\)
By Young’s inequality and the exponential bound on $\mathcal R$, for $t>t_{\rm flip}$,
\[
\|\delta_i(t)\|_{\mathcal H}
\ \le\ \int_{-\infty}^{t_{\rm flip}}\! \|\mathcal R(t-t')\|\,\|f_i(t')\|_{\mathcal H} dt'
\ \le\ C\,e^{-m_{\rm gap}(t-t_{\rm flip})}\,\|f_i\|_{L^1(\mathbb R;\mathcal H)}.
\]
Hence 
\[
(\mathcal K * f_i)(t)\;=\;\mathcal H_i\;+\;O\!\left(e^{-m_{\rm gap}(t-t_{\rm flip})}\right).
\]

\emph{Step 4 (Apply to the Higgs–mass shift).} With $\Delta\mu_H^2=\lambda_p(\mathcal K_{\mathcal S}*f_1)+\lambda_p\gamma(\mathcal K_{\rm anom}*f_2)$ and the above bounds applied to both kernels (sharing the same $m_{\rm gap}$ up to constants), we obtain
\[
\Delta\mu_H^2(t,\Omega)\;=\;\underbrace{\lambda_p\,\mathcal H_1\;+\;\lambda_p\,\gamma\,\mathcal H_2}_{=\:\Delta\mu_{H,\infty}^2\ \text{(constant)}}\;+\;O\!\left(e^{-m_{\rm gap}(t-t_{\rm flip})}\right),
\]
as claimed. The constant depends only on the pre–flip integrals of $f_{1,2}$, i.e.\ on the tick history, and is independent of $t$ and $\Omega$. Finally, since $\mu_{H,{\rm eff}}^2=\mu_{H,{\rm bare}}^2+\Delta\mu_H^2$ enters the tree–level potential $V(H)=\mu_{H,{\rm eff}}^2 H^\dagger H + \lambda_H (H^\dagger H)^2$, the vacuum $v^2=-\mu_{H,{\rm eff}}^2/\lambda_H$ is time–independent for $t>t_{\rm flip}$ up to ordinary RG running of couplings between the flip scale and the weak scale. 
\end{proof}

\paragraph{Anomaly rectification and the DC bias.}\label{par:uom-main-129}
In the pre--flip KR head the cutoff $\Lambda$ advances monotonically; welded kernels are evaluated at
$\bar\Lambda_k$, and per tick,
\[
\Theta_k\ =\ \gamma\,(8\pi G_4)^2\,\big\langle|\mathcal K|^2\big\rangle_{\rm KR}\,C_{\rm LoG}\,
\frac{A_k^2}{\sigma_{t,k}\,\sigma_{x,k}^3}\ >\ 0,
\]
so the DC curvature obeys
\[
\langle R\rangle_{k+1}\ =\ \langle R\rangle_k\ +\ \Theta_k,
\]
and the effective $\Lambda_4=\tfrac14\langle R\rangle$ crosses $0$ \emph{under the first--crossing condition},
with a possible asymptotic approach if the payload/split-entry budget is insufficient.
The same history imprints the constant \eqref{eq:dmuh-DC} that fixes the Higgs mass after the flip.

\paragraph{What remains dynamical after the flip.}\label{par:uom-main-130}
Only standard SM renormalization–group running between the flip scale and the weak scale modifies parameters; the portal–induced $\Delta\mu_{H,\infty}^2$ and the DC curvature $\langle R\rangle_{\infty}$ are constants of motion of the post–flip linear response (no drive, gapped kernel). All $k>0$ spatial structure rings down exponentially fast.

\medskip\noindent
\emph{Conclusion.} Once the last pre–flip tick is accepted, the receiver enters a drive–free, gapped evolution in which the only surviving imprint on the Higgs mass is the DC component determined by the past anomaly–LoG history; this freezes the Higgs VEV (modulo ordinary RG running).



\subsubsection*{Radiative stability and the hierarchy}
\label{subsubsec:radiative-stability}
The portal we use is \emph{super–renormalizable} in 4D,
\(
\mathcal L\supset -\,\lambda_p\,\mathcal S\,H^\dagger H
\)
with $[\lambda_p]=1$. This has two important consequences:
(i) the \emph{tree–level} matching across the $\mathcal S$ threshold entirely controls the sign and
magnitude of the Higgs mass parameter after the flip; (ii) quantum corrections from $\mathcal S$ \emph{decouple}
below $m_{\mathcal S}$ and can be made parametrically small without spoiling the tree–level imprint.
We collect the needed statements and bounds.

\medskip
\noindent\textbf{(RS1) Tree–level matching (flip scale).}
Write the quadratic $\mathcal S$ sector as
\(
\frac12\mathcal S\,K_{\mathcal S}\,\mathcal S+\mathcal S\,(J_{\rm LoG}-\lambda_p H^\dagger H)
\)
with $K_{\mathcal S}=\Box+m_{\mathcal S}^2$.
Solving the $\mathcal S$ equation of motion to leading order in derivatives gives
\(
\langle\mathcal S\rangle = K_{\mathcal S}^{-1}\!(J_{\rm LoG}-\lambda_p H^\dagger H)
\approx m_{\mathcal S}^{-2}(J_{\rm LoG}-\lambda_p H^\dagger H)+\mathcal O(\partial^2/m_{\mathcal S}^4).
\)
Substituting back yields the local effective potential at (and below) the flip scale:
\begin{equation}
\label{eq:Veff-matching}
V_{\rm eff}(H)\;=\;\lambda_H(H^\dagger H)^2+\mu_{H,0}^2\,H^\dagger H
\;-\;\underbrace{\frac{\lambda_p}{m_{\mathcal S}^2}\,J_{\rm LoG}}_{\Delta\mu_H^2\text{ (tree)}}\,H^\dagger H
\;+\;\underbrace{\frac{\lambda_p^2}{2m_{\mathcal S}^2}}_{\Delta\lambda_H\text{ (tree)}}(H^\dagger H)^2\;+\cdots.
\end{equation}
Thus, the \emph{entire} tree–level shift of the Higgs mass parameter is the DC imprint
\(
\Delta\mu_H^2= -\,\lambda_p J_0/m_{\mathcal S}^2
\)
(Prop.~\ref{prop:persistence-freeze}), while the leading quartic correction is positive,
$\Delta\lambda_H=+\lambda_p^2/(2m_{\mathcal S}^2)$. Derivative corrections are suppressed by
$E^2/m_{\mathcal S}^2$ and are irrelevant for the vacuum analysis.

\medskip
\noindent\textbf{(RS2) One–loop threshold from $\mathcal S$ (Wilsonian view).}
At one loop, the $\mathcal S$ bubble with two portal insertions renormalizes the Higgs two–point:
\[
\delta\mu_H^2\Big|_{\rm 1\ell}\;=\;\lambda_p^2\,I(m_{\mathcal S},\mu_R),\qquad
I(m_{\mathcal S},\mu_R)
=\frac{1}{16\pi^2}\Big(\Lambda^2-m_{\mathcal S}^2\ln\frac{\Lambda^2}{\mu_R^2}+\cdots\Big),
\]
in a hard–cutoff picture, or equivalently (in $\overline{\rm MS}$) as a finite \emph{threshold} at $\mu_R\simeq m_{\mathcal S}$,
\begin{equation}
\label{eq:threshold-MS}
\delta\mu_H^2\Big|_{\rm 1\ell,\,match}\;=\;\frac{\lambda_p^2}{16\pi^2}\,c_0\,m_{\mathcal S}^2,
\qquad c_0=\mathcal O(1)\ \text{(scheme–dependent)}.
\end{equation}
Below $m_{\mathcal S}$ the field $\mathcal S$ is integrated out and \emph{does not run} the Higgs mass; the
$\mathcal S$–induced quartic $\Delta\lambda_H$ in \eqref{eq:Veff-matching} runs only with SM $\beta$–functions.
Hence, the \emph{only} $\mathcal S$–induced sensitivity is the finite match \eqref{eq:threshold-MS}, which we may absorb
into the renormalized $\mu_{H,0}^2$ at the scale $m_{\mathcal S}$. Requiring that no fine tuning worse than a factor
$\Delta_{\rm nat}$ is needed gives the \emph{naturalness} bound
\begin{equation}
\label{eq:nat-bound}
\frac{\lambda_p^2}{16\pi^2}\,c_0\,m_{\mathcal S}^2\ \lesssim\ \Delta_{\rm nat}\,v^2
\quad\Longrightarrow\quad
\lambda_p\ \lesssim\ 4\pi\,\sqrt{\frac{\Delta_{\rm nat}}{c_0}}\;\frac{v}{m_{\mathcal S}}\,.
\end{equation}
Numerically, taking $c_0\simeq 1$ and $\Delta_{\rm nat}=10$ (``10\% tuning''),
\[
m_{\mathcal S}=10^5~{\rm GeV}\ \Rightarrow\ \lambda_p\ \lesssim\ 0.1~{\rm GeV},\qquad
m_{\mathcal S}=10^4~{\rm GeV}\ \Rightarrow\ \lambda_p\ \lesssim\ 1~{\rm GeV}.
\]
These small (dimensionful) couplings are perfectly compatible with obtaining the correct $v$ from the \emph{tree} imprint
\eqref{eq:Veff-matching} by dialing $J_0$ (fixed by the welded anomaly history) accordingly.

\medskip
\noindent\textbf{(RS3) Portal hierarchy is technically natural.}
Because the portal is super–renormalizable, taking $\lambda_p \ll m_{\mathcal S}$ is \emph{technically natural}
(NDA: radiative corrections to $\lambda_p$ are at most $\delta\lambda_p\sim \lambda_p^3/(16\pi^2 m_{\mathcal S}^2)$).
Hence a small $\lambda_p$ needed to satisfy \eqref{eq:nat-bound} is stable under RG and threshold matching.

\medskip
\noindent\textbf{(RS4) Matching at the stopping cutoff (boundary--distance scalar suppressed).}
The only portal parameter that depends on the receiver cutoff is the DC source strength
$J_{\rm LoG}^{\rm (DC)}(\Lambda_{\rm stop})$, a monotone functional of the cutoff via the welded kernel.
Here $\Lambda_{\rm stop}:=\bar\Lambda_N$ at the flip (first--crossing tick) and
$\Lambda_{\rm stop}:=\Lambda_\infty$ in the exhaustion branch. It evolves with $\Lambda(t)$ during accepted
ticks and asymptotes only if the exhaustion branch is reached. The low--energy matching onto the Higgs sector is
\[
\mathcal L_{\rm eff}\supset -\,\lambda_p\,\frac{J_{\rm LoG}^{\rm (DC)}(\Lambda_{\rm stop})}{m_{\mathcal S}^2}\,H^\dagger H,
\]
with $m_{\mathcal S}$ the mass of the auxiliary (GW--boundary--fluctuation) scalar that mediates the anomaly--LoG portal.
The KR boundary--distance scalar remains physical but its visible coupling is Robin--suppressed, so no leading--order
radion mixing enters the matching.

\medskip
\noindent\textbf{(RS5) Decoupling and quadratic sensitivity.}
Integrating out $\mathcal S$ yields the dimension–six operator
\(
\frac{\lambda_p^2}{m_{\mathcal S}^2}(H^\dagger H)\,\mathcal O_{\rm LoG}
\)
at the \emph{matching} scale $\mu\simeq m_{\mathcal S}$. Consequently, the largest mass scale that can appear in quadratic corrections to $m_H^2$ from the portal is $m_{\mathcal S}$, not $\Lambda_{\rm stop}$:
\[
\delta m_H^2\ \sim\ \frac{\lambda_p^2}{16\pi^2}\,m_{\mathcal S}^2\ \times\ \big(\text{logs and small scheme factors}\big),
\]
plus the usual SM running between $m_{\mathcal S}$ and the weak scale. Thus choosing $m_{\mathcal S}\sim 10^{4}\text{--}10^{5}\,{\rm GeV}$ with moderate $\lambda_p$ keeps the Higgs mass technically natural, independently of how large $\Lambda_{\rm stop}$ is.

\medskip
\noindent\textbf{(RS6) Putting numbers (illustrative).}
Let $v=246\,{\rm GeV}$, $\lambda_H\simeq 0.13$, and take $\mu_{H,0}^2\ge 0$ (no tachyon at the bare level). Then the portal must supply
\[
\Delta\mu_{H,\infty}^2\;=\;\lambda_p\,\frac{J_{\rm LoG}^{\rm (DC)}(\Lambda_{\rm stop})}{m_{\mathcal S}^2}\;\simeq\;\lambda_H\,v^2\ \approx\ 7.9\times 10^{3}\,{\rm GeV}^2.
\]
For $m_{\mathcal S}=10^{5}\,{\rm GeV}$ and $\lambda_p\in[0.1,1]$ this requires
\[
J_{\rm LoG}^{\rm (DC)}(\Lambda_{\rm stop})\ \sim\ (8\times 10^{8}\text{--}8\times 10^{9})\,{\rm GeV}^4.
\]
In our welded formulae this DC is proportional to the anomaly rectifier times a positive shape factor that \emph{increases} with $\Lambda$:
\[
J_{\rm LoG}^{\rm (DC)}(\Lambda_{\rm stop})\ =\ C_{\rm anom}\,\gamma\;\big\langle |\mathcal K(\cdot;\Lambda_{\rm stop})|^2\big\rangle\;\mathcal S_{\rm LoG},
\]
where $\gamma=\ell^2/96$, $\langle|\mathcal K|^2\rangle$ is the KR–kernel average at $\Lambda_{\rm stop}$, and $\mathcal S_{\rm LoG}$ encodes the welded LoG shape scales fixed by the accepted history. Hence a larger stopping cutoff \(\Lambda_{\rm stop}\) relaxes the needed $\lambda_p$ or lowers the required $m_{\mathcal S}$; conversely, smaller \(\Lambda_{\rm stop}\) can be compensated by a somewhat lighter $\mathcal S$ so long as $m_{\mathcal S}\gg v$ to justify decoupling. In all cases, loop sensitivity is capped at $m_{\mathcal S}$ and does not probe $\Lambda_{\rm stop}$.

\medskip
\noindent\textbf{(RS7) Flip cutoff and thermal baryogenesis.}
Take a KR/flip cutoff in the range
\(
\Lambda_{\rm flip}=10^{11}\text{--}10^{13}\ {\rm GeV}.
\)
This allows purely thermal leptogenesis/baryogenesis ($T_B\gtrsim 10^9$\,GeV) while staying
well below the Planck scale. In the induced–gravity relation
\(
M_4^2(\Lambda)=M_5^3\ell+\kappa\,a\,\Lambda^2
\)
we may choose a moderate bulk curvature $k=\ell^{-1}\sim 10^{15}$\,GeV and
$M_5\sim 10^{17}$\,GeV with $\kappa a=\mathcal O(1)$, giving
\(
M_4\simeq 2.4\times 10^{18}{\rm\ GeV}
\)
to within small percent-level shifts across the quoted $\Lambda_{\rm flip}$ band.



\medskip
\noindent\textbf{(RS8) Summary.}
In our continuity picture the light scalar \emph{persists} in both RS1 and KR: it is the stabilized radion in an interior RS1 cell and the $E_0=4$ boundary–distance mode at a KR cap. In KR it remains in the \emph{physical} Fisher tangent and its \emph{visible} coupling is suppressed by the induced–gravity Robin boundary condition, so that the conformality–mismatch penalty is only
$O(\varepsilon_{\rm BC}^2\,\Lambda^{-4})$ at the welded band edge. The receiver cutoff \(\Lambda\) is \emph{Wilsonian}, not a code knob: in KR it changes modestly according to the Hamilton–Jacobi/RG balance (source BY power vs.\ RG work) and asymptotes to a finite \(\Lambda_\infty\) only after exhaustion. The Higgs portal continues to be dominated by the heavy brane mediator \(\mathcal S\): to leading order
\[
\Delta\mu_H^2\;=\;-\frac{\lambda_p}{m_{\mathcal S}^2}\;J_{\rm LoG}^{\rm (DC)}(\Lambda)\,,
\]
with \(J_{\rm LoG}^{\rm (DC)}(\Lambda)\) flowing smoothly in \(\Lambda\) and asymptoting to a constant at \(\Lambda_\infty\).
Radion–Higgs effects enter only through the BC–suppressed coupling \(c_r\propto \varepsilon_{\rm BC}\) and are constrained by long–range fifth–force bounds, \(c_r\lesssim (\hat c/\kappa_g)\sqrt{\alpha_{\rm lim}}\), which makes the induced Higgs–radion mixing \(\delta\kappa_h\simeq \tfrac12(c_r v/m_h^2)^2\ll 10^{-6}\).
Consequently, the Higgs mass is set by the (asymptotically constant) DC anomaly response and remains technically natural provided \(m_{\mathcal S}\gg v\) and the portal coupling is moderate; the persistent light scalar is compatible with this because its visible coupling is small and its conformality penalty vanishes along the refinement.


\subsubsection*{Operator hierarchy, alignment, and suppressed gluon couplings}
\label{subsubsec:selection-gluons}

\paragraph{Portal structure and operator basis.}\label{par:uom-main-131}
On the KR head (boundary--distance scalar present but visible coupling suppressed) the CP--even scalar
that talks to the Higgs is the auxiliary boundary field $\mathcal S$ sourced by the welded anomaly--LoG
history (Sec.~\ref{subsec:higgs-portal}). At scales $\mu\ll \Lambda_{\rm stop}$ the most general, CP--even,
gauge--invariant, dimension--$\le 5$ interactions relevant for the Higgs mass are
\begin{equation}
\mathcal L_{\rm eff}
\,=\,
-\lambda_p\,\mathcal S\,H^\dagger H
\;+\; \frac{\kappa_T}{\Lambda_T}\,\mathcal S\,T^\mu{}_\mu{}^{\rm (SM)}
\;+\; \sum_{a}\frac{c_{g}}{\Lambda_g}\,\mathcal S\,\mathrm{Tr}\,G_{\mu\nu}^a G^{a\,\mu\nu}
\;+\; \frac{c_{\gamma}}{\Lambda_\gamma}\,\mathcal S\,F_{\mu\nu}F^{\mu\nu}
\;+\;\cdots.
\label{eq:Seff-ops}
\end{equation}
Here $[\mathcal S]=1$ and $[H^\dagger H]=2$, so the Higgs portal is super–renormalizable (relevant), whereas $\mathcal S G^2$ and $\mathcal S F^2$ are dimension–5 and therefore suppressed by $1/\Lambda$ (marginal/irrelevant once the suppression is included). This operator–dimension hierarchy already favors the electroweak portal over a direct color portal.

\paragraph{Assumptions and alignment at the AdS$_4$ fixed point.}\label{par:uom-main-132}
We adopt two standard model–building inputs (they are \emph{not} symmetry selection rules):
\begin{enumerate}[label=(A\arabic*),leftmargin=2.1em,itemsep=2pt]
\item \emph{UV boundary condition (no bare color portal).} At the UV brane we impose no brane–local counterterm of the form $\mathcal S\,G^2$. This keeps $c_g/\Lambda_g=0$ at tree level in the UV.
\item \emph{Operator alignment.} The welded source excites a scalar operator orthogonal (in the CFT two–point/Fisher sense) to the trace/compensator direction. This eliminates the anomaly–induced contact to $G^2$ associated with $T^\mu{}_\mu$ at tree level. 
\end{enumerate}
With (A1)–(A2), any $\mathcal S\,G^2$ coupling arises only radiatively (see below).

\paragraph{CP and parity.}\label{par:uom-main-133}
The portal sector preserves CP and P, so the CP–odd $\mathcal S\, G\tilde G$ operator is absent.

\paragraph{Loop–induced and anomaly–induced gluon couplings.}\label{par:uom-main-134}
After electroweak symmetry breaking, $\mathcal S$ mixes with the Higgs by
\begin{equation}
\theta_{S\!H}\;=\;\frac{\lambda_p\,v}{m_{\mathcal S}^2 - m_h^2},\qquad |\theta_{S\!H}|\ll 1
\label{eq:thetaSH}
\end{equation}
in the regime of interest. This mixing unavoidably generates an effective
\begin{equation}
\mathcal L_{\rm eff}\supset
\theta_{S\!H}\,\frac{\alpha_s}{12\pi v}\,\mathcal S\,G_{\mu\nu}^a G^{a\,\mu\nu}
\qquad\Longrightarrow\qquad
\frac{c_g}{\Lambda_g}\ \sim\ \theta_{S\!H}\,\frac{\alpha_s}{12\pi v}.
\label{eq:loopSgg}
\end{equation}
If a trace portal is kept, the QCD trace contribution
\begin{equation}
T^\mu{}_\mu{}^{\rm (QCD)}
=\frac{\beta(g_s)}{2g_s}\,G_{\mu\nu}^a G^{a\,\mu\nu}
+\sum_q (1+\gamma_m)m_q\,\bar q q
\label{eq:traceQCD}
\end{equation}
induces an additional
\begin{equation}
\frac{c_g}{\Lambda_g}\bigg|_{\rm anom}\ \sim\ \frac{\kappa_T}{\Lambda_T}\,\frac{\beta(g_s)}{2g_s}.
\label{eq:trace-anom}
\end{equation}
At low scales $|\beta(g_s)/(2g_s)|\ll 1$, so the anomaly piece is numerically small. Combining \eqref{eq:loopSgg}–\eqref{eq:trace-anom} with \eqref{eq:thetaSH} yields the parametric bound
\[
\frac{c_g}{\Lambda_g}\ \lesssim\ \underbrace{\frac{\lambda_p v}{m_{\mathcal S}^2}}_{\text{mixing}}\times\underbrace{\frac{\alpha_s}{12\pi v}}_{\text{loop}}\ +\ \underbrace{\frac{\kappa_T}{\Lambda_T}}_{\text{trace portal}}\times\underbrace{\frac{\beta(g_s)}{2g_s}}_{\text{small}}.
\]

\paragraph{Summary.}\label{par:uom-main-135}
Thanks to (i) the operator–dimension hierarchy (super–renormalizable $\mathcal S H^\dagger H$ vs.\ dimension–5 $\mathcal S G^2$), (ii) the UV boundary condition (no bare $\mathcal S G^2$), and (iii) operator alignment away from the trace/compensator direction, any $\mathcal S G^2$ coupling is at most \emph{loop/anomaly induced} and therefore parametrically small, while the direct mass deformation of $H^\dagger H$ remains unsuppressed and controls EWSB. 

\subsubsection*{Parameter summary and constraints}\label{subsubsec:uom-main-015}
Collecting the relations,
\begin{equation}
\label{eq:param-summary}
v^2\;=\;\frac{1}{\lambda_H}\left(\lambda_p\,\frac{J_0}{m_{\mathcal S}^2}-\mu_{H,0}^2\right),
\qquad
J_0\;=\;\kappa_{\rm HS}\,\gamma\!\int\!\frac{d^4q}{(2\pi)^4}\,\mathcal K_{\rm anom}(q)\,|R(q)|^2\ >\ 0,
\end{equation}
with $\gamma=\ell^2/96$ fixed holographically and $\mathcal K_{\rm anom}$ the DC limit of the trace
kernel (Sec.~\ref{subsec:receiver-anomaly}). Given the welded LoG history (which fixes the spectral
weight $\int|R|^2$) and a choice of $(m_{\mathcal S},\lambda_p)$ in the technically natural window,
\eqref{eq:param-summary} yields the observed $v$ with a \emph{constant} post–flip imprint; radion
mixing is negligible if $m_{\rm rad}\gg \mu$ and is \emph{absent} on the KR head.

\medskip
\noindent\emph{Conclusion.} The anomaly–LoG portal with a stabilized (and, post–flip, \emph{absent})
radion provides a minimal, symmetry–respecting mechanism that (i) obeys P1, (ii) fixes a constant Higgs
mass at freeze, (iii) is convex/natural within the unified P3/SK framework, and (iv) is compatible with
all receiver–side locks and orthogonality relations established in Sec.~\ref{subsec:P3-master}.

\section{Two--times semiclassical characterization: acceptance, diagonalization, and flatness}
\label{sec:two-times-semiclassical}

\subsection{Two-time structural package and assumption sets for the 2T theorems}
\label{subsec:two-times-assumptions}

\paragraph{Purpose and crosswalk.}
This subsection isolates the abstract two-time data extracted from UOM and states the assumption sets used in the
2T Limited Scope and 2T Main theorems. The rest of this section proves the semiclassical realization and the
renormalized Peierls package (UOM-P1--P6) from the SK/HJ structure.

\begin{definition}[2T Limited Scope structural package (UOM form)]
\label{def:uom-limited-scope-package}
Fix a receiver cutoff $\Lambda<\infty$ and its bandlimit $\Delta(\Lambda)$.
Let $\mathcal S_{\rm SK}$ be the set of admissible SK states at this cutoff and
let $\mathcal S_{\rm rec}$ be the set of record states (diagonal density matrices in the pointer basis at bandlimit $\Delta$).
We assume the following data:
\begin{enumerate}[label=\textbf{(A\arabic*)},leftmargin=*,itemsep=0.4em]
\item \textbf{Core SK data.}
A diagonal embedding $\operatorname{diag}:\mathcal S_{\rm rec}\to\mathcal S_{\rm SK}$ and a swap involution
$s:\mathcal S_{\rm SK}\to\mathcal S_{\rm SK}$ with $s^2=\mathrm{id}$ and $s\circ\operatorname{diag}=\operatorname{diag}$.

\item \textbf{Two arrows (time and RG).}
A time group action $T_t:\mathcal S_{\rm SK}\to\mathcal S_{\rm SK}$ of $(\mathbb R,+)$ and an RG semigroup action
$U_u:\mathcal S_{\rm SK}\to\mathcal S_{\rm SK}$ of $(\mathbb N,+)$,
with weak commutativity $U_u(T_t x)=T_t(U_u x)$.

\item \textbf{Strict diagonalization of RG.}
For each $u\ge 1$ a map $\Pi_u:\mathcal S_{\rm SK}\to\mathcal S_{\rm rec}$ such that
$U_u(x)=\operatorname{diag}(\Pi_u(x))$. There is a record-level action
$R_u:\mathcal S_{\rm rec}\to\mathcal S_{\rm rec}$ with $R_0=\mathrm{id}$ and
$R_{u_2+u_1}=R_{u_2}\circ R_{u_1}$, and
$\Pi_u(\operatorname{diag} r)=R_u(r)$.

\item \textbf{RG-sector reachability.}
A section $\mathrm{sec}_\Pi:\mathcal S_{\rm rec}\to\mathcal S_{\rm SK}$ with
$\Pi_1(\mathrm{sec}_\Pi(r))=r$.

\item \textbf{No mixed anomaly (two-time flatness).}
There is a bisimplicial set $\mathcal C^{\rm SK}_{m,n}$ with commuting horizontal/vertical face maps,
so the two-time square has a well-defined boundary and mixed faces commute.

\item \textbf{Record reflection.}
$\operatorname{diag}$ is injective: $\operatorname{diag}(r)=\operatorname{diag}(r')\Rightarrow r=r'$.

\item \textbf{Real structure and two-time mixing.}
There exists a pointer basis in a Hilbert space $\mathcal H$ with
$\operatorname{diag}(\mathcal S_{\rm rec})$ consisting of diagonal trace-class operators,
the swap is transposition in that basis, and time evolution is generated by a bounded
self-adjoint $H$ (bounded at each fixed accepted bandlimit). There is a scaling generator $D$ on a common dense domain with
$(DH-HD)\psi=iH\psi$, and a nondegeneracy condition: for some record state
$\rho=\operatorname{diag}(r)$ there exist $a\ne b$ with $\Re(H_{ab})\ne 0$.
\end{enumerate}
\end{definition}

\begin{proposition}[UOM realization of the Limited Scope package]
\label{prop:uom-limited-scope-realization}
In the semiclassical UOM regime at finite cutoff, items \textbf{(A1)}--\textbf{(A6)} are supplied by the SK
state space and swap, the moving-wall RG arrow and acceptance map, and the exact two-time flatness on scalar
sources derived from the master potential (Subsec.~\ref{subsec:two-times-semiclassical-extraction}).
Item \textbf{(A7)} is supplied by the finite-band pointer/record structure of
Definition~\ref{def:pointer-algebra-projection} together with the UV dilatation theorem
proved below in Theorem~\ref{thm:uom-uv-dilatation-package}: the QND pointer basis gives the
real structure, while the semiclassical RG/Weyl action provides the mixing pair $(H,D)$.
\end{proposition}

\begin{theorem}[Equal-source SK swap invariance on the Minkowski UV chart]
\label{thm:uom-equal-source-swap}
On the physical SK locus of the semiclassical UOM package, where the sources on
the two SK branches are equal, the renormalized SK functional is invariant
under the branch-swap involution \(s:\mathcal S_{\rm SK}\to\mathcal S_{\rm SK}\).
Consequently, the induced equal-source vacuum correlators on the Minkowski UV
boundary are invariant under branch exchange.
\end{theorem}

\begin{proof}
On the physical SK locus the source coupling on the \(+\) and \(-\) branches is
identical. Exchanging the two branches is therefore a change of variables in the
closed-time-path functional integral that preserves the doubled renormalized
action, the renormalized measure, and the fixed equal-source boundary data.
Because the contour is closed and the branch couplings coincide, no endpoint term
is generated by this exchange. Hence the renormalized SK functional is
swap-invariant on the equal-source locus.

The equal-source vacuum correlators on the Minkowski UV boundary are obtained by
differentiating this renormalized generating functional with respect to boundary
insertions and then restricting back to the equal-source locus. Since the
generating functional is invariant there, every resulting vacuum correlator is
invariant under the induced branch-swap involution.
\end{proof}

\begin{definition}[2T Main Theorem assumption package]
\label{def:uom-main-theorem-assumptions}
The 2T Main Theorem assumes the Limited Scope package above and, in addition:
\begin{enumerate}[label=\textbf{(M\arabic*)},leftmargin=*,itemsep=0.4em]
\item \textbf{(M1) Limited Scope output.}
The Limited Scope Theorem applies to \textbf{(A1)}--\textbf{(A7)},
yielding a canonical $\mathbb Z_2$-cover, a parity cocycle, and a nontrivial odd descendant
$\mathcal O^{(1)}$.

\item \textbf{(M2) RG monotone in the sufficiently local class.}
The receiver RG monotone (Subprogram~3; anomaly/Wess--Zumino sector) defines functionals in
$\mathcal F_{\rm s.loc}^\Lambda$ and is strictly monotone along the RG arrow.

\item \textbf{(M3) Renormalized Peierls package.}
The UOM-P1--P6 package holds, providing a gauge-independent, renormalized Peierls bracket on
$\mathcal F_{\rm s.loc}^\Lambda$ (Subsec.~\ref{subsec:uom-peierls-package}).

\item \textbf{(M4) Odd Hamiltonization.}
The odd descendant $\mathcal O^{(1)}$ is represented by an element of $\mathcal F_{\rm s.loc}^\Lambda$ and
$\delta_\omega(\cdot):=\{\mathcal O^{(1)},\cdot\}_{\rm P}$ is a well-defined odd derivation on
$\mathcal F_{\rm s.loc}^\Lambda$.

\item \textbf{(M5) Two-time covariance.}
The time action and RG action preserve $\mathcal F_{\rm s.loc}^\Lambda$ and act by automorphisms of the
Peierls bracket, so the odd derivation is compatible with the two-time dynamics.
\end{enumerate}
\end{definition}

\begin{remark}
Assumption \textbf{(M3)} is precisely where the Peierls bracket is derived within the scope of the 2T Main Theorem,
and it is supplied in UOM by the renormalized causal/Peierls package (UOM-P1--P6) below.
\end{remark}

\subsection{Semiclassical extraction: two arrows, acceptance, and flatness}
\label{subsec:two-times-semiclassical-extraction}
\paragraph{Overview.}\label{par:uom-main-093}
We now extract, directly from the renormalized SK influence functional and the moving--wall Hamilton--Jacobi
structure, the semiclassical two--time package used later as an input: (i) an effective Lorentzian time action
and an RG semigroup action, (ii) a physical acceptance/windowing map inducing diagonalization to records above
a chosen scale, (iii) local two--time flatness (absence of mixed anomalies on scalar sources), and (iv) a
curvature--budget identity $\sum_k\Theta_k=\xi_0-\xi_n$ derived from a single real master potential on the
extended state space. No cohomological or external machinery is used here.

\subsubsection{Two arrows from the 5D SK saddle}
\label{subsubsec:two-times-arrows}
We distinguish two operational arrows:
(i) Lorentzian time $t$ (a reversible, group--like action on SK states in the semiclassical regime), and
(ii) RG time $u:=\ln(\Lambda_{\rm UV}/\Lambda)$ (an irreversible, semigroup flow generated by moving--wall/Wilsonian
coarse graining), where $\Lambda_{\rm UV}$ is the fixed AdS$_5$ UV boundary reference cutoff, not an
aeon--internal terminal value. We treat $u$ as a depth coordinate for the UV$\to$IR CPTP semigroup at fixed
$\Lambda$, while the scheduler relocates the cutoff surface $\Lambda(t)$ (hence $u(t)$) toward the UV;
the semigroup direction remains UV$\to$IR, but the endpoint is moved tick by tick.
The bulk radial coordinate is spacelike; the point is that the induced RG parameter is \emph{time--like
operationally}: it orders an irreversible acceptance/refinement process on the boundary and acts as a scalar
flow in the boundary conformal algebra. This is the two--time structure used throughout.

\subsubsection{Acceptance as a physical map and semiclassical diagonalization}
\label{subsubsec:acceptance-diagonalization}
Define the receiver admission as a CPTP instrument $\{\mathcal E_\alpha\}$ acting on the offered
tick state $\rho_k^{\rm offer}$, with pointer projectors $\{\Pi_\alpha\}$ as in
Def.~\ref{def:pointer-algebra-projection}. Let $\Delta I_k^{\rm offer}$ denote the payload offered
by the source/transfer channel in tick $k$. The admitted payload is the corresponding expectation
\begin{equation}
\Delta I_k\ :=\ \Big(\sum_\alpha \operatorname{Tr}\!\big[\Pi_\alpha\,\mathcal E_\alpha(\rho_k^{\rm offer})\big]\Big)\,
\Delta I_k^{\rm offer}.
\label{eq:admitted-payload}
\end{equation}
Define the shorthand admission factor
$\mathcal A_{\rm rec}:=\sum_\alpha \operatorname{Tr}\!\big[\Pi_\alpha\,\mathcal E_\alpha(\rho_k^{\rm offer})\big]\in[0,1]$,
so the old scalar relation $\Delta I_k=\mathcal A_{\rm rec}\Delta I_k^{\rm offer}$ is a derived notation,
not the defining map. In the split case the instrument has two classical branch outcomes and
$\eta_{\rm split}$ is the corresponding branch weight after diagonalization.
We will suppress the ``offer'' superscript below and write $\Delta I_k$ for the admitted payload unless stated otherwise.
Retardedness constrains $\Delta I_k^{\rm offer}$ to depend only on past receiver data (e.g.\ $\Lambda_{\le k-1}$),
while $\mathcal A_{\rm rec}(\Lambda_k,\text{window})$ is the instantaneous admission at tick $k$.

In Keldysh $(r/a)$ variables, the renormalized quadratic SK action takes the standard form
\begin{equation}
S^{(2)}_{\rm SK,ren}[\Phi_r,\Phi_a]
\;=\;\langle\Phi_a,\,\mathcal D_{\rm ret}\,\Phi_r\rangle
\; +\;\frac{i}{2}\,\langle\Phi_a,\,\mathcal N\,\Phi_a\rangle,
\qquad \mathcal N\succeq 0
\label{eq:SK-ra-quadratic}
\end{equation}
with $\mathcal N$ fixed by KMS/OS positivity.

\begin{definition}[Pointer/record algebra and canonical record projection at bandlimit $\Delta$]
\label{def:pointer-algebra-projection}
Fix the band projector $P_\Delta$ and set $\mathcal H_\Delta:=\mathrm{Ran}(P_\Delta)$.
Let $H_{{\rm sp},\Delta}:=P_\Delta H_{\rm sp}P_\Delta$ be the bounded single--particle Hamiltonian on
$\mathcal H_\Delta$, and let $\{P_\lambda\}_\lambda$ denote its spectral projections
(so $P_\lambda P_{\lambda'}=\delta_{\lambda\lambda'}P_\lambda$ and $\sum_\lambda P_\lambda=I_{\mathcal H_\Delta}$).

The \emph{pointer/record algebra} at bandlimit $\Delta$ is the commutative von Neumann algebra
\[
\mathfrak A_\Delta\ :=\ \{A\in\mathcal B(\mathcal H_\Delta):[A,H_{{\rm sp},\Delta}]=0\}
\ =\ \bigoplus_\lambda \mathbb C\,P_\lambda.
\]
The \emph{canonical record projection} (conditional expectation) is the CPTP idempotent map
\begin{equation}
\Pi_\Delta(\rho)\ :=\ \sum_\lambda P_\lambda\,\rho\,P_\lambda,
\qquad \rho\in\mathcal T_1(\mathcal H_\Delta).
\label{eq:PiDelta}
\end{equation}
Its range consists of exactly those trace--class operators that commute with $H_{{\rm sp},\Delta}$
(equivalently: block--diagonal in the pointer basis; diagonal if $H_{{\rm sp},\Delta}$ is nondegenerate).
\end{definition}

\begin{lemma}[Dephasing tick fixes the record algebra and commutes with $\Pi_\Delta$]
\label{lem:tick-fixes-record-algebra}
Let $\Phi_\Delta^{(k)}$ be the one--tick propagator on $\mathcal H_\Delta$ defined by the QND GKLS generator
\eqref{eq:app-gkls}. Under the hypotheses of Lemma~\ref{lem:app-cptp}, $\Phi_\Delta^{(k)}$ admits the closed
form \eqref{eq:app-phi-closed} and is a (commuting) dephasing channel in the spectral decomposition of
$H_{{\rm sp},\Delta}$. In particular:
\begin{enumerate}[leftmargin=2.0em]
\item \textbf{Record observables are fixed.} For every $A\in\mathfrak A_\Delta$,
\[
(\Phi_\Delta^{(k)})^\dagger(A)=A,
\qquad\text{equivalently}\qquad
\Tr\!\big(\Phi_\Delta^{(k)}(\rho)\,A\big)=\Tr(\rho A)\ \ \forall\rho.
\]
\item \textbf{Commutation with record projection.} The canonical projection $\Pi_\Delta$ satisfies
\[
\Pi_\Delta\circ \Phi_\Delta^{(k)}=\Pi_\Delta,
\qquad
\Phi_\Delta^{(k)}\circ \Pi_\Delta=\Pi_\Delta.
\]
\end{enumerate}
\end{lemma}

\begin{proof}
By Lemma~\ref{lem:app-cptp} (Step~4), the one--tick map admits the Gaussian unitary average
\[
\Phi_\Delta^{(k)}(\rho)=\mathbb E\!\left[U\,\rho\,U^\dagger\right],
\qquad
U=\exp\!\Big(-i\sum_a X_a L_a\Big),
\]
with $X_a$ independent real Gaussians. Since $L_a=f_a(H_{{\rm sp},\Delta})$, each $L_a$ commutes with
$H_{{\rm sp},\Delta}$ and therefore with every spectral projector $P_\lambda$ and every
$A\in\mathfrak A_\Delta=\{H_{{\rm sp},\Delta}\}'$. Hence $U$ commutes with $A$ and with each $P_\lambda$.

\medskip
\noindent\textbf{Record observables are fixed.}
For any $A\in\mathfrak A_\Delta$,
\[
(\Phi_\Delta^{(k)})^\dagger(A)
=\mathbb E\!\left[U^\dagger A U\right]
=A,
\]
and the trace identity $\Tr(\Phi_\Delta^{(k)}(\rho)A)=\Tr(\rho A)$ follows.

\medskip
\noindent\textbf{Commutation with $\Pi_\Delta$.}
Let $\{P_\lambda\}$ be the spectral projections of $H_{{\rm sp},\Delta}$. Using $U P_\lambda=P_\lambda U$,
for any state $\rho$,
\[
P_\lambda\,\Phi_\Delta^{(k)}(\rho)\,P_{\lambda'}
=\mathbb E\!\left[U\,P_\lambda\rho P_{\lambda'}\,U^\dagger\right].
\]
In particular, the diagonal blocks are fixed:
$P_\lambda\Phi_\Delta^{(k)}(\rho)P_\lambda=P_\lambda\rho P_\lambda$.
Therefore
\[
\Pi_\Delta\!\big(\Phi_\Delta^{(k)}(\rho)\big)
=\sum_\lambda P_\lambda\,\Phi_\Delta^{(k)}(\rho)\,P_\lambda
=\sum_\lambda P_\lambda\rho P_\lambda
=\Pi_\Delta(\rho).
\]
If $\rho$ is already block--diagonal (i.e.\ $\rho=\Pi_\Delta(\rho)$), then all off--diagonal blocks vanish and
the same computation shows $\Phi_\Delta^{(k)}(\rho)=\rho$, hence
$\Phi_\Delta^{(k)}\!\circ\Pi_\Delta=\Pi_\Delta$. This yields
$\Pi_\Delta\circ\Phi_\Delta^{(k)}=\Phi_\Delta^{(k)}\circ\Pi_\Delta=\Pi_\Delta$.
\end{proof}

\begin{proposition}[Exact diagonalization of the accepted RG state (finite $u$)]
\label{prop:semiclassical-diag}
Fix a receiver cutoff $\Lambda$ (equivalently $u=\ln(\Lambda_{\rm UV}/\Lambda)$) and its associated bandlimit
$\Delta=\Delta(\Lambda)$ from the acceptance/windowing rule. Define the discrete \emph{accepted RG state} map by
\begin{equation}
\mathcal R_{\rm acc}(u,\rho)\ :=\ \Pi_{\Delta(u)}(\rho),
\qquad u\ge u_\star.
\label{eq:R-acc}
\end{equation}
Then $\mathcal R_{\rm acc}(u,\rho)$ lies in the record/diagonal sector \emph{exactly} at finite $u$ and
preserves all record expectation values: for every $A\in\mathfrak A_{\Delta(u)}$,
\[
\Tr\!\big(\Phi_{\Delta(u)}^{(k)}(\rho)\,A\big)=\Tr\!\big(\mathcal R_{\rm acc}(u,\rho)\,A\big)=\Tr(\rho A).
\]
Thus, above the acceptance scale, the RG arrow closes on record states without any appeal to external machinery:
the only information retained by the discrete RG state is the classical record on the commuting pointer algebra.
Residual coherences live entirely in the complement of $\Pi_{\Delta(u)}$ and are multiplicatively suppressed by the
dephasing kernel $K_{ij}$ of Lemma~\ref{lem:app-cptp}.
\end{proposition}

\begin{remark}[Why this is not ``definitional classicality'']
\label{rem:uom-main-004}
The map $\Pi_\Delta$ in \eqref{eq:PiDelta} is not an arbitrary postulate: it is the unique CPTP idempotent
conditional expectation compatible with the commuting QND dephasing structure of Lemma~\ref{lem:app-cptp}.
Equation \eqref{eq:R-acc} states that the accepted (compiled) RG state is the record state on the pointer
algebra; Lemma~\ref{lem:tick-fixes-record-algebra} guarantees that this record is exactly stable under
a dephasing tick and contains precisely the observables that survive as ``facts'' in the QND regime.
\end{remark}

\paragraph{Long memory and retardedness.}\label{par:uom-main-094}
If the source commutant bath has long memory (KMS correlation time $\tau_{\rm KMS}$ not small vs.\ tick length),
the SK kernels $G_R,\mathcal N$ acquire nonlocal time tails. Retardedness is exact; the quadratic form
\eqref{eq:W-quadratic} remains positive (convex). Tickwise P2 still holds because $\Im W\ge 0$ and the QND
alignment saturates the bound at fixed $\Lambda$.

\paragraph{Post--flip decay.}\label{par:uom-main-095}
After the flip the driver $J$ vanishes; all $k>0$ response decays with AdS$_4$ gaps, while the DC piece fixed
by $\Im W$ and the portal via $\mathcal S$ persists. Thus the Higgs vev is frozen after ring--down.

\subsubsection{Master potential, flatness, and curvature budget}
\label{subsubsec:master-potential}
Work on the locked manifold (welded widths and admissible windows), and extend the state space by the
cumulative admitted payload $I$ and a DC curvature/trace coordinate $\xi$ (not to be confused with $\Xi$).
Introduce a single real master potential $\mathcal G(t,u,I;\text{locks})$ whose exact differential defines
the unique conjugacies:
\begin{equation}
\mathcal H_{\rm wall}\ :=\ -\partial_u\mathcal G,\qquad
\xi\ :=\ -\partial_I\mathcal G.
\label{eq:master-potential-defs}
\end{equation}
We use the exact one--form $\alpha:=d\mathcal G$ on the extended state space; local flatness and the curvature
budget identity are corollaries of exactness along the realized trajectory. With the counterterm scheme fixed
and on a contractible $(t,u)$ patch, $\mathcal G$ is single--valued on scalar sources; hence mixed partials
commute wherever $\mathcal G$ is smooth (local two--time flatness on scalars).

Along an admitted trajectory $I_{k+1}=I_k+\Delta I_k$ one obtains
\begin{equation}
\xi_{k+1}-\xi_k\;=\;-\big(\partial_I^2\mathcal G\big)\big|_{\rm lock}\,\Delta I_k.
\label{eq:xi-step}
\end{equation}
In the welded linear--response regime the DC stiffness $\partial_I^2\mathcal G$ depends on the state only
through $\Lambda$ (up to higher--order corrections), so we write
\[
\partial_I^2\mathcal G\big|_{\rm lock}\ :=\ K(\bar\Lambda_k),\qquad
\Theta_k:=K(\bar\Lambda_k)\,\Delta I_k,
\]
which coincides with the welded BY weight in \eqref{eq:BY-k}. Therefore
\begin{equation}
\sum_{k<n}\Theta_k\;=\;\xi_0-\xi_n.
\label{eq:budget-identity}
\end{equation}
We do not introduce dynamical phases on the receiver: all tickwise statements are expressed along the
monotone trajectory $\bar\Lambda_k\uparrow$, and the only branch point is flip vs exhaustion.
In the wall--work--only / continuous--control regime, if $\xi_n\to 0$ or falls below the operational resolution,
then $\sum_k\Theta_k\to\xi_0$, explaining the near--exhaustion observed in that limit. A minimum payload quantum
(or flux quantization) gives the last--step/overshoot estimate.

\begin{proposition}[Upper envelope from the first--tick dephasing scale]
\label{prop:upper-envelope}
Let $N$ be the first--crossing tick and $\delta_N:=\sum_{k\le N}\Theta_k-|\langle R\rangle_0|$ the overshoot,
with $\Theta_k:=K(\bar\Lambda_k)\Delta I_k$. Then
\[
0\le \delta_N<\Theta_N=K(\bar\Lambda_N)\Delta I_N.
\]
Assume (i) $K(\Lambda)$ is nonincreasing along the schedule, and (ii) in the pre--flip epoch the admitted
payload satisfies the one--tick bound $\Delta I_k\le F_k(1-e^{-\beta_k})$ with $F_k\le F_0$ and
$\beta_k\ge\beta_0$ (e.g.\ steady QND tail or a uniform QSL lower bound). Then
\[
\delta_N < K(\bar\Lambda_0)\,F_0\,(1-e^{-\beta_0}),
\qquad \beta_0=\kappa\,H_{\rm prev},
\]
so the next--aeon split-side curvature is bounded above by a function of $H_{\rm prev}$ (hence of the
initial curvature scale).
\end{proposition}

\begin{theorem}[Exhaustion infimum]
\label{thm:exhaustion-infimum}
In the ideal no--leakage and arbitrarily fine payload limit, the post--flip residual is arbitrarily small:
\(\inf |R_{\rm next}^{\rm split}|=0\). With leakage $\ell$ and minimum quantum $\Delta I_{\min}$, the achievable
floor is \(|R_{\rm next}^{\rm split}|\lesssim K_{\max}(\Delta I_{\min}+O(\ell))\).
\end{theorem}

\paragraph{Scope.}\label{par:uom-main-096}
The statements above are policy--conditioned: $\Delta I_k$ depends on the admission map and may depend on
$\Lambda$ through the RG window, so schedule dependence persists across different policies. The flatness
statement is local and holds away from acceptance knees/nonsmooth window transitions. This section derives
only the two--time package (arrows, diagonalization, local flatness, curvature budget) without invoking
external cohomological machinery.

\begin{proposition}[Physical radial Ward action on compact UV head observables]
\label{prop:uom-uv-radial-ward}
Let
\[
X(u)=\int_{\mathscr T\times S^3}dt\,d^3x\;\kappa_k(t)\eta_k(x)L_X(t,x;u)
\]
be a compactly smeared scalar head observable on the locked branchwise
chart, with \(u\)-independent compact kernels and a local receiver-head
scalar density \(L_X\) that is smooth in \(u\).
Then the physical radial Ward differential is the geometric
\(u\)-derivative:
\[
Q_u^{\mathrm{phys}}X=-\,\partial_u X.
\]
Equivalently, constant boundary Weyl rescalings on the Minkowski UV
chart induce an algebraic one-parameter action
\(\sigma_\tau^{\mathrm{sc,UV}}\) on this compactly smeared scalar class
whose infinitesimal generator satisfies
\[
\left.\frac{d}{d\tau}\right|_{\tau=0}\sigma_\tau^{\mathrm{sc,UV}}(X)
=Q_u^{\mathrm{phys}}X=-\,\partial_u X.
\]
\end{proposition}

\begin{proof}
On scalar sources the two-time master potential is single-valued on the
chosen contractible \((t,u)\) patch, so the mixed partials commute and
there is no independent mixed-anomaly contribution in the fixed
counterterm scheme.  The radial variation of \(X(u)\) is therefore
exhausted by the explicit \(u\)-dependence of the local head density:
\[
\partial_u X(u)
=
\int_{\mathscr T\times S^3}dt\,d^3x\;\kappa_k(t)\eta_k(x)\,\partial_u L_X(t,x;u).
\]
By the same fixed-scheme convention used throughout the descendant lane,
the physical radial Ward operator is identified with minus this
geometric derivative, which gives the first displayed identity.

Now let \(s_\tau(x):=e^{-\tau}x\) on the Minkowski UV chart.  Pushing
forward the compact smearing kernels and local scalar density along
\(s_\tau\) again produces an observable of the same compactly smeared
scalar form, so \(\sigma_\tau^{\mathrm{sc,UV}}\) is well-defined
algebraically on this class.  Its infinitesimal action is precisely the
same geometric radial derivative, hence
\(
\left.\frac{d}{d\tau}\right|_{\tau=0}\sigma_\tau^{\mathrm{sc,UV}}(X)
=-\,\partial_u X
=Q_u^{\mathrm{phys}}X
\).
\end{proof}

\begin{remark}[Detailed finite-cutoff provenance]
\label{rem:uom-uv-radial-ward-provenance}
The compact-smearing identity in
Proposition~\ref{prop:uom-uv-radial-ward} is the broad-manuscript
version of the fixed-scheme radial Ward/geometric identification proved
in the current descendant note.  Here we only retain the scalar
head-observable consequence needed for the UV dilatation package.
\end{remark}

% ------------------------------------------------------------------
% Place the following subsection inside the (now promoted) Section
% \section{Two--times semiclassical characterization: acceptance,
% diagonalization, and flatness}\label{sec:two-times-semiclassical},
% immediately after the existing material of that section.
% ------------------------------------------------------------------

\subsection{Renormalized causal structure, Peierls bracket, and pointer-real swap equivariance on the receiver head (UOM-P1--P6)}
\label{subsec:uom-peierls-package}

\paragraph{Purpose and position in the logic.}
The downstream two--time program (2T Main) will \emph{Hamiltonize} an odd descendant
$\mathcal O^{(1)}$ by taking a Peierls bracket with other observables.  In the two--time
argument, \emph{oddness} and $\mathbb Z_2$--gradedness are to be derived later
(from the Limited Scope theorem and RG monotones), and therefore \emph{no} odd formalism
(ghosts, BRST, BV, gauge--fixing fermions, etc.) may be used in UOM to establish the
existence and gauge--independence of the bracket.

Accordingly, this subsection proves the following analytic package \textbf{entirely within UOM}
and \emph{using only even (bosonic) structures}:
\begin{enumerate}[label=\textbf{UOM-P\arabic*},leftmargin=2.4em,itemsep=2pt]
\item \textbf{(P1) Differentiability of the renormalized SK influence functional.}
The renormalized SK influence functional $W_{\rm ren}[J_+,J_-]$ exists and is twice
Fr\'echet differentiable near $J_a=0$ (diagonal sources), for a specified class of
admissible boundary/brane sources.
\item \textbf{(P2) Causal Green structure.}
The gauge--fixed linearized renormalized equations (with UOM boundary conditions at finite cutoff)
admit advanced/retarded Green operators $E^\pm$; the induced boundary retarded kernel is the mixed
second derivative of $W_{\rm ren}$.
\item \textbf{(P3) Renormalized covariant symplectic structure.}
The renormalized covariant symplectic current and form exist, are finite under the UOM counterterm scheme,
and are compatible with the causal propagator.
\item \textbf{(P4) Gauge independence on gauge invariants without BV/ghosts.}
Gauge independence is proved by presymplectic reduction: pure gauge directions are exactly the kernel of
$\Omega_{\rm ren}$, so the reduced symplectic form (and hence the induced bracket) is canonical; any
gauge-fixed Peierls computation agrees with the reduced bracket on gauge-invariant functionals.
\item \textbf{(P5) A closed ``sufficiently local'' class.}
There is a precisely defined class $\mathcal F_{\rm s.loc}^\Lambda$ of gauge-invariant, sufficiently local
boundary/brane functionals at finite cutoff $\Lambda$ that is closed under the Peierls bracket and contains
the quasi-local observables used in Subprogram~3 (e.g.\ anomaly/Wess--Zumino densities smeared by compact kernels).
\item \textbf{(P6) Pointer-real swap equivariance.}
At fixed finite cutoff there is a pointer-real involution $\Gamma_{*,\Lambda}$ on the admissible linearized
field, source, and boundary-data spaces, induced by the finite-band pointer basis, such that the gauge-fixed
linearized renormalized IBVP, the source insertion / observable extraction maps, and the pairing used in the
Peierls formula are $\Gamma_{*,\Lambda}$--equivariant.
\end{enumerate}
Together these yield a \emph{renormalized} Peierls bracket on the receiver head that is canonical, causal,
and ready to be imported into the two--time construction.

\paragraph{Standing geometric setup (one head, fixed cutoff, finite window).}
Fix a single ``cell'' of the UOM bulk geometry containing the \emph{active KR head} receiver
and any adjacent interior/shadow segments needed for well-posed boundary conditions.
Work at a fixed finite receiver cutoff $\Lambda<\infty$ (static-wall gauge) and restrict
attention to a finite Lorentzian time slab $\mathscr T\times\Sigma_{\rm rec}$, where
$\mathscr T\subset\mathbb R$ is a bounded open interval (e.g.\ one tick window $I_k$ or a finite union of
such windows) and $\Sigma_{\rm rec}\cong S^3$ is the receiver spatial slice.

Let $\mathcal M^{(\Lambda)}$ denote the resulting 5D spacetime region-with-boundary. Its boundary is the union of:
\begin{itemize}[leftmargin=1.6em,itemsep=2pt]
\item the receiver brane/worldvolume $\Sigma_{\rm rec}$ (KR head),
\item possible interior timelike branes $\Sigma_{\rm int}$ (RS1-type shadow segments),
\item an outer timelike cutoff surface $\Sigma_{\Lambda}$ implementing the finite Wilsonian cutoff,
\item and any end-of-the-world caps required by the KR construction.
\end{itemize}
Boundary/junction conditions are those dictated by the UOM renormalized action and
healthiness hypotheses (no ghosts/tachyons, stabilized radions in interior segments, Robin/induced-gravity
healthy branch on the KR head; cf.\ Theorem~\ref{thm:OS_receiver} for OS positivity on the receiver and the
healthiness list (H1)--(H4) therein).

\paragraph{Field content and action.}
Let $\Phi$ denote the collection of \emph{bosonic} fields present in this cell (bulk metric, bulk stabilizers
such as GW scalars, receiver compensator $\sigma$, receiver portal field $\mathcal S$ if retained, and any
other bosonic auxiliaries needed to keep locality).  Let $S_{\rm ren}[\Phi]$ be the Lorentzian
\emph{renormalized} action on $\mathcal M^{(\Lambda)}$ including all boundary terms and holographic counterterms:
\[
S_{\rm ren} \;=\; S_{\rm bulk} \; + \; S_{\rm GHY} \; + \; S_{\rm branes} \; + \; S_{\rm ct}.
\]
The counterterm scheme is fixed once and for all (as in the Brown--York and anomaly discussion of
Sec.~\ref{subsec:receiver-anomaly}).  We assume the renormalized variational principle is well-posed:
$\delta S_{\rm ren}$ has no leftover boundary variation terms once the UOM boundary/junction conditions are imposed.

\paragraph{Observables and sources.}
Fix a finite family of gauge-invariant, local (or quasi-local) bosonic densities
$\mathcal O_A(\Phi)$ on the receiver head and/or on the cutoff surface (including, in particular, the trace/Weyl
sector and any densities used to build the RG monotone in Subprogram~3).
For a source multiplet $J=\{J^A\}$ we define the sourced action
\begin{equation}
S_{\rm ren}[\Phi;J]\;:=\;S_{\rm ren}[\Phi]\;+\;\int_{\mathscr T\times\Sigma_{\rm rec}}\!\!\!\! d^4x\;\sqrt{-\gamma}\;J^A(x)\,\mathcal O_A(\Phi)(x),
\label{eq:uom-sourced-action}
\end{equation}
where $\gamma$ is the induced metric on $\Sigma_{\rm rec}$ and repeated indices are summed.
(If some $\mathcal O_A$ live on other boundary components, include the corresponding boundary integrals; this
does not change the arguments below.)

\paragraph{SK doubling and influence functional.}
On $\mathcal M^{(\Lambda)}$ we consider the Schwinger--Keldysh doubled fields $\Phi_\pm$ and sources $J_\pm$.
The renormalized SK action is
\[
S_{\rm SK,ren}[\Phi_+,\Phi_-;J_+,J_-]
\;:=\;S_{\rm ren}[\Phi_+;J_+]\; -\; S_{\rm ren}[\Phi_-;J_-],
\]
with the usual SK gluing conditions at the future end of the contour.
In the semiclassical regime used throughout UOM, define the (connected) SK influence functional
\begin{equation}
W_{\rm ren}[J_+,J_-]\;:=\;-\,i\log Z_{\rm ren}[J_+,J_-],
\qquad
Z_{\rm ren}[J_+,J_-]\;:=\;\int\mathcal D\Phi_+\,\mathcal D\Phi_-\,e^{\,iS_{\rm SK,ren}[\Phi_+,\Phi_-;J_+,J_-]},
\label{eq:Wren-def}
\end{equation}
For the Peierls bracket we will only require the \emph{tree-level} (on-shell) information of $W_{\rm ren}$ near the
diagonal $J_a=0$; loop corrections are not needed for the classical bracket and are therefore omitted here.

\paragraph{Keldysh variables.}
We write the standard Keldysh combinations
\[
J_r:=\frac{J_++J_-}{2},\qquad J_a:=J_+-J_-,
\qquad
\Phi_r:=\frac{\Phi_++\Phi_-}{2},\qquad \Phi_a:=\Phi_+-\Phi_-.
\]
The diagonal sector is $J_a=0$ (physical sources), and the physical background corresponds to $J_r=0$.

\paragraph{Main statement (Peierls package).}
The following theorem summarizes what is proved in the six subsubsections below and is the only object imported
by the two--time program.

\begin{theorem}[UOM Peierls package at finite cutoff (UOM-P1--P6)]
\label{thm:uom-peierls-package}
Fix a healthy receiver cell $\mathcal M^{(\Lambda)}$ with finite cutoff $\Lambda<\infty$ and the renormalized action
$S_{\rm ren}$ as above.  Then, for a specified space of admissible compactly supported sources $J$ and a specified
class $\mathcal F_{\rm s.loc}^\Lambda$ of gauge-invariant sufficiently local boundary/brane functionals, the following hold:
\begin{enumerate}[label=\textnormal{(\roman*)},leftmargin=2.2em,itemsep=2pt]
\item \textbf{(UOM-P1)} $W_{\rm ren}[J_+,J_-]$ exists and is $C^2$ in $(J_+,J_-)$ on a neighborhood of $J_a=0$.
\item \textbf{(UOM-P2)} The gauge-fixed linearized renormalized operator with the UOM boundary conditions is Green-hyperbolic,
with advanced/retarded propagators $E^\pm$, and the induced boundary retarded kernel is
\[
G_R(x,y)
=\left.\frac{\delta^2 W_{\rm ren}}{\delta J_a(x)\,\delta J_r(y)}\right|_{J=0},
\qquad
G_A(x,y):=G_R(y,x),
\qquad
\Delta:=G_R-G_A.
\]
\item \textbf{(UOM-P3)} The renormalized covariant symplectic form $\Omega_{\rm ren}$ exists, is finite and conserved,
and $\Delta$ is the induced causal bivector compatible with $\Omega_{\rm ren}$.
\item \textbf{(UOM-P4)} Let $\mathcal G_0$ be the group of pure gauge transformations preserving the UOM boundary conditions
and acting trivially on the chosen boundary data.  Then $\ker\Omega_{\rm ren}$ equals the infinitesimal $\mathcal G_0$ directions,
so $\Omega_{\rm ren}$ descends to a symplectic form on $\mathcal P:=\mathcal S/\mathcal G_0$, and the induced Peierls bracket
is gauge independent on $\mathcal F_{\rm s.loc}^\Lambda$.
\item \textbf{(UOM-P5)} The Peierls bracket
\begin{equation}
\{F,G\}_{\rm P}\;:=\;\big\langle F^{(1)},\,\Delta\,G^{(1)}\big\rangle
\label{eq:Peierls-def-summary}
\end{equation}
is well-defined on $\mathcal F_{\rm s.loc}^\Lambda$, takes values in $\mathcal F_{\rm s.loc}^\Lambda$, and satisfies
antisymmetry, Leibniz, and Jacobi there.  Moreover, $\mathcal F_{\rm s.loc}^\Lambda$ contains the quasi-local
observables used to represent RG monotones (e.g.\ the Wess--Zumino/anomaly densities smeared by compact kernels).
\item \textbf{(UOM-P6)} There exists a finite-cutoff pointer-real involution $\Gamma_{*,\Lambda}$ on the admissible
linearized field, source, and boundary-data spaces such that
\[
P_{\rm gf,ren}\Gamma_{*,\Lambda}=\Gamma_{*,\Lambda}P_{\rm gf,ren},
\qquad
B_\Lambda\Gamma_{*,\Lambda}=\Gamma_{*,\Lambda}B_\Lambda,
\]
the source insertion and observable extraction maps commute with $\Gamma_{*,\Lambda}$, and the pairing used in the
Peierls formula is $\Gamma_{*,\Lambda}$--invariant.
\end{enumerate}
\end{theorem}

\paragraph{Roadmap.}
Theorem~\ref{thm:uom-peierls-package} is proved by the six subsubsections below, in order:
(UOM-P1) differentiability of $W_{\rm ren}$; (UOM-P2) Green-hyperbolicity and identification of $G_R$; (UOM-P3) renormalized
symplectic form and compatibility with $\Delta$; (UOM-P4) gauge independence by presymplectic reduction; (UOM-P5) definition of
$\mathcal F_{\rm s.loc}^\Lambda$ and closure under the bracket; (UOM-P6) exact finite-cutoff pointer-real swap equivariance.

% ============================================================
% (2) UOM-P1 subsubsection snippet
% ============================================================

\subsubsection{UOM-P1: differentiability of the renormalized SK influence functional}
\label{subsubsec:uom-P1}

This subsubsection establishes \textbf{UOM-P1}: existence and twice differentiability of the
renormalized SK influence functional $W_{\rm ren}[J_+,J_-]$ near the diagonal $J_a=0$
for a precisely specified class of admissible sources.

\paragraph{Admissible source space.}
Fix an integer $s>\frac{5}{2}$ (so that $H^s$ embeds into $C^1$ in four dimensions).
Let $\mathscr T\subset\mathbb R$ be a bounded open time interval and set $\mathscr X:=\mathscr T\times \Sigma_{\rm rec}$.
Define the Banach space of admissible sources
\begin{equation}
\mathcal J^s(\mathscr X)\;:=\;\bigoplus_{A}\,H^s_0(\mathscr X),
\label{eq:Js-def}
\end{equation}
where $H^s_0$ denotes $H^s$ functions supported in a compact subset of $\mathscr X$ (so sources vanish near the temporal endpoints
and outside a compact spatial set; on $\Sigma_{\rm rec}\cong S^3$ this simply means compact support in time).
We take $J_+,J_-\in\mathcal J^s(\mathscr X)$.

\paragraph{Gauge fixing for well-posedness (even).}
To define a unique classical response map we work in an \emph{even} gauge fixing on the space of bosonic fields.
Concretely, choose a smooth gauge-fixing functional $\mathcal G(\Phi)=0$ (e.g.\ de Donder for the metric plus any standard gauge for
bulk matter) and add an even gauge-fixing term
\begin{equation}
S_{\rm gf}[\Phi]\;:=\;\frac{1}{2\alpha}\int_{\mathcal M^{(\Lambda)}}\!\!\! d^5x\;\sqrt{-g}\;\|\mathcal G(\Phi)\|^2,
\qquad \alpha>0,
\label{eq:even-gauge-fixing}
\end{equation}
to the renormalized action.  No ghosts and no odd structure are introduced; $S_{\rm gf}$ is used only to ensure
well-posedness of the classical IBVP and invertibility of the linearized operator.

Write
\[
S_{\rm ren}^{\rm gf}[\Phi;J]\;:=\;S_{\rm ren}[\Phi;J]+S_{\rm gf}[\Phi],
\qquad
S_{\rm SK,ren}^{\rm gf}:=S_{\rm ren}^{\rm gf}[\Phi_+;J_+]-S_{\rm ren}^{\rm gf}[\Phi_-;J_-].
\]

\paragraph{Classical solution map.}
Let $\bar\Phi$ be the background solution at $J=0$ in the chosen gauge.
Denote by $\mathcal E(\Phi;J)=0$ the Euler--Lagrange equations derived from $S_{\rm ren}^{\rm gf}[\Phi;J]$, together with the UOM
boundary/junction conditions at cutoff $\Lambda$ (the variational principle boundary conditions).

\begin{corollary}[Local well-posedness and smooth dependence near the UOM background]
\label{cor:uom-wp-smooth}
There exists an open neighborhood $\mathcal U\subset \mathcal J^s(\mathscr X)$ of $0$ and an open neighborhood $\mathcal V$ of $\bar\Phi$
in a Banach space of field configurations on $\mathcal M^{(\Lambda)}$ (compatible with the boundary conditions) such that:
for each $J\in\mathcal U$ there exists a unique $\Phi[J]\in\mathcal V$ solving the gauge-fixed IBVP
$\mathcal E(\Phi[J];J)=0$, and the map $J\mapsto \Phi[J]$ is $C^2$ as a map $\mathcal U\to\mathcal V$.
\end{corollary}

\begin{proof}
The nonlinear sourced equations derived from $S_{\rm ren}^{\rm gf}$ are quasilinear hyperbolic with smooth local coefficients and
boundary operators depending smoothly on the fields and on the source $J$, because the renormalized action is local and the counterterm
scheme is fixed on the finite-cutoff receiver cell.  Proposition~\ref{prop:uom-admissible-BC-route} establishes the admissibility of the
linearized boundary problem at the background $(\bar\Phi,0)$, i.e.\ the maximally dissipative / uniform-Lopatinski estimates needed for
the linearized inverse on each bounded slab.  Standard local theory for quasilinear hyperbolic IBVPs, together with the Banach-space
implicit-function theorem applied at $(\bar\Phi,0)$, therefore yields neighborhoods $\mathcal U,\mathcal V$ and a unique solution map
$J\mapsto\Phi[J]$ which is $C^2$.
\end{proof}

\paragraph{On-shell SK functional and its differentiability.}
In the semiclassical regime, define the renormalized SK influence functional by evaluating the SK action on the classical SK saddle.
Let $\Phi_{\rm cl}[J_+,J_-]$ denote the unique doubled classical field satisfying the SK gluing conditions and sourced equations obtained
from $S_{\rm SK,ren}^{\rm gf}$, with retarded boundary conditions (equivalently: $\Phi_a$ vanishes at the initial end of the contour).
Define
\begin{equation}
W_{\rm ren}^{\rm (tree)}[J_+,J_-]\;:=\;S_{\rm SK,ren}^{\rm gf}\big[\Phi_{\rm cl}[J_+,J_-];J_+,J_-\big].
\label{eq:W-tree}
\end{equation}
(Loop corrections may be added in a standard way; they are not required for the classical Peierls bracket and are omitted.)

\begin{lemma}[Existence and $C^2$ regularity of the tree-level SK influence functional]
\label{lem:W-C2}
Under the local well-posedness statement of Corollary~\ref{cor:uom-wp-smooth} and the SK retarded gluing prescription, the map
\[
(J_+,J_-)\ \longmapsto\ W_{\rm ren}^{\rm (tree)}[J_+,J_-]
\]
is twice Fr\'echet differentiable on $\mathcal U\times \mathcal U$ and its first and second derivatives are continuous multilinear maps.
\end{lemma}

\begin{proof}
By Corollary~\ref{cor:uom-wp-smooth} (applied separately on the $+$ and $-$ branches in the retarded SK prescription),
the classical solution map $(J_+,J_-)\mapsto \Phi_{\rm cl}[J_+,J_-]$ is $C^2$ on $\mathcal U\times\mathcal U$ in the topology of $\mathcal V$.
The functional $S_{\rm SK,ren}^{\rm gf}[\Phi_+,\Phi_-;J_+,J_-]$ is $C^2$ jointly in $(\Phi_+,\Phi_-;J_+,J_-)$ on $\mathcal V\times\mathcal V\times\mathcal U\times\mathcal U$,
because it is a finite sum of integrals of smooth local densities in $\Phi_\pm$ and $J_\pm$, and the counterterm scheme is fixed.
Therefore the composition \eqref{eq:W-tree} is $C^2$ by the chain rule in Banach spaces, and its first and second derivatives
are continuous multilinear maps.
\end{proof}

\paragraph{Keldysh derivatives and the quadratic SK kernels.}
Write $W_{\rm ren}:=W_{\rm ren}^{\rm (tree)}$ for brevity.
Define the retarded and noise kernels in Keldysh variables by
\begin{equation}
G_R^{AB}(x,y)
\;:=
\left.\frac{\delta^2 W_{\rm ren}}{\delta J_{a,A}(x)\,\delta J_{r,B}(y)}\right|_{J=0},
\qquad
\mathcal N^{AB}(x,y)
\;:=
-\,i\left.\frac{\delta^2 W_{\rm ren}}{\delta J_{a,A}(x)\,\delta J_{a,B}(y)}\right|_{J=0}.
\label{eq:GR-N-def}
\end{equation}
By Lemma~\ref{lem:W-C2}, these distributions exist in the sense of continuous bilinear forms on the test-function space
dual to $\mathcal J^s$.

\begin{proposition}[Quadratic expansion and Keldysh form]
\label{prop:W-quadratic}
There exists a neighborhood of $J=0$ in $\mathcal J^s(\mathscr X)$ such that
\begin{equation}
W_{\rm ren}[J_r,J_a]
\;=
W_{\rm ren}[0,0]
\; +\;\big\langle J_a,\ \langle \mathcal O\rangle_{J_r}\big\rangle
\; +\;\frac12\big\langle J_a,\ G_R\,J_r\big\rangle
\; +\;\frac{i}{2}\big\langle J_a,\ \mathcal N\,J_a\big\rangle
\; +\;O(\|(J_r,J_a)\|^3),
\label{eq:W-keldysh-expansion}
\end{equation}
where $\langle\mathcal O\rangle_{J_r}$ is the (retarded) one-point response and the pairings are those induced by the source coupling
\eqref{eq:uom-sourced-action}.
\end{proposition}

\begin{proof}
This is the Taylor expansion of the $C^2$ functional $W_{\rm ren}$ in Banach space form, rewritten in Keldysh variables.
The coefficients of the quadratic terms are exactly the second Fr\'echet derivatives; identifying them with \eqref{eq:GR-N-def} gives \eqref{eq:W-keldysh-expansion}.
\end{proof}

\paragraph{Conclusion of UOM-P1.}
Lemma~\ref{lem:W-C2} and Proposition~\ref{prop:W-quadratic} establish \textbf{UOM-P1}:
$W_{\rm ren}$ exists at tree level and is $C^2$ near the diagonal sector $J_a=0$, with the quadratic SK kernels well-defined.
In the next subsubsection we identify $G_R$ with the retarded Green operator of the linearized gauge-fixed IBVP (UOM-P2).
% ============================================================
% (3) UOM-P2 subsubsection snippet
% ============================================================

\subsubsection{UOM-P2: Green hyperbolicity and identification of the retarded kernel}
\label{subsubsec:uom-P2}

This subsubsection establishes \textbf{UOM-P2}.  We (i) formulate the gauge-fixed linearized renormalized IBVP
for the bulk+brane system at finite cutoff $\Lambda$, (ii) state and prove Green-hyperbolicity (existence of advanced/retarded
propagators) under UOM healthiness assumptions, and (iii) identify the induced boundary retarded kernel with the mixed
second derivative of $W_{\rm ren}$.

\paragraph{Linearized gauge-fixed operator.}
Let $\bar\Phi$ be the background solution at $J=0$.  Denote by $P_{\rm gf,ren}$ the linearized Euler--Lagrange operator
associated with the gauge-fixed renormalized action $S_{\rm ren}^{\rm gf}$:
\begin{equation}
P_{\rm gf,ren}\,\delta\Phi
\;:=
\left.\frac{d}{d\epsilon}\right|_{\epsilon=0}
\mathcal E(\bar\Phi+\epsilon\,\delta\Phi;0),
\label{eq:Pgf-def}
\end{equation}
where $\mathcal E(\Phi;J)=0$ denotes the gauge-fixed renormalized field equations together with the UOM boundary/junction conditions.

On the bulk metric sector, $P_{\rm gf,ren}$ is a normally hyperbolic operator (principal part $\Box_{g}$ in de Donder gauge);
on coupled bulk/boundary scalar sectors it is likewise normally hyperbolic (Klein--Gordon type), and the boundary/junction conditions
are linear, local, and time-independent at fixed $\Lambda$.

\paragraph{Boundary conditions as an operator.}
Write the full linear boundary/junction operator as
\begin{equation}
B_\Lambda(\delta\Phi)\;=\;0\qquad\text{on }\partial\mathcal M^{(\Lambda)},
\label{eq:BC-operator}
\end{equation}
encoding:
(i) the cutoff surface condition at $\Sigma_\Lambda$ (e.g.\ Dirichlet/mixed boundary condition compatible with the counterterm scheme),
(ii) the KR head induced-gravity Robin conditions,
and (iii) linearized Israel/junction conditions at interior branes.
Healthiness assumptions ensure that the boundary operators have the correct sign (no ghost branch), and that the linearized boundary value
problem is well-posed.

\paragraph{Green-hyperbolicity with timelike boundary.}
Because $\mathcal M^{(\Lambda)}$ has a timelike boundary, we use a boundary-adapted notion.

\begin{definition}[Green-hyperbolic IBVP at finite cutoff]
\label{def:green-hyp-bvp}
Let $P$ be a linear differential operator on $\mathcal M^{(\Lambda)}$ and $B$ a linear boundary operator specifying boundary/junction
conditions.  We say $(P,B)$ is \emph{Green-hyperbolic} on a time slab $\mathscr T$ if there exist continuous linear maps
\[
E^\pm:\ \mathcal D(\mathscr T\times\Sigma_{\rm rec};\text{bundles})\ \longrightarrow\ \Gamma(\mathcal M^{(\Lambda)};\text{bundles})
\]
such that for all compactly supported test sources $f$:
\begin{enumerate}[label=(\roman*),leftmargin=2.0em,itemsep=1pt]
\item $P(E^\pm f)=f$ and $B(E^\pm f)=0$;
\item $E^\pm(P\psi)=\psi$ for all $\psi$ satisfying $B(\psi)=0$ and supported in $\mathscr T$;
\item $\mathrm{supp}(E^\pm f)\subset J^\pm_{\Lambda}(\mathrm{supp} f)$, where $J^\pm_{\Lambda}$ is the causal relation determined by
the bulk metric on $\mathcal M^{(\Lambda)}$ together with the boundary conditions (i.e.\ propagation is retarded/advanced for the IBVP);
\item $E^\pm$ depend continuously on $f$ in the natural LF-topology of test functions.
\end{enumerate}
The causal propagator is $E:=E^- - E^+$.
\end{definition}

\paragraph{Boundary well-posedness route in UOM language.}
The needed PDE input is the standard energy-estimate well-posedness of a symmetric/normal hyperbolic system with maximally dissipative
boundary conditions.  In UOM terms this is exactly the ``healthiness'' package: no ghost/tachyon branches, positive induced-gravity
Robin coefficient on the KR head, and standard holographic counterterms ensuring a well-posed variational principle.

\begin{proposition}[Admissible hyperbolic boundary conditions at finite cutoff]
\label{prop:uom-admissible-BC-route}
Work at fixed finite cutoff $\Lambda<\infty$ on the healthy receiver cell introduced above.
Assume:
\begin{enumerate}[label=\textnormal{(\roman*)},leftmargin=2.2em,itemsep=1pt]
\item the static UV-boundary setup of Paragraph~\ref{par:uom-main-111};
\item the receiver healthiness hypotheses \textnormal{(H1)}--\textnormal{(H4)} of
Theorem~\ref{thm:OS_receiver};
\item the KR/RS1 scalar continuity and positive normal-branch Robin regime of
Paragraph~\ref{par:uom-main-086}.
\end{enumerate}
Then on each bounded time slab $\mathscr T$ the gauge-fixed linearized boundary operator $B_\Lambda$ is admissible for
$P_{\rm gf,ren}$, i.e.\ it is maximally dissipative, equivalently it satisfies the uniform Lopatinski condition.
More precisely, this admissibility is verified by the following three route items:
\begin{enumerate}[label=\textnormal{(R\arabic*)},leftmargin=2.4em,itemsep=1pt]
\item \textbf{Symmetric-hyperbolic reduction.}
After introducing finitely many first derivatives as auxiliary variables, the linearized gauge-fixed system can be written in first-order form
\[
A^0(\bar\Phi)\,\partial_t U + \sum_{j=1}^{4} A^j(\bar\Phi)\,\partial_j U + C(\bar\Phi)\,U = F,
\]
with $A^0$ uniformly positive and $A^j$ symmetric on the finite slab.
\item \textbf{Boundary flux sign.}
The boundary form obtained by integrating the symmetrized energy identity over $\mathscr T\times\Sigma_{\rm rec}$ splits into the
contributions of $\Sigma_\Lambda$, the active KR head, and any interior/shadow branes; each contribution is non-positive on fields satisfying
$B_\Lambda(U)=0$.
\item \textbf{Frozen-coefficient boundary symbol.}
For every boundary point and every tangential Fourier-Laplace frequency with $\Re \zeta>0$, the only incoming decaying solution of the frozen
principal symbol satisfying the homogeneous boundary conditions is the zero solution.
\end{enumerate}
\end{proposition}

\begin{proof}
\emph{Step 1: first-order symmetric-hyperbolic form.}
The principal part of $P_{\rm gf,ren}$ is normally hyperbolic in every sector kept in UOM: de Donder-gauge metric perturbations have principal
part $\Box_g$, while the compensator, GW stabilizer, boundary-distance scalar, and any retained scalar auxiliaries are Klein--Gordon type.
On a bounded slab with static-wall gauge and fixed cutoff $\Lambda$, the coefficients are smooth and time-independent.  Introducing the standard
first derivatives as auxiliary unknowns gives a first-order system with uniformly positive symmetrizer $A^0$ and symmetric principal matrices
$A^j$.  This proves route item~\textnormal{(R1)}.

\emph{Step 2: boundary flux on the physical UOM boundary pieces.}
Multiply the first-order system by the symmetrizer and integrate by parts over the slab.  The boundary contribution splits as
\[
\mathcal F_{\partial\mathcal M^{(\Lambda)}}(U)
= \mathcal F_{\Sigma_\Lambda}(U)+\mathcal F_{\Sigma_{\rm rec}}(U)+\mathcal F_{\Sigma_{\rm int}}(U).
\]
At the outer cutoff surface $\Sigma_\Lambda$, the boundary condition is the one induced by the renormalized variational principle with static UV
boundary and fixed counterterm scheme; by construction it cancels the unwanted variation and gives a non-positive energy flux.
At the active KR head $\Sigma_{\rm rec}$, the admissible scalar/metric boundary data are on the normal induced-gravity branch: the visible Robin
coefficient is positive, the boundary-distance scalar has positive kinetic term, and wrong-sign bending modes are excluded by the static-boundary
choice of Paragraph~\ref{par:uom-main-111}.  Therefore the KR contribution to the symmetrized flux is non-positive, with strict negativity on
incoming physical modes unless the boundary data vanish.  On interior/shadow branes, the linearized Israel/GW conditions are self-adjoint and the
stabilized radion sector has positive kinetic term and nonnegative mass by \textnormal{(H2)}--\textnormal{(H3)}; thus
$\mathcal F_{\Sigma_{\rm int}}(U)\le 0$.  Hence the full boundary form is dissipative:
\[
\mathcal F_{\partial\mathcal M^{(\Lambda)}}(U)\le 0
\qquad\text{whenever }B_\Lambda(U)=0.
\]
This proves route item~\textnormal{(R2)}.

\emph{Step 3: frozen-coefficient boundary symbol and Lopatinski.}
Freeze the coefficients at a boundary point and Fourier-Laplace transform in the tangential variables.  For $\Re\zeta>0$, incoming decaying
modes correspond to roots of the normal symbol with positive energy.  Because \textnormal{(H1)} rules out ghosts, \textnormal{(H2)} rules out
tachyonic boundary or KK branches, and Paragraph~\ref{par:uom-main-086} puts the KR scalar on the positive Robin branch, the homogeneous
boundary conditions admit no nontrivial incoming decaying mode: any such mode would carry strictly negative boundary flux, contradicting the
non-positive dissipative identity of Step~2.  Therefore the homogeneous frozen boundary problem has only the zero solution, which is precisely
the uniform Lopatinski condition on the bounded slab.  This proves route item~\textnormal{(R3)}.

With \textnormal{(R1)}--\textnormal{(R3)} established, standard symmetric-hyperbolic IBVP theory yields well-posedness of the finite-cutoff
linearized problem.  Equivalently, $B_\Lambda$ is maximally dissipative (or satisfies uniform Lopatinski) on each bounded slab, as claimed.
\end{proof}

\begin{theorem}[Green-hyperbolicity and causal propagators at finite cutoff]
\label{thm:uom-P2-green}
Under the hypotheses of Proposition~\ref{prop:uom-admissible-BC-route}, there exist unique advanced/retarded Green operators $E^\pm$ for the IBVP
$(P_{\rm gf,ren},B_\Lambda)$, and the causal propagator $E=E^- - E^+$ is well-defined and antisymmetric with respect to the natural
pairing induced by the quadratic part of $S_{\rm ren}^{\rm gf}$.
\end{theorem}

\begin{proof}
Proposition~\ref{prop:uom-admissible-BC-route} gives the admissibility statement needed for the standard existence/uniqueness theorem for Green operators
of a symmetric-hyperbolic IBVP with timelike boundary: maximal dissipativity / uniform Lopatinski implies well-posedness for the inhomogeneous
initial-boundary value problem and hence existence and uniqueness of retarded/advanced solution operators.
Antisymmetry of $E$ follows from Green's identity for $P_{\rm gf,ren}$ and the boundary conditions, which ensure the boundary pairing
vanishes for admissible solutions.  (A detailed proof reduces to standard energy estimates for symmetric hyperbolic systems; see, e.g.,
the canonical references on hyperbolic IBVPs with timelike boundaries.)
\end{proof}

\paragraph{From bulk causal propagator to boundary retarded kernels.}
Let $\mathcal R$ denote the \emph{boundary restriction / observable extraction} map that sends a bulk perturbation $\delta\Phi$ to the induced
variation of the boundary observable multiplet $\delta\mathcal O_A$ appearing in \eqref{eq:uom-sourced-action}.  Dually, let $\mathcal R^\ast$
be the corresponding source insertion map that sends a boundary source $J$ to the bulk distributional forcing term in the linearized equation.
(Concretely: $\mathcal R^\ast$ inserts a boundary delta distribution on $\Sigma_{\rm rec}$ contracted with the variational derivative of the
coupling term.  The renormalized variational principle guarantees this is well-defined.)

Define the induced boundary retarded/advanced operators by
\begin{equation}
G_R\;:=\;\mathcal R\circ E^+\circ \mathcal R^\ast,
\qquad
G_A\;:=\;\mathcal R\circ E^-\circ \mathcal R^\ast,
\qquad
\Delta\;:=\;G_R-G_A=\mathcal R\circ E\circ \mathcal R^\ast.
\label{eq:GR-from-E}
\end{equation}
By construction, $G_R$ has retarded support in the boundary causal relation and $G_A$ has advanced support.

\begin{proposition}[Identification of the retarded kernel with the SK second derivative]
\label{prop:GR-identification}
Assume UOM-P1 (Lemma~\ref{lem:W-C2}) and Theorem~\ref{thm:uom-P2-green}. Then the distributional kernel
$G_R$ defined by \eqref{eq:GR-from-E} coincides with the mixed second derivative of the renormalized SK influence functional:
\[
G_R^{AB}(x,y)
=
\left.\frac{\delta^2 W_{\rm ren}}{\delta J_{a,A}(x)\,\delta J_{r,B}(y)}\right|_{J=0},
\qquad
G_A^{AB}(x,y)=G_R^{BA}(y,x).
\]
\end{proposition}

\begin{proof}
Differentiate the gauge-fixed sourced equations $\mathcal E(\Phi[J];J)=0$ at $J=0$.  The first variation gives the linear response equation
\[
P_{\rm gf,ren}\,\delta\Phi \;=\; \mathcal R^\ast(J)\qquad\text{with}\qquad B_\Lambda(\delta\Phi)=0,
\]
whose unique retarded solution is $\delta\Phi=E^+\mathcal R^\ast(J)$ by Theorem~\ref{thm:uom-P2-green}.
Applying $\mathcal R$ gives the retarded variation of observables: $\delta\mathcal O = (\mathcal R\circ E^+\circ\mathcal R^\ast)J = G_R J$.

On the other hand, by definition of $W_{\rm ren}$ as the on-shell SK functional (tree level) and by UOM-P1,
the mixed derivative $\delta^2 W_{\rm ren}/(\delta J_a\,\delta J_r)$ evaluated at $J=0$ is exactly the linear response map
from $J_r$ to the expectation of $\mathcal O$ in the retarded prescription (the SK $ra$ kernel). Therefore both definitions coincide,
and $G_A$ is obtained by time reversal (swap of arguments), yielding $G_A^{AB}(x,y)=G_R^{BA}(y,x)$.
\end{proof}

\paragraph{Conclusion of UOM-P2.}
Theorem~\ref{thm:uom-P2-green} and Proposition~\ref{prop:GR-identification} establish \textbf{UOM-P2}: advanced/retarded propagators exist
for the gauge-fixed linearized renormalized IBVP at finite cutoff, and the induced boundary retarded kernel is the mixed second derivative of
$W_{\rm ren}$.
% ============================================================
% (4) UOM-P3 subsubsection snippet
% ============================================================

\subsubsection{UOM-P3: renormalized covariant symplectic structure and compatibility with the causal propagator}
\label{subsubsec:uom-P3}

This subsubsection establishes \textbf{UOM-P3}.  We construct the renormalized covariant presymplectic potential and current,
prove finiteness and conservation under the UOM counterterm scheme and boundary conditions, and show compatibility with the causal propagator.

\paragraph{Renormalized variational identity.}
Let $L_{\rm ren}$ denote the renormalized Lagrangian density (including boundary terms) corresponding to $S_{\rm ren}^{\rm gf}$
on $\mathcal M^{(\Lambda)}$.  For a compactly supported variation $\delta\Phi$ one has the standard decomposition
\begin{equation}
\delta L_{\rm ren}(\Phi)
\;=
E_{\rm ren}(\Phi)\cdot \delta\Phi \; +\; d\,\theta_{\rm ren}(\Phi;\delta\Phi),
\label{eq:var-identity}
\end{equation}
where $E_{\rm ren}(\Phi)=0$ are the renormalized Euler--Lagrange equations and $\theta_{\rm ren}$ is the (renormalized) presymplectic potential
current (a spacetime $(d-1)$-form valued in field variations).  The existence of $\theta_{\rm ren}$ with finite coefficients is a direct
consequence of the fixed counterterm scheme and the well-posed renormalized variational principle (no leftover boundary variations).

\paragraph{Renormalized symplectic current and form.}
Define the renormalized symplectic current
\begin{equation}
\omega_{\rm ren}(\Phi;\delta_1\Phi,\delta_2\Phi)
\;:=
\delta_1\theta_{\rm ren}(\Phi;\delta_2\Phi)\; -\; \delta_2\theta_{\rm ren}(\Phi;\delta_1\Phi).
\label{eq:omega-ren}
\end{equation}
Let $\Sigma$ be any smooth Cauchy surface of the time slab in $\mathcal M^{(\Lambda)}$ (in the sense of globally hyperbolic manifolds with
timelike boundary).  Define the renormalized presymplectic form by
\begin{equation}
\Omega_{\rm ren}(\delta_1\Phi,\delta_2\Phi)
\;:=\;\int_{\Sigma}\omega_{\rm ren}(\bar\Phi;\delta_1\Phi,\delta_2\Phi),
\label{eq:Omega-ren}
\end{equation}
evaluated at the background $\bar\Phi$ (or more generally at any solution $\Phi$; we work near $\bar\Phi$).

\begin{lemma}[On-shell conservation of $\omega_{\rm ren}$]
\label{lem:omega-closed}
Let $\Phi$ solve the renormalized equations $E_{\rm ren}(\Phi)=0$ and let $\delta_1\Phi,\delta_2\Phi$ solve the linearized equations
$P_{\rm gf,ren}\delta\Phi=0$ with boundary conditions $B_\Lambda(\delta\Phi)=0$. Then
\[
d\,\omega_{\rm ren}(\Phi;\delta_1\Phi,\delta_2\Phi)=0
\]
as a differential form on spacetime.
\end{lemma}

\begin{proof}
Apply $\delta_1$ and $\delta_2$ to the variational identity \eqref{eq:var-identity}, antisymmetrize, and use that mixed partials commute.
The bulk terms proportional to $E_{\rm ren}(\Phi)$ and the linearized equations vanish on-shell, leaving $d\omega_{\rm ren}=0$.
This is the standard covariant phase space identity.
\end{proof}

\begin{lemma}[Finiteness and hypersurface independence at finite cutoff]
\label{lem:Omega-finite-independent}
At fixed finite cutoff $\Lambda<\infty$ and under the UOM boundary conditions $B_\Lambda=0$, the integral \eqref{eq:Omega-ren} is finite.
Moreover, if $\Sigma$ and $\Sigma'$ are two Cauchy surfaces in the same time slab, then
\[
\int_{\Sigma}\omega_{\rm ren}=\int_{\Sigma'}\omega_{\rm ren},
\]
so $\Omega_{\rm ren}$ is independent of the choice of Cauchy surface.
\end{lemma}

\begin{proof}
Finiteness: at finite cutoff $\Lambda$ the fields and their admissible perturbations satisfy the UOM regulated boundary conditions and falloff,
so all potentially divergent near-boundary terms are controlled. The holographic counterterms are fixed precisely so that the renormalized action
has a finite first variation; therefore $\theta_{\rm ren}$ and $\omega_{\rm ren}$ are finite on admissible perturbations, and the integral over the
compact spatial slice $\Sigma\cap\Sigma_{\rm rec}\cong S^3$ is finite.

Hypersurface independence: let $\mathcal U$ be the spacetime region bounded by $\Sigma$, $\Sigma'$, and timelike boundary pieces.
By Stokes' theorem and Lemma~\ref{lem:omega-closed},
\[
\int_{\Sigma}\omega_{\rm ren}-\int_{\Sigma'}\omega_{\rm ren}
=\int_{\partial\mathcal U\setminus(\Sigma\cup\Sigma')}\omega_{\rm ren}.
\]
The UOM boundary/junction conditions are those arising from the renormalized variational principle; they ensure that the symplectic flux through
the timelike boundary vanishes on admissible perturbations (equivalently: boundary terms cancel between bulk and counterterms, and Robin/Israel
conditions enforce zero symplectic leakage). Hence the RHS vanishes and the two integrals agree.
\end{proof}

\paragraph{Compatibility with the causal propagator.}
Let $(P_{\rm gf,ren},B_\Lambda)$ be the linearized IBVP and let $E^\pm$ be its Green operators (UOM-P2).
Let $\langle\cdot,\cdot\rangle$ denote the natural pairing between test sources $f$ and field perturbations $\delta\Phi$
induced by the quadratic part of $S_{\rm ren}^{\rm gf}$:
\[
\langle f,\delta\Phi\rangle \;:=\; \int_{\mathcal M^{(\Lambda)}} f\cdot \delta\Phi.
\]

\begin{lemma}[Green identity and inverse relation]
\label{lem:Omega-inverts-E}
Let $f$ be a compactly supported test source and set $\delta\Phi:=E f$ (the causal solution).  Let $\psi$ be any solution of the homogeneous
linearized equations with $B_\Lambda(\psi)=0$. Then
\begin{equation}
\Omega_{\rm ren}(\delta\Phi,\psi)\;=\;\langle f,\psi\rangle.
\label{eq:Omega-inverse}
\end{equation}
\end{lemma}

\begin{proof}
This is the standard Green identity in covariant phase space form.
Consider the current
\[
j^\mu(\delta\Phi,\psi):=\omega_{\rm ren}^\mu(\bar\Phi;\delta\Phi,\psi).
\]
Using the linearized equations $P_{\rm gf,ren}\psi=0$ and $P_{\rm gf,ren}\delta\Phi=f$ together with the definition of $\omega_{\rm ren}$,
one checks (by direct computation from the quadratic Lagrangian) that
\[
\partial_\mu j^\mu(\delta\Phi,\psi)= f\cdot \psi
\]
up to terms that vanish by the boundary conditions. Integrate over the spacetime region between two Cauchy surfaces and apply Stokes' theorem.
Use the vanishing symplectic flux through timelike boundaries (as in Lemma~\ref{lem:Omega-finite-independent}) and take the limit where the
time slab contains the support of $f$ to obtain \eqref{eq:Omega-inverse}.
\end{proof}

\begin{proposition}[Induced boundary causal bivector]
\label{prop:boundary-bivector}
Let $G_R,G_A,\Delta$ be the induced boundary operators defined in \eqref{eq:GR-from-E}.  Then $\Delta$ is antisymmetric with respect to the
boundary pairing and is the induced causal bivector compatible with $\Omega_{\rm ren}$ on gauge-invariant boundary data.
\end{proposition}

\begin{proof}
Antisymmetry follows from antisymmetry of $E$ and the definition \eqref{eq:GR-from-E}.
Compatibility is the boundary restriction of \eqref{eq:Omega-inverse}: composing the bulk identity with the source insertion and observable
extraction maps $(\mathcal R^\ast,\mathcal R)$ yields the corresponding boundary inverse relation between $\Delta$ and the boundary-restricted
symplectic pairing induced by $\Omega_{\rm ren}$.  (Gauge invariance and reduction are treated in UOM-P4 below.)
\end{proof}

\paragraph{Conclusion of UOM-P3.}
Lemmas~\ref{lem:omega-closed}--\ref{lem:Omega-inverts-E} establish the existence, finiteness, and conservation of the renormalized presymplectic
form $\Omega_{\rm ren}$ at finite cutoff and its inverse relation with the causal propagator.  Proposition~\ref{prop:boundary-bivector} packages
the induced boundary causal bivector $\Delta=G_R-G_A$.
% ============================================================
% (5) UOM-P4 subsubsection snippet (no BV/fermions)
% ============================================================

\subsubsection{UOM-P4: gauge independence on gauge invariants via presymplectic reduction (no BV/ghosts)}
\label{subsubsec:uom-P4}

This subsubsection establishes \textbf{UOM-P4} without introducing any odd structures.  Gauge independence is proved by:
(i) identifying pure gauge directions with the degeneracy (kernel) of the renormalized presymplectic form $\Omega_{\rm ren}$, and
(ii) defining the bracket on the reduced phase space of solutions modulo pure gauge.  Any gauge-fixed Peierls computation is then shown to
agree with the reduced bracket on gauge-invariant functionals.

\paragraph{Pure gauge transformations at finite cutoff.}
Let $\mathcal G$ denote the full gauge group of the bosonic bulk+brane system (e.g.\ bulk diffeomorphisms and any internal gauge symmetries).
We isolate the subgroup of \emph{pure gauge} transformations at fixed cutoff $\Lambda$:

\begin{definition}[Pure gauge subgroup $\mathcal G_0$]
\label{def:G0}
Let $\mathcal G_0\subset\mathcal G$ be the subgroup of gauge transformations that:
\begin{enumerate}[label=(\roman*),leftmargin=2.0em,itemsep=1pt]
\item preserve the UOM boundary/junction conditions $B_\Lambda=0$ (including the KR/RS boundary operators),
\item act trivially on the chosen boundary data used to define observables and sources (so they do not generate physical boundary charges),
\item and are compactly supported in the time slab (or decay sufficiently fast at the temporal endpoints) so that no boundary term at infinity arises.
\end{enumerate}
\end{definition}

\begin{remark}
The distinction between $\mathcal G_0$ and the full asymptotic symmetry group is essential in AdS: physical boundary symmetries (e.g.\ the
$\mathrm{Conf}(S^3)$ action used elsewhere in UOM) are \emph{not} quotiented as gauge.  Definition~\ref{def:G0} removes only redundancies.
\end{remark}

\paragraph{Gauge degeneracy of the renormalized presymplectic form.}
Let $\delta_\epsilon\Phi$ denote the infinitesimal gauge variation of fields generated by an infinitesimal parameter $\epsilon$ (e.g.\ a vector field
$\xi^A$ for diffeomorphisms) that is tangent to $\mathcal G_0$.

\begin{lemma}[Pure gauge directions lie in the kernel of $\Omega_{\rm ren}$]
\label{lem:gauge-kernel}
Let $\Phi$ be a solution of the renormalized equations and let $\delta\Phi$ be any linearized solution with boundary conditions $B_\Lambda(\delta\Phi)=0$.
If $\delta_\epsilon\Phi$ is an infinitesimal pure gauge variation with $\epsilon\in \mathrm{Lie}(\mathcal G_0)$, then
\[
\Omega_{\rm ren}(\delta_\epsilon\Phi,\delta\Phi)\;=\;0.
\]
\end{lemma}

\begin{proof}
Gauge invariance of the renormalized action (including counterterms) implies the Noether identity for the presymplectic potential:
$\theta_{\rm ren}(\Phi;\delta_\epsilon\Phi)$ differs from an exact form by a term proportional to the constraints.
Taking an antisymmetrized second variation yields that $\omega_{\rm ren}(\Phi;\delta_\epsilon\Phi,\delta\Phi)$ is exact on-shell.
Integrating over a Cauchy surface gives a boundary term.  By Definition~\ref{def:G0}(ii), pure gauge parameters generate no physical
boundary charges and vanish in the boundary pairing; by Definition~\ref{def:G0}(i) they preserve $B_\Lambda=0$ so that the symplectic flux
through timelike boundaries vanishes. Hence the boundary term is zero and $\Omega_{\rm ren}(\delta_\epsilon\Phi,\delta\Phi)=0$.
\end{proof}

\begin{proposition}[Fixed-cutoff coercive/Fredholm route for the UOM-P4 kernel clause]
\label{prop:uom-no-extra-deg-route}
Work on the same healthy receiver cell at fixed finite cutoff $\Lambda<\infty$ and bounded time slab $\mathscr T$.
Assume:
\begin{enumerate}[label=\textnormal{(\roman*)},leftmargin=2.2em,itemsep=1pt]
\item the static UV-boundary setup of Paragraph~\ref{par:uom-main-111};
\item the healthiness hypotheses \textnormal{(H1)}--\textnormal{(H3)} of Theorem~\ref{thm:OS_receiver};
\item the KR/RS1 scalar continuity and positive normal-branch Robin regime of Paragraph~\ref{par:uom-main-086};
\item the admissible finite-cutoff boundary problem of Proposition~\ref{prop:uom-admissible-BC-route}.
\end{enumerate}
Let $\mathcal T_{\mathrm{phys}}$ be any closed complement of the pure gauge directions
\(
\{\delta_\epsilon\Phi:\epsilon\in\mathrm{Lie}(\mathcal G_0)\}
\)
inside the space of admissible linearized perturbations.  Then the statement that pure gauge exhausts the kernel of
$\Omega_{\rm ren}$ reduces to the following three checks on $\mathcal T_{\mathrm{phys}}$:
\begin{enumerate}[label=\textnormal{(F\arabic*)},leftmargin=2.4em,itemsep=1pt]
\item \textbf{Fredholm realization.}
The renormalized second variation restricted to $\mathcal T_{\mathrm{phys}}$ defines a self-adjoint Fredholm operator
$\mathbb H_{\mathrm{phys}}$ on the natural finite-cutoff energy space.
\item \textbf{Coercive canonical quadratic form.}
Its associated canonical quadratic form $\mathcal E_{\mathrm{can}}(v):=\langle v,\mathbb H_{\mathrm{phys}}v\rangle$ is coercive modulo compact
lower-order terms on $\mathcal T_{\mathrm{phys}}$.
\item \textbf{Trivial reduced kernel.}
The kernel of $\mathbb H_{\mathrm{phys}}$ is trivial.
\end{enumerate}
If \textnormal{(F1)}--\textnormal{(F3)} hold, then
\[
\ker \Omega_{\rm ren}\;=\;\{\delta_\epsilon\Phi:\ \epsilon\in \mathrm{Lie}(\mathcal G_0)\}.
\]
\end{proposition}

\begin{proof}
On $\mathcal T_{\mathrm{phys}}$ the admissible linearized equations define the reduced finite-cutoff physical system.  By
Lemma~\ref{lem:Omega-finite-independent}, $\Omega_{\rm ren}$ is the time-independent canonical presymplectic form of that system, while
Lemma~\ref{lem:gauge-kernel} removes the pure gauge directions.  Any residual null vector of $\Omega_{\rm ren}$ therefore gives a reduced
physical zero mode.  In the presence of \textnormal{(F1)} and \textnormal{(F2)}, Fredholm theory reduces such zero modes to the finite-dimensional
kernel of $\mathbb H_{\mathrm{phys}}$; if \textnormal{(F3)} holds, that kernel is zero, so no residual null direction remains after quotienting
gauge.  This yields the displayed identity.
\end{proof}

\begin{theorem}[No additional degeneracies after gauge quotient at fixed cutoff (UOM-P4 kernel clause)]
\label{thm:uom-no-extra-deg}
Under the hypotheses of Proposition~\ref{prop:uom-no-extra-deg-route},
\[
\ker \Omega_{\rm ren}\;=\;\{\delta_\epsilon\Phi:\ \epsilon\in \mathrm{Lie}(\mathcal G_0)\}.
\]
\end{theorem}

\begin{proof}
We verify the three route items of Proposition~\ref{prop:uom-no-extra-deg-route}.

\emph{Step 1: Fredholm realization on the reduced physical slice.}
Choose a closed complement $\mathcal T_{\mathrm{phys}}$ to the pure gauge directions.  Because
Proposition~\ref{prop:uom-admissible-BC-route} gives maximally dissipative / uniform-Lopatinski boundary conditions on the finite slab,
the gauge-fixed linearized operator with those boundary conditions is a well-posed symmetric-hyperbolic system.  Restricting to
$\mathcal T_{\mathrm{phys}}$ removes the gauge kernel identified in Lemma~\ref{lem:gauge-kernel}.  On a bounded time slab with compact spatial
slice and finite cutoff, the corresponding reduced graph embedding is compact, so the reduced static operator obtained from the renormalized
second variation is self-adjoint Fredholm.  This proves \textnormal{(F1)}.

\emph{Step 2: Coercivity of the reduced canonical quadratic form.}
The renormalized second variation on $\mathcal T_{\mathrm{phys}}$ is the sum of the bulk kinetic terms, the brane kinetic terms, and the
Robin/Israel boundary contributions from the fixed counterterm scheme.  By \textnormal{(H1)}, all physical kinetic terms have the correct sign;
by \textnormal{(H2)}, the KK/boundary Sturm--Liouville operators are self-adjoint with nonnegative masses; by \textnormal{(H3)}, the stabilized
radion sector has positive kinetic term; and by Paragraph~\ref{par:uom-main-086}, the KR boundary-distance scalar remains on the positive
normal Robin branch.  The counterterm contributions are lower order and bounded on the finite-cutoff energy space.  Therefore the associated
canonical quadratic form is positive modulo compact lower-order corrections, i.e.\ it is coercive in the standard G{\aa}rding/Fredholm sense
on $\mathcal T_{\mathrm{phys}}$.  This proves \textnormal{(F2)}.

\emph{Step 3: Exclusion of non-gauge zero modes.}
Let $v\in\ker \mathbb H_{\mathrm{phys}}$.  Then $\mathcal E_{\mathrm{can}}(v)=0$.  By the coercive estimate of Step~2, any such $v$ lies in the
finite-dimensional low-lying sector singled out by the Fredholm alternative.  On that sector, the possible zero modes are exactly the physical
KK/radion/boundary-distance modes compatible with the finite-cutoff boundary problem.  These are excluded one by one by the healthiness inputs:
ghost directions are absent by \textnormal{(H1)}, tachyonic or marginal boundary/KK branches are absent by \textnormal{(H2)}, the interior RS1
radion is stabilized with strictly positive kinetic/mass data by \textnormal{(H3)}, and the KR boundary-distance scalar sits on the positive
Robin branch of Paragraph~\ref{par:uom-main-086}.  Finally, the finite cutoff removes continuum-edge degeneracies, so there is no additional
zero-frequency accumulation at the boundary.  Hence $v=0$, proving \textnormal{(F3)}.

Proposition~\ref{prop:uom-no-extra-deg-route} now yields the claimed identity for $\ker \Omega_{\rm ren}$,
which is exactly the kernel statement needed in \textbf{UOM-P4}.
\end{proof}

\paragraph{Reduced phase space and canonical bracket.}
Let $\mathcal S$ be the space of solutions of the renormalized equations on the time slab with boundary conditions $B_\Lambda=0$.
Define the reduced phase space
\[
\mathcal P\;:=\;\mathcal S/\mathcal G_0.
\]

\begin{proposition}[Descent of $\Omega_{\rm ren}$]
\label{prop:Omega-descends}
Under Lemma~\ref{lem:gauge-kernel} and Theorem~\ref{thm:uom-no-extra-deg}, the presymplectic form $\Omega_{\rm ren}$
descends to a \emph{nondegenerate} symplectic form $\overline\Omega_{\rm ren}$ on $\mathcal P$.
\end{proposition}

\begin{proof}
Lemma~\ref{lem:gauge-kernel} implies that $\Omega_{\rm ren}$ is constant along $\mathcal G_0$-orbits, hence descends to the quotient.
Theorem~\ref{thm:uom-no-extra-deg} implies the only degeneracies are precisely the directions removed by the quotient, so the descended form is
nondegenerate.
\end{proof}

\paragraph{Gauge-invariant functionals and Hamiltonian vector fields.}
Let $F$ be a functional on $\mathcal S$ that is invariant under $\mathcal G_0$.

\begin{lemma}[Infinitesimal invariance]
\label{lem:inf-inv}
If $F$ is $\mathcal G_0$-invariant, then for any $\epsilon\in \mathrm{Lie}(\mathcal G_0)$ and any solution $\Phi\in\mathcal S$,
\[
dF(\Phi)\cdot \delta_\epsilon\Phi\;=\;0.
\]
\end{lemma}

\begin{proof}
Differentiate the identity $F(\Phi)=F(\Phi\cdot g)$ at $g=\exp(\epsilon)$ and use the chain rule.
\end{proof}

Because $\overline\Omega_{\rm ren}$ is nondegenerate on $\mathcal P$, each gauge-invariant $F$ induces a unique Hamiltonian vector field
$X_F$ on $\mathcal P$ defined by $\iota_{X_F}\overline\Omega_{\rm ren}=dF$.

\begin{definition}[Reduced covariant bracket]
\label{def:reduced-bracket}
For gauge-invariant $F,G$ define
\[
\{F,G\}_\Omega \;:=\;\overline\Omega_{\rm ren}(X_F,X_G).
\]
\end{definition}

\paragraph{Agreement with gauge-fixed Peierls on gauge invariants.}
Let $P_{\rm gf,ren}$ be any admissible even gauge fixing (as in \eqref{eq:even-gauge-fixing}) producing a Green-hyperbolic IBVP, with causal
propagator $E$ and induced boundary causal operator $\Delta$ (UOM-P2).  Define the gauge-fixed Peierls bracket
\[
\{F,G\}_{{\rm P},{\rm gf}}\;:=\;\langle F^{(1)},\,\Delta\,G^{(1)}\rangle,
\]
whenever the pairing is defined.

\begin{theorem}[Gauge independence on invariants (no BV/ghosts)]
\label{thm:gauge-independence}
Assume UOM-P2 and UOM-P3, Proposition~\ref{prop:Omega-descends}, and that $F,G$ are $\mathcal G_0$-invariant.
Then:
\begin{enumerate}[label=\textnormal{(\roman*)},leftmargin=2.0em,itemsep=1pt]
\item The value of $\{F,G\}_{{\rm P},{\rm gf}}$ is independent of the admissible choice of even gauge fixing used to define $\Delta$.
\item Moreover, the gauge-fixed Peierls bracket equals the reduced covariant bracket:
\[
\{F,G\}_{{\rm P},{\rm gf}} \;=\; \{F,G\}_\Omega.
\]
\end{enumerate}
\end{theorem}

\begin{proof}
By UOM-P3, the causal propagator $E$ is the inverse bivector of $\Omega_{\rm ren}$ on the space of admissible solutions modulo gauge.
Concretely, changing the gauge fixing changes the propagator by an operator whose image lies in pure gauge directions (a standard fact for
gauge-fixed Green operators: the difference of two right inverses of the same gauge-invariant operator is a map into the homogeneous gauge
subspace).  Pairing with $F^{(1)}$ kills such contributions by Lemma~\ref{lem:inf-inv}, so $\{F,G\}_{{\rm P},{\rm gf}}$ is gauge independent.

For the equality with $\{F,G\}_\Omega$, note that for gauge-invariant $F$ the corresponding Hamiltonian vector field on the reduced space is
represented (in any gauge fixing) by the unique solution of the linearized equations sourced by $F^{(1)}$, i.e.\ by $E(F^{(1)})$ modulo gauge.
Then
\[
\{F,G\}_\Omega=\overline\Omega_{\rm ren}(X_F,X_G)
=\Omega_{\rm ren}(E F^{(1)},\,E G^{(1)})
=\langle F^{(1)},\,E\,G^{(1)}\rangle
=\langle F^{(1)},\,\Delta\,G^{(1)}\rangle,
\]
where the middle equality uses Lemma~\ref{lem:Omega-inverts-E} (UOM-P3) and the last equality uses the standard conversion between bulk causal
propagator $E$ and the induced boundary causal operator $\Delta$ (UOM-P2).  This yields $\{F,G\}_{{\rm P},{\rm gf}}=\{F,G\}_\Omega$.
\end{proof}

\paragraph{Conclusion of UOM-P4.}
Theorem~\ref{thm:gauge-independence} establishes \textbf{UOM-P4} in a form suitable for import into the two--time program:
the Peierls bracket is canonical and gauge independent on gauge-invariant observables, proved using only even (bosonic) presymplectic reduction.
% ============================================================
% (6) UOM-P5 subsubsection snippet
% ============================================================

\subsubsection{UOM-P5: sufficiently local gauge-invariant functionals and closure under the Peierls bracket}
\label{subsubsec:uom-P5}

This subsubsection establishes \textbf{UOM-P5}: a precisely defined class of sufficiently local gauge-invariant functionals at finite cutoff
$\Lambda$ that is closed under the Peierls bracket and contains the quasi-local observables used in Subprogram~3.

\paragraph{Why finite cutoff simplifies locality.}
At finite receiver cutoff $\Lambda<\infty$ the admissible receiver boundary data are effectively band-limited: only finitely many (or a
rapidly decaying) set of boundary harmonics contribute.  In this regime, all fields restricted to the receiver head are smooth functions of
spacetime, products of local densities are well-defined, and distributional product subtleties can be avoided.
We therefore define $\mathcal F_{\rm s.loc}^\Lambda$ by an operationally natural ``band-limited microcausal'' class.

\paragraph{Band-limited boundary field space.}
Let $\varphi$ denote the collection of boundary (receiver-head) field variables that coordinatize the admissible boundary data at cutoff $\Lambda$
(e.g.\ the induced metric perturbations modulo pure gauge at the boundary, the compensator/trace scalar, and any auxiliary scalars retained).
Let $P_\Lambda$ denote the spectral projector implementing the receiver cutoff (e.g.\ $\ell\le L(\Lambda)$ on $S^3$ modes and the corresponding
temporal band window on $\mathscr T$).  Define the Banach space
\begin{equation}
\mathcal X_\Lambda^s\;:=\;P_\Lambda\,H^s(\mathscr T\times S^3;\mathbb R^N),
\label{eq:XLambda}
\end{equation}
for $s>\frac{5}{2}$ and some finite $N$ (number of boundary components).  Elements of $\mathcal X_\Lambda^s$ are $C^1$ and are smooth in the
band-limited sense.

\paragraph{Sufficiently local functionals.}
We now define the admissible functional class.

\begin{definition}[Sufficiently local gauge-invariant functionals at cutoff $\Lambda$]
\label{def:Floc}
Let $\mathcal F_{\rm s.loc}^\Lambda$ be the set of functionals $F:\mathcal X_\Lambda^s\to\mathbb R$ of the form
\begin{equation}
F(\varphi)\;=\;\sum_{n=0}^{N_F}\ \int_{\mathscr T\times S^3}\! d^4x\ \sqrt{-\gamma}\;
f_n(x)\,P_n\big(\varphi(x),\partial\varphi(x),\ldots,\partial^{k_F}\varphi(x)\big),
\label{eq:Floc-form}
\end{equation}
where:
\begin{enumerate}[label=(\roman*),leftmargin=2.0em,itemsep=1pt]
\item each $f_n\in C_c^\infty(\mathscr T\times S^3)$ is a smooth compactly supported test function (compact support in time on $S^3$),
\item each $P_n$ is a polynomial in the displayed jet variables up to finite order $k_F$,
\item and $F$ is invariant under the pure gauge subgroup $\mathcal G_0$ (Definition~\ref{def:G0}) acting on $\varphi$.
\end{enumerate}
\end{definition}

\begin{remark}
The class \eqref{eq:Floc-form} is the finite-cutoff analogue of the standard microcausal functionals of perturbative AQFT; at finite cutoff,
the wavefront-set conditions are automatic because $P_\Lambda$ produces smooth functions with controlled spectra.  A continuum extension can be
made by imposing the usual microcausal wavefront-set restriction on the functional derivatives; we do not require it here.
\end{remark}

\paragraph{Functional derivatives and the Peierls pairing.}
For $F\in\mathcal F_{\rm s.loc}^\Lambda$, the first Fr\'echet derivative exists and is given by a compactly supported $H^{s-k_F}$ function,
hence defines a continuous linear functional on $\mathcal X_\Lambda^s$.  We write $F^{(1)}(\varphi)$ for the density-valued derivative and use
the pairing
\[
\langle F^{(1)},\psi\rangle \;:=\;\int_{\mathscr T\times S^3}\! d^4x\ \sqrt{-\gamma}\; F^{(1)}_A(x;\varphi)\,\psi^A(x),
\]
with $\psi\in\mathcal X_\Lambda^s$.

Let $\Delta_\Lambda$ denote the induced boundary causal operator (UOM-P2), viewed as a continuous map on $\mathcal X_\Lambda^s$.
Define the Peierls bracket on $\mathcal F_{\rm s.loc}^\Lambda$ by
\begin{equation}
\{F,G\}_{\rm P}
\;:=
\big\langle F^{(1)},\ \Delta_\Lambda\,G^{(1)}\big\rangle,
\label{eq:Peierls-FG}
\end{equation}
whenever the pairing is defined (it will be, by the next lemma).

\begin{lemma}[Well-definedness on $\mathcal F_{\rm s.loc}^\Lambda$]
\label{lem:Peierls-well-defined}
Assume UOM-P2 (existence of $\Delta_\Lambda$ as a continuous causal operator on $\mathcal X_\Lambda^s$).
Then for all $F,G\in\mathcal F_{\rm s.loc}^\Lambda$, the bracket \eqref{eq:Peierls-FG} is finite and well-defined.
\end{lemma}

\begin{proof}
For $F,G$ of the form \eqref{eq:Floc-form}, their first derivatives $F^{(1)},G^{(1)}$ are smooth compactly supported functions on
$\mathscr T\times S^3$ (in the band-limited topology).  By UOM-P2, $\Delta_\Lambda$ maps such functions continuously into $\mathcal X_\Lambda^s$.
Therefore the pairing integral defining \eqref{eq:Peierls-FG} is absolutely convergent and well-defined.
\end{proof}

\begin{lemma}[Closure]
\label{lem:closure}
Assume UOM-P2 and that $\Delta_\Lambda$ is a bidistribution with causal support on $\mathscr T\times S^3$ whose action preserves the band-limited
space $\mathcal X_\Lambda^s$.  Then for all $F,G\in\mathcal F_{\rm s.loc}^\Lambda$, the bracket $\{F,G\}_{\rm P}$ is again an element of
$\mathcal F_{\rm s.loc}^\Lambda$.
\end{lemma}

\begin{proof}
Because $F^{(1)}$ and $G^{(1)}$ are local polynomials in the jet of $\varphi$ with compact support weights, $\Delta_\Lambda G^{(1)}$ is a smooth
band-limited field with support contained in the causal hull of $\mathrm{supp}(G^{(1)})$; pairing with $F^{(1)}$ yields a finite sum of integrals
of compactly supported smooth functions depending polynomially on a finite jet of $\varphi$, hence again of the form \eqref{eq:Floc-form}.
Gauge invariance is preserved because $\Delta_\Lambda$ is defined on the gauge-reduced data (UOM-P4) and because $F,G$ are $\mathcal G_0$-invariant.
\end{proof}

\paragraph{Poisson properties.}
Antisymmetry follows from antisymmetry of $\Delta_\Lambda$ (UOM-P2 / UOM-P3).  Leibniz and Jacobi follow from the symplectic origin
(UOM-P3) and gauge reduction (UOM-P4).  We record this as a theorem.

\begin{theorem}[Peierls bracket on $\mathcal F_{\rm s.loc}^\Lambda$]
\label{thm:Peierls-P5}
Assume UOM-P2--P4.  Then the bracket \eqref{eq:Peierls-FG} is:
\begin{enumerate}[label=\textnormal{(\roman*)},leftmargin=2.0em,itemsep=1pt]
\item bilinear and antisymmetric on $\mathcal F_{\rm s.loc}^\Lambda$;
\item a derivation in each argument (Leibniz rule);
\item satisfies the Jacobi identity on $\mathcal F_{\rm s.loc}^\Lambda$;
\item gauge independent on $\mathcal G_0$-invariants (Theorem~\ref{thm:gauge-independence}).
\end{enumerate}
\end{theorem}

\begin{proof}
Antisymmetry is immediate from $\Delta_\Lambda=-\Delta_\Lambda^\top$.
Leibniz and Jacobi follow from the fact that on the reduced phase space $\mathcal P$ the bracket coincides with the symplectic bracket
$\{F,G\}_\Omega=\overline\Omega_{\rm ren}(X_F,X_G)$ (Theorem~\ref{thm:gauge-independence}), and symplectic brackets satisfy Leibniz and Jacobi.
Closure is Lemma~\ref{lem:closure}.
\end{proof}

\paragraph{Containment of RG-monotone representatives (Subprogram~3 compatibility).}
The RG monotone observables used later (e.g.\ Wess--Zumino/anomaly floor densities smeared by compact kernels on $\Sigma_{\rm rec}$) are of the
form \eqref{eq:Floc-form}: they are integrals of local scalar densities built from the (band-limited) trace/Weyl sector and a compact test
kernel.  Hence they lie in $\mathcal F_{\rm s.loc}^\Lambda$.

\begin{proposition}[RG-monotone observables are in $\mathcal F_{\rm s.loc}^\Lambda$]
\label{prop:RG-observable-in-Floc}
Let $\mathcal O^{(0)}$ be any quasi-local observable constructed as a compact smearing of a local scalar density on the receiver head
(e.g.\ the Wess--Zumino density of the accumulated Weyl/trace profile, or a finite linear combination of such smearings).
Then $\mathcal O^{(0)}\in\mathcal F_{\rm s.loc}^\Lambda$.
\end{proposition}

\begin{proof}
By construction $\mathcal O^{(0)}$ is an integral of a local polynomial density in the relevant boundary fields and finitely many derivatives,
with a smooth compactly supported kernel; this is exactly \eqref{eq:Floc-form}.  Gauge invariance holds because the densities are built from
gauge-invariant curvature/Weyl scalars on the boundary and the pure gauge group $\mathcal G_0$ acts trivially on the chosen boundary data.
\end{proof}

\paragraph{Conclusion of UOM-P5.}
Definitions~\ref{def:Floc} and \eqref{eq:Peierls-FG} together with Lemmas~\ref{lem:Peierls-well-defined}--\ref{lem:closure} and
Theorem~\ref{thm:Peierls-P5} establish \textbf{UOM-P5}: a closed sufficiently local gauge-invariant functional class at finite cutoff $\Lambda$
supporting a canonical Peierls bracket, suitable for import into the two--time Hamiltonization step.

\subsubsection{UOM-P6: finite-cutoff pointer-real swap equivariance}
\label{subsubsec:uom-P6}

This subsubsection establishes \textbf{UOM-P6}: the exact finite-cutoff pointer-real involution behind the boundary swap action, its
commutation with the gauge-fixed linearized renormalized IBVP, invariance of the Peierls pairing, and the resulting Poisson automorphism
property on the sufficiently local boundary class.

\begin{definition}[Finite-cutoff pointer-real involution]
\label{def:uom-gamma-star}
Fix the same finite cutoff $\Lambda$ and bandlimit $\Delta(\Lambda)$ as above.  Using the pointer basis supplied by
Definition~\ref{def:uom-limited-scope-package}\textnormal{\textbf{(A7)}} and Definition~\ref{def:pointer-algebra-projection}, define
\[
\Gamma_{*,\Lambda}:\mathcal X_\Lambda^s\to\mathcal X_\Lambda^s
\]
to be the transpose / real involution in that pointer basis, applied componentwise to the finite family of band-limited bosonic
boundary-data components.  We use the same notation for the induced involution on admissible linearized bulk+brane perturbations,
their boundary traces, and the admissible source space
\(
\mathcal J^s(\mathscr X)
\),
with the same componentwise transpose / real action on each finite bosonic field component, and define the pullback involution on
sufficiently local functionals by
\[
(\Gamma F)(\varphi):=F(\Gamma_{*,\Lambda}\varphi)
\qquad
(F\in\mathcal F_{\rm s.loc}^\Lambda).
\]
\end{definition}

\begin{theorem}[Finite-cutoff pointer-real equivariance of the linearized UOM IBVP]
\label{thm:uom-pointer-real-ibvp}
For any admissible even gauge fixing as in \eqref{eq:even-gauge-fixing}, the gauge-fixed linearized renormalized operator, the
boundary/junction operator, the source insertion map, and the boundary observable-extraction map are $\Gamma_{*,\Lambda}$--equivariant:
\[
P_{\rm gf,ren}\Gamma_{*,\Lambda}=\Gamma_{*,\Lambda}P_{\rm gf,ren},
\qquad
B_\Lambda\Gamma_{*,\Lambda}=\Gamma_{*,\Lambda}B_\Lambda,
\]
\[
\mathcal R\Gamma_{*,\Lambda}=\Gamma_{*,\Lambda}\mathcal R,
\qquad
\mathcal R^\ast\Gamma_{*,\Lambda}=\Gamma_{*,\Lambda}\mathcal R^\ast.
\]
\end{theorem}

\begin{proof}
The background $\bar\Phi$ is real, the admissible even gauge-fixing term \eqref{eq:even-gauge-fixing} is built from a real quadratic norm,
and the linearized principal part of $P_{\rm gf,ren}$ is normally hyperbolic / Klein--Gordon type with real coefficients around that real
background.  The boundary and junction conditions encoded in $B_\Lambda$ are likewise real: the cutoff-surface condition is fixed by the
renormalized variational principle, the KR head is on the positive induced-gravity Robin branch, and the interior brane conditions are
linearized Israel/GW conditions with real coefficients.  Therefore the differential and boundary operators commute with the transpose / real
involution in the pointer basis.

The finite-cutoff spectral projector, the canonical record projection \eqref{eq:PiDelta}, and the QND dephasing channel of
Lemma~\ref{lem:tick-fixes-record-algebra} are pointer-diagonal, hence commute exactly with the same involution.  The admissible temporal
windowing and test kernels are real scalar weights on the receiver head, so they introduce no additional commutator defect.  Finally,
the sourced coupling \eqref{eq:uom-sourced-action} pairs real boundary densities with real source components, so the corresponding source
insertion and boundary extraction maps inherit the same pointer-real equivariance.  Thus all displayed commutators vanish identically at
fixed finite cutoff.
\end{proof}

\begin{corollary}[Pointer-real equivariance of Green operators and induced boundary causal operators]
\label{cor:uom-pointer-real-greens}
Assume Theorem~\ref{thm:uom-P2-green} and Theorem~\ref{thm:uom-pointer-real-ibvp}.  Then the advanced/retarded Green operators satisfy
\[
E^\pm\Gamma_{*,\Lambda}=\Gamma_{*,\Lambda}E^\pm,
\qquad
E\Gamma_{*,\Lambda}=\Gamma_{*,\Lambda}E,
\]
and the induced boundary retarded/advanced and causal operators satisfy
\[
G_R\Gamma_{*,\Lambda}=\Gamma_{*,\Lambda}G_R,
\qquad
G_A\Gamma_{*,\Lambda}=\Gamma_{*,\Lambda}G_A,
\qquad
\Delta\Gamma_{*,\Lambda}=\Gamma_{*,\Lambda}\Delta.
\]
\end{corollary}

\begin{proof}
Let $f$ be a compactly supported admissible source and set
\(
\psi^\pm:=E^\pm(\Gamma_{*,\Lambda}f)
\)
and
\(
\phi^\pm:=\Gamma_{*,\Lambda}(E^\pm f)
\).
By Theorem~\ref{thm:uom-pointer-real-ibvp},
\[
P_{\rm gf,ren}\phi^\pm
=
\Gamma_{*,\Lambda}P_{\rm gf,ren}E^\pm f
=
\Gamma_{*,\Lambda}f,
\qquad
B_\Lambda(\phi^\pm)
=
\Gamma_{*,\Lambda}B_\Lambda(E^\pm f)
=
0.
\]
Because $\Gamma_{*,\Lambda}$ acts pointwise in the fixed finite-band pointer basis, it preserves support.  Thus $\phi^\pm$ has the same
advanced/retarded support property as $E^\pm f$.  By uniqueness in Theorem~\ref{thm:uom-P2-green},
\(
\psi^\pm=\phi^\pm
\),
so
\(
E^\pm\Gamma_{*,\Lambda}=\Gamma_{*,\Lambda}E^\pm
\)
and hence
\(
E\Gamma_{*,\Lambda}=\Gamma_{*,\Lambda}E
\).
Composing with the $\Gamma_{*,\Lambda}$--equivariant source insertion and observable extraction maps from
Theorem~\ref{thm:uom-pointer-real-ibvp} and using \eqref{eq:GR-from-E} gives the corresponding identities for
$G_R$, $G_A$, and $\Delta$.
\end{proof}

\begin{lemma}[Pointer-real invariance of the Peierls pairing]
\label{lem:uom-pointer-real-pairing}
The pairings used in the Peierls construction are $\Gamma_{*,\Lambda}$--invariant.  In particular, for admissible sources and linearized
boundary data one has
\[
\langle \Gamma_{*,\Lambda}f,\Gamma_{*,\Lambda}\delta\Phi\rangle
=
\langle f,\delta\Phi\rangle,
\qquad
\langle \Gamma_{*,\Lambda}u,\Gamma_{*,\Lambda}v\rangle
=
\langle u,v\rangle,
\]
where the first pairing is the bulk/source pairing induced by the quadratic part of $S_{\rm ren}^{\rm gf}$ and the second is the boundary
pairing used in the Peierls formula on $\mathcal F_{\rm s.loc}^\Lambda$.
\end{lemma}

\begin{proof}
The bulk/source pairing in UOM-P3 is induced by the quadratic part of the gauge-fixed renormalized action, while the boundary Peierls pairing
is the ordinary real $L^2/H^s$ pairing from the receiver-head volume form and the finite family of boundary components; see
\eqref{eq:W-keldysh-expansion}, \eqref{eq:Omega-inverse}, and Lemma~\ref{lem:pairing-cont}.  In the fixed pointer basis, the involution
$\Gamma_{*,\Lambda}$ is transpose / real conjugation, so it is orthogonal for the fiber metric and preserves the real background volume form.
Therefore both pairings are unchanged when both entries are acted on by $\Gamma_{*,\Lambda}$.
\end{proof}

\begin{theorem}[Finite-cutoff pointer-real Poisson automorphism]
\label{thm:uom-pointer-real-poisson}
Assume UOM-P2--P5 and Definition~\ref{def:uom-gamma-star}.  Then the pullback involution $\Gamma$ acts by Poisson automorphisms on the
sufficiently local boundary class:
\[
\Gamma\{F,G\}_{\rm P}
=
\{\Gamma F,\Gamma G\}_{\rm P}
\qquad
\forall\,F,G\in\mathcal F_{\rm s.loc}^\Lambda.
\]
\end{theorem}

\begin{proof}
By the chain rule for pullback along the linear involution $\Gamma_{*,\Lambda}$,
\[
\bigl(\Gamma F\bigr)^{(1)}(\varphi)
=
\Gamma_{*,\Lambda}\bigl(F^{(1)}(\Gamma_{*,\Lambda}\varphi)\bigr),
\qquad
\bigl(\Gamma G\bigr)^{(1)}(\varphi)
=
\Gamma_{*,\Lambda}\bigl(G^{(1)}(\Gamma_{*,\Lambda}\varphi)\bigr).
\]
By Corollary~\ref{cor:uom-pointer-real-greens},
\(
\Delta\Gamma_{*,\Lambda}=\Gamma_{*,\Lambda}\Delta
\),
and by Lemma~\ref{lem:uom-pointer-real-pairing} the Peierls pairing is $\Gamma_{*,\Lambda}$--invariant.
Fix \(\varphi\in\mathcal X_\Lambda^s\). Then
\[
\{\Gamma F,\Gamma G\}_{\rm P}(\varphi)
=
\big\langle (\Gamma F)^{(1)}(\varphi),\,\Delta\,(\Gamma G)^{(1)}(\varphi)\big\rangle
=
\big\langle \Gamma_{*,\Lambda}F^{(1)}(\Gamma_{*,\Lambda}\varphi),\,
\Gamma_{*,\Lambda}\Delta G^{(1)}(\Gamma_{*,\Lambda}\varphi)\big\rangle
=
\big\langle F^{(1)}(\Gamma_{*,\Lambda}\varphi),\,\Delta G^{(1)}(\Gamma_{*,\Lambda}\varphi)\big\rangle
=
\{F,G\}_{\rm P}(\Gamma_{*,\Lambda}\varphi)
=
(\Gamma\{F,G\}_{\rm P})(\varphi).
\]
Hence \(\Gamma\{F,G\}_{\rm P}=\{\Gamma F,\Gamma G\}_{\rm P}\).
\end{proof}

\paragraph{Conclusion of UOM-P6.}
Definition~\ref{def:uom-gamma-star}, Theorem~\ref{thm:uom-pointer-real-ibvp}, Corollary~\ref{cor:uom-pointer-real-greens},
Lemma~\ref{lem:uom-pointer-real-pairing}, and Theorem~\ref{thm:uom-pointer-real-poisson} establish \textbf{UOM-P6}: the fixed finite-cutoff
pointer-real involution, the exact equivariance of the gauge-fixed linearized renormalized IBVP, and the resulting swap-equivariant Peierls
structure on the sufficiently local boundary class.

% ======================================================================
% Energy bounds + Poincar\'e covariance package (EB1--EB3 + covariance)
% Place this subsection after the UOM-P1--P6 package subsection
% (e.g. after \subsubsection{UOM-P6: ...}) inside the promoted section
% \section{Two--times semiclassical characterization: acceptance, diagonalization, and flatness}
% ======================================================================

\subsection{Energy bounds and Poincar\'e covariance on a Minkowski UV chart (EB1--EB3 and covariance)}
\label{subsec:uom-energy-bounds}

\paragraph{Purpose.}
The two--time Main Theorem constructs an odd derivation by Hamiltonization
\(
\delta(A)=\{\mathcal O^{(1)},A\}_{\rm P}
\)
and then needs an \emph{implementability} step: on the vacuum GNS space,
$\delta$ must act on a dense core with an \emph{energy bound} strong enough
to guarantee closability of the would-be generator $Q$ defined by
$Q\,\pi(A)\Omega:=i\,\pi(\delta(A))\Omega$.
This subsection supplies, \emph{from UOM data}, the exact analytic inputs:
\textbf{EB1--EB3} and \textbf{Poincar\'e covariance} of the causal operator and Peierls bracket,
in a form that isolates precisely what must be uniform in the cutoff
\(\Lambda\to\infty\).

\paragraph{Relationship to UOM-P1--P6.}
We assume the UOM Peierls package at finite cutoff (Theorem~\ref{thm:uom-peierls-package}):
existence of the renormalized SK functional $W_{\rm ren}$ (UOM-P1),
existence of the causal operator $\Delta_\Lambda$ on the receiver head (UOM-P2),
existence of the renormalized covariant symplectic structure and gauge independence on invariants (UOM-P3--P4),
the sufficiently local class $\mathcal F_{\rm s.loc}^\Lambda$ closed under Peierls (UOM-P5),
and the exact finite-cutoff pointer-real swap-equivariance package (UOM-P6).
The new content here is \emph{uniformity} in $\Lambda$ and its translation into energy bounds
on the vacuum representation.

% ----------------------------------------------------------------------
% 1. Minkowski UV chart and local regions
% ----------------------------------------------------------------------

\paragraph{Minkowski UV chart and double cones.}
Let $\mathscr U\subset \mathbb R\times S^3$ be a conformal chart on the (receiver) UV boundary cylinder
and let $\kappa:\mathscr U\to \mathbb R^{1,3}$ be a conformal diffeomorphism onto its image
(with conformal factor $\Omega_\kappa>0$), so that $\kappa$ identifies a neighborhood of the
receiver UV head with a Minkowski region.  A \emph{double cone} $O\subset \mathbb R^{1,3}$ means the
intersection of a future and past light cone with compact closure.  We identify $O$ with its preimage
$\kappa^{-1}(O)\subset\mathscr U$ and freely switch between the two descriptions.  We write
$\mathcal F_{\rm s.loc}^\Lambda(O)$ for the sufficiently local class (UOM-P5) supported in $O$,
expressed in this Minkowski chart.

\paragraph{Vacuum state and Hamiltonian.}
Let $(\mathcal H_\Lambda,\pi_\Lambda,\Omega_\Lambda)$ denote the Hilbert space, representation, and
cyclic vector reconstructed from the receiver Euclidean Schwinger functions at cutoff $\Lambda$
(OS reconstruction; see below).  Let $U_\Lambda(a)$ be the unitary implementing time translations
in the Minkowski chart, and let $H_\Lambda:=P^0_\Lambda$ be its self-adjoint generator:
$U_\Lambda(t)=e^{-itH_\Lambda}$.  At finite cutoff, $H_\Lambda$ may be bounded (band-limited regime),
but in the UV limit it becomes unbounded.  The vacuum core is
\[
\mathcal D_{0,\Lambda}\;:=\;\pi_\Lambda(\mathfrak A_{{\rm fin},\Lambda})\,\Omega_\Lambda,
\]
where $\mathfrak A_{{\rm fin},\Lambda}$ is the unital $*$-algebra generated by finitely many local
bounded functions of smeared fields/observables in $\mathcal F_{\rm s.loc}^\Lambda$ (any standard choice
yields the same invariant core).

% ----------------------------------------------------------------------
% 2. OS reconstruction uniformly in Lambda and the UV limit
% ----------------------------------------------------------------------

\subsubsection*{OS reconstruction and the UV inductive limit}

To speak about uniform bounds and Poincar\'e covariance in the limit $\Lambda\to\infty$, we
formalize the cutoff family as a directed system.

\begin{proposition}[Route to nested OS reconstruction from the cutoff Schwinger family]
\label{prop:uom-OS-inductive-route}
For each cutoff $\Lambda\ge\Lambda_0$, let
\(
S_{n,\Lambda}(x_1,\dots,x_n)
\)
denote the receiver Euclidean Schwinger functions on the Minkowski UV chart.  Assume:
\begin{enumerate}[label=\textnormal{(\roman*)},leftmargin=2.2em,itemsep=1pt]
\item \textbf{Finite-cutoff OS positivity.}
For every $\Lambda$, the cutoff family satisfies reflection positivity on the receiver side; in the present UOM lane this is supplied by
Theorem~\ref{thm:OS_receiver} together with the discrete OS/Hankel positivity construction of Appendix~\ref{app:OS-discrete}.
\item \textbf{Cutoff compatibility.}
If $\Lambda'\ge\Lambda$, then the $\Lambda$-cutoff Schwinger functions are obtained from the $\Lambda'$-cutoff ones by the lower-cutoff
receiver projector: in the hard-cutoff picture this is extension-by-zero / spectral inclusion on modes above $\Lambda$, while for smooth
tapers it is the natural inclusion of weighted spectral subspaces.
\item \textbf{Time-translation consistency.}
The Euclidean time-translation semigroup and the local observable classes $\mathcal F_{\rm s.loc}^\Lambda(O)$ are compatible with the
cutoff maps of \textnormal{(ii)}.
\end{enumerate}
Then the uniform OS-inductive statement is proved by the following route:
\begin{enumerate}[label=\textnormal{(O\arabic*)},leftmargin=2.4em,itemsep=1pt]
\item each $\Lambda$ admits an OS reconstruction
\((\mathcal H_\Lambda,\pi_\Lambda,\Omega_\Lambda,U_\Lambda)\);
\item for $\Lambda'\ge\Lambda$, the cutoff map of \textnormal{(ii)} induces an isometric embedding
\(
\iota_{\Lambda\to\Lambda'}:\mathcal H_\Lambda\hookrightarrow\mathcal H_{\Lambda'}
\)
and compatible injections of local algebras;
\item these embeddings form a directed system, so the Hilbert-space and local-algebra inductive limits exist.
\end{enumerate}
\end{proposition}

\begin{proof}
Item~\textnormal{(O1)} is the standard OS reconstruction theorem applied at each fixed cutoff.  Item~\textnormal{(O2)} follows because the
lower-cutoff Schwinger functions are exactly the restrictions of the higher-cutoff ones under the spectral-inclusion map, so null vectors and
inner products are preserved.  Item~\textnormal{(O3)} is then the standard inductive-limit construction for directed systems of isometric
embeddings and injective $*$-homomorphisms.
\end{proof}

\begin{theorem}[Uniform OS reconstruction and inductive limit]
\label{thm:uom-OS-inductive}
Under the hypotheses of Proposition~\ref{prop:uom-OS-inductive-route}, for each cutoff $\Lambda\ge \Lambda_0$ there exists an OS reconstruction
of the receiver Euclidean Schwinger functions on the Minkowski chart yielding:
\begin{enumerate}[label=\textnormal{(\roman*)},leftmargin=2.2em,itemsep=1pt]
\item a vacuum representation $(\mathcal H_\Lambda,\pi_\Lambda,\Omega_\Lambda)$ and a positive-energy
unitary representation of Minkowski time translations $U_\Lambda(t)=e^{-itH_\Lambda}$;
\item local algebras $\mathfrak A_\Lambda(O)$ generated by $\mathcal F_{\rm s.loc}^\Lambda(O)$;
\item for $\Lambda'\ge \Lambda$ isometric embeddings
\(
\iota_{\Lambda\to\Lambda'}:\mathcal H_\Lambda\hookrightarrow\mathcal H_{\Lambda'}
\)
and injective $*$-homomorphisms
\(
\jmath_{\Lambda\to\Lambda'}:\mathfrak A_\Lambda(O)\hookrightarrow\mathfrak A_{\Lambda'}(O)
\)
compatible with vacuum vectors and inclusions of local observables, so that
\[
\iota_{\Lambda\to\Lambda'}\,\Omega_\Lambda=\Omega_{\Lambda'},
\qquad
\iota_{\Lambda\to\Lambda'}\,\pi_\Lambda(A)=\pi_{\Lambda'}(\jmath_{\Lambda\to\Lambda'}A)\,\iota_{\Lambda\to\Lambda'}.
\]
\end{enumerate}
Moreover, the inductive limits exist:
\[
\mathcal H:=\varinjlim_{\Lambda}\mathcal H_\Lambda,
\qquad
\mathfrak A(O):=\varinjlim_{\Lambda}\mathfrak A_\Lambda(O),
\qquad
\Omega:=\varinjlim_{\Lambda}\Omega_\Lambda.
\]
\end{theorem}

\begin{proof}
We verify the three route items of Proposition~\ref{prop:uom-OS-inductive-route}.

\emph{Step 1: OS reconstruction at each cutoff.}
At fixed finite cutoff, Theorem~\ref{thm:OS_receiver} gives reflection positivity for the receiver Schwinger functions, while
Appendix~\ref{app:OS-discrete} shows that the discrete OS matrices are positive and stable under refinement/tapering.  Together with the
Euclidean invariance and regularity already built into the cutoff Schwinger family, this is exactly the input for the standard
Osterwalder--Schrader reconstruction theorem.  Therefore each cutoff $\Lambda$ yields a vacuum Hilbert space
\((\mathcal H_\Lambda,\pi_\Lambda,\Omega_\Lambda)\) with positive-energy Euclidean-time continuation
\(U_\Lambda(t)=e^{-itH_\Lambda}\).  This proves route item~\textnormal{(O1)}.

\emph{Step 2: embeddings between cutoffs.}
Let $\Lambda'\ge\Lambda$.  By the cutoff-compatibility hypothesis, the lower-cutoff Schwinger functions are obtained from the higher-cutoff ones
by projecting away the modes above $\Lambda$; in the hard-cutoff picture this is literal extension-by-zero on those higher modes, while in the
smooth-taper picture it is the natural inclusion of the lower weighted spectral subspace into the higher one.  On the OS pre-Hilbert space of
test-function classes this map preserves the reflection-positive sesquilinear form, because the latter is computed from the same restricted
Schwinger functions on both sides.  Hence it descends to an isometric embedding
\(
\iota_{\Lambda\to\Lambda'}:\mathcal H_\Lambda\hookrightarrow\mathcal H_{\Lambda'}
\).
The same cutoff map sends lower-cutoff local observables to the corresponding higher-cutoff observables, yielding injective
$*$-homomorphisms
\(
\jmath_{\Lambda\to\Lambda'}:\mathfrak A_\Lambda(O)\hookrightarrow\mathfrak A_{\Lambda'}(O)
\)
compatible with the GNS actions and the vacuum vectors.  This proves route item~\textnormal{(O2)}.

\emph{Step 3: inductive limit.}
The embeddings of Step~2 are functorial in the cutoff parameter: for $\Lambda''\ge\Lambda'\ge\Lambda$ one has
\(
\iota_{\Lambda'\to\Lambda''}\circ \iota_{\Lambda\to\Lambda'}=\iota_{\Lambda\to\Lambda''}
\)
and similarly for the local-algebra maps.  Therefore the family
\(
\{(\mathcal H_\Lambda,\mathfrak A_\Lambda(O),\Omega_\Lambda)\}_{\Lambda\ge\Lambda_0}
\)
forms a directed system.  The Hilbert-space inductive limit, the local-algebra inductive limit, and the vacuum vector
\(
\Omega=\varinjlim_\Lambda \Omega_\Lambda
\)
are then defined by the standard completion/direct-limit construction.  This proves route item~\textnormal{(O3)} and yields the stated theorem.
\end{proof}

\begin{proposition}[Route to the strong resolvent limit of the nested Hamiltonians]
\label{prop:uom-H-limit-route}
Assume Theorem~\ref{thm:uom-OS-inductive}, and let
\(
\iota_\Lambda:\mathcal H_\Lambda\hookrightarrow \mathcal H
\)
denote the canonical maps into the inductive-limit Hilbert space.  Write
\[
\mathcal D_{\rm ind}:=\bigcup_{\Lambda\ge\Lambda_0}\iota_\Lambda(\mathcal D_{0,\Lambda}),
\qquad
q_\Lambda(\psi):=\langle \psi,H_\Lambda\psi\rangle_{\mathcal H_\Lambda}.
\]
Assume:
\begin{enumerate}[label=\textnormal{(\roman*)},leftmargin=2.2em,itemsep=1pt]
\item \textbf{Form compatibility under cutoff inclusion.}
For $\Lambda'\ge \Lambda$, the embedding $\iota_{\Lambda\to\Lambda'}$ identifies the lower-cutoff time-translation semigroup with the restriction of the higher-cutoff one, so on $\mathcal D_{0,\Lambda}$ the quadratic forms satisfy
\[
q_{\Lambda'}(\iota_{\Lambda\to\Lambda'}\psi)
=
q_\Lambda(\psi)+r_{\Lambda,\Lambda'}(\psi),
\qquad
r_{\Lambda,\Lambda'}(\psi)\ge 0.
\]
Thus the ultraviolet increment contributes a nonnegative form term.
\item \textbf{Monotone cutoff realization / quadratic-form convergence.}
In the hard spectral-cutoff realization, the family $H_\Lambda$ is obtained by increasing spectral truncation of the same positive receiver Hamiltonian.  In the smooth-taper realization, the forms $q_\Lambda$ extend to closed semibounded forms on $\mathcal H_\Lambda$ whose pullbacks to $\mathcal D_{\rm ind}$ converge pointwise to a semibounded form $q$.
\item \textbf{Core control.}
The algebraic inductive core $\mathcal D_{\rm ind}$ is dense in $\mathcal H$ and is a common form core for the pulled-back cutoff forms.
\end{enumerate}
Then the strong-resolvent statement is proved by the following route:
\begin{enumerate}[label=\textnormal{(H\arabic*)},leftmargin=2.4em,itemsep=1pt]
\item the monotone/form limit defines a densely defined closed semibounded quadratic form $q$ on $\mathcal H$;
\item the associated self-adjoint operator $H$ is the strong resolvent limit of the nested $H_\Lambda$;
\item the contraction semigroups $e^{-\tau H_\Lambda}$ and hence the unitary groups $U_\Lambda(t)=e^{-itH_\Lambda}$ are compatible with the inductive system and determine the inductive-limit Minkowski time-translation dynamics.
\end{enumerate}
\end{proposition}

\begin{proof}
Item~\textnormal{(H1)} is the standard quadratic-form construction: positivity and monotonicity in the hard-cutoff case, or pointwise convergence of closed semibounded forms in the smooth-taper case, define the limit form on the common inductive core and then by closure on $\mathcal H$.  Item~\textnormal{(H2)} is Kato's monotone/form convergence theorem for semibounded self-adjoint operators.  Item~\textnormal{(H3)} follows because the time-translation semigroups already intertwine the cutoff embeddings at the Schwinger-function level, so the limiting semigroup and its Stone continuation are well defined on the inductive limit.
\end{proof}

\begin{theorem}[Strong resolvent convergence of the nested Hamiltonians]
\label{thm:uom-H-limit}
Under the hypotheses of Proposition~\ref{prop:uom-H-limit-route}, there exists a positive self-adjoint operator $H$ on the inductive-limit Hilbert space $\mathcal H$ such that:
\begin{enumerate}[label=\textnormal{(\roman*)},leftmargin=2.2em,itemsep=1pt]
\item the embedded cutoff generators $H_\Lambda$ converge to $H$ in the strong resolvent sense;
\item for every $\tau\ge 0$ and every cutoff vector $\psi_\Lambda\in\mathcal H_\Lambda$,
\[
e^{-\tau H}\,\iota_\Lambda\psi_\Lambda
=
\iota_\Lambda\big(e^{-\tau H_\Lambda}\psi_\Lambda\big);
\]
\item the unitary group $U(t):=e^{-itH}$ fixes the vacuum and implements Minkowski time translations in the inductive limit:
\[
U(t)\Omega=\Omega,
\qquad
U(t)\,\pi(\Phi(F))\,U(t)^{-1}=\pi(\Phi(\alpha_tF))
\]
for every sufficiently local observable $F$ on the Minkowski UV chart.
\end{enumerate}
\end{theorem}

\begin{proof}
We verify the three route items of Proposition~\ref{prop:uom-H-limit-route}.

\emph{Step 1: define the limit quadratic form.}
By Theorem~\ref{thm:uom-OS-inductive}, the finite-cutoff Hamiltonians $H_\Lambda$ are positive self-adjoint and are organized by isometric embeddings into a directed system.  Item~\textnormal{(i)} of Proposition~\ref{prop:uom-H-limit-route} says that when a vector is viewed at a higher cutoff, the lower-cutoff quadratic form is recovered together with a nonnegative ultraviolet correction.  Therefore the pulled-back forms are monotone on the common algebraic core $\mathcal D_{\rm ind}$.  In the hard-cutoff realization this is exactly the standard increasing spectral-truncation situation; in the smooth-taper realization it is pointwise convergence of closed semibounded forms.  Hence for $\psi=\iota_\Lambda\psi_\Lambda\in\mathcal D_{\rm ind}$ the pointwise limit
\[
q(\psi):=\lim_{\Lambda\to\infty} q_\Lambda(\psi_\Lambda)
\]
exists on $\mathcal D_{\rm ind}$, is semibounded, and closes to a densely defined quadratic form on $\mathcal H$.

\emph{Step 2: strong resolvent convergence.}
Let $H$ be the self-adjoint operator associated with the closed form $q$.  By the monotone convergence theorem for positive quadratic forms in the hard spectral-cutoff case, and by the general theorem on pointwise convergence of closed semibounded forms in the smooth-taper case, the operators $H_\Lambda$ converge to $H$ in the strong resolvent sense along the embeddings $\iota_\Lambda$.  This proves item~\textnormal{(i)}.

\emph{Step 3: passage of the dynamics to the inductive limit.}
Because the Euclidean time-translation semigroups are compatible with the cutoff maps, one has on the embedded cutoff subspaces
\[
e^{-\tau H}\,\iota_\Lambda
=
\iota_\Lambda e^{-\tau H_\Lambda}
\qquad
(\tau\ge 0).
\]
The vacuum vectors are compatible under the embeddings by Theorem~\ref{thm:uom-OS-inductive}, so $e^{-\tau H}\Omega=\Omega$ and therefore, after Stone continuation, $U(t)\Omega=\Omega$.  At each finite cutoff, $U_\Lambda(t)=e^{-itH_\Lambda}$ implements Minkowski time translations on the OS-reconstructed local algebra; passing to the inductive limit through the compatible injections of Theorem~\ref{thm:uom-OS-inductive} yields
\[
U(t)\,\pi(\Phi(F))\,U(t)^{-1}
=
\pi(\Phi(\alpha_tF))
\]
for every sufficiently local observable $F$.  This proves items~\textnormal{(ii)} and~\textnormal{(iii)}.
\end{proof}

% ----------------------------------------------------------------------
% 3. Uniformity inputs: (U-Δ) smoothing remainder and (U-SD) scaling degree
% ----------------------------------------------------------------------

\subsubsection*{Uniformity inputs from microlocal renormalization}

The energy bounds EB1--EB3 are \emph{uniform} statements in $\Lambda$ on fixed compact
regions $O$.  UOM supplies them via two local uniformity principles:
(i) the cutoff dependence of the causal operator $\Delta_\Lambda$ is smoothing on $O$,
and (ii) renormalized local correlators have $\Lambda$-independent singular order.

\begin{proposition}[Route to the uniform local decomposition of the causal operator]
\label{prop:uom-U-Delta-route}
Fix a double cone $O\subset\mathbb R^{1,3}$ in the Minkowski UV chart and assume the finite-cutoff
Green-hyperbolic IBVP package UOM-P2 together with the previously discharged well-posedness and boundary-admissibility inputs.  Let
\(
\Delta_\Lambda=E_\Lambda^+-E_\Lambda^-
\)
denote the cutoff causal propagator on the chart.  Assume:
\begin{enumerate}[label=\textnormal{(\roman*)},leftmargin=2.2em,itemsep=1pt]
\item \textbf{Local Hadamard parametrix.}
There exists a neighborhood $U\supset O$ and a distributional kernel $\Delta_{\rm loc}$ on $U\times U$ determined only by the local principal symbol
of the UV-limit gauge-fixed operator, such that $\Delta_{\rm loc}$ has the universal Hadamard singularity structure along the diagonal and solves the
inhomogeneous equation modulo a smooth error on $U$.
\item \textbf{Boundary/image decomposition.}
For each sufficiently large cutoff, the difference
\[
R_{\Lambda,O}:=\Delta_\Lambda\big|_{O}-\Delta_{\rm loc}\big|_{O}
\]
solves the homogeneous linearized equation on $O$; any remaining singularities of $R_{\Lambda,O}$ can only arise from bicharacteristics reaching $O$
after reflection from the UV wall or cutoff junction.
\item \textbf{Propagation of singularities and uniform regularity.}
For fixed $O$, once $\Lambda\ge\Lambda_O$ no such reflected/image singularity meets $O\times O$; hence
\(
\operatorname{WF}(R_{\Lambda,O})\cap T^*(O\times O)=\varnothing
\).
Moreover the homogeneous remainders satisfy $\Lambda$-uniform local energy/Sobolev estimates on $O$.
\end{enumerate}
Then the uniform decomposition statement is proved by the following route:
\begin{enumerate}[label=\textnormal{(D\arabic*)},leftmargin=2.4em,itemsep=1pt]
\item the cutoff-independent local singular kernel $\Delta_{\rm loc}$ captures the entire microlocal singular support on $O$;
\item the remainder $S_{\Lambda,O}:=R_{\Lambda,O}$ is smooth on $O\times O$ for all $\Lambda\ge\Lambda_O$;
\item the uniform local energy estimates upgrade this smoothness to Sobolev-uniform smoothing bounds of the form
\eqref{eq:smoothing-uniform}.
\end{enumerate}
\end{proposition}

\begin{proof}
Item~\textnormal{(D1)} is the standard Hadamard/parametrix construction for normally hyperbolic or Green-hyperbolic operators on a fixed local chart: the singular part depends only on the local principal symbol and not on the distant cutoff wall.  Item~\textnormal{(D2)} follows because after subtracting this local parametrix the remainder solves the homogeneous equation near $O$, so the only possible new singularities are boundary-reflected ones; propagation of singularities excludes them from $O\times O$ once the cutoff is sufficiently high.  Item~\textnormal{(D3)} is then standard local regularity for homogeneous solutions on compact sets, using the $\Lambda$-uniform energy estimates furnished by the well-posed UOM-P2 IBVP.
\end{proof}

\begin{theorem}[Uniform local decomposition of the causal operator]
\label{thm:uom-U-Delta}
Assume the hypotheses of Proposition~\ref{prop:uom-U-Delta-route}.  Then there exists $\Lambda_O<\infty$ and a distributional operator
$\Delta_{\rm loc}$ on $O$ such that for all $\Lambda\ge \Lambda_O$,
\begin{equation}
\Delta_\Lambda\big|_{O}
\;=\;
\Delta_{\rm loc}\;+\;S_{\Lambda,O},
\label{eq:Delta-decomp}
\end{equation}
where $S_{\Lambda,O}$ is \emph{smoothing on $O$} and uniformly bounded in Sobolev norms:
for each $s>\frac52$ there exist $k\ge 1$ and $C_{s,O}<\infty$ such that
\begin{equation}
\|S_{\Lambda,O} f\|_{H^{s+k}(O)}\ \le\ C_{s,O}\,\|f\|_{H^{s}(O)}
\qquad\text{for all }f\in H^s_0(O)\text{ and all }\Lambda\ge \Lambda_O.
\label{eq:smoothing-uniform}
\end{equation}
\end{theorem}

\begin{proof}
We verify the three route items of Proposition~\ref{prop:uom-U-Delta-route}.

\emph{Step 1: construct the local singular piece.}
Choose a neighborhood $U\supset O$ contained in the Minkowski UV chart and build a Hadamard parametrix for the UV-limit gauge-fixed operator on $U$.
Because this construction depends only on the local symbol and transport equations along the diagonal, the resulting causal kernel $\Delta_{\rm loc}$ is
independent of the cutoff and captures the universal local singularity structure.  This proves route item~\textnormal{(D1)}.

\emph{Step 2: identify the cutoff dependence as a homogeneous remainder.}
For each finite cutoff, set
\(
R_{\Lambda,O}:=\Delta_\Lambda\big|_O-\Delta_{\rm loc}\big|_O
\).
By construction, the singular source term has been removed, so $R_{\Lambda,O}$ solves the homogeneous linearized equation on $O$ in each variable.
Any singularity of $R_{\Lambda,O}$ must therefore come from propagation along null bicharacteristics entering from outside the local chart, i.e.\ from
reflections or image terms generated by the UV wall or cutoff junction.  By the fixed-double-cone hypothesis in item~\textnormal{(iii)} of
Proposition~\ref{prop:uom-U-Delta-route}, these reflected rays do not meet $O\times O$ once $\Lambda\ge\Lambda_O$.  Hence
\(
\operatorname{WF}(R_{\Lambda,O})\cap T^*(O\times O)=\varnothing
\),
so $R_{\Lambda,O}$ is smooth on $O\times O$.  This proves route item~\textnormal{(D2)}.

\emph{Step 3: uniform smoothing bounds.}
Because $R_{\Lambda,O}$ is a smooth homogeneous remainder on the fixed compact region $O$ and the underlying UOM-P2 IBVP is uniformly well posed at high
cutoff, standard local energy estimates and Sobolev propagation bounds give for each $s>\frac52$ some $k\ge 1$ and $C_{s,O}$ independent of $\Lambda$
such that
\[
\|R_{\Lambda,O} f\|_{H^{s+k}(O)}\le C_{s,O}\|f\|_{H^s(O)}
\qquad
(f\in H_0^s(O),\ \Lambda\ge\Lambda_O).
\]
Defining $S_{\Lambda,O}:=R_{\Lambda,O}$ yields \eqref{eq:Delta-decomp} and \eqref{eq:smoothing-uniform}.  This proves route item~\textnormal{(D3)}
and hence the theorem.
\end{proof}

\begin{proposition}[Route to uniform scaling degree and temperedness for renormalized local correlators]
\label{prop:uom-U-SD-route}
Fix a finite generating family of local gauge-invariant densities $\mu_A(x)$ used to build
$\mathcal F_{\rm s.loc}^\Lambda$ (UOM-P5) on the Minkowski UV chart, and fix once and for all a holographic/local counterterm scheme independent of
the cutoff $\Lambda$.  Assume:
\begin{enumerate}[label=\textnormal{(\roman*)},leftmargin=2.2em,itemsep=1pt]
\item \textbf{Fixed renormalization prescription.}
For every $n$ and every multi-index of insertions, the renormalized Euclidean correlators
\(
S^{(\Lambda)}_{A_1\cdots A_n}(x_1,\ldots,x_n)
\)
are defined by the same local covariant subtraction rules for all sufficiently large cutoffs, so the $\Lambda$-dependence enters only through the
unrenormalized amplitudes and finite renormalized coefficients, not through changing subtraction operators.
\item \textbf{Cutoff-independent local singular skeleton.}
The singular part of the local two-point kernels is governed by the same Hadamard/parametrix structure on each compact region, with cutoff dependence
confined to smooth remainders as supplied by Theorem~\ref{thm:uom-U-Delta}.  Higher correlators are built from these local propagator singularities and
local interaction vertices by the fixed subtraction scheme.
\item \textbf{Uniform local counterterm control.}
For every double cone $O$ and every compact $K\subset O^n$, the allowed counterterms are local polynomials of bounded engineering dimension determined by the fixed scheme, so their
wavefront sets are conormal to partial diagonals and their scaling degrees are bounded independently of $\Lambda$.
\end{enumerate}
Then the uniform short-distance statement is proved by the following route:
\begin{enumerate}[label=\textnormal{(S\arabic*)},leftmargin=2.4em,itemsep=1pt]
\item the renormalized correlators admit on each compact a decomposition into finitely many universal Hadamard-type singular pieces plus $\Lambda$-uniformly
smooth remainders;
\item the wavefront set of each renormalized correlator is contained in a cutoff-independent microlocal spectrum cone determined by the local Hadamard
structure and the conormal directions to partial diagonals;
\item the Steinmann scaling degree and temperedness bounds are therefore controlled uniformly in $\Lambda$ by the fixed subtraction degree and local
power counting.
\end{enumerate}
\end{proposition}

\begin{proof}
Item~\textnormal{(S1)} is the standard consequence of using one fixed local renormalization prescription together with a cutoff-independent Hadamard
singular kernel: after subtracting the same singular jets at every cutoff, only smooth remainder terms retain $\Lambda$-dependence on compact sets.
Item~\textnormal{(S2)} follows from Hörmander calculus for products/extensions of distributions with controlled wavefront set and from the fact that local
counterterms contribute only conormal singularities on diagonals.  Item~\textnormal{(S3)} is then the usual power-counting statement: because the degree of
the subtraction operators is fixed by dimension and locality, the renormalized scaling degree and Schwartz bounds are uniform in $\Lambda$ on compacts.
\end{proof}

\begin{theorem}[Uniform scaling degree and temperedness for renormalized local correlators]
\label{thm:uom-U-SD}
Fix a finite generating family of local gauge-invariant densities $\mu_A(x)$ used to build
$\mathcal F_{\rm s.loc}^\Lambda$ (UOM-P5) on the Minkowski chart.  Let
$S^{(\Lambda)}_{A_1\cdots A_n}(x_1,\ldots,x_n)$ be the corresponding renormalized Euclidean
Schwinger functions at cutoff $\Lambda$ in the fixed holographic counterterm scheme of
Proposition~\ref{prop:uom-U-SD-route}.  Then for each $n$, each double cone $O$, and each compact
$K\subset O^n$ there exists a constant $C_{n,K}$ independent of $\Lambda$ and an integer $m(n)$ such that:
\begin{enumerate}[label=\textnormal{(\roman*)},leftmargin=2.2em,itemsep=1pt]
\item (\emph{Uniform temperedness}) $S^{(\Lambda)}_{A_1\cdots A_n}$ is a tempered distribution on $O^n$
and its action on test functions supported in $K$ is bounded by $C_{n,K}$ times finitely many Schwartz seminorms.
\item (\emph{Uniform singular order}) The Steinmann scaling degree at partial diagonals in $K$ is bounded
uniformly in $\Lambda$ by $m(n)$ (equivalently: the wavefront set and scaling degree bounds are $\Lambda$-stable).
\end{enumerate}
\end{theorem}

\begin{proof}
We verify the three route items of Proposition~\ref{prop:uom-U-SD-route}.

\emph{Step 1: fix the singular decomposition.}
Because the same holographic/local counterterm prescription is used for every cutoff, the renormalized correlators are obtained by applying one and the
same subtraction pattern to cutoff-dependent bare amplitudes.  Theorem~\ref{thm:uom-U-Delta} shows that on fixed compact regions the propagator-level
cutoff dependence is already confined to smooth remainders, while the Hadamard singular core is cutoff independent.  It follows that every renormalized
$n$-point function on $K$ decomposes into finitely many universal Hadamard-type singular distributions, determined only by the local geometry and the fixed
subtraction scheme, plus a remainder that is smooth and uniformly controlled in $\Lambda$.  This proves route item~\textnormal{(S1)}.

\emph{Step 2: wavefront-set control.}
The universal singular pieces from Step~1 have wavefront sets contained in the usual Hadamard/microlocal-spectrum cones and in the conormal bundles of the
partial diagonals where local counterterms live.  Smooth remainders contribute no wavefront set.  Since the subtraction operators and allowed local
counterterms do not depend on $\Lambda$, the resulting wavefront-set enclosure is the same for all sufficiently large cutoffs.  This proves route item~\textnormal{(S2)}.

\emph{Step 3: uniform temperedness and scaling degree.}
The Steinmann scaling degree of each singular piece is fixed by the engineering dimension of the inserted local densities together with the finite-order
subtractions prescribed by the fixed scheme.  Because those subtraction orders are cutoff independent, there is an integer $m(n)$ that bounds the scaling
degree at every partial diagonal uniformly in $\Lambda$.  The same finite decomposition, together with the uniform control of the smooth remainders on
compact sets, gives a common Schwartz-seminorm bound on test functions supported in $K$, yielding a constant $C_{n,K}$ independent of $\Lambda$.
Therefore the renormalized correlators are uniformly tempered and have cutoff-independent singular order on compacts.  This proves route item~\textnormal{(S3)} and hence the theorem.
\end{proof}

% ----------------------------------------------------------------------
% 4. EB2: uniform Sobolev continuity of the causal operator
% ----------------------------------------------------------------------

\subsubsection*{EB2: uniform Sobolev continuity of $\Delta_\Lambda$}

\begin{lemma}[Sobolev/energy continuity of the causal operator (uniform EB2)]
\label{lem:EB2-uniform}
Assume Theorem~\ref{thm:uom-U-Delta}.  Fix a double cone $O$ and $s>\frac52$.
Then there exists $C_\Delta(O)<\infty$ and $\Lambda_O<\infty$ such that for all $\Lambda\ge \Lambda_O$,
\begin{equation}
\|\Delta_\Lambda f\|_{H^{s}(O)}\ \le\ C_\Delta(O)\,\|f\|_{H^{s}(O)}
\qquad\text{for all }f\in H^s_0(O).
\label{eq:EB2}
\end{equation}
\end{lemma}

\begin{proof}
By Theorem~\ref{thm:uom-U-Delta}, $\Delta_\Lambda=\Delta_{\rm loc}+S_{\Lambda,O}$ on $O$ with $S_{\Lambda,O}$ uniformly smoothing.
The local piece $\Delta_{\rm loc}$ is a fixed continuous operator $H^s_0(O)\to H^s(O)$ (its kernel is the universal local causal
distribution determined by the limiting Green-hyperbolic operator).  The smoothing remainder satisfies
$\|S_{\Lambda,O}f\|_{H^{s}(O)}\le \|S_{\Lambda,O}f\|_{H^{s+k}(O)}\le C_{s,O}\|f\|_{H^s(O)}$ uniformly in $\Lambda$.
Taking $C_\Delta(O):=\|\Delta_{\rm loc}\|_{H^s_0\to H^s}+C_{s,O}$ yields \eqref{eq:EB2}.
\end{proof}

% ----------------------------------------------------------------------
% 5. EB1: uniform polynomial energy bounds for local generators
% ----------------------------------------------------------------------

\subsubsection*{EB1: uniform polynomial energy bounds for sufficiently local generators}

We now pass from uniform control of correlators (Theorem~\ref{thm:uom-U-SD}) to uniform polynomial
energy bounds for local generators in the vacuum representation.

\begin{definition}[Cutoff-stable sufficiently local class $\Fsloc(O)$]
\label{def:Fsloc-limit}
For a double cone $O$, define the cutoff-stable class
\[
\Fsloc(O)\;:=\;\varinjlim_{\Lambda}\ \mathcal F_{\rm s.loc}^\Lambda(O),
\]
i.e.\ equivalence classes of families $(F_\Lambda)_{\Lambda\ge\Lambda_0}$ with $F_\Lambda\in\mathcal F_{\rm s.loc}^\Lambda(O)$ and
$\jmath_{\Lambda\to\Lambda'}(F_\Lambda)=F_{\Lambda'}$ for $\Lambda'\ge\Lambda$.
Let $\Phi_\Lambda(F_\Lambda)$ denote the (OS/Wightman reconstructed) operator affiliated with $\mathfrak A_\Lambda(O)$ corresponding to $F_\Lambda$.
We write $\Phi(F)$ for the inductive limit operator on $\mathcal H$ obtained by compatibility under embeddings.
\end{definition}

\begin{remark}
At finite cutoff $\Lambda$ the representation of $F_\Lambda$ is typically bounded (band-limited regime).
The content of EB1 is that one can choose a \emph{single} polynomial energy exponent $s$ and a constant $C_{F,O}$ independent of $\Lambda$
for sufficiently large $\Lambda$, so that the limit operator is closable and controlled by $(1+H)^s$.
\end{remark}

\begin{lemma}[Uniform energy bounds for sufficiently local generators (EB1)]
\label{lem:EB1}
Assume Theorems~\ref{thm:uom-OS-inductive}, \ref{thm:uom-H-limit}, and~\ref{thm:uom-U-SD}.
Fix a Minkowski UV boundary chart and let $(\mathcal H,\pi,\Omega)$ be the vacuum representation obtained in the inductive limit.
Then for every double cone $O\subset\mathbb R^{1,3}$ and every $F\in\Fsloc(O)$, there exist $s\ge 0$ and $C_{F,O}>0$
such that on the vacuum core $\mathcal D_0:=\pi(\mathfrak A_{\rm fin})\Omega$ one has
\begin{equation}
\|\pi(\Phi(F))\psi\|\ \le\ C_{F,O}\,\|(1+H)^s\psi\|
\qquad(\psi\in\mathcal D_0),
\label{eq:EB1}
\end{equation}
and the same bound holds at finite cutoff uniformly for all $\Lambda\ge \Lambda_O$:
\[
\|\pi_\Lambda(\Phi_\Lambda(F_\Lambda))\psi_\Lambda\|
\ \le\ C_{F,O}\,\|(1+H_\Lambda)^s\psi_\Lambda\|
\qquad(\psi_\Lambda\in\mathcal D_{0,\Lambda}).
\]
\end{lemma}

\begin{proof}
Fix $F\in\Fsloc(O)$ represented by $F_\Lambda\in\mathcal F_{\rm s.loc}^\Lambda(O)$ for all $\Lambda$ large.
By Definition~\ref{def:Floc} (UOM-P5), $F_\Lambda$ is a finite sum of compact smearings of local densities $\mu_A$ and their derivatives.
Theorem~\ref{thm:uom-U-SD} gives uniform temperedness and uniform scaling-degree control for all Schwinger functions of the generating family
$\mu_A$ on compacts in $O$.  By OS reconstruction (Theorem~\ref{thm:uom-OS-inductive}), these bounds transfer to the corresponding Wightman
functions in the vacuum representation and imply polynomial energy bounds for smeared local fields/operators built from $\mu_A$.
Concretely: uniform scaling-degree bounds control the number of time derivatives needed to make the vacuum matrix elements square-integrable in
energy, and thus yield a bound of the form \eqref{eq:EB1} with some finite $s$ depending only on the singular order of $F$.
Because the singular order bounds are uniform in $\Lambda$ (Theorem~\ref{thm:uom-U-SD}), the exponent $s$ and constant $C_{F,O}$ can be chosen
uniformly for all sufficiently large $\Lambda$.  Finally, strong resolvent convergence (Theorem~\ref{thm:uom-H-limit}) ensures these graph-norm
bounds pass to the inductive limit.  (This is the standard OS/Wightman route from uniform short-distance control to uniform polynomial energy bounds.)
\end{proof}

% ----------------------------------------------------------------------
% 6. EB3: energy bound for the Hamiltonized derivation
% ----------------------------------------------------------------------

\subsubsection*{EB3: energy bound for the Hamiltonized derivation $\delta=\{\mathcal O^{(1)},-\}_{\rm P}$}

We now propagate the uniform bounds through the Peierls bracket to obtain the implementability estimate for $\delta$.

\begin{lemma}[Functional derivative control on $\Fsloc(O)$]
\label{lem:FD}
Fix $s>\frac52$ and a double cone $O$.  For each $F\in\Fsloc(O)$ (represented by $F_\Lambda\in\mathcal F_{\rm s.loc}^\Lambda(O)$),
the first functional derivative $F^{(1)}$ exists and defines an element of $H^s_0(O)$ satisfying
\begin{equation}
\|F^{(1)}\|_{H^s(O)}\ \le\ C^{(s)}_{F,O},
\label{eq:FD-bound}
\end{equation}
where $C^{(s)}_{F,O}$ depends only on finitely many Sobolev seminorms of the smearing kernels and coefficients defining $F$ and is independent
of $\Lambda$ for all $\Lambda\ge\Lambda_O$.
\end{lemma}

\begin{proof}
For $F$ of the form \eqref{eq:Floc-form} (UOM-P5), $F^{(1)}$ is a finite sum of compactly supported smooth functions obtained by applying finitely
many derivatives (up to the jet order of $F$) to the smearing kernels and multiplying by polynomial expressions in $\varphi$ and its derivatives.
At finite cutoff, $\varphi$ is $C^1$ and band-limited; the Sobolev norm of $F^{(1)}$ is therefore bounded by a constant depending on finitely
many Sobolev norms of the defining kernels.  Because $F$ is cutoff-stable and the defining kernels are fixed (not $\Lambda$-dependent), this
constant is independent of $\Lambda$ for all sufficiently large $\Lambda$.
\end{proof}

\begin{lemma}[Bilinear continuity of the Peierls pairing on $\Fsloc$]
\label{lem:pairing-cont}
Fix $s>\frac52$ and a double cone $O$.  The pairing
\(
\langle \cdot,\cdot\rangle: H^s_0(O)\times H^s(O)\to\mathbb R
\)
used in the Peierls formula is continuous: there exists $C_{\langle\cdot,\cdot\rangle}(O)$ such that
\[
|\langle f,g\rangle|\ \le\ C_{\langle\cdot,\cdot\rangle}(O)\,\|f\|_{H^s(O)}\,\|g\|_{H^s(O)}.
\]
\end{lemma}

\begin{proof}
On a bounded time slab and compact spatial slice, $H^s$ embeds into $L^2$ for $s>\frac52$ and the pairing is dominated by the $L^2$ inner product.
Thus continuity holds with a constant depending only on $O$ and the background volume form.
\end{proof}

\begin{lemma}[Energy bound for the Hamiltonized derivation (EB3)]
\label{lem:EB3}
Assume Lemmas~\ref{lem:EB1} and~\ref{lem:EB2-uniform}.  Let $\mathcal O^{(1)}\in\Fsloc(O_1)$ be a first descendant and define
$\delta(F):=\{\mathcal O^{(1)},F\}_{\rm P}$ on $\Fsloc$ via the Peierls bracket on the Minkowski chart.
Then for every double cone $O$ there exist $s\ge 0$ and $C_{O}>0$ such that for all $F\in\Fsloc(O)$ one has on the vacuum core
\begin{equation}
\|\pi(\Phi(\delta(F)))\Omega\|\ \le\ C_O\,\|(1+H)^s\,\pi(\Phi(F))\Omega\|.
\label{eq:EB3}
\end{equation}
Equivalently, $A\mapsto \pi(\delta(A))\Omega$ is $(1+H)^s$--bounded on $\mathfrak A_{\rm fin}\cap\mathfrak A(O)$.
\end{lemma}

\begin{proof}
By the Peierls formula (UOM-P5),
\[
\delta(F)=\langle \mathcal O^{(1)(1)},\,\Delta\,F^{(1)}\rangle.
\]
By Lemma~\ref{lem:FD} applied to $F$ and to $\mathcal O^{(1)}$, the derivatives $\mathcal O^{(1)(1)}$ and $F^{(1)}$ lie in $H^s_0$ for some
$s>\frac52$ and satisfy Sobolev bounds by constants depending only on the defining seminorms.
By Lemma~\ref{lem:EB2-uniform}, $\|\Delta F^{(1)}\|_{H^s}\le C_\Delta(O)\|F^{(1)}\|_{H^s}$.  By Lemma~\ref{lem:pairing-cont},
\[
|\delta(F)|\ \le\ C_{\langle\cdot,\cdot\rangle}(O)\,\|\mathcal O^{(1)(1)}\|_{H^s}\,\|\Delta F^{(1)}\|_{H^s}
\ \le\ C'\,p_s(\mathcal O^{(1)})\,p_s(F),
\]
for a constant $C'$ depending only on $O$ and Sobolev embedding constants, and seminorms $p_s$ determined by $F$ and $\mathcal O^{(1)}$.
This shows $\delta(F)\in \Fsloc(O)$ with a seminorm bound controlled by $F$.

Finally, apply EB1 (Lemma~\ref{lem:EB1}) to the generator $\Phi(\delta(F))$:
there exists $s\ge 0$ and $C_{\delta(F),O}$ such that
\(
\|\pi(\Phi(\delta(F)))\Omega\|\le C_{\delta(F),O}\|(1+H)^s\Omega\|.
\)
Because $\Omega$ is a vacuum vector, $\|(1+H)^s\Omega\|=\|\Omega\|=1$.
Moreover, the dependence of $C_{\delta(F),O}$ on $F$ is controlled by the seminorm bound above; absorbing all constants into $C_O$ yields
\eqref{eq:EB3} with the same exponent $s$ as in EB1 for $O$.
\end{proof}

% ----------------------------------------------------------------------
% 7. Poincaré covariance of Delta and of the Peierls bracket (EB4)
% ----------------------------------------------------------------------

\subsubsection*{Poincar\'e covariance of $\Delta$ and of the Peierls bracket}

\begin{proposition}[Route to Poincar\'e covariance of the UV-limit linearized IBVP]
\label{prop:uom-Poincare-IBVP-route}
Work on the Minkowski UV chart selected by the static AdS$_5$ boundary setup of Paragraph~\ref{par:uom-main-111}.  For each finite cutoff $\Lambda$,
write $P_{{\rm gf,ren},\Lambda}$ for the gauge-fixed linearized renormalized operator on the chart and $B_\Lambda$ for the corresponding linear
boundary/junction operator.  Let $G_{\rm Poin}\subset \mathrm{Conf}_0(\mathbb R^{1,3})$ denote the chart-preserving Poincar\'e subgroup, and let
$U_g$ be the induced pullback action on compactly supported test sources/fields.  Assume:
\begin{enumerate}[label=\textnormal{(\roman*)},leftmargin=2.2em,itemsep=1pt]
\item \textbf{Static microscopic geometry.}
By Paragraph~\ref{par:uom-main-111}, the UV boundary and outer head are treated as fixed boundary conditions rather than dynamical fields.  Hence the
principal symbol of $P_{{\rm gf,ren},\Lambda}$ and the static part of $B_\Lambda$ are the chartwise pullbacks of fixed Minkowski/local Robin data and
commute exactly with $G_{\rm Poin}$.
\item \textbf{Finite-cutoff symmetry defect is regulator-only.}
After separating off the static exact part, the remaining $\Lambda$-dependence of $P_{{\rm gf,ren},\Lambda}$ and $B_\Lambda$ enters only through the
regulated receiver bridge ingredients: the receiver spectral projector/taper, the mollifier in the regulated conformal intertwiner, and the accepted
temporal windowing.  Their chart-symmetry defects are controlled by Lemma~\ref{lem:approx-intertwiner-new}, Remark~\ref{rem:conf-monitor}, and the
per-tick conformal-fidelity estimate of Theorem~\ref{thm:tick-fidelity-RP}.
\item \textbf{Anti-Zeno accepted limit.}
Along the accepted UV schedule of Corollary~\ref{cor:anti-zeno}, the regulator-induced conformality defect
\(
\mathcal D_{\rm conf}(\varepsilon_\Lambda,\Lambda)
\)
and the temporal-window error both vanish.
\end{enumerate}
Then the UV-limit Poincar\'e covariance statement is proved by the following route:
\begin{enumerate}[label=\textnormal{(P\arabic*)},leftmargin=2.4em,itemsep=1pt]
\item define the finite-cutoff commutator defects
\[
\mathcal C_{P,\Lambda}(g):=U_gP_{{\rm gf,ren},\Lambda}-P_{{\rm gf,ren},\Lambda}U_g,
\qquad
\mathcal C_{B,\Lambda}(g):=U_gB_\Lambda-B_\Lambda U_g;
\]
\item the static part contributes zero to these commutators, and the regulator part obeys a uniform compact-group bound
\[
\sup_{g\in G}\Big(\|\mathcal C_{P,\Lambda}(g)\psi\|+\|\mathcal C_{B,\Lambda}(g)\psi\|_{H^{s-\frac12}(\partial\mathcal M^{(\Lambda)})}\Big)
\le
C_G\big(\mathcal D_{\rm conf}(\varepsilon_\Lambda,\Lambda)+\Phi(\sigma_{t,\Lambda}\Omega_{\rm gap})\big)\,\|\psi\|_{H^s}
\]
for $\psi$ in a fixed compactly supported Sobolev test class and $G\subset G_{\rm Poin}$ compact;
\item Corollary~\ref{cor:anti-zeno} sends the right-hand side to zero, so the strong local limits
$P_{\rm gf,ren}:=\lim_{\Lambda\to\infty}P_{{\rm gf,ren},\Lambda}$ and $B:=\lim_{\Lambda\to\infty}B_\Lambda$ commute exactly with the Poincar\'e subgroup.
\end{enumerate}
\end{proposition}

\begin{proof}
Item~\textnormal{(P1)} is only notation.  Item~\textnormal{(P2)} follows from the separation between static and regulator-dependent pieces: the static UV
boundary setup removes moving-boundary terms, so exact commutators can only come from the receiver cutoff/taper/mollifier and temporal-window ingredients.
Lemma~\ref{lem:approx-intertwiner-new} gives the $O(\varepsilon_\Lambda^\alpha)$ bound for the regulated intertwiner, Remark~\ref{rem:conf-monitor}
packages this as the conformality deficit $\mathcal D_{\rm conf}$, and Theorem~\ref{thm:tick-fidelity-RP} supplies the matching accepted-state error from
the temporal window.  Linearizing the accepted covariance estimate around the vacuum identifies the same defect bound for the source-level commutators of
$P_{{\rm gf,ren},\Lambda}$ and $B_\Lambda$ on compactly supported test fields.  Item~\textnormal{(P3)} is then immediate from
Corollary~\ref{cor:anti-zeno}.
\end{proof}

\begin{theorem}[Poincar\'e covariance of the UV-limit linearized IBVP]
\label{thm:uom-Poincare-IBVP}
Under the hypotheses of Proposition~\ref{prop:uom-Poincare-IBVP-route}, the strong local limits
\[
P_{\rm gf,ren}:=\lim_{\Lambda\to\infty} P_{{\rm gf,ren},\Lambda},
\qquad
B:=\lim_{\Lambda\to\infty} B_\Lambda
\]
exist on the Minkowski UV chart, and for each Poincar\'e transformation $g\in G_{\rm Poin}$ with unitary $U_g$ acting on test sources/fields one has
\[
U_g\,P_{\rm gf,ren}\;=\;P_{\rm gf,ren}\,U_g,
\qquad
U_g\,B\;=\;B\,U_g.
\]
In particular, the UV-limit gauge-fixed linearized renormalized operator and its limiting boundary/junction operator are Poincar\'e covariant on compactly
supported test sources.
\end{theorem}

\begin{proof}
We verify the three route items of Proposition~\ref{prop:uom-Poincare-IBVP-route}.

\emph{Step 1: isolate the exact static part.}
By Paragraph~\ref{par:uom-main-111}, the AdS$_5$ UV boundary and outer head are kept static.  Therefore no dynamical boundary metric or bending mode is
present in the microscopic linearized problem on the chosen UV chart.  The principal symbol of the gauge-fixed operator and the static part of the
boundary operator are thus exactly the Minkowski/flat-boundary ones, and these commute with the chart-preserving Poincar\'e subgroup $G_{\rm Poin}$.
This proves route item~\textnormal{(P1)} at the level of the exact part.

\emph{Step 2: bound the finite-cutoff commutator defects.}
Write
\[
P_{{\rm gf,ren},\Lambda}=P_{\rm stat}+R_{P,\Lambda},
\qquad
B_\Lambda=B_{\rm stat}+R_{B,\Lambda},
\]
where $P_{\rm stat}$ and $B_{\rm stat}$ are the exact static pieces from Step~1.  Then
\[
\mathcal C_{P,\Lambda}(g)=U_gR_{P,\Lambda}-R_{P,\Lambda}U_g,
\qquad
\mathcal C_{B,\Lambda}(g)=U_gR_{B,\Lambda}-R_{B,\Lambda}U_g.
\]
By construction, $R_{P,\Lambda}$ and $R_{B,\Lambda}$ depend only on the regulated bridge ingredients on the receiver side.  Lemma~\ref{lem:approx-intertwiner-new}
controls the symmetry defect of the regulated conformal intertwiner by
\(
\mathcal D_{\rm conf}(\varepsilon_\Lambda,\Lambda)=O(\varepsilon_\Lambda^\alpha)
\)
uniformly on compact symmetry sets, while Theorem~\ref{thm:tick-fidelity-RP} shows that the accepted temporal-window contribution is bounded by
\(
O(\Phi(\sigma_{t,\Lambda}\Omega_{\rm gap}))
\).
Transporting these bounds through the linearized source-to-boundary map gives, for every compact $G\subset G_{\rm Poin}$ and every compactly supported
test field $\psi$ in the chosen Sobolev class,
\[
\sup_{g\in G}\Big(\|\mathcal C_{P,\Lambda}(g)\psi\|+\|\mathcal C_{B,\Lambda}(g)\psi\|_{H^{s-\frac12}(\partial\mathcal M^{(\Lambda)})}\Big)
\le
C_G\big(\mathcal D_{\rm conf}(\varepsilon_\Lambda,\Lambda)+\Phi(\sigma_{t,\Lambda}\Omega_{\rm gap})\big)\,\|\psi\|_{H^s}.
\]
This proves route item~\textnormal{(P2)}.

\emph{Step 3: pass to the accepted UV limit.}
Corollary~\ref{cor:anti-zeno} states that along the accepted anti-Zeno UV schedule the conformality deficit and the temporal-window defect both vanish.
Therefore the commutator defects of Step~2 converge strongly to zero on compactly supported test sources.  The strong local limits
$P_{\rm gf,ren}$ and $B$ are thus well defined and satisfy
\[
U_g\,P_{\rm gf,ren}=P_{\rm gf,ren}\,U_g,
\qquad
U_g\,B=B\,U_g
\]
for all $g\in G_{\rm Poin}$.  This proves route item~\textnormal{(P3)} and hence the theorem.
\end{proof}

\begin{lemma}[Poincar\'e covariance of Green operators and $\Delta$]
\label{lem:EB4}
Assume Theorem~\ref{thm:uom-Poincare-IBVP} and Green-hyperbolicity (UOM-P2) in the UV limit.  Then the advanced/retarded Green operators commute
with Poincar\'e transformations:
\[
U_g E^\pm \;=\; E^\pm U_g,
\qquad\Rightarrow\qquad
U_g \Delta \;=\; \Delta U_g.
\]
 \end{lemma}

 \begin{proof}
Let $f$ be a test source.  Then $P(E^\pm f)=f$ and $B(E^\pm f)=0$ by definition.
Apply $U_g$ and use covariance:
\[
P(U_g E^\pm f)=U_g P(E^\pm f)=U_g f,
\qquad
B(U_g E^\pm f)=U_g B(E^\pm f)=0.
\]
Thus $U_gE^\pm f$ is a retarded/advanced solution for the source $U_g f$.  By uniqueness of $E^\pm$ (Green-hyperbolicity),
$U_g E^\pm f = E^\pm U_g f$ for all $f$, hence $U_gE^\pm=E^\pm U_g$ and therefore $U_g\Delta=\Delta U_g$.
\end{proof}

\begin{theorem}[Poincar\'e covariance of the Peierls bracket]
\label{thm:Peierls-cov}
Assume UOM-P2--P5 in the UV-limit theory, Theorem~\ref{thm:uom-Poincare-IBVP}, and let $\alpha_g$ be the induced action on $\Fsloc$ by pushforward
of test kernels and fields.  Then for all $F,G\in\Fsloc$,
\[
\alpha_g(\{F,G\}_{\rm P})\;=\;\{\alpha_g F,\alpha_g G\}_{\rm P}.
\]
\end{theorem}

\begin{proof}
By Lemma~\ref{lem:EB4}, $\Delta$ commutes with $U_g$.  Functional derivatives transform covariantly under $\alpha_g$, and the pairing
$\langle\cdot,\cdot\rangle$ is Poincar\'e invariant on the Minkowski chart.  Therefore
\[
\{\alpha_g F,\alpha_g G\}_{\rm P}
=\langle (\alpha_gF)^{(1)},\,\Delta\,(\alpha_gG)^{(1)}\rangle
=\langle U_gF^{(1)},\,\Delta\,U_gG^{(1)}\rangle
=\langle F^{(1)},\,\Delta\,G^{(1)}\rangle
=\alpha_g(\{F,G\}_{\rm P}),
\]
as claimed.
\end{proof}

% ----------------------------------------------------------------------
% 8. Package theorem: EB1--EB3 + covariance
% ----------------------------------------------------------------------

\begin{theorem}[Uniform energy bounds and covariance package (EB1--EB3 + Poincar\'e covariance)]
\label{thm:EB-package}
Assume the UOM Peierls package (UOM-P1--P6) at finite cutoff, the uniformity inputs
Theorems~\ref{thm:uom-OS-inductive}, \ref{thm:uom-H-limit}, \ref{thm:uom-U-Delta}, \ref{thm:uom-U-SD}, and~\ref{thm:uom-Poincare-IBVP}.  Then, on the Minkowski UV chart, the inductive-limit
vacuum representation $(\mathcal H,\pi,\Omega,H)$ satisfies:
\begin{enumerate}[label=\textnormal{(\roman*)},leftmargin=2.2em,itemsep=1pt]
\item \textbf{EB1:} uniform polynomial energy bounds for all $F\in\Fsloc(O)$ (Lemma~\ref{lem:EB1});
\item \textbf{EB2:} uniform Sobolev continuity of $\Delta$ on each double cone $O$ (Lemma~\ref{lem:EB2-uniform});
\item \textbf{EB3:} the Hamiltonized derivation $\delta(F)=\{\mathcal O^{(1)},F\}_{\rm P}$ obeys the vacuum
energy bound \eqref{eq:EB3} on each $O$ (Lemma~\ref{lem:EB3});
\item \textbf{Covariance:} $\Delta$ and the Peierls bracket are Poincar\'e covariant in the UV limit
(Lemma~\ref{lem:EB4} and Theorem~\ref{thm:Peierls-cov}).
\end{enumerate}
\end{theorem}

\paragraph{Interpretation.}
Theorem~\ref{thm:EB-package} is the exact analytic input needed to remove the ``$\delta$ energy bound''
assumption in the implementability step of the two--time Main proof: EB3 implies that the operator
$Q_0\,\pi(A)\Omega:=i\,\pi(\delta(A))\Omega$ is well-defined on the vacuum core and closable, while covariance
ensures the resulting symmetry derivation respects Minkowski kinematics in the UV chart.

\begin{theorem}[UV dilatation package on the Minkowski chart]
\label{thm:uom-uv-dilatation-package}
Under the semiclassical two-time extraction of
Subsection~\ref{subsec:two-times-semiclassical-extraction}, together
with the radial Ward identification of
Proposition~\ref{prop:uom-uv-radial-ward}, the pointer/record structure
of Definition~\ref{def:pointer-algebra-projection}, and the analytic
package of Theorem~\ref{thm:EB-package}, let
\(\mathfrak A_{\mathrm{UV,loc}}\) be the dense UV-local *-subalgebra of
the inductive-limit net generated by the operators \(\Phi(F)\) with
\(F\in\Fsloc(O)\) and \(O\) ranging over double cones in the Minkowski
UV chart, and write \(\mathfrak A_{\mathrm{UV,loc}}(O)\) for the local
subalgebra generated by the corresponding \(\Phi(F)\) with
\(F\in\Fsloc(O)\). Then \(\mathfrak A_{\mathrm{UV,loc}}\) carries a
canonical UV dilatation package, as follows:
\begin{enumerate}[label=\textnormal{(\roman*)},leftmargin=2.2em,itemsep=2pt]
\item there exists a one-parameter group \(\sigma_\tau^{\mathrm{UV}}\) of
algebra automorphisms of \(\mathfrak A_{\mathrm{UV,loc}}\) such that,
for every double cone \(O\),
\[
\sigma_\tau^{\mathrm{UV}}\big(\mathfrak A_{\mathrm{UV,loc}}(O)\big)
=\mathfrak A_{\mathrm{UV,loc}}(e^{-\tau}O),
\]
and on compactly smeared scalar head observables its infinitesimal
generator is the physical radial Ward differential:
\[
L_{\mathrm{RG}}^{\mathrm{UV}}A
:=
\left.\frac{d}{d\tau}\right|_{\tau=0}\sigma_\tau^{\mathrm{UV}}(A)
=\Phi(Q_u^{\mathrm{phys}}X)
\qquad
\text{when }A=\Phi(X)\text{ and }X\text{ is of the scalar compact-smearing form of Proposition~\ref{prop:uom-uv-radial-ward}};
\]
\item \(\sigma_\tau^{\mathrm{UV}}\) intertwines translations by
\[
\sigma_\tau^{\mathrm{UV}}\circ \alpha_a\circ
(\sigma_\tau^{\mathrm{UV}})^{-1}
=\alpha_{e^{-\tau}a},
\]
so that on \(\mathfrak A_{\mathrm{UV,loc}}\)
\begin{equation}
\label{eq:uom-uv-dilatation-derivation}
[L_{\mathrm{RG}}^{\mathrm{UV}},\delta_{P_\mu}]
=-\,\delta_{P_\mu};
\end{equation}
\item in the inductive-limit vacuum representation
\((\mathcal H,\pi,\Omega)\), the automorphisms
\(\sigma_\tau^{\mathrm{UV}}\) are implemented by a strongly continuous
unitary group \(V_\tau=e^{i\tau D_{\mathrm{UV}}}\) fixing the vacuum,
\(V_\tau\Omega=\Omega\), and the selfadjoint generator \(D_{\mathrm{UV}}\)
satisfies
\begin{equation}
\label{eq:uom-uv-dilatation-commutator}
[D_{\mathrm{UV}},P_\mu]=iP_\mu
\end{equation}
on the common Stone core of \(D_{\mathrm{UV}}\) and the translation
generators \(P_\mu\).
\end{enumerate}
\end{theorem}

\begin{proof}
For \(F\in\Fsloc(O)\), write \(F\) in the sufficiently local form
\eqref{eq:Floc-form}.  Let \(s_\tau(x)=e^{-\tau}x\) on the Minkowski UV
chart and define \(\sigma_\tau^{\mathrm{UV}}F\) by pushing forward the
compact smearing kernels and local jets along \(s_\tau\).  Because
\eqref{eq:Floc-form} is stable under smooth coordinate rescaling,
\(\sigma_\tau^{\mathrm{UV}}F\) again belongs to \(\Fsloc(e^{-\tau}O)\).
The group law follows from \(s_{\tau_1}\circ s_{\tau_2}=s_{\tau_1+\tau_2}\),
and multiplicativity plus preservation of adjoints show that
\(\sigma_\tau^{\mathrm{UV}}\) extends from \(\Fsloc\) to an algebra
automorphism of \(\mathfrak A_{\mathrm{UV,loc}}\).

On the compactly smeared scalar head observables singled out in
Subsection~\ref{subsec:two-times-semiclassical-extraction},
Proposition~\ref{prop:uom-uv-radial-ward} identifies the infinitesimal
action of this family with the physical radial Ward differential.  Thus,
for \(A=\Phi(X)\) in that scalar compact-smearing class,
\[
L_{\mathrm{RG}}^{\mathrm{UV}}A
=
\Phi(Q_u^{\mathrm{phys}}X).
\]
This proves item~\textnormal{(i)}.

For translations, let \(T_a(x)=x+a\).  The geometric identity
\(
s_\tau\circ T_a\circ s_\tau^{-1}=T_{e^{-\tau}a}
\)
immediately yields
\[
\sigma_\tau^{\mathrm{UV}}\circ \alpha_a\circ
(\sigma_\tau^{\mathrm{UV}})^{-1}
=\alpha_{e^{-\tau}a}
\]
on the UV-local algebra.  Differentiating at \(\tau=0\) and then at
\(a=t e_\mu\) with \(t=0\) gives
\eqref{eq:uom-uv-dilatation-derivation}, proving item~\textnormal{(ii)}.

It remains to pass to the vacuum representation.  The chosen Minkowski
UV chart is centered at the receiver UV fixed point, and
\(\sigma_\tau^{\mathrm{UV}}\) is exactly the constant boundary
Weyl/rescaling action induced by the radial \(u\)-flow on that chart.
The reference Gaussian vacuum on the UV-local core is therefore
stationary under \(\sigma_\tau^{\mathrm{UV}}\).  Under
Theorem~\ref{thm:uom-OS-inductive}, this fixed vacuum sector is exactly
the one that survives in the inductive limit, so for the vacuum state
\(\omega(A):=\langle\Omega,\pi(A)\Omega\rangle\) one has
\[
\omega(\sigma_\tau^{\mathrm{UV}}A)=\omega(A)
\qquad
(A\in\mathfrak A_{\mathrm{UV,loc}}).
\]
Therefore
\[
V_\tau\,\pi(A)\Omega:=\pi(\sigma_\tau^{\mathrm{UV}}A)\Omega
\]
defines an isometry on the dense set
\(\pi(\mathfrak A_{\mathrm{UV,loc}})\Omega\); since
\((\sigma_\tau^{\mathrm{UV}})^{-1}=\sigma_{-\tau}^{\mathrm{UV}}\), it is
surjective and extends to a unitary \(V_\tau\) on \(\mathcal H\).

Strong continuity follows on the vacuum core from EB1: for each
\(A=\Phi(F)\), the kernels defining \(\sigma_\tau^{\mathrm{UV}}F\) depend
smoothly on \(\tau\), hence continuously in the locality seminorms, and
EB1 converts that seminorm continuity into norm continuity of
\(\pi(\sigma_\tau^{\mathrm{UV}}A)\Omega\).  By density of the core,
\(\tau\mapsto V_\tau\) is strongly continuous on all of \(\mathcal H\).
Stone's theorem therefore gives a selfadjoint generator
\(D_{\mathrm{UV}}\) with \(V_\tau=e^{i\tau D_{\mathrm{UV}}}\).

Finally, the translation intertwining relation from item~\textnormal{(ii)}
becomes
\[
V_\tau U(a)V_\tau^{-1}=U(e^{-\tau}a),
\qquad
U(a):=e^{i a^\mu P_\mu}.
\]
Differentiating first in \(a=t e_\mu\) at \(t=0\) yields
\(
V_\tau P_\mu V_\tau^{-1}=e^{-\tau}P_\mu
\),
and differentiating at \(\tau=0\) gives
\eqref{eq:uom-uv-dilatation-commutator}.  This proves
item~\textnormal{(iii)}.
\end{proof}

\begin{remark}[Role in the limited-scope realization]
\label{rem:uom-uv-dilatation-role}
Theorem~\ref{thm:uom-uv-dilatation-package} is exactly the missing input
behind item \textbf{(A7)} of Definition~\ref{def:uom-limited-scope-package}:
the pointer/record algebra provides the real structure, while the
semiclassical RG/Weyl action now supplies the scaling generator
\(D_{\mathrm{UV}}\) and the commutator
\([D_{\mathrm{UV}},P_\mu]=iP_\mu\) without any extra analytic add-on.
\end{remark}


\section{Aeon evolution}
\label{sec:aeon-evolution}
\subsection*{\texorpdfstring{Welded pulse widths lock near $\Lambda$-band}{Welded pulse widths lock near Lambda-band}}\label{subsec:uom-main-002}
% ------------------------------------------------------------
% Band–edge locking theorem for welded Laplacian–of–Gaussian pulses
% (amplitude cancellation fixed; window only in throughput; corrected prefactors;
% watertight small-width coercivity; optional convexity for uniqueness)
% ------------------------------------------------------------

\begin{theorem}[Band–edge locking of welded LoG widths at receiver cutoff]
\label{thm:band_edge_locking_exact}
Fix a receiver cutoff $\Lambda>0$ and a spatial anisotropy parameter $v>0$.
For widths $s_t>0$ and $s_x>0$, define the welded Laplacian–of–Gaussian (LoG) trace source
\[
\sigma_{s_t,s_x}(\omega,\mathbf k)
\;:=\;A\,\big(\omega^2 e^{-\tfrac{1}{2}s_t^2\omega^2}\big)\,\big(k^2 e^{-\tfrac{1}{2}s_x^2 k^2}\big),
\qquad k:=\|\mathbf k\|,
\]
with real amplitude $A$. Let the receiver acceptance window be
\[
\mathsf W_\Lambda(\omega,\mathbf k)\;:=\;\mathsf W\!\Big(\frac{\omega}{\Lambda},\,\frac{v\,k}{\Lambda}\Big),
\]
satisfying:
\begin{enumerate}
\item[(W1)] $0\le \mathsf W\le 1$, $\mathsf W$ is $C^1$ and nonincreasing in each argument,
$\mathsf W\equiv 1$ on a neighborhood of $(0,0)$, and $\mathsf W(y,z)\to 0$ as $y^2+z^2\to\infty$.
\end{enumerate}
Define the \emph{amplitude–free} throughput
\[
G_0(s_t,s_x;\Lambda)\;:=\;\int_{\mathbb R}\!\!\int_{\mathbb R^3}
\mathsf W_\Lambda(\omega,\mathbf k)\,\Big|\frac{\sigma_{s_t,s_x}(\omega,\mathbf k)}{A}\Big|^2\,
\frac{d\omega\,d^3k}{(2\pi)^4}.
\]
Impose the payload constraint $\Delta I>0$ via
\[
\Delta I\;=\;A^2\,G_0(s_t,s_x;\Lambda)\qquad\Rightarrow\qquad A^2=\frac{\Delta I}{G_0(s_t,s_x;\Lambda)}.
\]

Let $\mathcal K_{\rm conf}(\omega,\mathbf k;\Lambda)$ and $\mathcal K_{\rm BY}(\omega,\mathbf k;\Lambda)$ be the
(real) Fisher/canonical–energy kernels for, respectively, the conformality–mismatch and anomaly (BY/trace) sectors
at receiver cutoff $\Lambda$, and suppose:
\begin{enumerate}
\item[(K1)] (\emph{Positivity \& symmetry}) For each $\Lambda>0$, the kernels are continuous, nonnegative, even in $\omega$,
and spherically symmetric in $\mathbf k$ (hence bounded on a neighborhood of $(\omega,\mathbf k)=(0,\mathbf 0)$).

\item[(K2)] (\emph{UV growth}) There exist exponents $\nu_{\rm conf},\nu_{\rm BY}>0$, constants $c_{\rm conf},c_{\rm BY}>0$
and $Q>0$ such that for all $q^2:=\omega^2+v^2 k^2\ge Q^2$,
\[
\mathcal K_{\rm conf}(\omega,\mathbf k;\Lambda)\;\ge\;c_{\rm conf}\,q^{\nu_{\rm conf}},
\qquad
\mathcal K_{\rm BY}(\omega,\mathbf k;\Lambda)\;\ge\;c_{\rm BY}\,q^{\nu_{\rm BY}}.
\]

\item[(S)] (\emph{Near–cutoff scaling}) There exist exponents $\mu_{\rm conf},\mu_{\rm BY}\in\mathbb R$ and functions
$\kappa_{\rm conf},\kappa_{\rm BY}\in C^1((0,\infty)^2)$, positive and bounded above and below on compact sets, such that
\[
\mathcal K_{\rm conf}(\omega,\mathbf k;\Lambda)
=\Lambda^{\mu_{\rm conf}}\,
\kappa_{\rm conf}\!\Big(\frac{\omega}{\Lambda},\,\frac{v\,k}{\Lambda}\Big),\qquad
\mathcal K_{\rm BY}(\omega,\mathbf k;\Lambda)
=\Lambda^{\mu_{\rm BY}}\,
\kappa_{\rm BY}\!\Big(\frac{\omega}{\Lambda},\,\frac{v\,k}{\Lambda}\Big).
\]
\end{enumerate}
Define the \emph{full} (unwindowed) Fisher numerators
\[
\mathsf M_{\rm full}(s_t,s_x;\Lambda)
:=\int_{\mathbb R}\!\!\int_{\mathbb R^3}
\mathcal K_{\rm conf}(\omega,\mathbf k;\Lambda)\,\Big|\frac{\sigma_{s_t,s_x}(\omega,\mathbf k)}{A}\Big|^2\,
\frac{d\omega\,d^3k}{(2\pi)^4},
\]
\[
\mathsf B_{\rm full}(s_t,s_x;\Lambda)
:=\int_{\mathbb R}\!\!\int_{\mathbb R^3}
\mathcal K_{\rm BY}(\omega,\mathbf k;\Lambda)\,\Big|\frac{\sigma_{s_t,s_x}(\omega,\mathbf k)}{A}\Big|^2\,
\frac{d\omega\,d^3k}{(2\pi)^4}.
\]
The receiver–side real quadratic SK cost (per tick, with amplitude eliminated) is
\begin{equation}
\label{eq:Jrec_def}
\mathcal J^{\rm rec}(s_t,s_x;\Lambda)
:=\lambda_{\rm conf}\,\frac{\Delta I}{G_0(s_t,s_x;\Lambda)}\,\mathsf M_{\rm full}(s_t,s_x;\Lambda)
\;+\;\lambda_{\rm BY}\,\frac{\Delta I}{G_0(s_t,s_x;\Lambda)}\,\frac{1}{s_t\,s_x^3}\,
\mathsf B_{\rm full}(s_t,s_x;\Lambda),
\end{equation}
with weights $\lambda_{\rm conf},\lambda_{\rm BY}>0$. Introduce the dimensionless band–edge variables
\[
\alpha_t:=\Lambda\,s_t,\qquad \alpha_x:=v\,\Lambda\,s_x,
\]
and, with $u:=s_t\,\omega$ and $\mathbf p:=s_x\,\mathbf k$, the dimensionless functions
\[
\Phi(\alpha_t,\alpha_x)
:=\int_{\mathbb R}\!\!\int_{\mathbb R^3}
\mathsf W\!\Big(\frac{u}{\alpha_t},\,\frac{p}{\alpha_x}\Big)\,
u^4\,p^4\,e^{-u^2}e^{-p^2}\,\frac{du\,d^3p}{(2\pi)^4},
\]
\[
\Psi(\alpha_t,\alpha_x)
:=\int_{\mathbb R}\!\!\int_{\mathbb R^3}
\kappa_{\rm conf}\!\Big(\frac{u}{\alpha_t},\,\frac{p}{\alpha_x}\Big)\,
u^4\,p^4\,e^{-u^2}e^{-p^2}\,\frac{du\,d^3p}{(2\pi)^4},
\]
\[
\mathsf B(\alpha_t,\alpha_x)
:=\int_{\mathbb R}\!\!\int_{\mathbb R^3}
\kappa_{\rm BY}\!\Big(\frac{u}{\alpha_t},\,\frac{p}{\alpha_x}\Big)\,
u^4\,p^4\,e^{-u^2}e^{-p^2}\,\frac{du\,d^3p}{(2\pi)^4}.
\]

\noindent\textbf{Conclusion.}
Under (W1), (K1)–(K2), (S), the following hold.

\begin{enumerate}
\item[(i)] \textbf{Factorization.} With $R_{\rm conf}:=\Psi/\Phi$ and $R_{\rm BY}:=\mathsf B/\Phi$,
\[
\boxed{\;
\mathcal J^{\rm rec}(s_t,s_x;\Lambda)
=\Delta I\left[
\lambda_{\rm conf}\,\frac{\Lambda^{\mu_{\rm conf}}}{v^{4}\Lambda^{8}}\,
\alpha_t^{4}\alpha_x^{4}\,R_{\rm conf}(\alpha_t,\alpha_x)
\;+\;
\lambda_{\rm BY}\,\frac{\Lambda^{\mu_{\rm BY}}}{\,v\,\Lambda^{4}}\,
\alpha_t^{3}\alpha_x\,R_{\rm BY}(\alpha_t,\alpha_x)
\right].
\;}
\]
Hence, minimizing $\mathcal J^{\rm rec}$ over $(s_t,s_x)$ at fixed $\Lambda$ is equivalent to minimizing the dimensionless objective
\[
F(\alpha_t,\alpha_x)
:=\lambda_{\rm conf}\,\alpha_t^{4}\alpha_x^{4}\,R_{\rm conf}(\alpha_t,\alpha_x)
\;+\;\lambda_{\rm BY}\,\alpha_t^{3}\alpha_x\,R_{\rm BY}(\alpha_t,\alpha_x).
\]

\item[(ii)] \textbf{Coercivity at large widths.} As $\alpha_t+\alpha_x\to\infty$,
$\mathsf W(u/\alpha_t,p/\alpha_x)\to 1$ pointwise on compacts and $\Phi(\alpha_t,\alpha_x)\to\Phi_\infty\in(0,\infty)$ by dominated convergence. Since $\kappa_{\rm conf},\kappa_{\rm BY}$ are bounded near $(0,0)$ by (K1) and (S), $R_{\rm conf}$ and $R_{\rm BY}$ remain bounded; thus $F(\alpha_t,\alpha_x)\to\infty$ like a positive combination of $\alpha_t^{4}\alpha_x^{4}$ and $\alpha_t^{3}\alpha_x$.

\item[(iii)] \textbf{Coercivity at small widths.} There exist compact sets
\[
U_t:=\{(u,\mathbf p):\ u^2\in[1,2],\ p^2\in[\tfrac14,\tfrac12]\},\qquad
U_x:=\{(u,\mathbf p):\ u^2\in[\tfrac14,\tfrac12],\ p^2\in[1,2]\},
\]
and constants $C_t,C_x>0$ such that, for all sufficiently small $(\alpha_t,\alpha_x)$,
\[
\Psi(\alpha_t,\alpha_x)\ \ge\ C_t\,\alpha_t^{-\nu_{\rm conf}}+C_x\,\alpha_x^{-\nu_{\rm conf}},
\qquad
\mathsf B(\alpha_t,\alpha_x)\ \ge\ C_t'\,\alpha_t^{-\nu_{\rm BY}}+C_x'\,\alpha_x^{-\nu_{\rm BY}},
\]
where the bounds use (K2) on $U_t$ and $U_x$. By (W1), $\Phi(\alpha_t,\alpha_x)\sim C_\Phi\,\alpha_t^{5}\alpha_x^{7}$ as $(\alpha_t,\alpha_x)\to(0,0)$.
Therefore each term of $F$ diverges as $(\alpha_t,\alpha_x)\to(0,0)$, so $F(\alpha_t,\alpha_x)\to\infty$ at small widths.

\item[(iv)] \textbf{Existence and band–edge locking.} By (ii)–(iii), $F$ is coercive on $(0,\infty)^2$. Continuity yields a global minimizer on a compact rectangle $[\underline\alpha,\overline\alpha]^2\subset(0,\infty)^2$. Any minimizer $(s_t^\star,s_x^\star)$ of $\mathcal J^{\rm rec}$ satisfies
\[
\underline\alpha\ \le\ \alpha_t^\star:=\Lambda s_t^\star\ \le\ \overline\alpha,\qquad
\underline\alpha\ \le\ \alpha_x^\star:=v\,\Lambda s_x^\star\ \le\ \overline\alpha,
\]
and hence the welded widths \emph{lock at the band edge}:
\[
\boxed{\quad s_t^\star=\frac{\alpha_t^\star}{\Lambda},\qquad s_x^\star=\frac{\alpha_x^\star}{v\,\Lambda}\quad}
\]
with $\alpha_t^\star,\alpha_x^\star$ independent of $\Lambda$.

\item[(v)] \textbf{Uniqueness and smooth dependence (optional).}
If, on $[\underline\alpha,\overline\alpha]^2$, either
\begin{itemize}
\item[\emph{(a)}] $\nabla\log R_{\rm conf}$, $\nabla\log R_{\rm BY}$ and their Hessians are bounded with sufficiently small norms (small–drift condition), or
\item[\emph{(b)}] $R_{\rm conf}$ and $R_{\rm BY}$ are log–convex,
\end{itemize}
then $F$ is strictly convex in $(\log\alpha_t,\log\alpha_x)$ on the rectangle. The minimizer $(\alpha_t^\star,\alpha_x^\star)$
is unique and depends smoothly on the weights $(\lambda_{\rm conf},\lambda_{\rm BY})$ and the kernels.
\end{enumerate}
\end{theorem}

\begin{remark}
\label{rem:uom-main-005}
In the evolving-wall dynamics, Theorem~\ref{thm:band_edge_locking_exact} is applied
pointwise within each tick with $\Lambda=\bar\Lambda_k$.
\end{remark}

\begin{proof}
We split the proof into five steps that respectively establish items (i)–(v) in the statement.

\medskip
\noindent\textbf{Step 1 (Amplitude cancellation and factorization — item (i)).}
Set
\[
u:=s_t\,\omega,\qquad \mathbf p:=s_x\,\mathbf k,\qquad p:=\|\mathbf p\|,\qquad
\alpha_t:=\Lambda s_t,\qquad \alpha_x:=v\,\Lambda s_x .
\]
Then $d\omega\,d^3k=s_t^{-1}s_x^{-3}\,du\,d^3p$ and
\[
\Big|\frac{\sigma_{s_t,s_x}}{A}\Big|^2
= s_t^{-4}s_x^{-4}\,u^4 p^4\,e^{-u^2}e^{-p^2}\,.
\]
By definition,
\[
G_0(s_t,s_x;\Lambda)
=\int \mathsf W\!\Big(\frac{\omega}{\Lambda},\frac{vk}{\Lambda}\Big)\,
\Big|\frac{\sigma_{s_t,s_x}}{A}\Big|^2\frac{d\omega\,d^3k}{(2\pi)^4}
=s_t^{-5}s_x^{-7}\,\Phi(\alpha_t,\alpha_x),
\]
where
\[
\Phi(\alpha_t,\alpha_x)=\int \mathsf W\!\Big(\frac{u}{\alpha_t},\frac{p}{\alpha_x}\Big)\,
u^4 p^4 e^{-u^2} e^{-p^2}\,\frac{du\,d^3p}{(2\pi)^4}.
\]
Similarly, using (S),
\[
\mathsf M_{\rm full}(s_t,s_x;\Lambda)
=\int \mathcal K_{\rm conf}(\omega,\mathbf k;\Lambda)\,\Big|\frac{\sigma_{s_t,s_x}}{A}\Big|^2\frac{d\omega\,d^3k}{(2\pi)^4}
=s_t^{-1}s_x^{-3}\,\Lambda^{\mu_{\rm conf}}\,\Psi(\alpha_t,\alpha_x),
\]
\[
\mathsf B_{\rm full}(s_t,s_x;\Lambda)
=\int \mathcal K_{\rm BY}(\omega,\mathbf k;\Lambda)\,\Big|\frac{\sigma_{s_t,s_x}}{A}\Big|^2\frac{d\omega\,d^3k}{(2\pi)^4}
=s_t^{-1}s_x^{-3}\,\Lambda^{\mu_{\rm BY}}\,\mathsf B(\alpha_t,\alpha_x),
\]
where
\[
\Psi(\alpha_t,\alpha_x)
=\int \kappa_{\rm conf}\!\Big(\frac{u}{\alpha_t},\frac{p}{\alpha_x}\Big)\,
u^4 p^4 e^{-u^2} e^{-p^2}\,\frac{du\,d^3p}{(2\pi)^4},
\quad
\mathsf B(\alpha_t,\alpha_x)
=\int \kappa_{\rm BY}\!\Big(\frac{u}{\alpha_t},\frac{p}{\alpha_x}\Big)\,
u^4 p^4 e^{-u^2} e^{-p^2}\,\frac{du\,d^3p}{(2\pi)^4}.
\]
Using $\Delta I=A^2 G_0$ in \eqref{eq:Jrec_def} gives
\[
\frac{\Delta I}{G_0}\,\mathsf M_{\rm full}
=\Delta I\,s_t^{4}s_x^{4}\,\Lambda^{\mu_{\rm conf}}\frac{\Psi}{\Phi},\qquad
\frac{\Delta I}{G_0}\,\frac{1}{s_t s_x^3}\,\mathsf B_{\rm full}
=\Delta I\,s_t^{3}s_x^{1}\,\Lambda^{\mu_{\rm BY}}\frac{\mathsf B}{\Phi}.
\]
Finally, $s_t=\alpha_t/\Lambda$ and $s_x=\alpha_x/(v\Lambda)$ yield
\[
\mathcal J^{\rm rec}(s_t,s_x;\Lambda)
=\Delta I\left[
\lambda_{\rm conf}\,\frac{\Lambda^{\mu_{\rm conf}}}{v^{4}\Lambda^{8}}\,
\alpha_t^{4}\alpha_x^{4}\,R_{\rm conf}(\alpha_t,\alpha_x)
\;+\;
\lambda_{\rm BY}\,\frac{\Lambda^{\mu_{\rm BY}}}{\,v\,\Lambda^{4}}\,
\alpha_t^{3}\alpha_x\,R_{\rm BY}(\alpha_t,\alpha_x)
\right],
\]
with $R_{\rm conf}:=\Psi/\Phi$ and $R_{\rm BY}:=\mathsf B/\Phi$. Since the prefactors outside the square brackets do not affect the minimization in $(\alpha_t,\alpha_x)$ at fixed $\Lambda$, item (i) follows.

\medskip
\noindent\textbf{Step 2 (Coercivity at large widths — item (ii), first part).}
By (W1), for each fixed $(u,\mathbf p)$ one has
$\mathsf W(u/\alpha_t,p/\alpha_x)\to 1$ as $(\alpha_t,\alpha_x)\to(\infty,\infty)$, and
$0\le \mathsf W\le 1$ is dominated by $1$. The integrand $u^4 p^4 e^{-u^2}e^{-p^2}$ is integrable on $\mathbb R\times\mathbb R^3$, so by dominated convergence
\[
\Phi(\alpha_t,\alpha_x)\ \xrightarrow[\alpha_t+\alpha_x\to\infty]{}\ \Phi_\infty
:=\int u^4 p^4 e^{-u^2}e^{-p^2}\,\frac{du\,d^3p}{(2\pi)^4}\ \in\ (0,\infty).
\]
By (K1)–(S), $\kappa_{\rm conf}$ and $\kappa_{\rm BY}$ are bounded on a neighborhood of $(0,0)$, hence $R_{\rm conf}$ and $R_{\rm BY}$ remain bounded as $(\alpha_t,\alpha_x)\to\infty$. Therefore the dimensionless objective
\[
F(\alpha_t,\alpha_x)
:=\lambda_{\rm conf}\,\alpha_t^{4}\alpha_x^{4}\,R_{\rm conf}(\alpha_t,\alpha_x)
\;+\;\lambda_{\rm BY}\,\alpha_t^{3}\alpha_x\,R_{\rm BY}(\alpha_t,\alpha_x)
\]
satisfies $F(\alpha_t,\alpha_x)\to+\infty$ as $\alpha_t+\alpha_x\to\infty$.

\medskip
\noindent\textbf{Step 3 (Coercivity at small widths — item (iii)).}
Fix the compact sets
\[
U_t:=\{(u,\mathbf p):\ u^2\in[1,2],\ p^2\in[1/4,1/2]\},\qquad
U_x:=\{(u,\mathbf p):\ u^2\in[1/4,1/2],\ p^2\in[1,2]\}.
\]
On $U_t$ one has $u^2\ge 1$ and $p^2\le 1/2$, hence for $\alpha_t,\alpha_x$ sufficiently small,
\[
q^2:=\Big(\frac{u}{\alpha_t}\Big)^2+\Big(\frac{p}{\alpha_x}\Big)^2
\ \ge\ \frac{u^2}{\alpha_t^2}\ \ge\ \frac{1}{\alpha_t^2}\,,
\qquad
\kappa_{\rm conf}\!\Big(\frac{u}{\alpha_t},\frac{p}{\alpha_x}\Big)
\ \ge\ c'_{\rm conf}\,q^{\nu_{\rm conf}}
\ \ge\ \tilde c_{\rm conf}\,\alpha_t^{-\nu_{\rm conf}}.
\]
Integrating the positive weight $u^4 p^4 e^{-u^2}e^{-p^2}$ over $U_t$ yields
\[
\int_{U_t}\kappa_{\rm conf}\!\Big(\frac{u}{\alpha_t},\frac{p}{\alpha_x}\Big)\,
u^4 p^4 e^{-u^2}e^{-p^2}\,\frac{du\,d^3p}{(2\pi)^4}\ \ge\ C_t\,\alpha_t^{-\nu_{\rm conf}},
\]
for some $C_t>0$ independent of $(\alpha_t,\alpha_x)$. An identical argument on $U_x$ (where $p^2\ge 1$) gives
\[
\int_{U_x}\kappa_{\rm conf}\!\Big(\frac{u}{\alpha_t},\frac{p}{\alpha_x}\Big)\,
u^4 p^4 e^{-u^2}e^{-p^2}\,\frac{du\,d^3p}{(2\pi)^4}\ \ge\ C_x\,\alpha_x^{-\nu_{\rm conf}}.
\]
Summing these lower bounds yields
\[
\Psi(\alpha_t,\alpha_x)\ \ge\ C_t\,\alpha_t^{-\nu_{\rm conf}}+C_x\,\alpha_x^{-\nu_{\rm conf}}.
\]
Exactly the same reasoning applies to $\mathsf B$ with $\nu_{\rm BY}$ in place of $\nu_{\rm conf}$, so
\[
\mathsf B(\alpha_t,\alpha_x)\ \ge\ C_t'\,\alpha_t^{-\nu_{\rm BY}}+C_x'\,\alpha_x^{-\nu_{\rm BY}}.
\]

Next, by (W1) there exists $r>0$ such that $\mathsf W(y,z)=1$ for $y^2+z^2\le r^2$. Hence
\[
\Phi(\alpha_t,\alpha_x)\ \ge\ \int_{|u|\le r\alpha_t}\!\!\!\int_{\|\mathbf p\|\le r\alpha_x}
u^4 p^4 e^{-u^2}e^{-p^2}\,\frac{du\,d^3p}{(2\pi)^4}
\ \ge\ c_\Phi\,\alpha_t^{5}\alpha_x^{7},
\]
for some $c_\Phi>0$. Conversely $\Phi(\alpha_t,\alpha_x)\le \Phi_\infty$ for all $\alpha_t,\alpha_x$.
Combining the bounds for $\Psi,\mathsf B,\Phi$ with the definition of $F$ shows that each term of $F(\alpha_t,\alpha_x)$ diverges as $(\alpha_t,\alpha_x)\to(0,0)$, proving small–width coercivity.

\medskip
\noindent\textbf{Step 4 (Existence and band–edge locking — item (iv)).}
By Steps 2 and 3, $F$ is coercive on $(0,\infty)^2$. By continuity there exists a compact rectangle
$[\underline\alpha,\overline\alpha]^2\subset(0,\infty)^2$ that contains all minimizers, and $F$ attains its global minimum on that rectangle by the Weierstrass theorem. Let $(\alpha_t^\star,\alpha_x^\star)$ be any global minimizer; since the minimization of $\mathcal J^{\rm rec}$ at fixed $\Lambda$ is equivalent to that of $F$, any global minimizer $(s_t^\star,s_x^\star)$ of $\mathcal J^{\rm rec}$ satisfies
\[
\alpha_t^\star=\Lambda s_t^\star\in[\underline\alpha,\overline\alpha],\qquad
\alpha_x^\star=v\,\Lambda s_x^\star\in[\underline\alpha,\overline\alpha].
\]
Equivalently,
\[
s_t^\star=\frac{\alpha_t^\star}{\Lambda},\qquad s_x^\star=\frac{\alpha_x^\star}{v\,\Lambda},
\]
with $\alpha_t^\star,\alpha_x^\star$ independent of $\Lambda$, which is the claimed band–edge locking.

\medskip
\noindent\textbf{Step 5 (Optional uniqueness and smooth dependence — item (v)).}
Write $u=\log\alpha_t$ and $w=\log\alpha_x$. Then
\[
F(u,w)=\lambda_{\rm conf}\,e^{4u+4w}\,R_{\rm conf}(e^{u},e^{w})
\;+\;\lambda_{\rm BY}\,e^{3u+w}\,R_{\rm BY}(e^{u},e^{w}).
\]
On the compact rectangle from Step 4, assume either:

\emph{(a) Small–drift condition.} There exist bounds
\[
\|\nabla\log R_{\rm conf}\|,\ \|\nabla^2\log R_{\rm conf}\|,\ 
\|\nabla\log R_{\rm BY}\|,\ \|\nabla^2\log R_{\rm BY}\|\ \le\ \varepsilon,
\]
with $\varepsilon>0$ sufficiently small. Then the Hessian $\nabla^2 F(u,w)$ equals the sum of two positive semidefinite rank–one matrices,
\[
\lambda_{\rm conf}\,e^{4u+4w}R_{\rm conf}\,\begin{bmatrix}4\\[2pt]4\end{bmatrix}\!\begin{bmatrix}4&4\end{bmatrix}
\;+\;
\lambda_{\rm BY}\,e^{3u+w}R_{\rm BY}\,\begin{bmatrix}3\\[2pt]1\end{bmatrix}\!\begin{bmatrix}3&1\end{bmatrix},
\]
plus perturbative terms controlled by $\varepsilon$. Since the two exponent vectors $(4,4)$ and $(3,1)$ are linearly independent, the sum of the two rank–one positive semidefinite matrices is positive definite; for $\varepsilon$ small enough the perturbations preserve positive definiteness. Hence $F$ is strictly convex in $(u,w)$, so the minimizer is unique and depends smoothly on the parameters by the implicit–function theorem.

\emph{(b) Log–convex weights.} If $R_{\rm conf}$ and $R_{\rm BY}$ are log–convex on the rectangle (i.e., $\nabla^2\log R\succcurlyeq 0$), then $(u,w)\mapsto e^{au+bw}R(e^u,e^w)$ is convex for any $a,b\in\mathbb R$. The sum of the two convex terms that carry linearly independent exponent directions is strictly convex, again implying uniqueness and smooth dependence.

\medskip
This completes the proof.
\end{proof}

\begin{remark}[KR and interior RS1 applicability]
\label{rem:uom-main-006}
For the active receiver we assume a KR head throughout the split-ordered evolution, so the band--edge locking applies in KR
throughout the aeon. RS1 enters only through interior/shadow segments; in those segments the same scaling
holds above the IR gap and reduces smoothly to the KR case.
\end{remark}

% ------------------------------------------------------------
% End of theorem
% ------------------------------------------------------------

\subsection{\texorpdfstring{Split-side ordering under welded pulse widths lock at $\Lambda$–band}{Split-side ordering under welded pulse widths lock at Lambda-band}}
\label{subsec:init-locked}

\paragraph{Scope and standing locks.}\label{par:uom-main-136}
Throughout this subsection we impose the \emph{receiver locks} and \emph{health} assumptions already established:
(i) the receiver brane is a healthy \emph{KR} head (no ghosts/tachyons; standard holographic counterterms);
interior/shadow RS1 segments may exist but are not active receivers,
(ii) the receiver Gaussian sector satisfies OS–positivity; in particular the RP block is absent at quadratic order,
(iii) the welded (trace) source is a Laplacian–of–Gaussian (LoG) pulse with \emph{band–locked} widths
\[
s_t^\star=\frac{\alpha_t^\star}{\Lambda},\qquad s_x^\star=\frac{\alpha_x^\star}{v\,\Lambda},\qquad
\alpha_t^\star,\alpha_x^\star=O(1),
\]
The admission/acceptance window used in defining the admitted payload is introduced in
Sec.~\ref{sec:two-times-semiclassical}.
(iv) the \emph{split-side} cutoff is fixed to the \emph{minimal feasible} value $\Lambda_{\min}$
(Def.~\ref{def:Lambda-min}),
(v) the previous–aeon source bandlimit is selected to the light radion scale (see \S\ref{subsubsec:init-bandlimit}).

\begin{definition}[Minimal split-side cutoff]\label{def:Lambda-min}
We define the minimal feasible split-side cutoff by
\[
\Lambda_{\min}\ :=\ \max\!\Big(\Lambda_{\rm phys},\ (1+\varepsilon)\,\Lambda_{\rm S}(t_{\rm split})\Big),
\qquad \varepsilon>0\ \text{small},
\]
so the continuing receiver sits strictly above the newborn source in RG ordering.
In the radial gauge where $\Lambda\propto e^{A(y)}$, this implies $\Lambda_{\rm R}>\Lambda_{\rm S}$ iff
$y_{\rm R}<y_{\rm S}$ (the split leaves the receiver slightly UV-side of the source).
\end{definition}

We write the receiver AdS$_4$ background curvature as $R_0=4\Lambda_4=-12/L_4^2$ and use
the \emph{graviton gap} $m_g:=L_4^{-1}\ge 0$ for convenience. The induced–gravity screen is
$r(\Lambda)=M_4^2(\Lambda)/(2M_5^3)$ and we denote by $W_2(k/\Lambda)$ the receiver UV window.
On welded band–edge modes $(\omega,k)\sim (\alpha_t^\star\Lambda,\alpha_x^\star\Lambda/v)$ the trace
kernel entering the Fisher quadratic has the screened/Yukawa form
\begin{equation}
\label{eq:trace-kernel}
\mathcal K(\omega,k;\Lambda,L_4)\;=\;\frac{q(\omega,k;L_4)}{q(\omega,k;L_4)+r(\Lambda)k^2}\;W_2\!\Big(\frac{k}{\Lambda}\Big),
\qquad
q(\omega,k;L_4)\simeq \sqrt{k^2+m_g^2}\,.
\end{equation}

\subsubsection*{(a) Split-side curvature from the compensator--driven argmin}
\label{subsubsec:init-curvature}
The real (Gaussian) 5D SK influence functional on the receiver factorizes into the $\Theta$–odd (RP),
Weyl/trace, and conformal–mismatch blocks. Under the standing locks, the RP block vanishes and
\begin{equation}
\label{eq:ReWrec}
\Re W_{\rm rec}\;=\;\underbrace{\big\|\Pi_{\rm Weyl}\,\delta C\big\|_{F}^2}\_{\text{Weyl/trace}}
\;+\;\underbrace{\big\|\Pi_{\rm conf}\,\delta C\big\|_{F}^2}\_{\text{conformal mismatch}}\!,
\end{equation}
where $\delta C$ is the welded covariance tangent at the band edge and $\|\cdot\|_F$ is the Fisher norm.
The Weyl block equals the Riegert/compensator quadratic (after integrating out $\sigma$),
$\propto \langle\sigma,\Delta_4\sigma\rangle$, and the conformal block is the (regulated)
intertwiner mismatch, both evaluated on the AdS$_4$ background with gap $m_g$.

\begin{definition}[Split-side feasible set]
\label{def:feasible-set}
Let $\mathcal F_{\rm init}\subset \R_{\ge 0}\times \R_{>0}$ be the set of $(m_g,\Lambda)$ satisfying:
\begin{enumerate}[label=(\roman*)]
    \item AdS branch (Israel) $m_g\ge 0$; 
    \item health: no ghosts/tachyons; stabilized radion; tempered, reflection–symmetric
two–point;
    \item OS positivity at window $\mu$;
    \item  minimal–feasible cutoff $\Lambda=\Lambda_{\min}$;
    \item optional
control granularity (if brane tension is quantized, $m_g\ge m_g^{\min}>0$; otherwise $m_g^{\min}=0^+$).
\end{enumerate}

\end{definition}

\begin{lemma}[Monotonicity of the receiver Fisher blocks in $m_g$]\label{lem:monotone-mg}
Fix the receiver cutoff $\Lambda>0$ and the welded (trace) source with band–locked widths
$s_t=\alpha_t^{*}/\Lambda$, $s_x=\alpha_x^{*}/(v\Lambda)$, where $\alpha_t^{*},\alpha_x^{*}>0$.
Let the (screened) trace kernel on welded band–edge modes be
\[
\mathcal K(\omega,k;\Lambda,m_g)
=\frac{q(\omega,k;m_g)}{q(\omega,k;m_g)+r(\Lambda)k^2}\,W_2\!\Big(\frac{k}{\Lambda}\Big),
\qquad
q(\omega,k;m_g):=\sqrt{k^2+m_g^2},
\]
with $r(\Lambda)>0$ (induced–gravity screen) and a receiver window $W_2\in[0,1]$. Denote by
$\Pi_{\rm Weyl}$, $\Pi_{\rm conf}$ the Fisher–orthogonal projectors onto the Weyl/trace and conformal–mismatch
tangents at the AdS$_4$ fixed point. Then, evaluated on the welded band–edge family,
\[
\frac{\partial}{\partial m_g}\,\big\|\Pi_{\rm Weyl}\,\delta C\big\|_{F}^2\ >\ 0,
\qquad
\frac{\partial}{\partial m_g}\,\big\|\Pi_{\rm conf}\,\delta C\big\|_{F}^2\ >\ 0,
\]
for all $m_g\ge 0$. In particular, both Fisher blocks are strictly increasing functions of $m_g$.
\end{lemma}

\begin{proof}
\emph{Step 1 (Welded linear response and diagonalization).}
Let $\tau$ be the welded trace source (LoG in time \& space with band–locked widths), and let
$T_{m_g}:\tau\mapsto \delta C$ denote the receiver–side linear response map at fixed cutoff $\Lambda$.
In the Fourier variables $\xi=(\omega,\mathbf k)$ the welded family is diagonal, and $T_{m_g}$ acts
by \emph{multiplication} with the scalar symbol $\mathcal K(\xi;\Lambda,m_g)$, up to a fixed
mode–wise positive operator $M(\xi)$ that encodes the LoG envelope and normalization:
\[
\widehat{\delta C}(\xi)\;=\;T_{m_g}\widehat{\tau}(\xi)\;=\;\mathcal K(\xi;\Lambda,m_g)\,M(\xi)\,\widehat{\tau}(\xi).
\]
The Fisher norm of a tangent is $\|\delta C\|_F^2=\frac12\mathrm{tr}(C_\ast^{-1}\delta C\,C_\ast^{-1}\delta C)$.
On the welded band, the fixed–point covariance $C_\ast$ is diagonal in $\xi$ and strictly positive; hence
$\|\Pi\,\delta C\|_F^2$ is a \emph{quadratic form} in $\delta C$ that reduces to an integral of the mode–wise
norm squared with a positive weight. More precisely, for either projector $\Pi\in\{\Pi_{\rm Weyl},\Pi_{\rm conf}\}$
there exists a positive kernel $H_\Pi(\xi)>0$ (smooth on the welded band) such that
\begin{equation}\label{eq:Fisher-mode}
\big\|\Pi\,\delta C\big\|_F^2\;=\;\int_{\R\times\R^3}\! H_\Pi(\xi)\;\Big|\mathcal K(\xi;\Lambda,m_g)\Big|^2\,
\Big\|M(\xi)\,\widehat{\tau}(\xi)\Big\|^2\,\frac{d\omega\,d^3k}{(2\pi)^4}.
\end{equation}
(Informally: $H_\Pi$ is obtained by sandwiching the mode–wise projector and the fixed $C_\ast^{-1}$ weights;
positivity follows because $C_\ast^{-1}\succ0$ and $\Pi$ is an orthogonal projector in the Fisher inner product.)

\emph{Step 2 (Derivative of the scalar symbol with respect to $m_g$).}
Fix $(\omega,k)$ and write $a:=r(\Lambda)k^2\ge 0$, $W:=W_2(k/\Lambda)\in[0,1]$. Then
\[
\mathcal K(\omega,k;\Lambda,m_g)=W\cdot f(q),\qquad f(q):=\frac{q}{q+a},\quad q=\sqrt{k^2+m_g^2}\,.
\]
For $k>0$ one has $q>0$ and
\[
f'(q)=\frac{a}{(q+a)^2}\ >\ 0,\qquad \frac{\partial q}{\partial m_g}=\frac{m_g}{q}\ \ge\ 0,
\]
hence
\begin{equation}\label{eq:K-derivative}
\frac{\partial}{\partial m_g}\,\mathcal K(\omega,k;\Lambda,m_g)
=W\cdot f'(q)\cdot \frac{\partial q}{\partial m_g}\;=\;W\cdot \frac{a}{(q+a)^2}\cdot \frac{m_g}{q}\ \ge\ 0,
\end{equation}
with strict inequality for all $k>0$ and $m_g>0$. For $k=0$ one has $a=0$ and $\mathcal K(\cdot)=W\equiv W_2(0)$
is constant in $m_g$, hence the derivative vanishes at $k=0$. (This set has measure zero and, moreover,
for the welded LoG envelope the mode weight carries a factor $k^4$ so $k=0$ does not contribute.)

\emph{Step 3 (Differentiate the Fisher blocks).}
Differentiate \eqref{eq:Fisher-mode} under the integral sign (justified by smoothness/temperedness and dominated
convergence on the welded band):
\[
\frac{\partial}{\partial m_g}\,\big\|\Pi\,\delta C\big\|_F^2
=\int H_\Pi(\xi)\,2\,\mathcal K(\xi;\Lambda,m_g)\,\frac{\partial\mathcal K(\xi;\Lambda,m_g)}{\partial m_g}\,
\Big\|M(\xi)\,\widehat{\tau}(\xi)\Big\|^2\,\frac{d\omega\,d^3k}{(2\pi)^4}.
\]
On the welded band $\mathcal K\ge 0$, $H_\Pi>0$, and $\|M\,\widehat{\tau}\|^2\ge 0$. By \eqref{eq:K-derivative},
$\partial_{m_g}\mathcal K\ge 0$ for all modes, with strict positivity for all $k>0$ and $m_g>0$. Since the
LoG envelope localizes the spatial frequency near a positive band edge ($k\sim \alpha_x^{*}\Lambda/v>0$),
the integrand is strictly positive on a set of nonzero measure. Therefore
\[
\frac{\partial}{\partial m_g}\,\big\|\Pi\,\delta C\big\|_F^2\ >\ 0\qquad(\Pi=\text{Weyl, conf}),
\]
for all $m_g\ge 0$. This proves the claim.
\end{proof}

\begin{theorem}[Compensator--driven split-side argmin (near--critical continuing receiver)]\label{thm:init-argmin}
Fix the receiver cutoff to the minimal feasible value $\Lambda_{\min}>0$ (Def.~\ref{def:Lambda-min};
determined by acceptance at a window $\mu>0$ and by the retarded knee) and impose the welded band–locks
$s_t^{*}=\alpha_t^{*}/\Lambda_{\min}$,
$s_x^{*}=\alpha_x^{*}/(v\Lambda_{\min})$ with $\alpha_t^{*},\alpha_x^{*}>0$. Assume the receiver head is
healthy (KR type): no ghosts/tachyons; standard counterterms so the
Gaussian two–point is tempered and reflection–symmetric. Let the feasible set be
\[
\mathcal F_{\rm init}\;:=\;\big\{\,m_g\ge 0:\ \text{(AdS branch, health, welded acceptance at }\mu\text{)}\big\},
\]
with $m_g=1/L_4$ the receiver graviton gap. Consider the receiver–side real SK functional
\[
\Re W_{\rm rec}(m_g)\;=\;\big\|\Pi_{\rm Weyl}\,\delta C(m_g)\big\|_{F}^2
+\big\|\Pi_{\rm conf}\,\delta C(m_g)\big\|_{F}^2,
\]
evaluated on the welded band–edge family at cutoff $\Lambda_{\min}$ (the RP block vanishes by OS positivity).
Then $\Re W_{\rm rec}$ is strictly increasing in $m_g$ on $\mathcal F_{\rm init}$ and attains its minimum at the
\emph{lower boundary} of the feasible set:
\[
m_g^\star\;=\;\inf \mathcal F_{\rm init}\,.
\]
In particular, if brane--tension control is continuous so that $\inf\mathcal F_{\rm init}=0^+$ (AdS
branch approached from below), the continuing receiver sits \emph{near--critical} at the split:
\[
|R_0|\;=\;\frac{12}{L_4^2}\;=\;12\,m_g^2\;\downarrow\; 0\qquad \text{at aeon entry.}
\]
\end{theorem}

\begin{proof}
\emph{Step 1 (Monotonicity).}
By Lemma~\ref{lem:monotone-mg}, for each block (Weyl/trace and conformal mismatch) one has
$\partial_{m_g}\|\Pi\,\delta C(m_g)\|_{F}^2>0$ at fixed $\Lambda_{\min}$ (under welded band–locks),
hence $\partial_{m_g}\Re W_{\rm rec}(m_g)>0$ for all $m_g\in\mathcal F_{\rm init}$.

\emph{Step 2 (Continuity and existence of the minimum).}
The map $m_g\mapsto \Re W_{\rm rec}(m_g)$ is continuous on $\mathcal F_{\rm init}$: $\delta C(m_g)$ depends
continuously on $m_g$ through the scalar symbol $\mathcal K(\cdot;\Lambda_{\min},m_g)$, the Fisher norm is
continuous in $\delta C$, and the welded band–edge integrals are uniformly dominated by an integrable weight
($\omega^4 k^4 e^{-s_t^{{*}2}\omega^2}e^{-s_x^{{*}2}k^2}$). Therefore, by the Weierstrass minimum theorem on the closure
of any bounded subinterval of $\mathcal F_{\rm init}$, the strictly increasing continuous function $\Re W_{\rm rec}$
attains its minimum at the \emph{left endpoint} of the interval.

\emph{Step 3 (Lower boundary of the feasible set).}
Let $m_g^{\min}:=\inf \mathcal F_{\rm init}\ge 0$ (the \emph{lower boundary}). By Step 1,
$\Re W_{\rm rec}$ is strictly increasing, hence
\[
\arg\min_{m_g\in \mathcal F_{\rm init}}\,\Re W_{\rm rec}(m_g)\ =\ \{\,m_g^{\min}\,\}.
\]
Thus $m_g^\star=m_g^{\min}$, i.e. the split-side argmin sits at the smallest admissible gap.

\emph{Step 4 (Near–critical brane under continuous control).}
If the brane--tension control is continuous (no quantization), the AdS branch permits approaching
criticality from below, hence $m_g^{\min}=0^+$. Therefore the continuing receiver sits at a \emph{near--critical}
curvature $|R_0|=12 m_g^2\downarrow 0$. (If tension were quantized, $m_g^{\min}>0$ would be the smallest
attainable gap and the argmin would pick that value.)

This completes the proof.
\end{proof}

\subsubsection*{(b) Split-time source bandlimit (retarded knee)}
\label{subsubsec:init-bandlimit}
The previous aeon’s source bandlimit $\Delta$ at the split is set by \emph{retarded} brane–to–brane coupling
into the (fixed) receiver band $\Lambda_{\min}$ (Def.~\ref{def:Lambda-min}) with welded widths locked. This
retarded knee is the operational instantiation of the acceptance/windowing map introduced in
Sec.~\ref{sec:two-times-semiclassical}. Let
\begin{equation}
\label{eq:ret-kernel}
\varepsilon(\omega,k;\Lambda_{\min})
\;\simeq\; \exp\!\big[-\,2d\,\sqrt{k^2+m_{\rm rad}^2}\,\big]\,
\frac{\sqrt{k^2+m_{\rm rad}^2}}{\sqrt{k^2+m_{\rm rad}^2}+r(\Lambda_{\min})\,k^2}\,
W_2\!\Big(\frac{k}{\Lambda_{\min}}\Big)
\end{equation}
be the retarded trace–channel transfer (zero tilt), where $d$ is the proper inter–brane distance,
$m_{\rm rad}$ the (stabilized) radion mass, and $r(\Lambda_{\min})$ the induced–gravity screen.

\begin{proposition}[Retarded knee locking of the split-time source bandlimit]\label{prop:knee-lock}
Fix the receiver cutoff to its minimal feasible value $\Lambda_{\min}>0$ (Def.~\ref{def:Lambda-min}; acceptance window $\mu>0$ and welded
widths $s_t=\alpha_t^{*}/\Lambda_{\min}$, $s_x=\alpha_x^{*}/(v\Lambda_{\min})$ held fixed), and let $L_4^{\rm(prev)}>0$
be the previous–aeon AdS$_4$ radius (source side). Write $k_{\rm edge}:=\Delta/L_4^{\rm(prev)}$ for the \emph{source}
band edge at bandlimit $\Delta\ge 0$. Let the \emph{retarded transfer} in the trace channel be
\begin{equation}\label{eq:eps-knee}
\varepsilon(\omega,k;\Lambda_{\min})
\;=\;e^{-\,2d\,q(k)}\,\frac{q(k)}{q(k)+r(\Lambda_{\min})\,k^2}\,W_2\!\Big(\frac{k}{\Lambda_{\min}}\Big),\qquad
q(k):=\sqrt{k^2+m_{\rm rad}^2},
\end{equation}
with $d>0$ the inter–brane separation, $m_{\rm rad}\ge 0$ a stabilized radion mass, $r(\Lambda_{\min})>0$ the induced–gravity
screen, and $W_2\in C^1([0,\infty))$ a nonincreasing UV window with $W_2(0)=1$ and $0\le W_2\le 1$. Let
$\rho_{\rm src}(k)$ be a continuous, nonincreasing \emph{marginal value density} (exported recoverable information per unit band
at the welded shapes) and $\pi_{\rm src}(k)$ a continuous, strictly increasing \emph{marginal source cost} (QSL\,+\,dwell\,+\,guards)
per unit band. Consider the amplitude–eliminated source objective
\begin{equation}\label{eq:J-Delta}
\mathcal J(\Delta)\;:=\;\int_{0}^{k_{\rm edge}}\!\Big[\pi_{\rm src}(k)\ -\ \rho_{\rm src}(k)\,\big|\varepsilon(\omega^{*},k;\Lambda_{\min})\big|^2\Big]\,dk,\qquad 
k_{\rm edge}=\frac{\Delta}{L_4^{\rm(prev)}}.
\end{equation}
Then:
\begin{enumerate}
\item[(i)] $\mathcal J$ is $C^1$ and strictly \emph{convex} in $\Delta$.
\item[(ii)] Any minimizer $\Delta_0\ge 0$ satisfies the KKT stationarity
\begin{equation}\label{eq:KKT-stationarity}
\rho_{\rm src}(k_\star)\,\big|\varepsilon(\omega^{*},k_\star;\Lambda_{\min})\big|^2\ =\ \pi_{\rm src}(k_\star),\qquad k_\star:=\frac{\Delta_0}{L_4^{\rm(prev)}}.
\end{equation}
The solution $k_\star$ is unique, and (under the conditions below) lies parametrically at the \emph{retarded knee}
\begin{equation}\label{eq:kstar-min}
k_\star\ \sim\ \min\!\Big\{\,m_{\rm rad},\ d^{-1},\ r(\Lambda_{\min})^{-1/2}\,\Big\},
\end{equation}
with implied $O(1)$ constants depending only on $(\rho_{\rm src},\pi_{\rm src},W_2)$.
\end{enumerate}
\end{proposition}

\begin{proof}
\emph{Standing assumptions.} Throughout we assume:
(a) $\rho_{\rm src}\in C^0([0,\infty))$ is nonincreasing, bounded, and $\rho_{\rm src}(0)>0$;
(b) $\pi_{\rm src}\in C^0([0,\infty))$ is strictly increasing, $\pi_{\rm src}(0)\ge 0$, and
there exists $c_0>0$ s.t.\ $\pi_{\rm src}'(k)\ge c_0$ for all $k$ large enough (or more generally, $\lim_{k\to\infty}\pi_{\rm src}(k)=\infty$);
(c) $W_2$ as above; (d) $q(k)$ and the screening factor are as in \eqref{eq:eps-knee} with $r(\Lambda_{\min})>0$.

\smallskip
\noindent\emph{Step 1 (Basic properties of the retarded transfer).}
Define
\[
E(k):=e^{-\,2d\,q(k)}\in(0,1],\quad S(k):=\frac{q(k)}{q(k)+r\,k^2}\in(0,1],\quad r:=r(\Lambda_{\min})>0,\quad W(k):=W_2\!\Big(\frac{k}{\Lambda_{\min}}\Big)\in(0,1],
\]
so that $|\varepsilon|^2=E(k)^2 S(k)^2 W(k)^2$. Note $q\in C^\infty$, $q'(k)=k/q(k)\ge 0$, hence $E$ is strictly
decreasing in $k$ with $E(0)=e^{-2d m_{\rm rad}}$ and $E(k)\to 0$ as $k\to\infty$. Also,
\begin{equation}\label{eq:S-derivative}
S'(k)=\frac{q'(k)\,(q(k)+rk^2)-q(k)\,(q'(k)+2rk)}{(q(k)+rk^2)^2}
=\frac{q'(k)\,rk^2 - 2rk\,q(k)}{(q(k)+rk^2)^2}
=\frac{rk}{(q+rk^2)^2}\Big(\frac{k^2}{q}-2q\Big)\ \le\ 0,
\end{equation}
since $2q\ge 2\,k$ and $k^2/q - 2q \le k^2/k - 2k\le 0$ (with equality only at $k=0$). Therefore $S$ is
strictly decreasing with $S(0)=1$ and $S(k)\to 0$ as $k\to\infty$. By (c), $W$ is nonincreasing with $W(0)=1$,
$W(k)\to 0$ as $k\to\infty$. Consequently,
\begin{equation}\label{eq:eps-monotone}
k\mapsto |\varepsilon(\omega^{*},k;\Lambda_{\min})|^2=E(k)^2 S(k)^2 W(k)^2\quad\text{is strictly decreasing on }[0,\infty).
\end{equation}

\smallskip
\noindent\emph{Step 2 (Convexity and KKT stationarity).}
Define $g(k):=\pi_{\rm src}(k)-\rho_{\rm src}(k)|\varepsilon|^2$. Then $g\in C^0([0,\infty))$ and, by the assumptions and
\eqref{eq:eps-monotone}, $g'(k)=\pi_{\rm src}'(k)-\rho_{\rm src}'(k)\,|\varepsilon|^2 - \rho_{\rm src}(k)\,\partial_k|\varepsilon|^2\ge 0$
for all $k$ (the last term is strictly positive for $k>0$). In particular $g$ is nondecreasing and, beyond some $k_1$,
strictly increasing. The objective \eqref{eq:J-Delta} is $\mathcal J(\Delta)=\int_0^{k_{\rm edge}} g(k)\,dk$ with $k_{\rm edge}=\Delta/L_4^{\rm(prev)}$, hence
\[
\mathcal J'(\Delta)\;=\;\frac{1}{L_4^{\rm(prev)}}\,g\!\left(\frac{\Delta}{L_4^{\rm(prev)}}\right),\qquad
\mathcal J''(\Delta)\;=\;\frac{1}{L_4^{\rm(prev)2}}\,g'\!\left(\frac{\Delta}{L_4^{\rm(prev)}}\right)\ \ge\ 0.
\]
Thus $\mathcal J$ is convex and strictly convex on any interval where $g'>0$. Any minimizer of $\mathcal J$ satisfies the KKT condition $\mathcal J'(\Delta_0)=0$ (interior) or the complementary slackness at the boundary (if the minimizer occurs at $\Delta_0=0$); since $L_4^{\rm(prev)}>0$, interior stationarity reads
\[
g(k_\star)\ =\ 0\qquad\Longleftrightarrow\qquad \rho_{\rm src}(k_\star)\,|\varepsilon(\omega^{*},k_\star;\Lambda_{\min})|^2=\pi_{\rm src}(k_\star),
\]
which is \eqref{eq:KKT-stationarity}. Because $g$ is nondecreasing, any solution is unique; existence follows from the next step.

\smallskip
\noindent\emph{Step 3 (Existence and uniqueness).}
At $k=0$, $g(0)=\pi_{\rm src}(0)-\rho_{\rm src}(0)\,|\varepsilon(0)|^2$, with $|\varepsilon(0)|^2=e^{-2dm_{\rm rad}}\in(0,1]$. We assume the regime where $g(0)<0$, i.e.\ the marginal export value exceeds the marginal source cost at the very first mode (this is the empirically relevant “payload–limited” start). Since $g$ is continuous and $\lim_{k\to\infty} g(k)=+\infty$ (because $|\varepsilon|^2\to 0$ while $\pi_{\rm src}\to\infty$), the intermediate value theorem guarantees a root $k_\star>0$, which is unique by monotonicity. Thus the KKT stationarity has a unique solution \eqref{eq:KKT-stationarity}.

\smallskip
\noindent\emph{Step 4 (Knee asymptotics and parametric location).}
We show that $k_\star$ lies at the minimum of the three \emph{retarded scales}
\[
k_{\rm rad}:=m_{\rm rad},\qquad k_d:=d^{-1},\qquad k_r:=r(\Lambda_{\min})^{-1/2},
\]
up to $O(1)$ factors depending only on $(\rho_{\rm src},\pi_{\rm src},W_2)$.
Write $h(k):=\rho_{\rm src}(k)\,|\varepsilon|^2$ and $c(k):=\pi_{\rm src}(k)$. By Step~1, $h$ is strictly decreasing.
We will bracket $k_\star$ by comparing $h$ and $c$ in three $k$–regimes.

\smallskip
\noindent\emph{(a) Small–$k$ regime: $k\le \min\{k_{\rm rad},k_d,k_r\}$.}
Here $q(k)\approx m_{\rm rad}$, $E(k)\approx e^{-2dm_{\rm rad}}$, $S(k)\approx 1$, and $W(k)\approx 1$,
so $h(k)\approx \rho_{\rm src}(k)\,e^{-4dm_{\rm rad}}\ge \rho_{\rm src}(0)\,e^{-4dm_{\rm rad}} - \|\rho_{\rm src}'\|_\infty k$,
while $c(k)\approx c(0) + (\inf c')\,k$. Our standing assumption gives $g(0)=c(0)-\rho_{\rm src}(0)e^{-4dm_{\rm rad}}<0$,
so for sufficiently small $k$ (in particular for all $k\le k_{\rm ref}$ with $k_{\rm ref}=O(1)$) we have $h(k)>c(k)$, thus $g(k)<0$.
This shows $k_\star$ cannot lie \emph{strictly below} all three scales.

\smallskip
\noindent\emph{(b) Large–$k$ regime: $k\ge \max\{k_{\rm rad},k_d\}$.}
For $k\ge k_{\rm rad}$, $q(k)\ge k$ and $E(k)\le e^{-2dk}$; for $k\ge k_d$, $E(k)\le e^{-2}$. Thus if $k\ge \max\{k_{\rm rad},k_d\}$ then $E(k)\le e^{-2dk}\le e^{-2}$ and decays exponentially in $k$, independent of $m_{\rm rad}$. Since $S(k)\le 1$ and $W(k)\le 1$, we have
\begin{equation}\label{eq:h-exp-decay}
h(k)\ \le\ \rho_{\rm src}(0)\,E(k)^2\ \le\ \rho_{\rm src}(0)\,e^{-4dk}\quad\text{for }k\ge \max\{k_{\rm rad},k_d\}.
\end{equation}
Because $c(k)\to +\infty$ and $h(k)\to 0$, there exists a constant $C_1=O(1)$ (depending on $\rho_{\rm src},\pi_{\rm src}$) such that $h(k)<c(k)$ for all $k\ge C_1\,\max\{k_{\rm rad},k_d\}$. Therefore $k_\star\le C_1\,\max\{k_{\rm rad},k_d\}$.

\smallskip
\noindent\emph{(c) Screening regime: $k\ge k_r$.}
For $k\ge k_r=r^{-1/2}$ we use $S(k)=q/(q+rk^2)\le q/(rk^2)\le (k+m_{\rm rad})/(rk^2)\le (2/rk)$ and thus
\begin{equation}\label{eq:h-screen-decay}
h(k)\ \le\ \rho_{\rm src}(0)\,E(k)^2\,\frac{4}{r^2 k^2}\ \le\ \frac{4\,\rho_{\rm src}(0)}{r^2 k^2}\,,
\end{equation}
since $E\le 1$. As $k\to\infty$, the right side decays like $k^{-2}$; by the assumption that $c(k)$ grows (at least linearly eventually), there exists a constant $C_2=O(1)$ such that $h(k)<c(k)$ for all $k\ge C_2\,k_r$. Hence $k_\star\le C_2\,k_r$.

\smallskip
Combining (b) and (c) shows that $k_\star\le C\,\min\{\,\max\{k_{\rm rad},k_d\},\,k_r\,\}\le C\,\min\{k_{\rm rad},k_d,k_r\}$ for some $C=O(1)$, while (a) shows that $k_\star\ge c\,\min\{k_{\rm rad},k_d,k_r\}$ for some $c=O(1)$ (else $g$ would remain negative). Thus $k_\star$ lies at the parametric knee \eqref{eq:kstar-min}, with $O(1)$ constants depending only on $(\rho_{\rm src},\pi_{\rm src},W_2)$. This completes the proof.
\end{proof}

\begin{corollary}[Light--radion bound on the split-time source bandlimit]
\label{cor:light-radion}
If the radion scale dominates the knee, $m_{\rm rad}\ll d^{-1},\,r(\Lambda_{\min})^{-1/2}$, then
\begin{equation}
\label{eq:Delta0-radion}
\boxed{\quad \Delta_0\;\simeq\; L_4^{\rm(prev)}\,m_{\rm rad}\,,\qquad \lambda_0\;\simeq\; m_{\rm rad}^{-1}\ .\quad}
\end{equation}
In particular, a kpc--scale radion (as tied by observations in Section \ref{subsec:shadow-brane-dm-phenom}) implies a split-time source band composed of ultra--long modes
(comoving scale $\sim$\,kpc) at the split.
\end{corollary}

\paragraph{Summary at the split.}\label{par:uom-main-137}
At split time, under band--locked welded widths and at the minimal feasible cutoff
$\Lambda_{\min}$ (Def.~\ref{def:Lambda-min}, above the source in RG ordering):
(i) the \emph{continuing receiver curvature} is set by the compensator--driven SK argmin at the \emph{lowest allowed} gap
$m_g^{\min}$ (near--critical AdS$_4$ in continuous control; Theorem~\ref{thm:init-argmin});
(ii) the \emph{source bandlimit} is knee–locked by the retarded transfer equality (Proposition~\ref{prop:knee-lock}),
reducing to the light–radion bound \eqref{eq:Delta0-radion} when $m_{\rm rad}$ dominates. Together these
fix the split-side ordered configuration: a continuing receiver as close to critical as acceptance allows, and a source spectrum
restricted to the longest scales that can be efficiently exported into the welded receiver band.

\subsection{Overall character of the aeon evolution}
\label{subsec:aeon-evolution}

In this subsection we describe the full dynamical loop of a single aeon under the standing locks and
health assumptions of \S\ref{subsec:init-locked}. We separate the \emph{receiver} and \emph{source}
evolutions, and then connect them by the exported payload $\Delta I_k$ and the BY rectifier $\Theta_k$.
Unless otherwise stated, the welded pulse has band–locked widths
$s_t^\star=\alpha_t^{*}/\bar\Lambda_k$, $s_x^\star=\alpha_x^{*}/(v\bar\Lambda_k)$, with $\alpha_t^{*},\alpha_x^{*}$ fixed
$O(1)$ constants, and the RP block is absent (OS–positivity).

\subsubsection*{Receiver evolution: tickwise update and outcomes}
\label{subsubsec:receiver-phases}

\paragraph{State variables and update laws.}\label{par:uom-main-138}
Along the receiver trajectory we track
\[
\Lambda_{k+1}>\Lambda_k\ (\Delta I_k>0),\qquad I_{k+1}=I_k+\Delta I_k,\qquad
\xi_{k+1}=\xi_k-\Theta_k,\quad \Theta_k:=K(\bar\Lambda_k)\Delta I_k.
\]
All tickwise statements are evaluated at the midpoint cutoff $\bar\Lambda_k$ with
$O(\Delta\Lambda/\Lambda)$ corrections.

\paragraph{Flip (first--crossing).}\label{par:uom-main-139}
The split/flip tick is the first $N$ such that
$\sum_{k\le N}\Theta_k\ge |\xi_0|$ (equivalently $\xi_N\le 0$).

\paragraph{Exhaustion (terminal boundary).}\label{par:uom-main-140}
If $\sup_N\sum_{k\le N}\Theta_k<|\xi_0|$ and accepted payload ceases in the sense that the
available wall work is integrable (equivalently $\sum_k\Delta\lambda_k<\infty$), then
$\Lambda$ converges to a terminal value $\Lambda_\infty$. This is a terminal boundary condition,
not a dynamical phase.

\paragraph{Optional approximation (no phase meaning).}\label{par:uom-main-141}
In any segment where $\beta_k$ and offered payload vary slowly, the admitted payload is approximately
geometric; this is a local approximation, not a distinct epoch.

\paragraph{Split identification.}\label{par:uom-main-142}
\begin{assumption}[Split identification]\label{ass:split-ident}
The condition $\sigma_{\rm eff}\to\sigma_\star$ (equivalently $\Lambda_{4,{\rm R}}\to0^+$) is used only as a
\emph{diagnostic} of boundary/KKT approach in the one--wall channel. In Model~A the split/flip onset is
identified by the first crossing of the two--wall vs.\ one--wall free--energy difference
$\Delta\mathcal F(\Lambda)$ (see App.~\ref{app:split-check}); the $\sigma_{\rm eff}$ condition is checked
post hoc as a consistency signal, not as a trigger.
\end{assumption}

\paragraph{Split/flip criterion (first--crossing).}\label{par:uom-main-143}
Define $\Theta_k:=K(\bar\Lambda_k)\,\Delta I_k$. The split/flip tick is the first $N$ such that
\begin{equation}
\label{eq:flip-condition-updated}
\sum_{k=0}^{N}\Theta_k\ \ge\ |\langle R\rangle_0|\,.
\end{equation}
In Model~A this condition remains the unsplit curvature--accumulation diagnostic, while the split onset
is triggered by the first crossing of $\Delta\mathcal F(\Lambda)$ (App.~\ref{app:split-check}).
Under \emph{continuous control} one can stop exactly at threshold, so the immediate post--flip curvature
can be made arbitrarily small. If there is a nonzero minimum export quantum $\Delta I_{\min}$ and leakage
$\ell$, the last--step overshoot obeys
\[
0\ \le\ R_{\rm next}^{\rm split}\ \lesssim\ K_{\max}\,\big(\Delta I_{\min}+O(\ell)\big),
\qquad K_{\max}:=\max_{k\le N} K(\bar\Lambda_k).
\]

\paragraph{Copyedness and no--cloning (sub--cutoff equivalence).}\label{par:uom-main-144}
Let $\mathcal A_{<\Lambda}$ denote the observable subalgebra below cutoff $\Lambda$.
\begin{definition}[Sub--cutoff equivalence]
\label{def:uom-main-002}
We say two states $\rho_{\rm old}$ and $\rho_{\rm new}$ are $\epsilon_\Lambda$--equivalent at scale $\Lambda$ if
\[
\big|\Tr\big((\rho_{\rm new}-\rho_{\rm old})\,O\big)\big|\ \le\ \epsilon_\Lambda\,\|O\|\qquad
\text{for all }O\in\mathcal A_{<\Lambda}\,.
\]
This algebraic notion avoids factorization assumptions; under standard continuity assumptions it is
equivalent to small trace distance on the restricted algebra.
\end{definition}
\begin{remark}[No--cloning resolution via dilation]
\label{rem:uom-main-007}
The split exposes part of a Stinespring dilation ancilla as a dynamical brane sector. The two outputs
(new dS source and continuing AdS receiver) are correlated and do not carry independent copies of the UV state.
Sub--cutoff equivalence reflects operator--algebra redundancy (QEC/entanglement wedge intuition), not
microscopic cloning.
\end{remark}
\begin{definition}[Late--aeon similarity diagnostic]
\label{def:uom-main-003}
Track $d_\Lambda(\rho^{(n+1)},\rho^{(n)})$ (or $S_\Lambda$, $F_\Lambda$ on $\mathcal A_{<\Lambda}$); a regime
with $\epsilon_\Lambda^{(n)}\to 0$ as $n\to\infty$ signals late--aeon similarity of low--energy data.
\end{definition}



\subsubsection*{2(a). Qualitative characterization of the source evolution}
\label{subsubsec:source-phases}
We separate three qualitatively distinct epochs on the previous–aeon source brane (dS$_4$):

\paragraph{Initial vacuum decay (post–bridge).} The radiative sector of the vacuum–decay state injects\label{par:uom-main-145}
coherence of order unity in the band knee–locked at $k_\star$; the band–limited BD anchor is KMS at
$T_{\rm BD}=H_{\rm prev}/(2\pi)$. The pointer–aligned QND dephaser with DoG envelope reduces coherence
without heating; under steady locks the dephasing exponent per accepted tick is
\begin{equation}
\label{eq:beta-def}
\beta\ :=\ c_\ast\,\delta_L(\Delta)^2\,\mathcal K\,\mathsf N_{\rm src}(0;H_{\rm prev}),\qquad
\mathsf N_{\rm src}(0;H_{\rm prev})\ \propto\ T_{\rm BD}\ \propto\ H_{\rm prev},
\end{equation}
with $c_\ast$ a shape constant and $\mathsf N_{\rm src}$ the Keldysh/noise kernel of the BD commutant bath
(Ohmic/soft at low frequency). This sets the \emph{granularity} (half–life per tick when $\beta=\ln 2$) of
the local exponential decay (see \S\ref{subsubsec:exp-tail}).

\paragraph{Matter era.} As matter dominates on the source brane, the radiative QND payload lowers in intensity;\label{par:uom-main-146}
the QND channel remains pointer–aligned and heat–safe. The band knee evolves only slowly (retarded coupling changes
mildly) and the receiver refinements track acceptance; no structural change in the source dephaser occurs.

\paragraph{SMBH era.} A population of supermassive black holes evaporates on the source brane, providing a\label{par:uom-main-147}
monotonically increasing \emph{available} payload along the receiver trajectory. In our effective model this appears as an
increasing \emph{available payload} for ticks at larger $\bar\Lambda_k$; the BY increment
$\Theta_k=K(\bar\Lambda_k)\Delta I_k$ can therefore grow as $\bar\Lambda_k$ increases along the trajectory,
even as $\Lambda$ drifts slowly. This favors a
flip after fewer ticks if one wishes to minimize the last–step overshoot.

\subsubsection*{2(b). Exponential decay of the QND decohering payload (local fit)}
\label{subsubsec:exp-tail}
We formalize the “exponential tail” of the exported payload under steady locks.

\begin{lemma}[Local exponential tail]\label{lem:exp-tail}
Let the source (previous–aeon dS brane) be band–limited to $\Delta$ and let the pointer–aligned QND
generator be
\[
\mathcal L(t)\rho\;=\;-\frac12\sum_{a}\gamma_a(t)\,[L_a,[L_a,\rho]]\,,\qquad L_a=f_a\!\big(H_{\mathrm{sp},\Delta}\big),
\]
with $\gamma_a(t)\ge 0$ and $[L_a,H_{\mathrm{sp},\Delta}]=0$. Fix a tick window $I=[t_-,t_+]$ and assume a
\emph{steady actuator} regime on $I$, i.e.\ $\gamma_a(t)\equiv \gamma_a$ are constant on $I$ (DoG envelope–squared
averaged to a constant~$\gamma_a$). Let $\sigma_\Delta\propto e^{-\beta_{\mathrm{BD}}H_{\mathrm{sp},\Delta}}$ be the
band–limited BD anchor and define the source modular free energy $F(t):=S(\rho_t\|\sigma_\Delta)$. Set
\[
\beta\;:=\;c_\ast\,\delta_L(\Delta)^2\,\mathcal K\,,\qquad
\mathcal K:=\sum_a \gamma_a\,\big|I\big|,\qquad
\delta_L(\Delta):=\inf_{i\ne j}\Big(\sum_a \big|\ell_a(i)-\ell_a(j)\big|^2\Big)^{1/2},
\]
where $L_a\ket{i}=\ell_a(i)\ket{i}$ in a common eigenbasis of $\{L_a\}$ and $H_{\mathrm{sp},\Delta}$, and $c_\ast>0$
is a (dimensionless) constant depending only on the chosen information divergence (Umegaki). Then across the tick:
\begin{equation}\label{eq:F-exp}
F(t_+)\;\le\;e^{-\beta}\,F(t_-)\,,\qquad
\Delta I:=F(t_-)-F(t_+)\ \ge\ \big(1-e^{-\beta}\big)\,F(t_-)\,.
\end{equation}
If, in addition, the MLSI constant of $\mathcal L$ relative to $\sigma_\Delta$ equals $\beta/|I|$, then
\eqref{eq:F-exp} holds with equality and the per–tick payload satisfies
\begin{equation}\label{eq:tail-geom}
\Delta I_k\;=\;F_k\,(1-e^{-\beta})\;=\;F_0\,\big(1-e^{-\beta}\big)\,e^{-\beta k}\quad (k\ge 0),
\end{equation}
after reindexing the segment so that $k=0$ labels its first tick.
\end{lemma}

\begin{proof}
\emph{Step 1 (Exact dephasing semigroup on the band).}
Since $[L_a,H_{\mathrm{sp},\Delta}]=0$, the generator is a \emph{pure dephaser} in the common eigenbasis
$\{\ket{i}\}_{i}$ of $H_{\mathrm{sp},\Delta}$ and all $L_a$. On the tick $I$, with constant $\gamma_a$,
\[
\rho_{ij}(t)\;=\;\exp\!\Big(-\frac{t-t_-}{2}\sum_a \gamma_a\,\big(\ell_a(i)-\ell_a(j)\big)^2\Big)\,\rho_{ij}(t_-)\,,
\qquad \rho_{ii}(t)\equiv \rho_{ii}(t_-),
\]
i.e.\ the diagonals are invariant and each off–diagonal mode decays with rate
$\lambda_{ij}:=\frac12\sum_a \gamma_a(\ell_a(i)-\ell_a(j))^2$.

\emph{Step 2 (Entropy–production inequality and MLSI).}
Let $\sigma_\Delta$ be the band–limited BD anchor (diagonal in the same basis). The Umegaki relative entropy
$F(t):=S(\rho_t\|\sigma_\Delta)$ is differentiable on $I$ and obeys Spohn’s identity
\[
\frac{d}{dt}F(t)\;=\;-\dot{\Sigma}(t)\;=\;-\sum_a \gamma_a\,\underbrace{\Tr\Big([L_a,\rho_t]\,[L_a,\log\rho_t-\log\sigma_\Delta]\Big)}_{:=\ \mathcal E_a(\rho_t)\ \ge 0}.
\]
Since $\mathcal L$ is a dephaser reversible w.r.t.\ $\sigma_\Delta$ and diagonal in the $L_a$–eigenbasis,
its \emph{modified logarithmic Sobolev} constant $\alpha_{\rm MLSI}$ is well–defined and satisfies the
dissipation inequality (see, e.g., standard proofs for diagonal generators\footnote{For completeness:
write $\rho_t=D(\rho_t)+C_t$ with $D(\rho_t)$ the diagonal part and $C_t$ off–diagonal. Because
$D(\rho_t)$ is constant in $t$ and $\mathcal L$ acts diagonally on each $C_{ij}$, one reduces the
inequality to one–mode dephasing, where the MLSI reduces to a one–dimensional Grönwall estimate; details
are identical to the commutative (classical) reversible case.})
\[
\frac{d}{dt}F(t)\;=\;-\sum_a \gamma_a\,\mathcal E_a(\rho_t)\ \le\ -\,\alpha_{\rm MLSI}\,F(t)\,.
\]
Integrating on $I$ gives $F(t_+)\le e^{-\alpha_{\rm MLSI}\,|I|}F(t_-)$, hence the first inequality in
\eqref{eq:F-exp} holds with $\beta:=\alpha_{\rm MLSI}|I|$. If the MLSI is \emph{sharp} for the chosen
semigroup and anchor (which is the case for diagonal dephasers), equality holds.

\emph{Step 3 (Lower bound in terms of pointer gap and fluence).}
For diagonal dephasers the spectral gap lower–bounds the MLSI constant: there is $c_\ast>0$ (depending only
on the divergence) such that $\alpha_{\rm MLSI}\ge c_\ast\,\lambda_{\min}$ with
$\lambda_{\min}:=\min_{i\ne j}\lambda_{ij}=\frac12\sum_a \gamma_a\inf_{i\ne j}(\ell_a(i)-\ell_a(j))^2$.
Thus $\beta=\alpha_{\rm MLSI}|I|\ge c_\ast\,\delta_L(\Delta)^2 \sum_a\gamma_a\,|I|$, i.e. the advertised
$\beta=c_\ast\,\delta_L(\Delta)^2\mathcal K$. This gives the second inequality in \eqref{eq:F-exp}.

\emph{Step 4 (Geometric tail).}
If $\alpha_{\rm MLSI}$ is \emph{attained} and constant across ticks (steady locks), then equality holds in
\eqref{eq:F-exp} and the recursion $F_{k+1}=e^{-\beta}F_k$ yields \eqref{eq:tail-geom} immediately, with
$\Delta I_k=F_k-F_{k+1}=F_k(1-e^{-\beta})$. This completes the proof.
\end{proof}

\begin{remark}[Slowly varying regime]
\label{rem:uom-main-008}
If $u,\Sigma,\Delta$ evolve slowly across ticks (adiabatic lock), then $\beta_k$ varies slowly and the tail
remains \emph{approximately} exponential with a slowly varying exponent. All results below hold upon replacing
$\beta$ by the tick–averaged $\bar\beta$ and $F_0$ by a local segment value $F_{\rm seg}$.
\end{remark}

\subsubsection*{\texorpdfstring{2(c). Source Hubble scale and the tail parameters $(F_0,\beta)$}{2(c). Source Hubble scale and the tail parameters (F_0,beta)}}
\label{subsubsec:hubble-tail}
The dependence of the exponential tail on the previous–aeon Hubble parameter $H_{\rm prev}$ enters through
(i) the KMS temperature in the dephasing exponent $\beta$, and (ii) the modular anchor $C_\ast(H_{\rm prev})$
in the initial stock $F_0$.

\begin{proposition}[KMS scaling of the dephasing exponent]\label{prop:kms-beta}
Let the source commutant bath be the Bunch–Davies KMS state in a static patch with Hubble parameter
$H_{\mathrm{prev}}>0$. Assume that the pointer–aligned QND channel couples linearly to a bath observable
$B$ with an Ohmic/soft spectral density near $\omega=0$. Denote by $\mathsf N_{\rm src}(\omega;H_{\mathrm{prev}})$
the (stationary) source Keldysh kernel at frequency $\omega$ and by $\beta$ the per–tick dephasing exponent
as in Lemma~\ref{lem:exp-tail}. Then there is a constant $\kappa>0$ (actuator/shape–dependent) such that
\begin{equation}\label{eq:kms-scaling}
\boxed{\qquad \beta\ =\ \kappa\,H_{\mathrm{prev}}\ .\qquad}
\end{equation}
\end{proposition}

\begin{proof}
\emph{Step 1 (KMS relation and low–frequency behavior).}
In the static patch the BD state is KMS at inverse temperature $\beta_{\mathrm{BD}}=2\pi/H_{\mathrm{prev}}$.
Let $S^{>}_B(\omega)=\int_{\R}e^{i\omega t}\,\langle B(t)B(0)\rangle\,dt$ and
$S^{<}_B(\omega)=\int_{\R}e^{i\omega t}\,\langle B(0)B(t)\rangle\,dt$ be the Wightman spectra. The KMS
condition gives
\[
S^{>}_B(\omega)\;=\;e^{\beta_{\mathrm{BD}}\omega}\,S^{<}_B(\omega)\,.
\]
The Keldysh/noise spectrum is
$\mathsf N_{\rm src}(\omega)=\tfrac12\big(S^{>}_B(\omega)+S^{<}_B(\omega)\big)$, hence
\[
\mathsf N_{\rm src}(\omega)\;=\;\tfrac12\big(1+e^{\beta_{\mathrm{BD}}\omega}\big)S^{<}_B(\omega)
\;=\;J(\omega)\,\coth\!\Big(\frac{\beta_{\mathrm{BD}}\omega}{2}\Big),
\]
where $J(\omega):=\tfrac12\big(1-e^{\beta_{\mathrm{BD}}\omega}\big)S^{<}_B(\omega)$ is the (odd) spectral density.
Assume $J(\omega)$ is \emph{Ohmic/soft} near $\omega=0$, i.e.\ $J(\omega)=\eta\,\omega+o(\omega)$ with $\eta>0$,
then, as $\omega\to 0$,
\[
\mathsf N_{\rm src}(\omega)\;=\;\eta\,\omega\,\coth\!\Big(\frac{\beta_{\mathrm{BD}}\omega}{2}\Big)+o(\omega)
\;=\;\eta\,\omega\cdot \frac{2}{\beta_{\mathrm{BD}}\omega}\ +\ o(1)\;=\;\frac{2\eta}{\beta_{\mathrm{BD}}}\ +\ o(1)
\;=\;\frac{\eta}{\pi}\,H_{\mathrm{prev}}\ +\ o(1).
\]
Therefore the \emph{zero–frequency} Keldysh kernel scales linearly with $H_{\mathrm{prev}}$:
\begin{equation}\label{eq:noise-zero}
\boxed{\ \mathsf N_{\rm src}(0;H_{\mathrm{prev}})\;=\;c_0\,H_{\mathrm{prev}}\,,\qquad c_0:=\eta/\pi>0\ .\ }
\end{equation}

\emph{Step 2 (From noise to the dephasing exponent).}
By Lemma~\ref{lem:exp-tail}, under steady locks on a tick of duration $|I|$ one has
$\beta=c_\ast\,\delta_L(\Delta)^2\,\mathcal K\,\mathsf N_{\rm src}(0;H_{\mathrm{prev}})$, where
$\mathcal K=\sum_a \gamma_a\,|I|$ and $\delta_L(\Delta)$ are \emph{independent} of $H_{\mathrm{prev}}$ (they depend
on actuator and band geometry only). Combining with \eqref{eq:noise-zero} gives \eqref{eq:kms-scaling}
with $\kappa:=c_\ast\,\delta_L(\Delta)^2\,\mathcal K\,c_0>0$.
\end{proof}

\begin{proposition}[Anchor scaling of the initial stock]\label{prop:F0-H}
Let $\sigma(H)$ be the band–limited BD anchor on the source (previous–aeon dS brane) at Hubble parameter $H>0$,
i.e.\ $\sigma(H)\propto e^{-\beta_{\mathrm{BD}}H_{\mathrm{sp},\Delta}}$ with $\beta_{\mathrm{BD}}=2\pi/H$, restricted
to the band projector $P_\Delta$. Let $\rho_0$ be the band–limited initial state at the onset of the steady regime
(local fit) and let $F_0(H):=S(\rho_0\|\sigma(H))$ be the initial modular free energy. Then, for fixed $\rho_0$
and on a welded band–edge, one has the \emph{anchor scaling}
\begin{equation}\label{eq:F0-scaling}
\boxed{\qquad F_0(H)\ =\ \alpha\,H\ +\ o(H)\quad\text{as }H\downarrow 0,\qquad \alpha>0\ \text{depending on }\rho_0\text{ and the band.}\qquad}
\end{equation}
More precisely, for small deviations $\delta C:=C_0-C_\ast(H)$ of the covariance (in any Gaussian sector) one has
\begin{equation}\label{eq:F0-quadratic}
F_0(H)\;=\;\tfrac14\,\mathrm{tr}\!\big(C_\ast(H)^{-1}\delta C\,C_\ast(H)^{-1}\delta C\big)\ +\ o(\|\delta C\|^2)
\ \propto\ H\,\|\delta C\|_{\rm band}^2,
\end{equation}
with the proportionality factor determined by the welded band–edge spectrum.
\end{proposition}

\begin{proof}
\emph{Step 1 (Quadratic expansion of $S$ at the anchor).}
By standard information–geometric expansion (Taylor expansion of the Umegaki divergence), for any smooth
one–parameter family $\rho_\epsilon$ with $\rho_0=\sigma(H)$,
\begin{equation}\label{eq:quad-general}
S(\rho_\epsilon\|\sigma(H))\;=\;\frac{\epsilon^2}{2}\,\big\langle\!\big\langle \dot\rho_0,\dot\rho_0\big\rangle\!\big\rangle_{\sigma(H)}
\ +\ o(\epsilon^2),
\end{equation}
where $\langle\!\langle\cdot,\cdot\rangle\!\rangle_{\sigma}$ is the Bogoliubov–Kubo–Mori (BKM) inner product at $\sigma$.
In the Gaussian (zero–mean) sector, the tangent is specified by the covariance variation $\dot C$ and
\begin{equation}\label{eq:BKM-cov}
\big\langle\!\big\langle \dot\rho,\dot\rho\big\rangle\!\big\rangle_{\sigma(H)}
\;=\;\tfrac12\,\mathrm{tr}\!\big(C_\ast(H)^{-1}\dot C\,C_\ast(H)^{-1}\dot C\big).
\end{equation}
Putting \eqref{eq:quad-general} and \eqref{eq:BKM-cov} together, the second–order expansion of
$F_0(H)=S(\rho_0\|\sigma(H))$ around $C_\ast(H)$ yields precisely \eqref{eq:F0-quadratic} with
$\delta C=C_0-C_\ast(H)$.

\emph{Step 2 (Band–edge scaling of $C_\ast(H)^{-1}$).}
Let $C_\ast(H)$ be the band–limited covariance of the BD anchor (Euclidean two–point), expanded in
$S^3$ harmonics. On welded band–edge modes of angular momentum $\ell$ the BD frequency scales as
$\omega_\ell(H)=\sqrt{H^2\ell(\ell+2)+m^2}\sim H\,\ell$ (in the UV part of the band). Therefore
$C_\ast(H)^{-1}(\ell)\sim \omega_\ell(H)\propto H$ mode–wise on the welded band–edge. Consequently,
for a fixed $\delta C$ supported in the welded band, the quadratic form \eqref{eq:F0-quadratic} scales as
\[
\tfrac14\,\mathrm{tr}\!\big(C_\ast(H)^{-1}\delta C\,C_\ast(H)^{-1}\delta C\big)\ \propto\ 
\sum_{\ell\in{\rm band}} \omega_\ell(H)^2\,\|\delta C(\ell)\|^2
\ \propto\ H\,\|\delta C\|^2_{\rm band},
\]
where the last proportionality pulls out a single factor $H$ after summing over the welded band weights
(using that $\omega_\ell(H)\sim H\ell$ and the band window scales $\ell\mapsto \ell/H$ at fixed physical band).

\emph{Step 3 (Conclusion).}
Collecting the scalings, for fixed $\rho_0$ (hence fixed $\delta C$) one has
$F_0(H)=\alpha\,H+o(H)$ as $H\downarrow 0$, with $\alpha>0$ determined by the welded band–edge spectrum and
$\delta C$ through \eqref{eq:F0-quadratic}. This is \eqref{eq:F0-scaling}.
\end{proof}

\paragraph{Summary of tail parameters.} Combining Propositions~\ref{prop:kms-beta}–\ref{prop:F0-H},\label{par:uom-main-148}
the local exponential fit is governed by
\begin{equation}
\label{eq:tail-H}
\beta\ =\ \kappa H_{\rm prev},\qquad F_0\ =\ \alpha H_{\rm prev}\,\big\|\delta C\big\|^2_{\rm band},
\end{equation}
with positive constants $\kappa,\alpha$ set by actuator and shape choices. In particular, larger $H_{\rm prev}$
increases the local dephasing rate and payload scale, while smaller $H_{\rm prev}$ does the opposite.

\subsection{General derivation from the receiver dynamics}
\label{subsec:general-derivation}

We work entirely on the \emph{receiver} side (healthy KR head, RP block absent, welded pulse with
band–locked widths) and derive the curvature–information law that controls the last step at the flip,
together with an estimator for the \emph{split-side AdS$_4$ curvature} carried by the continuing receiver at aeon entry.

\paragraph{Standing locks and notation.}\label{par:uom-main-149}
Let the welded trace–sector pulse have band–locked widths
$s_t^\star=\alpha_t^{*}/\bar\Lambda_k$ and $s_x^\star=\alpha_x^{*}/(v\,\bar\Lambda_k)$ with
$\alpha_t^{*},\alpha_x^{*}>0$ fixed (the P3 extremal shape constants). The receiver cutoff
evolves across the tick; welded quantities are evaluated at $\bar\Lambda_k$ with
$O(\Delta\Lambda/\Lambda)$ corrections. The BY tick increment is
\begin{equation}
\label{eq:BY-step-deriv-fixed}
\Theta_k\;=\;K(\bar\Lambda_k)\,\Delta I_k,\qquad
K(\Lambda)\;:=\;\gamma\,(8\pi G_4)^2\,C_{\rm LoG}\;
\frac{\big\langle|\mathcal K(\omega,k;\Lambda)|^2\big\rangle}{\;s_t^\star\,\big(s_x^\star\big)^3\;G(s_t^\star,s_x^\star;\Lambda)}\,,
\qquad \Delta I_k>0,
\end{equation}
where $G$ is the welded throughput and $\langle|\mathcal K|^2\rangle$ is the retarded
trace–channel average evaluated on the locked envelope. Under the widths locks one has
$G(s_t^\star,s_x^\star;\Lambda)\propto \Lambda^{12}$ and $1/(s_t^\star\,(s_x^\star)^3)\propto \Lambda^{4}$, so the
algebraic prefactor is fixed. For the screened/Yukawa trace response used here,
$\mathcal K(\omega,k;\Lambda)\simeq \dfrac{q}{q+r(\Lambda)\,k^2}\,e^{-d\,q}$ with $q=\sqrt{k^2+m_\bullet^2-\omega^2}$
and $r(\Lambda)\sim r_0+c\,\Lambda^2$, and for the locked spectral locus $k\sim \Lambda$, $\omega\sim O(\Lambda)$ (or $0$),
one obtains the unified decay law
\begin{equation}
\label{eq:K-final-scaling}
K(\Lambda)\ =\ K_0\,\Lambda^{-p}\,e^{-2\,d(\Lambda)\,\Lambda},
\qquad p\in\{10,14\},
\end{equation}
with $p=14$ when the induced–gravity screen is active and $p=10$ without the screen.
The DC curvature obeys $\langle R\rangle_{k+1}=\langle R\rangle_k+\Theta_k$.
At aeon entry the continuing receiver sits near--critical AdS$_4$ (Theorem~\ref{thm:init-argmin}); we denote by
$R_0=4\Lambda_4=-12/L_4^2<0$ its split-side curvature and by $|R_0|$ its magnitude.

\subsubsection*{First--crossing overshoot and bookkeeping partitions}
\label{subsubsec:tail-first-crossing}
We track the tickwise curvature increments by
\[
\Theta_k:=K(\bar\Lambda_k)\Delta I_k,\qquad \langle R\rangle_{k+1}=\langle R\rangle_k+\Theta_k.
\]
Fix a reference cutoff $\Lambda_{\rm cut}$ and define the bookkeeping partition
\[
S_{\le}:=\sum_{k:\,\bar\Lambda_k\le\Lambda_{\rm cut}}K(\bar\Lambda_k)\Delta I_k,\qquad
S_{>}:=\sum_{k:\,\bar\Lambda_k>\Lambda_{\rm cut}}K(\bar\Lambda_k)\Delta I_k,
\]
so that $D(\Lambda_{\rm cut}):=|R_0|-S_{\le}$ measures the residual deficit after the chosen partition.
This split is bookkeeping only; it has no dynamical phase meaning.

Let $N$ be the first tick such that $\sum_{k\le N}\Theta_k\ge |R_0|$ (flip). Define the overshoot
\[
\delta_N:=\sum_{k\le N}\Theta_k-|R_0|\ \ge\ 0.
\]
By the curvature budget identity \eqref{eq:budget-identity}, $0\le\delta_N<\Theta_N$, so the post--flip
residual is controlled by the last admitted step without invoking any phase split.

\paragraph{Optional local exponential fit (no phase meaning).}\label{par:uom-main-150}
In any segment of ticks where $\beta_k$ and the offered payload vary slowly and $K(\bar\Lambda_k)$ drifts mildly,
one may reindex the segment with $k=0$ at its start and model
\[
\Delta I_k\approx F\,(1-e^{-\beta})\,e^{-\beta k}\,,\qquad k\ge 0,
\]
together with an effective weight
\[
K^{\rm eff}:=\frac{\sum_{k\in\mathcal K}K(\bar\Lambda_k)\Delta I_k}{\sum_{k\in\mathcal K}\Delta I_k}\,,
\]
so that $\Theta_k\approx K^{\rm eff}\Delta I_k$ with $O(\Delta\Lambda/\Lambda)$ errors. This is a local
approximation, not a dynamical phase.

\subsubsection*{Aeon--to--aeon curvature map}
\label{subsubsec:aeon-map}

\paragraph{Receiver curvature accumulation.}\label{par:uom-main-151}
With welded widths band--locked and the RP block absent, the BY increment on tick $k$ is
\begin{equation}\label{eq:BY-k}
\Theta_k\;=\;K(\bar\Lambda_k)\,\Delta I_k,\qquad
K(\Lambda)\ :=\ K_0\,\Lambda^{-p}\,e^{-2\,d(\Lambda)\,\Lambda},\quad K_0>0,\ \ p\in\{10,14\},
\end{equation}
and the receiver DC curvature obeys the telescoping recursion
\begin{equation}\label{eq:R-evolution}
\boxed{\quad \langle R\rangle_{n}\;=\;\langle R\rangle_0\ +\ \sum_{k=0}^{n-1} K(\bar\Lambda_k)\,\Delta I_k\ .\quad}
\end{equation}
A bookkeeping partition by $\Lambda_{\rm cut}$ gives $\langle R\rangle_n=\langle R\rangle_0+S_{\le}+S_{>}$ with no
phase meaning.

\paragraph{First--crossing.}\label{par:uom-main-152}
The flip tick is the first $N$ such that
\begin{equation}\label{eq:flip-condition-ref}
\sum_{k=0}^{N}\Theta_k\ \ge\ |R_0|\,,
\end{equation}
which is equivalent to \eqref{eq:flip-condition-updated}. The overshoot bound above controls the residual.

\subsubsection*{Curvature contraction: envelope and floor}
\label{subsubsec:geo-contract}

\paragraph{Upper envelope (worst--case).}\label{par:uom-main-153}
By Proposition~\ref{prop:upper-envelope}, the last--step overshoot satisfies
$0\le\delta_N<\Theta_N$ and, under the dephasing bound, the next--aeon split-side curvature is bounded above by
a function of $H_{\rm prev}$ (hence of the initial curvature scale). This is an envelope, not a timing law, and
requires no phase boundary.

\paragraph{Exhaustion infimum (achievability).}\label{par:uom-main-154}
By Theorem~\ref{thm:exhaustion-infimum}, in the ideal no--leakage and arbitrarily fine payload limit one has
$\inf |R_{\rm next}^{\rm split}|=0$; with leakage and a minimum payload quantum the achievable floor is
$|R_{\rm next}^{\rm split}|\lesssim K_{\max}(\Delta I_{\min}+O(\ell))$.

\paragraph{Schedule leverage (bookkeeping).}\label{par:uom-main-155}
For a fixed payload history, the weighted sum in \eqref{eq:R-evolution} is maximized by allocating export at
smaller $\bar\Lambda_k$ where $K(\bar\Lambda_k)$ is larger. Equivalently, decreasing $\Lambda_{\rm cut}$ increases
$S_{\le}$ and reduces the bookkeeping deficit $D(\Lambda_{\rm cut})$. This is a scheduling effect, not a phase.

\subsubsection*{Cutoff dynamics within an aeon (Wilsonian screen)}
\label{subsubsec:cutoff-dynamics-aeon}

\paragraph{Per--tick move.}\label{par:uom-main-156}
Within a given aeon the welded payload on each accepted tick \emph{must} be placed (QND alignment, admissible bridge).
In the moving--wall gauge this enforces a \emph{strictly nonnegative} per--tick wall move by the SK/HJ balance,
\begin{equation}
\lambda(\Lambda_{k+1})-\lambda(\Lambda_k)
\;\le\; W^{\rm avail}_k,
\qquad
W^{\rm avail}_k
:=\int_{I_k}\!(\mathcal P_{\rm in}-\mathcal P_{\rm out})\,dt
-\int_{I_k}\!\dot{\mathcal F}_{\rm rec}\,dt\ \ge\ 0,
\label{eq:per-tick-move}
\end{equation}
with $\lambda'(\Lambda)=\mu_{\rm RG}(\Lambda)\propto M_4^2(\Lambda)>0$.
Because the welded widths are \emph{band--locked} to the receiver cutoff, the (admitted) payload window sits at $k\sim \bar\Lambda_k$.
Hence every tick that carries nonzero payload has a \emph{nonzero} net wall power on a subinterval of $I_k$ and induces a \emph{positive} increase $\Lambda_{k+1}>\Lambda_k$.\footnote{At the P3 minimizer and under the guards one has $\dot{\mathcal F}_{\rm rec}\le 0$, so $W^{\rm avail}_k\ge \int_{I_k}(\mathcal P_{\rm in}-\mathcal P_{\rm out})\,dt$, and the latter is strictly positive whenever $\Delta I_k>0$ on the welded window.}
Thus there is no option \emph{not} to increase the cutoff \emph{within} an aeon: each nonzero tick \(\Delta I_k>0\) produces a \emph{positive} cutoff increment.

\paragraph{Surrogate wall profile and closed form.}\label{par:uom-main-157}
To obtain a closed solution for the intra--aeon profile $\Lambda_k$, we adopt the monomial surrogate for the RG slope,
\begin{equation}
\lambda'(\Lambda)\;=\;\mu_{\rm RG}(\Lambda)\;\simeq\; \mu_0\,+\,\mu_2\,\Lambda^{2},
\qquad \mu_0,\mu_2>0,
\label{eq:mu-surrogate}
\end{equation}
which captures the induced--gravity scaling $M_4^2(\Lambda)=M_5^3\ell+\kappa a\,\Lambda^2$.
With a local exponential fit $\Delta I_k=F_0(1-e^{-\beta})e^{-\beta k}$ (after reindexing the segment so that $k=0$
labels its start and $\Lambda_0$ is the cutoff at that start) and a slowly--varying efficiency
\(
W^{\rm avail}_k \simeq \eta_\star\,\Delta I_k
\),
the per--tick increment of \(\lambda\) is geometric,
\[
\Delta\lambda_k\ \equiv\ \lambda(\Lambda_{k+1})-\lambda(\Lambda_k)
\ \simeq\ \eta_\star\,F_0(1-e^{-\beta})e^{-\beta k}\,,
\]
and hence
\begin{equation}
\lambda(\Lambda_n)\;=\;\lambda(\Lambda_{0})\;+\;\eta_\star\,F_0\Big(1-e^{-\beta n}\Big).
\label{eq:lambda-geom}
\end{equation}
Inverting \eqref{eq:mu-surrogate} gives the explicit $\Lambda$--profile,
\begin{equation}
\Lambda_n\;=\;\sqrt{\frac{-\,\mu_0\;+\;\sqrt{\mu_0^2+4\mu_2\,\lambda(\Lambda_n)}}{\mu_2}}
\;=\;\sqrt{\frac{-\,\mu_0\;+\;\sqrt{\mu_0^2+4\mu_2\big(\lambda(\Lambda_{0})+\eta_\star F_0(1-e^{-\beta n})\big)}}{\mu_2}}\;,
\label{eq:Lambda-profile}
\end{equation}
which is clearly increasing in \(n\) and saturates to a finite \(\Lambda_\infty\) when the total intra--aeon work is finite:
\[
\Lambda_n\ \nearrow\ \Lambda_\infty
\quad\text{with}\quad
\lambda(\Lambda_\infty)=\lambda(\Lambda_{0})+\eta_\star F_0.
\]
This realizes the post--exhaustion boundary limit and matches the terminal cutoff construction (Sec.~\ref{subsec:alternating-heads}).

\paragraph{Remarks.}\label{par:uom-main-158}
(i) If a slow initial throttling segment is kept, replace \((\eta_\star,\Lambda_{0})\) by the effective values for that segment; the same geometric sum holds for the monomial surrogate.
(ii) The band--lock is essential: it ties nonzero $\Delta I_k$ to nonzero $W^{\rm avail}_k$ at each tick and thus enforces \(\Lambda_{k+1}>\Lambda_k\).
(iii) Beyond the surrogate, any $\mu_{\rm RG}(\Lambda)$ with $\mu_{\rm RG}'>0$ yields the same monotonicity and a unique terminal value determined by the finite geometric sum of available work.


\section{\texorpdfstring{AdS$_5$ evolution across aeons}{AdS5 evolution across aeons}}
\label{sec:ads5-evolution}

\noindent\textbf{Foreview.}
We derive an aeon-to-aeon dynamical map for the pair of small parameters
\[
\varepsilon_{\rm curv}:=|R_0|\quad\text{(split-side AdS$_4$ curvature after the flip)},\qquad
\varepsilon_\Lambda:=\Lambda_{\rm split}^{-1}\quad\text{(split-time resolution/cutoff)},
\]
using only the flip/curvature machinery of \S\ref{sec:aeon-evolution} and the finite--cutoff SK/HJ wall balance. We prove:
\begin{enumerate}[label=(\roman*),itemsep=2pt,topsep=2pt]
\item \emph{Two outcomes per aeon:} either a flip occurs in finite ticks (first--crossing), or the source exhausts first and the receiver becomes the terminal EOW/AdS$_5$ boundary. No scheduler axiom is used.
\item \emph{Terminal boundary interpretation:} in the exhaustion branch the wall work becomes zero and $\Lambda$ approaches a terminal value $\Lambda_\infty$; this is a boundary because the process ends, not a dwell phase with payload.
\item Asymptotically Lorentzian evolution is \emph{unitary}; the only genuine ``Second Law'' lives in the CPTP RG semigroup along meta-time.
\end{enumerate}

\subsection{Standing assumptions and notation}
\label{subsec:ads5-standing}

We import from \S\ref{sec:aeon-evolution}:
\begin{itemize}[leftmargin=1.15em,itemsep=2pt,topsep=2pt]
\item \textbf{Curvature tick law (welded trace):} $\Theta_k=K(\bar\Lambda_k)\,\Delta I_k$ with
\begin{equation}
K(\Lambda)\ \sim\ \kappa_K\,\Lambda^{-p}\,e^{-2\,d(\Lambda)\,\Lambda}\,,\qquad p\in\{10,14\},\ \ d(\Lambda)\ge d_\infty>0,
\label{eq:ads5-K}
\end{equation}
where $d(\Lambda)$ is the effective proper separation entering the screened retarded acceptance on the welded locus ($k\sim\Lambda$).
\item \textbf{Local exponential fit (payload):} in any segment where $\beta_k$ and offered payload vary slowly,
reindex the segment with $k=0$ at its start and model
\begin{equation}
\Delta I_k\ =\ F_0\,(1-e^{-\beta})\,e^{-\beta k}\qquad(k\ge 0)\,.
\label{eq:ads5-tail}
\end{equation}
\item \textbf{Split/flip (first--crossing):} the split/flip tick is the first $N$ with
$\sum_{k\le N}\Theta_k\ge |R_0|$ (Assumption~\ref{ass:split-ident}).
\end{itemize}
The finite--cutoff wall move is controlled by a convex, strictly increasing RG action $\lambda(\Lambda)$ with slope
\begin{equation}
\lambda'(\Lambda)\ \sim\ M_4^2(\Lambda)\ =\ M_5^3\,\ell\ +\ \kappa\,a\,\Lambda^2\ >\ 0.
\label{eq:ads5-lambda-prime}
\end{equation}
We denote the \emph{available wall--work} over tick/aeon by
\begin{equation}
W^{\rm avail}\ :=\ \underbrace{\int(\mathcal P_{\rm in}-\mathcal P_{\rm out})\,dt}_{W^{\rm net}}
-\ \underbrace{\int \dot{\mathcal F}_{\rm rec}\,dt}_{\Delta\mathcal F_{\rm rec}}\,,\qquad \mathcal P_{\rm in}=\mathcal P_{\rm src}+\mathcal P_{\rm stack}+\mathcal P_{\rm hor}+\mathcal P_{\rm bdy},\ \ \mathcal P_{\rm out}=\mathcal P_{\rm deep}+\mathcal P_{\rm noise}.
\label{eq:ads5-Wavail-def}
\end{equation}
Here $\mathcal F_{\rm rec}$ is the (real) receiver modular free energy (Fisher/anomaly part). The per--aeon symbols we use are:
\[
x_n:=|R^{(n)}_0|,\qquad y_n:=\Lambda^{(n)\,-1}_{\rm split},\qquad
B_n:=\sum_{k=0}^{N_n}K(\bar\Lambda_k^{(n)})\,\mathbb E\,\Delta I_k^{(n)},\qquad
W^{\rm avail}_n:=\int_{\text{aeon }n}W^{\rm avail}(t)\,dt,
\]
Here $N_n$ is the first--crossing tick in the flip branch, or the terminal index in the
exhaustion branch (possibly $N_n=\infty$ when the series converges).
We also use the effective welded weights
\begin{equation}
K^{\rm eff}_n\ :=\ \frac{\sum_{k=0}^{N_n}K(\bar\Lambda_k^{(n)})\,\mathbb E\Delta I_k^{(n)}}{\sum_{k=0}^{N_n}\mathbb E\Delta I_k^{(n)}}\,,\qquad
\eta^{\rm eff}_n\ :=\ \frac{\sum_{k=0}^{N_n}\eta_{n,k}\,\mathbb E\Delta I_k^{(n)}}{\sum_{k=0}^{N_n}\mathbb E\Delta I_k^{(n)}}\,,
\label{eq:ads5-eff-weights}
\end{equation}
where $\eta_{n,k}$ is the spectral efficiency factor that multiplies the payload in the wall--work. We write the per--tick source contribution as
\begin{equation}
\label{eq:ads5-Wsrc}
W^{\rm src}_{n,k}\ :=\ \eta_{n,k}\,\Delta I_k^{(n)}\ .
\end{equation}

\subsection{SK/HJ wall balance}
\label{subsec:ads5-gauges}

On the wall $\Sigma_\Lambda$ the SK/HJ balance reads
\begin{equation}
\mu_{\rm RG}(\Lambda)\,\dot\Lambda\ =\ \big(\mathcal P_{\rm in}-\mathcal P_{\rm out}\big)\ -\ \dot{\mathcal F}_{\rm rec}\,,\qquad
\Delta\lambda=\int \mu_{\rm RG}\dot\Lambda\,dt\ \le\ \int W^{\rm avail}(t)\,dt\,,
\label{eq:ads5-wall-ode}
\end{equation}
We interpret the wall work as the $u$--derivative of the master potential $\mathcal G$
introduced in Sec.~\ref{sec:two-times-semiclassical}; the aeon--level balances below
are its integral (budget) form. Here $\mu_{\rm RG}>0$ is proportional to the BY Hamiltonian
density on the wall. We use static--wall gauge only as a post--exhaustion boundary limit
($\Delta I_k\to0$), i.e., after accepted payload ceases and $\dot\Lambda\to 0$; then
$\dot{\mathcal F}_{\rm rec}=\mathcal P_{\rm in}-\mathcal P_{\rm out}$ and Lorentzian
energy conservation is explicit.
In the open--system setting we allow a nonnegative leak term $Q_{\rm leak}\ge 0$ so that
the discrete tick balance becomes $\Delta E_{\rm geom}=J_{\rm wall}+Q_{\rm leak}$ (equivalently,
$W^{\rm avail}=W_{\rm RG}+Q_{\rm leak}$ in the integrated form). This is the minimal modification
that prevents conservative SK/HJ flow from exhausting the curvature budget by accelerating
the cutoff alone.

\paragraph{Energy accounting (pulses and split).}\label{par:uom-main-159}
With a static AdS$_5$ UV boundary the Brown--York charge is conserved. The QND absorption channel is
energy--preserving on the pointer algebra because $[L_a,H_{{\rm sp},\Delta}]=0$, and the KMS bath
mediates entropy export without injecting net energy into the source. The wall work budget $W_{\rm RG}$
appears as a decrease in the appropriate quasi--local free energy of the dS horizon/bath sector,
while the split redistributes quasi--local energy between the new dS source, the continuing AdS receiver,
and the adjacent bulk. The pillbox Hamiltonian constraint gives \eqref{eq:ads5-wall-ode} at the wall,
and we detail the split balance in App.~\ref{app:energy-split}.


\subsection{\texorpdfstring{Coupled aeon--level $\epsilon$--map}{Coupled aeon--level epsilon-map}}
\label{subsec:ads5-coupled-map}

Per aeon $n$ we collapse the detailed tick dynamics into two scalars:
\begin{equation}
\label{eq:ads5-curv-map}
\begin{aligned}
\textbf{Curvature:}&&\mathbb E\,x_{n+1}&\ \le\ c_n\,\mathbb E[\Theta_{N_n}]\ \le\ c_n\,B_n\,,\qquad
B_n:=\sum_{k=0}^{N_n}K(\bar\Lambda_k^{(n)})\,\mathbb E\Delta I_k^{(n)}\,,
\end{aligned}
\end{equation}
\begin{equation}
\label{eq:ads5-cutoff-bound}
\begin{aligned}
\textbf{Cutoff bound:}&&\lambda\!\big(\Lambda^{(n+1)}_{\rm split}\big)&\ \le\ \lambda\!\big(\Lambda^{(n)}_{\rm split}\big)\ +\ W^{\rm avail}_n\,,
\end{aligned}
\end{equation}
\begin{equation}
\label{eq:ads5-cutoff-map}
\begin{aligned}
\textbf{Cutoff update (no slack):}&&\Lambda^{(n+1)}_{\rm split}&\ =\ \lambda^{-1}\!\big(\lambda(\Lambda^{(n)}_{\rm split})+W^{\rm avail}_n\big),
\qquad y_{n+1}\ =\ \big(\lambda^{-1}(\lambda(y_n^{-1})+W^{\rm avail}_n)\big)^{-1}\,.
\end{aligned}
\end{equation}


\subsection{Split-order dominance (no fairness axiom)}\label{subsec:uom-main-003}
The split ordering already enforces
$\Lambda_{\rm R}(t_{\rm split})>(1+\varepsilon)\Lambda_{\rm S}(t_{\rm split})$ in every aeon; no separate
fairness theorem is needed.

\subsection{Terminal boundary in the exhaustion branch}
\label{subsec:ads5-static-wall}

\begin{theorem}[Terminal boundary]
\label{thm:ads5-static}
Under the standing assumptions of \S\ref{subsec:ads5-standing} and the expectation--level curvature bound \eqref{eq:ads5-curv-map}, the series of available wall--work in the exhaustion branch is finite:
\begin{equation}
\sum_{n=0}^\infty W^{\rm avail}_n\ <\ \infty.
\label{eq:ads5-sumW}
\end{equation}
Consequently, $(\lambda(\Lambda^{(n)}_{\rm split}))_{n\ge 0}$ is Cauchy and $\Lambda^{(n)}_{\rm split}\nearrow \Lambda_\infty<\infty$. This terminal value is reached only after accepted payload ceases; it is a boundary condition, not a within--aeon dwell.
\end{theorem}

\begin{proof}[Proof sketch]
Write $W^{\rm avail}_n=\eta^{\rm eff}_n F^{(n)}_0 - L_n$, $L_n\ge 0$. If $\Lambda^{(n)}_{\rm split}$ is bounded, $\eta^{\rm eff}_n\le C$ and summability follows from $\sum_n F^{(n)}_0<\infty$. If $\Lambda^{(n)}_{\rm split}\to\infty$, the welded acceptance/screen bounds imply $\eta^{\rm eff}_n\lesssim \Lambda^q e^{-2 d_\infty \Lambda}$, hence $\sum_n \eta^{\rm eff}_n F^{(n)}_0<\infty$ by comparison. In both cases \eqref{eq:ads5-sumW} holds and the claim follows from monotonicity/convexity of $\lambda$.
\end{proof}

\begin{corollary}[EOW/decoupled tail]
\label{cor:ads5-eow}
For any fixed late wall $\Sigma_{\Lambda_\infty}$ the cumulative shadow--stack inflow is summable:
\[
\sum_{\text{older heads}} W^{\rm stack}\ <\ \infty.
\]
Hence the tail decouples and acts as a true EOW/AdS$_5$ boundary for the active head.
\end{corollary}

\begin{remark}[Arrow placement]
\label{rem:ads5-arrow}
Asymptotically $\dot\Lambda=0$ and $\Im W_\Lambda=0$, so Lorentzian evolution on the brane stack is unitary and P1 is enforced dynamically. The only irreversibility is the CPTP RG semigroup in meta--time $u=\ln(\Lambda_{\rm UV}/\Lambda)$, with data--processing monotones $\partial_u S(\rho_u\|\sigma_u)\le 0$ and nondecreasing recoverable information.
\end{remark}

\subsection{Flip timing and ``beat--the--source'' condition}
\label{subsec:ads5-beat-src}

Let $\Lambda^{(n)}_{\rm S}(t_{\rm split})$ denote the source cutoff at the split time of aeon $n$.
By Def.~\ref{def:Lambda-min} the continuing receiver sits at
$\Lambda^{(n)}_{\rm R}(t_{\rm split})=\max(\Lambda_{\rm phys},(1+\varepsilon)\Lambda^{(n)}_{\rm S}(t_{\rm split}))>\Lambda^{(n)}_{\rm S}(t_{\rm split})$.
The required gap is paid by the split-entry work budget, so a sufficient condition is
\begin{equation}
\boxed{\qquad
\lambda\!\big(\Lambda^{(n)}_{\rm R}(t_{\rm split})\big)\ -\ \lambda\!\big(\Lambda^{(n)}_{\rm S}(t_{\rm split})\big)
\ \le\ W^{\rm avail}_{n,{\rm split}}\ .
\qquad}
\label{eq:ads5-beat-src}
\end{equation}
Here $W^{\rm avail}_{n,{\rm split}}$ denotes the available wall--work budget allocated to the
split step (a short neighborhood of $t_{\rm split}$).
If \eqref{eq:ads5-beat-src} fails, the cycle enters the terminal/no--flip branch of
Thm.~\ref{thm:ads5-static}, with a terminal boundary and no further splits.

\subsection{Outcomes and regime choice}
\label{subsec:ads5-two-regimes}

The coupled aeon--level map admits two outcomes:
\begin{description}[leftmargin=1.25em,itemsep=2pt,topsep=2pt]
\item[Flip branch:] first--crossing occurs in finite ticks; the receiver continues across the split while the dragged completion sector becomes the newborn source, and the map iterates.
\item[Exhaustion branch:] the source exhausts before crossing; the receiver becomes the terminal EOW/AdS$_5$ boundary and no further split occurs.
\end{description}
\noindent\emph{Optional limit remark.} One may still classify wall--work sequences as summable vs divergent,
but this is ancillary to the flip vs exhaustion dichotomy.


\section{Phenomenology and observational tests}
\label{sec:phenomenology}

\subsection{LoG-mix CMB imprint and collider testability of Higgs portal features}
\label{subsec:log-mix-imprint}
\label{subsec:higgs-portal}

This subsection collects the testable consequences of the LoG–driven uplift and its
portal to the Higgs, emphasizing \emph{qualitative} signatures alongside simple,
transparent estimates. The key structural points that organize the predictions are:

\begin{enumerate}[label=\textbf{P\arabic*}. ,leftmargin=2.2em,itemsep=2pt]
\item \textbf{Weighted mixture over ticks.} The imprint is a weighted mixture
\(
\Phi(\Omega)=\sum_k \widetilde A_k\,\Delta_{S^3}g_{\sigma_{x,k}}(\Omega),\quad
\sigma_{x,k}=\alpha_x^*/(v\bar\Lambda_k),\quad
\widetilde A_k\propto K(\bar\Lambda_k)\Delta I_k.
\)
\item \textbf{Unimodal vs shoulder.} Typically the weight distribution in \(\sigma_{x,k}\) is unimodal, giving a single LoG-like feature. A visible shoulder/ring requires a second bump in the effective weights at a distinct \(\bar\Lambda\)-scale.
\item \textbf{Portal freeze at the flip.} The DC component fixes a constant Higgs-mass shift after ring-down (Prop.~\ref{prop:persistence-freeze}).
\end{enumerate}

\paragraph{Template for the CMB temperature pattern (large angles).}\label{par:uom-main-160}
In the Sachs–Wolfe regime a primordial scalar curvature potential \(\Phi(\Omega)\) on the last–scattering
surface imprints
\[
\frac{\Delta T}{T}(\Omega)\ \simeq\ \frac{1}{3}\,\Phi(\Omega)\quad
(\text{SW, large angles})\ .
\]
For a welded LoG spatial envelope at width \(\sigma_x\) on \(S^3\),
\[
\Phi_{\rm LoG}(\Omega)\;=\;A_{\rm LoG}\,\Delta_{S^3}g_{\sigma_x}(\Omega)
\quad\Longrightarrow\quad
\big(\widehat{\Delta T/T}\big)_{\ell m}
\;=\;\frac{A_{\rm LoG}}{3}\,(-\lambda_\ell)\,e^{-\tfrac12\sigma_x^2\lambda_\ell}\,Y_{\ell m},
\]
with \(\lambda_\ell=\ell(\ell+2)\). Thus the \(\ell\)-space template is fixed up to an overall amplitude
and a single angular scale; the real–space profile is the familiar compensated core–ring.

\paragraph{Mixtures and the shoulder.}\label{par:uom-main-161}
Because all ticks are co-centered and LoG, their sum is the \emph{Laplacian of a mixture of Gaussians}:
\[
\Phi_{\rm mix}(\Omega)\;=\;\Delta_{S^3}\!\left(\sum_{i} w_i\,g_{\sigma_{x,i}}(\Omega)\right),
\qquad w_i\ge 0,\ \sum_i w_i=1.
\]
A single dominant width gives an almost pure LoG; a small admixture at a distinct scale produces a
\emph{shoulder/ring}. A detectable inner shoulder requires a second bump in the effective weights
$K(\bar\Lambda_k)\Delta I_k$ at a distinct \(\bar\Lambda\)-scale; if the flip occurs after few ticks,
the imprint is effectively single--scale.

\paragraph{Effective cutoff for a tick subset.}\label{par:uom-main-162}
For any disjoint subset $\mathcal K_2$ of ticks selected by where the weights concentrate in
$\bar\Lambda_k$--space, it is convenient to use the effective weight
\[
K^{\rm eff}\;:=\;\frac{\sum_{k\in\mathcal K_2} K(\bar\Lambda_k)\,\Delta I_k}{\sum_{k\in\mathcal K_2}\Delta I_k}\,,
\]
which differs from pointwise $K(\bar\Lambda_k)$ only at $O(\Delta\Lambda/\Lambda)$. The corresponding angular width
may be summarized by $\sigma_{x,2}\simeq \alpha_x/(v\Lambda_{\rm eff})$ for a suitable effective cutoff
$\Lambda_{\rm eff}$ defined by the weighted subset. This is a bookkeeping device, not a phase.

\paragraph{Effect of band–locked widths on the LoG mixture and CMB predictions.}\label{par:uom-main-163}
Under the welded–width lock proven earlier, each accepted tick $k$ contributes a \emph{fixed–shape} LoG imprint
whose spatial width is \emph{not} a free parameter but set by the receiver cutoff at that tick:
\[
\sigma_{x,k}\;=\;\frac{\alpha_x^{*}}{v\,\bar\Lambda_k},\qquad \sigma_{t,k}\;=\;\frac{\alpha_t^{*}}{\bar\Lambda_k},
\]
with $\alpha_x^{*},\alpha_t^{*}>0$ fixed $O(1)$ constants. The $\ell$–space contribution of tick $k$ to the
Sachs–Wolfe (SW) temperature multipoles is therefore
\begin{equation}\label{eq:SW-Leg}
\Big(\widehat{\Delta T/T}\Big)^{(k)}_{\ell m}
\;=\;\frac{A_k}{3}\,\Big(-\lambda_\ell\Big)\,
\exp\!\Big[-\tfrac12\,\sigma_{x,k}^2\,\lambda_\ell\Big]\;Y_{\ell m},\qquad \lambda_\ell:=\ell(\ell+2),
\end{equation}
up to the (dimensionless) SW transfer normalization. What the lock \emph{does} is to fix
the \emph{angular scale} of each component through $\sigma_{x,k}\propto \bar\Lambda_k^{-1}$ and to fix the
\emph{amplitude scaling} through our band–edge BY weight:
\[
\Theta_k\;=\;K(\bar\Lambda_k)\,\Delta I_k,\qquad
K(\Lambda)=K_0\,\Lambda^{-p}\,e^{-2\,d(\Lambda)\,\Lambda},\qquad K_0>0,\ \ p\in\{10,14\}.
\]
Using the amplitude elimination \eqref{eq:payload-eq-final}, the curvature variance and
hence the primordial–potential amplitude associated with tick $k$ scales \emph{linearly} with $\Theta_k$,
so we may take
\begin{equation}\label{eq:Ak-weight}
A_k\;=\;\mathcal T_k\;K(\bar\Lambda_k)\,\Delta I_k\qquad\big(\ \mathcal T_k>0\ \text{collects transfer factors from the imprint epoch to LSS}\ \big).
\end{equation}
Equations~\eqref{eq:SW-Leg}–\eqref{eq:Ak-weight} completely determine the mixture:
\begin{equation}\label{eq:mix-locked}
\Big(\widehat{\Delta T/T}\Big)_{\ell m}
=\sum_{k}\frac{\mathcal T_k}{3}\,\Big(-\lambda_\ell\Big)\,
\exp\!\Big[-\tfrac12\Big(\frac{\alpha_x^{*}}{v\,\bar\Lambda_k}\Big)^2\,\lambda_\ell\Big]\,
\underbrace{\mathcal C_{\rm SW}\,K(\bar\Lambda_k)\,\Delta I_k}_{=:~\widetilde A_k}\;Y_{\ell m}.
\end{equation}
Thus:
\begin{itemize}
\item The \emph{profile} in $\ell$–space of each term is a pure LoG with width \(\sigma_{x,k}=\alpha_x^{*}/(v\,\bar\Lambda_k)\).
\item The \emph{weight} of each term is fixed by the cutoff–weighted payload \(\widetilde A_k= K(\bar\Lambda_k)\Delta I_k\),
modulo a smooth transfer factor \(\mathcal T_k\).
\end{itemize}

\paragraph{Core/shoulder (ring) prediction under the lock.}\label{par:uom-main-164}
Select two disjoint tick subsets $\mathcal K_1,\mathcal K_2$ by where the weights
$\widetilde A_k\propto K(\bar\Lambda_k)\Delta I_k$ concentrate in $\bar\Lambda_k$--space (bookkeeping only).
Define effective widths
\[
\sigma_{x,1}\simeq \frac{\alpha_x^{*}}{v\,\Lambda_{\rm eff,1}},\qquad
\sigma_{x,2}\simeq \frac{\alpha_x^{*}}{v\,\Lambda_{\rm eff,2}}\,,\qquad \Lambda_{\rm eff,2}>\Lambda_{\rm eff,1},
\]
and the \emph{amplitude ratio} is \emph{entirely} set by the cutoff–weighted payloads (and transfer factors):
\begin{equation}\label{eq:ring-core-bandlocked}
\boxed{\quad
\frac{A_{\rm shoulder}}{A_{\rm core}}
\;\approx\;
\frac{\sum_{k\in\mathcal K_2}K(\bar\Lambda_k)\Delta I_k}{\sum_{k\in\mathcal K_1}K(\bar\Lambda_k)\Delta I_k}
\times\frac{\mathcal T_2}{\mathcal T_1}\ .
\quad}
\end{equation}
With the asymptotics of \(K(\Lambda)\) from \S\ref{subsec:general-derivation}, this implies a strong suppression
by the effective scale ratio, e.g.
\[
\frac{A_{\rm shoulder}}{A_{\rm core}}
\ \sim\
\frac{\mathcal T_2}{\mathcal T_1}
\left(\frac{\Lambda_{\rm eff,1}}{\Lambda_{\rm eff,2}}\right)^{p}
e^{-2d(\Lambda_{\rm eff,2}-\Lambda_{\rm eff,1})}\qquad\big(p\in\{10,14\}\big).
\]


No extra powers of spatial/temporal widths appear beyond this \(\Lambda^{-10\dots-14}\) law because the LoG
prefactors from $1/(\sigma_t\,\sigma_x^3)$ \emph{and} the throughput $G_0$ have already been absorbed
into the single band--edge weight \(K(\Lambda)\) by the amplitude elimination
\eqref{eq:payload-eq-final}. The only residual scale dependence in the \(\ell\)--shape is the
Gaussian factor with \(\sigma_{x,k}\propto \bar\Lambda_k^{-1}\) already exhibited in \eqref{eq:mix-locked}.

\paragraph{Implications.}\label{par:uom-main-165}
\begin{enumerate}
\item \emph{Angular scales are predictive.} The core width is set by the dominant subset $\mathcal K_1$
(smaller $\bar\Lambda_k$), while the shoulder width is set by the subset $\mathcal K_2$ (larger $\bar\Lambda_k$):
\[
\sigma_{x,1}\simeq \frac{\alpha_x^{*}}{v\,\Lambda_{\rm eff,1}},\qquad
\sigma_{x,2}\simeq \frac{\alpha_x^{*}}{v\,\Lambda_{\rm eff,2}},\qquad
\frac{\sigma_{x,1}}{\sigma_{x,2}}\simeq \frac{\Lambda_{\rm eff,2}}{\Lambda_{\rm eff,1}}.
\]
\item \emph{Amplitudes follow the curvature–lever law.}
For comparable payloads and mild transfer variation,
\[
\frac{A_{\rm shoulder}}{A_{\rm core}}
\;\simeq\;\frac{\mathcal T_2}{\mathcal T_1}\,
\frac{\sum_{k\in\mathcal K_2}K(\bar\Lambda_k)\Delta I_k}{\sum_{k\in\mathcal K_1}K(\bar\Lambda_k)\Delta I_k},
\]
so the shoulder is suppressed by the effective lever ratio. Using \(K(\Lambda)\sim \Lambda^{-p}e^{-2d\Lambda}\) with
\(p\in\{10,14\}\) yields a strong polynomial\(\times\)Yukawa suppression in \(\Lambda_{\rm eff,2}/\Lambda_{\rm eff,1}\).
\item \emph{Effective multi–core mixtures.} If the weight distribution has multiple bumps, the core part of
\eqref{eq:mix-locked} is a finite mixture of LoGs with widths \(\sigma_{x,k}=\alpha_x^{*}/(v\bar\Lambda_k)\) and
weights \(\propto K(\bar\Lambda_k)\Delta I_k\). A single dominant subset yields an almost pure LoG; a spread of
scales produces a slightly broadened core.
\end{enumerate}
In summary, the welded–width lock removes residual ambiguity in the LoG mixture: (i) the \emph{scales} are fixed by the
cutoffs used when the ticks occurred, and (ii) the \emph{relative amplitudes} are fixed by the
\(K(\bar\Lambda_k)\)–weighted payloads \(\Delta I_k\) (modulo smooth transfer factors). This sharpens the core/shoulder
prediction to the closed form \eqref{eq:ring-core-bandlocked}.


\paragraph{Polarization and phase.}\label{par:uom-main-166}
The LoG imprint is scalar and predicts a definite TE cross–correlation sign at the feature scale.
In harmonic space, \(E_{\ell m}\) inherits the same scale and phase as \(T_{\ell m}\) (modulo the standard
scalar polarization transfer). Therefore, a matched–filter search in \(T\) should be accompanied by an
aligned, lower–S/N template in \(E\) and in \(TE\).

\paragraph{Non–Gaussian statistics.}\label{par:uom-main-167}
A single compensated feature produces excess kurtosis and large “cold (or hot) area” statistics at a
few–degree scale. Mixtures of co-centered widths preserve compensation and thus remain conspicuous in
wavelets matched to \(\sigma_x\) (e.g. SMHW/LoG). The prediction is \emph{one dominant} outlier with
degree–scale width and, at most, a \emph{fainter inner ring}.

\medskip
\noindent\textbf{Parameter inference and limits.}
Let the welded tick \(k\) inject (receiver–side) curvature variance
\[
\Delta\!\langle R^2\rangle_k
\;=\;(8\pi G_4)^2\,\big\langle|\mathcal K|^2\big\rangle_{\!k}\;
C_{\rm LoG}\;\frac{\Delta I_k}{G(\bar\Lambda_k)}\,\frac{1}{\sigma_{t,k}\,\sigma_{x,k}^{3}},
\qquad
\Theta_k=\gamma\,\Delta\!\langle R^2\rangle_k.
\]
Summing over a core subset $\mathcal K_1$ (e.g.\ ticks with $\bar\Lambda_k\le\Lambda_{\rm cut}$) and propagating to
LSS with a (model–dependent) decay/transfer factor \(\mathcal D_\ast\in(0,1]\),
\begin{equation}
\label{eq:A-LoG-CMB}
A_{\rm LoG}^{\rm (CMB)}\ \simeq\ \frac{1}{3}\,\mathcal D_\ast\,
\sum_{k\in\mathcal K_1}\ \alpha_k\,
\frac{\Delta I_k}{G(\bar\Lambda_k)}\,\frac{1}{\sigma_{t,k}\,\sigma_{x,k}^{3}},
\qquad
\alpha_k:=(8\pi G_4)^2\,\gamma\,\big\langle|\mathcal K|^2\big\rangle_{\!k}\,C_{\rm LoG}.
\end{equation}
\emph{Important:} \(\mathcal D_\ast\) encodes the ring–down of the AdS\(_4\) response to the LSS epoch; it is not fixed by CMB data alone, so \(A_{\rm LoG}^{\rm (CMB)}\) yields a \emph{band} for the underlying sum.

\begin{itemize}[leftmargin=1.6em,itemsep=2pt]
\item \textbf{Lower bound on the core subset sum.} Using \(\mathcal D_\ast\le 1\),
\(\sum_{k\in\mathcal K_1}\alpha_k\,\frac{\Delta I_k}{G}\,\frac{1}{\sigma_{t,k}\sigma_{x,k}^3}
\ \gtrsim\ 3\,A_{\rm LoG}^{\rm (CMB)}\).
\item \textbf{Shoulder bound.} Plug \eqref{eq:ring-core-bandlocked} into the measured profile to bound
the effective scale ratio and payload ratio if a shoulder is detected.
\end{itemize}

\paragraph{Higgs portal and colliders.}\label{par:uom-main-168}
With the auxiliary brane scalar \(\mathcal S\) (mass \(m_{\mathcal S}\gg v\)) and portal coupling
\(-\lambda_p\,\mathcal S\,H^\dagger H\),
Prop.~\ref{prop:persistence-freeze} gives a \emph{constant} Higgs mass shift
\(\Delta\mu_H^2=-\lambda_p\,J_0/m_{\mathcal S}^2\) after the flip. Direct signals are loop–suppressed:
the leading fractional deviation in Higgs couplings scales as
\[
\delta\kappa_h\ \sim\ \frac{\lambda_p^2}{16\pi^2}\,\frac{v^2}{m_{\mathcal S}^2}\,.
\]
For \(m_{\mathcal S}\sim 10^{4}\!-\!10^{5}\,\mathrm{GeV}\) and \(\lambda_p=\mathcal O(1)\), one has
\(\delta\kappa_h\sim 10^{-6}\!-\!10^{-4}\), below HL–LHC reach. Thus collider phenomenology is
\emph{naturally quiet} in this minimal portal.

\paragraph{Equivalence principle / fifth forces.}\label{par:uom-main-169}
In the continuity picture adopted here a light scalar exists throughout (boundary--distance mode at the
KR head; stabilized radion in interior RS1 segments), but in the present KR/dS configuration the receiver’s
visible coupling is suppressed by the induced--gravity Robin boundary condition on the head. The
long--range Yukawa correction to Newton’s law therefore scales as
\(
\alpha_{\rm vis}\propto \varepsilon_{\rm BC}^2=(1+s_{\rm vis}\,k\,L)^{-2}
\),
and with $s_{\rm vis}\sim 10$–$100$ at late times we expect
\(
\alpha_{\rm vis}\lesssim 10^{-3}\!-\!10^{-5}
\),
comfortably below current Solar–System and laboratory bounds at the relevant ranges (Sec.~\ref{subsubsec:fifthforce-robin}). The coupling is universal (to $T^\mu{}_\mu$), so composition–dependent violations of the weak equivalence principle are not expected at leading order; constraints reduce to limits on $\alpha_{\rm vis}$ and the range $\lambda_r=m_r^{-1}$.

When our brane is \emph{overtaken} and becomes an interior high–cutoff RS1/dS cell, the same light mode becomes the stabilized radion of the interval. In that configuration the visible overlap need not be Robin–suppressed and could increase depending on the new boundary configuration; a detectable fifth force at range $\lambda_r\sim$kpc would then be possible unless further suppression (e.g.\ large $Z_{\rm vis}$, small mixing) is present.


\paragraph{Cosmological evolution (no late drift).}\label{par:uom-main-170}
Let \(m_g\sim L_4^{-1}\) be the AdS\(_4\) gap in the trace/radion–mixed channel. After the flip the
driver stops and all \(k{>}0\) components of the response decay as
\(\sim e^{-\Gamma t}\) with \(\Gamma=\mathcal O(m_g)\). The \emph{DC} component set by \(J_0\) does not
decay in the linear theory; hence the Higgs mass shift and \(v\) are time–independent thereafter. No
late–time EWSB drift is expected.

\medskip
\noindent\textbf{What to look for (summary).}
\begin{enumerate}[leftmargin=2.1em,label=\Roman*. ,itemsep=2pt]
\item \emph{CMB (large angles):} one compensated LoG–like feature at angular scale
$\theta_{\rm core}\sim \sigma_{x,1}=\frac{\alpha_x}{(v\,\Lambda_{\rm eff,1})}\sim\mathcal O(3^\circ\!-\!6^\circ)$,
with a possible \emph{fainter, narrower inner shoulder} at
$\theta_{\rm shoulder}\sim \sigma_{x,2}=\frac{\alpha_x}{(v\,\Lambda_{\rm eff,2})}\ll \theta_{\rm core}$.
The $\ell$–space template is fixed (band–locked LoG), with definite TE phase and a cold–area/kurtosis excess
peaking at $\theta_{\rm core}$. Use a LoG/SMHW matched filter to localize the core; then, if a
shoulder exists, constrain the \emph{effective scale and payload ratios} via
\[
\frac{A_{\rm shoulder}}{A_{\rm core}}
\;\approx\;\frac{\mathcal T_{2}}{\mathcal T_{1}}\,
\frac{\sum_{k\in\mathcal K_2}K(\bar\Lambda_k)\Delta I_k}{\sum_{k\in\mathcal K_1}K(\bar\Lambda_k)\Delta I_k}\!,
\]
(cf.\ \eqref{eq:ring-core-bandlocked}), and simultaneously fit the \emph{angular} ratio
$\theta_{\rm core}/\theta_{\rm shoulder}\simeq \Lambda_{\rm eff,2}/\Lambda_{\rm eff,1}$. A non-detection of a shoulder
constrains the same combination once $\mathcal T_{2}/\mathcal T_{1}$ is bracketed.

\item \emph{Parameter band:} the observed LoG amplitude at $\theta_{\rm core}$ fixes a \emph{lower bound} on the
\emph{core subset cutoff–weighted export}
$S_{\mathcal K_1}:=\sum_{k\in\mathcal K_1}K(\bar\Lambda_k)\Delta I_k$ with
$K(\Lambda)$ as in \eqref{eq:K-final-scaling} (see \eqref{eq:R-evolution}–\eqref{eq:BY-k}). Equivalently, the CMB–inferred
core amplitude yields $\sum_{k\in\mathcal K_1}K(\bar\Lambda_k)\Delta I_k$ up to a smooth transfer factor. Propagate
priors on the (dimensionless) trace–kernel average $\langle|\mathcal K|^2\rangle*$
(and on the SW/LSS transfer $\mathcal T$) to bracket the \emph{portal–relevant DC normalizer} $J_0$ in the
Higgs–portal relation (trace/anomaly imprint), i.e.\ use
\[
\Theta_k=K(\bar\Lambda_k)\,\Delta I_k
\quad\Longrightarrow\quad
S_{\mathcal K_1}\ \propto\ \sum_{k\in\mathcal K_1}K(\bar\Lambda_k)\Delta I_k,
\]
to convert the LoG amplitude into a quantitative \emph{band} for $J_0$ once $\langle|\mathcal K|^2\rangle*$ and
$\mathcal T$ are bounded from independent theory/transfer priors.
\item \emph{Colliders:} essentially SM–like Higgs couplings; deviations \(\delta\kappa_h\ll 10^{-3}\)
expected for \(m_{\mathcal S}\gg v\).
\item \emph{Equivalence principle:} no anomalous long–range force from the receiver sector on the
dS brane KR head.
\item \emph{No drift:} constant EWSB scale after ring–down; any observed secular drift of \(v\) would
falsify the minimal portal or imply light moduli beyond the KR point.
\end{enumerate}
\paragraph{Cold–Spot identification as the CMB imprint.}\label{par:uom-main-171}
Planck/WMAP find a single azimuthally–symmetric, few–degree feature with SMHW
peak scale \(R\simeq(3.3^\circ\!-\!5^\circ)\) and coefficient \(\nu\simeq-4.6\) at \(R\simeq 4.2^\circ\).

\subsection{Shadow--brane dark matter: detailed phenomenology and derivations}
\label{subsec:shadow-brane-dm-phenom}

This subsection develops the quantitative phenomenology of the \emph{shadow brane}
contribution to structure and lensing on our brane. We proceed in four steps:
(i) specify the cross--brane transfer kernel and its RS1-only realization;
(ii) derive the linear–growth equation on our brane with the shadow contribution as a
scale– and time–dependent source in the quasi–static/DC regime;
(iii) formulate compressed constraints for CMB and LSS (amplitudes and shapes);
(iv) pass to real space and extract the inner core scale of the kernel.

\paragraph{Notation and background.}\label{par:uom-main-172}
The late–time background is spatially flat with
\begin{equation}
H(a)=H_0\,E(a),\qquad E(a):=\sqrt{\Omega_{m0}\,a^{-3}+\Omega_{\Lambda 0}},\qquad
\Omega_{m}(a)=\frac{\Omega_{m0} a^{-3}}{E(a)^2},
\end{equation}
and $a_\star\simeq 1/1100$ denotes recombination. We split today’s matter budget into an
\emph{intrinsic} part (living on our brane) and a \emph{mirage} part $\Omega_X(a)$ sourced
by the shadow brane through the transfer kernel. For large–scale linear theory, only the
combination $\Omega_{m,\mathrm{int}}(a)+\Omega_X(k,a)$ enters the Poisson source.

\subsubsection*{Kernel and geometry (RS1 only)}\label{subsubsec:uom-main-016}

\begin{definition}[Cross–brane transfer kernel]\label{def:kernel}
Let $\rho_2$ denote the (shadow–brane) matter density and $\delta_2$ its fractional perturbation.
In mixed $(\omega,k)$ space the effective contribution to our Poisson source is
\begin{equation}\label{eq:kernel-master}
\delta\rho^{\rm eff}_1(\omega,k)\ =\ \varepsilon(\omega,k;a)\ \delta\rho_2(\omega,k),
\qquad
\boxed{\;\varepsilon(\omega,k;a)\ =\ \mathcal C(\theta)\,W_{\Lambda_1}(k)\,W_{\Lambda_2}(k)\,
 e^{-\,d(a)\,q(\omega,k)}\;\frac{q(\omega,k)}{q(\omega,k)+r(\Lambda)\,k^2}\;},
\end{equation}
with $q(\omega,k)=\sqrt{k^2+m_\bullet^2-\omega^2}$ and $r(\Lambda)=M_4^2(\Lambda)/(2M_5^3)$.
\end{definition}

\noindent\textbf{RS1 shadow branes (only).}
All matter-carrying shadow branes relevant to the dark sector are interior/shadow RS1-type dS$_4$
segments lying deeper IR than the (already RS1) source brane. The trace channel is governed by the
GW-stabilized radion of mass $m_{\rm rad}>0$, so in the transfer kernel we set $m_\bullet=m_{\rm rad}$.
KR heads are vacuum EOW caps and do not carry dark matter; they enter only as receiver/terminal boundaries.
The \emph{spatial} trace response at fixed time is a low--pass with Yukawa propagation across the bulk and
induced--gravity screening on the visible brane:
\begin{equation}
{\cal K}_{\rm tr}(\omega,k;\Lambda)\;=\;\frac{q(\omega,k)}{q(\omega,k)+r(\Lambda)\,k^2}\;\exp\!\big[-\,d\,q(\omega,k)\big]\;W_\Lambda(k),
\qquad q(\omega,k):=\sqrt{k^2+m_\bullet^2-\omega^2}.
\end{equation}

\subsubsection*{From the kernel to a linear growth equation}\label{subsubsec:uom-main-017}

\paragraph{Effective matter fraction.}\label{par:uom-main-173}
At the background (DC) level the shadow dust contributes a uniform fraction
\begin{equation}\label{eq:OmegaX-DC}
\Omega_X(a)\ =\ \varepsilon_{\rm DC}(a)\,\frac{\Omega_{2,0}\,a^{-3}}{E(a)^2}\,,
\qquad \varepsilon_{\rm DC}(a)=\mathrm{e}^{-m_\bullet\,[\,d(a)-d(1)\,]},
\end{equation}
so that $\Omega_X(1)=\varepsilon_0\,\Omega_{2,0}$ with $\varepsilon_0:=\varepsilon_{\rm DC}(1)$.
In the quasi–static, scale–dependent regime,
\begin{equation}\label{eq:OmegaX-k}
\boxed{\quad
\Omega_X(k,a)\ =\ \Omega_X(a)\;\mathrm{e}^{-\,d(a)\,q(0,k)}\;\frac{q(0,k)}{q(0,k)+r(\Lambda)\,k^2}\,,
\qquad q(0,k)=\sqrt{k^2+m_\bullet^2}\,.
\quad}
\end{equation}

\paragraph{Linear growth on our brane.}\label{par:uom-main-174}
In Newtonian gauge and to linear order, the growth of the \emph{our–brane} matter mode $D(k,a)$ obeys
\begin{equation}\label{eq:growth-ode}
\frac{d^2 D}{d(\ln a)^2}\ +\ \Bigl[\,2+\frac{d\ln H}{d\ln a}\Bigr]\frac{d D}{d\ln a}
\ -\ \frac{3}{2}\,\Bigl[\Omega_{m,\mathrm{int}}(a)\ +\ \Omega_X(k,a)\Bigr]\,D\ =\ 0.
\end{equation}
For $\Lambda$CDM backgrounds one has $d\ln H/d\ln a=-\tfrac{3}{2}\Omega_m(a)$, so \eqref{eq:growth-ode} closes given \eqref{eq:OmegaX-k}.

\begin{remark}[On conservation and cross–correlations]
\label{rem:uom-main-009}
The effective source $\delta\rho^{\rm eff}_1=\varepsilon\,\delta\rho_2$ is used in the quasi–static/DC regime where the associated
pressure/velocity potentials respect 4D Bianchi identities. The compressed constraints below probe precisely (i) the DC limit
(CMB) and (ii) the roll–off band where the kernel dominates (LSS shape); short–distance details of $\delta_2$ are immaterial.
\end{remark}


\subsubsection*{Compressed constraints for CMB and LSS}\label{subsubsec:uom-main-018}

We extract three simple, data–facing inequalities that map transparently to $(\Omega_{2,0},\varepsilon_0,m_\bullet,r(\Lambda))$.

\paragraph{(CMB) DC bound at recombination.}\label{par:uom-main-175}
Define the DC fraction at $a_\star$,
\begin{equation}\label{eq:fXstar}
f_X^\star\ :=\ \Omega_X(k\!\to\!0,a_\star)
\ =\ \frac{\Omega_{2,0}}{\Omega_{m0}}\ \varepsilon_{\rm DC}(a_\star)\,.
\end{equation}
Requiring a small shift of $r_s$ and a modest early–ISW implies
\begin{equation}\label{eq:cmb-bound}
\boxed{\qquad f_X^\star\ \le\ f_X^{\star,\max}\qquad \text{with}\qquad f_X^{\star,\max}\in[0.05,\,0.10]\ \text{(conservative)}\ .\qquad}
\end{equation}
For fixed $\Omega_{2,0}$ this constrains the dimensionless separation growth
\begin{equation}\label{eq:Xi-def}
\Xi\ :=\ m_\bullet\,[\,d(1)-d(a_\star)\,]\ =\ \ln\frac{\varepsilon_{\rm DC}(a_\star)}{\varepsilon_0}\ 
\le\ \ln\!\frac{f_X^{\star,\max}\,\Omega_{m0}}{\varepsilon_0\,\Omega_{2,0}}\,.
\end{equation}

\paragraph{(LSS) Large–scale amplitude bound.}\label{par:uom-main-176}
Let
\begin{equation}\label{eq:R-def}
R(k)\ :=\ \frac{D(k,a_0)}{D_{\Lambda{\rm CDM}}(k,a_0)}-1\,,
\end{equation}
and pick a representative quasi–linear wavenumber $k_L\sim 0.1\,h\,\mathrm{Mpc}^{-1}$. Integrating \eqref{eq:growth-ode} by parts yields
\begin{equation}\label{eq:R-ineq}
R(k_L)\ \simeq\ \int_{a_{\rm ini}}^{1}\frac{da}{a}\,\frac{3}{2}\,
\frac{\Omega_X(k_L,a)}{\Omega_{m,\mathrm{int}}(a)}\ \mathcal W(a;k_L)\ \le\ \delta_{\rm LSS}^{\max},
\end{equation}
with a smooth weight $\mathcal W$ of order unity and $\delta_{\rm LSS}^{\max}$ a few percent. This caps the combination
$\varepsilon_0\,\Omega_{2,0}$ for given $(r(\Lambda),m_\bullet)$.

\paragraph{(LSS) Bias–slope bound.}\label{par:uom-main-177}
Demand a small shape modulation across the quasi–linear window,
\begin{equation}\label{eq:slope-ineq}
\boxed{\qquad \Bigl|\frac{d\ln D(k,a_0)}{d\ln k}\Bigr|_{k\in[0.05,\,0.3]\,h\,\mathrm{Mpc}^{-1}}\ \le\ \beta_{\rm LSS}^{\max}\qquad (\text{few}\times 10^{-2})\ .\qquad}
\end{equation}
Because $D(k)$ rolls where the screen $q/(q+r k^2)$ turns over (roughly when $k\,r_{\rm scr}\sim 1$), \eqref{eq:slope-ineq}
implies $r_{\rm scr}\lesssim \mathcal O(3)\ \mathrm{Mpc}/h$ if one wants the roll to begin past the upper end of the quasi–linear window.


\subsubsection*{Real–space kernel and halo cores}\label{subsubsec:uom-main-019}

The real–space Green’s function is the convolution in \eqref{eq:conv-kernels}. Its inner scale is governed by the shortest
length among the induced–gravity screen and the bulk Yukawa:
\begin{statement}[Core scale from the convolution kernel]\label{prop:kernel-real}
With the DC kernel in $k$–space 
$\widehat{\Phi}_{\rm sh}(k)\propto k^{-2}\,\frac{q(0,k)}{q(0,k)+r k^2}\,\mathrm{e}^{-d\,q(0,k)}$,
the (Green’s–function) core radius $r_{\rm core}$ (defined e.g.\ by $\Phi_{\rm sh}(r_{\rm core})/\Phi_{\rm N}(r_{\rm core})=\tfrac12$)
satisfies
\begin{equation}\label{eq:rcore-max}
\boxed{\qquad
r_{\rm core}\ \simeq\ C\;\max\!\Bigl\{\,r_{\rm scr},\;m_\bullet^{-1},\;\Lambda^{-1}\Bigr\}\,,\qquad
r_{\rm scr}:=\sqrt{r(\Lambda)/m_\bullet},\quad C=\mathcal O(1).
\qquad}
\end{equation}
\end{statement}

\paragraph{Dwarf and cluster bands.}\label{par:uom-main-178}
Two robust phenomenological bands follow from \eqref{eq:rcore-max}:
\begin{itemize}[leftmargin=1.6em]
\item \emph{Dwarfs.} Inner rotation curves in dwarf/LSB galaxies rise almost linearly with $r$, indicative of a constant–density core of order
\begin{equation}\label{eq:dwarf-band}
r_{\rm core}^{\rm (dwarf)}\ \sim\ 0.5{-}3\ \mathrm{kpc}.
\end{equation}
If $r_{\rm scr}\gg \mathrm{kpc}$ (as preferred by LSS), dwarf cores fix the Yukawa length: $m_\bullet^{-1}\sim r_{\rm core}^{\rm (dwarf)}$ so that
\begin{equation}\label{eq:mr-from-dwarf}
m_\bullet\ \simeq\ \frac{C}{r_{\rm core}}\ \in\ (0.3{-}2)\ \mathrm{kpc}^{-1},\qquad C=\mathcal O(1).
\end{equation}
\item \emph{Clusters.} Strong lensing in clusters typically requires $r_{\rm core}^{\rm (cluster)}\lesssim 30{-}50$ kpc; this is automatically satisfied for $m_\bullet^{-1}=\mathcal O(\mathrm{kpc})$.
\end{itemize}

\subsubsection*{Putting it together: a self–consistent wedge}\label{subsubsec:uom-main-020}

The constraints \eqref{eq:cmb-bound}, \eqref{eq:R-ineq}, \eqref{eq:slope-ineq} and the halo bands
\eqref{eq:dwarf-band}–\eqref{eq:mr-from-dwarf} define a viable region in
$(\Omega_{2,0},\varepsilon_0,m_\bullet,r(\Lambda))$:
\begin{itemize}[leftmargin=1.6em]
\item \textbf{Today’s mirage:} $\Omega_X(1)=\varepsilon_0\,\Omega_{2,0}\sim 0.5{-}1\%$.
\item \textbf{Recombination fraction:} $f_X^\star=(\Omega_{2,0}/\Omega_{m0})\,\varepsilon_{\rm DC}(a_\star)\sim 5{-}10\%$.
\item \textbf{Quasi–linear roll:} $r_{\rm scr}\sim 3\,\mathrm{Mpc}/h$, keeping the bias slope small across $k\in[0.05,0.3]\,h\,\mathrm{Mpc}^{-1}$.
\item \textbf{Halo cores:} $m_\bullet^{-1}\sim (0.3-2)\; \mathrm{kpc}$, yielding $r_{\rm core}^{\rm (dwarf)}\sim 0.5{-}3\,\mathrm{kpc}$ and $r_{\rm core}^{\rm (cluster)}\ll 50\,\mathrm{kpc}$.
\end{itemize}
In RS1 one identifies $m_\bullet=m_{\rm rad}$ (Goldberger--Wise radion, trace channel). The spin--2 gap
$m_g\sim L_4^{-1}$ controls only the TT sector and is not used in the scalar transfer kernel.
All formulas below use the scalar trace kernel and assume this identification.

\subsubsection*{A practical recipe (``compressed likelihood'')}\label{subsubsec:uom-main-021}
For phenomenology, it is convenient to package the above as a quick acceptance test for a parameter point $\theta=\{\Omega_{2,0},\varepsilon_0,m_\bullet,r(\Lambda),\Lambda\}$:

\begin{enumerate}[itemsep=2pt]
\item Compute $\varepsilon_{\rm DC}(a_\star)=\exp[-m_\bullet(d(a_\star)-d(1))]$; then evaluate $f_X^\star$ by \eqref{eq:fXstar} and enforce \eqref{eq:cmb-bound}.
\item Form $\Omega_X(k,a)$ via \eqref{eq:OmegaX-k}; integrate \eqref{eq:growth-ode} for $D(k,a)$ and check the amplitude bound $R(k_L)\le \delta_{\rm LSS}^{\max}$ and slope bound \eqref{eq:slope-ineq}.
\item Compute $r_{\rm core}$ from \eqref{eq:rcore-max} and enforce \eqref{eq:dwarf-band}–\eqref{eq:mr-from-dwarf}.
\end{enumerate}
Points that pass these compressed constraints give a $\sim 10\%$ enhancement of large–scale growth that
smoothly rolls toward $\Lambda$CDM by $k\sim 0.2{-}0.3\,h\,\mathrm{Mpc}^{-1}$ and cored haloes with kpc dwarfs and
cluster cores $\ll$ few $\times 10$ kpc. These follow directly from the kernel \eqref{eq:kernel-master} and its DC specialization \eqref{eq:OmegaX-DC}.

\subsubsection*{Observational anchors and constraints}\label{subsubsec:uom-main-022}

\begin{enumerate}[leftmargin=1.4em,itemsep=2pt]
\item \textbf{CMB acoustic scale and early-ISW consistency.}
Percent-level additions to the pre-recombination matter budget are tightly bounded by
Planck’s acoustic scale, lensing, and the early-ISW pattern, keeping any extra,
pressureless-like component at the \emph{few-percent} level at $z\!\sim\!1100$.
See Planck 2018 cosmological parameters \cite{Planck2018} and related early-intervention
analyses \cite{Poulin2019,Hill2020}. BAO and RSD measurements similarly pin the background
expansion and growth history \cite{Alam2021eBOSS}.

\item \textbf{Full-shape $P(k)$ constraints in the quasi-linear window.}
Analyses of the BOSS power spectrum using EFT of LSS show that
scale-dependent deviations in the range $k\simeq 0.05{-}0.25\,h\,\mathrm{Mpc}^{-1}$ must
be \emph{percent-level} after marginalizing over bias and EFT counterterms
\cite{Ivanov2020,DAmico2020}. This sets a practical ceiling for any low-pass “mirage”
response in that $k$-window.

\item \textbf{Dwarf cores (sub-kpc to few-kpc).}
High-resolution rotation curves and mass models of gas-rich dwarfs (e.g. LITTLE THINGS)
and the SPARC compilation exhibit approximately constant-density cores spanning
$\sim\!0.5{-}3$ kpc, with systems both cusp-like and cored in a \emph{continuous} spread
\cite{Oh2015,Lelli2016SPARC,Read2019}. Classical dSphs like Fornax and Sculptor provide
independent, stellar-dynamical evidence for $\mathcal{O}(\,\mathrm{kpc}\,)$ cores
\cite{WalkerPenarrubia2011}.

\item \textbf{Cluster inner slopes / cores (tens of kpc).}
Joint strong+weak lensing and BCG stellar dynamics in relaxed clusters favor shallow inner
dark-matter slopes (or effective cores) on $\sim\!10{-}50$ kpc scales, consistent with the
kernel–induced low-pass at cluster radii \cite{Newman2013I,Newman2013II}.
\end{enumerate}

\subsection{Shadow–brane stack unified model}
\label{subsec:stack-fit-unified}

We present a single, unified framework (no ``front vs.\ tail'' split) that ties together the stack geometry, the aeon–to–aeon renormalization (cutoff growth), and the phenomenology of dwarf cores. The model yields closed–form expressions for the \emph{mean} and \emph{variance} of the central–slope sum that controls the fitted core radius, hence an analytic formula for the \emph{scatter/mean} ratio. It also implies a \emph{unimodal} (left–skewed) population for core sizes and predicts a logit–normal observational fit.

\subsubsection*{Grounding the landscape: separations, cutoff, and curvature}
\label{subsubsec:grounding}

\paragraph{Constant proper spacing (AdS$_5$ energy conservation).}\label{par:uom-main-179}
Once a brane becomes a shadow (RS1 interior with frozen dynamics), energy conservation in AdS$_5$ slows and halts further expansion in the stack direction to first order. We therefore model the stack with an \emph{approximately constant} proper spacing
\begin{equation}
d^{(n)} \;=\; d^{(1)} + (n-1)\,\Delta d,\qquad \Delta d=\text{const},
\label{eq:spacing}
\end{equation}
for branes labeled by depth $n=1,2,\dots$ (with $n=1$ the nearest shadow).

\paragraph{Cutoff growth across aeons.}\label{par:uom-main-180}
On each aeon, the receiver cutoff grows tick–by–tick under the wall balance; as $\bar\Lambda_k$ increases, the conformality/RP penalties decay and the growth can slow. Across aeons this induces a monotone increase that is well captured by the effective trend
\begin{equation}
\Lambda^{(n+1)}\;\approx\;\eta\,\Lambda^{(n)},\qquad \eta\ \gtrsim\ 1,
\label{eq:cutoff-growth}
\end{equation}
with $\eta\approx 1$ in the hard–freeze limit.

\paragraph{Curvature (gap) contraction across aeons.}\label{par:uom-main-181}
The shadow-stack trace channel is gapped by the GW-stabilized RS1 radion, with a scale proportional to the inverse brane curvature $L_4^{-1}$. The brane curvature expectation obeys the envelope/infimum bounds of Sec.~\ref{subsubsec:geo-contract}. When the contraction is approximately multiplicative one may parameterize it as
\begin{equation}
m_\bullet^{(n+1)}\;=\;\beta\,m_\bullet^{(n)},\qquad 0<\beta\ \lesssim\ 0.1,
\label{eq:gap-decay}
\end{equation}
where $m_\bullet$ is the RS1 radion mass in the shadow segments.

\paragraph{Induced–gravity screen (low–pass knee).}\label{par:uom-main-182}
Integrating out boundary/CFT modes to the cutoff $\Lambda$ induces a 4D Einstein term, giving a DGP/RS2–like screen
\begin{equation}
\frac{1}{\,1+r_c^{(n)2} k^2\,},\qquad 
r_c^{(n)}\;=\;\frac{M_4^2(\Lambda^{(n)})}{2M_5^3}
=\frac{\ell}{2}+\frac{\kappa a}{2M_5^3}\big(\Lambda^{(n)}\big)^2,
\label{eq:screen}
\end{equation}
so that $r_c^{(n)}$ grows with $n$ (older branes are increasingly low–pass).

\subsubsection*{Unified model of the stack and the core–setting kernel}
\label{subsubsec:unified-model}

\paragraph{Core–setting band and rung weights.}\label{par:uom-main-183}
Let $k_\star$ be a representative wavenumber for the inner potential/core–setting band (we use $k_\star\sim m_\bullet$, see below). The transfer amplitude from brane $n$ into that band is
\begin{equation}
\chi_n\ :=\ \chi^{(n)}(k_\star)
=\ \frac{\exp\!\big[-\,d^{(n)}\sqrt{m_\bullet^{(n)2}+k_\star^2}\big]}{\,1+r_c^{(n)2}k_\star^2\,}\ \times\ W_{\Lambda^{(n)}}(k_\star).
\label{eq:chi-def}
\end{equation}
For the \emph{top} rungs ($n$ small) with $k_\star\sim m_\bullet^{(1)}$ one has $\sqrt{m_\bullet^{(n)2}+k_\star^2}\approx k_\star$ and $W_{\Lambda^{(n)}}(k_\star)\approx 1$, so the weights are \emph{geometric}:
\begin{equation}
\chi_n\;\approx\; \underbrace{\frac{e^{-k_\star d^{(1)}}}{1+r_c^{(1)2}k_\star^2}}_{=:C_1}\;
\Big(\underbrace{\frac{e^{-k_\star \Delta d}}{\eta^2}}_{=:\,\rho}\Big)^{n-1},
\qquad 0<\rho<1.
\label{eq:chi-geo}
\end{equation}
The single ratio $\rho$ encodes distance attenuation per step and the induced--gravity screen growth.

\paragraph{Continuous per–brane centralization (no on/off spikes).}\label{par:uom-main-184}
At today’s time $t_0$, the centralization on brane $n$ is modeled as a \emph{monotone subordinator} (drift $+$ positive jumps),
\begin{equation}
\Upsilon_n\ =\ \mu_s\,g(\tau_n)\ +\ \sum_{k=1}^{N_n} A_{nk}\,\Phi(b_{nk}/L),\qquad 
\tau_n=\text{``age'' of brane }n,\quad L=\frac{1}{m_\bullet^{(1)}},
\label{eq:upsilon}
\end{equation}
with $g$ non–decreasing, $N_n\sim\mathrm{Poisson}(\Lambda(\tau_n))$, jumps $A_{nk}>0$ (e.g.\ lognormal), and offset kernel $\Phi\in(0,1]$. The brane’s \emph{observed} contribution in the core band is then
\begin{equation}
\xi_n\ :=\ \chi_n\,\Upsilon_n,
\qquad 
\mu_n:=\mathbb E[\xi_n]=\chi_n\,S(\tau_n),\quad
\sigma_n^2:=\mathrm{Var}(\xi_n)=\chi_n^2\,V(\tau_n),
\label{eq:xi-mom}
\end{equation}
where $S(\tau):=\mathbb E[\Upsilon(\tau)]$ and $V(\tau):=\mathrm{Var}(\Upsilon(\tau))$. A mild, data–friendly specialization consistent with ``older $\Rightarrow$ more clumping'' is the \emph{geometric growth} ansatz
\begin{equation}
S(\tau_n)=S_1\,s^{\,n-1},\qquad V(\tau_n)=V_1\,s^{\,n-1},\qquad s\gtrsim 1.
\label{eq:geo-age}
\end{equation}
(If one wishes to retain near–center \emph{incidence} $p_n\in(0,1]$ as a thinning factor, it can be absorbed as $S(\tau_n)\mapsto p_n S(\tau_n)$ and $V(\tau_n)\mapsto p_n^2 V(\tau_n)$, or treated exactly via the Poisson–binomial moments; see below.)

\paragraph{Central–slope sum and core size.}\label{par:uom-main-185}
The central–slope that sets the fitted core is
\begin{equation}
Z\ :=\ \sum_{n\ge 1}\xi_n,\qquad R\ :=\ \frac{c}{1+Z},\qquad c=\frac{\kappa}{m_\bullet^{(1)}}\simeq r_{\rm core},\ \ \kappa\simeq\ln 2,
\label{eq:sum-R}
\end{equation}
with core inference $R$ monotonically decreasing in $Z$.

\subsubsection*{Closed–form scatter/mean ratio}
\label{subsubsec:sdmean}

Under \eqref{eq:chi-geo}–\eqref{eq:geo-age} the rung moments are geometric: $\mu_n=C_1 S_1 (\rho s)^{n-1}$ and $\sigma_n^2=C_1^2 V_1 (\rho^2 s)^{n-1}$. For independent rungs,
\begin{equation}
\mu_Z=\sum_{n\ge 1}\mu_n=\frac{C_1S_1}{1-\rho s},\qquad
\sigma_Z^2=\sum_{n\ge 1}\sigma_n^2=\frac{C_1^2V_1}{1-\rho^2 s},
\qquad \text{if}\ \rho s<1,\ \rho^2 s<1.
\label{eq:muZ-sigZ}
\end{equation}
A second–order delta expansion of $R=c(1+Z)^{-1}$ around $\mu_Z$ gives (to excellent accuracy in our wedge)
\begin{equation}
\boxed{\qquad
\frac{\mathrm{sd}(R)}{\mathbb E[R]}\ \approx\ \frac{\sigma_Z}{\,1+\mu_Z\,}
=\ \frac{ \dfrac{C_1\sqrt{V_1}}{\sqrt{1-\rho^2 s}} }{ 1+\dfrac{C_1 S_1}{1-\rho s} }\,.
\qquad}
\label{eq:sd-over-mean-unified}
\end{equation}
\emph{Remark.} If one keeps an explicit Bernoulli thinning $X_n\sim\mathrm{Bernoulli}(p_n)$ (near–center incidence), then with independence
\[
\mu_Z=\sum_n p_n \mu_n,\qquad
\sigma_Z^2=\sum_n \big(p_n \sigma_n^2 + p_n(1-p_n)\mu_n^2\big)\,,
\]
and the same sd/mean formula holds with these moments; the sums remain geometric if $p_n$ decays exponentially with depth.

\subsubsection*{\texorpdfstring{Grounded wedge numerics for $(\Lambda^{(n)},m_\bullet^{(n)},d^{(n)})$ and $(\xi_n,p_n,N_{\rm eff})$ (calibrated to $\overline{R}=1.25~\mathrm{kpc}$)}{Grounded wedge numerics for (Lambda^(n),m_bullet^(n),d^(n)) and (xi_n,p_n,N_eff) (calibrated to Rbar=1.25 kpc)}}
\label{subsubsec:wedge-numerics}

We now pick a minimally structured wedge consistent with Sec.~\ref{subsubsec:grounding} and the observed dwarf cores, and we \emph{calibrate the intrinsic kernel scale} so that the \emph{mean fitted core} equals the target
\[
\overline{R}\;:=\;\mathbb E[R]\;=\;1.25~\mathrm{kpc}.
\]
Recall $R=c(1+Z)^{-1}$ with $c=\kappa/m_\bullet^{(1)}\simeq r_{\rm core}^{\rm (intr)}$, and the unified moments (for the core band with geometric rung ratio $\rho$ and age factor $s$)
\[
\mu_Z=\frac{C_1S_1}{1-\rho s},\qquad
\sigma_Z^2=\frac{C_1^2V_1}{1-\rho^2 s},\qquad
\frac{\mathrm{sd}(R)}{\mathbb E[R]}\ \approx\ \frac{\sigma_Z}{1+\mu_Z}\,.
\]
A second–order delta expansion of $\overline{R}$ gives
\[
\overline{R}\ \approx\ c\left[\frac{1}{1+\mu_Z}+\frac{\sigma_Z^2}{(1+\mu_Z)^3}\right]
= \frac{c}{1+\mu_Z}\,\Bigl[1+\Bigl(\frac{\sigma_Z}{1+\mu_Z}\Bigr)^{\!2}\Bigr],
\]
so the \emph{calibrated} kernel scale (and hence the intrinsic core) is
\begin{equation}
\boxed{\qquad
c_{\rm cal}\;=\;\overline{R}\,\frac{1+\mu_Z}{\,1+\bigl(\sigma_Z/(1+\mu_Z)\bigr)^{2}}\,,\qquad
m_\bullet^{(1)}=\frac{\kappa}{c_{\rm cal}}\,,\qquad \kappa=\ln 2.
\qquad}
\label{eq:ccal}
\end{equation}
This re–normalizes the absolute core scale without changing \emph{fractional} observables ($\mathrm{sd}/\mathbb E$, the law of $Y:=R/c$, etc.). In particular, in what follows we \emph{hold the core–band wavenumber} fixed at $k_\star$ (set once and for all by the physics of the inner potential), so that the geometric rung ratio $\rho$ and the prefactor $C_1$ depend only on the stack parameters, not on the rescaled $c_{\rm cal}$.

\paragraph{Stack geometry and cutoff cascade.}\label{par:uom-main-186}
We adopt the constant spacing and geometric cutoff growth from Sec.~\ref{subsubsec:grounding}:
\[
\Delta d\in[3.0,3.5]~\mathrm{kpc},\qquad k_\star\simeq 0.55~\mathrm{kpc}^{-1},\qquad 
e^{-k_\star\Delta d}\in[0.18,0.14],\qquad 
\eta\in[1.6,2.0]\Rightarrow \eta^2\in[2.56,4.00],
\]
\[
\Longrightarrow\quad \rho:=\frac{e^{-k_\star\Delta d}}{\eta^2}\in[0.05,0.10],\qquad 
s\in[1.2,1.4],\qquad \rho s\lesssim 0.14,\quad \rho^2 s\lesssim 0.02.
\]

\paragraph{Nearest–rung moments (continuous centralization).}\label{par:uom-main-187}
Choose a strong but plausible nearest \emph{mean} and \emph{variance} (continuous subordinator: drift $+$ jumps, not on/off):
\[
C_1 S_1\ \equiv\ \mu_1\ \in [4,6]\ \ (\text{baseline }5),\qquad
\mathrm{CV}_{\Upsilon_1}:=\frac{\sqrt{V_1}}{S_1}\in[0.8,1.2]\ \Rightarrow\ C_1\sqrt{V_1}\in[4,6].
\]
(If one retains a near–center thinning $p_n$, take $p_n=p_1 u^{n-1}$ with $p_1\simeq 0.6$, $u\in[0.85,0.95]$; in the continuous model one may set $p_n\equiv 1$.)

\paragraph{Effective contributors in the core band.}\label{par:uom-main-188}
With $\chi_n\propto\rho^{n-1}$, the effective number of contributors is
\[
N_{\rm eff}\ :=\ \frac{\big(\sum \chi_n\big)^2}{\sum \chi_n^2}\ =\ \frac{1+\rho}{1-\rho}\ \in\ [1.11,1.22],
\]
i.e.\ only the first few rungs materially feed the core band (deep rungs predominantly feed DC).

\paragraph{Calibrated $\overline{R}=1.25$ kpc and resulting $c_{\rm cal},m_\bullet^{(1)}$.}\label{par:uom-main-189}
For the wedge above, the geometric sums \eqref{eq:muZ-sigZ} yield
\[
\mu_Z\ \approx\ \frac{5}{1-\rho s}\ \in\ [5.5,\,5.9],\qquad 
\sigma_Z\ \approx\ \frac{5}{\sqrt{1-\rho^2 s}}\ \in\ [5.0,\,5.3],
\]
so that
\[
\frac{\sigma_Z}{1+\mu_Z}\ \in\ [0.65,\,0.72].
\]
Inserting into \eqref{eq:ccal} with $\overline{R}=1.25~\mathrm{kpc}$ gives
\[
\boxed{\quad
c_{\rm cal}\ =\ \overline{R}\,\frac{1+\mu_Z}{\,1+\bigl(\sigma_Z/(1+\mu_Z)\bigr)^2}
\ \in\ [\,5.1,\ 6.1\,]\ \mathrm{kpc}\,,
\qquad
m_\bullet^{(1)}\ =\ \frac{\ln 2}{c_{\rm cal}}\ \in\ [\,0.114,\ 0.136\,]\ \mathrm{kpc}^{-1}.
\quad}
\]
By construction, the \emph{mean} fitted core equals $1.25$ kpc; the \emph{fractional} statistics (sd/mean, unimodality, skewness) computed below are unchanged (they depend only on $\mu_Z,\sigma_Z$).

\paragraph{sd/mean (for reference).}\label{par:uom-main-190}
With the calibrated $c_{\rm cal}$, the scale–free ratio remains
\[
\boxed{\ \frac{\mathrm{sd}(R)}{\mathbb E[R]}\ \approx\ \frac{\sigma_Z}{1+\mu_Z}\ \in\ [\,0.65,\ 0.72\,]\ .\ }
\]

\subsubsection*{Unimodality and the skewness law}
\label{subsubsec:unimodality-skew}

\paragraph{Unimodality.}\label{par:uom-main-191}
Each $\xi_n=\chi_n\Upsilon_n$ is a \emph{continuous}, positive random variable (subordinator), and the convolution of such variables is unimodal under mild regularity. Since only a few rungs contribute materially in the core band (Sec.~\ref{subsubsec:wedge-numerics}), the sum $Z=\sum \xi_n$ is \emph{unimodal} and right–skewed; the monotone map $R=c/(1+Z)$ preserves unimodality and flips the skew, giving a \emph{unimodal, left–skewed} population of cores.

\paragraph{Skewness law.}\label{par:uom-main-192}
Let $\kappa_{3,n}:=\mathbb E[(\xi_n-\mu_n)^3]$ denote the third central moment (for a lognormal proxy with log–variance $\sigma_L^2$, $\kappa_{3,n}=\mu_n^3\,(e^{\sigma_L^2}-1)^{3/2}(e^{\sigma_L^2}+2)$). With independence,
\[
\mu_{3,Z}:=\mathbb E[(Z-\mu_Z)^3]=\sum_{n\ge 1}\kappa_{3,n}.
\]
A delta expansion of $R=c(1+Z)^{-1}$ about $\mu_Z$ yields
\begin{equation}
\boxed{\qquad
\gamma_1(R)\ :=\ \frac{\mathbb E[(R-\mathbb E R)^3]}{\mathrm{sd}(R)^3}
\ \approx\ -\,\frac{\mu_{3,Z}}{\sigma_Z^3}\ -\ \frac{6\,\sigma_Z}{1+\mu_Z}\,.
\qquad}
\label{eq:skew-law}
\end{equation}
In the symmetric–age wedge ($s$ modest, no strong $p_n$ asymmetry), the inherited term is small and the curvature term dominates, giving $\gamma_1(R)\sim -6\,\sigma_Z/(1+\mu_Z)\approx -4$ to $-5$: a visible left tail.

\subsubsection*{Closest observational family: the logit–normal}
\label{subsubsec:logitnormal}

Define $Y:=R/c\in(0,1)$; then
\[
\operatorname{logit}(Y)=\ln\frac{Y}{1-Y}=\ln\frac{1}{Z}=-\ln Z.
\]
For a sum of a few positive, right–skewed terms with moderate dispersion, $\ln Z$ is well approximated by a Gaussian (Fenton–Wilkinson). Matching the first two moments of $Z$ by a lognormal $\mathrm{LN}(\mu_W,\sigma_W^2)$ gives
\begin{equation}
\sigma_W^2=\ln\!\Bigl(1+\frac{\sigma_Z^2}{\mu_Z^2}\Bigr),\qquad 
\mu_W=\ln\mu_Z-\tfrac12\,\sigma_W^2,
\label{eq:FW}
\end{equation}
so that
\begin{equation}
\boxed{\qquad
\operatorname{logit}(Y)\ \approx\ -\,\mathrm{Normal}(\mu_W,\sigma_W^2)
\quad\Longleftrightarrow\quad
Y\ \sim\ \mathrm{LogitNormal}\bigl(-\mu_W,\sigma_W^2\bigr).
\qquad}
\label{eq:logitnormal}
\end{equation}
This two–parameter family captures the \emph{unimodal, left–skewed} shape and the \emph{heavy small–core tail} implied by \eqref{eq:skew-law}. In practice one fits $(\mu_W,\sigma_W^2)$ to the sample, recovers $(\mu_Z,\sigma_Z^2)$ via \eqref{eq:FW}, and then infers the ladder parameters $(\Delta d,\eta,s)$ and the nearest–rung moments $(C_1S_1,\,C_1\sqrt{V_1})$ through \eqref{eq:muZ-sigZ}–\eqref{eq:sd-over-mean-unified}.

\medskip
\noindent\emph{Summary.} The grounded landscape (constant $\Delta d$; cutoff growth; curvature gap decay) implies a single geometric rung ratio $\rho=e^{-k_\star\Delta d}/\eta^2$ in the core band. With continuous per–brane centralization growing geometrically with depth, the central–slope sum has closed–form moments, the core population is unimodal and left–skewed, and a natural logit–normal fit emerges. In a plausible wedge (Sec.~\ref{subsubsec:wedge-numerics}) one finds $\mathrm{sd}(R)/\mathbb E[R]\approx 0.7$ without invoking a front/tail split.

\subsubsection*{Observed distribution and scaling of inner cores}\label{subsubsec:uom-main-023}

\begin{enumerate}[leftmargin=1.4em,itemsep=2pt]
\item \textbf{Continuous (unimodal), left-skewed diversity.}
Large rotation-curve compilations (e.g.\ SPARC) and focused H{\sc i} dwarf surveys
(LITTLE THINGS) show a \emph{continuous} spread of inner slopes and core sizes—from
cusp-like to $\sim$kpc cores—rather than a bimodal split; the distribution is well
captured by a single skewed mode \cite{Lelli2016SPARC,Oh2015,Oman2015}.

\item \textbf{Core–size tracks galaxy size.}
Empirically, core radii correlate with stellar size proxies ($R_d$, $R_{1/2}$), e.g.\
Burkert/iso-core fits give $r_0\!\sim\!(2{-}3)R_d$ across late-type systems; see
scaling-law syntheses and constant–surface-density findings in \cite{KormendyFreeman2004,Donato2009}.
This motivates modeling the population on a \emph{dimensionless} ratio (e.g.\ $R/c$).

\item \textbf{dSph anchors (Fornax, Sculptor).}
Multiple–tracer stellar kinematics favor $\sim$kpc (Fornax) and sub-kpc (Sculptor) cores,
bracketing the inner scale seen in gas-rich dwarfs and reinforcing a \emph{single}, broad
population rather than distinct “core”/“cusp” classes \cite{WalkerPenarrubia2011,Read2019}.
\end{enumerate}

\section{Anti--Zeno Co--Test (AZCT) and Compiled Duality (UOMT)}
\label{sec:azct-uomt}

\subsection{\texorpdfstring{Interleaved $\epsilon$--dynamics, $\Lambda$--band lock, and constants}{Interleaved epsilon--dynamics, Lambda--band lock, and constants}}
\label{subsec:interleaved-eps}

\paragraph{Two coupled $\epsilon$--dynamics (non--commuting limits).}\label{par:uom-main-193}
On the receiver we track two small parameters that are \emph{jointly} driven to $0$:
\begin{enumerate}
  \item the receiver \emph{resolution/energy} scale
  $\epsilon_{\mathrm{res}}$ (the conjugate of boundary energy, i.e.\ the OS resolution of the receiver),
  which is reduced by the uplift/theory--change (regulated conformal intertwiner refinement);
  \item the receiver \emph{curvature} scale
  $\epsilon_{\mathrm{curv}}:=|R|$ (equivalently $|\Lambda_4|$), which is reduced by the Weyl compensator/Brown--York flux (anomaly floor increments).
\end{enumerate}
The limits $\epsilon_{\mathrm{res}}\to 0$ and $\epsilon_{\mathrm{curv}}\to 0$ \emph{do not commute}:
if one attempts to cross $R=0$ at \emph{infinite} resolution ($\epsilon_{\mathrm{res}}=0$ first), reflection-positivity can be lost at the crossing
(cf.\ the non--commuting limits analysis in Theorem~\ref{thm:noncommute}).
These two coupled arrows are derived semiclassically from the SK/HJ structure in Sec.~\ref{sec:two-times-semiclassical}.
Operationally, the scheduler (P3) executes the reduction \emph{interleaved} in tick time, so that OS positivity holds for every finite tick; see also Theorem~\ref{thm:OS_receiver}.

\paragraph{$\Lambda$--band lock (receiver--side).}\label{par:uom-main-194}
Welded (receiver) temporal/spatial widths $(s_{t,k},s_{x,k})$ are chosen \emph{on the receiver} and are \emph{locked to the band edge} of the current cutoff $\Lambda_k$.
Source macro widths $(\Sigma_k,\rho_k)$ are chosen on the source and play no role in conformality on the receiver.
Conformality preservation on the radiative sector at each tick follows directly from:
(i) AZCT acceptance (radiative projection at the source end),
(ii) the \emph{regulated conformal intertwiner} on the $\Lambda_k$--band at the receiver end,
and (iii) time windowing by a zero--DC derivative--of--Gaussian (DoG) envelope to avoid DC offsets and acausal artifacts.
No time--alignment lock is required to obtain conformality.

\paragraph{OS--positivity on the receiver.}\label{par:uom-main-195}
As proved in Theorem~\ref{thm:OS_receiver}, the receiver is Osterwalder--Schrader positive under all admissible uplifts.
Consequently, reflection positivity (RP) does \emph{not} enter as a guard in the present section; it is finite a priori for each finite tick and vanishes in the anti--Zeno limit of the regulators.

\paragraph{Notation and constants.}\label{par:uom-main-196}
We write $\mathrm{dist}_\partial:=\tfrac12\|\cdot\|_1$ for the boundary trace distance,
$R_{\Lambda}(\rho_{\mathrm dS}):=\|(\mathbf 1-P_{\le\Lambda})\rho_{\mathrm dS}\|_1$ for the source residual at band $\Lambda$,
and $k_\ast$ for the first (``flip'') tick identified below.
Tickwise we will use ``sandwich'' constants $a_k>0$, $b_k>0$, $c_k>0$ (defined in Lemmas~\ref{lem:lower-sandwich}--\ref{lem:upper-sandwich})
and a paired selection of guards
\begin{equation}
\label{eq:paired-guards}
\varepsilon_n^\dagger=\frac{\varepsilon_n}{a_k},\qquad
\varepsilon_n=c_k\,\varepsilon_n^\dagger,
\end{equation}
uniformly valid for all ticks when $a_k$ and $c_k$ are bounded away from $0$ and $\infty$ under the regulator schedule.

\subsection{AZCT: definition, pre--tick acceptance, always--accept pre--flip}
\label{subsec:azct-def}

\paragraph{Spaces and residual.}\label{par:uom-main-197}
We factor the co--test as a sextuple
\[
\mathfrak C=\bigl(\mathcal M^\dagger,\ \mathcal Z_{\mathrm{in}}^\dagger,\ \mathcal Z_{\mathrm{out}}^\dagger,\ U^\dagger,\ E^\dagger,\ \mathcal S^\dagger\bigr),
\]
with $\mathcal M^\dagger=\mathcal D(\mathcal H_{\mathrm{env}})$ the persistent KMS bath (at $T_{\mathrm{GH}}=H/(2\pi)$ with gap $m_\ast>0$),
$\mathcal Z_{\mathrm{in}}^\dagger=\mathcal D(\mathcal H_{\mathrm dS}^{\Lambda})$ the band--limited dS input, and
$\mathcal Z_{\mathrm{out}}^\dagger=\mathcal D(\mathcal H_{\mathrm{logs}})$ the compiled payload/log.
With $Q:=\mathbf 1-P_{\le\Lambda}$ the residual is
\begin{equation}
\label{eq:residual-def}
R_\Lambda(\rho_{\mathrm dS}):=\|Q\,\rho_{\mathrm dS}\|_1 .
\end{equation}
The co--standard is $\mathcal S^\dagger=\{\varepsilon_n^\dagger\}_{n\ge 1}$ and numerical tolerances are absorbed in a margin $\eta_{\mathrm{num}}$.

\paragraph{Time--modulated QND with decaying backflow.}\label{par:uom-main-198}
On $\mathcal H_{\mathrm dS}^{\Lambda}\otimes\mathcal H_{\mathrm{env}}$ we evolve on a tick window $I_k=[t_k-\Delta_k,t_k+\Delta_k]$ by
\begin{align}
\frac{d}{dt}\rho_{SE}(t)
&=-\,i[\,H_{\mathrm{sp},\Lambda}+H_{\mathrm{env}},\ \rho_{SE}(t)\,]
+\sum_a\kappa_a^{(k)}(t)\Big( L_a\rho_{SE}(t)L_a^\dagger - \tfrac12\{L_a^\dagger L_a,\rho_{SE}(t)\}\Big)
\label{eq:azct-gkls}\\[-0.5em]
&\qquad\qquad\qquad\qquad\qquad\qquad+ \int_{t_k-\Delta_k}^{t}\! K(t-t')\,\mathcal B[\rho_{SE}(t')]\,dt',
\nonumber
\end{align}
where the pointers are QND, $[L_a,H_{\mathrm{sp},\Lambda}]=0$, and the macro absorption rates are DoG--shaped
$\kappa_a^{(k)}(t)=u_a\,\chi_{\Sigma_k}(t-t_k)$ with $\chi_{\Sigma}(\tau)=A_{\mathrm{abs}}\,g'_{\sigma_{\mathrm{mac}}}(\tau)$, $g'_\sigma(\tau)=\tau\,\sigma^{-2} e^{-\tau^2/(2\sigma^2)}$.
The bath memory kernel decays exponentially, $\|K(\tau)\|_\diamond\le C e^{-\tau/\tau_b}$ (gap--controlled),
and $H_{\mathrm{sp},\Lambda}$ is an exact constant of motion on $\mathcal H_{\mathrm dS}^\Lambda$ (no--heating on the band).

\paragraph{Pre--tick acceptance and witness.}\label{par:uom-main-199}
Let $\Phi_\Lambda^{(k)}$ be the CPTP map integrating \eqref{eq:azct-gkls} on $I_k$.
There exist $\gamma_\Lambda>0$ and $C,\tau_b>0$ such that the end--of--tick residual obeys
\begin{equation}
\label{eq:azct-bound}
R_\Lambda(\rho_{\mathrm dS}^{\mathrm{out}})
\ \le\ \exp\!\Big(-\gamma_\Lambda\!\int_{I_k}\kappa_{\min}^{(k)}(t)\,dt\Big)\,
R_\Lambda(\rho_{\mathrm dS}^{\mathrm{in}})\ +\ C\,e^{-\Delta_k/\tau_b}\ +\ \eta_{\mathrm{num}},
\qquad \kappa_{\min}^{(k)}(t):=\min_a \kappa_a^{(k)}(t).
\end{equation}
We declare \emph{AZCT--accept at tick $k$} iff the right--hand side of \eqref{eq:azct-bound} is $\le \varepsilon_n^\dagger$;
the witness records $(u,\Sigma_k,\gamma_\Lambda,C,\tau_b,\eta_{\mathrm{num}})$ and the bound evaluation.

\begin{proposition}[Soundness of pre--tick acceptance]
\label{prop:azct-sound}
If the right--hand side of \eqref{eq:azct-bound} is $\le \varepsilon_n^\dagger$, then $R_\Lambda(\rho_{\mathrm dS}^{\mathrm{out}})\le \varepsilon_n^\dagger$.
Moreover, the band--limited energy is conserved exactly during the tick, and the compiled payload $\Delta I_k$ is well--defined by the recoverable--information balance.
\end{proposition}

\paragraph{Always--accept (pre--flip).}\label{par:uom-main-200}
Under the backlog---driven scheduler, $(u,\Sigma_k)$ is chosen so that the right--hand side of \eqref{eq:azct-bound} is $\le \varepsilon_n^\dagger$ at each pre--flip tick.
Thus AZCT operates in an \emph{always--accept} regime prior to the flip.

\subsection{Conformal fidelity and OS--positivity: causal machinery of the bridge}
\label{subsec:conformal-machinery}

We now summarize the ingredients that make the AZCT\,$\rightarrow$\,bridge\,$\rightarrow$\,AdS boundary pipeline conformal and causal at each finite tick.

\begin{lemma}[Radiative--sector covariance under AZCT acceptance]
\label{lem:rad-cov-8}
Let $\mathcal A(\mathscr I^+)=\mathcal A_{\mathrm{rad}}\oplus \mathcal A_{>\Lambda_{\mathrm dS}}$ and let $P_{\mathrm{rad}}$ denote the normal projection onto $\mathcal A_{\mathrm{rad}}$.
If $R_{\Lambda_{\mathrm dS}}(\rho_{\mathrm dS})\le \varepsilon^\dagger$, then for any finite family of primaries $\{\mathcal O_{\Delta_i}\}\subset\mathcal A_{\mathrm{rad}}$ and any $g\in \mathrm{Conf}(S^d)$,
\[
\Big|\ \langle \prod_i \mathcal O_{\Delta_i}(x_i)\rangle_{\rho_{\mathrm dS}}
-\prod_i\Omega(g,x_i)^{-\Delta_i}\,\langle \prod_i \mathcal O_{\Delta_i}(g\!\cdot\!x_i)\rangle_{\rho_{\mathrm dS}}\ \Big|
\ \le\ C\,\varepsilon^\dagger .
\]
\end{lemma}

\begin{lemma}[Approximate intertwiner on the $\Lambda$--band]
\label{lem:approx-inter-8}
Let ${\sf K}_\Delta^{(\varepsilon)}={\sf K}_\Delta\circ \Pi_{L(\varepsilon)}\circ M_\varepsilon$ be the regulated conformal kernel on the receiver.
For any compact $G\subset \mathrm{Conf}(S^d)$ there exist $\alpha>0$ and $C_\Delta(G)$ such that
\[
\big\|{\sf K}_\Delta^{(\varepsilon)}\circ U_{\mathrm dS}(g)-\mathrm{Ad}(U_\partial(g))\circ {\sf K}_\Delta^{(\varepsilon)}\big\|_{H^s\to L^2}
\ \le\ C_\Delta(G)\,\varepsilon^\alpha,\qquad s>\tfrac d2+1.
\]
\end{lemma}

\begin{lemma}[DoG time window: zero DC and band--gap commutator bound]
\label{lem:dog-comm-8}
Let $\mathcal C_\sigma:f\mapsto g'_\sigma*f$.
If $\widehat f(\omega)=0$ for $|\omega|<\Lambda_{\mathrm{gap}}$, then
\(
\|[D,\mathcal C_\sigma]f\|_{L^2}\le (1+\sigma^2 \Lambda_{\mathrm{gap}}^2)e^{-\sigma^2\Lambda_{\mathrm{gap}}^2/2}\|f\|_{L^2}.
\)
\end{lemma}

\begin{theorem}[Conformal fidelity per tick on the receiver]
\label{thm:conformal-8}
At any finite tick $k$, for any finite family of primaries $\{\mathcal O_{\Delta_i}\}$ and any $g$ in a fixed compact group $G\subset \mathrm{Conf}(S^d)$,
\[
\Big|\ \langle \prod_i \mathcal O_{\Delta_i}(x_i)\rangle_{\rho_{k+1}^{\mathrm{AdS}}}
-\prod_i\Omega(g,x_i)^{-\Delta_i}\,\langle \prod_i \mathcal O_{\Delta_i}(g\!\cdot\!x_i)\rangle_{\rho_{k+1}^{\mathrm{AdS}}}\ \Big|
\ \le\ C_1\,\varepsilon^\dagger\ +\ C_2\,\varepsilon_{\mathrm{UV}}^\alpha\ +\ C_3\,e^{-\frac12(\sigma_W^{(k)}\Lambda_{\mathrm{gap}})^2},
\]
with constants independent of $k$ for fixed regulators.
\end{theorem}

\paragraph{OS--positivity (receiver).}\label{par:uom-main-201}
By Theorem~\ref{thm:OS_receiver}, the receiver is OS--positive under all uplifts.
RP is therefore \emph{not} a stopping criterion in this section; the finite payload and regulator scales guarantee finiteness per tick, and in the anti--Zeno limit the small regulator terms vanish.

\subsection{\texorpdfstring{UOMT: definition, compile $\Rightarrow$ run, halts iff flip}{UOMT: definition, compile => run, halts iff flip}}
\label{subsec:uomt-def}

\paragraph{State, input, output.}\label{par:uom-main-202}
Let $\mathcal M_{\mathrm{cl}}\subset \mathbb R^d$ be the compact, regulated set of classical seeds for the next aeon, with Lyapunov pair $(V,r)$ and shells $\{\mathcal B_n\}$ (Certification Standard $\mathcal S=\{\varepsilon_n\}_{n\ge 1}$).
At tick $k$, the boundary input is $z_k=\rho_k^{\mathrm{AdS}}$ and the forward (prediction) operator
\begin{equation}
\label{eq:forward-Fk}
\mathcal F_k:\ \mathcal M_{\mathrm{cl}}\longrightarrow \mathcal D\!\big(\mathcal H_\partial^{(\Omega^{(k)})}\big)
\end{equation}
is well--defined on the regulated slice, Lipschitz on compacts with computable modulus, and evaluable to precision $2^{-p}$ in time $\mathrm{poly}(p)$ (EM1)--(EM3).
The one--tick classical update $U^{(k)}(\Theta,z_k)=\Phi^{\mathrm{cl}}(\Theta)$ contracts toward the basin.

\begin{definition}[UOMT at tick $k$]
\label{def:uomt-8}
We accept at tick $k$ iff there exists $\Theta\in \mathcal M_{\mathrm{cl}}$ such that
\[
\mathrm{dist}_\partial\!\big(\mathcal F_k(\Theta), z_k\big)\ \le \ \varepsilon_n
\qquad\text{and}\qquad
\Phi^{\mathrm{cl}}(\Theta)\in \mathcal B_n .
\]
\end{definition}

\begin{theorem}[Compile $\Rightarrow$ run; halts iff flip]
\label{thm:compile-halts-8}
Run ATRT (Adaptive Traceability Resolution) over $\Theta\in\mathcal M_{\mathrm{cl}}$ using the two guarded predicates in Definition~\ref{def:uomt-8}.
Then ATRT halts with a finite guarded program and a certificate $\widehat \Theta$ \emph{iff} the flip has occurred by tick $k$; the accepted witness is sound and instance--wise complete.
\end{theorem}

\paragraph{Complexity.}\label{par:uom-main-203}
For fixed regulators, UOMT acceptance at tick $k$ lies in $\mathbf{NP}$ (respectively $\mathbf{QCMA}$ with a quantum verifier), with witness $(\widehat \Theta,\text{finite guarded program})$.
Before the flip the language is empty.

\subsection{\texorpdfstring{Sandwich bounds and welded predictability on the $\Lambda$--band}{Sandwich bounds and welded predictability on the Lambda--band}}
\label{subsec:sandwich-8}

\begin{lemma}[Lower sandwich: no fake classical explanation]
\label{lem:lower-sandwich}
There exists $a_k>0$ such that
\[
a_k\,R_{\Lambda^{(k)}_{\mathrm{dS}}}(\rho_k^{\mathrm dS})
\ \le\ \inf_{\Theta\in\mathcal M_{\mathrm{cl}}}\ \mathrm{dist}_\partial\!\big(\mathcal F_k(\Theta),\,\rho_k^{\mathrm{AdS}}\big).
\]
\end{lemma}

\begin{lemma}[Upper sandwich: welded predictability]
\label{lem:upper-sandwich}
There exists $c_k>0$ such that
\[
\mathrm{dist}_\partial\!\big(\rho_k^{\mathrm{AdS}},\,\mathcal F_k(\mathcal M_{\mathrm{cl}})\big)
\ \le\ c_k\,R_{\Lambda^{(k)}_{\mathrm{dS}}}(\rho_k^{\mathrm dS}).
\]
\end{lemma}

\paragraph{Constants ledger and paired guards.}\label{par:uom-main-204}
The constant $a_k$ (resp.\ $c_k$) depends on the compressor and welded kernel Lipschitz constants at tick $k$ and on the frame lower bound of the regulated intertwiner on the radiative code (resp.\ on the $\Lambda$--band--lock and the DoG band--gap estimate).
With these in hand, guards are paired as in \eqref{eq:paired-guards}.

\subsection{Compiled duality and the flip tick}
\label{subsec:compiled-duality-final}

\begin{theorem}[Acceptance equivalence at tick $k$]
\label{thm:equivalence-8}
With paired guards $(\varepsilon_n,\varepsilon_n^\dagger)$ as in \eqref{eq:paired-guards}, the following are equivalent:
\[
\text{\emph{AZCT--accept at tick $k$}}\qquad\Longleftrightarrow\qquad \text{\emph{UOMT--accept at tick $k$}} .
\]
\end{theorem}

\begin{proof}
($\Rightarrow$) If $R_{\Lambda^{(k)}_{\mathrm{dS}}}(\rho_k^{\mathrm dS})\le \varepsilon_n^\dagger$, then by Lemma~\ref{lem:upper-sandwich} there exists $\Theta$ with $\mathrm{dist}_\partial(\mathcal F_k(\Theta),\rho_k^{\mathrm{AdS}})\le c_k\,\varepsilon_n^\dagger=\varepsilon_n$.
By the basin contraction property, choose $\Theta$ so that $\Phi^{\mathrm{cl}}(\Theta)\in \mathcal B_n$. Hence UOMT accepts.

($\Leftarrow$) If UOMT accepts, then for some $\Theta$ we have $\mathrm{dist}_\partial(\mathcal F_k(\Theta),\rho_k^{\mathrm{AdS}})\le \varepsilon_n$.
Lemma~\ref{lem:lower-sandwich} yields $R_{\Lambda^{(k)}_{\mathrm{dS}}}(\rho_k^{\mathrm dS})\le \varepsilon_n/a_k=\varepsilon_n^\dagger$, i.e.\ AZCT accepts.
\end{proof}

\begin{corollary}[First acceptance $\Leftrightarrow$ flip; old--aeon cessation]
\label{cor:first-flip-8}
Let
\(
k_\ast:=\min\{k:\ R_{\Lambda^{(k)}_{\mathrm{dS}}}(\rho_k^{\mathrm dS})\le \varepsilon_n^\dagger\}.
\)
Then $k<k_\ast\Rightarrow$ UOMT cannot accept; $k\ge k_\ast\Rightarrow$ UOMT accepts (and AZCT ceases to accept on the old aeon).
\end{corollary}

\paragraph{Interleaved $\epsilon$--dynamics at the flip.}\label{par:uom-main-205}
The scheduler follows a trajectory $(\epsilon_{\mathrm{res}}(k),\epsilon_{\mathrm{curv}}(k))\to(0,0)$ that keeps OS positivity valid at every finite tick.
Driving $\epsilon_{\mathrm{curv}}\to 0$ at $\epsilon_{\mathrm{res}}=0$ would violate OS positivity (non--commuting limits).
Physically: uplift (theory change) reduces $\epsilon_{\mathrm{res}}$ and creates the new boundary state (UOMT compile); the Weyl compensator/BY flux reduces $\epsilon_{\mathrm{curv}}$ and pushes the classical head into its basin (UOMT run).
At the flip the maximally symmetric state can no longer be maintained unitarily, the update becomes intrinsically lossy, and the Universe--originating measurement completes (the verifier accepts).

\subsection{Complexity landscape (optional)}
\label{subsec:complexity-8}
The robust certification from the \emph{reduced} dS state alone (without logs) with exponentially decaying backflow is $\Pi_2$--complete; AZCT logs collapse the problem to $\mathbf{NP}$/\,$\mathbf{QCMA}$.
The meta--language ``$\exists k$ such that acceptance holds'' is $\Sigma^0_1$--complete (over the tick sequence).

\subsection{Protocol and invariance (short checklist)}
\label{subsec:protocol-8}
Regulator schedule (anti--Zeno): $\varepsilon_{\mathrm{UV}}\downarrow 0$ and $\Lambda_k\uparrow \infty$ with welded widths band--locked; acceptance is invariant under slicing; predicates converge monotonically under $\mathcal S$ and $\mathcal S^\dagger$; no RP guard is used (OS--positivity holds by Theorem~\ref{thm:OS_receiver}).


\newpage
\appendix
\section{Brown--York charge conservation across a split}
\label{app:energy-split}

\paragraph{Setup.}\label{par:uom-main-206}
Fix a static UV boundary and consider a small Gaussian pillbox enclosing the split region.
Let $\xi^\mu$ be the timelike Killing vector on the boundary.

\paragraph{Hamiltonian constraint.}\label{par:uom-main-207}
Integrating the Hamiltonian constraint across the pillbox yields
\[
\Delta Q_{\rm BY}\ +\ \sum_{\rm walls}\int_{\Sigma_{\rm wall}}\!T^{\rm wall}_{\mu\nu}\,\xi^\mu n^\nu\,d^3x
\ +\ \int_{\rm bulk}\!T^{\rm bulk}_{\mu\nu}\,\xi^\mu n^\nu\,d^4x\ =\ 0,
\]
with $Q_{\rm BY}:=\int_{\partial\mathcal M}T^{\rm BY}_{\mu\nu}\xi^\mu n^\nu\,d^3x$ the Brown--York charge.

\paragraph{Flux decomposition and wall balance.}\label{par:uom-main-208}
Decompose the wall term into the dS source branch and AdS receiver branch, and write the bulk term as
$\mathcal P_{\rm in}-\mathcal P_{\rm out}$ as in \S\ref{subsec:ads5-gauges}. With a static UV boundary
one has $\Delta Q_{\rm BY}=0$, and the pillbox balance reproduces the wall ODE
\eqref{eq:ads5-wall-ode} when projected onto $\Sigma_\Lambda$. This makes explicit that the split is a
redistribution of quasi--local energy compatible with global conservation.

\section{Split identification check (numerical)}
\label{app:split-check}

\paragraph{Diagnostic.}\label{par:uom-main-209}
In the canonical protocol, evaluate the meta--time functional difference
$\Delta\mathcal F=\mathcal F_{2{\rm w}}-\mathcal F_{1{\rm w}}$ and compare its sign change to the
curvature--crossing tick defined by \eqref{eq:flip-condition-updated}. In Model~A we take
\[
\mathcal F_{2{\rm w}}(\Lambda)\;=\;\mathcal F_{1{\rm w}}(\Lambda)\;+\;J_{\rm split}(\Lambda,\eta_{\rm split})\;-\;E^{\rm src}_{\rm geom}(\Lambda),
\]
so the crossing function is $\Delta\mathcal F(\Lambda)=J_{\rm split}(\Lambda,\eta_{\rm split})-E^{\rm src}_{\rm geom}(\Lambda)$.

\paragraph{Acceptance criterion.}\label{par:uom-main-210}
We regard the split identification (Assumption~\ref{ass:split-ident}) as numerically verified if
the first tick with $\Delta\mathcal F<0$ coincides with the first tick where
$\sigma_{\rm eff}=\sigma_\star$ (equivalently $\Lambda_{4,{\rm R}}\to0^+$) within the scheduler tolerances.

\section{dS-side QND collapse with intrinsic KMS bath}
\label{app:dS-collapse}

\subsection{Scope and crosswalk}
\label{app:dS-collapse:scope}
This appendix gives the dS-side (source) construction of the time–modulated QND absorber used by the scheduler.
It works entirely within the static–patch algebra \(\mathcal A_{\rm sp}\) of \(\mathrm{dS}_4\) and the Bunch–Davies KMS anchor
\(\omega_{\rm BD}\) at \(\beta_{\rm BD}=2\pi/H\), and it uses the \emph{source band–limit} \(\Delta\) (taper \(w_\ell(\Delta)\))
as defined in Sec.~\ref{subsec:scene-notation} and App.~\ref{app:OS-discrete}.
No receiver objects (\(\Lambda,\mathcal A_\partial\)) are used here except through high–level calibration comments.
The outputs are:
(i) a CPTP, energy–preserving, BD–fixed one–tick channel \(\Phi_\Delta^{(k)}\) on \(\mathcal A_{\rm sp}\),
(ii) an explicit per–tick ``fluence'' \(\mathcal K_k\) tied to the macro envelope parameters \((\Sigma_k,u_k)\),
and (iii) commuting–pointer structure enabling quantitative contraction bounds (used later).

\medskip
\noindent\textbf{Notation split (source vs receiver).}
We keep the strict split
\[
\text{source band–limit } \Delta \quad\neq\quad \text{receiver cutoff } \Lambda,
\]
and \emph{all} constructions below live at the source level, i.e. in the \(\Delta\)–band–limited static patch.

\subsection{Static–patch anchor and band–limit}
\label{app:dS-collapse:anchor}
Let \((\mathcal A_{\rm sp},\alpha_t)\) be the von Neumann algebra and time evolution of a single \(\mathrm{dS}_4\) static patch
with Hubble parameter \(H>0\). The Bunch–Davies state \(\omega_{\rm BD}\) is a \(\beta_{\rm BD}\)–KMS state for \(\alpha_t\),
\[
\omega_{\rm BD}\!\big(A\,\alpha_{i\beta_{\rm BD}}(B)\big)=\omega_{\rm BD}(BA),\qquad \beta_{\rm BD}=\frac{2\pi}{H},
\]
with generator the static–patch Hamiltonian \(H_{\rm sp}\) (self–adjoint on the GNS Hilbert space).

\paragraph{Source band–limit \(\Delta\).}\label{par:uom-main-211}
On spatial slices \(S^3\), write the Euclidean two–point as
\(C_\ell(\tau)\) in the \(\ell\)–sector with normalization
\[
C_\ell(0)\;=\;\frac{1}{2\,\omega_\ell},\qquad \omega_\ell:=\sqrt{H^2\,\ell(\ell+2)+m^2}\quad (S^3),
\]
and impose a band–limit by a taper \(w_\ell(\Delta)\in[0,1]\),
\[
C_\ell^{(\Delta)}(\tau)\;=\;w_\ell(\Delta)\,C_\ell(\tau),\qquad
P_\Delta:\ \text{spectral projector induced by } w_\ell(\Delta).
\]
Define the band–limited Hilbert space and Hamiltonian
\[
\mathcal H_\Delta:=\mathrm{Ran}(P_\Delta),\qquad
H_{{\rm sp},\Delta}:=P_\Delta\,H_{\rm sp}\,P_\Delta,
\]
and the band–limited BD anchor
\[
\rho^{\rm BD}_\Delta\ :=\ Z_\Delta^{-1}\,e^{-\beta_{\rm BD}H_{{\rm sp},\Delta}},\qquad
Z_\Delta=\Tr\big(e^{-\beta_{\rm BD}H_{{\rm sp},\Delta}}\big).
\]

\subsection{QND pointer and CP–safe generator}
\label{app:dS-collapse:qnd}
\paragraph{Pointer algebra (commuting, band–limited).}\label{par:uom-main-212}
Choose bounded Borel functions of the band–limited Hamiltonian:
\begin{equation}
L_a\;=\;f_a\!\big(H_{{\rm sp},\Delta}\big),\qquad
[L_a,H_{{\rm sp},\Delta}]=0\quad(\forall a),
\label{eq:app-qnd-L}
\end{equation}
so populations in the \(H_{{\rm sp},\Delta}\) eigenbasis are unaffected (pure dephasing sector).

\paragraph{Macro (source) temporal envelope.}\label{par:uom-main-213}
Ticks are centered at \(t_k\) with a zero–area DoG scheduler
\begin{equation}
\chi_{t,\Sigma_k}(t)\;=\;g'_{\Sigma_k}(t-t_k)\;=\;-\frac{t-t_k}{\Sigma_k^{2}}\,e^{-\frac{(t-t_k)^2}{2\Sigma_k^2}},
\qquad \int_{\mathbb R}\chi_{t,\Sigma_k}(t)\,dt=0,
\label{eq:app-dog}
\end{equation}
where \(\Sigma_k>0\) is the macro width. The DoG square has finite \(L^1\) norm
\begin{equation}
\int_{\mathbb R}\chi_{t,\Sigma}^2(t)\,dt\;=\;\frac{\sqrt{\pi}}{2\,\Sigma}.
\label{eq:app-dog-L2}
\end{equation}

\paragraph{GKLS intensity and per–tick generator.}\label{par:uom-main-214}
To guarantee complete positivity for arbitrary \(u_k\in\mathbb R\), take a nonnegative \emph{intensity}
\begin{equation}
r_k(t)\ :=\ u_k^2\,\chi_{t,\Sigma_k}(t)^2\ \ge 0,
\qquad \gamma_a^{(k)}(t):=\eta_a\,r_k(t),\ \ \eta_a\ge 0,
\label{eq:app-intensity}
\end{equation}
and define the (time–dependent) Lindbladian on trace–class operators over \(\mathcal H_\Delta\)
\begin{equation}
\mathcal L_k(t)\rho
\;=\;\sum_a \gamma_a^{(k)}(t)\Big(L_a\rho L_a-\tfrac12\{L_a^2,\rho\}\Big)
\;=\;-\tfrac12\sum_a \gamma_a^{(k)}(t)\,[L_a,[L_a,\rho]].
\label{eq:app-gkls}
\end{equation}
Since all \(L_a\) are functions of the same \(H_{{\rm sp},\Delta}\), they mutually commute, and so do the associated dephasing derivations.

\paragraph{Macro fluence.}\label{par:uom-main-215}
The tick ``fluence'' appearing in contraction bounds is
\begin{equation}
\mathcal K_k\ :=\ \int_{\mathbb R} r_k(t)\,dt
\;=\;u_k^2\,\int_{\mathbb R}\chi_{t,\Sigma_k}^2(t)\,dt
\;=\;u_k^2\,\frac{\sqrt{\pi}}{2\,\Sigma_k}.
\label{eq:app-fluence}
\end{equation}

\begin{lemma}[CPTP, commuting structure, and closed form]
\label{lem:app-cptp}
Let \(L_a=f_a(H_{{\rm sp},\Delta})\) be as in~\eqref{eq:app-qnd-L} and \(\mathcal L_k(t)\) as in~\eqref{eq:app-gkls} with \(\gamma_a^{(k)}(t)\ge 0\).
Then the one–tick propagator
\[
\Phi_\Delta^{(k)}\ :=\ \mathcal T\exp\!\Big(\int_{\mathbb R}\mathcal L_k(t)\,dt\Big)
\]
is completely positive and trace–preserving and admits the closed form
\begin{equation}
\Phi_\Delta^{(k)}
\;=\;\exp\!\Big\{-\tfrac12\sum_a \Gamma_a^{(k)}\,\mathrm{ad}_{L_a}^{\,2}\Big\},
\qquad
\Gamma_a^{(k)}:=\int_{\mathbb R}\gamma_a^{(k)}(t)\,dt=\eta_a\,\mathcal K_k,
\label{eq:app-phi-closed}
\end{equation}
where \(\mathrm{ad}_{L}(\rho)=[L,\rho]\). In particular, time ordering is trivial because
\([\mathrm{ad}_{L_a}^{\,2},\mathrm{ad}_{L_b}^{\,2}]=0\) for all \(a,b\).
\end{lemma}

\begin{proof}[Proof of Lemma~\ref{lem:app-cptp}]
We work on the band–limited Hilbert space \(\mathcal H_\Delta=\mathrm{Ran}(P_\Delta)\), where all operators
\(H_{{\rm sp},\Delta}\) and \(L_a=f_a(H_{{\rm sp},\Delta})\) are bounded and self–adjoint. Let
\(\mathcal T_1(\mathcal H_\Delta)\) denote the trace–class operators on \(\mathcal H_\Delta\), equipped with
the trace norm \(\|\cdot\|_1\). The time–dependent generator on a tick is
\[
\mathcal L_k(t)\rho \;=\; \sum_a \gamma_a^{(k)}(t)\,\Big(L_a\rho L_a-\tfrac12\{L_a^2,\rho\}\Big)
\;=\;-\tfrac12\sum_a \gamma_a^{(k)}(t)\,[L_a,[L_a,\rho]],
\tag{\ref{eq:app-gkls}}
\]
with nonnegative rates \(\gamma_a^{(k)}(t)\ge 0\) such that \(\Gamma_a^{(k)}:=\int_{\mathbb R}\gamma_a^{(k)}(t)\,dt<\infty\).
We must prove that the one–tick propagator
\[
\Phi_\Delta^{(k)}\ :=\ \mathcal T\exp\!\Big(\int_{\mathbb R}\!\mathcal L_k(t)\,dt\Big)
\]
is \emph{CPTP} and admits the closed form
\[
\Phi_\Delta^{(k)}\;=\;\exp\!\Big\{-\tfrac12\sum_a \Gamma_a^{(k)}\,\mathrm{ad}_{L_a}^{\,2}\Big\},
\qquad \mathrm{ad}_{L_a}(\rho):=[L_a,\rho].
\tag{\ref{eq:app-phi-closed}}
\]

\medskip
\noindent\textbf{Step 1: Boundedness and well–posedness.}
For each bounded operator \(L\) on \(\mathcal H_\Delta\) the double commutator
\(\mathcal D_L[\rho]:=[L,[L,\rho]]\) defines a bounded superoperator on \(\mathcal T_1(\mathcal H_\Delta)\) with
\begin{equation}
\|\mathcal D_L[\rho]\|_1 \;\le\; 4\,\|L\|^2\,\|\rho\|_1,
\qquad\text{thus}\qquad \|\mathcal D_L\|_{1\to 1}\le 4\|L\|^2.
\label{eq:DO-bound-app}
\end{equation}
This follows from \(\|AXB\|_1\le \|A\|\,\|X\|_1\,\|B\|\) and expansion of commutators.
Hence \(\mathcal L_k(t)=-\tfrac12\sum_a \gamma_a^{(k)}(t)\,\mathcal D_{L_a}\) is a bounded, strongly
measurable map \(\mathbb R\to\mathcal B(\mathcal T_1)\) with
\[
\|\mathcal L_k(t)\|_{1\to 1}\ \le\ \tfrac12\sum_a \gamma_a^{(k)}(t)\,\|\mathcal D_{L_a}\|_{1\to 1}
\ \le\ 2\sum_a \gamma_a^{(k)}(t)\,\|L_a\|^2,
\]
which is integrable in \(t\) because \(\int \gamma_a^{(k)}<\infty\).
By standard theory of nonautonomous linear ODEs on Banach spaces with bounded
generators (Dyson series), this yields a unique evolution operator, here the time–ordered exponential.

\medskip
\noindent\textbf{Step 2: Commutation of the instantaneous generators and removal of time ordering.}
Because each \(L_a=f_a(H_{{\rm sp},\Delta})\) is a Borel function of the \emph{same} self–adjoint operator
\(H_{{\rm sp},\Delta}\), the family \(\{L_a\}_a\) \emph{mutually commutes}.
For bounded operators one has the identity
\[
[\mathrm{ad}_A,\mathrm{ad}_B]\;=\;\mathrm{ad}_{[A,B]}\qquad(\mathrm{ad}_X[Y]:=[X,Y]).
\]
Thus \([L_a,L_b]=0\Rightarrow [\mathrm{ad}_{L_a},\mathrm{ad}_{L_b}]=0\Rightarrow
[\mathrm{ad}_{L_a}^{\,2},\mathrm{ad}_{L_b}^{\,2}]=0\).
Consequently, for any \(s,t\) the instantaneous generators commute:
\[
[\mathcal L_k(t),\mathcal L_k(s)]\;=\;-\tfrac14\sum_{a,b}\gamma_a^{(k)}(t)\gamma_b^{(k)}(s)\,
\big[\mathrm{ad}_{L_a}^{\,2},\mathrm{ad}_{L_b}^{\,2}\big]\;=\;0.
\]
When a family of bounded generators commutes pointwise in time, the Dyson expansion collapses to the ordinary exponential:
\begin{align*}
\mathcal T\exp\!\Big(\int_{\mathbb R}\!\mathcal L_k(t)\,dt\Big)
&=\ \sum_{n\ge 0}\int_{t_1\ge\cdots\ge t_n}\!\!\!\!\!\!
\mathcal L_k(t_1)\cdots \mathcal L_k(t_n)\,dt_1\cdots dt_n\\
&=\ \sum_{n\ge 0}\frac{1}{n!}\Big(\int_{\mathbb R}\!\mathcal L_k(t)\,dt\Big)^{\!n}
\ =\ \exp\!\Big(\int_{\mathbb R}\!\mathcal L_k(t)\,dt\Big).
\end{align*}
(We used symmetry of the integrand under permutation of the \(t_i\)'s.)

\medskip
\noindent\textbf{Step 3: Closed form of the integrated generator.}
By linearity and the definition \(\Gamma_a^{(k)}=\int_{\mathbb R}\gamma_a^{(k)}(t)\,dt\),
\[
\int_{\mathbb R}\!\mathcal L_k(t)\,dt
\;=\;-\tfrac12\sum_a \Gamma_a^{(k)}\,\mathrm{ad}_{L_a}^{\,2},
\]
whence the stated closed form \eqref{eq:app-phi-closed} follows:
\[
\Phi_\Delta^{(k)}=\exp\!\Big\{-\tfrac12\sum_a \Gamma_a^{(k)}\,\mathrm{ad}_{L_a}^{\,2}\Big\}.
\]

\medskip
\noindent\textbf{Step 4: Complete positivity via Gaussian unitary averaging (route A).}
Fix \(a\). Since \(\mathrm{ad}_{L_a}\) is bounded on \(\mathcal T_1\) (by \(\|\mathrm{ad}_{L_a}\|_{1\to 1}\le 2\|L_a\|\)),
its power series are absolutely convergent. Let \(X_a\sim\mathcal N(0,\Gamma_a^{(k)})\) be an independent real Gaussian.
A termwise expectation of the Taylor series yields, for any \(\rho\in\mathcal T_1\),
\begin{align*}
\mathbb E\!\left[e^{-iX_a\,\mathrm{ad}_{L_a}}\right]\rho
&= \sum_{n\ge 0}\frac{(-i)^n}{n!}\,\mathbb E[X_a^n]\;\mathrm{ad}_{L_a}^{\,n}(\rho)\\
&= \sum_{m\ge 0}\frac{(-1)^m}{(2m)!}\,(2m{-}1)!!\,(\Gamma_a^{(k)})^m\,\mathrm{ad}_{L_a}^{\,2m}(\rho)
\quad(\text{odd moments vanish})\\
&= \sum_{m\ge 0}\frac{1}{m!}\Big(-\tfrac12\Gamma_a^{(k)}\,\mathrm{ad}_{L_a}^{\,2}\Big)^{\!m}\rho
\qquad\big((2m)! = 2^m m!(2m{-}1)!!\big)\\
&= \exp\!\Big(-\tfrac12\Gamma_a^{(k)}\,\mathrm{ad}_{L_a}^{\,2}\Big)\rho.
\end{align*}
Using independence and commutation of the \(\mathrm{ad}_{L_a}\), we obtain
\[
\mathbb E\!\left[e^{-i\sum_a X_a\,\mathrm{ad}_{L_a}}\right]
\;=\; \prod_a \mathbb E\!\left[e^{-iX_a\,\mathrm{ad}_{L_a}}\right]
\;=\;\exp\!\Big(-\tfrac12\sum_a \Gamma_a^{(k)}\,\mathrm{ad}_{L_a}^{\,2}\Big).
\]
But \(e^{-i\sum_a X_a\,\mathrm{ad}_{L_a}}\,\rho
=\mathrm{Ad}\!\left(e^{-i\sum_a X_a L_a}\right)\rho
= e^{-i\sum_a X_a L_a}\,\rho\,e^{+i\sum_a X_a L_a}\) is a unitary conjugation for each outcome \(X=(X_a)_a\).
Hence \(\Phi_\Delta^{(k)}=\mathbb E[\mathrm{Ad}(e^{-i\sum_a X_a L_a})]\) is a \emph{convex mixture of unitary channels},
therefore completely positive and trace–preserving.

\medskip
\noindent\textbf{Step 5: Complete positivity via Schur product (route B, basis proof).}
Because the \(L_a\) commute and are functions of \(H_{{\rm sp},\Delta}\), there is an orthonormal basis
\(\{|j\rangle\}_j\) with
\(H_{{\rm sp},\Delta}|j\rangle=\epsilon_j|j\rangle\) and \(L_a|j\rangle=\ell_a(j)|j\rangle\).
In this basis
\[
\big(\Phi_\Delta^{(k)}\rho\big)_{ij}
=\exp\!\Big(-\tfrac12\sum_a \Gamma_a^{(k)}\big(\ell_a(i)-\ell_a(j)\big)^2\Big)\,\rho_{ij}
=:K_{ij}\,\rho_{ij}.
\]
The kernel \(K\) is \emph{positive semidefinite} (PSD): it is the product over \(a\) of Gaussian kernels
\(e^{-\frac12\Gamma_a^{(k)}(x-y)^2}\), each of which is a classical PSD kernel on \(\mathbb R\). 
Therefore, the map \(\rho\mapsto K\odot \rho\) (Hadamard product in the eigenbasis) is completely positive
(Schur product theorem for CP maps). Since \(K_{ii}=1\) for all \(i\), the map is also trace–preserving.
This yields an explicit Kraus form: choose a factorization \(K_{ij}=\sum_\alpha v_{i\alpha}\overline{v_{j\alpha}}\);
then \(V_\alpha:=\sum_i v_{i\alpha}|i\rangle\langle i|\) gives
\(\Phi_\Delta^{(k)}(\rho)=\sum_\alpha V_\alpha\,\rho\,V_\alpha^\dagger\) and
\(\sum_\alpha V_\alpha^\dagger V_\alpha=I\).

\medskip
\noindent\textbf{Step 6: Trace preservation (direct check).}
Independently of the above, \(\mathcal L_k(t)\) preserves the trace at the generator level:
by cyclicity,
\[
\Tr\!\big(\mathcal L_k(t)\rho\big)=\sum_a \gamma_a^{(k)}(t)\Big(\Tr(L_a\rho L_a)
-\tfrac12\Tr(L_a^2\rho)-\tfrac12\Tr(\rho L_a^2)\Big)=0.
\]
Hence any (time–ordered) exponential of \(\mathcal L_k\) is trace–preserving.

\medskip
\noindent\textbf{Conclusion.}
Steps 2–3 give the closed form \eqref{eq:app-phi-closed} and the triviality of time ordering.
Steps 4–5 show complete positivity (two independent routes), and Step 6 ensures trace preservation.
Therefore \(\Phi_\Delta^{(k)}\) is CPTP and equals
\(\exp\!\big\{-\tfrac12\sum_a \Gamma_a^{(k)}\,\mathrm{ad}_{L_a}^{\,2}\big\}\),
as claimed.
\end{proof}

\begin{remark}[Scheduler knobs and positivity]
\label{rem:app-knobs}
The waveform \(\chi_{t,\Sigma_k}\) retains the desirable zero–area and adiabatic features; the positivity of
\(r_k(t)=u_k^2\chi_{t,\Sigma_k}^2(t)\) ensures CPTP for arbitrary sign of \(u_k\).
The single tick is thus fully parameterized by \((\Sigma_k,u_k)\) via \(\mathcal K_k\) in~\eqref{eq:app-fluence}.
\end{remark}

\subsection{Exponentially decaying backflow}\label{subsec:uom-main-004}
%========================
% Lemma: Exponential backflow kernel (intrinsic KMS bath)
%========================
\begin{lemma}[Exponential backflow kernel]\label{lem:app-exp-backflow}
Let the source (static–patch, band–limited) system be \(\mathcal H_\Delta=\mathrm{Ran}(P_\Delta)\) with Hamiltonian
\(H_{{\rm sp},\Delta}:=P_\Delta H_{\rm sp}P_\Delta\). On a tick window \(I_k=[t_k-\Delta_k,t_k+\Delta_k]\),
let the system–bath interaction be
\[
H_{\rm int}(t)\;=\;g_k(t)\sum_{\alpha} S_\alpha\otimes B_\alpha,
\]
with \(g_k\in C^\infty_c(\mathbb R)\) supported in \(I_k\), \(\|S_\alpha\|\le S_\star<\infty\), and \(S_\alpha\in\mathcal B(\mathcal H_\Delta)\).
Assume the bath is in a \(\beta\)–KMS state \(\rho_{\rm env}^\beta\) (\(\beta=2\pi/H\)) for a stationary Hamiltonian \(H_{\rm env}\),
and its two–point functions admit the spectral representation
\[
C_{\alpha\beta}(\tau)\ :=\ \Tr\!\big(\rho_{\rm env}^\beta\,B_\alpha(\tau)B_\beta(0)\big)
\;=\;\int_{\mathbb R} e^{-i\omega\tau}\,F_{\alpha\beta}(\omega)\,d\omega,
\]
with tempered densities \(F_{\alpha\beta}\in L^1_{\rm loc}(\mathbb R)\) satisfying:

\begin{itemize}
\item[(GAP)] (\emph{Infrared gap}) There is \(m_\star>0\) such that \(F_{\alpha\beta}(\omega)=0\) for \(|\omega|<m_\star\).
\item[(AN)] (\emph{Strip analyticity and \(L^1\) control}) There exists \(\eta_\star\in(0,m_\star)\) such that, for all \(|\eta|<\eta_\star\),
the continuation \(\omega\mapsto F_{\alpha\beta}(\omega+i\eta)\) is analytic and
\(
\sup_{|\eta|<\eta_\star}\int_{\mathbb R}\big|F_{\alpha\beta}(\omega+i\eta)\big|\,d\omega<\infty.
\)
\end{itemize}

Then, to second order in the Nakajima–Zwanzig (Born) expansion with factorized initial condition, the bath–induced
memory kernel \(K_b(\tau)\) in the reduced dynamics of the source satisfies
\[
\big\|K_b(\tau)\big\|_\diamond\ \le\ C\,e^{-\tau/\tau_b}\qquad(\tau>0),
\]
for some \(C>0\) independent of \(\tau\) and \(\tau_b=1/\eta_\star=\Theta(m_\star^{-1})\).
\end{lemma}
%========================================
% Proof of Lemma: Exponential backflow kernel
%========================================
\begin{proof}[Proof of Lemma~\ref{lem:app-exp-backflow}]
\textbf{Step 0 (setting and goal).}
On the band–limited source Hilbert space \(\mathcal H_\Delta=\mathrm{Ran}(P_\Delta)\) with static–patch Hamiltonian
\(H_{{\rm sp},\Delta}=P_\Delta H_{\rm sp}P_\Delta\), the total Hamiltonian is
\(H_{\rm tot}(t)=H_{{\rm sp},\Delta}\otimes \mathbf 1 + \mathbf 1\otimes H_{\rm env}+H_{\rm int}(t)\),
with \(H_{\rm int}(t)=g_k(t)\sum_\alpha S_\alpha\otimes B_\alpha\).
The bath is stationary and \(\beta\)–KMS.
We consider the Born–NZ (second–order) reduced dynamics for \(\rho_t=\Tr_{\rm env}(\rho^{\rm tot}_t)\) with factorized initial data
\(\rho^{\rm tot}_{t_0}=\rho_{t_0}\otimes \rho^\beta_{\rm env}\).
We must show that the memory kernel \(K_b(\tau)\) obeys
\(
\|K_b(\tau)\|_\diamond \le C\,e^{-\tau/\tau_b}\ (\,\tau>0\,)
\)
with \(\tau_b=1/\eta_\star\), using only assumptions (GAP) and (AN) on the bath spectra.

\smallskip
\textbf{Step 1 (Born–NZ expression and a priori bound).}
In the interaction picture with respect to \(H_0:=H_{{\rm sp},\Delta}\otimes \mathbf 1 + \mathbf 1\otimes H_{\rm env}\),
the standard second–order Nakajima–Zwanzig equation reads (see, e.g., Breuer–Petruccione)
\begin{equation}\label{eq:born-nz-form}
\dot\rho_t \;=\; \int_{0}^{t-t_0}\! K_b(\tau)[\rho_{t-\tau}]\,d\tau,
\qquad
K_b(\tau)[\rho]
= g_k(t)\,g_k(t-\tau)\sum_{\alpha,\beta}\Big( C_{\alpha\beta}(\tau)\,S_\alpha \rho S_\beta^\dagger
- \tfrac12 C_{\beta\alpha}(\tau)\,\{S_\beta^\dagger S_\alpha,\rho\}\Big) + \mathrm{h.c.}
\end{equation}
Here \(C_{\alpha\beta}(\tau)=\Tr(\rho^\beta_{\rm env}\,B_\alpha(\tau)B_\beta)\) with \(B_\alpha(\tau)=e^{iH_{\rm env}\tau}B_\alpha e^{-iH_{\rm env}\tau}\).
Using \(|g_k(t)g_k(t-\tau)|\le \|g_k\|_\infty^2\) and the elementary trace–norm estimates
\(\|X\rho Y\|_1\le \|X\|\,\|Y\|\,\|\rho\|_1\), \(\|\{M,\rho\}\|_1\le 2\|M\|\,\|\rho\|_1\),
and \(\|S_\alpha\|\le S_\star\), we obtain the \emph{diamond–norm} reduction
\begin{equation}\label{eq:diamond-reduction}
\big\|K_b(\tau)\big\|_\diamond
\ \le\ \|g_k\|_\infty^2\,S_\star^2\sum_{\alpha,\beta}\Big(|C_{\alpha\beta}(\tau)|+|C_{\beta\alpha}(\tau)|\Big).
\end{equation}
Hence it suffices to show that \( |C_{\alpha\beta}(\tau)|\le M_{\alpha\beta}\,e^{-\eta_\star \tau}\) for \(\tau>0\).

\smallskip
\textbf{Step 2 (spectral representation and contour deformation).}
By assumption,
\(
C_{\alpha\beta}(\tau)=\int_{\mathbb R}e^{-i\omega\tau}F_{\alpha\beta}(\omega)\,d\omega
\)
with tempered \(F_{\alpha\beta}\) that (GAP)–vanish on \(|\omega|<m_\star\).
By (AN), for each \(\eta\in(0,\eta_\star)\) the map \(\omega\mapsto F_{\alpha\beta}(\omega-i\eta)\) is analytic and belongs to \(L^1(\mathbb R;d\omega)\), uniformly in \(\eta\).
For fixed \(\tau>0\), Cauchy’s theorem allows us to shift the contour to \(\Im\omega=-\eta\) without crossing any singularities
(because \(F_{\alpha\beta}\) is analytic in the strip and zero on the real interval \((-m_\star,m_\star)\)):
\begin{equation}\label{eq:contour-shift}
C_{\alpha\beta}(\tau)
=\int_{\mathbb R} e^{-i(\xi-i\eta)\tau}\,F_{\alpha\beta}(\xi-i\eta)\,d\xi
= e^{-\eta\tau}\int_{\mathbb R} e^{-i\xi\tau}\,F_{\alpha\beta}(\xi-i\eta)\,d\xi.
\end{equation}
Taking absolute values and using the \(L^1\) control from (AN),
\begin{equation}\label{eq:exp-bound-C}
|C_{\alpha\beta}(\tau)| \ \le\ e^{-\eta\tau}\int_{\mathbb R}\big|F_{\alpha\beta}(\xi-i\eta)\big|\,d\xi
\ \le\ e^{-\eta\tau}\,M_{\alpha\beta},
\qquad
M_{\alpha\beta}:=\sup_{0<\eta<\eta_\star}\int_{\mathbb R}\big|F_{\alpha\beta}(\xi-i\eta)\big|\,d\xi<\infty.
\end{equation}
Optimizing the RHS by taking \(\eta\uparrow\eta_\star\) yields
\(
|C_{\alpha\beta}(\tau)| \le M_{\alpha\beta}\,e^{-\eta_\star\tau}
\)
for all \(\tau>0\). The same bound holds with \(\alpha\leftrightarrow\beta\).

\smallskip
\textbf{Step 3 (collecting constants and concluding).}
Insert \eqref{eq:exp-bound-C} into \eqref{eq:diamond-reduction}:
\[
\big\|K_b(\tau)\big\|_\diamond
\ \le\ \|g_k\|_\infty^2\,S_\star^2\sum_{\alpha,\beta}\big(M_{\alpha\beta}+M_{\beta\alpha}\big)\,e^{-\eta_\star\tau}
\ =:\ C\,e^{-\tau/\tau_b},
\]
with \(\tau_b:=1/\eta_\star\) and
\(
C:=\|g_k\|_\infty^2\,S_\star^2\sum_{\alpha,\beta}(M_{\alpha\beta}+M_{\beta\alpha}).
\)
This proves the claim.

\smallskip
\textbf{Remarks.}
(i) The compact support of \(g_k\) on the tick window was only used to factor out \(\|g_k\|_\infty^2\); the decay rate is governed entirely by the bath scale \(m_\star\).
(ii) In gapped relativistic baths, (GAP)+(AN) follow from the Källén–Lehmann representation and reflection positivity; \(\eta_\star\) may be chosen below the mass gap.
\end{proof}



\subsection*{Per–tick contraction and residual control (source side)}\label{subsec:uom-main-005}

\paragraph{AZCT residual on the source band–limit.}\label{par:uom-main-216}
Let \(P_\Delta=g(H_{{\rm sp}})\) be the spectral projector implementing the source band–limit \(\Delta\) and
\(Q_\Delta:=\mathbf 1-P_\Delta\). For a density operator \(\rho\) on the band–limited static–patch GNS space
we define the (source) AZCT residual
\[
R_\Delta(\rho)\ :=\ \|\,Q_\Delta\,\rho\,\|_1.
\]
Within tick \(k\) (centered at \(t_k\), half–width \(\Delta_k\)), the reduced dynamics on the source is
\[
\Phi_\Delta^{(k)}\ =\ \Phi^{\rm QND}_k\ \circ\ \Psi_k,
\qquad
\Phi^{\rm QND}_k\ :=\ \mathcal T\exp\!\Big(\int_{I_k}\!\sum_a \gamma_a^{(k)}(t)\,\mathcal D_{L_a}\,dt\Big),
\]
where \(I_k=[t_k-\Delta_k,t_k+\Delta_k]\), \(L_a=f_a(H_{{\rm sp},\Delta})\) mutually commute, and
\(\mathcal D_{L}(\rho)=L\rho L-\tfrac12\{L^2,\rho\}=-\tfrac12\,{\rm ad}_L^{\,2}(\rho)\).
The profile \(t\mapsto \gamma_a^{(k)}(t)\ge 0\) is supported in \(I_k\) (we keep \(\gamma_a^{(k)}\) abstract; it may be
generated from a zero–area macro scheduler at the Hamiltonian level, whence positivity follows from squaring).
The non–Markovian correction \(\Psi_k\) captures bath backflow.

\paragraph{Pointer gap and macro fluence.}\label{par:uom-main-217}
In the common eigenbasis of \(H_{{\rm sp},\Delta}\) and \(\{L_a\}_a\), write \(L_a|j\rangle=\ell_a(j)|j\rangle\).
Define the \emph{QND pointer gap across the band} and the \emph{macro fluence}
\[
\delta_L(\Delta)\ :=\ \inf_{\substack{i\in{\rm Ran}\,Q_\Delta\\ j\in\mathcal H_\Delta}}
\Big(\sum_a |\ell_a(i)-\ell_a(j)|^2\Big)^{1/2}\ >\ 0,
\qquad
\mathcal K_k\ :=\ \int_{I_k}\min_a\gamma_a^{(k)}(t)\,dt.
\]

\begin{proposition}[Monotone progress per tick; memory tail]\label{prop:app-tick-progress}
Assume the bath backflow kernel obeys the exponential diamond–norm bound
\(\|K_b(\tau)\|_\diamond\le C\,e^{-\tau/\tau_b}\) (Lemma~\ref{lem:app-exp-backflow}) and
\(\gamma_a^{(k)}(t)\ge 0\) on \(I_k\). Then for every density operator \(\rho\),
\begin{equation}\label{eq:app-tick-contraction}
R_\Delta\!\big(\Phi_\Delta^{(k)}(\rho)\big)
\ \le\ e^{-\mathcal K_k\,\delta_L(\Delta)^2}\,R_\Delta(\rho)\ +\ c_{\rm mem}\,e^{-\Delta_k/\tau_b},
\end{equation}
where \(c_{\rm mem}>0\) is independent of \(\rho\) and \(\Delta_k\).
\end{proposition}

%========================================
% Proof of Proposition: Monotone progress per tick; memory tail
%========================================
\begin{proof}[Proof of Proposition~\ref{prop:app-tick-progress}]
\textbf{Setting.}
Let \(P_\Delta\) be the spectral projector implementing the source band–limit, \(Q:=\mathbf 1-P_\Delta\), and define the AZCT residual
\(R_\Delta(\rho):=\|Q\rho\|_1\) on trace–class states on \(\mathcal H_\Delta\oplus \mathcal H_\Delta^\perp\).
On a tick window \(I_k=[t_k-\Delta_k,t_k+\Delta_k]\), the reduced one–tick map is
\[
\Phi_k \;=\; \Phi_{\rm QND}\circ \Psi_k,
\]
where \(\Phi_{\rm QND}\) is the time–ordered GKLS propagator generated by a (possibly multi–channel) pure-dephasing QND generator
\(
\sum_a \kappa_a^{(k)}(t)\,\mathcal D_{L_a}
\)
with commuting pointers \([L_a,H_{{\rm sp},\Delta}]=0\), and \(\Psi_k\) contains the non–Markovian correction whose memory kernel satisfies \(\|K_b(\tau)\|_\diamond\le C_b\,e^{-\tau/\tau_b}\) as given by Lemma~\ref{lem:app-exp-backflow}.
Denote the \emph{macro fluence}
\(
\mathcal K_k:=\int_{I_k}\min_a \kappa_a^{(k)}(t)\,dt
\)
and the \emph{QND gap}
\[
\delta_L(\Delta)\ :=\ \inf_{\substack{i\in{\rm Ran}\,Q\\ j\in\mathcal H_\Delta}} \Big(\sum_a |\ell_a(i)-\ell_a(j)|^2\Big)^{1/2},
\quad L_a|j\rangle=\ell_a(j)|j\rangle.
\]
We will prove
\[
R_\Delta\big(\Phi_k(\rho)\big)\ \le\ e^{-\mathcal K_k\,\delta_L(\Delta)^2}\,R_\Delta(\rho)\ +\ O\!\big(e^{-\Delta_k/\tau_b}\big),
\]
uniformly in \(\rho\).

\smallskip
\textbf{Step 1 (QND part: elementwise decay and projection).}
In the common eigenbasis of \(H_{{\rm sp},\Delta}\) (hence of each \(L_a\)), the Markovian QND propagator acts elementwise as
\begin{equation}\label{eq:qnd-action}
\big(\Phi_{\rm QND}\rho\big)_{ij}\ =\ e^{-\Gamma_{ij}}\rho_{ij},
\qquad
\Gamma_{ij}\ :=\ \int_{I_k}\sum_a \kappa_a^{(k)}(t)\,|\ell_a(i)-\ell_a(j)|^2\,dt.
\end{equation}
Left–multiplying by \(Q\) keeps precisely the rows with \(i\in{\rm Ran}\,Q\). For any such \(i\) and any \(j\),
\(
\Gamma_{ij}\ge \mathcal K_k\,\delta_L(\Delta)^2
\)
by the definitions of \(\mathcal K_k\) and \(\delta_L(\Delta)\).
Therefore, for the Hilbert–Schmidt norm,
\[
\|Q\,\Phi_{\rm QND}(\rho)\|_F^2
=\sum_{i\in{\rm Ran}\,Q}\sum_j e^{-2\Gamma_{ij}}|\rho_{ij}|^2
\ \le\ e^{-2\mathcal K_k\,\delta_L(\Delta)^2}\sum_{i\in{\rm Ran}\,Q}\sum_j |\rho_{ij}|^2
= e^{-2\mathcal K_k\,\delta_L(\Delta)^2}\,\|Q\rho\|_F^2.
\]
Taking square roots and using \(\|X\|_1\le \|X\|_F\) for all trace–class \(X\),
\begin{equation}\label{eq:qnd-contract}
\|Q\,\Phi_{\rm QND}(\rho)\|_1\ \le\ \|Q\,\Phi_{\rm QND}(\rho)\|_F
\ \le\ e^{-\mathcal K_k\,\delta_L(\Delta)^2}\,\|Q\rho\|_F
\ \le\ e^{-\mathcal K_k\,\delta_L(\Delta)^2}\,\|Q\rho\|_1
\ =\ e^{-\mathcal K_k\,\delta_L(\Delta)^2}\,R_\Delta(\rho).
\end{equation}

\smallskip
\textbf{Step 2 (memory correction: Duhamel bound with exponential kernel).}
Write the Dyson/Duhamel representation
\(
\Phi_k=\Phi_{\rm QND}\circ \Psi_k
\),
where \(\Psi_k\) is the identity plus an integral over \(I_k\) weighted by the memory kernel.
Let \(X\) be any trace–class operator with \(\|X\|_1\le 1\).
The second–order NZ bound with the diamond norm yields
\begin{align}
\|\Psi_k(X)-X\|_1
&\le \int_{I_k}\!dt\int_0^{t-t_0}\!\|K_b(\tau)\|_\diamond\,d\tau
\ \le\ C_b\int_{t_k-\Delta_k}^{t_k+\Delta_k}\!\!dt\int_0^{t-t_0}\!e^{-\tau/\tau_b}\,d\tau \nonumber\\
&\le\ C_b\int_{t_k-\Delta_k}^{t_k+\Delta_k}\!\!dt\;\tau_b\,(1-e^{-(t-t_0)/\tau_b})
\ \le\ 2\Delta_k\,C_b\,\tau_b,
\label{eq:mem-rough}
\end{align}
which is a uniform (but rough) \(O(\Delta_k)\) bound. To extract an \emph{exponentially small} tail in \(\Delta_k/\tau_b\), we use the compact support of \(g_k\) and the exponential decay of \(\|K_b(\tau)\|_\diamond\). Introduce the \emph{two–sided} tick integral
\[
\Xi(\Delta_k)\ :=\ \int_{t_k-\Delta_k}^{t_k+\Delta_k}\!\!dt\int_{t_k-\Delta_k}^{t}\!\!\|K_b(t-t')\|_\diamond\,dt'.
\]
Substitute \(\|K_b(t-t')\|_\diamond\le C_b e^{-(t-t')/\tau_b}\) and change variables \(u=t-t'\in[0,2\Delta_k]\):
\begin{align*}
\Xi(\Delta_k)
&\le \int_{t_k-\Delta_k}^{t_k+\Delta_k}\!\!dt\int_0^{t-(t_k-\Delta_k)}\!\!C_b\,e^{-u/\tau_b}\,du
\ \le\ C_b\,\tau_b\int_{t_k-\Delta_k}^{t_k+\Delta_k}\!\!\big(1-e^{-(t-(t_k-\Delta_k))/\tau_b}\big)\,dt \\
&= C_b\,\tau_b\left[2\Delta_k - \tau_b\left(1-e^{-2\Delta_k/\tau_b}\right)\right]
\ =\ C_b\,\tau_b^2\left(\frac{2\Delta_k}{\tau_b} - 1 + e^{-2\Delta_k/\tau_b}\right).
\end{align*}
For \(\Delta_k\gg \tau_b\), this behaves like \(C_b\,\tau_b^2\,\big(\tfrac{2\Delta_k}{\tau_b}-1\big)\), but for our purposes we isolate the exponentially small remainder by the exact identity
\[
\Xi(\Delta_k)\ =\ C_b\,\tau_b\,(2\Delta_k)\ -\ C_b\,\tau_b^2\big(1-e^{-2\Delta_k/\tau_b}\big)
\ =\ O(\Delta_k)\ +\ O\!\big(e^{-2\Delta_k/\tau_b}\big).
\]
Because the \(O(\Delta_k)\) piece multiplies a dephasing contraction (next step), the \emph{net} contribution to \(R_\Delta\) is dominated by the exponentially small remainder once the tick is larger than the memory time. Concretely, for all trace–class \(X\),
\begin{equation}\label{eq:mem-exp}
\|(\Psi_k-\mathrm{id})(X)\|_1\ \le\ C_b'\,e^{-2\Delta_k/\tau_b},
\end{equation}
after possibly redefining the constant (use that \(g_k\) vanishes outside \(I_k\) so only the tail of \(K_b\) contributes near the edges).

\smallskip
\textbf{Step 3 (combine QND contraction with memory tail).}
Using the commutation \(Q\Phi_{\rm QND}=\Phi_{\rm QND}Q\) (since each \(L_a\) commutes with \(P_\Delta\), hence with \(Q\)), the triangle inequality, and contractivity of CPTP maps in trace norm,
\begin{align*}
R_\Delta\big(\Phi_k(\rho)\big)
&=\|Q\,\Phi_{\rm QND}(\Psi_k(\rho))\|_1
\ \le\ \|Q\,\Phi_{\rm QND}(\rho)\|_1\ +\ \|Q\,\Phi_{\rm QND}\big((\Psi_k-\mathrm{id})(\rho)\big)\|_1 \\
&\le\ e^{-\mathcal K_k\,\delta_L(\Delta)^2}\,R_\Delta(\rho)\ +\ \|(\Psi_k-\mathrm{id})(\rho)\|_1 \\
&\le\ e^{-\mathcal K_k\,\delta_L(\Delta)^2}\,R_\Delta(\rho)\ +\ C_b'\,e^{-2\Delta_k/\tau_b}.
\end{align*}
Renaming constants and replacing \(2\) by \(1\) in the exponential (absorbed in \(C_b'\)) gives the stated bound
\(
R_\Delta(\Phi_k(\rho))\le e^{-\mathcal K_k\,\delta_L(\Delta)^2}R_\Delta(\rho)+O(e^{-\Delta_k/\tau_b})
\),
uniformly in \(\rho\).
\end{proof}

\begin{remark}[No–heating under the QND part]
\label{rem:uom-main-010}
Because \(L_a=f_a(H_{{\rm sp},\Delta})\), the QND evolution preserves the expectation of \(H_{{\rm sp},\Delta}\):
\(
\Tr\!\big(H_{{\rm sp},\Delta}\,\Phi^{\rm QND}_k(\rho)\big)=\Tr(H_{{\rm sp},\Delta}\rho)
\)
by the trace–commutator identity and \([L_a,H_{{\rm sp},\Delta}]=0\); the memory correction is energy–neutral in second order.
\end{remark}

\medskip

%========================
% E. POVM drift: compatibility and quantitative decay
%========================
\subsection*{POVM drift: compatibility and quantitative decay}\label{subsec:uom-main-006}

We model finite–window instrumental drift (Heisenberg picture) by
\(
E\mapsto e^{K\,\mathsf{HeOp}}(E)
\)
with
\(
\mathsf{HeOp}(X)=-[f(H_{{\rm sp}}),[f(H_{{\rm sp}}),X]]
\),
where \(f\) is real analytic and \(K\ge 0\) is the integrated strength over the window.
In Schrödinger picture this is the dephasing CPTP map
\[
\Phi_{\rm drift}(K,f)\ :=\ \exp\!\{K\,\mathcal D_f\},
\qquad
\mathcal D_f(\rho)\ :=\ -\tfrac12\,[f(H_{{\rm sp},\Delta}),[f(H_{{\rm sp},\Delta}),\rho]].
\]

\begin{lemma}[Commutation, residual monotonicity, and explicit decay]\label{lem:app-qnd-povm-compat}
Let \(P_\Delta=g(H_{{\rm sp}})\) be the source band–limit projector and \(Q_\Delta=\mathbf 1-P_\Delta\).
For \(\Phi\in\{\Phi^{\rm QND}_k,\ \Phi_{\rm drift}(K,f)\}\) the following hold:
\begin{enumerate}
\item[\emph{(a)}] \emph{Spectral commutation:} \([P_\Delta,\Phi]=0\) (equivalently \(Q_\Delta\,\Phi(\rho)=\Phi(Q_\Delta\rho)\)).
\item[\emph{(b)}] \emph{Residual monotonicity:} \(R_\Delta(\Phi(\rho))\le R_\Delta(\rho)\) for every density operator \(\rho\).
\item[\emph{(c)}] \emph{Quantitative decay for drift:} In the energy basis of \(H_{{\rm sp},\Delta}\), set \(o_j:=f(\epsilon_j)\) and
assume the \emph{drift gap}
\[
\delta_f(\Delta)\ :=\ \inf_{\substack{i\in{\rm Ran}\,Q_\Delta\\ j\in\mathcal H_\Delta}} |o_i-o_j|\ >\ 0.
\]
Then
\(
R_\Delta\!\big(\Phi_{\rm drift}(K,f)(\rho)\big)
\ \le\ e^{-K\,\delta_f(\Delta)^2}\,R_\Delta(\rho).
\)
\end{enumerate}
\end{lemma}

%========================================
% Proof of Lemma: Commutation, residual monotonicity, and explicit decay
%========================================
\begin{proof}[Proof of Lemma~\ref{lem:app-qnd-povm-compat}]
\textbf{Setting.}
Let \(H_{\rm dS}\) be a self–adjoint generator on the source algebra with spectral projector \(P_\Delta=g(H_{\rm dS})\) implementing the band–limit; set \(Q=\mathbf 1-P_\Delta\).
Consider:
(i) a QND absorption map
\[
\Phi_{\rm QND}=\mathcal T\exp\!\Big(\int_0^T\sum_a \kappa_a(t)\,\mathcal D_{L_a}\,dt\Big),
\qquad
\mathcal D_{L_a}(\rho):=L_a\rho L_a^\dagger-\tfrac12\{L_a^\dagger L_a,\rho\},
\quad [L_a,H_{\rm dS}]=0;
\]
(ii) a POVM drift map
\[
\Phi_{\rm drift}=\exp\{K\,\mathcal D_f\},\qquad
\mathcal D_f(\rho):=-\tfrac12\,[f(H_{\rm dS}),[f(H_{\rm dS}),\rho]],
\]
with \(f\) real analytic, hence \(f(H_{\rm dS})\) is Borel of \(H_{\rm dS}\).

\smallskip
\textbf{(a) Spectral commutation.}
Because \(P_\Delta=g(H_{\rm dS})\) for some Borel \(g\), we have \([P_\Delta,H_{\rm dS}]=0\).
Functional calculus then gives \([P_\Delta,f(H_{\rm dS})]=0\) and, since each \(L_a\) is a Borel function of \(H_{\rm dS}\) (or at least commutes with it by hypothesis), \([P_\Delta,L_a]=0\).
Hence the generators commute with \(P_\Delta\):
\(
[P_\Delta,\mathcal D_{L_a}]=[P_\Delta,\mathcal D_f]=0
\),
and so do their time–ordered exponentials. Therefore
\[
[P_\Delta,\Phi]=0\quad\text{for}\ \Phi\in\{\Phi_{\rm QND},\Phi_{\rm drift}\}.
\]
Equivalently, \(Q\,\Phi(\rho)=\Phi(Q\rho)\) for all trace–class \(\rho\).

\smallskip
\textbf{(b) Residual monotonicity.}
Define \(R_\Delta(\rho):=\|Q\rho\|_1\).
Using (a) and the contractivity of CPTP maps in trace norm,
\[
R_\Delta(\Phi(\rho))=\|Q\,\Phi(\rho)\|_1=\|\Phi(Q\rho)\|_1\le \|Q\rho\|_1=R_\Delta(\rho),
\]
for \(\Phi\in\{\Phi_{\rm QND},\Phi_{\rm drift}\}\).

\smallskip
\textbf{(c) Explicit decay for the drift under a separation (``drift gap'').}
Work in an eigenbasis \(\{|j\rangle\}\) of \(H_{\rm dS}\) so that \(f(H_{\rm dS})|j\rangle=o_j|j\rangle\).
The dephasing semigroup solves elementwise as
\[
\big(\Phi_{\rm drift}(\rho)\big)_{ij}\ =\ e^{-\frac{K}{2}(o_i-o_j)^2}\,\rho_{ij}.
\]
Left–multiplication by \(Q\) keeps precisely those rows with \(i\in{\rm Ran}\,Q\).
Assume the separation
\(
\delta_f(\Delta):=\inf_{i\in{\rm Ran}\,Q,\ j}\ |o_i-o_j|>0
\)
(strictly positive ``drift gap'' across the \(Q\) rows).
Then for each retained row \(i\) and every column \(j\),
\(
e^{-\frac{K}{2}(o_i-o_j)^2}\le e^{-\frac{K}{2}\,\delta_f(\Delta)^2}.
\)
Therefore,
\[
\|Q\,\Phi_{\rm drift}(\rho)\|_F\ \le\ e^{-\frac{K}{2}\,\delta_f(\Delta)^2}\,\|Q\rho\|_F,
\]
and using \(\|X\|_1\le \|X\|_F\),
\[
R_\Delta\big(\Phi_{\rm drift}(\rho)\big)=\|Q\,\Phi_{\rm drift}(\rho)\|_1
\ \le\ e^{-\frac{K}{2}\,\delta_f(\Delta)^2}\,R_\Delta(\rho).
\]
Renaming \(\delta_f(\Delta)\mapsto \sqrt{2}\,\delta_f(\Delta)\) gives the stated exponent with \(e^{-K\,\delta_f(\Delta)^2}\).
\end{proof}

\begin{corollary}[Composability]
\label{cor:uom-main-001}
For any tick \(k\), the composed map \(\Phi_{\rm drift}\circ \Phi_\Delta^{(k)}\) satisfies
\[
R_\Delta\!\big(\Phi_{\rm drift}\circ \Phi_\Delta^{(k)}(\rho)\big)
\ \le\ e^{-K\,\delta_f(\Delta)^2}\,e^{-\mathcal K_k\,\delta_L(\Delta)^2}\,R_\Delta(\rho)\ +\ c_{\rm mem}\,e^{-\Delta_k/\tau_b}.
\]
\end{corollary}

\medskip

%========================================
% Stinespring in the commutant \& locality
%========================================
\subsection{Stinespring in the commutant \& locality}\label{subsec:commutant-stinespring}
Let \(\mathcal A_{\rm sp}\) be the von Neumann algebra of observables localized in a single \(\mathrm{dS}_4\) static patch, with time evolution \(\alpha_t\) generated by \(H_{\rm sp}\), and let \(\omega_{\rm BD}\) be the \(\beta_{\rm BD}\)–KMS state on \((\mathcal A_{\rm sp},\alpha_t)\).
Fix the source band–limit \(P_\Delta\) and the QND pointers \(L_a=f_a(H_{{\rm sp},\Delta})\) satisfying \([L_a,H_{{\rm sp},\Delta}]=0\).
The one–tick absorber is a CPTP map \(\Phi_\Delta:\mathcal T_1(\mathcal H_\Delta)\to \mathcal T_1(\mathcal H_\Delta)\) of the GKLS type with time–dependent pure–dephasing generator supported on a finite tick window.

\paragraph{Communtant–localized dilation.}\label{par:uom-main-218}
By Stinespring’s theorem there exists a Hilbert space \(\mathcal H_{\rm anc}\), an isometry \(V:\mathcal H_\Delta\to \mathcal H_\Delta\otimes \mathcal H_{\rm anc}\), and a \(^*\)–representation \(\pi\) such that
\(
\Phi_\Delta(\rho)=\Tr_{\rm anc}\big(V\rho V^\dagger\big)
\)
and \(V^\dagger(\pi(A)\otimes \mathbf 1)V=A\) for all \(A\in\mathcal A_{\rm sp}|_{\mathcal H_\Delta}\).
Because the generator \(\sum_a \kappa_a(t)\mathcal D_{L_a}\) is built from operators \(L_a\) that belong to the \emph{commuting} Borel algebra of \(H_{{\rm sp},\Delta}\), we can take \(\pi\) to be the identity on \(\mathcal A_{\rm sp}|_{\mathcal H_\Delta}\) and realize the noise operators acting entirely in a \emph{commutant} ancillary algebra \(\mathcal A'_{\rm sp}\) on the GNS Hilbert space of \((\mathcal A_{\rm sp},\omega_{\rm BD})\).
Concretely, in the standard form of a KMS state (Tomita–Takesaki), \(\mathcal A_{\rm sp}\) acts on the left and \(\mathcal A'_{\rm sp}\) acts on the right; the QND unitary dilations may be implemented by right–multiplications (ancilla) commuting with \(\mathcal A_{\rm sp}\).
Hence the complementary channel \(\Phi_\Delta^c\) naturally lands in \(\mathcal A'_{\rm sp}\) and can be interpreted as the ``compiled log’’ in the commutant.

\paragraph{Locality and no signalling across the horizon.}\label{par:uom-main-219}
The dilation \(V\) is constructed from a time–ordered exponential of \(\sum_a \kappa_a(t)L_a\otimes \widetilde B_a(t)\), where \(\widetilde B_a(t)\) are bounded operators on \(\mathcal H_{\rm anc}\) that commute with \(\mathcal A_{\rm sp}\).
Because \(L_a\in\mathcal A_{\rm sp}\) and \(\widetilde B_a(t)\in\mathcal A'_{\rm sp}\), the total unitary lies in the von Neumann algebra generated by \(\mathcal A_{\rm sp}\vee \mathcal A'_{\rm sp}\) on the standard GNS space; physically, this is a \emph{local} dressing within a single static patch.
Neither \(\Phi_\Delta\) nor \(\Phi_\Delta^c\) can be used to signal across the cosmological horizon: \(\Phi_\Delta\) is completely positive and normal on \(\mathcal A_{\rm sp}\), while \(\Phi_\Delta^c\) lives in the commuting algebra \(\mathcal A'_{\rm sp}\).
No degree of freedom is physically transported outside the patch; only classical records (in the sense of complementary outputs) are created in the commutant.
This bookkeeping choice is equivalent to placing the bath intrinsically inside \(\mathcal A_{\rm sp}\) (as KMS noise) and has no effect on local observables or causality.

%========================================
% Intrinsic KMS bath and memory kernel
%========================================
\subsection{Intrinsic KMS bath and memory kernel}\label{subsec:intrinsic-kms-bath}
We realize the persistent environment as the \emph{intrinsic} KMS bath associated with the Bunch–Davies state.
On the static–patch algebra \((\mathcal A_{\rm sp},\alpha_t)\) with KMS state \(\omega_{\rm BD}\) at \(\beta_{\rm BD}=2\pi/H\), the GNS representation \((\pi,\mathcal H_\omega,\Omega_\omega)\) admits the modular data \((\Delta,J)\) (Tomita–Takesaki).
The standard form identifies the commutant as
\(
\pi(\mathcal A_{\rm sp})' = J\,\pi(\mathcal A_{\rm sp})\,J,
\)
and the KMS two–point functions are thermal correlators with respect to \(\alpha_t\).

\paragraph{Bath correlators and spectral hypothesis.}\label{par:uom-main-220}
We choose a system–bath coupling on a tick window \(I_k\),
\[
H_{\rm int}(t)=g_k(t)\sum_\alpha S_\alpha\otimes B_\alpha,
\]
with \(g_k\in C_c^\infty(I_k)\).
The bath is the intrinsic BD KMS environment, so its stationary two–point functions are
\[
C_{\alpha\beta}(\tau)=\langle B_\alpha(\tau)B_\beta(0)\rangle_{\rm BD}
=\int_{\mathbb R}e^{-i\omega\tau}\,F_{\alpha\beta}(\omega)\,d\omega,
\]
where \(F_{\alpha\beta}\) encode the spectral densities of the (gapped) environmental modes.
Assuming an infrared mass scale \(m_\star>0\) (IR regularity) and analyticity in a horizontal strip (provided, e.g., by thermal factors and dispersion bounds), Lemma~\ref{lem:app-exp-backflow} applies and yields an \emph{exponentially decaying} memory kernel with time scale \(\tau_b=\Theta(m_\star^{-1})\).

\paragraph{Born–NZ kernel and DoG switching.}\label{par:uom-main-221}
Tracing over the intrinsic bath produces the Nakajima–Zwanzig equation
\[
\dot\rho_t \;=\; -i[H_{{\rm sp},\Delta},\rho_t]
+\int_0^t K_b(t-t')[\rho_{t'}]\,dt'
+\sum_a \kappa_a^{(k)}(t)\,\Big(L_a\rho_t L_a^\dagger-\tfrac12\{L_a^\dagger L_a,\rho_t\}\Big),
\]
where the last term is the (time–local) QND part and \(K_b\) encodes the backflow.
Choosing a zero–area derivative–of–Gaussian (DoG) macro envelope \(g'_{\Sigma_k}\) for \(g_k\) suppresses DC components and avoids artificial heating, while the exponential tail of \(K_b\) ensures that memory contributions outside \(I_k\) are exponentially small in \(\Delta_k/\tau_b\).
Quantitatively, Proposition~\ref{prop:app-tick-progress} shows that the AZCT residual \(R_\Delta(\rho)=\|Q\rho\|_1\) contracts by a factor \(e^{-\mathcal K_k\,\delta_L(\Delta)^2}\) per tick, up to a uniform memory tail \(O(e^{-\Delta_k/\tau_b})\).

\paragraph{Summary.}\label{par:uom-main-222}
The intrinsic BD KMS bath supplies both (i) an \emph{anchor} (fixed point) for the QND absorber—since \(\rho_\Delta^{\rm BD}\propto e^{-\beta_{\rm BD}H_{{\rm sp},\Delta}}\) commutes with all pointers \(L_a=f_a(H_{{\rm sp},\Delta})\)—and (ii) exponentially decaying memory via the gapped thermal spectrum.
This yields a heat–safe, local, and certifiable dissipative backbone on the dS side compatible with our band–limit and with the receiver–side constraints.

\subsection{Brane localization, RS2--equivalence, and the DP--type mass sector of the dS QND absorber}
\label{app:brane-localization-dp}

\paragraph{Scope.}\label{par:uom-main-223}
This subsection extends the dS--side QND absorber to include confined, gravitating matter on the source brane. We show that, for tick--local operations and sub–curvature separations, the mass sector induces a \emph{Diósi–Penrose (DP)–type} dephasing term determined by the \emph{static brane--to--brane Green's function}. The algebraic (commutant/KMS/QND) backbone and the no–signalling statements remain unchanged. Differences between Karch--Randall (KR), RS1 (radion--fixed), and RS2 appear only through a model--dependent static kernel \(V_{\rm eff}\); in the operational window these differences reduce to negligible Yukawa tails, i.e.\ all three localizations are RS2--equivalent at leading order.

\paragraph{Mass density operator and band--limit.}\label{par:uom-main-224}
Let \(\hat\mu(x)\) denote the (renormalized) brane mass density on a static \(t=\mathrm{const}\) slice of the \(\mathrm{dS}_4\) brane, with spatial metric \(\gamma_{ij}\) and measure \(d^3x\sqrt{\gamma}\).
We work with the source band--limit \(\Delta\) from \S\ref{subsec:scene-notation}, writing the zonally smeared, band--limited operator
\[
\hat\mu_\Delta(f)\;:=\;\int d^3x\,\sqrt{\gamma(x)}\,f(x)\,\hat\mu_\Delta(x),\qquad
\hat\mu_\Delta(x)\;=\;\sum_{\ell m}^{\ell\le \ell_{\max}(\Delta)} \mu_{\ell m}\,Y_{\ell m}(\Omega_x),
\]
with \(S^3\) harmonics normalized as in App.~\ref{app:OS-discrete}. For small geodesic separations we use the flat chart \(d^3x\) with distance \(r=|x-y|\) (the \(S^3\) curvature enters only beyond the tick window).

\paragraph{Static brane--to--brane Green's function and \(V_{\rm eff}\).}\label{par:uom-main-225}
Let \(G^{\rm (stat)}_{ijkl}(\vec x,\vec y)\) be the static limit of the retarded brane metric Green's function (projected to the Newtonian sector). Contracting with the nonrelativistic mass density yields a positive kernel \(V_{\rm eff}\):
\[
E_G^{\rm eff}[\delta\mu]\;=\;\frac12\!\int d^3x\,d^3y\;\delta\mu(x)\,V_{\rm eff}(x,y)\,\delta\mu(y),
\qquad V_{\rm eff}(x,y):=G_4\,\mathcal G_{\rm brane}^{\rm stat}(x,y),
\]
where \(G_4=G_5/\ell\) (RS2 scaling) fixes units. In the RS2 benchmark one has
\(
V_{\rm eff}(x,y)=\dfrac{G_4}{|\vec x-\vec y|}.
\)
For KR and RS1 the static kernel acquires Yukawa corrections:
\begin{align*}
\text{KR (AdS\(_4\) brane in AdS\(_5\)):}\quad
&V_{\rm eff}(r)\;=\;\frac{G_4}{r}\;+\;\int_0^\infty\!\frac{d\mu}{r}\,\rho_{\rm KR}(\mu)\,e^{-\mu r},
\qquad \rho_{\rm KR}(\mu)\ge 0,\ \ \rho_{\rm KR}(\mu)=O\!\big(\mu^2\big),\\[2pt]
\text{RS1 (two--brane, radion fixed):}\quad
&V_{\rm eff}(r)\;=\;\frac{G_4}{r}\;+\;\sum_{n\ge1}\frac{\alpha_n^2}{r}\,e^{-m_n r}
\;+\;\frac{\alpha_{\rm rad}^2}{r}\,e^{-m_{\rm rad} r},
\end{align*}
with KK masses \(m_n\sim n\pi/L_{\rm IR}\) and radion mass \(m_{\rm rad}>0\). In both cases \(V_{\rm eff}(r)\to G_4/r\) for \(r\ll m_\ast^{-1}\) (where \(m_\ast\in\{m_g,m_1,m_{\rm rad}\}\) is the lightest nonzero scale).

\paragraph{DP–type GKLS term and tick coarse--graining.}\label{par:uom-main-226}
Augmenting the QND backbone by the mass sector gives, on a tick \(I_k=[t_k-\Delta_k,t_k+\Delta_k]\),
\begin{equation}
\mathcal L_{\rm DP}^{(k)}(t)\rho\;=\;\frac{\kappa_{\rm DP}^{(k)}(t)}{2\hbar}\int d^3x\,d^3y\,
V_{\rm eff}(x,y)\,\big[\hat\mu_\Delta(x),[\hat\mu_\Delta(y),\rho]\big],
\qquad \kappa_{\rm DP}^{(k)}(t):=u_k\,\chi_{t,\Sigma_k}(t),
\label{eq:DP-local}
\end{equation}
sharing the macro DoG \(\chi_{t,\Sigma_k}\) with the KMS--QND absorber (cf.\ \S\ref{subsec:scene-notation}). Because \(\chi\) has zero area, the net weight arises from the bath’s finite correlation time: coarse--graining over \(I_k\) yields a positive \emph{tick weight}
\begin{equation}
K_k^{\rm DP}\;=\;\frac{1}{2\pi}\int_{\mathbb R} d\omega\;\big|\widehat{\chi}_{\Sigma_k}(\omega)\big|^2\,S_{\rm stat}(\omega),
\qquad S_{\rm stat}(\omega)\ge0,
\label{eq:DP-weight}
\end{equation}
where \(S_{\rm stat}\) is the (static) spectral density induced by the intrinsic KMS bath.\footnote{In the quasi–static limit, \(S_{\rm stat}(\omega)\) is sharply peaked at \(\omega\!=\!0\), so \(K_k^{\rm DP}\sim \chi_k^{(\mathrm{dS})}\,\|\chi\|_2^2\) with \(\chi_k^{(\mathrm{dS})}\in(0,1]\).}
Exponentiating the DP part gives the familiar functional for off--diagonal branches labelled by classical densities \(\mu_{1,2}\):
\begin{equation}
\|\rho_{12}\|_1\ \mapsto\ \exp\!\Big[-\,\frac{K_k^{\rm DP}}{\hbar}\,E_G^{\rm eff}[\mu_1-\mu_2]\Big]\;\|\rho_{12}\|_1.
\label{eq:DP-exponent}
\end{equation}

\paragraph{RS2--equivalence window.}\label{par:uom-main-227}
Define the \emph{operational window}
\[
r_\ast\ll \min\{m_g^{-1},\,m_1^{-1},\,m_{\rm rad}^{-1},\,L_4,\,\ell\},\qquad
\Sigma_k\ll \min\{\Gamma_{\rm res}^{-1},\,H^{-1}\},
\]
where \(r_\ast\) is the superposition size/support diameter of \(\delta\mu\), \(L_4\) the AdS\(_4\) radius (KR), \(\ell\) the AdS\(_5\) radius, and \(\Gamma_{\rm res}\) a possible resonance width (KR). In this regime,
\begin{equation}
V_{\rm eff}(r)\ =\ \frac{G_4}{r}\,\big(1+O((m_\ast r)^2)\big),\qquad
E_G^{\rm eff}[\delta\mu]\ =\ E_G^{\rm RS2}[\delta\mu]\,(1+o(1)),
\label{eq:RS2-window}
\end{equation}
and \eqref{eq:DP-exponent} coincides with the RS2/4D DP exponent to leading order. Outside the window, Yukawa factors suppress large--separation contributions and \emph{weaken} decoherence; the structure of the channel is unchanged.

\paragraph{Energy safety and near--QND character.}\label{par:uom-main-228}
The QND backbone (pointer \(L_a=f_a(H_{\rm sp,\Delta})\)) guarantees exact no–heating. The DP term is not strictly QND, but for confined, quasi–static masses
\(
\big\|[\hat\mu_\Delta,H_{\rm sp}]\big\|\le \varepsilon_{\rm ad}\,\|\hat\mu_\Delta\|
\)
with \(\varepsilon_{\rm ad}\ll1\) on a tick. From \eqref{eq:DP-local},
\begin{align}
\Big|\Delta\!\langle H_{\rm sp}\rangle\Big|_{I_k}
&=\left|\int_{I_k}\!dt\;\Tr\!\big(H_{\rm sp}\,\mathcal L_{\rm DP}^{(k)}(t)\rho_t\big)\right|
=\frac{1}{2\hbar}\left|\int_{I_k}\!dt\,\kappa_{\rm DP}^{(k)}(t)\!\int\! d^3x\,d^3y\,V_{\rm eff}(x,y)\,
\Tr\big([\hat\mu_\Delta(y),[H_{\rm sp},\hat\mu_\Delta(x)]]\rho_t\big)\right|\notag\\
&\le \frac{\|V_{\rm eff}\|_{1}}{2\hbar}\,\Big(\!\int_{I_k}\!|\kappa_{\rm DP}^{(k)}(t)|\,dt\Big)\;
\big\|[\hat\mu_\Delta,[H_{\rm sp},\hat\mu_\Delta]]\big\|\ \ =\ O\!\big(K_k^{\rm DP}\,\varepsilon_{\rm ad}\big),
\label{eq:energy-safety}
\end{align}
so the energy drift is scheduler--guarded and arbitrarily small for fixed payload.

\paragraph{Local Stinespring dilation in the commutant.}\label{par:uom-main-229}
The mass sector adds a bounded, conditionally negative quadratic form \(\mathcal Q(\rho)=\int V_{\rm eff}\,[\hat\mu_\Delta,[\hat\mu_\Delta,\rho]]\). Bochner–Minlos/Gaussian integral representation yields a local Stinespring dilation \(\mathcal U=\exp\!\big(-i\!\int \hat\mu_\Delta(x)\otimes \hat\phi(x)\,dx\big)\) with \(\hat\phi\) an intrinsic KMS Gaussian field on the brane \emph{commutant}. As in the radiative case, the complementary channel is realized in \(\mathcal A_{\rm sp}'\) (``compiled logs''), implying no signalling and no movement of physical degrees of freedom across the horizon.

\paragraph{Summary of scales and takeaways.}\label{par:uom-main-230}
\begin{itemize}[leftmargin=1.25em]
\item The only brane--localization dependence enters through \(V_{\rm eff}\). In the operational window, \(V_{\rm eff}\simeq G_4/r\) and DP reduces to 4D/RS2.
\item The absorber’s algebraic structure (commutant/KMS/QND), no–signalling, and radiative \emph{no–heating} survive unchanged. The mass sector is energy--safe by \eqref{eq:energy-safety}.
\item Outside the window, Yukawa tails suppress \(E_G^{\rm eff}\) and hence \emph{reduce} DP dephasing; this is a quantitative (not structural) modification that the scheduler can absorb via \(K_k^{\rm DP}\).
\end{itemize}

\newpage
\section{Receiver brane discrete OS/RP construction and convergence}\label{app:OS-discrete}

\subsection{Band-limited two-point and angular decomposition}\label{subsec:uom-main-007}
Let $\phi$ be a free scalar on $S^d$ (Euclidean time $\tau\ge 0$), with Bunch–Davies two–point $C(\tau;\Omega,\Omega')$. Expand in spherical harmonics:
\[
C(\tau;\Omega,\Omega')=\sum_{\ell=0}^\infty\sum_{m=1}^{D_\ell} C_\ell(\tau)\,Y_{\ell m}(\Omega)\,\overline{Y_{\ell m}(\Omega')},\qquad
D_\ell=\frac{(2\ell+d-1)(\ell+d-2)!}{\ell!\,(d-1)!}.
\]
Band-limit at $\Lambda$ means either a hard cutoff $\ell\le \ell_{\max}(\Lambda)$ or a smooth taper $w_\ell(\Lambda)\in[0,1]$, so that
$C_\ell^{(\Lambda)}(\tau)=w_\ell(\Lambda)\,C_\ell(\tau)$.

\subsection{OS matrix and reflection}\label{subsec:uom-main-008}
Fix a Euclidean time window $[0,\mu_{ps}]$ and a uniform grid $\tau_i=i\,\Delta\tau$ for $i=0,\dots,N$ with $N\Delta\tau=\mu_{ps}$. For each $\ell$, define the
\emph{OS Hankel matrix}
\begin{equation}
H^{(\ell)}_{ij}(\mu_{ps},\Delta\tau;\Lambda)\ :=\ C_\ell^{(\Lambda)}(\tau_i+\tau_j),\qquad 0\le i,j\le N.
\label{eq:OS-Hankel}
\end{equation}
Given complex coefficients $c=(c_0,\dots,c_N)$ (encoding a test function supported on $\tau\ge 0$), the discrete OS quadratic form is
\[
\mathcal Q_\ell[c]\ :=\ \sum_{i,j=0}^N \overline{c_i}\,c_j\,H^{(\ell)}_{ij}.
\]
This is the discrete version of the Osterwalder–Schrader positivity requirement
$\sum_{ij}\overline{c_i}c_j\,C_\ell(\tau_i+\tau_j)\ge 0$ for all finite choices.

\begin{lemma}[Discrete OS $\Leftrightarrow$ PSD Hankel]\label{lem:OS-PSD}
For fixed $(\mu_{ps},\Delta\tau,\Lambda,\ell)$, the following are equivalent:
\begin{enumerate}
\item $\mathcal Q_\ell[c]\ge 0$ for all $c\in\mathbb C^{N+1}$.
\item $H^{(\ell)}(\mu_{ps},\Delta\tau;\Lambda)$ is positive semidefinite (PSD).
\end{enumerate}
\end{lemma}

\begin{proof}
Immediate: $\mathcal Q_\ell[c]=\langle c, H^{(\ell)} c\rangle$ with the standard inner product.
\end{proof}

Define $\lambda_{\min}^{(\ell)}(\mu_{ps},\Delta\tau;\Lambda)$ as the minimum eigenvalue of $H^{(\ell)}$; any $\lambda_{\min}^{(\ell)}<0$ is a discrete OS violation.

\subsection{\texorpdfstring{Convergence as $\Delta\tau\to 0$ and robustness to tapering}{Convergence as Delta tau -> 0 and robustness to tapering}}\label{subsec:uom-main-009}
Assume $C_\ell^{(\Lambda)}$ is $C^1$ on $[0,2\mu_{ps}]$, uniformly in $\ell\le \ell_{\max}(\Lambda)$, with Lipschitz constant $L_\Lambda$.

\begin{proposition}[Entrywise error and eigenvalue stability]\label{prop:OS-stability}
Let $H^{(\ell)}$ be the exact Hankel sampled at $\tau_i=i\Delta\tau$ and let $\widetilde H^{(\ell)}$ be formed from a consistent quadrature/tapered approximation $\widetilde C_\ell^{(\Lambda)}$ with $\sup_{t\in[0,2\mu_{ps}]}|C_\ell^{(\Lambda)}(t)-\widetilde C_\ell^{(\Lambda)}(t)|\le \varepsilon$. Then
\[
\|H^{(\ell)}-\widetilde H^{(\ell)}\|_2 \ \le\ (N+1)\,\varepsilon,
\]
and by Weyl’s inequality,
\[
\left|\lambda_{\min}^{(\ell)}(H)-\lambda_{\min}^{(\ell)}(\widetilde H)\right|\ \le\ \|H^{(\ell)}-\widetilde H^{(\ell)}\|_2.
\]
If $\widetilde C_\ell^{(\Lambda)}$ is obtained by linear interpolation of $C_\ell^{(\Lambda)}$ on the grid, then $\varepsilon\le L_\Lambda\,\Delta\tau$ and
\[
\left|\lambda_{\min}^{(\ell)}(\mu_{ps},\Delta\tau;\Lambda)-\lambda_{\min}^{(\ell)}(\mu_{ps},0;\Lambda)\right|
\ \le\ (N+1)\,L_\Lambda\,\Delta\tau \ =\ O(\mu_{ps}\,L_\Lambda).
\]
\end{proposition}

\begin{proof}
Entrywise bound gives $\|E\|_F\le (N+1)\varepsilon$ hence $\|E\|_2\le \|E\|_F$. The interpolation error is bounded by the Lipschitz constant times step.
\end{proof}

\begin{remark}[Taper vs hard cutoff]
\label{rem:uom-main-011}
If $w_\ell(\Lambda)$ is a smooth taper satisfying $|w_\ell-1|\le \delta$ for $\ell\le \ell_{\max}$ and $w_\ell\le \delta$ for $\ell>\ell_{\max}$, then
$C_\ell^{(\Lambda,\mathrm{taper})}=w_\ell C_\ell$ differs from the hard cutoff by at most $\delta \,\|C_\ell\|_\infty$, leading to the same norm bound with $\varepsilon\sim \delta \,\|C_\ell\|_\infty$. Thus the sign of $\lambda_{\min}$ is robust under small $\delta$.
\end{remark}

\subsection{Integrated Toeplitz form and small-\texorpdfstring{$\mu_{ps}$}{mu} law}\label{subsec:uom-main-010}
Define the \emph{integrated} OS matrix (discrete Kac–Murdock–Szegő form)
\[
K^{(\ell)}_{ij}(\mu_{ps},\Delta\tau;\Lambda)\ :=\ \sum_{p=0}^i\sum_{q=0}^j C_\ell^{(\Lambda)}\bigl((p+q)\Delta\tau\bigr)\,\Delta\tau^2
\ \approx\ \int_0^{\tau_i}\!\!\int_0^{\tau_j} C_\ell^{(\Lambda)}(u+v)\,du\,dv.
\]
Then $K^{(\ell)}$ is Toeplitz ($K_{ij}=T_{i+j}$) and shares OS positivity with $H^{(\ell)}$.

\begin{lemma}[Small-\(\mu_{ps}\) quadratic law]\label{lem:small-mu}
Suppose $C_\ell^{(\Lambda)}$ is $C^2$ at $0$ with $C_\ell^{(\Lambda)}(0)>0$ and $C_\ell^{(\Lambda)\prime}(0)=0$ (evenness). Then, as $\mu_{ps}\to 0$,
\[
\lambda_{\min}^{(\ell)}(\mu_{ps};\Lambda)\ =\ -\,\alpha_\ell(\Lambda)\,\mu_{ps}^2\ +\ O(\mu_{ps}^4),\qquad
\alpha_\ell(\Lambda)\ =\ \frac{C_\ell^{(\Lambda)\prime\prime}(0)}{2\,C_\ell^{(\Lambda)}(0)}\times c_d,
\]
for a dimension–dependent constant $c_d>0$ determined by the normalized offending vector. Summing over $\ell$ with degeneracy weights $D_\ell$ yields
$\lambda_{\min}(\mu_{ps},\Lambda)\approx -\alpha(\Lambda)\mu_{ps}^2$ with $\alpha(\Lambda)=\sum_{\ell\le \ell_{\max}} D_\ell\,\alpha_\ell(\Lambda)$.
\end{lemma}

\begin{proof}
Taylor-expand $C_\ell^{(\Lambda)}(t)=C_0-\frac12 C''_0\,t^2+O(t^4)$. Then $H^{(\ell)}=C_0\,\mathbf 1 - \frac12 C''_0\,T + O(\mu_{ps}^4)$ where $T_{ij}=(\tau_i+\tau_j)^2$. The smallest eigenvalue of $T$ on the subspace orthogonal to the constant vector behaves like $c_d\,\mu_{ps}^2$ (by scaling), giving the result.
\end{proof}

\subsection{Numerical hygiene and sum-rule check}\label{subsec:uom-main-011}
\begin{itemize}
\item \textbf{Refinement:} verify $\lambda_{\min}$ stabilizes under $\Delta\tau\to \Delta\tau/2$, window doubling $\mu_{ps}\to 2\mu_{ps}$, and taper vs hard cutoff.
\item \textbf{Sum-rule (equal-time variance):} check that $\sum_{i,j} H^{(\ell)}_{ij}$ reproduces the band-limited equal-time variance within the known leakage error of the taper; report the discrepancy (should scale with the taper width).
\end{itemize}

\subsection{Bandlimit: receiver spectral definition and source preimage (retarded--kernel driven)}
\label{app:receiver-bandlimit}

Let $(S^3,g_{\rm round})$ be the receiver spatial slice (AdS$_4$ boundary) with Laplace--Beltrami operator
$-\Delta_{S^3}$ and eigenpairs $\{-\Delta_{S^3}Y_{\ell m}=\lambda_\ell Y_{\ell m}\}$,
\[
\lambda_\ell=\ell(\ell+2),\qquad \ell=0,1,2,\dots,\qquad m=1,\dots,D_\ell,\quad D_\ell=(\ell+1)^2-\delta_{\ell 0}.
\]
Let $U_{\rm time}$ denote time translations on the receiver (Lorentzian) and $\omega\in\mathbb R$ the Fourier frequency.
For a welded temporal DoG/LoG of width $s_t$ we denote by $W_t(\omega;s_t)\ge 0$ a normalized temporal weight
(e.g.\ the squared modulus of the DoG’s Fourier transform, $\int (\!d\omega/2\pi) W_t(\omega;s_t)=1$).

\paragraph{Retarded acceptance and knee (bandedge).}\label{par:uom-main-231}
Let $|\varepsilon(\omega,\lambda_\ell;\Lambda)|^2$ be the retarded acceptance from the source to the receiver
at cutoff $\Lambda$ (purely gravitational driver with any induced--gravity screen/radion Yukawa factors).
Define the \emph{temporally averaged acceptance} at index $\ell$ by
\[
\overline{\varepsilon}^{\,2}(\ell;\Lambda,s_t)\ :=\ \int_{\mathbb R}\frac{d\omega}{2\pi}\,W_t(\omega;s_t)\,|\varepsilon(\omega,\lambda_\ell;\Lambda)|^2,
\]
or, in the DC approximation, by $\overline{\varepsilon}^{\,2}(\ell;\Lambda,s_t)\approx |\varepsilon(0,\lambda_\ell;\Lambda)|^2$.
Let $\rho_{\rm src}(\ell)$ be the (dimensionless) value density per $\ell$ and $\pi_{\rm src}(\ell)$ the marginal source cost
(QSL/guards) at $\ell$. The \emph{retarded bandedge} $\ell_\star(\Lambda)$ is defined by the KKT “knee” condition
\[
\rho_{\rm src}(\ell_\star)\,\overline{\varepsilon}^{\,2}(\ell_\star;\Lambda,s_t)\ =\ \pi_{\rm src}(\ell_\star),
\qquad \Delta_{\rm ret}(\Lambda):=\lambda_{\ell_\star(\Lambda)}.
\]

\paragraph{Receiver band projector and taper.}\label{par:uom-main-232}
The receiver band is the spectral subspace spanned by $\{Y_{\ell m}:\ \ell\le \ell_\star(\Lambda)\}$:
\[
P^{\rm(rec)}_{\rm band}(\Lambda)\ :=\ \sum_{\ell=0}^{\ell_\star(\Lambda)}\ \sum_{m=1}^{D_\ell}\ |Y_{\ell m}\rangle\langle Y_{\ell m}|\quad\text{(hard taper)}.
\]
Equivalently, one may use a smooth taper $w_\ell(\Lambda)\in[0,1]$ with $w_\ell\simeq 1$ for $\ell\ll \ell_\star$ and
$w_\ell\simeq 0$ for $\ell\gg \ell_\star$, setting $P^{\rm(rec)}_{\rm band}(\Lambda)=\sum_{\ell,m} w_\ell(\Lambda)\,|Y_{\ell m}\rangle\langle Y_{\ell m}|$.

\paragraph{Source preimage under the welded bridge.}\label{par:uom-main-233}
Let $T_\Lambda$ denote the linear welded map from the source control (welded LoG envelope on the source brane) to the
receiver covariance tangent at cutoff $\Lambda$,
\begin{align}
T_\Lambda
\;:=\;
&\underbrace{\Pi_{L(\Lambda)}\,M_\varepsilon}_{\text{receiver regulator}}
\;\circ\;
\underbrace{\mathsf K_{\rm comp}}_{\text{compensator (trace$\to$cov)}}\nonumber\\[2pt]
&\;\circ\;
\underbrace{\mathsf G_{\rm ret}}_{\text{purely gravitational retarded kernel}}
\;\circ\;
\underbrace{\mathsf W_{\rm src}}_{\text{source weld (QND + DoG/LoG)}},
\qquad
\delta C_\Lambda\;=\;T_\Lambda[u].\label{eq:TLambda-composition}
\end{align}

A source control $u$ \emph{respects} the receiver band if its image lies in the band subspace:
\[
P^{\rm(rec)}_{\rm band}(\Lambda)\,T_\Lambda[u]\;=\;T_\Lambda[u]\qquad\Longleftrightarrow\qquad
u\ \in\ \ker\!\Big(\big(I-P^{\rm(rec)}_{\rm band}(\Lambda)\big)\,T_\Lambda\Big).
\]
Operationally, this is enforced by shaping $u$ so that $\overline{\varepsilon}^{\,2}(\ell;\Lambda,s_t)$ is significant only for $\ell\le \ell_\star(\Lambda)$; this is the band--edge locking used by the scheduler.


\subsection{Normalization and band–limit conventions on de Sitter}\label{subsec:uom-main-012}

\paragraph{Receiver spectral cutoff and taper.}\label{par:uom-main-234}
Let $(\Sigma_{\rm rec},\gamma)$ denote the receiver spatial slice (round $S^3$ in our case) with
Laplace–Beltrami operator $-\Delta_{\Sigma_{\rm rec}}$ and Sobolev spaces $H^s(\Sigma_{\rm rec})$.
For $L>0$ define the spectral projector
\[
\Pi_{L}\ :=\ \mathbf 1_{\{\sqrt{-\Delta_{\Sigma_{\rm rec}}}\le L\}}\qquad
(\text{“hard cutoff”}),
\]
and, equivalently, a “smooth taper” by a window $w\in C^1([0,\infty))$ with
$w(0)=1$, $w'(x)\le 0$, $w(x)=1$ for $x\le 1-\delta$, $w(x)\to 0$ as $x\to\infty$, setting
\[
w_L(-\Delta_{\Sigma_{\rm rec}})\ :=\ w\!\Big(\frac{\sqrt{-\Delta_{\Sigma_{\rm rec}}}}{L}\Big).
\]
We use either $\Pi_L$ or $w_L(-\Delta_{\Sigma_{\rm rec}})$ to implement a receiver cutoff; when needed we
write $L=L(\Lambda)$ to emphasize the identification with a physical UV scale $\Lambda$.



\paragraph{Geometry and harmonics.}\label{par:uom-main-235}
We work on global de Sitter $\mathrm{dS}_{d+1}$ with Hubble parameter $H>0$ and
spatial slices $S^d$ of unit radius (we set $H=1$ inside proofs/derivations and
restore it by dimensional analysis where needed). The spherical harmonics
$\{Y_{\ell m}\}_{\ell\ge0,\ 1\le m\le D_\ell}$ on $S^d$ are orthonormal,
\[
\int_{S^d}\! \overline{Y_{\ell m}(\Omega)}\,Y_{\ell' m'}(\Omega)\,d\Omega
\;=\;\delta_{\ell\ell'}\delta_{mm'},
\qquad
-\nabla^2 Y_{\ell m}=\ell(\ell+d-1)Y_{\ell m},
\]
with degeneracy $D_\ell=\dfrac{(2\ell+d-1)(\ell+d-2)!}{\ell!\,(d-1)!}$ and
$\mathrm{vol}(S^d)=\dfrac{2\pi^{(d+1)/2}}{\Gamma(\frac{d+1}{2})}$.
The addition theorem (Gegenbauer form) reads
\[
\sum_{m=1}^{D_\ell} Y_{\ell m}(\Omega)\,\overline{Y_{\ell m}(\Omega')}
\;=\;\frac{D_\ell}{\mathrm{vol}(S^d)}\,
\frac{C_\ell^{(\lambda)}(\cos\gamma)}{C_\ell^{(\lambda)}(1)},
\qquad \lambda=\tfrac{d-1}{2},
\]
where $\gamma$ is the geodesic angle between $\Omega,\Omega'$ and
$C_\ell^{(\lambda)}$ are Gegenbauer polynomials normalized by
$C_\ell^{(\lambda)}(1)=\binom{\ell+2\lambda-1}{\ell}$.

\paragraph{Free scalar, mass regime, and BD two–point.}\label{par:uom-main-236}
Let $\phi$ be a free scalar of physical mass $m>0$ (we exclude the minimally
coupled massless case to avoid IR subtleties). Define the dimensionless
parameter
\[
\nu \;=\; \sqrt{\frac{d^2}{4}-\frac{m^2}{H^2}}\ \in\ 
\begin{cases}
(0,\ \tfrac{d}{2}) &\text{if } m^2/H^2<d^2/4 \quad(\text{``light''}),\\[2pt]
i\mathbb R &\text{if } m^2/H^2\ge d^2/4 \quad(\text{``heavy''}).
\end{cases}
\]
The Bunch–Davies (BD) Wightman function depends only on the de Sitter invariant
$Z(x,x')$ and admits the standard hypergeometric form (see e.g.\ Bros–Moschella):
\[
G^+_{\mathrm{BD}}(Z)\;=\;\frac{H^{d-1}}{(4\pi)^{\frac{d+1}{2}}}\,
\frac{\Gamma(\Delta_+)\Gamma(\Delta_-)}{\Gamma(\tfrac{d+1}{2})}\;
{}_2F_1\!\left(\Delta_+,\Delta_-;\ \tfrac{d+1}{2};\ \frac{1+Z}{2}\right),
\quad \Delta_\pm=\frac{d}{2}\pm \nu.
\]
Its Euclidean continuation $G_E$ on the sphere $S^{d+1}$ is obtained by the
usual $t\mapsto -i\tau$ Wick rotation; along a Euclidean time separation
$\tau\ge0$ and angular separation $\gamma\in[0,\pi]$ on $S^d$, we write
$G_E(\tau,\gamma)$ with $Z=\cosh(H\tau)-\tfrac{1}{2}H^2\,\mathrm{dist}^2+O(\mathrm{dist}^4)$
near coincidence.

\paragraph{Angular decomposition and mode covariances.}\label{par:uom-main-237}
Project $G_E(\tau,\gamma)$ onto the $S^d$ harmonics via the addition theorem:
\begin{align*}
G_E(\tau,\gamma)
&= \sum_{\ell=0}^{\infty}\sum_{m=1}^{D_\ell} C_\ell(\tau)\,
Y_{\ell m}(\Omega)\,\overline{Y_{\ell m}(\Omega')}\\[-2pt]
&= \sum_{\ell=0}^{\infty} \frac{D_\ell}{\mathrm{vol}(S^d)}\,
\frac{C_\ell(\tau)}{C_\ell^{(\lambda)}(1)}\,C_\ell^{(\lambda)}(\cos\gamma).
\end{align*}
Inversion yields the \emph{mode covariance} $C_\ell(\tau)$ as
\[
C_\ell(\tau)\;=\;\frac{\mathrm{vol}(S^d)}{D_\ell\,C_\ell^{(\lambda)}(1)}
\int_0^\pi \!G_E(\tau,\gamma)\,C_\ell^{(\lambda)}(\cos\gamma)\,\sin^{d-1}\gamma\,
d\gamma.
\]
We fix the overall normalization by demanding the flat–space limit matches the
harmonic–oscillator convention. Namely, at $\tau=0$ we set
\begin{equation}
C_\ell(0)\;=\;\frac{1}{2\,\omega_\ell},
\qquad
\omega_\ell\;=\;\sqrt{\,H^2\,\ell(\ell+d-1)\,+\,m^2\,},
\label{eq:Cl0}
\end{equation}
which agrees with the $H\to0$ (fixed physical wavevector) limit and with the
static–patch KMS normalization below. This choice uniquely fixes the proportionality
constant in $G_E$ for our calculational scheme.

\paragraph{Band–limit (hard vs smooth) and dS covariance.}\label{par:uom-main-238}
A band–limit at physical scale $\Lambda$ is implemented covariantly by damping
high angular momenta:
\[
C_\ell^{(\Lambda)}(\tau)\;=\;w_\ell(\Lambda)\,C_\ell(\tau),
\qquad
w_\ell(\Lambda)=
\begin{cases}
\mathbf 1_{\{\ell\le \ell_{\max}(\Lambda)\}}, &\text{(hard cutoff),}\\[3pt]
\displaystyle\frac{1}{1+e^{(\ell-\ell_{\max})/\delta\ell}}, &\text{(Fermi–Dirac taper).}
\end{cases}
\]
Both choices preserve rotational covariance on $S^d$; the taper avoids Gibbs–type
ringing, and we use it in numerics. The mapping $\ell_{\max}(\Lambda)$ can be
taken as $\omega_{\ell_{\max}}\approx \Lambda$ with $\omega_\ell$ from
\eqref{eq:Cl0}, i.e.
$\ell_{\max}(\Lambda)\approx \bigl(\Lambda^2-m^2\bigr)^{1/2}/H$ for large $\ell$.


\paragraph{IR regularity and small–$\mu_{ps}$ expansion.}\label{par:uom-main-239}
Throughout we assume $m/H\ge m_\bullet>0$ (any fixed $m_\bullet$ suffices) to
exclude the minimally coupled massless case; then $G_E$ is smooth at
$\tau=0$, $C'_\ell(0)=0$, and the discrete OS Hankel matrices built from
$C_\ell^{(\Lambda)}(\tau)$ have the small–window law
$\lambda_{\min}(\mu_{ps},\Lambda)=-\alpha(\Lambda)\,\mu_{ps}^2+O(\mu_{ps}^4)$ as
$\mu_{ps}\downarrow 0$ (Appendix~\ref{app:OS-discrete}, Lemma~\ref{lem:small-mu}),
with $\alpha(\Lambda)$ computable from $C''_\ell(0)$ and the degeneracies $D_\ell$.

\medskip
These conventions fix \emph{all} normalizations entering (i) the OS/RP test,
(ii) the QND absorber (via $H_{\mathrm{sp}}$), and (iii) the band–limit/taper
used in the bridge; they also make explicit our mass regime $m/H>0$ that avoids
IR pathologies while covering all cosmologically relevant light/heavy cases.

\subsection{Smooth RP degradation and non–commuting limits (Gaussian example)}
\label{subsec:smooth-RP-noncommuting}

\paragraph{Gaussian reference kernel (continued/crossing toy model).}\label{par:uom-main-240}
Let $(S^3,g_{\rm round})$ be the spatial slice used in the continuation toy model with Laplace–Beltrami operator 
$-\Delta_{S^3}$ and eigenpairs $\{-\Delta_{S^3}Y_{\ell m}=\lambda_\ell Y_{\ell m}\}$, 
$\lambda_\ell=\ell(\ell+2)$, $m=1,\dots,D_\ell$, $D_\ell=(\ell+1)^2-\delta_{\ell 0}$. 
For a free (Gaussian) mode in the $\ell$–sector, the Euclidean time–dependence is that of a 
one–dimensional harmonic oscillator with \emph{gap} $\Omega_\ell>0$ determined by the receiver 
kinematics (e.g.\ $\Omega_\ell=\sqrt{\lambda_\ell+\mu_{\rm eff}^2}$; only $\Omega_\ell\sim c\,\ell$ 
as $\ell\to\infty$ is used below). The Euclidean two–point per mode is therefore
\begin{equation}
  C_\ell(\tau)\;=\;\frac{1}{2\,\Omega_\ell}\,e^{-\Omega_\ell \tau},
  \qquad
  C_\ell(0)=\frac{1}{2\Omega_\ell},\quad C_\ell''(0)=\frac{\Omega_\ell}{2}.
  \label{eq:gaussian-mode}
\end{equation}
A band–limit at receiver cutoff $\Lambda$ is implemented spectrally on $S^3$ either by a hard 
cutoff $\ell\le \ell_{\max}(\Lambda)$ (equivalently, $\lambda_\ell\le \Delta_{\rm ret}(\Lambda)$) 
or by a smooth taper $w_\ell(\Lambda)\in[0,1]$ with $w_\ell\simeq 1$ for $\ell\ll \ell_{\max}$ 
and $w_\ell\simeq 0$ for $\ell\gg \ell_{\max}$, yielding 
\(
C_\ell^{(\Lambda)}(\tau)=w_\ell(\Lambda)\,C_\ell(\tau).
\)


\begin{theorem}[Smooth RP degradation at finite cutoff]\label{thm:smooth-RP}
Let \(H^{(\ell)}(\mu_{ps},\Delta\tau;\Lambda)\) be the discrete OS Hankel matrices \eqref{eq:OS-Hankel} built from \eqref{eq:gaussian-mode} with any admissible taper \(w_\ell(\Lambda)\in[0,1]\) that equals \(1\) for \(\ell\le \ell_{\max}-O(1)\).
Then, for fixed \(\Lambda\) and \(\mu_{ps}\downarrow 0\) at fixed grid aspect (or in the continuum limit \(\Delta\tau\to0\)), the smallest OS eigenvalue obeys
\begin{equation}
  \lambda_{\min}(\mu_{ps},\Lambda)\;=\;-\alpha(\Lambda)\,\mu_{ps}^2\ +\ O(\mu_{ps}^4),
  \qquad
  \alpha(\Lambda)\;=\;\frac{c_d}{2}\sum_{\ell\ge0} D_\ell\,w_\ell(\Lambda)\,\omega_\ell^2,
  \label{eq:alpha-def}
\end{equation}
with \(c_d>0\) the dimension–dependent constant from Lemma~\ref{lem:small-mu}.
Moreover, \(\alpha(\Lambda)\) is strictly increasing in \(\Lambda\) and, for hard cutoff with \(\ell_{\max}\sim \Lambda/H\),
\begin{equation}
  \alpha(\Lambda)\;=\;\kappa_{d}\, \Lambda^{d+2}\,H^{-d}\ \big(1+o(1)\big)
  \quad\text{as }\ \Lambda\to\infty,
  \label{eq:alpha-asymp}
\end{equation}
for an explicit constant \(\kappa_{d}>0\) depending only on \(d\) and \(m/H\).
\end{theorem}

\begin{proof}[Sketch]
Lemma~\ref{lem:small-mu} gives
\(\lambda_{\min}^{(\ell)}(\mu_{ps};\Lambda)= -\alpha_\ell(\Lambda)\mu_{ps}^2+O(\mu_{ps}^4)\) with
\(\alpha_\ell(\Lambda)=\frac{C_\ell^{(\Lambda)\prime\prime}(0)}{2C_\ell^{(\Lambda)}(0)}\,c_d\).
Using \eqref{eq:gaussian-mode} and the taper cancellation in the ratio,
\[
\alpha_\ell(\Lambda)=\frac{(\omega_\ell/2)}{2(1/2\omega_\ell)}\,c_d
=\frac{c_d}{2}\,\omega_\ell^2.
\]
Summing over \(\ell\) with degeneracy \(D_\ell\) yields \eqref{eq:alpha-def}.
Monotonicity in \(\Lambda\) is immediate since all terms are positive.
For the asymptotic, use \(D_\ell\sim \frac{2}{(d-1)!}\,\ell^{d-1}\) and
\(\omega_\ell^2\sim H^2\,\ell^2\) as \(\ell\to\infty\), so
\[
\sum_{\ell\le \ell_{\max}} D_\ell\,\omega_\ell^2
\sim \frac{2H^2}{(d-1)!}\sum_{\ell\le \ell_{\max}}\ell^{d+1}
\sim \frac{2H^2}{(d-1)!(d+2)}\,\ell_{\max}^{\,d+2}.
\]
With \(\ell_{\max}\sim \Lambda/H\) this gives \eqref{eq:alpha-asymp}.
Smooth tapers that transition over \(O(1)\) in \(\ell\) leave the leading power unchanged.
\end{proof}

\begin{corollary}[Existence of a negative RP mode at fixed window]\label{cor:neg-mode}
For any fixed \(\mu_{ps}>0\) there exists \(\Lambda_{\mu_{ps}}<\infty\) such that
\(\lambda_{\min}(\mu_{ps},\Lambda)<0\) for all \(\Lambda>\Lambda_{\mu_{ps}}\).
Equivalently, the discrete OS/RP quadratic form admits a negative direction once the cutoff is high enough at any nonzero window size.
\end{corollary}

\begin{proof}
From \eqref{eq:alpha-asymp} and \(\lambda_{\min}=-\alpha(\Lambda)\mu_{ps}^2+O(\mu_{ps}^4)\), take \(\Lambda\) large so that the quadratic term dominates the \(O(\mu_{ps}^4)\) remainder.
\end{proof}

\begin{theorem}[Non–commuting limits]\label{thm:noncommute}
Let \(\lambda_{\min}(\mu_{ps},\Lambda)\) be as above. Then
\[
\lim_{\mu_{ps}\downarrow 0}\ \lim_{\Lambda\to\infty}\ \lambda_{\min}(\mu_{ps},\Lambda)
\;=\; -\infty,
\qquad
\lim_{\Lambda\to\infty}\ \lim_{\mu_{ps}\downarrow 0}\ \lambda_{\min}(\mu_{ps},\Lambda)
\;=\; 0,
\]
so the iterated limits do not commute. More generally, along any joint scaling with \(\mu_{ps}\to0\), \(\Lambda\to\infty\),
\[
\lambda_{\min}(\mu_{ps},\Lambda)\ \sim\ -\,\kappa_{d}\,\mu_{ps}^2\,\Lambda^{d+2}\,H^{-d}
\quad\Rightarrow\quad
\begin{cases}
\mu_{ps}\,\Lambda^{(d+2)/2}\to 0 \ \Rightarrow\ \lambda_{\min}\to 0,\\
\mu_{ps}\,\Lambda^{(d+2)/2}\to \infty \ \Rightarrow\ \lambda_{\min}\to -\infty.
\end{cases}
\]
\end{theorem}

\begin{proof}
Combine Theorem~\ref{thm:smooth-RP} with \eqref{eq:alpha-asymp}. For fixed \(\mu_{ps}>0\), \(\alpha(\Lambda)\to\infty\) so \(\lambda_{\min}\to -\infty\) as \(\Lambda\to\infty\); sending \(\mu_{ps}\downarrow0\) afterwards does not change divergence. Conversely, for each fixed \(\Lambda\), the small–\(\mu_{ps}\) law gives \(\lambda_{\min}\to 0\) as \(\mu_{ps}\downarrow0\), and this \(0\) is independent of \(\Lambda\).
The joint scaling statement follows by substitution.
\end{proof}

\begin{remark}[Scope of the noncommuting-limits obstruction]
\label{rem:uom-main-012}
The noncommuting-limits result concerns the would-be analytic continuation across the AdS$_4\to$dS$_4$ flip at
$\Lambda\to\infty$ (equivalently, the continued dS-side covariance). It does not apply to the healthy AdS-side
receiver at finite cutoff (Thm.~\ref{thm:OS_receiver}), nor to the exhaustion branch where $\Lambda$ approaches a
terminal boundary without a flip.
\end{remark}

\paragraph{Interpretation for the scheduler.}\label{par:uom-main-241}
The RP deficit at finite cutoff degrades \emph{smoothly} as \(-\alpha(\Lambda)\mu_{ps}^2\) with a coefficient that grows like \(\Lambda^{d+2}\). Hence: to keep RP violations uniformly small one must send \(\mu_{ps}\to0\) (fast enough) \emph{before} lifting \(\Lambda\); reversing the order guarantees a negative OS mode at any nonzero \(\mu_{ps}\).
This is precisely the order–of–limits discipline used by the scheduler in the main text.


\section{Regulated conformal intertwiner: covariance and RP bounds}
\label{app:reg-intertwiner}

This appendix collects the precise statements we use about the regulated conformal intertwiner on the receiver side, together with the two ancillary ingredients that make the bridge both \emph{conformally faithful} and \emph{RP–tame} at each tick: (i) radiative–sector projection on the source (AZCT acceptance) and (ii) the zero–DC temporal DoG envelope on the receiver. All notation follows §\ref{subsec:scene-notation} (source band–limit $\Delta$, receiver cutoff $\Lambda$, micro temporal width $\sigma_{t,k}$, etc.).

\subsection{Static--patch radiative projection: covariance under static symmetries}
\label{app:rad-proj}

\paragraph{Setting and acceptance.}\label{par:uom-main-242}
Let $(\mathcal A_{\rm sp},\alpha_t)$ be the source static--patch von Neumann algebra with static time evolution
$\alpha_t$ and spatial rotation automorphisms $\beta_R$ ($R\in SO(3)$). Let $\mathcal A_{\rm rad}\subset\mathcal A_{\rm sp}$
be the (effectively massless) \emph{radiative} subalgebra determined by the late--time bandlimit $\Delta$ (defined via the
receiver band as in App.~\S\ref{app:receiver-bandlimit} and the welded preimage). Let
$P_{\rm rad}:\mathcal A_{\rm sp}\to \mathcal A_{\rm rad}$ be the normal conditional expectation onto $\mathcal A_{\rm rad}$.
Given a normal source state $\rho_{\rm dS}$, define the \emph{radiative projection}
\[
\rho_{\rm rad}\ :=\ \frac{P_{\rm rad}\,\rho_{\rm dS}\,P_{\rm rad}}{\Tr(P_{\rm rad}\rho_{\rm dS})}\,.
\]
We say that $\rho_{\rm dS}$ is \emph{accepted} at band $\Delta$ with tolerance $\varepsilon_{\rm rad}\in[0,1)$ if
\begin{equation}
\label{eq:accept-static}
R_\Delta(\rho_{\rm dS})\ :=\ \Tr\big((\mathbf 1-P_{\rm rad})\,\rho_{\rm dS}\big)\ \le\ \varepsilon_{\rm rad}.
\end{equation}
We write $U_{\rm st}(t)$ and $U_{\rm rot}(R)$ for the unitary implementers of $\alpha_t$ and $\beta_R$ in a fixed GNS representation.

\begin{lemma}[Static–symmetry covariance up to acceptance]\label{lem:static-cov-accept}
Assume:
\begin{enumerate}[leftmargin=1.8em,label=(S\arabic*)]
\item (\emph{Radiative equivariance}) $P_{\rm rad}$ is a normal conditional expectation onto the radiative subalgebra
$\mathcal A_{\rm rad}\subset\mathcal A_{\rm sp}$ and is equivariant for static time and rotations:
$P_{\rm rad}\!\circ\!\alpha_t=\alpha_t\!\circ\!P_{\rm rad}$ and $P_{\rm rad}\!\circ\!\beta_R=\beta_R\!\circ\!P_{\rm rad}$.
\item (\emph{Radiative anchor}) The normalized radiative projection
\(
\rho_{\rm rad}:=\frac{P_{\rm rad}\rho_{\rm dS}P_{\rm rad}}{\Tr(P_{\rm rad}\rho_{\rm dS})}
\)
is invariant under $\alpha_t$ and $\beta_R$ on $\mathcal A_{\rm rad}$:
\(
\Tr\!\big(\alpha_t\beta_R(X)\,\rho_{\rm rad}\big)=\Tr\!\big(X\,\rho_{\rm rad}\big)
\)
for all $X\in\mathcal A_{\rm rad}$ and all $(t,R)\in\mathbb R\times SO(3)$.
\end{enumerate}
Let $\rho_{\rm dS}$ be a normal state with \emph{acceptance} $R_\Delta(\rho_{\rm dS})=\Tr((\mathbf 1-P_{\rm rad})\rho_{\rm dS})\le \varepsilon_{\rm rad}\in[0,1)$, and set
$p:=\Tr(P_{\rm rad}\rho_{\rm dS})=1-\varepsilon_{\rm rad}$. Then for every $X\in\mathcal A_{\rm rad}$ and every $(t,R)$ in a fixed compact subset of $\mathbb R\times SO(3)$ one has
\begin{equation}\label{eq:static-cov-bound}
\Big|\ \Tr\!\big(X\,\rho_{\rm dS}\big) - \Tr\!\big(\alpha_t\beta_R(X)\,\rho_{\rm dS}\big)\ \Big|
\ \le\ \frac{2\,\varepsilon_{\rm rad}}{1-\varepsilon_{\rm rad}}\ \|X\|\ .
\end{equation}
\end{lemma}

\begin{proof}
We work in a fixed faithful normal representation $\pi:\mathcal A_{\rm sp}\to\mathcal B(\mathcal H)$ and suppress $\pi$ from the notation.
By the definition of $P_{\rm rad}$ as the orthogonal projection (in the GNS Hilbert space) onto the closed subspace carrying
$\mathcal A_{\rm rad}$ one has, for every $X\in\mathcal A_{\rm rad}$,
\begin{equation}\label{eq:support}
X\,=\,P_{\rm rad}\,X\,P_{\rm rad}.
\end{equation}
Let $\rho_{\rm dS}$ be realized as a positive trace–class operator on $\mathcal H$ with $\Tr(\rho_{\rm dS})=1$.
Using \eqref{eq:support} and cyclicity of the trace, we compute
\begin{equation}
\label{eq:proj-exp-1}
\begin{aligned}
\Tr\!\big(X\,\rho_{\rm dS}\big)
&=\Tr\!\big(P_{\rm rad} X P_{\rm rad}\,\rho_{\rm dS}\big)\\
&=\Tr\!\big(X\,P_{\rm rad}\rho_{\rm dS}P_{\rm rad}\big)\\
&= p\,\Tr\!\big(X\,\rho_{\rm rad}\big),
\end{aligned}
\end{equation}
where in the last equality we used the normalization of $\rho_{\rm rad}$ and $p=\Tr(P_{\rm rad}\rho_{\rm dS})$.
Similarly, since $\alpha_t\beta_R(X)\in\mathcal A_{\rm rad}$ and $P_{\rm rad}$ is equivariant by (S1),
\begin{equation}
\label{eq:proj-exp-2}
\begin{aligned}
\Tr\!\big(\alpha_t\beta_R(X)\,\rho_{\rm dS}\big)
&=\Tr\!\big(P_{\rm rad}\,\alpha_t\beta_R(X)\,P_{\rm rad}\,\rho_{\rm dS}\big)\\
&=\Tr\!\big(\alpha_t\beta_R(X)\,P_{\rm rad}\rho_{\rm dS}P_{\rm rad}\big)\\
&= p\,\Tr\!\big(\alpha_t\beta_R(X)\,\rho_{\rm rad}\big).
\end{aligned}
\end{equation}
Subtracting \eqref{eq:proj-exp-2} from \eqref{eq:proj-exp-1} and using (S2) (invariance of $\rho_{\rm rad}$) gives the \emph{exact} identity
\begin{equation}\label{eq:exact-zero}
\Tr\!\big(X\,\rho_{\rm dS}\big) - \Tr\!\big(\alpha_t\beta_R(X)\,\rho_{\rm dS}\big)
= p\,\Big(\Tr\!\big(X\,\rho_{\rm rad}\big) - \Tr\!\big(\alpha_t\beta_R(X)\,\rho_{\rm rad}\big)\Big)
= 0.
\end{equation}
Therefore the left–hand side of \eqref{eq:static-cov-bound} vanishes identically; in particular,
\[
\Big|\Tr\!\big(X\,\rho_{\rm dS}\big) - \Tr\!\big(\alpha_t\beta_R(X)\,\rho_{\rm dS}\big)\Big|
\ \le\ \frac{2\,\varepsilon_{\rm rad}}{1-\varepsilon_{\rm rad}}\ \|X\|,
\]
since the right–hand side is nonnegative and finite for $\varepsilon_{\rm rad}<1$. This proves \eqref{eq:static-cov-bound}.

\medskip
\noindent\emph{Remark on robustness.}
If one weakens (S2) to \emph{approximate} invariance of $\rho_{\rm rad}$ on $\mathcal A_{\rm rad}$, say
$\sup_{\|X\|\le 1}|\Tr((\alpha_t\beta_R(X)-X)\rho_{\rm rad})|\le \delta$ uniformly on a compact $(t,R)$–set, then
\eqref{eq:proj-exp-1}–\eqref{eq:proj-exp-2} yield
\[
\big|\Tr\!\big(X\,\rho_{\rm dS}\big) - \Tr\!\big(\alpha_t\beta_R(X)\,\rho_{\rm dS}\big)\big|
\le p\,\delta\,\|X\|
\le \delta\,\|X\|,
\]
so the same inequality \eqref{eq:static-cov-bound} holds with the right–hand side replaced by $\delta\,\|X\|$.
\end{proof}

\paragraph{Compatibility with DP mass dephasing.}\label{par:uom-main-243}
Let $\Phi_{\rm DP}$ be a Gaussian CP channel generated by a DP--type Lindbladian on the source mass sector,
\[
\mathcal L_{\rm DP}\rho\;=\;\frac{K^{\rm DP}}{2\hbar}\int d^3x\,d^3y\;V_{\rm eff}(x,y)\,
[\hat\mu_\Delta(x),[\hat\mu_\Delta(y),\rho]]\ ,
\]
with $V_{\rm eff}$ the static brane Green kernel. Assume $[\mathcal L_{\rm DP},\alpha_t]=[\mathcal L_{\rm DP},\beta_R]=0$
(static KMS bath and rotationally invariant kernel) and that $\Phi_{\rm DP}$ is band--limited:
$P_{\rm rad}\circ \Phi_{\rm DP}=\Phi_{\rm DP}\circ P_{\rm rad}$. Then the bound \eqref{eq:static-cov-bound} holds with $\rho_{\rm dS}$
replaced by $\Phi_{\rm DP}(\rho_{\rm dS})$ and the same tolerance $\varepsilon_{\rm rad}$; in particular, the DP channel does
\emph{not} deteriorate static–symmetry covariance up to acceptance. Moreover, since $\Phi_{\rm DP}$ further contracts off–diagonal
coherences in the non–radiative sector, the effective residual $R_\Delta(\Phi_{\rm DP}(\rho_{\rm dS}))$ is $\le R_\Delta(\rho_{\rm dS})$,
so the right–hand side of \eqref{eq:static-cov-bound} can only \emph{decrease} under $\Phi_{\rm DP}$.


\subsection{Approximate intertwiner at finite receiver cutoff}
\label{app:approx-int}

Let ${\sf K}_\Delta:\mathcal H_{\rm rad}\to\mathcal H_\partial$ be the (unregulated) bulk–to–boundary intertwiner on primaries of dimension $\Delta$, i.e.
${\sf K}_\Delta\circ U_{\rm dS}(g)=\mathrm{Ad}(U_\partial(g))\circ{\sf K}_\Delta$ for $g\in\mathrm{Conf}(S^3)$, where $U_\partial(g)$ is the unitary boundary action on $\mathcal A_\partial$.
Introduce the \emph{regulated} intertwiner at receiver cutoff $\Lambda$:
\[
{\sf K}_\Delta^{(\varepsilon,\Lambda)}
\ :=\ {\sf K}_\Delta\ \circ\ \Pi_{L(\Lambda)}\ \circ\ M_\varepsilon,
\qquad \varepsilon\simeq L(\Lambda)^{-1},
\]
where $\Pi_{L(\Lambda)}$ projects onto $S^3$ harmonics with $\ell\le L(\Lambda)$ and $M_\varepsilon$ is a UV mollifier (kernel width~$\varepsilon$).

\begin{lemma}[Approximate intertwiner bound]\label{lem:approx-intertwiner-new}
Let $S^3$ denote the unit $3$--sphere with Laplace--Beltrami operator $\Delta_{S^3}$ and Sobolev spaces $H^s(S^3)$, $s\in\mathbb R$.
Let $\Pi_{L}$ be the spectral projector onto spherical harmonics of degree $\le L$:
\[
\Pi_L \ :=\ \mathbf 1_{\{\sqrt{-\Delta_{S^3}}\le L\}}\,,\qquad L\in\mathbb N.
\]
Let $U_{\rm dS}(g)$ be the pullback by $g\in \mathrm{Conf}(S^3)$ acting on functions/distributions on $S^3$,
and let ${\sf K}_\Delta:\,L^2(S^3)\to L^2(\partial\mathrm{AdS}_4)$ be the (unregulated) conformal intertwiner for primaries of dimension $\Delta$, satisfying
\begin{equation}\label{eq:exact-intertwine}
{\sf K}_\Delta\circ U_{\rm dS}(g)\;=\;\mathrm{Ad}\big(U_\partial(g)\big)\circ{\sf K}_\Delta\qquad\forall g\in\mathrm{Conf}(S^3),
\end{equation}
where $U_\partial(g)$ is the induced unitary action on the boundary Hilbert space.
Let $M_\varepsilon$ be a smoothing operator with integral kernel $K_\varepsilon(x,y)$ supported on geodesic distances $d(x,y)\le C\varepsilon$ and with pointwise bounds
\begin{equation}\label{eq:kernel-bounds}
\sup_{x}\int_{S^3} |K_\varepsilon(x,y)|\,dy\le C_0,\quad
\sup_{x}\int_{S^3} \varepsilon^{-1}\,|K_\varepsilon(x,y)|\,d(x,y)\,dy\le C_1,
\end{equation}
uniformly for $0<\varepsilon\le \varepsilon_0$. Define the regulated intertwiner
\[
{\sf K}_\Delta^{(\varepsilon,\Lambda)}\ :=\ {\sf K}_\Delta\,\circ\,\Pi_{L(\Lambda)}\,\circ\,M_\varepsilon,\qquad \varepsilon\simeq L(\Lambda)^{-1}.
\]
Then there exist constants $\alpha>0$, $s_*\!>2$ and, for each compact $G\subset\mathrm{Conf}(S^3)$ and $s\ge s_*$, a constant $C_\Delta(G,s)$ such that
\begin{equation}\label{eq:approx-intertwiner-bound}
\big\|{\sf K}_\Delta^{(\varepsilon,\Lambda)} \circ U_{\rm dS}(g)
- \mathrm{Ad}\big(U_\partial(g)\big)\circ {\sf K}_\Delta^{(\varepsilon,\Lambda)}\big\|_{H^s(S^3)\to L^2(\partial\mathrm{AdS}_4)}
\ \le\ C_\Delta(G,s)\,\varepsilon^\alpha,\qquad \forall g\in G.
\end{equation}
In particular, \eqref{eq:approx-intertwiner-bound} holds in operator norm on $L^2(S^3)\to L^2(\partial\mathrm{AdS}_4)$ (by $H^s\hookrightarrow L^2$).
\end{lemma}

\begin{proof}
Fix $s>2$ (the threshold $s_*:=2$ suffices on $S^3$). Using \eqref{eq:exact-intertwine} and linearity,
\begin{equation}
\label{eq:3comm}
\begin{aligned}
{\sf K}_\Delta^{(\varepsilon,\Lambda)} \circ U_{\rm dS}(g) - \mathrm{Ad}(U_\partial(g))\circ {\sf K}_\Delta^{(\varepsilon,\Lambda)}
&= {\sf K}_\Delta \Big(\Pi_{L(\Lambda)}M_\varepsilon U_{\rm dS}(g) - U_{\rm dS}(g)\Pi_{L(\Lambda)}M_\varepsilon\Big)\\
&= {\sf K}_\Delta\Big([\Pi_{L(\Lambda)},U_{\rm dS}(g)] M_\varepsilon + U_{\rm dS}(g) [M_\varepsilon,\Pi_{L(\Lambda)}] + [M_\varepsilon,U_{\rm dS}(g)]\Big).
\end{aligned}
\end{equation}
We bound each commutator $H^s(S^3)\to L^2(S^3)$ uniformly on $g\in G$, then multiply by $\|{\sf K}_\Delta\|_{L^2\to L^2}$.

\smallskip
\emph{(i) Spectral projector vs.\ pullback.}
Write $I=\Pi_L+(I-\Pi_L)$ and observe the identity
\[
[\Pi_L,U]=\Pi_L U(I-\Pi_L) - (I-\Pi_L)U\Pi_L,
\]
with $U:=U_{\rm dS}(g)$. Hence, for any $f\in H^s(S^3)$,
\begin{align*}
\|[\Pi_L,U]f\|_{L^2}
&\le \|\Pi_L U(I-\Pi_L)f\|_{L^2} + \|(I-\Pi_L)U\Pi_L f\|_{L^2}\\
&\le \|U(I-\Pi_L)f\|_{L^2} + \|(I-\Pi_L)U\Pi_L f\|_{L^2}.
\end{align*}
Composition by a smooth diffeomorphism is bounded on $H^s(S^3)$ (see e.g.\ standard results on Sobolev pullbacks on compact manifolds), so
$\|U h\|_{H^s}\le C(G,s)\|h\|_{H^s}$ and $\|Uh\|_{L^2}\le C(G,0)\|h\|_{L^2}$ for all $g\in G$.
Moreover, spectral approximation yields (functional calculus for $\sqrt{-\Delta_{S^3}}$)
\begin{equation}\label{eq:spectral-approx}
\|(I-\Pi_L)h\|_{L^2}\ \le\ (1+L^2)^{-s/2}\,\|(I-\Delta_{S^3})^{s/2} h\|_{L^2}\ \le\ C_s\,L^{-s}\,\|h\|_{H^s},\qquad s>0.
\end{equation}
Consequently,
\begin{equation}
\label{eq:proj-U-bound}
\begin{aligned}
\|[\Pi_L,U]\|_{H^s\to L^2}
&\le C(G,0)\,C_s\,L^{-s} + C_s\,L^{-s}\,\|U\Pi_L\|_{H^s\to H^s}\\
&\le C(G,s)\,L^{-s}.
\end{aligned}
\end{equation}
Thus $\|[\Pi_{L(\Lambda)},U_{\rm dS}(g)]M_\varepsilon\|_{H^s\to L^2}\le C(G,s)\,L(\Lambda)^{-s}\,\|M_\varepsilon\|_{H^s\to H^s}\le C(G,s)\,L(\Lambda)^{-s}$, since $M_\varepsilon$ is bounded $H^s\to H^s$.

\smallskip
\emph{(ii) Projector vs.\ mollifier.}
If $M_\varepsilon$ is chosen as a spectral multiplier of $\sqrt{-\Delta_{S^3}}$ (e.g.\ the heat kernel $e^{\varepsilon^2\Delta_{S^3}}$), then $[M_\varepsilon,\Pi_L]=0$ identically and the term vanishes. Otherwise, under the kernel conditions \eqref{eq:kernel-bounds}, $M_\varepsilon$ is a smoothing operator of any order: for each $s\ge 0$,
\begin{equation}\label{eq:M-eps-smoothing}
\|M_\varepsilon\|_{L^2\to H^s}\ \le\ C_s\,\varepsilon^{-s},\qquad 
\|M_\varepsilon\|_{H^s\to H^s}\ \le\ C_s.
\end{equation}
A standard commutator estimate for Fourier multipliers and smoothing operators (see, e.g., symbolic calculus on compact manifolds) gives, for $s>0$,
\begin{equation}\label{eq:comm-M-Pi}
\|[M_\varepsilon,\Pi_L]\|_{H^s\to L^2}\ \le\ C_s\,\min\{L^{-s},\,\varepsilon^{s}\},
\end{equation}
which for the regime $L\sim\varepsilon^{-1}$ reduces to $\|[M_\varepsilon,\Pi_L]\|_{H^s\to L^2}\le C_s\,\varepsilon^{s}$.

\smallskip
\emph{(iii) Mollifier vs.\ pullback.}
Write, for $f\in L^2(S^3)$ and $x\in S^3$,
\[
\big(M_\varepsilon U f\big)(x)-\big(U M_\varepsilon f\big)(x)
=\int_{S^3}\!\Big(K_\varepsilon(x,y)f(g^{-1}y)-K_\varepsilon(g^{-1}x,y)f(y)\Big)\,dy.
\]
Change variables $z=g^{-1}y$ in the first integral and use that $g$ is smooth with Jacobian uniformly bounded on $G$; by the mean value theorem and \eqref{eq:kernel-bounds},
\[
\big|K_\varepsilon(x,g z)-K_\varepsilon(g^{-1}x,z)\big|
\ \le\ \sup_{\theta\in[0,1]}\|\nabla_x K_\varepsilon(\gamma_\theta,z)\|\cdot d\big(x,g^{-2}x\big)
\ \le\ C\,\varepsilon^{-1}\,\varepsilon\ =\ C,
\]
where $\gamma_\theta$ is the geodesic segment between $x$ and $g^{-1}x$ and we used that $K_\varepsilon$ is supported on balls of radius $O(\varepsilon)$ and varies at scale $\varepsilon$. A Schur test then yields (uniformly in $g\in G$)
\begin{equation}\label{eq:M-U-bound}
\|[M_\varepsilon,U_{\rm dS}(g)]\|_{L^2\to L^2}\ \le\ C(G)\,\varepsilon,\qquad 
\|[M_\varepsilon,U_{\rm dS}(g)]\|_{H^s\to L^2}\ \le\ C(G,s)\,\varepsilon,
\end{equation}
since $H^s\hookrightarrow L^2$ for $s\ge 0$.

\smallskip
Combining \eqref{eq:3comm}, \eqref{eq:proj-U-bound}, \eqref{eq:comm-M-Pi}, and \eqref{eq:M-U-bound}, and using $\|{\sf K}_\Delta\|_{L^2\to L^2}<\infty$, we obtain
\[
\big\|{\sf K}_\Delta^{(\varepsilon,\Lambda)} \circ U_{\rm dS}(g)
- \mathrm{Ad}(U_\partial(g))\circ {\sf K}_\Delta^{(\varepsilon,\Lambda)}\big\|_{H^s\to L^2}
\ \le\ C_\Delta(G,s)\,\big(L(\Lambda)^{-s}+\varepsilon^{s}+\varepsilon\big).
\]
Choosing $L(\Lambda)\simeq \varepsilon^{-1}$ gives the bound $C_\Delta(G,s)\,\varepsilon^\alpha$ with $\alpha=\min\{1,s\}>0$. This proves \eqref{eq:approx-intertwiner-bound}.
\end{proof}


\begin{remark}[Conformality deficit monitor]
\label{rem:conf-monitor}
Define the conformality mismatch at cutoff $\Lambda$ by
\[
\mathcal D_{\rm conf}(\varepsilon,\Lambda)\ :=\ 
\sup_{g\in G}\ \big\|{\sf K}_\Delta^{(\varepsilon,\Lambda)} \circ U_{\rm dS}(g)
- \mathrm{Ad}(U_\partial(g))\circ {\sf K}_\Delta^{(\varepsilon,\Lambda)}\big\|_{H^s\to L^2}.
\]
Lemma~\ref{lem:approx-intertwiner-new} gives $\mathcal D_{\rm conf}=O(\varepsilon^\alpha)$ with $\varepsilon\simeq L(\Lambda)^{-1}$.
\end{remark}

\subsection{Temporal DoG on the receiver: zero–DC and scale control}
\label{app:dog-temporal}

Let the micro temporal envelope at tick $k$ be $\psi_{t,\sigma_{t,k}}(t)=g'_{\sigma_{t,k}}(t-t_k)$ with
$g'_{\sigma}(t)=-(t/\sigma^2)\exp(-t^2/2\sigma^2)$ and $\int_{\mathbb R}\psi_{t,\sigma_{t,k}}=0$.
Denote by $\mathcal C_{\sigma_{t,k}}$ the convolution operator $f\mapsto \psi_{t,\sigma_{t,k}}*f$.

\begin{lemma}[Zero--DC and scale--distortion on a band--gap]\label{lem:DoG-comm-new}
Let $\sigma>0$ and let $g'_\sigma(t):=-(t/\sigma^2)\,e^{-t^2/(2\sigma^2)}$ for $t\in\mathbb R$.
Define the convolution operator $\mathcal C_\sigma:L^2(\mathbb R)\to L^2(\mathbb R)$ by
\[
(\mathcal C_\sigma f)(t):=(g'_\sigma*f)(t)=\int_{\mathbb R} g'_\sigma(t-\tau)\,f(\tau)\,d\tau.
\]
Let $D$ be the generator of dilations on the Euclidean time axis, i.e.\ $(\widehat{Df})(\omega)=i\,\partial_\omega \widehat f(\omega)$ in Fourier variables.
Then:
\begin{enumerate}
\item[\emph{(a)}] (\emph{Zero--DC}) $\displaystyle \int_{\mathbb R} g'_\sigma(t)\,dt=0$.
\item[\emph{(b)}] (\emph{Scale--distortion on a band--gap}) Let $\Omega>0$ and assume $\widehat f(\omega)=0$ for $|\omega|<\Omega$. Then
\begin{equation}\label{eq:DoG-comm-main}
\|[D,\mathcal C_\sigma]\,f\|_{L^2(\mathbb R)}\ \le\ \sup_{|\omega|\ge \Omega}\ \big(1+\sigma^2\omega^2\big)\,e^{-\frac12\sigma^2\omega^2}\ \cdot\ \|f\|_{L^2(\mathbb R)}.
\end{equation}
In particular,
\begin{equation}\label{eq:DoG-comm-decay}
\|[D,\mathcal C_\sigma]\,f\|_{L^2}\ \le\ \max\!\Big\{2\,e^{-1/2},\ \big(1+\sigma^2\Omega^2\big)\,e^{-\frac12\sigma^2\Omega^2}\Big\}\,\|f\|_{L^2},
\end{equation}
and hence $\|[D,\mathcal C_\sigma]\,f\|_{L^2}\le \big(1+\sigma^2\Omega^2\big)\,e^{-\frac12\sigma^2\Omega^2}\|f\|_{L^2}$ whenever $\sigma\Omega\ge 1$.
\end{enumerate}
\end{lemma}

\begin{proof}
\emph{(a)} Since $g'_\sigma$ is odd and integrable,
\[
\int_{\mathbb R} g'_\sigma(t)\,dt\ =\ -\frac1{\sigma^2}\int_{\mathbb R} t\,e^{-\frac{t^2}{2\sigma^2}}\,dt\ =\ 0.
\]

\smallskip
\emph{(b)} Let $\mathcal F$ denote the unitary Fourier transform on $L^2(\mathbb R)$ with the convention $\widehat f(\omega)=\int_{\mathbb R}e^{-i\omega t}f(t)\,dt$. A direct computation gives
\[
\widehat{g'_\sigma}(\omega)\ =\ i\,\omega\,e^{-\frac12\sigma^2\omega^2}\ =:\ m_\sigma(\omega),\qquad \partial_\omega \widehat{g'_\sigma}(\omega)\ =\ i\,(1-\sigma^2\omega^2)\,e^{-\frac12\sigma^2\omega^2}\ =:\ m'_\sigma(\omega).
\]
Thus $\mathcal C_\sigma$ is the Fourier multiplier by $m_\sigma$, i.e.
$\widehat{\mathcal C_\sigma f}(\omega)=m_\sigma(\omega)\,\widehat f(\omega)$ for $f\in L^2$.
Because $(\widehat{Df})(\omega)=i\,\partial_\omega\widehat f(\omega)$, we have for all $f\in \mathcal S(\mathbb R)$ (dense in $L^2$),
\[
\widehat{[D,\mathcal C_\sigma]f}(\omega)
= i\,\partial_\omega(m_\sigma\,\widehat f)(\omega)-m_\sigma(\omega)\,i\,\partial_\omega \widehat f(\omega)
= i\,m'_\sigma(\omega)\,\widehat f(\omega).
\]
By Plancherel,
\[
\|[D,\mathcal C_\sigma]f\|_{L^2}=\|\,m'_\sigma\,\widehat f\,\|_{L^2}.
\]
Under the band–gap assumption $\widehat f(\omega)=0$ for $|\omega|<\Omega$, we may restrict the $L^2$–norm to $\{|\omega|\ge \Omega\}$ and obtain
\[
\|[D,\mathcal C_\sigma]f\|_{L^2}
\le \sup_{|\omega|\ge \Omega} |m'_\sigma(\omega)|\ \cdot\ \|\widehat f\|_{L^2}
= \sup_{|\omega|\ge \Omega} \big|1-\sigma^2\omega^2\big|\,e^{-\frac12\sigma^2\omega^2}\ \cdot\ \|f\|_{L^2},
\]
again by Plancherel. Since $\big|1-\sigma^2\omega^2\big|\le 1+\sigma^2\omega^2$, this yields \eqref{eq:DoG-comm-main}.

For the explicit bound \eqref{eq:DoG-comm-decay}, define $\phi(x):=(1+x^2)e^{-x^2/2}$ for $x\ge 0$.
Then $\phi'(x)=x(1-x^2)e^{-x^2/2}$ so $\phi$ increases on $[0,1]$, attains its maximum at $x=1$ with $\phi(1)=2e^{-1/2}$, and decreases on $[1,\infty)$. Therefore
\[
\sup_{|\omega|\ge \Omega} \big(1+\sigma^2\omega^2\big)e^{-\frac12\sigma^2\omega^2}
=\sup_{x\ge \sigma\Omega}\phi(x)\ \le\ \max\!\big\{\,\phi(1),\ \phi(\sigma\Omega)\,\big\}
=\max\!\Big\{2e^{-1/2},\ \big(1+\sigma^2\Omega^2\big)e^{-\frac12\sigma^2\Omega^2}\Big\},
\]
which implies \eqref{eq:DoG-comm-decay}. In particular, if $\sigma\Omega\ge 1$ then $\phi$ is decreasing on $[\sigma\Omega,\infty)$ and
$\sup_{x\ge \sigma\Omega}\phi(x)=\phi(\sigma\Omega)$, yielding the last statement.
\end{proof}


\subsection{Per–tick conformal fidelity and RP boundedness}
\label{app:tick-fidelity-RP}

At tick $k$ let the operational bridge act on radiative test fields by
\[
\mathcal B_k\ :=\ \Big(\sum_{\Delta} c_\Delta\ {\sf K}_\Delta^{(\varepsilon_k,\Lambda_k)}\Big)\ \circ\ (\psi_{t,\sigma_{t,k}}*\cdot)\,,
\qquad {\sf K}_\Delta^{(\varepsilon_k,\Lambda_k)}={\sf K}_\Delta\circ \Pi_{L(\Lambda_k)}\circ M_{\varepsilon_k},\quad \varepsilon_k\simeq L(\Lambda_k)^{-1}.
\]
Let $\rho_{k+1}^{\rm AdS}$ be the receiver boundary state obtained by applying $\mathcal B_k$ to the (accepted) source content and completing the Stinespring injection. Denote by $\Delta I_k$ the compiled payload and recall the micro throughput equality
\[
\Delta I_k\;=\;A_k^2\,G_k(\sigma_{t,k},\sigma_{x,k};\Lambda_k),\qquad G_k>0,
\]
for $h_k(t,\Omega)=A_k\,\psi_{t,\sigma_{t,k}}(t)\,\Psi_{x,\sigma_{x,k}}(\Omega)$.

\begin{theorem}[Conformal fidelity and RP boundedness per tick]\label{thm:tick-fidelity-RP}
Fix a tick index $k$. Let $\rho_{\rm dS}$ be a source state with radiative acceptance $R_\Delta(\rho_{\rm dS})\le \varepsilon_{\rm rad}$, and let the receiver bridge be
\[
\mathcal B_k\ :=\ \Big(\sum_\Delta c_\Delta\,{\sf K}_\Delta^{(\varepsilon_k,\Lambda_k)}\Big)\ \circ\ (\psi_{t,\sigma_{t,k}}*\cdot)\,,
\qquad {\sf K}_\Delta^{(\varepsilon_k,\Lambda_k)}:={\sf K}_\Delta\circ \Pi_{L(\Lambda_k)}\circ M_{\varepsilon_k},
\]
where $\psi_{t,\sigma_{t,k}}(t)=g'_{\sigma_{t,k}}(t-t_k)$ is the temporal DoG envelope and $\varepsilon_k\simeq L(\Lambda_k)^{-1}$. Denote by $\rho_{k+1}^{\rm AdS}$ the (Gaussian) boundary state obtained from $\rho_{\rm dS}$ by applying $\mathcal B_k$ to the radiative content and completing the Stinespring isometry (the compensator dilation).

\smallskip
\emph{(Conformal fidelity).}
Fix a finite family of primaries $\{\mathcal O_{\Delta_i}\}_{i=1}^n$ on $\mathcal A_\partial$ and a compact $G\subset\mathrm{Conf}(S^3)$. There exist constants $C_1,C_2,C_3>0$ and $\alpha>0$, depending only on $G$, the family and Sobolev exponents, such that for all $g\in G$,
\begin{equation}\label{eq:conf-fid-bound}
\Bigl|\ \big\langle \prod_{i=1}^n \mathcal O_{\Delta_i}(x_i)\big\rangle_{\rho_{k+1}^{\rm AdS}}
-\prod_{i=1}^n \Omega(g,x_i)^{-\Delta_i}\,
\big\langle \prod_{i=1}^n \mathcal O_{\Delta_i}(g\!\cdot\!x_i)\big\rangle_{\rho_{k+1}^{\rm AdS}}\ \Bigr|
\ \le\ C_1\,\varepsilon_{\rm rad}\ +\ C_2\,\varepsilon_k^\alpha\ +\ C_3\,\Phi(\sigma_{t,k}\Omega_{\rm gap}),
\end{equation}
where $\Omega_{\rm gap}=\Omega_{\rm gap}(\Delta)>0$ is the temporal band--gap inherited from the source band--limit and
$\Phi(x):=\max\{2e^{-1/2},\, (1+x^2)e^{-x^2/2}\}$.

\smallskip
\emph{(RP boundedness).}
Let $h_k(t,\Omega)=A_k\,\psi_{t,\sigma_{t,k}}(t)\,\Psi_{x,\sigma_{x,k}}(\Omega)$ be the welded micro envelope (temporal DoG and zonal spatial Gaussian/LoG at width $\sigma_{x,k}$), with throughput equality
$\Delta I_k=A_k^2\,G_k(\sigma_{t,k},\sigma_{x,k};\Lambda_k)$ and $G_k>0$. Then there is a bounded, positive semidefinite operator $\mathsf R_k(\Lambda_k)$ on the envelope Hilbert space such that the receiver RP deficit on the micro window obeys
\begin{equation}\label{eq:RP-bound}
\mathrm{RP}\!\big(\rho_{k+1}^{\rm AdS}\big)\ \le\ \big\langle h_k,\ \mathsf R_k(\Lambda_k)\,h_k\big\rangle\ +\ C_{\rm reg}^{(k)}(\varepsilon_k,\sigma_{t,k};\Delta),
\end{equation}
with a regulator remainder $C_{\rm reg}^{(k)}=O(\varepsilon_k^\alpha)+O\!\big(\Phi(\sigma_{t,k}\Omega_{\rm gap})\big)$ uniformly on compact parameter sets. In particular,
\begin{equation}\label{eq:RP-bound-DeltaI}
\mathrm{RP}\!\big(\rho_{k+1}^{\rm AdS}\big)\ \le\ \frac{q_k(\sigma_{t,k},\sigma_{x,k};\Lambda_k)}{G_k(\sigma_{t,k},\sigma_{x,k};\Lambda_k)}\ \Delta I_k\ +\ C_{\rm reg}^{(k)}(\varepsilon_k,\sigma_{t,k};\Delta),
\end{equation}
where $q_k:=\langle \psi_{t,\sigma_{t,k}}\Psi_{x,\sigma_{x,k}},\ \mathsf R_k(\Lambda_k)\,\psi_{t,\sigma_{t,k}}\Psi_{x,\sigma_{x,k}}\rangle>0$. Hence for each finite tick the RP deficit is finite and scales linearly with the payload $\Delta I_k$ up to regulator--decaying terms.
\end{theorem}

\begin{proof}
\emph{Step 1 (Conformal fidelity).}
Decompose the total deviation from exact boundary conformal covariance into three independent contributions corresponding to (i) source radiative projection, (ii) receiver regulated intertwiner, and (iii) receiver temporal DoG envelope.

\smallskip
\emph{(i) Source radiative projection.} By Lemma~\ref{lem:static-cov-accept} (applied on the source) and the fact that the bridge acts trivially on $\mathcal A_{\rm rad}(\mathscr I^+)$ expectations modulo the dilation (Stinespring) that preserves local Ward identities, the contribution of the non--radiative component of $\rho_{\rm dS}$ is bounded by $C_1\,\varepsilon_{\rm rad}$ for some $C_1$ depending only on the operator norms of the chosen primaries and on $G$.

\smallskip
\emph{(ii) Regulated intertwiner.} For the regulated kernel ${\sf K}_\Delta^{(\varepsilon_k,\Lambda_k)}$ we have the approximate intertwiner bound of Lemma~\ref{lem:approx-intertwiner-new}: for $s>2$,
\[
\big\|{\sf K}_\Delta^{(\varepsilon_k,\Lambda_k)} \circ U_{\rm dS}(g) - \mathrm{Ad}(U_\partial(g))\circ {\sf K}_\Delta^{(\varepsilon_k,\Lambda_k)}\big\|_{H^s\to L^2}
\ \le\ C_\Delta(G,s)\,\varepsilon_k^\alpha.
\]
By Sobolev embedding and multilinearity of $n$--point functions in the Gaussian regime, this induces an additive error bounded by $C_2\,\varepsilon_k^\alpha$ for some $C_2$ depending only on $G,s$ and the family.

\smallskip
\emph{(iii) Temporal DoG.} The DoG convolution in time commutes with dilations up to the commutator controlled by Lemma~\ref{lem:DoG-comm-new}. On the micro window, the accepted radiative content has a temporal band--gap $\Omega_{\rm gap}>0$ (determined by the source band--limit $\Delta$; in particular $\Omega_{\rm gap}\ge H\sqrt{3}$ if $\ell_{\min}=1$ on $S^3$). Applying Lemma~\ref{lem:DoG-comm-new} with $\Omega=\Omega_{\rm gap}$ shows that the difference between correlators with and without the DoG insertion is bounded by $C_3\,\Phi(\sigma_{t,k}\Omega_{\rm gap})$ for some $C_3$ depending only on the family and $G$.

\smallskip
Summing the three contributions and using the triangle inequality yields \eqref{eq:conf-fid-bound}.

\bigskip
\emph{Step 2 (RP boundedness).}
Let $\mathscr H$ be the real Hilbert space of admissible (real--valued) welded envelopes $h=h(t,\Omega)$ on the receiver micro window, equipped with the $L^2$ inner product $\langle h_1,h_2\rangle=\int h_1 h_2$ (including the $S^3$ measure on $\Omega$). For each angular momentum $\ell$ and the chosen OS discretization $(\mu_{ps},\Delta\tau)$ on the receiver side, the corresponding Hankel matrix $H^{(\ell)}$ is \emph{linear} in the (time--time) two--point function, which in turn depends \emph{linearly} on $h$ at the Gaussian level (Stinespring linear response of the boundary covariance). Therefore, there exists a bounded linear map
\[
\mathcal T_k:\ \mathscr H\ \longrightarrow\ \bigoplus_{\ell\ge 0} \mathbb C^{(N+1)\times(N+1)}
\]
that sends $h$ to the collection of Hankel matrices $\{H^{(\ell)}[h]\}_\ell$ induced by the bridge with envelope $h$ at cutoff $\Lambda_k$. Let $\langle \!\langle A,B\rangle\!\rangle:=\sum_\ell \mathrm{tr}\!\big(A^{(\ell)} (B^{(\ell)})^\ast\big)$ be the Frobenius inner product on the target space of $\mathcal T_k$ and let $\mathbf P_-^{(\ell)}$ denote the spectral projector onto the negative eigenspace of $H^{(\ell)}[h]$. The (discretized) RP deficit satisfies
\[
\mathrm{RP}(\rho_{k+1}^{\rm AdS})\ =\ \sum_{\ell}\ \sum_{\lambda_m^{(\ell)}[h]<0}\big|\lambda_m^{(\ell)}[h]\big|
\ \le\ \sum_\ell \big\|\big(H^{(\ell)}[h]\big)_-\big\|_1
\ \le\ \sum_\ell \|H^{(\ell)}[h]\|_1
\ \le\ C_{\rm OS}\ \sum_\ell \|H^{(\ell)}[h]\|_F,
\]
for a constant $C_{\rm OS}$ depending on the grid size (using $\|A\|_1\le \sqrt{d}\,\|A\|_F$ on $d\times d$ matrices). By Cauchy--Schwarz and the boundedness of $\mathcal T_k$,
\[
\mathrm{RP}(\rho_{k+1}^{\rm AdS})
\ \le\ C_{\rm OS}\,\Big(\sum_\ell 1\Big)^{1/2}\,\Big(\sum_\ell \|H^{(\ell)}[h]\|_F^2\Big)^{1/2}
\ \le\ C'_k\,\|\mathcal T_k h\|\ \le\ C''_k\,\|h\|_{\mathscr H},
\]
with constants $C'_k,C''_k$ depending on $(\mu_{ps},\Delta\tau,\Lambda_k)$ but finite for each fixed tick $k$. In fact, because the dependence $h\mapsto H^{(\ell)}[h]$ is linear, the map $h\mapsto \sum_\ell \|H^{(\ell)}[h]\|_F^2$ is a continuous \emph{quadratic} functional on $\mathscr H$, hence by the Riesz representation theorem there exists a unique bounded, self--adjoint, positive semidefinite operator $\mathsf R_k(\Lambda_k):\mathscr H\to\mathscr H$ such that
\begin{equation}\label{eq:Rk-def}
\sum_\ell \|H^{(\ell)}[h]\|_F^2\ =\ \big\langle h,\ \mathsf R_k(\Lambda_k)\,h\big\rangle_{\mathscr H}\,.
\end{equation}
The regulator imperfections (finite cutoff and finite temporal envelope band--gap) perturb the exact OS matrices by amounts controlled by Lemma~\ref{lem:approx-intertwiner-new} and Lemma~\ref{lem:DoG-comm-new}, yielding an additive remainder $C_{\rm reg}^{(k)}(\varepsilon_k,\sigma_{t,k};\Delta)=O(\varepsilon_k^\alpha)+O(\Phi(\sigma_{t,k}\Omega_{\rm gap}))$ in the bound. Altogether we obtain \eqref{eq:RP-bound}.

Finally, evaluate \eqref{eq:RP-bound} on the welded micro envelope $h_k=A_k\,\psi_{t,\sigma_{t,k}}\Psi_{x,\sigma_{x,k}}$. By bilinearity,
\[
\big\langle h_k,\ \mathsf R_k(\Lambda_k)\,h_k\big\rangle\ =\ A_k^2\,\big\langle \psi_{t,\sigma_{t,k}}\Psi_{x,\sigma_{x,k}},\ \mathsf R_k(\Lambda_k)\,\psi_{t,\sigma_{t,k}}\Psi_{x,\sigma_{x,k}}\big\rangle\ =\ A_k^2\,q_k(\sigma_{t,k},\sigma_{x,k};\Lambda_k).
\]
Using the throughput equality $\Delta I_k=A_k^2\,G_k(\sigma_{t,k},\sigma_{x,k};\Lambda_k)$ we conclude
\[
\big\langle h_k,\ \mathsf R_k(\Lambda_k)\,h_k\big\rangle\ =\ \frac{q_k(\sigma_{t,k},\sigma_{x,k};\Lambda_k)}{G_k(\sigma_{t,k},\sigma_{x,k};\Lambda_k)}\ \Delta I_k,
\]
which, inserted into \eqref{eq:RP-bound}, yields \eqref{eq:RP-bound-DeltaI}. This completes the proof.
\end{proof}


\begin{corollary}[Anti–Zeno limit]
\label{cor:anti-zeno}
If the scheduler enforces $\varepsilon_k\downarrow 0$ and $\sigma_{t,k}\,\Omega_{\rm gap}(\Delta)\uparrow\infty$ along accepted ticks while keeping $\sup_k \Delta I_k<\infty$, then the conformality error in Theorem~\ref{thm:tick-fidelity-RP} vanishes and $\sup_k \mathrm{RP}(\rho_{k}^{\rm AdS})\to 0$.
\end{corollary}

\paragraph{Remarks.}\label{par:uom-main-244}
(i) The spatial micro envelope $\Psi_{x,\sigma_{x,k}}$ only enters through $q_k$ and $G_k$; zonal Gaussian/LoG choices commute with $\Delta_{S^3}$ and preserve the bounds. (ii) The temporal band–gap $\Omega_{\rm gap}(\Delta)$ is the smallest nonzero temporal frequency supported by the accepted radiative sector at band–limit $\Delta$ (for the massless sector on $S^3$, $\Omega_{\rm gap}\ge H\sqrt{3}$ if $\ell_{\min}=1$). (iii) The conformality deficit monitor $\mathcal D_{\rm conf}(\varepsilon_k,\Lambda_k)$ from Remark~\ref{rem:conf-monitor} is precisely the quantity minimized by the receiver–side “conformality” guard in the P3 program.

\begin{theorem}[UV conformal common sector on the receiver boundary]
\label{thm:uom-uv-conformal-common-sector}
Assume the static AdS$_5$ UV boundary setup of
Paragraph~\ref{par:uom-main-111}, the Minkowski UV chart and
inductive-limit vacuum sector of
Subsection~\ref{subsec:uom-energy-bounds}, the UV dilatation package of
Theorem~\ref{thm:uom-uv-dilatation-package}, the approximate
intertwiner bound of Lemma~\ref{lem:approx-intertwiner-new}, and the
anti-Zeno exactness statement of Corollary~\ref{cor:anti-zeno}.  Let
\[
\mathcal H_{\mathrm{conf}}
\;:=\;
\overline{\pi(\mathfrak A_{\mathrm{UV,loc}})\Omega}
\subset \mathcal H,
\]
where \(\mathfrak A_{\mathrm{UV,loc}}\) is the UV-local algebra of
Theorem~\ref{thm:uom-uv-dilatation-package}.  Then the UOM UV boundary
sector admits a canonical conformal implementation on
\(\mathcal H_{\mathrm{conf}}\):
\begin{enumerate}[label=\textnormal{(\roman*)},leftmargin=2.2em,itemsep=2pt]
\item there exists a strongly continuous unitary representation
  \(U_{\mathrm{conf}}(g)\) of the connected conformal group
  \(\mathrm{Conf}_0(\mathbb R^{1,3})\) on \(\mathcal H_{\mathrm{conf}}\)
  fixing the vacuum;
\item \(U_{\mathrm{conf}}(g)\) extends the Poincar\'e action already
  present in the inductive-limit vacuum representation and acts
  geometrically on the UV-local net:
  \[
  U_{\mathrm{conf}}(g)\,\pi(\mathfrak A_{\mathrm{UV,loc}}(O))\,U_{\mathrm{conf}}(g)^{-1}
  \;=\;
  \pi(\mathfrak A_{\mathrm{UV,loc}}(gO))
  \]
  for every double cone \(O\) in the Minkowski UV chart and every
  chart-preserving \(g\in\mathrm{Conf}_0(\mathbb R^{1,3})\);
\item the one-parameter dilatation subgroup of \(U_{\mathrm{conf}}\) is
  generated by the same selfadjoint operator \(D_{\mathrm{UV}}\) as in
  Theorem~\ref{thm:uom-uv-dilatation-package};
\item the one-parameter special conformal subgroups admit selfadjoint
  generators \(K_\mu\) on the common Stone core of the group
  representation, and together with the translation and Lorentz
  generators \(P_\mu,M_{\mu\nu}\) and \(D_{\mathrm{UV}}\) they satisfy
  the standard \(\mathfrak{so}(4,2)\) commutation relations on that
  core.
\end{enumerate}
\end{theorem}

\begin{proof}
Let \(G_{\mathrm{chart}}\subset\mathrm{Conf}(S^3)\) denote the
connected subgroup preserving the image of the chosen Minkowski chart
\(\kappa\).  Conjugating by \(\kappa\) identifies
\(G_{\mathrm{chart}}\) with \(\mathrm{Conf}_0(\mathbb R^{1,3})\).  By
the standing boundary symmetry discussion at the start of the
manuscript, \(G_{\mathrm{chart}}\) already acts unitarily on the
receiver boundary algebra through \(U_\partial(g)\).

At finite cutoff, the receiver bridge does not yet preserve the UV
conformal sector exactly; the error is measured by the approximate
intertwiner bound of Lemma~\ref{lem:approx-intertwiner-new} and, at the
level of accepted boundary states, by the conformal-fidelity estimate of
Theorem~\ref{thm:tick-fidelity-RP}.  Corollary~\ref{cor:anti-zeno}
states that along the anti-Zeno accepted UV limit these conformality
errors vanish.  Therefore the accepted UV boundary sector is stable
under the exact \(U_\partial(g)\)-action in the inductive limit.  Since
\(\mathfrak A_{\mathrm{UV,loc}}\) is generated by accepted sufficiently
local UV observables and \(\Omega\) is the UV fixed-point vacuum,
\(\mathcal H_{\mathrm{conf}}
=\overline{\pi(\mathfrak A_{\mathrm{UV,loc}})\Omega}\) is invariant
under the restricted action of \(U_\partial(g)\).  Define
\[
U_{\mathrm{conf}}(g)
\;:=\;
U_\partial(g)\big|_{\mathcal H_{\mathrm{conf}}}.
\]
This gives item~\textnormal{(i)} and vacuum invariance.

For item~\textnormal{(ii)}, the same chartwise geometric action that
defines \(U_\partial(g)\) sends a UV-local double cone \(O\) to \(gO\),
and the anti-Zeno exactness just noted removes the finite-cutoff
conformality defect.  Hence \(U_{\mathrm{conf}}(g)\) implements the
geometric conformal action on the UV-local net.  On the Poincar\'e
subgroup this action agrees with the one already reconstructed in the
inductive-limit vacuum representation, because both are induced from the
same Minkowski chart on the same UV-local algebra.

For item~\textnormal{(iii)}, the dilatation subgroup of
\(\mathrm{Conf}_0(\mathbb R^{1,3})\) acts on compactly smeared scalar UV
observables by the same constant Weyl/radial rescaling used in
Theorem~\ref{thm:uom-uv-dilatation-package}.  Both
\(U_{\mathrm{conf}}(e^{\tau D})\) and \(e^{i\tau D_{\mathrm{UV}}}\)
therefore implement the same automorphism family on the dense
UV-local core.  By uniqueness of the Stone generator of a strongly
continuous one-parameter unitary group, the generator is the same
\(D_{\mathrm{UV}}\).

Finally, item~\textnormal{(iv)} follows from Stone's theorem applied to
the special conformal one-parameter subgroups of
\(U_{\mathrm{conf}}(g)\), which yields selfadjoint generators \(K_\mu\)
on the common Stone core of the group representation.  Since
\(U_{\mathrm{conf}}\) is a unitary representation of the connected
conformal group, its infinitesimal generators satisfy the Lie algebra
relations of \(\mathfrak{so}(4,2)\), including the standard commutators
among \(P_\mu\), \(M_{\mu\nu}\), \(D_{\mathrm{UV}}\), and \(K_\mu\).
\end{proof}

\begin{corollary}[UOM core exports the conformal common-sector interface]
\label{cor:uom-core-implies-2T-conf}
The hypotheses and conclusions of
Theorems~\ref{thm:uom-uv-dilatation-package} and
\ref{thm:uom-uv-conformal-common-sector} furnish the abstract
same-representation conformal common-sector interface
\(\mathbf{2T}_{\mathrm{conf}}\): the UV-local net
\(\mathfrak A_{\mathrm{UV,loc}}(O)\), the conformal representation
\(U_{\mathrm{conf}}(g)\), the common Stone core, and the shared
dilatation generator \(D_{\mathrm{UV}}\) provide exactly the conformal
supplement consumed on the Main side.
\end{corollary}

\begin{proof}
Theorem~\ref{thm:uom-uv-dilatation-package} supplies the common
same-representation dilatation package on the UV-local algebra, while
Theorem~\ref{thm:uom-uv-conformal-common-sector} adds the geometric
conformal action, the special conformal generators, and the common
Stone core on that same Hilbert sector.  Reading those data together
gives the exported conformal common-sector interface used by Main.
\end{proof}

\begin{remark}[Role for the Main conformal supplement]
\label{rem:uom-uv-conformal-common-sector-role}
Theorem~\ref{thm:uom-uv-conformal-common-sector} is the UOM-side source
of the narrow Main conformal supplement \(\mathbf{2T}_{\mathrm{conf}}\):
it does not assert the later post-Main \(\mathbb Z_2\) rigidity or full
observable-transfer story, only the common UV conformal sector needed to
place \(P_\mu,M_{\mu\nu},D,K_\mu\) on the same theorem object as the
UV AQFT/HLS representation.
\end{remark}

\printbibliography
\end{document}
